% ===== begin preamble.tex =====
% MAT670 -- Packings, Lattices and Configurations -- shared preamble
\documentclass[11pt,openany]{book}
\usepackage[margin=1.1in,marginparwidth=0.9in]{geometry}
\usepackage{amsmath,amssymb,amsthm,mathtools}
\usepackage[T1]{fontenc}
\usepackage{enumitem}
\usepackage{xcolor}
\usepackage{tikz}
\usetikzlibrary{calc,arrows.meta,patterns,decorations.pathmorphing}
\usepackage{graphicx}
\usepackage{booktabs}
\usepackage[colorlinks=true,linkcolor=blue!50!black,citecolor=green!40!black,urlcolor=blue!50!black]{hyperref}

% ---------- theorem environments ----------
\theoremstyle{plain}
\newtheorem{theorem}{Theorem}[chapter]
\newtheorem{lemma}[theorem]{Lemma}
\newtheorem{proposition}[theorem]{Proposition}
\newtheorem{corollary}[theorem]{Corollary}
\newtheorem{conjecture}[theorem]{Conjecture}
\newtheorem{claim}[theorem]{Claim}
\theoremstyle{definition}
\newtheorem{definition}[theorem]{Definition}
\newtheorem{example}[theorem]{Example}
\newtheorem{exercise}[theorem]{Exercise}
\newtheorem{question}[theorem]{Question}
\newtheorem{problem}[theorem]{Problem}
\theoremstyle{remark}
\newtheorem{remark}[theorem]{Remark}

% ---------- editorial markup (first-pass) ----------
% \added{...}   material NOT in the handwritten notes (gap filled by editor)
% \fixed{...}   an error in the notes that was corrected; argument says what changed
% \unclear{...} a reading of the handwriting that is uncertain, or a gap not resolved
% \orig{...}    short note on what the original page said / showed (e.g. "sketch omitted")
\definecolor{addcol}{RGB}{0,100,160}
\definecolor{flagcol}{RGB}{190,60,0}
\newcommand{\added}[1]{{\color{addcol}#1}}
\newcommand{\fixed}[1]{\marginpar{\raggedright\scriptsize\color{flagcol}\textbf{Fix:} #1}}
\newcommand{\unclear}[1]{\marginpar{\raggedright\scriptsize\color{flagcol}\textbf{?} #1}}
\newcommand{\orig}[1]{\marginpar{\raggedright\scriptsize\color{gray}#1}}
\newenvironment{addedblock}{\par\color{addcol}\noindent\textsf{\small[Added]}\ }{\par}

% ---------- math macros ----------
\newcommand{\R}{\mathbb{R}}
\newcommand{\Z}{\mathbb{Z}}
\newcommand{\Q}{\mathbb{Q}}
\newcommand{\C}{\mathbb{C}}
\newcommand{\N}{\mathbb{N}}
\newcommand{\T}{\mathbb{T}}
\newcommand{\Sph}{\mathbb{S}}          % sphere S^{n-1}
\newcommand{\Lat}{\Lambda}             % a lattice
\newcommand{\vol}{\operatorname{vol}}
\newcommand{\covol}{\operatorname{covol}}
\newcommand{\dist}{\operatorname{dist}}
\newcommand{\diam}{\operatorname{diam}}
\newcommand{\conv}{\operatorname{conv}}
\newcommand{\interior}{\operatorname{int}}
\newcommand{\GL}{\mathrm{GL}}
\newcommand{\SL}{\mathrm{SL}}
\newcommand{\SO}{\mathrm{SO}}
\newcommand{\Orth}{\mathrm{O}}
\newcommand{\Conf}{\operatorname{Conf}}
\newcommand{\Hess}{\operatorname{Hess}}
\newcommand{\ind}{\operatorname{ind}}
\newcommand{\abs}[1]{\left\lvert #1 \right\rvert}
\newcommand{\norm}[1]{\left\lVert #1 \right\rVert}
\newcommand{\ip}[2]{\langle #1, #2\rangle}
\newcommand{\dens}{\delta}             % packing density
\newcommand{\cent}{\theta}             % center density
\newcommand{\kiss}{\tau}               % kissing number
\newcommand{\Ball}{B}                  % ball
% ===== end preamble.tex =====

\title{\textbf{Packings, Lattices and Configurations}\\[0.6em]
\large Lecture notes for MAT670, TU Graz, Summer Term 2016\\[0.3em]
\normalsize First typeset draft from the handwritten notes}
\author{Woden Kusner}
\date{Draft of \today}

\begin{document}
\frontmatter
\maketitle

\chapter*{About this draft}
\addcontentsline{toc}{chapter}{About this draft}

These notes are a first typeset pass at the handwritten notes for the course
\emph{Packings, Lattices and Configurations} (MAT670), taught in eleven two-hour lectures at TU Graz
in the summer term of 2016. The course description read:
\begin{quote}
This course will introduce a variety of different tools and models to explore some of the theory of
geometric configuration spaces that underlie physical systems. \dots\ By the end of the course, you
should have some familiarity with the geometry and topology of configuration spaces, Morse theory, the
geometry of lattices, maximal density arrangements, random packings and statistical mechanics.
\end{quote}

The structure, order, examples and choice of proofs follow the handwritten notes. Telegraphic notes have
been written out as prose, sketches have been redrawn, and gaps have been filled. The colour code is:

\begin{itemize}
\item Black text is a cleaned-up rendering of what is on the page.
\item \added{Blue text} was \emph{not} in the handwritten notes. It covers completed definitions and
  proofs, missing hypotheses, references, and remarks on results that appeared after 2016.
  Longer insertions are marked \added{\textsf{\small[Added]}}.
\item Orange margin notes \textbf{Fix:} mark places where the notes contained an error that has been
  corrected, saying what changed. \textbf{?} marks a doubtful reading of the handwriting,
  or a step that is still open.
\item Grey margin notes describe the original page (crossed-out text, omitted doodles, and so on).
\end{itemize}

\paragraph{The interludes.} Several chapters are entirely new and are set in black for readability, but
all of them are added material:
\begin{itemize}
\item[A] From Dynkin diagrams to Lie algebras and Lie groups (after Lecture 6);
\item[B] Lattices, quadratic forms and the geometry of numbers (after A);
\item[C] The torus, Diophantine approximation and aperiodic order (after B);
\item[D] Min-max, configuration spaces and braids (after Lecture 11);
\item[E] Random packings, entropy and information (after D);
\item[F] Computation and certified proof (after E);
\item[G] Packings in the laboratory (after F);
\item[H] Dynkin diagrams everywhere (at the end).
\end{itemize}

\paragraph{Extensions.} Sections titled ``Extension'' inside the lectures are also new material set in
black: Rogers' simplex bound (end of Lecture~2) and why root lattices pack well (end of Lecture~6).

\paragraph{Sources.} The scans used are \texttt{MAT670\_1}, \texttt{\_2}, \texttt{\_3},
\texttt{\_46}, \texttt{\_78}, \texttt{\_9X} and \texttt{\_XI} on the course web page.
\texttt{MAT670\_emsemble} is a bundled copy of Lectures 1--3.
\texttt{MAT670\_classification} is an earlier scan of Lectures 4--5, the first 31 pages of
\texttt{MAT670\_46} minus one recap page. Both were transcribed independently and used only as a
cross-check.

\paragraph{Known open items in this draft.}
\begin{itemize}
\item Lecture 3: the ``improved'' radial Blichfeldt gauge is stated in the notes without proof. No proof
  has been supplied here; a numerical search in low dimensions found no counterexample.
\item Lecture 10: the deformation theorem for $0<\theta<\pi/N$ had no proof in the notes. The proof given
  in blue is the editor's and should be checked. Pages 24--26 of the scan are mostly headings.
  The notes themselves mark this lecture ``needs to be rewritten, out of order / missing pieces''.
\item Lecture 11: the count of locally jammed configurations in the open-ended channel may be off at the
  ends, giving $\binom{\ell-k-1}{k}$ rather than $\binom{\ell-k+1}{k}$. It is flagged but not changed.
\end{itemize}

\tableofcontents

% ===== begin chapters/threads.tex =====
% Front matter: an associative guide to the notes. New material.
\chapter*{Threads through the notes}
\addcontentsline{toc}{chapter}{Threads through the notes}
\label{THR:ch}

These notes are not meant to be read in a straight line. The same few objects keep turning up in
different clothes, and following them is the most fun way through. Below are some of the threads,
each with its stops. Lectures are numbered $1$--$11$ and the new interludes are lettered A--H.

\begin{center}
\begin{tikzpicture}[x=0.56cm,y=0.52cm]
  \node[font=\small] at (0,1) {1};
  \node[font=\small] at (1,1) {2};
  \node[font=\small] at (2,1) {3};
  \node[font=\small] at (3,1) {4};
  \node[font=\small] at (4,1) {5};
  \node[font=\small] at (5,1) {6};
  \node[font=\small\bfseries,addcol] at (6,1) {A};
  \node[font=\small\bfseries,addcol] at (7,1) {B};
  \node[font=\small\bfseries,addcol] at (8,1) {C};
  \node[font=\small] at (9,1) {7};
  \node[font=\small] at (10,1) {8};
  \node[font=\small] at (11,1) {9};
  \node[font=\small] at (12,1) {10};
  \node[font=\small] at (13,1) {11};
  \node[font=\small\bfseries,addcol] at (14,1) {D};
  \node[font=\small\bfseries,addcol] at (15,1) {E};
  \node[font=\small\bfseries,addcol] at (16,1) {F};
  \node[font=\small\bfseries,addcol] at (17,1) {G};
  \node[font=\small\bfseries,addcol] at (18,1) {H};
  \node[anchor=east,font=\small] at (-0.7,0) {golden ratio};
  \draw[gray!30] (0,0)--(18,0);
  \draw[addcol,thick] (6,0)--(18,0);
  \fill[addcol] (6,0) circle (0.17);
  \fill[addcol] (8,0) circle (0.17);
  \fill[addcol] (13,0) circle (0.17);
  \fill[addcol] (15,0) circle (0.17);
  \fill[addcol] (18,0) circle (0.17);
  \node[anchor=east,font=\small] at (-0.7,-1) {the root system $A_{n-1}$};
  \draw[gray!30] (0,-1)--(18,-1);
  \draw[addcol,thick] (5,-1)--(18,-1);
  \fill[addcol] (5,-1) circle (0.17);
  \fill[addcol] (6,-1) circle (0.17);
  \fill[addcol] (9,-1) circle (0.17);
  \fill[addcol] (10,-1) circle (0.17);
  \fill[addcol] (14,-1) circle (0.17);
  \fill[addcol] (18,-1) circle (0.17);
  \node[anchor=east,font=\small] at (-0.7,-2) {Blichfeldt and Rogers};
  \draw[gray!30] (0,-2)--(18,-2);
  \draw[addcol,thick] (1,-2)--(7,-2);
  \fill[addcol] (1,-2) circle (0.17);
  \fill[addcol] (2,-2) circle (0.17);
  \fill[addcol] (5,-2) circle (0.17);
  \fill[addcol] (7,-2) circle (0.17);
  \node[anchor=east,font=\small] at (-0.7,-3) {$E_8$};
  \draw[gray!30] (0,-3)--(18,-3);
  \draw[addcol,thick] (5,-3)--(18,-3);
  \fill[addcol] (5,-3) circle (0.17);
  \fill[addcol] (6,-3) circle (0.17);
  \fill[addcol] (7,-3) circle (0.17);
  \fill[addcol] (8,-3) circle (0.17);
  \fill[addcol] (15,-3) circle (0.17);
  \fill[addcol] (16,-3) circle (0.17);
  \fill[addcol] (18,-3) circle (0.17);
  \node[anchor=east,font=\small] at (-0.7,-4) {the space of lattices};
  \draw[gray!30] (0,-4)--(18,-4);
  \draw[addcol,thick] (0,-4)--(8,-4);
  \fill[addcol] (0,-4) circle (0.17);
  \fill[addcol] (2,-4) circle (0.17);
  \fill[addcol] (3,-4) circle (0.17);
  \fill[addcol] (7,-4) circle (0.17);
  \fill[addcol] (8,-4) circle (0.17);
  \node[anchor=east,font=\small] at (-0.7,-5) {contacts and critical points};
  \draw[gray!30] (0,-5)--(18,-5);
  \draw[addcol,thick] (12,-5)--(17,-5);
  \fill[addcol] (12,-5) circle (0.17);
  \fill[addcol] (13,-5) circle (0.17);
  \fill[addcol] (14,-5) circle (0.17);
  \fill[addcol] (15,-5) circle (0.17);
  \fill[addcol] (16,-5) circle (0.17);
  \fill[addcol] (17,-5) circle (0.17);
  \node[anchor=east,font=\small] at (-0.7,-6) {entropy};
  \draw[gray!30] (0,-6)--(18,-6);
  \draw[addcol,thick] (2,-6)--(17,-6);
  \fill[addcol] (2,-6) circle (0.17);
  \fill[addcol] (9,-6) circle (0.17);
  \fill[addcol] (13,-6) circle (0.17);
  \fill[addcol] (15,-6) circle (0.17);
  \fill[addcol] (17,-6) circle (0.17);
  \node[anchor=east,font=\small] at (-0.7,-7) {circles and tori};
  \draw[gray!30] (0,-7)--(18,-7);
  \draw[addcol,thick] (6,-7)--(9,-7);
  \fill[addcol] (6,-7) circle (0.17);
  \fill[addcol] (7,-7) circle (0.17);
  \fill[addcol] (8,-7) circle (0.17);
  \fill[addcol] (9,-7) circle (0.17);
  \node[anchor=east,font=\small] at (-0.7,-8) {tetrahedra and icosahedra};
  \draw[gray!30] (0,-8)--(18,-8);
  \draw[addcol,thick] (13,-8)--(18,-8);
  \fill[addcol] (13,-8) circle (0.17);
  \fill[addcol] (15,-8) circle (0.17);
  \fill[addcol] (17,-8) circle (0.17);
  \fill[addcol] (18,-8) circle (0.17);
  \node[anchor=east,font=\small] at (-0.7,-9) {certified computation};
  \draw[gray!30] (0,-9)--(18,-9);
  \draw[addcol,thick] (0,-9)--(16,-9);
  \fill[addcol] (0,-9) circle (0.17);
  \fill[addcol] (13,-9) circle (0.17);
  \fill[addcol] (16,-9) circle (0.17);
  \foreach \x in {5.5,8.5,13.5} \draw[gray,dashed] (\x,1.5)--(\x,-9.5);
\end{tikzpicture}
\end{center}
Chapters are listed in book order (lectures in black, interludes in blue); a dot marks a stop on the thread.


\paragraph{The golden ratio.} Lecture~11 counts the jammed states of a narrow channel with binary words
having no two adjacent $1$s, and Fibonacci numbers appear. The most likely fraction of $1$s is
$(5-\sqrt5)/10$, and the entropy per site is $\log\varphi$. Interlude~E recognizes this as the Shannon
capacity of the ``golden mean'' channel. Interlude~C meets the Fibonacci word again: it is a Sturmian
word of slope $\varphi^{-2}$, the cutting sequence of an irrational line through the square lattice, its gaps obey the three-gap
theorem, and the golden ratio is the worst approximable number (Hurwitz). Folding the $E_8$ diagram
with $\tau=\varphi$ gives the non-crystallographic group $H_4$ and icosahedral quasicrystals
(Interlude~C). Interlude~A sends the binary icosahedral group to the extended $E_8$ diagram (McKay).
The icosahedron itself is not made of regular tetrahedra (Interlude~E).

\paragraph{The root system $A_{n-1}$.} It is built in Lecture~6 as the vectors $e_i-e_j$ in the
hyperplane $\sum x_i=0$. In Lecture~8 the same hyperplanes $x_i=x_j$ cut $\Conf(N,\R)$ into Weyl
chambers. In Interlude~A the roots are the eigenvalues of conjugation on the matrix units
$E_{ij}$; the maximal torus of $\mathrm{SU}(n)$ is $\R^{n-1}$ modulo the $A_{n-1}$ lattice; and its Weyl
alcove is the triangle of gaps between points on a circle from Lecture~7. Complexify the hyperplanes
and divide by $\Sigma_N$, and you get the braid group (Interlude~D).

\paragraph{Blichfeldt.} His 1929 gauge trick bounds the density in Lectures~2 and~3. Rogers' simplex bound (end of Lecture~2) replaces his
spreading of mass by a cutting of Voronoi cells, and is sharp in the plane, where it is Thue's theorem. His 1914 lemma
proves Minkowski's convex body theorem (Interlude~B). The Cohn--Elkies bound is a ``two-sided''
Blichfeldt gauge, with positivity demanded on both sides of the Fourier transform, and with Viazovska's
magic function it settles dimension $8$ (Interlude~B). The gauges for long bodies in Lecture~3 (capsules, cylinders) are the
subject of G.~Fejes T\'oth and W.~Kuperberg, \emph{Blichfeldt's density bound revisited} (Math.\ Ann.\ 1993),
which leads on to the cylinder bounds of A.~Bezdek and W.~Kuperberg (1990) and W.~Kusner (2014).

\paragraph{$E_8$.} It is constructed in Lecture~6 as $D_8\cup(D_8+h)$ and from the Hamming code. In
Interlude~A it is $248=8+240$: Cartan subalgebra plus root spaces, and its density $\pi^4/384$ can be
read off from $\det C=1$. In Interlude~B its theta series is the Eisenstein series $E_4$, so it has
$240\,\sigma_3(m)$ vectors of norm $2m$, and Viazovska's theorem makes it the densest packing in
dimension~$8$. In Interlude~C it folds to $H_4$, and in Interlude~E the Golay code builds its big
sibling, the Leech lattice.

\paragraph{Why root lattices?} They are within a few percent of Rogers' bound in dimensions $2$--$8$
(end of Lecture~6). Their density can be read off from the determinant of the Cartan matrix
(Interlude~A). The Weyl group's irreducibility makes them eutactic and the abundance of roots makes them
perfect, so Voronoi's theorem makes each one a local optimum. Their $O(n^2)$ roots are too few in high
dimension, where Interlude~B's random lattices take over.

\paragraph{The space of lattices.} It appears in Lecture~4 with Mahler's criterion, and in
Interlude~B with Hermite's constant, reduction of quadratic forms and Siegel's mean value theorem. From
the last of these comes the Minkowski--Hlawka lower bound $\dens\ge 2^{-d}\cdot2\zeta(d)$, which was the
lower bound quoted in Lectures~1 and~3. In Interlude~C dynamics on the space takes over: bounded orbits
correspond to badly approximable numbers (Dani), and Margulis's proof of the Oppenheim conjecture.

\paragraph{Contacts, forces and critical points.} The contact graphs and force balance of
Lectures~10--11 make the max-min distance function a ``min-type'' Morse function. Interlude~D turns this
into min-max: the Tammes problem as the top critical value, Betti numbers of configuration spaces as
lower bounds on the number of critical configurations, and kissing numbers. Interlude~E reads the same
force-balance condition as jamming, with Maxwell counting and isostaticity. The slack in the
icosahedral kissing configuration (Interlude~E) is what lets the jitterbug of Lecture~11 move.

\paragraph{Entropy.} Lecture~7 derives the Boltzmann distribution and the Gaussian fluctuations of a
gas in a box. Lecture~11 introduces Edwards' entropy of jammed states and solves hard rods. In
Interlude~E the partition function of hard spheres becomes the volume of a configuration space; the
parking problem has an exact answer; and maximum entropy, Shannon's channel capacity and sphere
packing in high dimensions turn out to be one circle of ideas, with Kabatiansky--Levenshtein
(Lecture~3) as a coding theorem.

\paragraph{Circles and tori.} Points on a circle (Lecture~7) have a triangle of gaps. Interlude~A finds
that triangle as the space of conjugacy classes of $\mathrm{SU}(3)$. Interlude~C cuts the torus along an
irrational line (``shredding the torus''), proves equidistribution, and finds at most three gaps.
Interlude~B counts lattice points in a circle and asks how wrong $\pi r^2$ can be (the Gauss circle
problem).

\paragraph{Tetrahedra and icosahedra.} Twelve balls around one fit at the vertices of an icosahedron
with room to spare. Lecture~11's jitterbug uses that room, and Interlude~E proves it with rotating
tropical rings. Interlude~H sees the same fact in a Coxeter diagram: $\{3,3,5\}$ is spherical, not flat,
so five regular tetrahedra close up around an edge only on $\Sph^3$, in the $600$-cell. In the FCC
lattice, octahedra fill exactly the gap the tetrahedra leave. In the laboratory, tetrahedral dice and
colloidal clusters show what happens in between (Interlude~G).

\paragraph{Certified computation.} Thue's triangle lemma (Lecture~1) is proved again in Interlude~F by
a branch and bound in exact rational arithmetic, using a Lagrange multiplier that is the force balance of
Lectures~10--11. The same chapter treats the case analyses that the Tammes and Kepler problems
need, and the certificates (intervals, sums of squares, formal proof)
that make them trustworthy.
% ===== end chapters/threads.tex =====

\mainmatter
\part{Packing bounds}
% ===== begin chapters/L01.tex =====
% ============================================================
% MAT670 -- Lecture 1 (source: MAT670_1, pp. 1--17)
% ============================================================
\chapter{Lecture 1: Packing problems}
\label{L1:chap}

\section{Course format}
\label{L1:sec:admin}
\orig{1.0 ``Admin'': contact, room, breaks omitted}
Evaluation is by short quizzes and a project: a talk of at least 20 minutes
together with a written summary. Updates are posted by email, through TU Online,
and on the course page \texttt{wkusner.github.io/MAT670/}.

\section{Mathematics of packings, lattices and configurations}
\label{L1:sec:overview}

The subject of this course sits at the meeting point of several areas:
\begin{itemize}[nosep]
  \item classical convex and discrete geometry, and the geometry of numbers;
  \item applied topology and Morse theory;
  \item statistical mechanics, combinatorics and information theory.
\end{itemize}
\orig{p.\,1: a diagram linking these areas; ``Rigidity'' crossed out there}
We will take various ideas from each of these, all of them dealing with
\emph{configurations} of points or bodies. Related topics include rigidity theory
and configuration spaces of linkages, information theory (codes), and the geometry
of numbers. Packing problems are my main motivation.

Packing problems are typically \emph{easy to state but hard to solve}. The same
is true of the geometric inequalities and the configuration spaces we will meet
along the way.

\begin{remark}
\added{Some landmarks, for orientation: the densest packing of congruent balls
in $\R^3$ (the Kepler conjecture, density $\pi/\sqrt{18}$) was proved by Hales
(announced 1998, published 2005; formal verification by the Flyspeck project 2014/2017).
In dimension $8$ the $E_8$ lattice packing was proved optimal by Viazovska (2016),
and in dimension $24$ the Leech lattice packing by
Cohn--Kumar--Miller--Radchenko--Viazovska (2016/2017). These dimension $8$ and $24$
results appeared during the course itself. In every other dimension
$d\ge 4$ the problem is still open.}
\end{remark}

\section{Packing problems}
\label{L1:sec:packing}

\subsection{Packing congruent bodies in a box}

We start in two dimensions, with the problem of packing circles (disks) in a box.
Let $B\subset\R^2$ be a box (a square, say) and $K$ a disk. We want to place $N$
congruent copies of $K$ in $B$ so as to maximize the \emph{density}, or volume fraction,
\[
  \frac{\vol\bigl(\bigcup_{i=1}^N K_i\bigr)}{\vol(B)} .
\]

\begin{center}
\begin{tikzpicture}[scale=0.8]
  \draw[thick] (0,0) rectangle (4,4);
  \foreach \x/\y in {1/1,3/1,1/3,3/3} \draw[fill=blue!10] (\x,\y) circle (1);
  \node[right] at (4,2) {$B$};
\end{tikzpicture}
\\[2pt]
{\small \added{Four congruent disks packed in a square $B$ (this $2\times2$ arrangement is optimal for $N=4$).}}
\end{center}

\begin{definition}[Packing in a box]
\label{L1:def:packing-box}
A collection $K_1,\dots,K_N$ of subsets of $\R^n$ is a \emph{packing of $K$ in $B$} if
\begin{enumerate}[nosep]
  \item the $K_i$ are congruent to $K$: $K_i=\varphi_i K$ with
        $\varphi_i\in\R^n\rtimes\SO(n)$ a rigid motion\added{;}
  \item $K_i\subset B$ for all $i$\added{;}
  \item the interiors are pairwise disjoint: $\interior K_i\cap\interior K_j=\emptyset$
        for $i\neq j$.
\end{enumerate}
\end{definition}
\orig{p.\,3 writes $\R^2\rtimes\SO(2)$}

\begin{example}
\label{L1:ex:radius}
Two equivalent formulations of the disk-in-a-box problem:
\begin{enumerate}[nosep]
  \item Find the largest scale $r$ such that $\bigcup_{i=1}^N \varphi_i(rK)\subset B$ is
        a packing of $B$\added{, where $K$ is the unit disk}. (Here $rK$ is in Minkowski notation:
        $rK=\{rx : x\in K\}$.)
  \item Find a configuration $\{x_1,\dots,x_N\}$ of $N$ points in $B$ such that
        \[
          d(x_i,x_j)\ \ge\ 2r \quad\text{for } i\ne j,
          \qquad\text{and}\qquad d(x_i,\partial B)\ \ge\ r \quad\text{for all } i,
        \]
        with $r$ as large as possible.
\end{enumerate}
\fixed{notes: $d(x_i,x_j)>r$; disks of radius $r$ need centres $2r$ apart}
\added{The points $x_i$ are the centres of the disks $\varphi_i(rK)$.}
\end{example}

These are hard problems, and there is an infinite variety of them: we can change
$N$, $K$ and $B$, the dimension, add other constraints, and so on.

One case of particular interest is $B=\R^2$: the circle packing problem in the
plane. Other natural choices are $B=\R^d$, or $B=\Sph^2$.
\added{(For $B=\R^d$ the number of bodies is infinite and ``density'' must be
defined by a limit; this is done next.)}

\subsection{Packings of the plane}
\label{L1:subsec:plane}
\orig{numbered ``1.2'' a second time in the notes}

Now take $B=\R^2$. We have an infinite collection of congruent disks $K_i$, now
normalized to have unit radius, with $\interior K_i\cap\interior K_j=\emptyset$ for $i\ne j$.
Write $P=\bigcup_i K_i$.

\begin{definition}[Upper and lower density]
\label{L1:def:density}
The \emph{upper density} of the packing $P$ is
\[
  \dens^+(P)\;=\;\limsup_{R\to\infty}\frac{\vol(P\cap R\,\mathbb{D}^2)}{\vol(R\,\mathbb{D}^2)},
\]
where $\mathbb{D}^2$ is the closed unit disk \added{centred at the origin}.
The \emph{lower density} $\dens^-(P)$ is defined in the same way with $\liminf$.
In general dimension,
\[
  \dens^+(P)\;=\;\limsup_{R\to\infty}\frac{\vol(P\cap R\,\mathbb{D}^n)}{\vol(R\,\mathbb{D}^n)}
\]
\added{with $\mathbb{D}^n$ the unit ball in $\R^n$}.
\end{definition}

\begin{remark}
These quantities need not be nice in general \added{(e.g.\ $\dens^+(P)\neq\dens^-(P)$
is possible, and then ``the'' density of $P$ is not defined)}. However:
\begin{itemize}[nosep]
  \item they do \emph{not} depend on the choice of origin \added{(moving the centre
        changes the numerator only in a boundary layer of volume $O(R^{n-1})$)};
  \item they \emph{do} depend, in general, on the shape of the averaging region
        $\mathbb{D}^n$ \added{(for periodic packings, though, any reasonable shape gives the same value)}.
\end{itemize}
\end{remark}

Does a packing of maximal density exist? We might look at this in more detail later.
\added{(It does: the supremum of $\dens^+$ over all packings is attained, by a compactness
argument; see e.g.\ Groemer's theorem in Conway--Sloane, \emph{Sphere Packings, Lattices
and Groups}, Ch.~1.)}
\unclear{existence deferred in the notes}

\subsection{Bounds on the density}

How can we compute bounds on the density? Clearly $\dens^+(P)\le 1$, which is a bad bound.

We can also construct lower bounds using highly structured packings. For example,
the unit balls centred at the points of $(2\Z)^d$ form a packing. Since this structure
is periodic, we can work with the fundamental domain.

\begin{exercise}[Easy]
\label{L1:ex:periodic-density}
Why can we compute the density on a fundamental domain? \added{Show that for a periodic
packing, with period lattice $\Lat$ and $m$ balls of radius $1$ centred in a fundamental
domain $F$ of $\Lat$, the density exists (that is, $\dens^+=\dens^-$) and equals
$m\,\vol(\Ball^d)/\vol(F)$.}
\end{exercise}

Recall the volume of the ball of radius $r$ in $\R^d$:
\[
  V_d(r)\;=\;\vol\bigl(\Ball^d(r)\bigr)\;=\;\frac{\pi^{d/2}}{\Gamma\!\left(\frac d2+1\right)}\,r^d .
\]
This is easier to remember in even dimension: $V_{2k}(r)=\dfrac{\pi^k}{k!}\,r^{2k}$.
For $(2\Z)^d$ the fundamental domain is the cube $[0,2)^d$ of volume $2^d$, containing one
ball, so the density is
\[
  \dens\bigl((2\Z)^d\bigr)\;=\;\frac{V_d(1)}{2^d}\;=\;\frac{\pi^{d/2}}{\Gamma\!\left(\frac d2+1\right)2^d}
  \;\overset{d=2k}{=}\;\frac{\pi^{k}}{k!\,2^{2k}} .
\]
\orig{notes write $(\tfrac d2)!$ for $\Gamma(\tfrac d2+1)$}

We could improve this by considering better periodic systems.

\begin{exercise}[Easy]
\label{L1:ex:periodic-approx}
Show that periodic packings can approximate a densest packing: \added{the supremum of
$\dens^+(P)$ over all packings $P$ equals the supremum over periodic packings.
(Hint: cut a large cube out of a near-optimal packing, remove the balls that meet its
boundary, and repeat periodically.)}
\end{exercise}

But which periodic systems should we use?

\begin{exercise}
\label{L1:ex:root-lattices}
Consider the lattices coming from the root systems $A_1\times\cdots\times A_1$
($d$ factors), \dots, $A_d$. Compute their packing densities and compare.
\unclear{p.\,6: ``Consider $A^1$, $A_1\times^d A_1$ \dots\ $A_d$? root systems''}
\added{(For reference: $A_1^d$ gives $\Z^d$, which is $(2\Z)^d$ after scaling, with
density $V_d(1)/2^d$. The lattice $A_d$ has minimal norm $\sqrt2$ and covolume
$\sqrt{d+1}$, hence density $V_d(1)\,2^{-d/2}/\sqrt{d+1}$. This gives
$\pi/\sqrt{12}$ for $d=2$ and $\pi/\sqrt{18}$ (FCC) for $d=3$.)}
\end{exercise}

\section{A better bound (non-constructive)}
\label{L1:sec:saturation}
\orig{\S1.3 in the notes}

\begin{definition}[Saturated packing]
\label{L1:def:saturated}
A packing $P=\bigcup_i K_i$ of congruent copies of $K$ in $\R^n$ is \emph{saturated} if
$P\cup\varphi K$ is not a packing for any $\varphi\in\R^n\rtimes\SO(n)$; that is,
no additional disk (ball) can be added.
\end{definition}

\begin{itemize}[nosep]
  \item Saturated packings certainly exist: \added{start from any packing and keep adding
        balls}. This is an ``infinite algorithm''. \added{It can be made rigorous with
        Zorn's lemma, or by adding balls greedily along an enumeration of a countable
        dense set of centres.}
  \item Saturation does not decrease $\dens^+(P)$ \added{since we only add balls}.
\end{itemize}
So for lower bounds on the best possible density we may assume that $P$ is saturated.

This is a surprisingly useful idea. Let $K$ be the unit ball, $K_i=\varphi_i K$ with centre $c_i$.

\begin{proposition}
\label{L1:prop:saturated-bound}
If $P$ is a saturated packing of unit balls in $\R^d$, then $\dens^+(P)\ge 2^{-d}$.
\end{proposition}

\begin{proof}
If $P$ is saturated, then for every $x\in\R^d$ there is an $i$ with $d(x,K_i)<1$.
\added{Otherwise $|x-c_i|\ge 2$ for all $i$, and the unit ball centred at $x$ could be added.}
Hence the doubled balls
\[
  2K_i\;:=\;\varphi_i\bigl(2\,\varphi_i^{-1}(K_i)\bigr)\;=\;\Ball(c_i,2)
\]
form a covering of $\R^d$. \added{Each $2K_i$ has $2^d$ times the volume of $K_i$, so the
fraction of $R\,\mathbb{D}^d$ covered by $\bigcup_i 2K_i$ is, up to a boundary term that
vanishes as $R\to\infty$, at most $2^d$ times the fraction covered by $P$. Since the
doubled balls cover everything,}
\[
  2^d\,\dens^+(P)\ \ge\ 1,\qquad\text{i.e.}\qquad \dens^+(P)\ \ge\ \frac1{2^d}. \qedhere
\]
\end{proof}
\orig{boxed correction: dilate about the centre, $(2\varphi_i^{-1}K)\varphi_i$}

\begin{remark}
\added{Compare with $(2\Z)^d$, of density $V_d(1)/2^d$. Since $V_d(1)<1$ exactly for
$d\ge13$, and $V_d(1)\to0$ superexponentially, the non-constructive bound $2^{-d}$ is
the better one in high dimensions.}
\end{remark}

We may go over some slight improvements in the future, but none seem satisfying.
\begin{remark}
\added{Known improvements of the $2^{-d}$ lower bound include the bound of
Minkowski and Hlawka
($2\zeta(d)\,2^{-d}$), Rogers and Ball ($c\,d\,2^{-d}$), Venkatesh 2013
($c\,d\log\log d\,2^{-d}$ along a sparse sequence of $d$), and, after the course,
Campos, Jenssen, Michelen and Sahasrabudhe 2023 ($c\,d\log d\,2^{-d}$) and
Klartag 2025 ($c\,d^2\,2^{-d}$). All of them are only polynomial improvements over $2^{-d}$.}
\end{remark}

The problem is that there is much more freedom in dimensions $d>3$!
\unclear{``$d>3$! even.'': perhaps ``even $d=4$''}
For example, take $(2\Z)^4$. The main diagonal of the cube $[0,2]^4$ has length
$2\sqrt4=4$, so the centre $(1,1,1,1)$ is at distance $2$ from every vertex.
\added{So a unit ball fits into every such hole. Adding them gives the
$D_4$ packing $(2\Z)^4\cup\bigl((2\Z)^4+(1,1,1,1)\bigr)$, which has twice the density,
$\pi^2/16$.}

\section{Upper bounds in dimension 2}
\label{L1:sec:upper2}
\orig{\S1.4}

The lower bound in $d=2$ is constructive: it comes from the hexagonal (triangular) lattice packing.

\begin{center}
\begin{tikzpicture}[scale=0.55]
  \clip (-0.6,-1.2) rectangle (10.6,6.4);
  \foreach \j in {-1,0,1,2,3,4} {
    \foreach \i in {-3,...,6} {
      \pgfmathsetmacro{\x}{2*\i+\j}
      \pgfmathsetmacro{\y}{sqrt(3)*\j}
      \draw[fill=blue!8] (\x,\y) circle (1);
      \fill (\x,\y) circle (1.5pt);
    }
  }
  \draw[very thick,red!70!black] (2,0) -- (4,0) -- (3,{sqrt(3)}) -- cycle;
  \draw[->,thick] (0,0) -- (2,0) node[midway,below] {$v_1$};
  \draw[->,thick] (0,0) -- (1,{sqrt(3)}) node[midway,left] {$v_2$};
\end{tikzpicture}
\\[2pt]
{\small Hexagonal packing of unit disks, basis $v_1=(2,0)$, $v_2=(1,\sqrt3)$ \added{(angle $\pi/3$)}, and an equilateral triangle of side $2$.}
\end{center}

\begin{example}[Hexagonal packing]
\label{L1:ex:hex}
Consider the equilateral triangle of side $2$ spanned by three mutually tangent unit disks.
Its height is $\sqrt{2^2-1^2}=\sqrt3$, so its area is $\sqrt3$. It contains three
sectors of angle $\pi/3$, i.e.\ half a unit disk, of area $\pi/2$. The plane is tiled by
such triangles, so
\[
  \dens_{\mathrm{hex}}\;=\;\frac{\pi/2}{\sqrt3}\;=\;\frac{\pi}{2\sqrt3}\;=\;\frac{\pi}{\sqrt{12}}\;\approx\;0.9069 .
\]
\end{example}

\subsection{Lattice packings}

For lattices, we want to maximize the minimum distance between two points of a unit lattice
\added{(a lattice $\Lat\subset\R^2$ of covolume $1$)}.
\orig{crossed out: ``minimize the det of the span for a unit lattice''}
This minimum distance can be arbitrarily small, but it is bounded above.

\begin{theorem}[\added{Lagrange}]
\label{L1:thm:lattice2}
Every lattice packing of congruent disks in $\R^2$ has density at most $\pi/\sqrt{12}$.
\added{Equality holds only for the hexagonal lattice.}
\end{theorem}

\begin{proof}
Scale so that $\Lat$ has covolume $1$, and let $2r$ be its minimal distance, so that disks of
radius $r$ centred at $\Lat$ form a packing of density $\pi r^2$. Let $v$ be a shortest
vector, $|v|=2r$. The lattice is the union of translates of the row $\Z v$, which lie on
parallel lines. Since the covolume is $1$, adjacent lines are at distance $h=1/(2r)$.
Note that this is the height of a fundamental parallelogram; the second basis vector
need not be orthogonal to $v$.

\begin{center}
\begin{tikzpicture}[scale=1.1]
  \draw (-0.5,0) -- (6.5,0);
  \draw (-0.5,1.9) -- (6.5,1.9);
  \foreach \x in {0,2,4,6} \fill (\x,0) circle (1.6pt);
  \fill (3,1.9) circle (1.6pt);
  \draw[dashed] (2,0) -- (3,1.9) -- (4,0);
  \draw[dotted] (3,0) -- (3,1.9);
  \draw[gray] (2,0) ++(0:2) arc (0:180:2);
  \node[below] at (3,0) {$2r$};
  \draw[<->] (6.3,0) -- (6.3,1.9) node[midway,right] {$h=\frac1{2r}$};
  \node[below] at (2,0) {\footnotesize $0$};
  \node[below] at (4,0) {\footnotesize $v$};
\end{tikzpicture}
\end{center}

\added{Translating along the row by multiples of $v$, we may take a lattice point $w$ in the
adjacent row whose orthogonal projection onto the first line lies within distance $r$ of
some lattice point there.} Since $|w-u|\ge 2r$ for all lattice points $u\ne w$,
\[
  h^2+r^2\;\ge\;(2r)^2,\qquad\text{so}\qquad \frac1{2r}\;=\;h\;\ge\;\frac{\sqrt3}{2}\,(2r)\;=\;\sqrt3\,r .
\]
Hence $r^2\le \frac1{2\sqrt3}$, i.e.\ $2r\le\sqrt{2/\sqrt3}$, and
\[
  \dens\;=\;\frac{\pi r^2}{1}\;\le\;\frac{\pi}{2\sqrt3}\;=\;\frac{\pi}{\sqrt{12}} .
\]
\added{Equality forces $h=\sqrt3 r$ with offset exactly $r$, which is the hexagonal lattice.}
\end{proof}
\orig{p.\,9 is an unnumbered duplicate of this computation}

\section{Partitioning space: Voronoi cells and Delaunay triangulations}
\label{L1:sec:voronoi}
\orig{\S1.5}

For packings that are not lattices, we need some analogous method to partition space.
Classically, Dirichlet--Voronoi diagrams work well.

\begin{definition}[Voronoi cell]
\label{L1:def:voronoi}
Given a packing with bodies $P_i$, the \emph{(Dirichlet--)Voronoi cell} associated with $P_i$ is
\[
  D_i\;=\;\{x\in\R^2 : d(x,P_i)<d(x,P_j)\ \text{for all } j\neq i\}.
\]
\end{definition}

When the $P_i$ are congruent disks and $d$ is the standard distance, this is the same as using
the distance to the centres of the disks. \added{Indeed $d(x,P_i)=|x-c_i|-1$ for $x\notin P_i$.}

Voronoi cells have nice properties:
\begin{itemize}[nosep]
  \item each cell is an intersection of half-spaces determined by the centres
        \added{($D_i=\bigcap_{j\ne i}\{x:|x-c_i|<|x-c_j|\}$)}, so it is convex;
  \item the cells partition space up to a boundary term \added{(a set of measure zero)}.
\end{itemize}

\begin{center}
\begin{tikzpicture}[scale=0.6]
  \clip (-3.2,-2.6) rectangle (5.2,4.3);
  \foreach \j in {-2,...,3} {
    \foreach \i in {-3,...,4} {
      \pgfmathsetmacro{\x}{2*\i+\j}
      \pgfmathsetmacro{\y}{sqrt(3)*\j}
      \draw[gray!60] (\x,\y) circle (1);
      \draw[blue!70!black,thick] (\x,\y) ++(30:1.1547) -- ++(150:1.1547) -- ++(210:1.1547)
         -- ++(270:1.1547) -- ++(330:1.1547) -- ++(30:1.1547) -- cycle;
      \draw[red!70!black,densely dashed] (\x,\y) -- ++(2,0);
      \draw[red!70!black,densely dashed] (\x,\y) -- ++(1,{sqrt(3)});
      \draw[red!70!black,densely dashed] (\x,\y) -- ++(-1,{sqrt(3)});
      \fill (\x,\y) circle (1.3pt);
    }
  }
\end{tikzpicture}
\\[2pt]
{\small \added{Hexagonal packing: Voronoi cells (solid hexagons) and Delaunay triangles (dashed).}}
\end{center}

The dual notion is the \emph{Delaunay triangulation}.
\begin{itemize}[nosep]
  \item It is not unique in general (e.g.\ four points at the vertices of a square).
  \item It is harder to define directly. Its characterization: the circumcircles of the
        Delaunay triangles are empty of other points (vertices).
  \item The circumcentre of a Delaunay triangle is a Voronoi vertex.
\end{itemize}
\begin{addedblock}
Precisely: for a locally finite set $X\subset\R^2$ (here, the centres), a triangle with vertices
in $X$ is \emph{Delaunay} if its open circumdisk contains no point of $X$. A Delaunay triangulation is
a triangulation of (the convex hull of) $X$ by Delaunay triangles. It is unique when no four
points of $X$ are concyclic.
\end{addedblock}

\subsection{Existence via lifting}

Consider a collection of points in the plane and lift them to the paraboloid
$z=x^2+y^2$.

\begin{lemma}
\label{L1:lem:lifting}
Let $a,b,c,d$ be points in the $(x,y)$-plane, and lift them to
$\hat a,\hat b,\hat c,\hat d$ on $z=x^2+y^2$ \added{(i.e.\ $\hat p=(p,|p|^2)$)}. Then $d$ lies
inside the circumcircle of $a,b,c$ if and only if $\hat d$ lies in the lower half-space of the
plane $\langle\hat a,\hat b,\hat c\rangle$ \added{spanned by $\hat a,\hat b,\hat c$}.
\end{lemma}

\begin{proof}
A plane tangent to the paraboloid at the point above $(p,q)$ has the form
\[
  z\;=\;2px+2qy-(p^2+q^2).
\]
Shifting it up by $h^2$ gives
\[
  z\;=\;2px+2qy-(p^2+q^2)+h^2 .
\]
\fixed{signs: notes have $+(p^2+q^2)$ and $(x+p)^2$, $(y+q)^2$}
It meets the paraboloid where
\[
  x^2+y^2\;=\;2px+2qy-(p^2+q^2)+h^2
  \quad\Longleftrightarrow\quad (x-p)^2+(y-q)^2\;=\;h^2 .
\]
So the plane through $(p,q,p^2+q^2)$, parallel to the tangent plane there, at height $h^2$
above it, projects its intersection with the paraboloid onto the circle of radius $h$
about $(p,q)$ in the $(x,y)$-plane.
\begin{addedblock}
Every non-vertical plane $z=\alpha x+\beta y+\gamma$ meeting the paraboloid in more than one point
is of this form, with $p=\alpha/2$, $q=\beta/2$ and $h^2=\gamma+p^2+q^2>0$. Apply this to the plane
through $\hat a,\hat b,\hat c$. The corresponding circle passes through $a,b,c$, so it is their
circumcircle. A point $\hat x$ of the paraboloid lies strictly below this plane iff
$x^2+y^2<2px+2qy-(p^2+q^2)+h^2$, i.e.\ iff $(x-p)^2+(y-q)^2<h^2$, i.e.\ iff $x$ lies strictly
inside the circumcircle.
\end{addedblock}
\end{proof}

\begin{corollary}[Existence of Delaunay triangulations]
\label{L1:cor:delaunay-exists}
The existence of the convex hull of the lifted points implies the existence of a Delaunay triangulation.
\added{More precisely: project the faces of the lower convex hull of $\{\hat x:x\in X\}$ to the plane.
By Lemma~\ref{L1:lem:lifting}, a lower face has no lifted point below its plane, so its projection has
an empty circumcircle. Triangulating any non-triangular faces (which come from concyclic points)
gives a Delaunay triangulation. For the infinite, locally finite point sets of a packing, apply this
to large finite pieces; the triangles near a fixed point stabilize.}
\end{corollary}

\section{Thue's theorem}
\label{L1:sec:thue}
\orig{\S1.6; a first ``1.6 Thm'' on p.\,13 is abandoned}

Let $P$ be a saturated packing of unit disks in $\R^2$ and $DT$ a Delaunay triangulation of its set of
centres. Every edge of $DT$ has length at least $2$, since the disks do not overlap.
\added{Moreover, every triangle of $DT$ has circumradius $R<2$. Its open circumdisk contains no
centres, so the circumcentre is at distance $\ge R$ from every centre. If $R\ge2$, a unit disk centred
there could be added, contradicting saturation.}

\begin{lemma}
\label{L1:lem:angle}
The largest angle $\theta$ of a triangle $ABC\in DT$ of a saturated packing satisfies
\[
  \frac\pi3\;\le\;\theta\;<\;\frac{2\pi}3 .
\]
\end{lemma}
\unclear{upper inequality written ambiguously; the proof gives strict ``$<$''}

\begin{proof}
The lower bound is obvious, since $\theta$ is the largest angle of a triangle.
Assume $\theta\ge 2\pi/3$. Let $A$ be the smallest angle. Then $A\le\frac12(\pi-\theta)\le\pi/6$,
so $\sin A\le\frac12$. By the circumradius formula, and since the side $BC$ opposite $A$ has
length at least $2$,
\[
  R\;=\;\frac12\,\frac{|BC|}{\sin A}\;\ge\;\frac12\cdot\frac{2}{1/2}\;=\;2 .
\]
\fixed{notes: ``$=4$''; $\frac12\cdot\frac{2}{1/2}=2$}
\added{This contradicts $R<2$.}
\end{proof}

\begin{lemma}
\label{L1:lem:triangle-density}
The density in a single Delaunay triangle is at most $\pi/\sqrt{12}$, with equality
if and only if the triangle is equilateral \added{(of side $2$)}.
\end{lemma}
\begin{addedblock}
Here the \emph{density} of a triangle $T\in DT$ means $\frac{\pi/2}{\operatorname{area}(T)}$:
the three sectors of the unit disks at the vertices of $T$ have angles summing to $\pi$, hence total area $\pi/2$.
\end{addedblock}

\begin{proof}
Let $B$ be the largest angle of $ABC$. Then
\[
  \operatorname{area}(ABC)\;=\;\tfrac12\,|AB|\,|CB|\,\sin B
  \;\ge\;\tfrac12\cdot2\cdot2\cdot\min_{B\in[\pi/3,\,2\pi/3)}\sin B
  \;=\;\tfrac12\cdot2\cdot2\cdot\tfrac{\sqrt3}{2}\;=\;\sqrt3,
\]
using Lemma~\ref{L1:lem:angle}, and the minimum is attained at $B=\pi/3$. So $\operatorname{area}(T)\ge\sqrt3$
for every $T\in DT$, and the density of $T$ is at most $\frac{\pi/2}{\sqrt3}=\frac{\pi}{2\sqrt3}=\frac{\pi}{\sqrt{12}}$.
\added{Equality requires $|AB|=|CB|=2$ and $B=\pi/3$, i.e.\ an equilateral triangle of side $2$.}
\end{proof}

\begin{theorem}[\added{Thue; Fejes T\'oth}]
\label{L1:thm:thue}
Every packing $P$ of unit disks in the plane satisfies $\dens^+(P)\le\pi/\sqrt{12}$.
\end{theorem}

\begin{proof}
\added{Saturating $P$ does not decrease $\dens^+$ (\S\ref{L1:sec:saturation}), so we may assume $P$ is saturated.}
For a finite union of Delaunay triangles, the density is
\[
  \sum_{T\in DT}\frac{\operatorname{area}(T)}{\operatorname{area}(DT)}\cdot\dens(T)
  \;\le\;\frac{\pi}{\sqrt{12}},
\]
a weighted average of the densities of the triangles, each at most $\pi/\sqrt{12}$
by Lemma~\ref{L1:lem:triangle-density}. Hence any finite union of Delaunay triangles has density
at most $\pi/\sqrt{12}$.

For a large disk $R\,\mathbb{D}^2$, most of its area is covered by the Delaunay triangles inside it,
so the error in the limit is a boundary term.
\begin{addedblock}
Every Delaunay triangle has circumradius $<2$, so diameter $<4$. Let $U_R$ be the union of the
triangles meeting $R\,\mathbb{D}^2$. Then $R\,\mathbb{D}^2\subset U_R\subset(R+4)\mathbb{D}^2$. The sectors
in the triangles of $U_R$ account for every disk centred in $R\,\mathbb{D}^2$. The disks meeting
$R\,\mathbb{D}^2$ but centred outside it lie in the annulus $(R+1)\mathbb{D}^2\setminus(R-1)\mathbb{D}^2$, of area $O(R)$. So
\[
  \vol(P\cap R\,\mathbb{D}^2)\;\le\;\frac{\pi}{\sqrt{12}}\,\vol(U_R)+O(R)
  \;\le\;\frac{\pi}{\sqrt{12}}\,\pi(R+4)^2+O(R).
\]
Dividing by $\pi R^2$ and letting $R\to\infty$ gives $\dens^+(P)\le\pi/\sqrt{12}$.
\end{addedblock}
\end{proof}
\orig{p.\,16 sketch: large concentric disks}

\begin{remark}
\added{Combined with Example~\ref{L1:ex:hex}, this shows that the hexagonal packing is the
densest packing of congruent disks in the plane. The statement goes back to Thue (1892, 1910).
The first proof generally accepted as complete is due to L.~Fejes T\'oth (1940); the argument above
follows his Delaunay-triangle approach. For lattice packings the result is Lagrange's
(Theorem~\ref{L1:thm:lattice2}).}
\end{remark}
% ===== end chapters/L01.tex =====
% ===== begin chapters/L02.tex =====
% MAT670 -- Lecture 2 (transcribed from MAT670_2, handwritten pp. 17--31)
\chapter{Lecture 2: The Fejes T\'oth inequality and Blichfeldt's bound}\label{L2:chap}

\section{Last time: Thue's theorem in the plane}\label{L2:sec:thue}
\orig{Notes pp.~17--31 (MAT670\_2).}
Last time we gave a modern take on Thue's result in $\R^2$:

\begin{theorem}[Thue]\label{L2:thm:thue}
The density of a packing of the plane by congruent circles is bounded by the density of the regular
hexagonal packing\added{, namely $\pi/\sqrt{12}\approx 0.9069$}.
\end{theorem}

\begin{center}
\begin{tikzpicture}[scale=0.9]
  \draw[thick] (0,0) circle (0.5);
  \foreach \k in {0,...,5} { \draw[thick] ({\k*60}:1) circle (0.5); }
\end{tikzpicture}
\end{center}

The idea of the proof was to use \emph{saturation} and the \emph{Delaunay triangulation}.

\begin{remark}
In the proof we left out the error analysis of the limits, etc.
\end{remark}

\begin{remark}
Rogers showed a density bound coming from the regular simplex.
\added{Precisely (C.~A.~Rogers, 1958): the density of any packing of unit balls in $\R^n$ is at most the
fraction of a regular simplex of edge length $2$ that is covered by the $n+1$ unit balls centred at its
vertices. For $n=2$ this is exactly $\pi/\sqrt{12}$, so Rogers' bound generalizes Thue's theorem. See Section~\ref{L2:sec:rogers} for a sketch.}
\end{remark}

% ---------------------------------------------------------------------------
\section{The Fejes T\'oth inequality for \texorpdfstring{$\Sph^2$}{S2}}\label{L2:sec:FT}

All distances on $\Sph^2$ are spherical (angular) distances.

\begin{theorem}[Fejes T\'oth\added{, 1943}]\label{L2:thm:FT}
For any $n\ge 3$ points on $\Sph^2$ there exists a pair with spherical distance
\begin{equation}\label{L2:eq:FT}
d \;\le\; \arccos\!\left(\frac{\cot^2(\omega)-1}{2}\right),
\qquad \omega=\frac{n}{n-2}\cdot\frac{\pi}{6}. \tag{$*$}
\end{equation}
\end{theorem}
\fixed{Notes have $\cot(\omega)$; must be $\cot^2(\omega)$.}

\begin{addedblock}
Check: for $n=12$, $\omega=\pi/5$, and $(\cot^2(\pi/5)-1)/2 = 1/\sqrt5$, so
$d\le\arccos(1/\sqrt5)\approx 63.43^\circ$, the edge of the regular icosahedron. Similarly $n=3,4,6$ give
$120^\circ$, $\arccos(-1/3)\approx109.47^\circ$ and $90^\circ$ (equilateral triangle on a great circle,
regular tetrahedron, regular octahedron). These are exactly the cases in which the bound is attained.
\end{addedblock}

\begin{remark}\label{L2:rem:capdensity}
The inequality can be rewritten as a density result for spherical caps of diameter $d$ (i.e.\ angular
radius $d/2$): a cap of angular radius $r$ has area $2\pi(1-\cos r)$, so the density of $n$ non-overlapping
caps of diameter $d$ is
\[
\frac{n\cdot 2\pi\bigl(1-\cos(d/2)\bigr)}{4\pi}
\added{\ \le\ \frac n2\Bigl(1-\cos\tfrac{d_n}{2}\Bigr)}
\;\longrightarrow\; \frac{\pi}{\sqrt{12}}
\added{\quad (n\to\infty),}
\]
\added{where $d_n$ is the right-hand side of \eqref{L2:eq:FT}.} So in the limit we recover Thue's bound.
\end{remark}

\begin{exercise}
Verify the limit in Remark~\ref{L2:rem:capdensity}. \added{(Hint: as $n\to\infty$, $d_n\to 0$ and
$\frac{\sqrt3}{4}d_n^2\sim \frac{4\pi}{2n}$; see also Exercise~\ref{L2:ex:Lhuilier}.)}
\end{exercise}

The ``flat'' values also converge, but it is not so nice to take the limit.
\unclear{``Flat values'': presumably the planar triangles of the convex hull instead of spherical ones.}

A natural question: is there a good \emph{lower} bound, i.e.\ can we construct a sequence of sets of
points that achieves this bound?
\begin{addedblock}
Equality in \eqref{L2:eq:FT} holds only for $n=3,4,6,12$. Asymptotically the bound is sharp: there are
configurations (e.g.\ suitable subdivisions of the icosahedron) whose cap density tends to
$\pi/\sqrt{12}$.
\end{addedblock}

\subsection{A lemma on spherical triangles}

To prove the Fejes T\'oth theorem we need the following lemma.

\begin{lemma}\label{L2:lem:triangle}
Let $ABC$ be a spherical triangle and let $\overline{AB}$ be its shortest side. Let $ABC'$ be the
equilateral spherical triangle drawn on $\overline{AB}$\added{, with $C'$ on the same side of the great
circle $AB$ as $C$}. If
\[
\operatorname{Area}(ABC) < \operatorname{Area}(ABC'),
\]
then the spherical radius of the circumcircle of $ABC$ (the circle in which the plane through $A,B,C$
meets $\Sph^2$) is greater than the length of $\overline{AB}$.
\end{lemma}

\begin{proof}
Write $a$ for the length of $\overline{AB}$.

\emph{Step 1: draw the picture.} Let $C'$ be such that $ABC'$ is equilateral, with $C'$ on the same side
of the great circle $AB$ as $C$. Let $\alpha,\beta,\gamma$ be the circles of radius $a$ centred at
$A,B,C'$ respectively\added{; thus $\alpha$ passes through $B$ and $C'$, $\beta$ through $A$ and $C'$,
and $\gamma$ through $A$ and $B$}. Let $A'$ be the second intersection point of $\alpha$ and $\gamma$
\added{(so that $AC'A'$ is equilateral)}, and $B'$ the second intersection point of $\beta$ and $\gamma$.
Let $-A,-B$ denote the antipodes of $A,B$.

\begin{figure}[ht]
\centering
\begin{tikzpicture}[scale=0.95]
  \begin{scope}
    \clip (-4.3,-2.5) rectangle (4.3,5.3);
    % admissible region for C: C' side of AB, below the Lexell circle, outside alpha and beta
    \fill[gray!18] (0,1.470) circle (3.522);
    \fill[white] (0,2.885) circle (2.885);
    \fill[white] (-2.683,-3.482) circle (4.396);
    \fill[white] (2.683,-3.482) circle (4.396);
    % circles
    \draw[thick] (0,1.470) circle (3.522);                 % great circle AB
    \draw[thick,dashed] (0,2.885) circle (2.885);          % Lexell circle
    \draw (-2.683,-3.482) circle (4.396);                  % alpha
    \draw (2.683,-3.482) circle (4.396);                   % beta
    \draw (0,0) circle (2.241);                            % gamma
  \end{scope}
  % equilateral triangle ABC' (great-circle arcs)
  \draw[red!70!black,densely dotted,thick] (-1.368,-1.775) -- (-1.122,-1.868) -- (-0.875,-1.941) -- (-0.626,-1.995) -- (-0.376,-2.031) -- (-0.125,-2.049) -- (0.125,-2.049) -- (0.376,-2.031) -- (0.626,-1.995) -- (0.875,-1.941) -- (1.122,-1.868) -- (1.368,-1.775)
     -- (1.212,-1.573) -- (1.067,-1.384) -- (0.930,-1.207) -- (0.800,-1.038) -- (0.676,-0.877) -- (0.557,-0.722) -- (0.441,-0.572) -- (0.328,-0.426) -- (0.218,-0.283) -- (0.109,-0.141) -- (0,0)
     -- (-0.109,-0.141) -- (-0.218,-0.283) -- (-0.328,-0.426) -- (-0.441,-0.572) -- (-0.557,-0.722) -- (-0.676,-0.877) -- (-0.800,-1.038) -- (-0.930,-1.207) -- (-1.067,-1.384) -- (-1.212,-1.573) -- cycle;
  % points
  \foreach \p/\n/\pos in {(-1.368,-1.775)/A/below left, (1.368,-1.775)/B/below right, (0,0)/C'/above,
      (-2.065,0.870)/A'/left, (2.065,0.870)/B'/right, (2.790,3.620)/-A/right, (-2.790,3.620)/-B/left}
      { \fill \p circle (1.6pt) node[\pos,font=\small] {$\n$}; }
  \fill (2.7,1.2) circle (1.6pt) node[right,font=\small] {$C$};
  % labels
  \node[font=\small] at (-4.05,1.0) {$\alpha$};
  \node[font=\small] at (4.05,1.0) {$\beta$};
  \node[font=\small] at (-1.75,1.65) {$\gamma$};
  \node[font=\small,align=center] at (0,4.25) {Lexell\\circle};
  \draw[thin,gray] (3.47,2.1) -- (3.95,2.4) node[right,font=\small,black] {great circle $AB$};
\end{tikzpicture}
\caption{The configuration of Lemma~\ref{L2:lem:triangle} (side $a=70^\circ$), drawn in stereographic
projection centred at $C'$, so circles on $\Sph^2$ appear as circles. Shaded: the points between the great circle $AB$ and the Lexell circle that lie outside $\alpha$ and
$\beta$ (the possible positions of $C$); they avoid the disc bounded by $\gamma$.}
\label{L2:fig:lexell}
\end{figure}
\orig{Hand sketch p.~19 has the great circle $AB$ as a straight line; redrawn exactly here.}

\emph{Step 2: show the picture is ``correct''.} Note that the spherical triangles $ABA'$, $ABB'$ and $ABC'$
all have equal area. \added{(For $ABA'$: it is isosceles with legs $a$ and apex angle $2\theta$, where
$\theta$ is the angle of $ABC'$; a direct computation with the formula
$\tan\frac E2=\frac{\sin\varphi\,\tan\frac b2\tan\frac c2}{1+\cos\varphi\,\tan\frac b2\tan\frac c2}$ for the
area $E$ of a triangle with sides $b,c$ and included angle $\varphi$, together with
$\tan^2\frac a2=1-2\cos\theta$, shows that both areas agree.)} Hence the circle in which the plane through
$\{A',B',C'\}$ meets the sphere is the locus of points $X$ with
$\operatorname{Area}(ABX)=\operatorname{Area}(ABC')$ (on that side of $AB$): this is the
\emph{Lexell circle theorem}. \added{(Lexell: for fixed $A,B$, the locus of apexes $X$ of triangles $ABX$
of fixed area is an arc of a small circle through $-A$ and $-B$.)}

By assumption, $\operatorname{Area}(ABC)<\operatorname{Area}(ABC')$, so $C$ lies between the Lexell circle
and the great circle $AB$. As $\overline{AB}$ is the shortest edge of $ABC$, $C$ does not lie in
\added{the open discs bounded by} $\alpha$ or $\beta$. It follows \added{(see
Figure~\ref{L2:fig:lexell})} that $C$ does not lie in \added{the disc bounded by} $\gamma$.
\added{(The part of the disc of $\gamma$ on the $C'$ side of $AB$ and below the Lexell circle is covered by
the discs of $\alpha$ and $\beta$.)}
\unclear{This covering step is read off the picture in the notes.}

But $C$ is on the $C'$ side of $AB$. \added{The circles through $A$ and $B$ form a pencil whose centres lie
on the perpendicular bisector of $AB$; on the $C'$ side, the region enclosed by such a circle grows with its
radius. Since $\gamma$ is the member of the pencil centred at $C'$ and $C$ lies outside it,} the
circumcircle of $ABC$ has radius greater than the radius of $\gamma$, which is the length of
$\overline{AB}$.
\end{proof}

\subsection{Proof of the Fejes T\'oth theorem}

The case $n=3$ is trivial\added{: the three sides of a spherical triangle have total length at most
$2\pi$, so one of them is at most $2\pi/3=\arccos(-1/2)$}. So let $n\ge 4$ and let $p_1,\dots,p_n\in\Sph^2$.

Without loss of generality the convex hull $\conv\{p_1,\dots,p_n\}$ contains the origin $0$
\added{in its interior}. Otherwise we could flow the points from the pole towards the equator
\added{of a hemisphere containing them} and increase all distances.
\begin{addedblock}
If all points lie in an open hemisphere with pole $N$, replace the colatitude $\phi$ (angle from $N$) of
each point by $\lambda\phi$, $\lambda>1$ close to $1$. From
$\cos d=\cos\phi_1\cos\phi_2+\sin\phi_1\sin\phi_2\cos\Delta$ one checks
$\frac{\partial}{\partial\lambda}\cos d\le-(\phi_1-\phi_2)\sin(\phi_1-\phi_2)\le0$, with strict
inequality for distinct points, so all mutual distances increase.
\end{addedblock}
\unclear{Degenerate case (points in a closed, not open, hemisphere) needs a small perturbation argument.}

\emph{Step 1: triangulate the convex hull.}
\added{The convex hull is a polytope whose vertices are among the $p_i$; after a small perturbation (which
changes distances only slightly) we may assume it is simplicial, with all $n$ points as vertices.}
By the Euler characteristic, it has $2n-4$ faces.

\begin{exercise}\label{L2:ex:euler}
Show that a triangulated convex polytope with $n$ vertices has $2n-4$ faces.
\emph{Solution sketch:} $v-e+f=2$ and $3f=2e$ for triangulations; hence $2n-2e+2f=4$, so
$2n-3f+2f=4$, i.e.\ $f=2n-4$.
\end{exercise}

The triangulation of $\conv\{p_1,\dots,p_n\}$ gives a \emph{spherical net} by radial projection to the
sphere: the edges become geodesic arcs\added{, and since $0$ is in the interior of the hull the $2n-4$
spherical triangles tile $\Sph^2$}.

\begin{center}
\begin{tikzpicture}[scale=1.1]
  % tetrahedron
  \coordinate (P1) at (-1.0,-0.45); \coordinate (P2) at (0.45,-0.6);
  \coordinate (P3) at (0.85,-0.15); \coordinate (P4) at (0.05,0.75);
  \draw[thick] (P1)--(P2)--(P3)--(P4)--cycle (P4)--(P2);
  \draw[dashed] (P1)--(P3);
  \draw[->] (1.4,0.05) -- (2.3,0.05);
  % spherical net
  \begin{scope}[xshift=3.6cm]
    \draw (0,0) circle (1);
    \draw[thick] (-0.83,-0.38) to[bend right=18] (0.40,-0.80);
    \draw[thick] (0.40,-0.80) to[bend right=18] (0.83,-0.28);
    \draw[thick] (0.83,-0.28) to[bend right=22] (0.05,0.96);
    \draw[thick] (0.05,0.96) to[bend right=22] (-0.83,-0.38);
    \draw[thick] (0.05,0.96) to[bend left=10] (0.40,-0.80);
    \draw[dashed] (-0.83,-0.38) to[bend left=12] (0.83,-0.28);
    \foreach \p in {(-0.83,-0.38),(0.40,-0.80),(0.83,-0.28),(0.05,0.96)} \fill \p circle (1.3pt);
  \end{scope}
\end{tikzpicture}
\end{center}

\emph{Step 2.} Now, for $\{p_1,\dots,p_n\}$, suppose that
\[
\operatorname{length}\bigl(\widehat{p_ip_j}\bigr) > d \qquad\text{for all } i\ne j,
\]
where $d=d_n$ is the right-hand side of \eqref{L2:eq:FT}.

\begin{exercise}[L'Huilier]\label{L2:ex:Lhuilier}
Show that $d$ in \eqref{L2:eq:FT} is the side length of an equilateral spherical triangle of area
$\dfrac{4\pi}{2n-4}$. \added{(One can use L'Huilier's formula, the spherical analogue of Heron's formula.
Alternatively: the angle $\theta$ of such a triangle satisfies $3\theta-\pi=\frac{2\pi}{n-2}$, i.e.\
$\theta=2\omega$, and the spherical law of cosines for angles gives
$\cos d=\frac{\cos\theta}{1-\cos\theta}=\frac{\cot^2\omega-1}{2}$.)}
\end{exercise}

Then the triangle $p_ip_jp_k$ of minimal area in the spherical net satisfies the hypotheses of the Lemma:
its area is small, but its edges are all longer than $d$. Here ``small'' means: less than or equal to
the area of the equilateral triangle of side length $d$\added{, since $2n-4$ triangles of total area $4\pi$
tile the sphere. As the area of an equilateral triangle increases with its side, and the shortest side of
$p_ip_jp_k$ is longer than $d$, the hypothesis of Lemma~\ref{L2:lem:triangle} holds with strict
inequality.}

Hence the circumcircle of $p_ip_jp_k$ has radius larger than $d$. But the net construction implies that
the circumcircle is empty\added{: the face $p_ip_jp_k$ lies in a supporting plane of the convex hull, so no
$p_l$ lies in the open cap cut off by that plane (and this cap is smaller than a hemisphere, since $0$ is on
the other side)}.

So we may place a new point at the centre \added{of this cap; it has distance greater than $d$ from every
$p_l$}. This new collection of $n+1$ points satisfies the same inequalities for the \emph{same}~$d$!
\added{Moreover the new hull still contains $0$ in its interior and has $2(n+1)-4$ faces, so its smallest
triangle again has area less than that of the equilateral triangle of side $d$, and we can repeat
indefinitely.}

This is absurd, e.g.\ by a volume bound: $\vol(\Sph^2)=4\pi$, and \added{the caps of radius $d/2$ around the
points are pairwise disjoint, so} we would have
\[
\sum_{i=1}^{N}\operatorname{area}\bigl(\operatorname{cap}(d/2)\bigr) < 4\pi \qquad\text{for all } N>0,
\]
\added{which is impossible since each cap has positive area.} \qed

% ---------------------------------------------------------------------------
\section{Higher-dimensional sphere packing: Blichfeldt's method}\label{L2:sec:blichfeldt}

Consider an early method of Blichfeldt (1929) to get density bounds in higher dimensions.
\added{(H.~F.~Blichfeldt, \emph{The minimum value of quadratic forms, and the closest packing of spheres},
Math.\ Ann.\ 101 (1929).)}

\begin{definition}\label{L2:def:convexbody}
A \emph{convex body} $C$ in $\R^n$ is a compact convex subset of $\R^n$ with non-empty interior.
\end{definition}
\orig{Crossed out: ``(so think $C=\mathbb{B}^n$)''.}

Also consider collections of isometries $\{\varphi_i\}_{i=1}^\infty$ of $\R^n$ such that
$\{\varphi_iC\}_{i=1}^\infty$ is a packing \added{(i.e.\ the interiors of the $\varphi_iC$ are pairwise
disjoint)}.

For now, $C=\Ball^n$ \added{is the closed unit ball}, and a packing is of the form
$\{\varphi_i\Ball^n\}_{i=1}^\infty$.

Consider replacing the ball with its characteristic function
\[
\chi(x)=\begin{cases}1, & |x|\le1,\\ 0, & |x|>1.\end{cases}
\]
So it has mass
\[
I(\chi)=\int_{\R^n}\chi(x)\,dx \;\added{=\vol(\Ball^n)},
\]
and a packing corresponds to the function
\[
x\longmapsto\sum_{i=1}^\infty\chi(\varphi_i^{-1}x)
\]
\added{(the characteristic function of $\bigcup_i\varphi_i\Ball^n$, up to the measure-zero set of boundary
points).} For a general convex body $C$ this works as well.

Furthermore, we can now replace $\chi$ with some other function $f$, and in general associate to a packing
the function
\[
x\longmapsto\sum_{i=1}^\infty f(\varphi_i^{-1}x).
\]
A density may be defined similarly to the packing density, where each object now has mass
\[
I(f)=\int_{\R^n}f(x)\,dx
\]
spread over $\operatorname{supp}(f)$. So:
\[
\Delta_f \;=\; \dens_{\text{packing}}\cdot\frac{I(f)}{\vol(C)}.
\]
\added{Here $\dens_{\text{packing}}$ is the density of the packing $\{\varphi_iC\}$ (Lecture~1), i.e.\ the
average number of copies per unit volume times $\vol(C)$; $\Delta_f$ is the average of
$\sum_i f(\varphi_i^{-1}x)$ over $x\in\R^n$.}

We can think of this as spreading out and renormalizing the characteristic function in some (nice,
hopefully) way: ``cut up the sphere and spread it out.''

\textbf{Insight.} There are functions $f$ for which the sums $\sum_{i=1}^\infty f(\varphi_i^{-1}x)$ are
uniformly pointwise bounded over $\R^n$ \emph{and} over all collections of isometries $\{\varphi_i\}$ such
that $\{\varphi_iC\}$ is a packing!

% ---------------------------------------------------------------------------
\section{Blichfeldt gauges}\label{L2:sec:gauge}

\begin{definition}\label{L2:def:gauge}
Given a convex body $C$ in $\R^n$, a \added{non-negative, integrable} function $f$ on $\R^n$ is a
\emph{Blichfeldt gauge} for $C$ if for any collection $\{\varphi_i\}_{i=1}^\infty$ of isometries of
$\R^n$ such that $\{\varphi_iC\}_{i=1}^\infty$ is a packing,
\[
\sum_{i=1}^\infty f(\varphi_i^{-1}x)\le 1 \qquad\text{for all } x\in\R^n.
\]
\end{definition}
\fixed{Notes: ``isometries of $\R^d$''; should be $\R^n$.}
\orig{Margin: ``add non-negative condition''.}

\begin{lemma}\label{L2:lem:gauge}
If $f$ is a Blichfeldt gauge for $C$, then the density of any packing of $\R^n$ by copies of $C$ satisfies
\[
\dens(C)\;\le\;\frac{\vol(C)}{I(f)}. \tag{$**$}
\]
\end{lemma}
\fixed{Notes have strict $<$; equality can occur (e.g.\ $f=\chi_{\interior C}$, $C$ a cube).}

\begin{proof}[Proof sketch]
By the definition, $\Delta_f\le 1$ \added{(the average of a function bounded by $1$ is at most $1$)}.
Since $\Delta_f=\dens\cdot I(f)/\vol(C)$, this gives $(**)$. \added{The limiting argument that makes this
precise is carried out for balls in Section~\ref{L2:sec:bound}; the same argument works for any $C$ and any
gauge with bounded support.}
\end{proof}

% ---------------------------------------------------------------------------
\section{Back to \texorpdfstring{$C=\Ball^n$}{C = B\^{}n}}\label{L2:sec:radial}

We can consider radially symmetric functions by an averaging argument: there is no benefit to anisotropy,
since $\SO(n)$ kills it.
\begin{addedblock}
If $f$ is a gauge for $\Ball^n$, so is $\bar f(x)=\int_{\SO(n)}f(gx)\,dg$ (Haar probability measure):
for each $g$, $\{\varphi_ig^{-1}\Ball^n\}=\{\varphi_i\Ball^n\}$ is the same packing, so
$\sum_i f(g\varphi_i^{-1}x)\le1$, and averaging over $g$ gives $\sum_i\bar f(\varphi_i^{-1}x)\le 1$.
Moreover $I(\bar f)=I(f)$, so $\bar f$ gives the same bound.
\end{addedblock}
So it suffices to consider functions of the distance $|x|$! Rankin does a complicated study to get a good
function \added{(R.~A.~Rankin, \emph{On the closest packing of spheres in $n$ dimensions}, Ann.\ of Math.\
48 (1947))}. We consider
\[
f_0(x)=\begin{cases}1-\tfrac12|x|^2, & |x|\le\sqrt2,\\[2pt] 0, & |x|>\sqrt2.\end{cases}
\]

We can show the following for a packing of unit radius balls.

\begin{lemma}\label{L2:lem:f0}
$f_0$ is a Blichfeldt gauge for $\Ball^n$.
\end{lemma}

\begin{proof}
Consider $m$ unit balls \added{in a packing} with centres $x_1,\dots,x_m\in\R^n$. Then clearly, for all
pairs,
\[
|x_i-x_j|^2\ge 4. \tag{$\dagger$}
\]
Summing $(\dagger)$ over $1\le i<j\le m$ (the cross terms cancel\added{, via the identity
$\sum_{i<j}|x_i-x_j|^2=m\sum_i|x_i|^2-\bigl|\sum_ix_i\bigr|^2$}), we get
\[
m\sum_{i=1}^m|x_i|^2-\Bigl|\sum_{i=1}^m x_i\Bigr|^2\;\ge\;2m(m-1),
\]
and therefore
\[
\sum_{i=1}^m|x_i|^2\;\ge\;2(m-1).
\]
\orig{Notes write the centres in coordinates $(a_i,b_i,\dots,n_i)$.}
The origin here is an arbitrary point. So, given $x\in\R^n$, consider the collection of the $m$ balls whose
centres lie within $\sqrt2$ of $x$ \added{(the only ones contributing to $\sum_if_0(x-x_i)$)}, with masses
defined by $f_0$ centred at the $x_i$. Writing $r_i=|x-x_i|$, we see that
\[
\sum_{i=1}^m\frac{2-r_i^2}{2}\;\le\;\frac{2m-2(m-1)}{2}\;=\;1,
\]
and so $f_0$ is a Blichfeldt gauge.
\end{proof}

% ---------------------------------------------------------------------------
\section{Blichfeldt's bound}\label{L2:sec:bound}

So we can integrate this, using spherical shells:
\[
I(f_0)=\int_0^{\sqrt2}\frac{2-r^2}{2}\,d\bigl(\vol(\Ball^n)\,r^n\bigr)
=\frac{2^{\frac n2+1}}{n+2}\,\vol(\Ball^n).
\]
\fixed{Notes first had $\frac{4\cdot2^{n/2}}{n+2}$; corrected in red to $\frac{2^{(n+2)/2}}{n+2}$.}

\begin{exercise}\label{L2:ex:integral}
Verify this integral using spherical shells. \added{($I(f_0)=n\vol(\Ball^n)\int_0^{\sqrt2}
(1-\tfrac{r^2}{2})r^{n-1}dr=\vol(\Ball^n)\,2^{n/2}\bigl(1-\tfrac{n}{n+2}\bigr)$.)}
\end{exercise}

\subsection*{Limit argument for a box}

Take a large box of side length $t$ with $K$ balls in it.
\orig{Notes write $|I|$ for $K$.}
\added{Their centres lie in the concentric box of side $t-2$, so the supports of the gauges
$f_0(x-x_i)$ lie in the concentric box of side} $t+2\sqrt2-2$. Since $f_0\ge 0$ and
$\sum_if_0(x-x_i)\le1$ everywhere, integrating over this larger box gives
\[
(t+2\sqrt2-2)^n\;\ge\;K\int_{\R^n}f_0(x)\,dx\;=\;\frac{K\,\vol(\Ball^n)\,2^{\frac n2+1}}{n+2}.
\]
Hence the density of these $K$ balls in the box $[0,t]^n$ satisfies
\[
\frac{K\vol(\Ball^n)}{t^n}\;\le\;\frac{n+2}{2^{\frac n2+1}}\Bigl(1+\frac{2\sqrt2-2}{t}\Bigr)^n,
\]
and letting $t\to\infty$:

\begin{theorem}[Blichfeldt's bound]\label{L2:thm:blichfeldt}
For any packing of $\R^n$ by unit balls,
\[
\dens\;\le\;\frac{n+2}{2^{\frac n2+1}} \;=\; \frac{n+2}{2}\,2^{-n/2}.
\]
\end{theorem}
\fixed{Notes had $\frac{n+2}{4\cdot2^{n/2}}$ (red correction in margin).}

One can do this explicitly for balls too; we are also still missing a boundary term.
\unclear{Probably: balls meeting the boundary of the box; they contribute $O(t^{n-1})$, negligible as $t\to\infty$.}

\begin{remark}
\added{For $n=2$ the bound is the trivial $\dens\le1$; for $n=3$ it gives $5/(4\sqrt2)\approx0.884$, compared
with $\pi/\sqrt{18}\approx0.7405$ for the FCC packing (Rogers' simplex bound gives $\approx0.7796$).
Blichfeldt's bound decays like $2^{-0.5n}$ up to polynomial factors; Rankin (1947) improved the constant,
and Kabatiansky--Levenshtein (1978) improved the exponent to $2^{-0.599n}$.}
\end{remark}

%======================================================================
\section{Extension: Rogers' simplex bound}
\label{L2:sec:rogers}

\begin{addedblock}
This section is an added extension, not in the handwritten notes. It is set in black for readability.
\end{addedblock}

The remark in Section~\ref{L2:sec:thue} mentioned Rogers' bound. Here is what it says and why it is
true, at the level of a sketch. It is the natural sequel to Blichfeldt: instead of spreading the mass of
each ball over a larger region, we cut space into pieces adapted to the packing and bound the density
in each piece.

\paragraph{The tightest local cluster.} The densest possible arrangement of $n+1$ unit balls is a
regular simplex of edge $2$ with a ball at each vertex: every pair touches. Let
\[
  \sigma_n=\frac{\vol\big(S\cap\bigcup_{v}\Ball(v,1)\big)}{\vol S}
\]
be the fraction of such a simplex $S$ that is covered by its $n+1$ vertex balls. If $\omega_n$ is the
fraction of the unit ball lying in the cone at a vertex of $S$ (its normalized solid angle), then
\[
  \sigma_n=\frac{(n+1)\,\omega_n\vol\Ball^n}{\vol S},\qquad \vol S=\frac{2^{n/2}\sqrt{n+1}}{n!}.
\]

\begin{theorem}[Rogers 1958]\label{L2:thm:rogers}
The upper density of any packing of unit balls in $\R^n$ is at most $\sigma_n$.
\end{theorem}

If regular simplices tiled space, Rogers' bound would be attained. They do for $n=2$, where
$\sigma_2=\pi/\sqrt{12}$ and the theorem is Thue's. They do not for $n\ge3$: already in $\R^3$ the
dihedral angle $\arccos\frac13$ does not divide $2\pi$, and twenty regular tetrahedra do not make an
icosahedron (Interlude~E). So for $n\ge3$ the bound is not attained (and the remark below records
better bounds).

\begin{center}
\renewcommand{\arraystretch}{1.15}
\begin{tabular}{rcccc}
\toprule
$n$ & Blichfeldt $\frac{n+2}{2^{n/2+1}}$ & Rogers $\sigma_n$ & best known packing & ratio best$/\sigma_n$\\
\midrule
$2$ & $1.0000$ & $0.9069$ & $0.9069$ ($A_2$) & $1.000$\\
$3$ & $0.8839$ & $0.7796$ & $0.7405$ ($A_3$) & $0.950$\\
$4$ & $0.7500$ & $0.6478$ & $0.6169$ ($D_4$) & $0.952$\\
$5$ & $0.6187$ & $0.5257$ & $0.4653$ ($D_5$) & $0.885$\\
$6$ & $0.5000$ & $0.4192$ & $0.3729$ ($E_6$) & $0.890$\\
$7$ & $0.3977$ & $0.3299$ & $0.2953$ ($E_7$) & $0.895$\\
$8$ & $0.3125$ & $0.2568$ & $0.2537$ ($E_8$) & $0.988$\\
\bottomrule
\end{tabular}
\end{center}
(The values of $\sigma_n$ were computed from the formula above, with $\omega_n$ evaluated as a Gaussian
orthant probability: $\omega_n$ is the probability that $n+1$ independent standard normals $Y_0,\dots,Y_n$
satisfy $Y_i>\bar Y$ for $i=1,\dots,n$.) Note how close the root lattices come, especially $A_3$, $D_4$ and $E_8$; the end
of Lecture~6 takes up why.

\paragraph{Sketch of the proof.}
\begin{enumerate}[label=(\arabic*)]
\item \emph{Reduce to Voronoi cells.} As in Lecture~1, we may assume the packing is saturated, and the
  density is (in the limit) an average of $\vol\Ball/\vol V(x)$ over the Voronoi cells $V(x)$. So it
  is enough to cut each cell into pieces, each of which is covered by the ball $\Ball(x,1)$ in a
  fraction at most $\sigma_n$.
\item \emph{Cut each cell into orthoschemes.} Pick a flag of faces of $V(x)$,
  $F_{n-1}\supset F_{n-2}\supset\dots\supset F_0$, with $\dim F_k=k$. Let $p_k$ be the foot of the
  perpendicular from $x$ to the affine hull of $F_k$. The simplex $\conv(x,p_{n-1},\dots,p_0)$ is an
  \emph{orthoscheme}: the steps $x\to p_{n-1}\to p_{n-2}\to\dots\to p_0$ are mutually orthogonal. Over all
  flags these simplices tile $V(x)$, at least when each foot $p_k$ lies in its face. (In general some
  pieces must be counted with signs, and handling this is part of Rogers' work.)
\item \emph{Each step is long.} A $k$-face $F_k$ consists of points equidistant from $x$ and at least
  $n-k$ further centres $y_1,\dots,y_{n-k}$, which can be chosen with $y_i-x$ linearly independent so
  that $\operatorname{aff}F_k$ is exactly the set of points equidistant from $x,y_1,\dots,y_{n-k}$. The foot $p_k$ is the circumcentre of the
  $(n-k)$-simplex $x,y_1,\dots,y_{n-k}$, so $\abs{x-p_k}$ is its circumradius. All its edges are at
  least $2$, and this forces the circumradius to be large. If $j+1$ points $v_0,\dots,v_j$ lie on a
  sphere of radius $r$ about the origin, then
  $\sum_{a<b}\abs{v_a-v_b}^2=(j+1)\sum_a\abs{v_a}^2-\abs{\sum_av_a}^2\le(j+1)^2r^2$. There are $j(j+1)/2$
  pairs, so some distance satisfies $\abs{v_a-v_b}^2\le 2(j+1)r^2/j$. Since every distance is at least $2$,
  \[
    \abs{x-p_k}\ \ge\ R_{n-k},\qquad R_j=\sqrt{\frac{2j}{j+1}},
  \]
  with equality for the regular simplex. Here $R_j$ is the circumradius of the regular $j$-simplex of edge $2$. In particular $R_1=1$, since a facet is at
  distance at least $1$, half the minimum distance $2$.
\item \emph{Monotonicity (the heart of the matter).} Among all orthoschemes with
  $\abs{x-p_k}\ge R_{n-k}$ for every $k$, the fraction covered by $\Ball(x,1)$ is largest when all of these
  are equalities. That extremal orthoscheme is one of the $(n+1)!$ congruent pieces of the barycentric
  subdivision of the regular simplex of edge $2$, so its covered fraction is exactly $\sigma_n$.
  Pushing a vertex of the orthoscheme outward adds volume that is farther from $x$, where the ball is
  thinner. Making this precise is Rogers' lemma, the technical core of his paper.
\end{enumerate}

For $n=2$ the picture is Thue's proof from Lecture~1, cut more finely. The faces are an edge at
distance $\ge1$ and a vertex at distance $\ge2/\sqrt3$, the extremal orthoscheme is the $30$--$60$--$90$
triangle that is one sixth of the equilateral triangle, and its covered fraction is $\pi/\sqrt{12}$.

\begin{remark}
\begin{enumerate}[nosep]
\item Asymptotically $\sigma_n\sim\frac{n}{e}\,2^{-n/2}$ (Rogers; Daniels). This beats Blichfeldt's
  $\frac{n+2}{2}2^{-n/2}$ only by the constant factor $2/e$. The Kabatiansky--Levenshtein bound of
  Lecture~3 is exponentially better in high dimensions.
\item The Fejes T\'oth inequality of Section~\ref{L2:sec:FT} is the same idea on the sphere $\Sph^2$.
  B\"or\"oczky (Acta Math.\ Acad.\ Sci.\ Hungar.\ 32 (1978)) extended the simplex bound to spheres and
  hyperbolic spaces of every dimension.
\item The Cohn--Elkies linear programming bound (Interlude~B) was, when it appeared (2003), the best known
  bound in dimensions $4$ through $36$, so in particular it improves on Rogers there. In dimensions $8$
  and $24$ it is sharp (Viazovska; Cohn--Kumar--Miller--Radchenko--Viazovska).
\item References: C.~A.~Rogers, \emph{The packing of equal spheres}, Proc.\ London Math.\ Soc.\ (3) 8
  (1958) 609--620; Conway--Sloane, \emph{Sphere Packings, Lattices and Groups}, Ch.~1.
\end{enumerate}
\end{remark}
% ===== end chapters/L02.tex =====
% ===== begin chapters/L03.tex =====
% MAT670 -- Lecture 3 (18.4.2016) -- transcribed from MAT670_3, pp. 32--48 of the handwritten notes
\chapter[Lecture 3: Blichfeldt gauges, Kabatjanskii--Levenshtein]{Lecture 3: Blichfeldt gauges and the Kabatjanskii--Levenshtein bound}
\label{L3:ch}

\noindent\emph{Lecture of 18 April 2016.}\orig{Notes pp.~32--48.}
Last time we proved the Fejes T\'oth inequality for point configurations on the sphere and
Blichfeldt's sphere packing bound. Today we make some remarks on both, look at the landscape of
known bounds for the packing density, prove the Kabatjanskii--Levenshtein-type bound via a Blichfeldt
gauge, and then extend the notion of a Blichfeldt gauge from balls to general convex bodies.

\added{Throughout, ``a packing of unit balls'' means a family of balls of radius $1$ in $\R^n$ with
pairwise disjoint interiors, i.e.\ with centers at mutual distance at least $2$.}

%======================================================================
\section{Some remarks}
\label{L3:sec:remarks}

\subsection*{Where is the Fejes T\'oth inequality sharp?}

Recall the Fejes T\'oth inequality: \added{if $n\ge 3$ points are placed on the unit sphere $\Sph^2\subset\R^3$,
then the minimal angular distance $d$ between two of them satisfies}
\begin{equation}\label{L3:eq:FT}
  d \;\le\; \arccos\!\left(\frac{\cot^2\omega-1}{2}\right),
  \qquad \omega=\frac{n}{n-2}\cdot\frac{\pi}{6}.
\end{equation}
\fixed{Notes have $\cot(\omega)$; the square is needed (check $n=12$).}%
The right-hand side is the angular edge length of the equilateral spherical triangle with area
$\frac{4\pi}{2n-4}$ on the unit sphere. \added{(A triangulation of the sphere with $n$ vertices has $2n-4$
triangular faces, so this is the area of a face when the $4\pi$ of total area is split evenly.)}

So the bound should be sharp exactly when the points are the vertices of an equilateral deltahedron
\added{(a polytope all of whose faces are congruent equilateral triangles) inscribed in the sphere}, i.e.\ for
\[
  n = 3,\ 4,\ 6,\ 12,
\]
and perhaps asymptotically as $n\to\infty$.
\begin{addedblock}
Check: $n=3$ gives $\omega=\pi/2$, $d\le\arccos(-\tfrac12)=120^\circ$ (an equilateral triangle on a great circle, the degenerate case);
$n=4$ gives $\cot^2\omega=\tfrac13$, $d\le\arccos(-\tfrac13)\approx109.47^\circ$ (regular tetrahedron);
$n=6$ gives $\omega=\pi/4$, $d\le 90^\circ$ (regular octahedron);
$n=12$ gives $\omega=\pi/5$, $d\le\arccos(1/\sqrt5)\approx 63.43^\circ$ (regular icosahedron).
Equality requires a triangulation by congruent equilateral triangles, so these are the only cases of equality;
as $n\to\infty$ the bound is asymptotically sharp.
\end{addedblock}
\orig{$n=3$ underlined in the notes.}

There are other proofs of the Fejes T\'oth inequality, but they are longer.

%======================================================================
\section{Aside: Blichfeldt's gauge revisited}
\label{L3:sec:aside}
\orig{Numbered ``3.2 aside''. Crossed out above: ``Again, there exists e.g.\ a trivial Blichfeldt gauge for $C$, given by $\chi_C$.''}

We saw last time that
\begin{equation}\label{L3:eq:f0}
  f_0(x)=\begin{cases} 1-\tfrac12\abs{x}^2, & \abs{x}\le\sqrt2,\\[2pt] 0, & \abs{x}>\sqrt2,\end{cases}
\end{equation}
is a Blichfeldt gauge for $\Ball^n$: \added{for every packing of unit balls with centers $c_1,c_2,\dots$,
$\sum_i f_0(x-c_i)\le 1$ for all $x\in\R^n$, and consequently the packing density satisfies
$\dens\le \vol(\Ball^n)/\int_{\R^n}f_0$.}
It is non-trivial for our purposes, that is,
\[
  \int_{\R^n} f_0(x)\,dV \;>\; \int_{\R^n}\chi_{\Ball^n}\,dV=\vol(\Ball^n)
  \qquad\text{(at least for } n>2).
\]
Indeed,
\begin{equation}\label{L3:eq:intf0}
  \int_{\R^n} f_0(x)\,dV=\frac{2^{1+\frac n2}}{n+2}\,\vol(\Ball^n).
\end{equation}
\begin{addedblock}
Proof: in polar coordinates, $dV = n\vol(\Ball^n)\,t^{n-1}\,dt\,d\sigma$ (with $\sigma$ the normalized surface measure), so
\[
  \int f_0 = n\vol(\Ball^n)\int_0^{\sqrt2}\Bigl(1-\frac{t^2}{2}\Bigr)t^{n-1}\,dt
  = \vol(\Ball^n)\Bigl(2^{n/2}-\frac{n}{2(n+2)}2^{(n+2)/2}\Bigr)
  = \frac{2\cdot 2^{n/2}}{n+2}\vol(\Ball^n).
\]
Hence $\int f_0>\vol(\Ball^n)$ iff $2^{1+n/2}>n+2$, which holds exactly for $n\ge3$; for $n=2$ there is equality and for
$n=1$ the gauge is worse than $\chi_{\Ball^1}$.
\end{addedblock}

\subsection*{From the gauge to the density bound}
Consider a packing $P$ of unit balls and a large cube $\mathrm{Box}(t)$ of side length $t$.
\begin{addedblock}
Split the balls meeting $\mathrm{Box}(t)$ into the \emph{inner} ones, whose centers lie in the concentric cube of side
$t-2\sqrt2$, and the \emph{boundary} ones. The support of $f_0(\cdot-c_i)$ for an inner center lies in $\mathrm{Box}(t)$, so
integrating $\sum_i f_0(x-c_i)\le 1$ over $\mathrm{Box}(t)$ gives
$N_{\mathrm{in}}\int f_0\le t^n$. The boundary centers lie in a shell of bounded width around $\partial\,\mathrm{Box}(t)$,
so there are at most $C\,t^{n-1}$ of them (by a volume argument, $C$ depending only on $n$).
\end{addedblock}
Therefore, for all $t$,
\[
  \frac{\vol\bigl(P\cap \mathrm{Box}(t)\bigr)}{\vol\bigl(\mathrm{Box}(t)\bigr)}
  \;\le\; \frac{n+2}{2^{1+\frac n2}} \;+\; \frac{C\,t^{n-1}}{t^n},
\]
\fixed{Notes: main term times $(1+\frac{\sqrt2-2}{t})^n$ (Lecture~2: $2\sqrt2-2$); not needed here, tends to $1$.}%
and letting $t\to\infty$,
\begin{equation}\label{L3:eq:blichfeldt}
  \dens(\Ball^n)\;\le\;\limsup_{t\to\infty}\frac{\vol(P\cap\mathrm{Box}(t))}{\vol(\mathrm{Box}(t))}\;\le\;\frac{n+2}{2^{1+\frac n2}}.
\end{equation}

\begin{remark}
We could equally well use
\[
  \dens=\limsup_{t\to\infty}\frac{\vol\bigcup\{\,\Ball\in P:\ \Ball\subseteq\mathrm{Box}(t)\,\}}{\vol\bigl(\mathrm{Box}(t)\bigr)},
\]
\added{counting only the balls entirely inside the box; since the balls meeting $\partial\,\mathrm{Box}(t)$ have total volume
$O(t^{n-1})$, this gives the same limit.}
\end{remark}

%======================================================================
\section{What does this look like?}
\label{L3:sec:landscape}
\orig{Section number overwritten: ``3.3'' / ``4''.}

The bound \eqref{L3:eq:blichfeldt} coming from $f_0$ gives the following density bounds
\added{(two-digit values as in the notes)}:
\begin{center}
\begin{tabular}{c|cccccccccc}
\toprule
$n$ & 1 & 2 & 3 & 4 & 5 & 6 & 7 & 8 & 9 & 10\\
\midrule
$\dfrac{n+2}{2^{1+n/2}}$ & 1.06 & 1 & .88 & .75 & .61 & .5 & .39 & .31 & .24 & .19\\
\bottomrule
\end{tabular}
\end{center}
\added{For comparison, the true optimal densities are $1$, $\pi/\sqrt{12}\approx .9069$, $\pi/\sqrt{18}\approx .7405$ in dimensions $1,2,3$,
and $\pi^4/384\approx .2537$ in dimension $8$.}

\subsection*{Asymptotic bounds}
\begin{center}
\begin{tabular}{p{0.44\textwidth}|p{0.44\textwidth}}
\textbf{Lower bounds} & \textbf{Upper bounds}\\
\midrule
$\dfrac{1}{2^n}$ \ (saturation) & $\dfrac{n+2}{2^{\frac n2+1}}$ \ (Blichfeldt)\\[10pt]
\quad$\vdots$ & \quad$\vdots$\\[4pt]
Ball (1992): $\dfrac{2n}{2^n}$ & $2^{-0.599n}$ \ (Kabatjanskii--Levenshtein)\\[10pt]
Vance, $n\equiv 0 \bmod 4$: $\dfrac{6}{e}\cdot\dfrac{n}{2^n}$ & \quad (Cohn, Elkies, Kumar, \dots)\\[10pt]
Venkatesh (sparse) \added{(2013)}: $\dfrac{1}{2}\cdot\dfrac{n\log\log n}{2^n}$ \ (!!!) & \\
\end{tabular}
\end{center}
\fixed{Notes: $e^{-\gamma}/2$; Venkatesh proves $\Delta_n\ge n\log\log n/2^{n+1}$, i.e.\ $1/2$.}%
\begin{addedblock}
Comments. (1) \emph{Saturation:} if a packing of unit balls is saturated (no further ball can be added), then the balls of
radius $2$ with the same centers cover $\R^n$, so $2^n\dens\ge 1$; saturated packings exist (e.g.\ by Zorn's lemma or a greedy
construction), hence $\dens(\Ball^n)\ge 2^{-n}$. (2) Ball (1992) proved $\dens\ge 2(n-1)\zeta(n)2^{-n}$ for lattice packings;
Vance (2011) proved $\dens\ge \frac{6n}{e\,2^n}$ for lattice packings when $4\mid n$; Venkatesh (\emph{A note on sphere packings in
high dimension}, IMRN 2013) proved $\dens\ge 65963\,n\,2^{-n}$ for all large $n$ and
$\dens\ge \frac12\,n\log\log n\,2^{-n}$ for infinitely many $n$ (the ``!!!'': the first improvement over linear growth of the
numerator). The $\log\log n$ comes from cyclotomic lattices: by Mertens' theorem $m/\varphi(m)$ can be as large as about
$e^{\gamma}\log\log m$. (3) \emph{Update after 2016:} Campos, Jenssen, Michelen and Sahasrabudhe (2023) proved
$\dens\ge(\frac12-o(1))\,n\log n\,2^{-n}$ for all $n$, and Klartag (2025) proved $\dens\ge c\,n^2 2^{-n}$ (even for lattice packings).
On the upper side, the exponent $0.599$ of Kabatjanskii--Levenshtein (1978) has only been improved by subexponential factors
(e.g.\ Cohn--Zhao 2014).
\end{addedblock}

\subsection*{Known good constructions}
There are known good constructions in low dimensions and special ones \added{(e.g.\ $D_4$, $E_8$, the Leech lattice $\Lambda_{24}$)}.
We need to talk about these at some point.\orig{Margin doodle (sad face) omitted.}

\subsection*{Exact results (upper bounds that are attained)}
\begin{itemize}
\item $d=1$: trivial.
\item $d=2$: we proved it. \added{Thue (1892, 1910); the first complete proof is due to L.~Fejes T\'oth (1940). The optimal lattice
packing was known to Lagrange and Gauss.}\unclear{Notes: ``(1892\dots ???) Thue''.}
\item $d=3$: Hales, twice, around 2000; a complicated geometric analysis. \added{Hales (with Ferguson), announced 1998, published in
Ann.\ of Math.\ 2005; the formal computer verification (Flyspeck project) was completed in 2014.}
\item $d=8$ and $d=24$: one month ago, via linear programming duality: one constructs a special (``magic'') auxiliary function.
\added{Viazovska (dimension 8, March 2016) and Cohn--Kumar--Miller--Radchenko--Viazovska (dimension 24, a week later);
both published in Ann.\ of Math.\ 2017. The LP bound itself is due to Cohn--Elkies (2003).}
\end{itemize}

\begin{remark}
Remember that the cubic lattice $\Z^n$ gives a packing of balls of radius $\frac12$ with density
\[
  \dens(\Z^n)=\vol\bigl(\tfrac12\Ball^n\bigr)=\frac{\pi^{n/2}}{2^n\,(\frac n2)!},
\]
\added{where $(\frac n2)!:=\Gamma(\frac n2+1)$. This decays superexponentially, much faster than the lower bound $2^{-n}$.}
\end{remark}

%======================================================================
\section{Kabatjanskii and Levenshtein}
\label{L3:sec:KL}

\begin{definition}\label{L3:def:M}
For $\varphi\in(0,\pi]$, let $M(d,\varphi)$ be the maximum number of points on $\Sph^{d-1}$ with pairwise
\added{angular} distance at least $\varphi$.
Equivalently, $M(d,\varphi)$ is the maximum number of points in $\R^d\setminus\{0\}$ such that the angle
\added{at the origin} spanned by any two of them is at least $\varphi$ \added{(project radially onto the sphere)}.
\end{definition}
\added{For example $M(d,\pi)=2$, and $M(d,\pi/3)=\kiss_d$ is the kissing number ($\kiss_2=6$, $\kiss_3=12$, $\kiss_4=24$, $\kiss_8=240$,
$\kiss_{24}=196560$).}

\begin{theorem}\label{L3:thm:KLgauge}
\added{In any dimension $d$}\orig{Notes: ``in high dimensions''.} and for $1<r<2$, the function
\[
  f_d(x)=\begin{cases} M\bigl(d,\arccos(1-\tfrac{2}{r^2})\bigr)^{-1}, & \abs{x}<r,\\[2pt] 0, & \abs{x}\ge r,\end{cases}
\]
is a Blichfeldt gauge for $\Ball^d$.
\end{theorem}

The key geometric fact is the following.

\begin{lemma}\label{L3:lem:angle}
Let $r<2$, let $x\in\R^d$ and let $y,z\in\R^d$ with $\abs{y-z}\ge 2$ and $\abs{x-y},\abs{x-z}\le r$. Then the angle $\theta$ at $x$
spanned by the segment $[y,z]$ satisfies $\cos\theta\le 1-\frac{2}{r^2}$, i.e.\ $\theta\ge\arccos(1-\frac2{r^2})$.
\end{lemma}

\begin{figure}[ht]
\centering
\begin{tikzpicture}[scale=1.5]
  \coordinate (X) at (0,0);
  \coordinate (Y) at (1.249,1.0);
  \coordinate (Z) at (1.249,-1.0);
  \draw[dashed,gray] (X) circle (1.6);
  \draw[thick] (X) -- (Y) node[midway,above left] {$r$};
  \draw[thick] (X) -- (Z) node[midway,below left] {$r$};
  \draw[thick] (Y) -- (Z) node[midway,right] {$2$};
  \draw (0.4,0.32) arc (38.66:-38.66:0.512);
  \node at (0.75,0) {$\theta$};
  \fill (X) circle (1.2pt) node[left] {$x$};
  \fill (Y) circle (1.2pt) node[above right] {$y$};
  \fill (Z) circle (1.2pt) node[below right] {$z$};
\end{tikzpicture}
\caption{The extremal case of Lemma~\ref{L3:lem:angle}: both centers at distance exactly $r$ from $x$ and at distance $2$ from each
other; then $\sin(\theta/2)=1/r$.}
\label{L3:fig:angle}
\end{figure}

\begin{proof}
First suppose both points are at distance exactly $r$ from $x$ (the worst case, see Figure~\ref{L3:fig:angle}).
Put $a=(y-x)/r$, $b=(z-x)/r$, unit vectors with $a\cdot b=\cos\theta$, and let $\ell=\abs{y-z}\ge2$. Then
\[
  \Bigl|\frac2r\Bigr|^2\le\Bigl|\frac{\ell}{r}\Bigr|^2=\abs{a-b}^2=2-2\,a\cdot b=2-2\cos\theta,
\]
so $\frac{4}{r^2}\le 2-2\cos\theta$, i.e.\ $\frac2{r^2}\le1-\cos\theta$, which is
\[
  \boxed{\cos\theta\le 1-\frac{2}{r^2}.}
\]
\begin{addedblock}
Why this is the worst case: in general, with $\alpha=\abs{x-y}$, $\beta=\abs{x-z}\in(0,r]$, the law of cosines gives
$\cos\theta=\frac{\alpha^2+\beta^2-\ell^2}{2\alpha\beta}\le h(\alpha,\beta):=\frac{\alpha^2+\beta^2-4}{2\alpha\beta}$, and
$\partial_\alpha h=\frac{\alpha^2-\beta^2+4}{2\alpha^2\beta}>0$ because $\beta\le r<2$. By symmetry $h$ is increasing in each variable,
so $h(\alpha,\beta)\le h(r,r)=1-\frac2{r^2}$. Equivalently, the apex angle of the isosceles triangle with legs $r$ and base $2$ is
$2\arcsin(1/r)$, and $\cos(2\arcsin\frac1r)=1-\frac2{r^2}$.
\end{addedblock}
\fixed{Side sketch says apex angle $2\arccos(1/r)$; it is $2\arcsin(1/r)$.}%
\end{proof}

\begin{proof}[Proof of Theorem~\ref{L3:thm:KLgauge}]
Take any packing of unit balls in $\R^d$ and $r<2$. For every $x\in\R^d$ there are at most $M\bigl(d,\arccos(1-2/r^2)\bigr)$ centers at
distance less than $r$ from $x$: by Lemma~\ref{L3:lem:angle}, the angle at $x$ spanned by any segment of length $\ge2$ whose ends are within
$r$ of $x$ is at least $\arccos(1-2/r^2)$, so the directions from $x$ to these centers form a configuration as in
Definition~\ref{L3:def:M}. \added{(If $x$ is itself a center, all other centers are at distance $\ge2>r$, and the count is $1$.)}
Hence $\sum_i f_d(x-c_i)\le 1$.
\end{proof}

\begin{corollary}\label{L3:cor:KL}
For $1\le r<2$,
\[
  \dens(\Ball^d)\le r^{-d}\,M\bigl(d,\arccos(1-\tfrac{2}{r^2})\bigr).
\]
Equivalently, substituting $r=1/\sin(\varphi/2)$,
\[
  \dens(\Ball^d)\le \sin\bigl(\tfrac\varphi2\bigr)^d\,M(d,\varphi),\qquad \frac\pi3<\varphi\le\pi.
\]
\end{corollary}
\begin{proof}
\added{$\int f_d=r^d\vol(\Ball^d)/M$, so $\dens\le\vol(\Ball^d)/\int f_d=r^{-d}M$. For the substitution: $1-\frac{2}{r^2}=1-2\sin^2\frac\varphi2=\cos\varphi$,
and $r\in[1,2)$ corresponds to $\varphi\in(\pi/3,\pi]$.}
\end{proof}
\added{Letting $r\to2$ recovers the simple bound $\dens(\Ball^d)\le \kiss_d\,2^{-d}$.}

For $d$ large, take $\varphi\approx63^\circ$; the Kabatjanskii--Levenshtein bound on spherical codes gives
\[
  \sin\bigl(\tfrac\varphi2\bigr)^d M(d,\varphi)\le 2^{-0.599d+o(d)},
\]
and hence
\[
  \dens(\Ball^d)\le 2^{-0.599d+o(d)}.
\]
\begin{addedblock}
More precisely, Kabatjanskii and Levenshtein (1978), using linear programming bounds for spherical codes, proved
$M(d,\varphi)\le 2^{d(c(\varphi)+o(1))}$ for $0<\varphi<\pi/2$, where with $s=\sin\varphi$,
\[
  c(\varphi)=\frac{1+s}{2s}\log_2\frac{1+s}{2s}-\frac{1-s}{2s}\log_2\frac{1-s}{2s}.
\]
The exponent $c(\varphi)+\log_2\sin(\varphi/2)$ is minimized at $\varphi^*\approx 62.9974^\circ$, where it equals $-0.5990\ldots$
(see Conway--Sloane, \emph{Sphere Packings, Lattices and Groups}, Ch.~1 and~9).
\end{addedblock}

%======================================================================
\section{Blichfeldt gauges for general bodies}
\label{L3:sec:general}
\orig{Section number overwritten, presumably 3.5.}

So we defined a Blichfeldt gauge for $\Ball^n$. Can we extend this? Abstractly:

\begin{definition}\label{L3:def:gauge}
Given a (convex) body $C$ in $\R^n$, a function $f:\R^n\to\R$ is a \emph{Blichfeldt gauge for $C$} if for any collection of
isometries $\{\varphi_i\}_{i=1}^\infty$ of $\R^n$ such that $\{\varphi_iC\}_{i=1}^\infty$ is a packing,
\[
  \sum_{i=1}^\infty f(\varphi_i^{-1}x)\le 1\qquad\forall x\in\R^n.
\]
\end{definition}
\added{As for balls, integrating over a large box shows that if $f$ is integrable and compactly supported, then every packing of
congruent copies of $C$ has density at most $\vol(C)/I(f)$, where $I(f):=\int_{\R^n}f$.}

\begin{remark}
One could add a non-negativity condition on $f$; what we really want is that the mass of $f$ is large: $I(f)\ge I(\chi_C)$.
\end{remark}
\begin{remark}
The index set could also be finite.
\end{remark}
\begin{remark}
The function $f$ could be much stranger than what we allow here.\unclear{Rest of the remark illegible: ``\dots accuracy first''?}
\end{remark}
\begin{exercise}[$\ast\ast\ast$]
Think about such functions distributionally.
\end{exercise}

There always exists a Blichfeldt gauge for $C$, e.g.\ $\chi_C$ \added{(the copies $\varphi_iC$ have disjoint interiors, so a point lies in
the interior of at most one of them; here one should take $\chi$ of the interior, or note that boundaries have measure zero)}.

\begin{lemma}\label{L3:lem:scaling}
If $f(x)$ is a Blichfeldt gauge for $C$, then $f(x/\lambda)$ is a Blichfeldt gauge for $\lambda C$, the $\lambda$-scaling of $C$.
\end{lemma}
\added{Indeed, a packing $\{\varphi_i(\lambda C)\}$ is the image under $x\mapsto\lambda x$ of the packing $\{\psi_iC\}$ with
$\psi_i(y)=\lambda^{-1}\varphi_i(\lambda y)$.}
For instance, $f_0$ (supported on the ball of radius $\sqrt2$) is a gauge for $\Ball^n$, and $f_0(x/\lambda)$ (supported on the ball of
radius $\lambda\sqrt2$) is a gauge for $\lambda\Ball^n$.\orig{Diagram of this scaling omitted.}

\begin{definition}\label{L3:def:inradius}
For a body $C$ define its \emph{inradius} $r(C)$ to be the radius of the largest sphere \added{(ball)} contained in $C$.
For $0\le\gamma\le r(C)$ define $C_\gamma$, the \emph{inner parallel body at depth $\gamma$}, to be the set of points
\[
  C_\gamma:=\{x\in C:\ x+\gamma\Ball^n\subseteq C\}.
\]
\end{definition}
\added{If $C$ is convex and compact, then $C_\gamma$ is convex, compact and non-empty for $\gamma\le r(C)$.}
We can define the straight-line distance from a point $x$ to a body $C$ by $d(x,C)$ \added{$=\min_{y\in C}\abs{x-y}$}.

\begin{figure}[ht]
\centering
\begin{tikzpicture}[scale=1.1]
  \draw[thick] (0,0) rectangle (4,2.4);
  \draw[dashed] (0.7,0.7) rectangle (3.3,1.7);
  \node at (2,1.2) {$C_\gamma$};
  \node at (0.35,2.1) {$C$};
  \draw[gray] (3.3,1.7) circle (0.7);
  \draw[->,gray] (3.3,1.7) -- ++(-150:0.7) node[midway,below,black] {\small$\gamma$};
  \fill (3.3,1.7) circle (1.2pt) node[below right] {$x_i$};
  \fill (6,3) circle (1.2pt) node[above right] {$x$};
  \draw[densely dotted] (6,3) -- (3.3,1.7) node[midway,above left] {$d(x,C_\gamma)$};
\end{tikzpicture}
\caption{The inner parallel body $C_\gamma$ and the nearest point $x_i\in C_\gamma$ to $x$; the ball $x_i+\gamma\Ball^n$ lies in $C$.}
\label{L3:fig:parallel}
\end{figure}

\begin{theorem}\label{L3:thm:parallel}
Let $C$ be a convex body, $0<\gamma\le r(C)$, and let $f$ be a radial Blichfeldt gauge for $\Ball^n$, i.e.\ $f(x)=F(\abs x)$.
Define
\[
  g(x)=F\Bigl(\frac{d(x,C_\gamma)}{\gamma}\Bigr).
\]
Then $g$ is a Blichfeldt gauge for $C$.
\end{theorem}

This is fairly easy to see from a picture (Figure~\ref{L3:fig:parallel}).

\begin{proof}
We would like to show that
\[
  \sum_{i=1}^\infty g(\varphi_i^{-1}x)\le1\qquad\forall x\in\R^n,\ \text{for all packings }\{\varphi_iC\}_{i=1}^\infty \text{ of } C.
\]
Consider a point $x\in\R^n$ and, for each $i$, a point $x_i\in\varphi_iC_\gamma$ at distance $d(x,\varphi_iC_\gamma)$ from $x$
\added{(it exists since $C_\gamma$ is compact)}. Then $x_i+\gamma\Ball^n$ is contained in $\varphi_iC$
\added{(note $\varphi_iC_\gamma=(\varphi_iC)_\gamma$)}, and so $\{x_i+\gamma\Ball^n\}_i$ is a packing.
\orig{Parenthetical: ``since $C$ is convex, things are well defined''.}%
Since $F(\abs{y}/\gamma)$ is a Blichfeldt gauge for $\gamma\Ball^n$ (Lemma~\ref{L3:lem:scaling}), and $x$ is at distance
$\abs{x-x_i}=d(x,\varphi_iC_\gamma)$ from $x_i$, we get
\[
  \sum_i g(\varphi_i^{-1}x)=\sum_i F\Bigl(\frac{d(x,\varphi_iC_\gamma)}{\gamma}\Bigr)=\sum_i F\Bigl(\frac{\abs{x-x_i}}{\gamma}\Bigr)\le 1,
\]
\added{using $d(\varphi_i^{-1}x,C_\gamma)=d(x,\varphi_iC_\gamma)$ because $\varphi_i$ is an isometry.}
\end{proof}

\subsection*{Is this useful?}
As with Blichfeldt gauges and other approaches: is this useful?\unclear{``is this useful?'' tentative; could be ``is there a catch?''} If we have a convex body like $A+\Ball^n$ \added{(a Minkowski sum
with a ball)}, it could be, provided there is enough high-dimensional mass.

\paragraph{A better gauge for balls.}
Consider
\[
  f(x)=\begin{cases} f_0(x), & \abs x>1,\\[2pt] 1-f_0\bigl(2-\abs{x}\bigr), & \abs x\le1,\end{cases}
\]
\added{with $f_0$ from \eqref{L3:eq:f0}, viewed as a function of the radius. Explicitly, $f=1$ for $\abs x\le2-\sqrt2$, $f=\frac12(2-\abs x)^2$ for
$2-\sqrt2\le\abs x\le1$, and $f=f_0$ beyond; $f$ is continuous and $f\ge f_0$ everywhere (Figure~\ref{L3:fig:profiles}).}
This leads to
\[
  \dens\le\Bigl(\frac{2}{d+2}\sqrt2^{\,d}\,(1+b_d)\Bigr)^{-1},\qquad
  b_d=\frac{1}{(\sqrt2)^d(d+1)}-(\sqrt2-1)^{d+1}\Bigl(1+\frac{\sqrt2}{d+1}\Bigr).
\]
\unclear{Notes do not verify that $f$ is a gauge; see the added paragraph below.}
\begin{addedblock}
The factor $\frac{2}{d+2}\sqrt2^{\,d}=\int f_0/\vol(\Ball^d)$ is Blichfeldt's, and $b_d=\int(f-f_0)\big/\int f_0$
(we checked this integral numerically against the formula for $d\le10$). The notes do not prove that $f$ is a Blichfeldt gauge.
The idea: a point $x$ lies in at most one ball of the packing, say at distance $s=\abs{x-c_0}\le1$ from its center, and all other centers
are at distance $\ge2-s$ from $x$; one has to show that their total $f_0$-contribution is at most $f_0(2-s)$, so that it is
compensated by replacing $f_0(s)$ with $1-f_0(2-s)$. Blichfeldt's inequality $\sum_i\abs{x-c_i}^2\ge 2(m-1)$ alone only gives the weaker
bound $s^2/2$ for this contribution, so a finer argument is needed. (A numerical search in dimensions $2,3,4$ found no violation.)
\end{addedblock}

\begin{figure}[ht]
\centering
\begin{tikzpicture}[xscale=3.4,yscale=2.2]
  \draw[->] (0,0) -- (1.7,0) node[right] {$\abs x$};
  \draw[->] (0,0) -- (0,1.2);
  \draw[thick,gray,dashed] plot[domain=0:1.41421,samples=60] (\x,{1-\x*\x/2});
  \draw[thick] (0,1) -- (0.58579,1);
  \draw[thick] plot[domain=0.58579:1,samples=40] (\x,{(2-\x)*(2-\x)/2});
  \draw[thick] plot[domain=1:1.41421,samples=40] (\x,{1-\x*\x/2});
  \draw[thick] (1.41421,0) -- (1.65,0);
  \foreach \t/\l in {0.58579/{2-\sqrt2},1/1,1.41421/{\sqrt2}} \draw (\t,0.03) -- (\t,-0.03) node[below] {\small$\l$};
  \node[left] at (0,1) {$1$};
  \node[gray] at (0.35,0.8) {\small$f_0$};
  \node at (0.85,0.95) {\small$f$};
\end{tikzpicture}
\caption{Radial profiles of Blichfeldt's gauge $f_0$ (dashed) and the improved function $f$ (solid); they agree for $\abs x\ge1$.}
\label{L3:fig:profiles}
\end{figure}

The resulting bounds compared with those from $f_0$ \added{(truncated to two digits)}:
\begin{center}
\begin{tabular}{c|cccccccccc}
\toprule
$d$ & 1 & 2 & 3 & 4 & 5 & 6 & 7 & 8 & 9 & 10\\
\midrule
bound from $f$ & 1 & .94 & .84 & .72 & .60 & .49 & .39 & .31 & .24 & .18\\
bound from $f_0$ & 1.06 & 1 & .88 & .75 & .61 & .5 & .39 & .31 & .24 & .19\\
\bottomrule
\end{tabular}
\end{center}
\orig{An arrow on p.~47 links $.94$ ($f$, $d=2$) with $1$ ($f_0$, $d=2$).}

\paragraph{Bodies with a lot of 2-dimensional mass.}
So this works for things with a lot of ``2-dimensional'' (or higher-dimensional) mass. For a long capsule in $\R^3$
(a segment plus a ball), the cross-sections are discs and Theorem~\ref{L3:thm:parallel} gives a bound of $.94+\text{error}$.
By contrast, for a long box \added{in the plane} the cross-sections are $1$-dimensional and we only get $1+\text{error}$: the method gives
nothing, since there is no nontrivial $1$-dimensional Blichfeldt gauge.
\unclear{Reading of the box sketch and of ``hence doesn't work as there is no 1d Blichfeldt gauge'' is tentative.}
\added{(In dimension $1$, intervals tile the line, so every gauge has $I(f)\le\vol(\Ball^1)$.)}

\begin{figure}[ht]
\centering
\begin{tikzpicture}[scale=0.9]
  \draw[thick] (0,1) -- (4,1) arc (90:-90:1) -- (0,-1) arc (270:90:1);
  \draw[very thick] (0,0) -- (4,0) node[midway,above] {$\ell$};
  \fill (0,0) circle (1.3pt) (4,0) circle (1.3pt);
  \draw[->,gray] (4,0) -- ++(-40:1) node[midway,below left,black] {\small$1$};
\end{tikzpicture}
\caption{An $\ell$-capsule: a segment of length $\ell$ plus a unit ball, drawn in a plane through the axis.}
\label{L3:fig:capsule}
\end{figure}

\begin{example}[$\ell$-capsules]\label{L3:ex:capsule}
Let $C=S_\ell+\Ball^d$, where $S_\ell$ is a segment of length $\ell$ \added{(the notes write ``$\ell\times\Ball^n$'')}.
\added{Then $r(C)=1$ and $C_1=S_\ell$. Take a radial gauge $F(\abs x)$ for $\Ball^d$ and $\gamma=1$, so $g(x)=F(d(x,S_\ell))$.
Write $I_k(F)=\int_{\R^k}F(\abs y)\,dy$. Points whose nearest point on $S_\ell$ is interior contribute a cylinder over $S_\ell$, and the two
ends together contribute a full ball's worth, so}
\[
  I(g)=\ell\cdot\frac{I_{d-1}(F)}{\vol(\Ball^{d-1})}\,\vol(\Ball^{d-1})+\frac{I_d(F)}{\vol(\Ball^d)}\,\vol(\Ball^d)
  =\ell\,I_{d-1}(F)+I_d(F),
\]
and
\[
  \vol(\ell\text{-capsule})=\ell\vol(\Ball^{d-1})+\vol(\Ball^d).
\]
Hence
\[
  \dens(\ell\text{-capsule})\le\frac{\ell\vol(\Ball^{d-1})+\vol(\Ball^d)}{\ell\,I_{d-1}(F)+I_d(F)}.
\]
\added{As $\ell\to\infty$ this tends to $\vol(\Ball^{d-1})/I_{d-1}(F)$, the $(d-1)$-dimensional ball bound for the same profile; for $d=3$ and
the improved $f$ above this is the value $.94$ from the table (assuming $f$ is a gauge in $\R^3$).}
\end{example}

\begin{addedblock}
This extension of Blichfeldt's method to ``longish'' bodies, such as long cylinders
$\Ball^{n-1}\times[0,h]$ and sausage-like solids $\Ball^n+[0,h]$, is due to G.~Fejes T\'oth and
W.~Kuperberg, \emph{Blichfeldt's density bound revisited}, Math.\ Ann.\ 295 (1993), 721--727. For
circular cylinders see also A.~Bezdek and W.~Kuperberg, Mathematika 37 (1990) (infinite cylinders:
density $\pi/\sqrt{12}$), and W.~Kusner, Discrete Comput.\ Geom.\ 51 (2014) (finite cylinders:
$\pi/\sqrt{12}+10/h$). Whether the improved radial gauge for the ball above is also taken from the 1993
paper (or from Blichfeldt 1929 or Rankin 1947) has not been checked against the original.
\end{addedblock}

\begin{remark}
\added{For a genuine cylinder $S_\ell\times\Ball^{d-1}$ (flat ends) the same computation works up to} end errors.
\end{remark}
% ===== end chapters/L03.tex =====

\part{Lattices and root systems}
% ===== begin chapters/L04a.tex =====
% MAT670 -- Lectures 4 and 5 (first part), transcribed from MAT670_46.pdf, pages 1--23
% Local drawing helpers for Dynkin diagrams (prefixed to avoid clashes).
% Vertices sit at (x,0); edges are drawn before vertices.
\newcommand{\LIVv}[1]{\filldraw[fill=white,draw=black] (#1) circle (0.09);}
\newcommand{\LIVe}[2]{\draw (#1) -- (#2);}
% double edge from (x,0) to (x+1,0), arrow pointing right (R) or left (L)
\newcommand{\LIVdR}[1]{\draw (#1,0.045) -- (#1+1,0.045); \draw (#1,-0.045) -- (#1+1,-0.045); \draw (#1+0.42,0.13) -- (#1+0.58,0) -- (#1+0.42,-0.13);}
\newcommand{\LIVdL}[1]{\draw (#1,0.045) -- (#1+1,0.045); \draw (#1,-0.045) -- (#1+1,-0.045); \draw (#1+0.58,0.13) -- (#1+0.42,0) -- (#1+0.58,-0.13);}
\newcommand{\LIVdN}[1]{\draw (#1,0.045) -- (#1+1,0.045); \draw (#1,-0.045) -- (#1+1,-0.045);}
% triple edge from (x,0) to (x+1,0), arrow pointing right / none
\newcommand{\LIVtR}[1]{\draw (#1,0.07) -- (#1+1,0.07); \draw (#1,0) -- (#1+1,0); \draw (#1,-0.07) -- (#1+1,-0.07); \draw (#1+0.4,0.16) -- (#1+0.6,0) -- (#1+0.4,-0.16);}
\newcommand{\LIVtN}[1]{\draw (#1,0.07) -- (#1+1,0.07); \draw (#1,0) -- (#1+1,0); \draw (#1,-0.07) -- (#1+1,-0.07);}

%======================================================================
\chapter[Lecture 4: Classification of root systems]{Lecture 4: The classification of (crystallographic) root systems}
\label{L04a:ch:L4}
\orig{Heading reads ``Lecture 4) (and 5?)''.}

%----------------------------------------------------------------------
\section{Motivation and lattice basics}
\label{L04a:sec:4.1}

In this lecture we examine \emph{root lattices}, coming from \emph{crystallographic} root systems.
One motivation is to find nice constructions of good packings in low dimensions. Another is to show
that such simple constructions do not really extend to high dimensions in a good way.
\orig{crossed out: ``there are not such simple constructions''}

The best known packings up through $\R^8$ are root lattices. In fact $A_1$, $A_2$, $A_3$ and $E_8$ are
optimal among \emph{all} packings in their dimensions.
\added{($A_1$: trivial; $A_2$: Thue, Fejes T\'oth 1940; $A_3$, the FCC packing: Hales 1998/2005, the Kepler
conjecture; $E_8$: Viazovska 2016. In dimension $24$ the Leech lattice, which is not a root lattice, is
optimal by Cohn--Kumar--Miller--Radchenko--Viazovska 2016.)}
\unclear{Box reads ``$A_1,A_2,A_3,E_8$ optimal as gen.\ packings''; ``$A_3$'' read from a blurred subscript.}

\begin{remark}
The classification we are about to give covers, in effect, half of the classification of complex
semisimple Lie algebras (equivalently, of complex semisimple Lie groups): the irreducible root systems
correspond to the simple ones.
\added{(See e.g.\ Humphreys, \emph{Introduction to Lie Algebras and Representation Theory}, Ch.~III, or
Fulton--Harris, \emph{Representation Theory}, Lecture~21.)}
\orig{Messy boxed aside, partly scribbled over; a book reference is unreadable.}
\end{remark}

\begin{definition}\label{L04a:def:lattice}
A \emph{lattice} $\Lat$ is a discrete additive subgroup of $\R^n$. For our purposes it will always be
co-compact, i.e.\ have full rank. Such a lattice has a $\Z$-basis $\{a_1,\dots,a_n\}$,
\added{$\Lat = \Z a_1\oplus\dots\oplus\Z a_n$,} that is at the same time an $\R$-basis of $\R^n$.
\end{definition}

\begin{example}
$\{(1,0)^T,(0,1)^T\}$ is a basis of the lattice $\Z^2\subset\R^2$.
\end{example}

A lattice can have many different bases. However, two bases $\{a_i\}$, $\{a_i'\}$ differ by an
\emph{integer} change of basis $K=(k_{ij})$ with integral inverse:
\[
  a_i' = \sum_{j=1}^n k_{ij} a_j, \qquad k_{ij}\in\Z
  \quad\Longrightarrow\quad \det K = \pm 1
\]
\added{(since $\det K$ and $\det K^{-1}$ are both integers whose product is $1$).}
So the volume of the torus $\R^n/\Lat$,
\[
  \vol(\R^n/\Lat) = \abs{\det\Lat} = \vol\Big(\Big\{\textstyle\sum_i t_i a_i : 0\le t_i\le 1\Big\}\Big),
\]
\added{i.e.\ the volume of the parallelepiped spanned by $a_1,\dots,a_n$, where $\det\Lat$ denotes the
determinant of the matrix with columns $a_1,\dots,a_n$,} is a clear invariant of $\Lat$.

The largest balls that can be placed at the lattice points without overlapping have radius
$\tfrac12\min_{\lambda\in\Lat\setminus\{0\}}\abs{\lambda}$, so each lattice gives a packing
\[
  \Big\{\, \varphi_i + \Ball\big(\tfrac12 \min_{\lambda\in\Lat\setminus\{0\}} \abs{\lambda}\big) \;:\; \varphi_i\in\Lat \Big\}.
\]
\begin{addedblock}
Its density is
$\dens(\Lat) = \vol \Ball^n\!\big(\tfrac12\lambda_1(\Lat)\big)\big/\abs{\det\Lat}$, where
$\lambda_1(\Lat)=\min_{\lambda\neq0}\abs{\lambda}$.
\end{addedblock}

Since scaling does not change the density, we can normalize to $\abs{\det\Lat}=1$ (i.e.\ consider
lattices $g\Z^n$ with $g\in\SL_n(\R)$).
\fixed{notes: ``(to $SL_n$, compact)''; the space is not compact}
\added{The space $\SL_n(\R)/\SL_n(\Z)$ of such lattices is \emph{not} compact (e.g.\ the lattices
$\operatorname{diag}(t,t^{-1},1,\dots,1)\Z^n$, $t\to0$, leave every compact set). What is true is
Mahler's compactness criterion: for each $\varepsilon>0$ the set of unimodular lattices with
$\lambda_1\ge\varepsilon$ is compact.}
Then there is clearly an upper bound on the length of the shortest vector among lattices with
$\abs{\det\Lat}=1$, and so a bound on the density.
\added{(Trivially, the balls are disjoint, so $\vol\Ball^n(\lambda_1/2)\le 1$; Minkowski's convex body
theorem gives the sharper $\vol\Ball^n(\lambda_1)\le 2^n$. Since lattices with very small $\lambda_1$
have small density, the supremum of the density can be taken over a set
$\{\lambda_1\ge\varepsilon\}$, which is compact by Mahler; so an optimal lattice exists in each
dimension.)}

So\dots\ what is the optimal lattice? The optimal lattice packings are known up to $n=8$
(Blichfeldt 1935 for $n=6,7,8$), and they are
\[
  A_1,\ A_2,\ A_3,\ D_4,\ D_5,\ E_6,\ E_7,\ E_8 .
\]
These are root lattices.
\orig{crossed out: ``still open''}
\added{(Earlier cases: Lagrange for $n=2$, Gauss for $n=3$, Korkine--Zolotarev for $n=4,5$. The only
other dimension where the optimal lattice is known is $n=24$: the Leech lattice, Cohn--Kumar 2009.)}

%----------------------------------------------------------------------
\section{Good candidates in low dimensions}
\label{L04a:sec:4.2}

\begin{definition}\label{L04a:def:integral}
A lattice $\Lat$ is \emph{integral} if $\ip{x}{y}\in\Z$ for all $x,y\in\Lat$.
A \emph{root lattice} is an integral lattice generated by roots, i.e.\ by elements of a root system
(Definition~\ref{L04a:def:rootsystem}).
\end{definition}

\begin{addedblock}
In lattice language one usually calls a vector $v$ of an integral lattice a root if
$\ip{v}{v}\in\{1,2\}$; then $2\ip{x}{v}/\ip{v}{v}\in\Z$ for all $x\in\Lat$, so the reflection $s_v$
(below) maps $\Lat$ to itself, and the roots of $\Lat$ form a root system. The irreducible lattices
generated by vectors of norm $2$ are exactly $A_n$ ($n\ge1$), $D_n$ ($n\ge4$), $E_6,E_7,E_8$ (Witt);
see Conway--Sloane, \emph{Sphere Packings, Lattices and Groups}, Ch.~4.
\end{addedblock}

\begin{definition}\label{L04a:def:rootsystem}
A \emph{root system} is a subset $\Delta\subset\R^n\setminus\{0\}$ satisfying:
\begin{enumerate}[label=\arabic*)]
\item $\Delta$ is finite and spans $\R^n$;
\item if $\alpha\in\Delta$, then $n\alpha\in\Delta \iff n=\pm1$ \quad(\emph{reduced});
\item $\Delta$ is invariant under the orthogonal reflection in the hyperplane $\alpha^\perp$, for any
  $\alpha\in\Delta$: writing
  \[
    \operatorname{proj}_\alpha\beta = \alpha\,\frac{\ip{\beta}{\alpha}}{\ip{\alpha}{\alpha}},\qquad
    s_\alpha\beta := \beta - 2\operatorname{proj}_\alpha\beta ,
  \]
  we require $s_\alpha\beta\in\Delta$ for all $\alpha,\beta\in\Delta$;
\item (\emph{crystallographic}) for all $\alpha,\beta\in\Delta$,
  \[
    n_{\beta\alpha} := \frac{2\ip{\beta}{\alpha}}{\ip{\alpha}{\alpha}}\in\Z .
  \]
\end{enumerate}
\added{The elements of $\Delta$ are called \emph{roots}.}
\end{definition}

\begin{center}
\begin{tikzpicture}[scale=1.1,>=Stealth]
  \draw[dashed,gray] (0,-0.4) -- (0,1.7) node[above] {\small$\alpha^\perp$};
  \draw[->,thick] (0,0) -- (2,0) node[right] {$\alpha$};
  \draw[->,thick] (0,0) -- (1.2,1.2) node[above right] {$\beta$};
  \draw[dotted] (1.2,1.2) -- (1.2,0);
  \draw[->,very thick,blue!60!black] (0,0) -- (1.2,0);
  \node[below] at (1.2,-0.05) {\small$\operatorname{proj}_\alpha\beta$};
  \draw[->,thick] (0,0) -- (-1.2,1.2) node[above left] {$s_\alpha\beta$};
  \draw[dotted] (-1.2,1.2) -- (1.2,1.2);
\end{tikzpicture}
\end{center}

Property 4) is the restriction that the projection of $\beta$ onto $\alpha$ is an integer or
half-integer multiple of $\alpha$; but it really forces $\Delta$ to give rise to a lattice
\added{(the $\Z$-span of $\Delta$ is discrete: it is the $\Z$-span of a fundamental system, which is a
basis of $\R^n$, see \S\ref{L04a:sec:4.3}).}

\begin{exercise}
Try building some other root systems, i.e.\ sets satisfying 1)--3), which are \emph{not}
crystallographic.
\added{(Hint: the $10$ unit vectors at angles $k\pi/5$ in $\R^2$; the $30$ vertices of an
icosidodecahedron in $\R^3$, i.e.\ type $H_3$.)}
\end{exercise}

\paragraph{Property 4) is strong.}
We have
\[
  n_{\beta\alpha} = \frac{2\ip{\beta}{\alpha}}{\ip{\alpha}{\alpha}} = \frac{2\abs{\beta}}{\abs{\alpha}}\cos\theta,
\]
where $\theta$ is the angle between $\alpha$ and $\beta$. Since both $n_{\beta\alpha}$ and $n_{\alpha\beta}$ are integers,
\[
  \boxed{\,n_{\beta\alpha}\,n_{\alpha\beta} = 4\cos^2\theta\in\Z\,}
  \qquad\Longrightarrow\qquad 4\cos^2\theta\in\{0,1,2,3,4\}.
\]
If $4\cos^2\theta=4$, then $\alpha$ and $\beta$ are parallel, so $\alpha=\beta$ or $\alpha=-\beta$
\added{by the reducedness axiom 2)}.

Assuming $\abs{\alpha}\ge\abs{\beta}$ (and $\alpha\neq\pm\beta$), this leaves the following table.
\begin{center}
\renewcommand{\arraystretch}{1.25}
\begin{tabular}{cccccc}
\toprule
$4\cos^2\theta$ & $n_{\beta\alpha}$ & $n_{\alpha\beta}$ & $\abs{\alpha}/\abs{\beta}$ & $\cos\theta$ & $\theta$\\
\midrule
$3$ & $+1$ & $+3$ & $\sqrt3$ & $+\sqrt3/2$ & $\pi/6$\\
$3$ & $-1$ & $-3$ & $\sqrt3$ & $-\sqrt3/2$ & $5\pi/6$\\
$2$ & $+1$ & $+2$ & $\sqrt2$ & $+\sqrt2/2$ & $\pi/4$\\
$2$ & $-1$ & $-2$ & $\sqrt2$ & $-\sqrt2/2$ & $3\pi/4$\\
$1$ & $+1$ & $+1$ & $1$ & $+1/2$ & $\pi/3$\\
$1$ & $-1$ & $-1$ & $1$ & $-1/2$ & $2\pi/3$\\
$0$ & $0$ & $0$ & (undetermined) & $0$ & $\pi/2$\\
\bottomrule
\end{tabular}
\end{center}
\added{(For $4\cos^2\theta\neq0$ the ratio follows from $n_{\alpha\beta}/n_{\beta\alpha}=\abs{\alpha}^2/\abs{\beta}^2$.)}

\begin{definition}\label{L04a:def:reducible}
A root system is \emph{reducible} (\emph{decomposable}) if $\Delta=\Delta_1\cup\Delta_2$
\added{with $\Delta_1,\Delta_2$ nonempty} such that $\ip{\alpha_1}{\alpha_2}=0$ for all
$\alpha_1\in\Delta_1$, $\alpha_2\in\Delta_2$; i.e.\ all the elements of one part are orthogonal to all
the elements of the other. Otherwise we call it \emph{irreducible} (\emph{indecomposable}).
\end{definition}

\begin{example}[Root systems in dimensions $1$ and $2$]\label{L04a:ex:rank2}
\leavevmode
\begin{center}
\begin{tikzpicture}[>=Stealth,scale=0.95]
  % A1
  \begin{scope}
    \draw[->,thick] (0,0) -- (1,0) node[right] {\small$\alpha$};
    \draw[->,thick] (0,0) -- (-1,0) node[left] {\small$-\alpha$};
    \fill (0,0) circle (1.2pt);
    \node at (0,-1.75) {$A_1$\quad($n=1$)};
  \end{scope}
  % A1 x A1
  \begin{scope}[xshift=4cm]
    \draw[->,thick] (0,0) -- (1,0) node[right] {\small$\alpha$};
    \draw[->,thick] (0,0) -- (-1,0) node[left] {\small$-\alpha$};
    \draw[->,thick] (0,0) -- (0,0.7) node[above] {\small$\beta$};
    \draw[->,thick] (0,0) -- (0,-0.7) node[below] {\small$-\beta$};
    \node at (0,-1.75) {$A_1\times A_1$\ (reducible), $\theta=\frac\pi2$};
  \end{scope}
  % A2
  \begin{scope}[xshift=9.4cm]
    \draw[gray!60] (0:1)--(60:1)--(120:1)--(180:1)--(240:1)--(300:1)--cycle;
    \foreach \a in {0,60,...,300} \draw[->,thick] (0,0) -- (\a:1);
    \node[right] at (0:1) {\small$\alpha$};
    \node[above left] at (120:1) {\small$\beta$};
    \node[above right] at (60:1) {\small$\alpha+\beta=s_\alpha\beta$};
    \node[left] at (180:1) {\small$-\alpha$};
    \node[below left] at (240:1) {\small$s_\alpha(-\beta)$};
    \node[below right] at (300:1) {\small$-\beta$};
    \node at (0,-1.75) {$A_2$, $\theta=\frac\pi3$};
  \end{scope}
\end{tikzpicture}

\medskip
\begin{tikzpicture}[>=Stealth,scale=0.95]
  % B2
  \begin{scope}
    \draw[gray!60] (-1,-1) rectangle (1,1);
    \foreach \v in {(1,0),(0,1),(-1,0),(0,-1),(1,1),(-1,1),(-1,-1),(1,-1)} \draw[->,thick] (0,0) -- \v;
    \node[right] at (1,0) {\small$\alpha$};
    \node[above left] at (-1,1) {\small$\beta$};
    \node[above] at (0,1) {\small$s_\beta\alpha$};
    \node at (0,-1.5) {$B_2$, $\theta=\frac\pi4$};
  \end{scope}
  % G2
  \begin{scope}[xshift=5cm]
    \foreach \a in {0,60,...,300} \draw[->,thick] (0,0) -- (\a:0.7);
    \foreach \a in {30,90,...,330} \draw[->,thick] (0,0) -- (\a:1.212);
    \draw[gray!60] (30:1.212)--(150:1.212)--(270:1.212)--cycle;
    \draw[gray!60] (90:1.212)--(210:1.212)--(330:1.212)--cycle;
    \node at (0,-1.5) {$G_2$, $\theta=\frac\pi6$};
  \end{scope}
\end{tikzpicture}
\end{center}
(Note: $\Delta$ is not a basis!)
\added{In $B_2$ the roots are $\pm e_1,\pm e_2,\pm e_1\pm e_2$; in $G_2$ there are $6$ short roots and
$6$ long roots of length ratio $\sqrt3$. In each case $\theta$ is the smallest angle between two roots.}
\end{example}

%----------------------------------------------------------------------
\section{Classification of root systems: preliminaries}
\label{L04a:sec:4.3}

From a root system $\Delta$ we can choose a (non-unique) subset, the \emph{simple roots}.

For each $\alpha\in\Delta$ consider $\alpha^\perp$, its (linear) orthogonal complement.
\orig{crossed out: ``there is a unique $\perp$ hyperplane''}
Since $\abs{\Delta}$ is finite, there exists $d\in\R^n$ such that $\ip{\alpha}{d}\neq0$ for all
$\alpha\in\Delta$ \added{(avoid the finitely many hyperplanes $\alpha^\perp$)}. Hence there is a
partition of $\Delta$ into
\[
  \text{positive roots } \Delta_d^+ = \{\alpha\in\Delta : \ip{\alpha}{d}>0\},\qquad
  \text{negative roots } \Delta_d^- = \{\alpha\in\Delta : \ip{\alpha}{d}<0\}
\]
\added{(note $\Delta_d^-=-\Delta_d^+$).}

\begin{definition}\label{L04a:def:simple}
A root $\alpha\in\Delta_d^+$ is \emph{simple} if it is not the sum of two other positive roots.
A set of simple roots is a \emph{fundamental system} of $\Delta$ \added{(we write $\Pi=\Pi_d$)}.
\end{definition}

\begin{center}
\begin{tikzpicture}[>=Stealth,scale=1.3]
  \fill[gray!12] (165:1.7) -- (-15:1.7) arc (-15:165:1.7) -- cycle;
  \draw[dashed,red!70!black] (165:1.8) -- (-15:1.8) node[right] {\small$d^\perp$};
  \foreach \a in {0,60,...,300} \draw[->,thick] (0,0) -- (\a:1.2);
  \node[right] at (0:1.2) {\small$\alpha$};
  \node[above left] at (120:1.2) {\small$\beta$};
  \node[above right] at (60:1.2) {\small$\alpha+\beta$};
  \draw[->,red!70!black,thick] (0,0) -- (75:0.9) node[above] {\small$d$};
  \node at (35:1.75) {\small$\Delta_d^+$};
  \node at (-110:1.4) {\small$\Delta_d^-$};
\end{tikzpicture}
\end{center}
\orig{Sketch on p.\,8, redrawn for $A_2$: simple roots $\alpha,\beta$.}

The choice of $d$ affects the fundamental system, but such systems are equivalent under the action of
the group of reflections $\langle s_\alpha\rangle_{\alpha\in\Delta}$
\added{(the \emph{Weyl group} $W$; see Proposition~\ref{L04a:prop:reconstruct}).}

\begin{theorem}\label{L04a:thm:basis}
A fundamental system of $\Delta$ is an $\R$-basis for $\R^n$.
\added{Moreover every positive root is a linear combination of simple roots with nonnegative integer
coefficients.}
\end{theorem}

\begin{proof}[Sketch]
\emph{Linear independence.} Suppose $\sum_i c_i\gamma_i = 0$ with $\gamma_i\in\Pi$. Partition the
indices into those with $c_i>0$ $(+)$, those with $c_i<0$ $(-)$, and those with $c_i=0$, and put
\added{$x := \sum_{c_i>0} c_i\gamma_i = \sum_{c_j<0}\abs{c_j}\gamma_j$.} Then
\[
  \ip{x}{x} = \sum_{c_i>0,\ c_j<0} c_i\,\abs{c_j}\,\ip{\gamma_i}{\gamma_j}\le 0
\]
by Lemma~\ref{L04a:lem:obtuse} (distinct simple roots make non-acute angles), so $x=0$.
\added{But then $0=\ip{x}{d}=\sum_{c_i>0}c_i\ip{\gamma_i}{d}$ with every $\ip{\gamma_i}{d}>0$, so no
$c_i$ is positive; likewise none is negative.}

\emph{Span.} Let $\gamma$ be orthogonal to all simple roots. Then by the orbit argument
(Proposition~\ref{L04a:prop:reconstruct}: every root is the image of a simple root under the group
generated by the simple reflections $s_\alpha$, $\alpha\in\Pi$, each of which fixes $\gamma$),
$\gamma$ is orthogonal to all roots, hence $\gamma=0$ since $\Delta$ spans.
\begin{addedblock}
A more elementary route: if $\alpha\in\Delta_d^+$ is not simple, then $\alpha=\beta+\beta'$ with
$\beta,\beta'\in\Delta_d^+$, and $\ip{\beta}{d},\ip{\beta'}{d}<\ip{\alpha}{d}$. Since $\Delta$ is
finite, induction on $\ip{\alpha}{d}$ shows every positive root is a sum of simple roots. Hence $\Pi$
spans $\operatorname{span}\Delta=\R^n$, and the integrality claim follows too.
\end{addedblock}
\end{proof}

\begin{proposition}\label{L04a:prop:reconstruct}
$\Delta$ can be reconstructed from a fundamental system via reflections.
\added{Precisely: if $W_\Pi$ denotes the group generated by $s_\gamma$, $\gamma\in\Pi$, then
$\Delta = W_\Pi\cdot\Pi$, and any two fundamental systems are related by an element of $W_\Pi$.}
\end{proposition}

\begin{proof}[Sketch]
Consider the \emph{Weyl chamber} of the fundamental system $\Pi$,
\[
  C(\Pi)=\{x\in\R^n : \ip{x}{\gamma}>0 \text{ for all }\gamma\in\Pi\}.
\]
\orig{crossed out: ``orbit $d$ to each\dots''}
Then $d\in C(\Pi)$, and moving $d$ within its chamber does not change the fundamental system.
Any other fundamental system is based on some $d'$ in another chamber, and there are finitely many
chambers \added{(the connected components of $\R^n\setminus\bigcup_{\alpha}\alpha^\perp$)}. Take a path
$d\to d'$ \added{in general position}; the walls it crosses give reflections $s_{\text{walls}}$, whose
composition takes $d$ to (the chamber of) $d'$, and hence takes the fundamental system of $d$ to that
of $d'$.
\begin{addedblock}
Details: if $d$ and $d'=s_\alpha d$ lie in adjacent chambers separated by the wall $\alpha^\perp$,
then $\Pi_{d'}=s_\alpha\Pi_d$, since $s_\alpha$ is an isometry permuting $\Delta$ with
$\ip{s_\alpha\beta}{s_\alpha d}=\ip{\beta}{d}$. The walls of $C(\Pi)$ are exactly the hyperplanes
$\gamma^\perp$, $\gamma\in\Pi$; writing each wall-crossing reflection in terms of the previous ones
($s_{w\gamma}=ws_\gamma w^{-1}$) shows the composite lies in $W_\Pi$. Finally every root $\alpha$ is
simple for a suitable $d'$ (take $d'$ very close to a generic point of $\alpha^\perp$, on the side where
$\ip{\alpha}{d'}>0$), so $\alpha=w\gamma$ for some $w\in W_\Pi$, $\gamma\in\Pi$. See Humphreys,
\emph{Introduction to Lie Algebras and Representation Theory}, \S10.3.
\end{addedblock}
\end{proof}

\begin{center}
\begin{tikzpicture}[>=Stealth,scale=1.6]
  % hyperplanes alpha^perp (90), beta^perp (30), (alpha+beta)^perp (150)
  \draw[magenta,thick] (90:1.5) -- (270:1.5) node[below] {\small$\alpha^\perp$};
  \draw[magenta,thick] (30:1.5) -- (210:1.5) node[below left] {\small$\beta^\perp$};
  \draw[magenta,thick] (150:1.5) -- (330:1.5) node[below right] {\small$(\alpha+\beta)^\perp$};
  \foreach \a in {0,60,...,300} \draw[->,thick] (0,0) -- (\a:1.15);
  \node[right] at (0:1.15) {\small$\alpha$};
  \node[left] at (180:1.15) {\small$s_\alpha(\alpha)$};
  \node[above left] at (120:1.15) {\small$\beta=s_\alpha(\alpha+\beta)$};
  \node[above right] at (60:1.15) {\small$\alpha+\beta=s_\alpha(\beta)$};
  \fill[red!70!black] (75:0.8) circle (1pt) node[right] {\small$d$};
  \fill[green!45!black] (105:0.8) circle (1pt) node[left] {\small$s_\alpha d$};
  \draw[dashed] (75:0.8) to[bend right=15] (105:0.8);
\end{tikzpicture}
\end{center}
\orig{Colour sketch on p.\,11, redrawn.}
For $A_2$: with $d$ as shown, $\Pi_d=\{\alpha,\beta\}$; reflecting across $\alpha^\perp$ gives
$s_\alpha d$, with fundamental system $\Pi_{s_\alpha d}=\{s_\alpha(\alpha),s_\alpha(\beta)\}=\{-\alpha,\alpha+\beta\}$.

\begin{lemma}\label{L04a:lem:diff}
If $\alpha,\beta\in\Delta$ are not parallel and $\ip{\alpha}{\beta}>0$, then $\alpha-\beta\in\Delta$.
\end{lemma}

\begin{proof}
We have $n_{\alpha\beta}=2\ip{\alpha}{\beta}/\ip{\beta}{\beta}>0$, and by the table
\added{($n_{\alpha\beta}n_{\beta\alpha}\in\{1,2,3\}$ with both factors positive integers)}
$n_{\alpha\beta}=1$ or $n_{\beta\alpha}=1$. If $n_{\alpha\beta}=1$, then
\[
  s_\beta(\alpha) = \alpha - 2\beta\,\frac{\ip{\alpha}{\beta}}{\ip{\beta}{\beta}}
  = \alpha-n_{\alpha\beta}\,\beta = \alpha-\beta\in\Delta.
\]
Similarly, if $n_{\beta\alpha}=1$, then $s_\alpha(-\beta)=-\beta+n_{\beta\alpha}\alpha=\alpha-\beta\in\Delta$.
\fixed{written $m_{\beta\alpha}$, $s_\beta(-\beta)$ (typos)}
\end{proof}

\begin{lemma}\label{L04a:lem:obtuse}
If $\alpha,\beta$ are distinct simple roots, then $\ip{\alpha}{\beta}\le0$.
\end{lemma}

\begin{proof}
If $\ip{\alpha}{\beta}>0$, then \added{(as $\alpha\neq\pm\beta$, both being positive)}
Lemma~\ref{L04a:lem:diff} gives $\gamma:=\alpha-\beta\in\Delta$. So $\gamma$ or $-\gamma$ lies in
$\Delta_d^+$, and then $\alpha=\beta+\gamma$ or $\beta=\alpha+(-\gamma)$ is a sum of two positive roots,
contradicting simplicity.
\end{proof}

\begin{proposition}\label{L04a:prop:decomp}
A fundamental system is decomposable \added{(into two nonempty mutually orthogonal subsets)} if and only
if the root system is decomposable.
\end{proposition}

\begin{proof}[Sketch]
Similar arguments.
\begin{addedblock}
If $\Delta=\Delta_1\cup\Delta_2$ orthogonally, then $\Pi=(\Pi\cap\Delta_1)\cup(\Pi\cap\Delta_2)$, and
both parts are nonempty because $\Pi$ spans $\R^n$. Conversely, if $\Pi=\Pi_1\cup\Pi_2$ orthogonally,
the reflections $s_\gamma$, $\gamma\in\Pi_1$, fix $\Pi_2$ pointwise and preserve
$\operatorname{span}\Pi_1$, and vice versa; so by
Proposition~\ref{L04a:prop:reconstruct} every root lies in $W_\Pi\Pi_1\subset\operatorname{span}\Pi_1$ or
in $W_\Pi\Pi_2\subset\operatorname{span}\Pi_2$, which are orthogonal.
\end{addedblock}
\end{proof}

%----------------------------------------------------------------------
\section{Classification via Coxeter graphs}
\label{L04a:sec:4.4}

Simple roots determine root systems (Proposition~\ref{L04a:prop:reconstruct}), so it is enough to
classify fundamental systems.

Elements $\alpha,\beta$ of a fundamental system are orthogonal or make an obtuse angle
(Lemma~\ref{L04a:lem:obtuse}). By the table, therefore,
\[
  \theta \in \Big\{\frac\pi2,\ \frac{2\pi}3,\ \frac{3\pi}4,\ \frac{5\pi}6\Big\},
  \qquad\text{that is,}\qquad n_{\alpha\beta}n_{\beta\alpha}=4\cos^2\theta\in\{0,1,2,3\}.
\]

\begin{definition}\label{L04a:def:coxeter}
The \emph{Coxeter graph} of $\Delta$ has a vertex for each simple root, and an edge of weight
$n_{\alpha\beta}\cdot n_{\beta\alpha}$ between $\alpha$ and $\beta$ \added{(drawn as that many edges;
no edge when $\alpha\perp\beta$)}.
\end{definition}

It is well defined \added{(independent of the choice of fundamental system)} by the Weyl chamber
argument, since fundamental systems are carried to fundamental systems by the reflections
$\langle s_\alpha\rangle$, \added{which are isometries and so preserve all the angles}.

\begin{definition}\label{L04a:def:dynkin}
The \emph{Dynkin diagram} is the Coxeter graph with arrows on the double and triple edges, pointing to
the shorter vector.
\end{definition}

\begin{remark}
In an irreducible system, the possible lengths are determined by the angles
\added{(by the table, the ratio of the lengths of two non-orthogonal roots is fixed by their angle;
connecting through the diagram fixes all lengths up to one overall scale)}. For a reducible system the
components can be rescaled independently.
\unclear{Margin note by $A_1\times A_1$, probably ``unrelated scale!''.}
\end{remark}

\begin{example}
\leavevmode
\begin{center}
\begin{tabular}{r@{\quad$\longrightarrow$\quad}l}
$A_1$ & \tikz[baseline=-3pt]{\LIVv{0,0}} \\[4pt]
$A_1\times A_1$ & \tikz[baseline=-3pt]{\LIVv{0,0}\LIVv{1,0}} \\[4pt]
$A_2$ & \tikz[baseline=-3pt]{\LIVe{0,0}{1,0}\LIVv{0,0}\LIVv{1,0}} \\[4pt]
$B_2$ & \tikz[baseline=-3pt]{\LIVdR{0}\LIVv{0,0}\LIVv{1,0}} \\[4pt]
$G_2$ & \tikz[baseline=-3pt]{\LIVtR{0}\LIVv{0,0}\LIVv{1,0}}
\end{tabular}
\end{center}
\end{example}

\orig{End of Lecture 4 (p.\,15).}

%======================================================================
\chapter[Lecture 5: Root systems (continued)]{Lecture 5: Classification of crystallographic root systems (continued)}
\label{L04a:ch:L5}

\paragraph{Recall.} A root system $\Delta\subset\R^n\setminus\{0\}$ satisfies
\begin{enumerate}[label=\arabic*)]
\item $\Delta$ is finite and spans $\R^n$;
\item $\alpha\in\Delta\implies(n\alpha\in\Delta\iff n=\pm1)$;
\item $\Delta$ is invariant under the reflections in the hyperplanes $\alpha^\perp$:
  $s_\alpha(\beta)\in\Delta$ for all $\alpha,\beta\in\Delta$;
\item $n_{\beta\alpha}:=2\ip{\beta}{\alpha}/\ip{\alpha}{\alpha}\in\Z$,
\end{enumerate}
and 4) implies $4\cos^2\theta_{\alpha\beta}\in\Z$, really $\in\{0,1,2,3,4\}$.

The plan: root lattices $\to$ root systems $\longleftrightarrow$ simple roots (last time)
$\longleftrightarrow$ Dynkin diagrams (today).

From \S\ref{L04a:sec:4.4}: the \emph{Coxeter graph} has a vertex for each simple root and an edge of
weight $n_{\alpha\beta}\cdot n_{\beta\alpha}$ between $\alpha$ and $\beta$. The \emph{Dynkin diagram} is
the Coxeter graph with arrows on the multi-edges pointing at the shorter root.
\orig{Page has ``(longer)'' with ``shorter'' written below.}

%----------------------------------------------------------------------
\section{Classification of Coxeter graphs}
\label{L04a:sec:5.1}

We ignore lengths.

\begin{definition}\label{L04a:def:admissible}
A linearly independent set of $n$ unit vectors $\{v_1,\dots,v_n\}$ spanning $\R^n$ is \emph{admissible}
if for all $i\neq j$
\[
  \ip{v_i}{v_j}\le0 \qquad\text{and}\qquad 4\ip{v_i}{v_j}^2 = 4\cos^2\theta_{ij}\in\{0,1,2,3\}.
\]
An \emph{admissible diagram} is the Coxeter graph of an admissible set
\added{(vertices $v_i$, with $4\ip{v_i}{v_j}^2$ edges between $v_i$ and $v_j$)}.
\end{definition}

\begin{center}
\fbox{\parbox{0.8\linewidth}{\centering A normalized set of simple roots is admissible.\\
(Note: the angles are the determining factor.)}}
\end{center}
\added{This follows from Theorem~\ref{L04a:thm:basis}, Lemma~\ref{L04a:lem:obtuse} and the table.}

The simple roots of an irreducible $\Delta$ are not all orthogonal
\added{--- more precisely, they cannot be split into two mutually orthogonal sets
(Proposition~\ref{L04a:prop:decomp})} --- so the Coxeter graph is connected. So we only need to classify
connected admissible diagrams.

\begin{theorem}\label{L04a:thm:classification}
The Dynkin diagram of an irreducible (crystallographic) root system is of one of the following types
($n$ = number of vertices).

\smallskip
\noindent\emph{Infinite families:}
\begin{center}
\begin{tabular}{lll}
$A_n$ & \tikz[baseline=-3pt]{\LIVe{0,0}{1,0}\draw[dotted] (1,0)--(2,0);\LIVe{2,0}{3,0}\foreach\x in{0,1,2,3}{\LIVv{\x,0}}} & $n\ge1$\\[6pt]
$B_n$ & \tikz[baseline=-3pt]{\LIVe{0,0}{1,0}\draw[dotted] (1,0)--(2,0);\LIVdR{2}\foreach\x in{0,1,2,3}{\LIVv{\x,0}}} & $n\ge2$\\[6pt]
$C_n$ & \tikz[baseline=-3pt]{\LIVe{0,0}{1,0}\draw[dotted] (1,0)--(2,0);\LIVdL{2}\foreach\x in{0,1,2,3}{\LIVv{\x,0}}} & $n\ge3$\\[6pt]
$D_n$ & \tikz[baseline=-3pt]{\LIVe{0,0}{1,0}\draw[dotted] (1,0)--(2,0);\LIVe{2,0}{2.8,0.45}\LIVe{2,0}{2.8,-0.45}\foreach\x in{0,1,2}{\LIVv{\x,0}}\LIVv{2.8,0.45}\LIVv{2.8,-0.45}} & $n\ge4$
\end{tabular}
\end{center}

\noindent\emph{Exceptional types:}
\begin{center}
\begin{tabular}{ll}
$E_6$ & \tikz[baseline=-3pt]{\foreach\x in{0,1,2,3}{\LIVe{\x,0}{\x+1,0}}\LIVe{2,0}{2,0.7}\foreach\x in{0,...,4}{\LIVv{\x,0}}\LIVv{2,0.7}}\\[12pt]
$E_7$ & \tikz[baseline=-3pt]{\foreach\x in{0,...,4}{\LIVe{\x,0}{\x+1,0}}\LIVe{2,0}{2,0.7}\foreach\x in{0,...,5}{\LIVv{\x,0}}\LIVv{2,0.7}}\\[12pt]
$E_8$ & \tikz[baseline=-3pt]{\foreach\x in{0,...,5}{\LIVe{\x,0}{\x+1,0}}\LIVe{2,0}{2,0.7}\foreach\x in{0,...,6}{\LIVv{\x,0}}\LIVv{2,0.7}}\\[6pt]
$F_4$ & \tikz[baseline=-3pt]{\LIVe{0,0}{1,0}\LIVdR{1}\LIVe{2,0}{3,0}\foreach\x in{0,...,3}{\LIVv{\x,0}}}\\[6pt]
$G_2$ & \tikz[baseline=-3pt]{\LIVtR{0}\LIVv{0,0}\LIVv{1,0}}
\end{tabular}
\end{center}
\added{Conversely, each of these diagrams is the Dynkin diagram of an irreducible root system, which is
unique up to isomorphism.}
\unclear{Parenthetical under $D_n$ reads roughly ``(each also unique)''.}
\added{(The ranges of $n$ avoid repetitions: $B_2=C_2$, $D_3=A_3$.)}
\end{theorem}

It will suffice to classify admissible diagrams. By definition, subsets of the vectors of an admissible
set satisfy the admissibility condition inside their span (the resulting diagram could be
disconnected); \added{that is, any subdiagram of an admissible diagram, obtained by deleting vertices,
is admissible.}

\begin{lemma}\label{L04a:lem:tree}
A connected admissible diagram is a tree.
\end{lemma}

\begin{proof}
Consider $v=\sum_{i=1}^n v_i$, where $\{v_i\}$ is admissible. The $v_i$ are linearly independent, so
$v\neq0$, and
\[
  0<\ip{v}{v} = \sum_i\ip{v_i}{v_i} + \sum_{i<j}2\ip{v_i}{v_j} = n+\sum_{i<j}2\ip{v_i}{v_j}.
\]
If $v_i$ is connected to $v_j$ in the Coxeter graph, then by the table
$2\ip{v_i}{v_j}\in\{-1,-\sqrt2,-\sqrt3\}$\added{, in particular $2\ip{v_i}{v_j}\le-1$}. Hence
$\sum_{i<j}2\ip{v_i}{v_j}$ has at most $n-1$ terms not equal to $0$, i.e.\ there are at most $n-1$
connected pairs of vertices. But by assumption the diagram is connected, so there are exactly $n-1$
pairs of connected vertices \added{(a connected graph on $n$ vertices has at least $n-1$ edges)}, and
it is a tree (maybe with multiple edges, but not too many).
\orig{a scribbled-out remark at the bottom of p.\,20}
\end{proof}

\begin{lemma}\label{L04a:lem:degree}
The degree of each vertex, counted with multiplicity, is at most $3$.
\end{lemma}

\begin{proof}
Fix a vertex $v$, connected to the vertices $v_1,\dots,v_k$.
\fixed{written ``fix vertex $v_1$'', then used as $v$}
Since the diagram is a tree \added{(Lemma~\ref{L04a:lem:tree}, applied to the subdiagram on
$v,v_1,\dots,v_k$)}, $\ip{v_i}{v_j}=0$ for $i\neq j$, so $\{v_1,\dots,v_k\}$ is orthonormal. Also,
$\{v_1,\dots,v_k,v\}$ is linearly independent (by admissibility; for simple roots this is Theorem~\ref{L04a:thm:basis}), so
the normalized projection
\[
  v_0 := \frac{\operatorname{proj}_{\langle v_1,\dots,v_k\rangle^\perp}(v)}{\norm{\operatorname{proj}_{\langle v_1,\dots,v_k\rangle^\perp}(v)}}
\]
is defined, and $\ip{v}{v_0}\neq0$. Now $\{v_0,\dots,v_k\}$ is an orthonormal set, and $v$ lies in its
span, so
\[
  v=\sum_{i=0}^k\ip{v}{v_i}v_i \qquad\text{and}\qquad \ip{v}{v}=\sum_{i=0}^k\ip{v}{v_i}^2=1 .
\]
Since $\ip{v}{v_0}^2\neq0$,
\[
  \sum_{i=1}^k 4\ip{v}{v_i}^2<4 .
\]
But $4\cos^2\theta = 4\ip{v}{v_i}^2$ is the number of edges from $v$ to $v_i$. So the degree of $v$,
counted with multiplicity, is at most $3$.
\end{proof}

\begin{corollary}\label{L04a:cor:G2}
\tikz[baseline=-3pt]{\LIVtN{0}\LIVv{0,0}\LIVv{1,0}} is the only connected admissible diagram with a
triple edge.
\end{corollary}

\begin{proof}[\added{Proof}]
\added{If $v$ and $w$ are joined by a triple edge, both already have degree $3$, so by
Lemma~\ref{L04a:lem:degree} neither is joined to any other vertex; by connectedness there are no
other vertices.}
\end{proof}

So from now on we consider only double and single edges.

\begin{definition}\label{L04a:def:chain}
A \emph{simple chain} is a collection of vertices connected by single edges:
\begin{center}
\begin{tikzpicture}
  \LIVe{0,0}{1,0}\draw[dotted] (1,0)--(2,0);\LIVe{2,0}{3,0}
  \foreach\x in{0,1,2,3}{\LIVv{\x,0}}
  \node[below] at (0,-0.1) {\small$v_1$};
  \node[below] at (1,-0.1) {\small$v_2$};
  \node[below] at (2,-0.1) {\small$v_{k-1}$};
  \node[below] at (3,-0.1) {\small$v_k$};
\end{tikzpicture}
\end{center}
\end{definition}

\begin{lemma}[Simple chain collapse]\label{L04a:lem:collapse}
In an admissible diagram, a simple chain represented by $\{v_1,\dots,v_k\}$ can be replaced by the single
vector $v=\sum_{i=1}^k v_i$, giving an admissible diagram.
\end{lemma}

\begin{proof}
We want to show that $v$ is a unit vector and that the collapsed diagram is admissible.
\added{(Linear independence of the new set is clear, since the $v_i$ and the remaining vectors are
linearly independent.)}
We have
\[
  \ip{v}{v} = k + \sum_{i<j}2\ip{v_i}{v_j},
\]
and $\ip{v_i}{v_j}=0$ for all $i<j$ except $j=i+1$ \added{(a tree has no other edges among the $v_i$)}.
For a single edge, recall $2\ip{v_i}{v_{i+1}}=-1$. Hence
\[
  \ip{v}{v} = k-(k-1) = 1 .
\]
Now consider a vertex $u$ not represented in the chain. It connects to at most one vertex in the chain
(by the tree property), say $v_j$ \added{(if none, $\ip{u}{v}=0$)}. Then
\[
  \ip{u}{v}=\sum_{i=1}^k\ip{u}{v_i}=\ip{u}{v_j},
\]
so the angles $\theta_{u,v}$ and $\theta_{u,v_j}$ are the same, and the collapsed set is admissible.
\orig{scribbled-out ``recall admissible\dots'' at the foot of p.\,23}
\end{proof}
% ===== end chapters/L04a.tex =====
% ===== begin chapters/L04b.tex =====
% ============================================================
% MAT670_46, pages 24--46.
% Pages 24--31: continuation of Lecture 5 (classification of Coxeter graphs).
% Pages 32--46: Lecture 6 (root systems -> root lattices; codes).
% ============================================================
\tikzset{L04bv/.style={circle,draw,fill=white,inner sep=0pt,minimum size=6pt},
         L04bdbl/.style={double,double distance=2pt},
         L04blab/.style={font=\scriptsize}}

\section{Classification of Coxeter graphs (continued)}
\label{L04b:sec:coxeter-cont}
\orig{continues Lecture 5 from p.~23}

Recall from the previous pages: a connected admissible diagram is a tree; each vertex
meets at most three edges counted with multiplicity; the only admissible diagram with a
triple edge is \tikz[baseline=-3pt]{\node[L04bv] (a) at (0,0) {}; \node[L04bv] (b) at (0.6,0) {};
\draw (a.north east)--(b.north west); \draw (a)--(b); \draw (a.south east)--(b.south west);}
($G_2$); and a simple chain $\{v_1,\dots,v_k\}$ may be collapsed to the single unit vector
$v=\sum_i v_i$, leaving an admissible diagram. From now on only single and double edges occur.

\subsection{Forbidden configurations}

\begin{proposition}\label{L04b:prop:forbidden}
Coxeter graphs of admissible diagrams containing any of the following subgraphs are forbidden:
\begin{center}
\begin{tikzpicture}[scale=0.8]
  % two double edges
  \foreach \i in {0,1,2,3,4} \node[L04bv] (a\i) at (\i,0) {};
  \draw[L04bdbl] (a0)--(a1); \draw (a1)--(a2); \draw[dotted,thick] (a2)--(a3);
  \draw[L04bdbl] (a3)--(a4);
  % double edge + fork
  \begin{scope}[xshift=6cm]
  \foreach \i in {0,1,2,3} \node[L04bv] (b\i) at (\i,0) {};
  \node[L04bv] (b4) at (3.8,0.5) {}; \node[L04bv] (b5) at (3.8,-0.5) {};
  \draw[L04bdbl] (b0)--(b1); \draw (b1)--(b2); \draw[dotted,thick] (b2)--(b3);
  \draw (b3)--(b4); \draw (b3)--(b5);
  \end{scope}
  % two forks
  \begin{scope}[xshift=11.5cm]
  \node[L04bv] (c0) at (-0.8,0.5) {}; \node[L04bv] (c1) at (-0.8,-0.5) {};
  \foreach \i in {0,1,2} \node[L04bv] (d\i) at (\i,0) {};
  \node[L04bv] (c4) at (2.8,0.5) {}; \node[L04bv] (c5) at (2.8,-0.5) {};
  \draw (c0)--(d0); \draw (c1)--(d0); \draw (d0)--(d1); \draw[dotted,thick] (d1)--(d2);
  \draw (d2)--(c4); \draw (d2)--(c5);
  \end{scope}
\end{tikzpicture}
\end{center}
Hence a connected admissible diagram (with no triple edge) contains at most one branch
point \emph{or} at most one double edge, but not both.
\end{proposition}

\begin{proof}
\added{In each case collapse the simple chain joining the two special features (possibly of
length zero) to a single vertex. By the collapsing lemma the result is again admissible, but the
new vertex meets edges of total multiplicity $2+2$, $2+1+1$, or $1+1+1+1$, i.e.\ $4$, contradicting
the bound of $3$ edges (with multiplicity) at each vertex. The same argument rules out two double
edges, two branch points, or a double edge and a branch point anywhere in the tree, since they are
always joined by a simple chain.}
\end{proof}

So, besides $G_2$, the only possible connected Coxeter graphs are:
\begin{center}
\begin{tikzpicture}[scale=0.85]
  % chain with one double edge
  \foreach \x/\n/\l in {0/u1/u_1,1/u2/u_2,2.6/up/u_p,3.6/vq/v_q,5.2/v2/v_2,6.2/v1/v_1}
     {\node[L04bv] (\n) at (\x,0) {}; \node[L04blab,below=2pt] at (\x,0) {$\l$};}
  \draw (u1)--(u2); \draw[dotted,thick] (u2)--(up); \draw[L04bdbl] (up)--(vq);
  \draw[dotted,thick] (vq)--(v2); \draw (v2)--(v1);
  % branched
  \begin{scope}[xshift=8.3cm]
  \foreach \x/\n/\l in {0/a1/u_1,1/a2/u_2,2.4/ap/u_{p-1}}
     {\node[L04bv] (\n) at (\x,0) {}; \node[L04blab,below=2pt] at (\x,0) {$\l$};}
  \node[L04bv] (x) at (3.3,0) {}; \node[L04blab,below=2pt] at (3.3,0) {$x$};
  \node[L04bv] (bq) at (4.1,0.6) {}; \node[L04blab,above=2pt] at (4.1,0.6) {$v_{q-1}$};
  \node[L04bv] (b1) at (5.5,0.6) {}; \node[L04blab,above=2pt] at (5.5,0.6) {$v_1$};
  \node[L04bv] (cr) at (4.1,-0.6) {}; \node[L04blab,below=2pt] at (4.1,-0.6) {$w_{r-1}$};
  \node[L04bv] (c1) at (5.5,-0.6) {}; \node[L04blab,below=2pt] at (5.5,-0.6) {$w_1$};
  \draw (a1)--(a2); \draw[dotted,thick] (a2)--(ap); \draw (ap)--(x);
  \draw (x)--(bq); \draw[dotted,thick] (bq)--(b1);
  \draw (x)--(cr); \draw[dotted,thick] (cr)--(c1);
  \end{scope}
\end{tikzpicture}
\end{center}
together with the simple chains \tikz[baseline=-3pt]{\foreach \i/\x in {1/0,2/0.5,3/1.3,4/1.8} \node[L04bv] (s\i) at (\x,0) {};
\draw (s1)--(s2); \draw[dotted,thick] (s2)--(s3); \draw (s3)--(s4);}, which give $A_n$.

\subsection{One double edge: \texorpdfstring{$B_n$, $C_n$, $F_4$}{Bn, Cn, F4}}

Label the diagram with one double edge as above: $u_1,\dots,u_p$, then the double edge
$u_p = v_q$, then $v_q,\dots,v_1$. Put
\[
  u=\sum_{i=1}^p i\,u_i,\qquad v=\sum_{j=1}^q j\,v_j .
\]
Since the $u_i$ form a simple chain, $2\ip{u_i}{u_{i+1}}=-1$ for $1\le i\le p-1$ and all other pairs
are orthogonal. Hence
\begin{align*}
  \ip{u}{u} &= \sum_{i=1}^p i^2\ip{u_i}{u_i} + \sum_{i<j} 2ij\,\ip{u_i}{u_j}
   = \sum_{i=1}^p i^2 + \sum_{i=1}^{p-1} i(i+1)\cdot\underbrace{2\ip{u_i}{u_{i+1}}}_{=-1}\\
  &= \sum_{i=1}^p i^2 - \sum_{i=1}^{p-1} i(i+1) = p^2 - \sum_{i=1}^{p-1} i = \frac{p(p+1)}{2},
\end{align*}
and similarly $\ip{v}{v} = q(q+1)/2$.

Since the double edge $u_p = v_q$ is the only non-orthogonal relation between the two chains,
$\ip{u}{v} = pq\,\ip{u_p}{v_q}$, and $4\ip{u_p}{v_q}^2=2$. Thus
\[
  \ip{u}{v}^2 = \frac{p^2q^2}{2}.
\]
Since $u$ and $v$ are not parallel \added{(the $u_i, v_j$ are linearly independent)}, the
Cauchy--Schwarz inequality is strict: $\ip{u}{v}^2 < \ip{u}{u}\ip{v}{v}$, i.e.
\[
  \frac{p^2q^2}{2} < \frac{p(p+1)}{2}\cdot\frac{q(q+1)}{2},
  \qquad\text{so}\qquad 2pq < (p+1)(q+1),
\]
which, for integers $p,q\ge 1$, is
\[
  \boxed{(p-1)(q-1) < 2.}
\]
Hence either $p=q=2$, which is $F_4$, or $p=1$ (or $q=1$) with the other arbitrary, which is
$B_n$ and $C_n$ \added{($n=p+q$; the Coxeter graphs of $B_n$ and $C_n$ coincide, and they are
distinguished only by the direction of the arrow in the Dynkin diagram)}.

\subsection{One branch point: \texorpdfstring{$D_n$, $E_6$, $E_7$, $E_8$}{Dn, E6, E7, E8}}

Finally, take the branched diagram above, with branch vertex $x$ and three arms
$u_1,\dots,u_{p-1}$, $v_1,\dots,v_{q-1}$, $w_1,\dots,w_{r-1}$ ($p,q,r\ge 2$), the indices
increasing towards $x$. As before put
\[
  u=\sum_{i=1}^{p-1} i\,u_i,\qquad v=\sum_{i=1}^{q-1} i\,v_i,\qquad w=\sum_{i=1}^{r-1} i\,w_i .
\]
Note that $u,v,w$ are mutually orthogonal vectors (no edges between different arms), and
$x$ is not in their span $\langle u,v,w\rangle$ \added{(by linear independence of all the vertices)}.
By the previous computation with $p$ replaced by $p-1$,
$\ip{u}{u} = p(p-1)/2$, and similarly for $v,w$.

Since $x$ is a unit vector not in $\langle u,v,w\rangle$, and $u/\norm{u},v/\norm{v},w/\norm{w}$ is an
orthonormal basis of that space, the squared length of the projection of $x$ onto it is strictly
less than $1$:
\[
  1 = \ip{x}{x} > \Big\langle x,\frac{u}{\norm{u}}\Big\rangle^2 + \Big\langle x,\frac{v}{\norm{v}}\Big\rangle^2
  + \Big\langle x,\frac{w}{\norm{w}}\Big\rangle^2 .
\]
Also $\ip{x}{u_i}=0$ unless $i=p-1$, and $4\ip{x}{u_{p-1}}^2 = 1$ (single edge). Then
\[
  \ip{x}{u}^2 = \Big(\sum_{i=1}^{p-1} i\,\ip{x}{u_i}\Big)^2 = (p-1)^2\ip{x}{u_{p-1}}^2 = \frac{(p-1)^2}{4}
\]
\fixed{notes: $\sum i^2\langle x,u_i\rangle$, square misplaced}%
and therefore
\[
  \Big\langle x,\frac{u}{\norm{u}}\Big\rangle^2 = \frac{(p-1)^2}{4}\cdot\frac{2}{p(p-1)}
   = \frac12\Big(1-\frac1p\Big).
\]
The same holds for $v$ (with $q$) and $w$ (with $r$). Then
\[
  2 > \Big(1-\frac1p\Big) + \Big(1-\frac1q\Big) + \Big(1-\frac1r\Big)
  \quad\Longrightarrow\quad \frac1p + \frac1q + \frac1r > 1,
\]
where $p,q,r\ge 2$ are integers. Assume $p\ge q\ge r\ge 2$. Then $3/r > 1$ forces $r=2$, and:
\begin{itemize}
  \item $q=2$: any $p\ge 2$ works \added{--- this is $D_{p+2}$};
  \item $q=3$: $\frac1q+\frac1r = \frac56$, so $\frac1p>\frac16$, i.e.\ $3\le p<6$
        \added{--- $p=3,4,5$ give $E_6,E_7,E_8$};
  \item $q\ge4$: \added{$\frac1p+\frac1q+\frac12 \le \frac14+\frac14+\frac12=1$}, so there are no solutions.
\end{itemize}
\added{(The diagram has $p+q+r-2$ vertices, which gives the subscripts.)}

\begin{addedblock}
\textbf{Summary.} The connected Coxeter graphs of admissible diagrams (equivalently, of irreducible
crystallographic root systems) are
$A_n$ ($n\ge1$), $B_n=C_n$ as graphs ($n\ge2$), $D_n$ ($n\ge4$), $E_6$, $E_7$, $E_8$, $F_4$, $G_2$.
\begin{center}
\begin{tikzpicture}[scale=0.62, every node/.style={font=\small}]
  % A_n
  \node[anchor=east] at (-0.4,0) {$A_n$};
  \foreach \i in {0,1,3,4} \node[L04bv] (A\i) at (\i,0) {};
  \draw (A0)--(A1); \draw[dotted,thick] (A1)--(A3); \draw (A3)--(A4);
  % B_n / C_n
  \node[anchor=east] at (-0.4,-1.2) {$B_n,C_n$};
  \foreach \i in {0,1,3,4} \node[L04bv] (B\i) at (\i,-1.2) {};
  \draw (B0)--(B1); \draw[dotted,thick] (B1)--(B3); \draw[L04bdbl] (B3)--(B4);
  % D_n
  \node[anchor=east] at (-0.4,-2.6) {$D_n$};
  \foreach \i in {0,1,3} \node[L04bv] (D\i) at (\i,-2.6) {};
  \node[L04bv] (Du) at (3.8,-2.1) {}; \node[L04bv] (Dd) at (3.8,-3.1) {};
  \draw (D0)--(D1); \draw[dotted,thick] (D1)--(D3); \draw (D3)--(Du); \draw (D3)--(Dd);
  % F4, G2
  \node[anchor=east] at (-0.4,-4.0) {$F_4$};
  \foreach \i in {0,1,2,3} \node[L04bv] (F\i) at (\i,-4.0) {};
  \draw (F0)--(F1); \draw[L04bdbl] (F1)--(F2); \draw (F2)--(F3);
  \node[anchor=east] at (-0.4,-5.2) {$G_2$};
  \node[L04bv] (G0) at (0,-5.2) {}; \node[L04bv] (G1) at (1,-5.2) {};
  \draw (G0)--(G1); \draw ([yshift=2pt]G0.east)--([yshift=2pt]G1.west);
  \draw ([yshift=-2pt]G0.east)--([yshift=-2pt]G1.west);
  % E6, E7, E8
  \begin{scope}[xshift=9cm]
  \foreach \n/\k/\y in {E_6/5/0, E_7/6/-2, E_8/7/-4} {
    \node[anchor=east] at (-0.4,\y) {$\n$};
    \foreach \i in {1,...,\k} \node[L04bv] (e\i) at (\i-1,\y) {};
    \foreach \i [evaluate=\i as \j using int(\i-1)] in {2,...,\k} \draw (e\i)--(e\j);
    \node[L04bv] (top) at (2,\y+0.9) {}; \draw (e3)--(top);
  }
  \end{scope}
\end{tikzpicture}
\end{center}
\end{addedblock}

\begin{remark}
It remains to show that each of
$A_n$, $B_n$, $C_n$, $D_n$, $E_6$, $E_7$, $E_8$, $F_4$, $G_2$ actually exists, i.e.\ to construct
the corresponding root systems. This is done in the next lecture; $G_2$ is left as an exercise
\added{(its $12$ roots were drawn in Lecture~4: check axioms 1)--4) directly)}.
\orig{p.~31: ``Existence of $A_n,\dots,F_4$: construct next time; $G_2$ -- ex.''}
\end{remark}

% ============================================================
\chapter{Lecture 6: Root systems and root lattices}
\label{L04b:ch:lecture6}

\section{Construction of crystallographic root systems}
\label{L04b:sec:construction}

\subsection*{\texorpdfstring{$A_n$}{An}}
In $A_n$ the simple roots are pairwise orthogonal or at angle $2\pi/3$.

Consider $x\in\Z^{n+1}$ such that $\sum_{i=1}^{n+1} x_i = 0$. These form an $n$-dimensional
\added{sublattice} $S$ of $\Z^{n+1}$. Let
\[
  \Delta = \{x\in S : \norm{x}=\sqrt2\}.
\]
That is, all but two coordinates are $0$, one coordinate is $+1$ and one is $-1$: the vectors
are the permutations of $(1,-1,0,\dots,0)$, so $|\Delta| = n(n+1)$. Choose the simple roots
\[
  \alpha_i = e_i - e_{i+1},\qquad i=1,\dots,n,
\]
\orig{diagram labelled $e_{i+1}-e_i$}
with matrix of coordinates and Dynkin diagram
\[
  \begin{bmatrix} 1 & -1 & & & \\ & 1 & -1 & & \\ & & \ddots & \ddots & \\ & & & 1 & -1\end{bmatrix}
  \qquad
  \begin{tikzpicture}[baseline=-3pt,scale=0.9]
    \foreach \x/\n/\l in {0/a/e_1-e_2,1.4/b/e_2-e_3,3.4/c/e_n-e_{n+1}}
      {\node[L04bv] (\n) at (\x,0) {}; \node[L04blab,below=2pt] at (\x,0) {$\l$};}
    \draw (a)--(b); \draw[dotted,thick] (b)--(c);
  \end{tikzpicture}
\]
To check the Dynkin diagram: the number of edges joining $\alpha_i$ and $\alpha_j$ is
$4\ip{\alpha_i}{\alpha_j}^2/(\norm{\alpha_i}^2\norm{\alpha_j}^2) = \ip{\alpha_i}{\alpha_j}^2\in\{1,0\}$,
and it is $1$ exactly when $j=i\pm1$.

Note that these roots generate the full sublattice $\{x\in\Z^{n+1}: \sum x_i=0\}$
\added{(any such $x$ equals $\sum_{k=1}^{n} (x_1+\dots+x_k)\,\alpha_k$)}. So the $A_n$ root lattice
is all of $S$. This description is maybe not so useful for visualizing, since $S$ sits inside
$\R^{n+1}$.

\begin{example}[$A_2$]
Simple roots $(-1,1,0)$ and $(0,-1,1)$, whose sum is $(-1,0,1)$; together with the negatives
$(1,-1,0)$, $(0,1,-1)$, $(1,0,-1)$ these are the six roots, lying in the plane
$x_1+x_2+x_3=0$. In that plane they form the familiar hexagonal configuration:
\begin{center}
\begin{tikzpicture}[scale=1.1]
  \foreach \a in {0,60,...,300} \draw[-{Stealth}] (0,0)--(\a:1);
  \draw[-{Stealth},very thick,addcol] (0,0)--(0:1) node[right,font=\scriptsize] {$\alpha_1$};
  \draw[-{Stealth},very thick,addcol] (0,0)--(120:1) node[above left,font=\scriptsize] {$\alpha_2$};
\end{tikzpicture}
\end{center}
\orig{p.~33: pink star in the plane $\sum x_i=0$}
\end{example}

\begin{example}[$A_3$]
The $12$ roots of $A_3$ are the $12$ nearest neighbours of the origin in the FCC lattice
(``the FCC system'') \added{--- the $A_3$ root lattice is the face-centred cubic lattice, with
kissing number $12$}.
\end{example}

\subsection*{\texorpdfstring{$B_n$}{Bn}}
$\Delta$ is the set of all integer vectors $x\in\Z^n$ of norm $1$ or norm $\sqrt2$, i.e.\
$\pm e_i$ and $\pm e_i\pm e_j$. So $|\Delta| = 2n + 4\binom n2 = 2n^2$. Simple roots:
\[
  \alpha_i = e_i - e_{i+1}\ (1\le i\le n-1),\qquad \text{short root } \alpha_n = e_n,
\]
\[
  \begin{bmatrix} 1 & -1 & & & \\ & 1 & -1 & & \\ & & \ddots & \ddots & \\ & & & 1 & -1\\ & & & & 1\end{bmatrix}
  \qquad
  \begin{tikzpicture}[baseline=-3pt,scale=0.9]
    \foreach \x/\n/\l in {0/a/e_1-e_2,1.4/b/e_2-e_3,3.2/c/e_{n-1}-e_n,4.6/d/e_n}
      {\node[L04bv] (\n) at (\x,0) {}; \node[L04blab,below=2pt] at (\x,0) {$\l$};}
    \draw (a)--(b); \draw[dotted,thick] (b)--(c); \draw[L04bdbl] (c)--(d);
    \draw[thick] (3.8,0.12)--(3.98,0)--(3.8,-0.12);
  \end{tikzpicture}
\]
Sadly, from this construction the associated root lattice is all of $\Z^n$ \added{(it contains
every $e_i$)}, so we get nothing new.

\subsection*{\texorpdfstring{$C_n$}{Cn}}
$\Delta$ is the set of all integer vectors of the form $2x$ with $x\in\Z^n$ of length $1$, or
$x\in\Z^n$ of length $\sqrt2$; i.e.\ $\pm 2e_i$ and $\pm e_i\pm e_j$. Again
$|\Delta| = 2n + 4\binom n2 = 2n^2$. Simple roots:
\[
  \alpha_i = e_i - e_{i+1}\ (1\le i\le n-1),\qquad \text{long root } \alpha_n = 2e_n,
\]
\[
  \begin{bmatrix} 1 & -1 & & & \\ & 1 & -1 & & \\ & & \ddots & \ddots & \\ & & & 1 & -1\\ & & & & 2\end{bmatrix}
  \qquad
  \begin{tikzpicture}[baseline=-3pt,scale=0.9]
    \foreach \x/\n/\l in {0/a/e_1-e_2,1.4/b/e_2-e_3,3.2/c/e_{n-1}-e_n,4.6/d/2e_n}
      {\node[L04bv] (\n) at (\x,0) {}; \node[L04blab,below=2pt] at (\x,0) {$\l$};}
    \draw (a)--(b); \draw[dotted,thick] (b)--(c); \draw[L04bdbl] (c)--(d);
    \draw[thick] (4.0,0.12)--(3.82,0)--(4.0,-0.12);
  \end{tikzpicture}
\]
\fixed{diagram drawn with a single last edge}
So the $C_n$ root lattice is the set of all integer vectors with \emph{even} coordinate sum.

\begin{example}[$B_2$ and $C_2$]
The two root systems, with the simple roots highlighted:
\begin{center}
\begin{tikzpicture}[scale=0.9]
  \foreach \v in {(1,0),(0,1),(-1,0),(0,-1),(1,1),(1,-1),(-1,1),(-1,-1)} \draw[-{Stealth},gray] (0,0)--\v;
  \draw[-{Stealth},very thick,addcol] (0,0)--(1,-1) node[below right,font=\scriptsize] {$\alpha_1=e_1-e_2$};
  \draw[-{Stealth},very thick,addcol] (0,0)--(0,1) node[above,font=\scriptsize] {$\alpha_2=e_2$};
  \node at (0,-1.7) {$B_2$};
  \begin{scope}[xshift=5cm]
  \foreach \v in {(2,0),(0,2),(-2,0),(0,-2),(1,1),(1,-1),(-1,1),(-1,-1)} \draw[-{Stealth},gray] (0,0)--\v;
  \draw[-{Stealth},very thick,addcol] (0,0)--(1,-1) node[below right,font=\scriptsize] {$\alpha_1=e_1-e_2$};
  \draw[-{Stealth},very thick,addcol] (0,0)--(0,2) node[above,font=\scriptsize] {$\alpha_2=2e_2$};
  \node at (0,-2.4) {$C_2$};
  \end{scope}
\end{tikzpicture}
\end{center}
\orig{redrawn; sketch showed only $\alpha_1,\alpha_2$}
\added{In both, the simple roots meet at angle $3\pi/4$; $C_2$ is $B_2$ rotated by $\pi/4$ and rescaled.}
\end{example}

\subsection*{\texorpdfstring{$D_n$}{Dn}}
$\Delta$ is the set of all integer vectors $x\in\Z^n$ of length $\sqrt2$, i.e.\ $\pm e_i\pm e_j$, so
\[
  |\Delta| = 4\binom n2 = 2n(n-1).
\]
\fixed{notes: $2\binom n2$}
Simple roots:
\[
  \alpha_i = e_i - e_{i+1}\ (1\le i\le n-1)\qquad\text{and}\qquad \alpha_n = e_{n-1}+e_n,
\]
\[
  \begin{bmatrix} 1 & -1 & & & \\ & \ddots & \ddots & & \\ & & 1 & -1 & \\ & & & 1 & -1\\ & & & 1 & 1\end{bmatrix}
  \qquad
  \begin{tikzpicture}[baseline=-3pt,scale=0.9]
    \foreach \x/\n/\l in {0/a/e_1-e_2,1.4/b/e_2-e_3,3.2/c/e_{n-2}-e_{n-1}}
      {\node[L04bv] (\n) at (\x,0) {}; \node[L04blab,below=2pt] at (\x,0) {$\l$};}
    \node[L04bv] (u) at (4.3,0.55) {}; \node[L04blab,right=2pt] at (4.3,0.55) {$e_{n-1}+e_n$};
    \node[L04bv] (d) at (4.3,-0.55) {}; \node[L04blab,right=2pt] at (4.3,-0.55) {$e_{n-1}-e_n$};
    \draw (a)--(b); \draw[dotted,thick] (b)--(c); \draw (c)--(u); \draw (c)--(d);
  \end{tikzpicture}
\]
This also generates the lattice $\{x\in\Z^n : \text{the coordinate sum is even}\}$, the same as
for $C_n$. (I can't draw $D_4$, sadly.)

\subsection*{\texorpdfstring{$F_4$}{F4}}
$\Delta$ is the set of vectors $x\in\R^4$ of length $1$ or $\sqrt2$ such that the coordinates of
$2x$ are all even or all odd integers. \added{Explicitly: $\pm e_i$ ($8$), $\pm e_i\pm e_j$ ($24$) and
$\tfrac12(\pm1,\pm1,\pm1,\pm1)$ ($16$), so $|\Delta|=48$.} Simple roots:
\[
  \alpha_1 = (1,-1,0,0),\quad \alpha_2=(0,1,-1,0),\quad \alpha_3=(0,0,1,0),\quad
  \alpha_4 = -\tfrac12(1,1,1,1),
\]
\begin{center}
  \begin{tikzpicture}[baseline=-3pt,scale=0.9]
    \foreach \x/\n in {0/1,1/2,2/3,3/4}
      {\node[L04bv] (f\n) at (\x,0) {}; \node[L04blab,below=2pt] at (\x,0) {$\alpha_\n$};}
    \draw (f1)--(f2); \draw[L04bdbl] (f2)--(f3); \draw (f3)--(f4);
    \draw[thick] (1.42,0.12)--(1.6,0)--(1.42,-0.12);
  \end{tikzpicture}
\end{center}
(check by hand). \added{Indeed $\ip{\alpha_1}{\alpha_2}=-1$ with both of length $\sqrt2$ (angle $2\pi/3$);
$\ip{\alpha_2}{\alpha_3}=-1$ with lengths $\sqrt2,1$ (angle $3\pi/4$, double edge);
$\ip{\alpha_3}{\alpha_4}=-\frac12$ with both of length $1$ (angle $2\pi/3$); all other pairs are orthogonal.}

The $F_4$ root lattice: all coordinates are integers, or all coordinates are halves of odd integers
--- no mixing. \added{That is, $\Z^4\cup(\Z^4+\tfrac12(1,1,1,1))$, which is $D_4^*$, a scaled copy of $D_4$.}

\subsection*{\texorpdfstring{$G_2$}{G2}}
We already drew this one (the $12$ roots form a ``star of David'', two hexagons of lengths $1$ and
$\sqrt3$). \orig{star sketch, see Lecture~4}
The $G_2$ root lattice coincides with the $A_2$ root lattice \added{(the hexagonal lattice)}.

% ------------------------------------------------------------
\section{The \texorpdfstring{$E_8$}{E8} root system}
\label{L04b:sec:E8}

\begin{definition}\label{L04b:def:E8}
The $E_8$ root system consists of the $|\Delta| = 240$ vectors in $\R^8$ of the forms
\[
  \pm e_i \pm e_j\ (i\neq j) \qquad\text{and}\qquad \tfrac12(\pm e_1 \pm \dots \pm e_8)
  \ \added{\text{with an even number of minus signs}}.
\]
\fixed{added the even-sign condition}
\added{There are $4\binom82=112$ of the first kind and $2^7=128$ of the second.}
\end{definition}

Simple roots:
\[
  \alpha_i = e_i-e_{i+1}\ (1\le i\le 6),\qquad \alpha_7 = e_6+e_7,\qquad \alpha_8 = -\tfrac12(1,\dots,1),
\]
\[
  \begin{bmatrix}
   1&-1&0&0&0&0&0&0\\ &1&-1&0&0&0&0&0\\ &&1&-1&0&0&0&0\\ &&&1&-1&0&0&0\\ &&&&1&-1&0&0\\
   &&&&&1&-1&0\\ &&&&&1&1&0\\ -\frac12&-\frac12&-\frac12&-\frac12&-\frac12&-\frac12&-\frac12&-\frac12
  \end{bmatrix}
  \qquad
  \begin{tikzpicture}[baseline=-3pt,scale=0.8]
    \foreach \i in {1,...,6} {\node[L04bv] (g\i) at (\i-1,0) {}; \node[L04blab,below=2pt] at (\i-1,0) {$\alpha_{\i}$};}
    \node[L04bv] (g7) at (4,1) {}; \node[L04blab,above=2pt] at (4,1) {$\alpha_7$};
    \node[L04bv] (g8) at (5,1) {}; \node[L04blab,above=2pt] at (5,1) {$\alpha_8$};
    \foreach \i/\j in {1/2,2/3,3/4,4/5,5/6,5/7,7/8} \draw (g\i)--(g\j);
  \end{tikzpicture}
\]
\added{(Check: $\ip{\alpha_7}{\alpha_5}=-1$, $\ip{\alpha_7}{\alpha_6}=0$, $\ip{\alpha_8}{\alpha_7}=-1$ and
$\ip{\alpha_8}{\alpha_i}=0$ for $i\le 6$; the arms at $\alpha_5$ give $(p,q,r)=(5,3,2)$.)}

This generates the lattice $E_8$ of all $x\in\R^8$ such that
\begin{enumerate}[label=\arabic*)]
  \item the coordinates are all integers or all half-integers \added{(halves of odd integers)}, and
  \item the sum of the coordinates is even.
\end{enumerate}

This is clearly related to $D_8$. Consider the ``deep holes'' of the $D_n$ lattice, i.e.\ the
half-integer points such as $h=\frac12(1,\dots,1)$. In dimension $8$,
\[
  \norm{\tfrac12(1,\dots,1)} = \sqrt2,
\]
and a typical shortest lattice vector has $\norm{(1,1,0,\dots,0)}=\sqrt2$. Similarly,
\[
  \norm{(1,1,0,\dots,0)-\tfrac12(1,\dots,1)} = \sqrt2 .
\]
So if we have $D_8$, there is another copy $D_8+h$ that fits into the holes of $D_8$ without
decreasing the minimal distance: $E_8 = D_8\cup(D_8+h)$.
\added{(For general $n$, $h$ is at distance $\sqrt n/2$ from $D_n$; this equals the minimal
distance $\sqrt2$ exactly when $n=8$.)}

\subsection*{\texorpdfstring{$E_7$, $E_6$}{E7, E6}}
Let $v=(1,1,\dots,1)$ and $w=-e_7-e_8$. Then
\[
  E_7 = E_8\cap v^\perp,\qquad E_6 = E_8\cap \langle w,v\rangle^\perp .
\]
\added{Here $\frac12 v=-\alpha_8$ and $w$ are roots with $\ip{\tfrac12 v}{w}=-1$, so they span an
$A_2$; thus $E_7$ and $E_6$ are the orthogonal complements of an $A_1$ and an $A_2$ in $E_8$.}

These descriptions are not so satisfying, but perhaps there is an easy-to-see construction given
by collapsing simple chains: replace $\alpha_1,\alpha_2$ by $\alpha_1+\alpha_2$, then
$\alpha_1+\alpha_2,\alpha_3$ by $\alpha_1+\alpha_2+\alpha_3$, and so on.
\begin{center}
\begin{tikzpicture}[scale=0.85]
  \foreach \x/\n in {0/1,1/2,2/3,3/4,4/5,5/7,6/8}
    {\node[L04bv] (h\n) at (\x,0) {}; \node[L04blab,below=2pt] at (\x,0) {$\alpha_\n$};}
  \node[L04bv] (h6) at (4,1) {}; \node[L04blab,above=2pt] at (4,1) {$\alpha_6$};
  \foreach \i/\j in {1/2,2/3,3/4,4/5,5/7,7/8,5/6} \draw (h\i)--(h\j);
  \draw (0,0.35)--(0,0.55)--(1,0.55)--(1,0.35);
  \draw (0,0.65)--(0,0.85)--(2,0.85)--(2,0.65);
\end{tikzpicture}
\end{center}
By the collapsing lemma, $\alpha_1+\alpha_2,\alpha_3,\dots,\alpha_8$ is a simple system of type $E_7$,
and $\alpha_1+\alpha_2+\alpha_3,\alpha_4,\dots,\alpha_8$ one of type $E_6$, \added{so the lattices they
generate are copies of the $E_7$, $E_6$ root lattices inside $E_8$.} This continues
\added{(collapsing suitable chains)} to
\[
  D_5,\quad D_4,\quad A_3,\quad A_2,\quad A_1 .
\]

\begin{exercise}
Find other constructions of $E_8$. \added{(For instance via the extended Hamming code and
Construction~A, \S\ref{L04b:sec:codes}.)}
\end{exercise}

% ------------------------------------------------------------
\section{Laminated lattices}
\label{L04b:sec:laminated}

Before (Lecture~4) we had the lattices
\[
  A_1,\ A_2,\ A_3,\ D_4,\ D_5,\ E_6,\ E_7,\ E_8
\]
as the best lattice packings, and best known packings, in dimensions $1$ to $8$. Somehow these
are all related to $E_8$ (as we just saw, they all sit inside $E_8$).

\begin{addedblock}
Optimality among lattices is proven in dimensions $1$--$8$ (Lagrange $n=2$, Gauss $n=3$,
Korkine--Zolotarev $n=4,5$, Blichfeldt $n=6,7,8$); among \emph{all} packings it is proven for
$n=2$ (Thue, Fejes T\'oth), $n=3$ (Hales) and $n=8$ (Viazovska, posted March 2016, during this course).
Similarly the Leech lattice is optimal in dimension $24$ (Cohn--Kumar among lattices;
Cohn--Kumar--Miller--Radchenko--Viazovska 2016/17 among all packings).
\end{addedblock}

Dimension $9$ is, as far as I know, still unknown.
\added{(This is still the case, to the editor's knowledge.)}

The best constructions for lattices come from an ``inductive'' process. Start with the even
integers $\Lambda_1 = 2\Z\subset\R^1$.

\begin{definition}\label{L04b:def:laminated}
The \emph{$n$-dimensional laminated lattice} $\Lambda_n$ is a lattice of maximal density among
all lattices in $\R^n$ with shortest vector of length $\norm{x}=2$ that contain $\Lambda_{n-1}$
as a sublattice.
\end{definition}
So $\Lambda_n$ is obtained by gluing: $\Lambda_n=\langle \Lambda_{n-1}, g\rangle$\added{, i.e.\ by stacking
layers $\Lambda_{n-1}+kg$, $k\in\Z$, with each new layer placed over the deepest holes of the previous one}.
\begin{center}
\begin{tikzpicture}[scale=0.45]
  % A1
  \foreach \i in {0,...,4} \fill (\i,0) circle (2.5pt);
  \node at (2,-1.6) {$A_1$};
  \draw[-{Stealth}] (5,0)--(6.5,0);
  % A2
  \begin{scope}[xshift=7.5cm]
    \foreach \j in {-2,...,2} \foreach \i in {0,...,4} {
      \pgfmathsetmacro{\xx}{\i+0.5*mod(abs(\j),2)} \fill (\xx,{0.866*\j}) circle (2.5pt); }
    \node at (2.2,-3.0) {$A_2$};
  \end{scope}
  \draw[-{Stealth}] (13.3,0)--(14.8,0);
  % A3: hexagonal layer + next layer over holes
  \begin{scope}[xshift=16cm]
    \foreach \j in {-2,...,2} \foreach \i in {0,...,4} {
      \pgfmathsetmacro{\xx}{\i+0.5*mod(abs(\j),2)} \fill (\xx,{0.866*\j}) circle (2.5pt); }
    \foreach \j in {-2,...,1} \foreach \i in {0,...,3} {
      \pgfmathsetmacro{\xx}{\i+0.5+0.5*mod(abs(\j),2)} \draw (\xx,{0.866*\j+0.2887}) circle (3.5pt); }
    \node at (2.2,-3.0) {$A_3$};
  \end{scope}
  \draw[-{Stealth}] (21.8,0)--(23.3,0);
  \node at (24.6,0) {$\cdots\ E_8$};
\end{tikzpicture}
\end{center}
\orig{p.~40 sketch; circles = next layer}
The laminated lattices $\Lambda_1,\dots,\Lambda_8$ coincide with the root lattices above
\added{(scaled to minimal length $2$)}.

\begin{remark}
Dimension $9$ is unusual: there is a continuous family of packings built from layers of $\Lambda_8$,
\fixed{notes: ``cts family of lattices''}
all of the same density as $\Lambda_9$. \added{These are the ``fluid'' packings; in general they are
\emph{not} lattice packings, so ``family of lattices'' should be read as ``family of packings'' (see
Conway--Sloane, \emph{Sphere Packings, Lattices and Groups}, Ch.~1 and Ch.~6).}
\end{remark}

% ------------------------------------------------------------
\section{Dimension 10: codes and Construction A}
\label{L04b:sec:codes}

In dimension $10$, a \emph{non-lattice} packing beats the best known lattices.

\begin{definition}\label{L04b:def:code}
A \emph{binary code} of length $n$ is a subset $C\subseteq\{0,1\}^n$ of the binary words of length $n$.
The elements of $C$ are the \emph{code words}. \added{The \emph{Hamming distance} between two words is the
number of positions in which they differ; $d(C)$ denotes the minimal distance between distinct code words.}
\end{definition}

\begin{definition}[Construction A]\label{L04b:def:constrA}
Given a binary code $C$ of length $n$, there is a packing
\fixed{notes: $|C|=n$}
\[
  P(C) = \{x\in\Z^n : x \bmod 2 \in C\}.
\]
\end{definition}
\begin{addedblock}
Two distinct points of $P(C)$ either differ by an element of $2\Z^n$ (distance $\ge2$) or reduce to
different code words, hence differ by $\pm1$ in at least $d(C)$ coordinates (distance $\ge\sqrt{d(C)}$).
So $P(C)$ is a packing of balls of radius $\frac12\min(2,\sqrt{d(C)})$ with $|C|$ centres per cube
$[0,2)^n$, and its centre density is $|C|\,2^{-n}\big(\tfrac12\min(2,\sqrt{d(C)})\big)^n$.
\end{addedblock}

This gives some relation between packing problems in high dimensions and error detecting /
error correcting codes. Think of a code word sent through a continuous channel and received
with an error of size up to $\varepsilon$: the centres of the $\varepsilon$-spheres of a packing correspond
to code words that are at distance $\ge 2\varepsilon$ apart. Dense packings correspond to large collections of
code words that are well separated.
\orig{sketch: $\varepsilon$-ball around a code word}

\begin{example}[Block codes]
The $(6,3)$ repetition code repeats the pattern: $101\mapsto101101$. The $(3,2)$ parity check code
(mod $2$) appends a parity bit: $11\mapsto 110$. We can see the latter on the Hamming cube: its code
words $000,011,101,110$ all have Hamming distance $2$, i.e.\ are $2$ bits apart.
\begin{center}
\begin{tikzpicture}[scale=1.3]
  \coordinate (000) at (0,0); \coordinate (100) at (1,0); \coordinate (010) at (0.45,0.35);
  \coordinate (001) at (0,1); \coordinate (110) at (1.45,0.35); \coordinate (101) at (1,1);
  \coordinate (011) at (0.45,1.35); \coordinate (111) at (1.45,1.35);
  \draw (000)--(100)--(101)--(001)--cycle; \draw (001)--(011)--(111)--(101); \draw (100)--(110)--(111);
  \draw[dashed] (000)--(010)--(110); \draw[dashed] (010)--(011);
  \foreach \p in {000,011,101,110} {\fill[addcol] (\p) circle (2.2pt);}
  \node[below left,font=\scriptsize] at (000) {$000$}; \node[above,font=\scriptsize] at (011) {$011$};
  \node[right,font=\scriptsize] at (101) {$101$}; \node[right,font=\scriptsize] at (110) {$110$};
\end{tikzpicture}
\end{center}
In general, parity check codes (all even-weight words) correspond to $D_n$, the lattice with even
coordinate sums.
\end{example}

Construction A thus links binary codes and lattices: $P(C)$ is a lattice if and only if $C$ is a
\emph{linear} code, that is, sums of code words (mod $2$) are code words.
The Hamming codes give $E_7$ and $E_8$. \added{More precisely (after scaling by $1/\sqrt2$): the extended
Hamming code $[8,4,4]$ gives $E_8$ ($16+14\cdot16=240$ vectors of norm $4$ in $P(C)$), and the $[7,3,4]$
simplex code, the dual of the Hamming code $[7,4,3]$, gives $E_7$ ($14+7\cdot16=126$).}

\subsection*{The code \texorpdfstring{$C_{10}$}{C10}}
So in dimension $10$ there is a special binary code $C_{10}$ \added{(Best's code)}, which is unique
\added{(Litsyn--Vardy)} and quite dense. It is based on a Gray code, or reflected binary code.

\begin{definition}
A \emph{Gray code} is a map $(\Z_2)^n\to(\Z_2)^n$, \added{a re-ordering of all binary words of length $n$},
where the successor of an element differs by a single bit flip.
\end{definition}
\[
  \begin{array}{ccc}
    \text{binary} & & \text{Gray}\\ \hline
    00 & 0 & 00\\ 01 & 1 & 01\\ 10 & 2 & 11\\ 11 & 3 & 10
  \end{array}
\]
There is a replacement procedure to convert between binary and Gray codes.

\begin{exercise}
Find the conversion binary $\leftrightarrow$ Gray, and iterate it to get the Gray code of length $n$.
\added{(Hint: $g = b \oplus (b\gg1)$; conversely $b_k = g_1\oplus\dots\oplus g_k$ from the leading bit.)}
\end{exercise}

In dimension $10$ \added{we look for} binary codes with minimum distance at least $4$. We construct
$C_{10}$ as follows. Take $(a,b,c,d,e)\in(\Z_4)^5$ satisfying
\[
  b,c,d\in\{\pm1\},\qquad a = c-d,\qquad e = b+c,
\]
\unclear{$\{\pm1\}$: a crossed-out entry; $\{\pm1\}$ fits the count $2^3$}
together with all cyclic permutations thereof. Then apply the Gray map
\[
  0\mapsto 00,\quad 1\mapsto 01,\quad 2\mapsto 11,\quad 3\mapsto 10.
\]
This gives $8\cdot5=40$ code words ($2^3$ choices of $(b,c,d)$ times $5$ cyclic shifts) of length $10$.

\begin{claim}
The Hamming distance between distinct code words of $C_{10}$ is at least $4$.
\end{claim}
\begin{addedblock}
The Gray map is an isometry from $\Z_4$ with the Lee metric ($d(0,2)=d(1,3)=2$, other distinct pairs
distance $1$) to $\Z_2^2$ with the Hamming metric, so it suffices to check Lee distance $\ge4$ in $(\Z_4)^5$.
(Editor: confirmed by direct computer check of all $\binom{40}{2}$ pairs; the $40$ words are distinct.)
\end{addedblock}

Moreover $A(10,4)=40$, i.e.\ $40$ is the maximal size of a binary code of length $10$ and minimal
distance $4$, so $C_{10}$ is optimal for dimension $10$.
\unclear{notes: ``LP bound $\Rightarrow<40.\ldots$''; $A(10,4)=40$ is due to Best (c.~1980)}
\begin{addedblock}
Construction A applied to $C_{10}$ gives the non-lattice packing $P_{10c}$: radius $1$ and $40$ centres
per $2^{10}$ volume, so centre density $40/2^{10}=5/128\approx0.0391$, larger than that of the best known
lattice $\Lambda_{10}$ (centre density $1/(16\sqrt3)\approx0.0361$). See Conway--Sloane, Ch.~5.
\end{addedblock}

% ------------------------------------------------------------
\section{An aside: lattice vs.\ non-lattice packings of other objects}
\label{L04b:sec:aside}

For other objects than balls, the question of lattice versus non-lattice packings looks different:
\begin{itemize}
  \item ellipsoid packings;
  \item polygon packings;
  \item honeycomb (HCB) and bubble problems. \unclear{``HCB'', ``Bubble prob.''?}
\end{itemize}
\added{For example, in the plane a lattice packing is optimal for any centrally symmetric convex body
(Fejes T\'oth, Rogers), but not for some non-symmetric ones; and congruent ellipsoids in $\R^3$ can be packed
more densely by non-lattice arrangements than by any lattice (Bezdek--Kuperberg). Hales proved the honeycomb
conjecture in 1999.} Density issues are subtle here, e.g.\ in $2$d versus $1$d.

On the same page there is a computation: with $a^2b=1$,
\[
  \begin{bmatrix} 1&&\\ &a&\\ &&a^2\end{bmatrix}
  \begin{pmatrix}1\\1\\1\end{pmatrix} = \begin{pmatrix}1\\a\\a^2\end{pmatrix},\qquad
  \begin{bmatrix} 1&&\\ &a&\\ &&a^2\end{bmatrix}
  \begin{pmatrix}a\\a\\b\end{pmatrix} = \begin{pmatrix}a\\a^2\\1\end{pmatrix}.
\]
\unclear{purpose unclear; maybe an ellipsoid-packing example}

% ------------------------------------------------------------
%======================================================================
\section{Extension: why root lattices pack well (in small dimensions)}
\label{L04b:sec:rootdense}

\begin{addedblock}
This section is an added extension, not in the handwritten notes. It is set in black for readability.
\end{addedblock}

Lecture~4 opened with two motivations: root lattices give good packings in low dimensions, and such
simple constructions do not extend to high dimensions. Here are some reasons for both.

\paragraph{1. The density is written in the diagram.} Normalize the roots to have norm $2$. Then the
Gram matrix of the simple roots is the Cartan matrix $C$, the shortest vectors are the roots, and
\[
  \dens(\text{root lattice})=\frac{\vol\Ball^n\big(1/\sqrt2\big)}{\sqrt{\det C}}
\]
(Interlude~A, Section~\ref{LIE:sec:readoff}). A dense lattice needs a large minimum and a small
covolume. Root lattices have both: the minimum is fixed at $\sqrt2$, and $\det C$ is tiny, namely
$n+1$ for $A_n$, $4$ for $D_n$, and $3,2,1$ for $E_6,E_7,E_8$. Against Rogers' simplex bound
(Section~\ref{L2:sec:rogers}):
\begin{center}
\renewcommand{\arraystretch}{1.15}
\begin{tabular}{clcccc}
\toprule
$n$ & lattice & $\det C$ & kissing number $\abs{\Delta}$ & density & density$/\sigma_n$\\
\midrule
$2$ & $A_2$ & $3$ & $6$ & $0.9069$ & $1.000$\\
$3$ & $A_3$ & $4$ & $12$ & $0.7405$ & $0.950$\\
$4$ & $D_4$ & $4$ & $24$ & $0.6169$ & $0.952$\\
$5$ & $D_5$ & $4$ & $40$ & $0.4653$ & $0.885$\\
$6$ & $E_6$ & $3$ & $72$ & $0.3729$ & $0.890$\\
$7$ & $E_7$ & $2$ & $126$ & $0.2953$ & $0.895$\\
$8$ & $E_8$ & $1$ & $240$ & $0.2537$ & $0.988$\\
\bottomrule
\end{tabular}
\end{center}
These are the densest lattices in their dimensions, and their kissing numbers $\abs{\Delta}$ are the
largest possible for lattices there. Among all arrangements of touching balls, $6,12,24,240$ are known to
be optimal in dimensions $2,3,4,8$; for $n=5,6,7$ the values $40,72,126$ are only the best known.

\paragraph{2. The roots are balanced: eutaxy from symmetry.} Let
$S=\sum_{\alpha\in\Delta}\alpha\alpha^T$. For any $w$ in the Weyl group, $w$ permutes $\Delta$, so
$wSw^{-1}=S$. For an irreducible root system the Weyl group acts irreducibly on $\R^n$, so each
eigenspace of the symmetric matrix $S$ is a nonzero $W$-invariant subspace and hence all of $\R^n$.
Thus $S=cI$, and taking traces gives $c=2\abs{\Delta}/n$:
\[
  \sum_{\alpha\in\Delta}\ip{\alpha}{x}^2=\frac{2\abs{\Delta}}{n}\,\abs{x}^2\qquad\text{for all }x.
\]
The minimal vectors pull equally in every direction (they form a \emph{spherical $2$-design}). In
Voronoi's language the lattice is \emph{strongly eutactic}.

\paragraph{3. The roots are plentiful: perfection.} A lattice is \emph{perfect} if its minimal vectors
determine its Gram matrix: the quadratic forms $x\mapsto\ip{\alpha}{x}^2$, $\alpha$ minimal, span the whole
$\frac{n(n+1)}2$-dimensional space of quadratic forms. For $A_n$, inside the hyperplane $\sum x_i=0$ of
$\R^{n+1}$, this is a pleasant exercise in graph Laplacians. A quadratic form on the hyperplane is the
same thing as a symmetric $(n+1)\times(n+1)$ matrix $Q$ with zero row sums, and every such matrix is a
weighted Laplacian,
\[
  Q=\sum_{i<j}(-q_{ij})\,(e_i-e_j)(e_i-e_j)^T,
\]
that is, a combination of the forms $(x_i-x_j)^2=\ip{e_i-e_j}{x}^2$ attached to the positive roots. So
$A_n$ is perfect, with exactly as many positive roots, $\frac{n(n+1)}2$, as the dimension requires.
$D_n$ ($n\ge4$), $E_6$, $E_7$ and $E_8$ are perfect too (Coxeter, \emph{Extreme forms}, 1951).

\begin{theorem}[Voronoi 1908]\label{L04b:thm:voronoi}
A lattice is \emph{extreme}, that is, a local maximum of the packing density among lattices, if and
only if it is perfect and eutactic. (The maximum is then strict up to similarity, since perfect lattices
are isolated.)
\end{theorem}

\noindent So \emph{every irreducible root lattice is a local optimum for lattice packing}. Symmetry
(irreducibility of $W$) gives eutaxy, and abundance of roots gives perfection. Voronoi also showed that
each dimension has only finitely many perfect forms up to similarity. There are $1,1,1,2,3,7,33,10916$ of
them in dimensions $1,\dots,8$ (the last count is due to Dutour Sikiri\'c, Sch\"urmann and Vallentin,
2007), so in principle the densest lattice can be found by a finite search. In dimension $2$ the only
perfect form is the hexagonal $A_2$, which recovers Lagrange's theorem from Lecture~1.

\paragraph{4. Why it stops working.} Two numbers are pinned down by the classification, and both are
fatal in high dimension.
\begin{itemize}
\item \emph{Too few roots.} Root lattices have $O(n^2)$ minimal vectors ($D_n$ has $2n(n-1)$), but
  the best kissing numbers grow exponentially in $n$. There is only a finite list of exceptional
  diagrams, so there is only a finite budget of exceptional symmetry.
\item \emph{An integral lattice cannot spread out.} The minimum stays at $\sqrt2$ while $\det C$ stays
  bounded, so Hermite's invariant $\min^2/\det^{1/n}=2/(\det C)^{1/n}$ stays near $2$. The best
  lattices have Hermite invariant growing linearly in $n$ (Minkowski--Hlawka; Interlude~B). Concretely,
  $\dens(D_n)=\vol\Ball^n(1/\sqrt2)/2$ decays like $(\pi e/n)^{n/2}$, up to polynomial factors. From $n=28$ on it is below even the trivial
  $2^{-n}$ of Lecture~1.
\end{itemize}

\paragraph{5. Filling the holes.} The first repair is to fill the deep holes, as in the construction of
$E_8$. The deep holes of $D_n$ include $h=(\frac12,\dots,\frac12)$, at distance $\sqrt n/2$ from the
lattice. For even $n\ge8$, $D_n^+=D_n\cup(D_n+h)$ is again a lattice with minimum $\sqrt2$, now of
covolume $1$, so
\[
  \dens(D_n^+)=\vol\Ball^n\big(1/\sqrt2\big)\qquad(n\ge8\text{ even}).
\]
At $n=8$ this is $E_8$, and the added balls exactly touch the old ones. Gluing (and its cousin,
lamination) is how many of the best known lattices up to dimension $24$ are built, culminating in the Leech
lattice. But it only postpones the decay: $\dens(D_{10}^+)\approx0.0797$, while the best known lattice in
dimension $10$, the laminated lattice $\Lambda_{10}$, has density about $0.0920$.

\section*{Projects and quiz}
\orig{admin (date, room) omitted}
Ideas for projects:
\begin{itemize}
  \item error correcting codes and information theory;
  \item other sphere packing bounds;
  \item entropy of jammed systems;
  \item translation of an older work on packings. \unclear{last item barely legible}
\end{itemize}
Format of the quiz: definitions, short answers, sketches.
% ===== end chapters/L04b.tex =====

% Interludes (new material): lettered so that lecture numbers stay aligned with chapter numbers
\begingroup\renewcommand{\thechapter}{A}\renewcommand{\theHchapter}{interludeA}
% ===== begin chapters/Lie.tex =====
% MAT670 -- Interlude after Lecture 6: from Dynkin diagrams to Lie algebras and Lie groups.
% Entirely new material (not in the handwritten notes). Labels prefixed LIE:.

\chapter[Interlude A: Lie algebras and Lie groups]{Interlude A: From Dynkin diagrams to Lie algebras and Lie groups}
\label{LIE:ch}

\begin{addedblock}
This interlude is new: it is not in the handwritten notes. It picks up the remark at the start of
Lecture~4, that the classification of root systems is ``half of the classification of complex
semisimple Lie algebras'', and follows it. To keep it readable it is set in black; all of it should be
read as added material.
\end{addedblock}

In Lectures~4--6 we took the root system axioms as given, classified the possible Dynkin diagrams, and
then built each root system by hand from lattices. Two questions were left hanging.
\begin{itemize}
\item Where do the axioms come from? Why should anyone care about a finite set of vectors closed under
  reflections, with the integrality condition $n_{\beta\alpha}\in\Z$?
\item The same root system $A_{n-1}$ appears twice in these notes, once as the vectors
  $e_i-e_j$ in the hyperplane $\sum x_i=0$ (Lecture~6) and, later, once as the walls $x_i=x_j$ of
  $\Conf(N,\R)$ (Lecture~8). Is there a single object behind both?
\end{itemize}
The answer to both is \emph{matrices}. In this interlude we find root systems inside matrix
algebras, see that the crystallographic axioms are forced by a small piece of linear algebra, and
then use the exponential map to pass from the algebra to groups of matrices. At the end the root
lattice comes back, this time sitting inside a group.

Throughout, $M_n$ denotes the $n\times n$ complex matrices.

%======================================================================
\section{The roots of \texorpdfstring{$A_{n-1}$}{A(n-1)} inside matrices}
\label{LIE:sec:sln}

\subsection*{Matrix units}
Let $E_{ij}$ be the matrix with a $1$ in position $(i,j)$ and zeros elsewhere. For $3\times3$ matrices,
for instance,
\[
  E_{12}=\begin{pmatrix}0&1&0\\0&0&0\\0&0&0\end{pmatrix},\qquad
  E_{23}=\begin{pmatrix}0&0&0\\0&0&1\\0&0&0\end{pmatrix},
\]
and multiplying them out,
\[
  E_{12}E_{23}=\begin{pmatrix}0&0&1\\0&0&0\\0&0&0\end{pmatrix}=E_{13},\qquad
  E_{23}E_{12}=0 .
\]
The pattern is easiest to see by watching what $E_{ij}$ does to the standard basis: it sends $e_j$ to
$e_i$ and kills every other $e_m$. Think of $E_{ij}$ as an \emph{arrow from $j$ to $i$}. A product
$E_{ij}E_{kl}$ first follows the arrow $l\to k$ and then the arrow $j\to i$. The arrows compose when
the head of the first meets the tail of the second ($k=j$), and otherwise the product is zero:
\begin{equation}\label{LIE:eq:units}
  E_{ij}E_{kl}=\delta_{jk}\,E_{il}.
\end{equation}
A diagonal matrix is a combination of the ``loops'' $E_{kk}$:
$h=\operatorname{diag}(h_1,\dots,h_n)=\sum_kh_kE_{kk}$. By \eqref{LIE:eq:units},
\[
  hE_{ij}=\sum_kh_kE_{kk}E_{ij}=h_iE_{ij},\qquad E_{ij}h=\sum_kh_kE_{ij}E_{kk}=h_jE_{ij}:
\]
multiplying by $h$ on the left scales row $i$, and on the right scales column $j$. With
$h=\operatorname{diag}(a,b,c)$:
\[
  hE_{12}=\begin{pmatrix}0&a&0\\0&0&0\\0&0&0\end{pmatrix},\qquad
  E_{12}h=\begin{pmatrix}0&b&0\\0&0&0\\0&0&0\end{pmatrix}.
\]

\subsection*{Symmetry first: conjugating by diagonal matrices}
What are the symmetries of the matrix algebra $M_n$? Any invertible $g$ gives a map
\[
  \operatorname{Ad}_g\colon M_n\to M_n,\qquad X\mapsto gXg^{-1},
\]
which respects both sums and products, since $g(XY)g^{-1}=(gXg^{-1})(gYg^{-1})$. This is the
\emph{adjoint action} of the group $\GL_n$ on its own matrices. (``Adjoint'' is used in the old sense
of ``attached'': the group acts on the space it lives in.) Each $\operatorname{Ad}_g$ is a linear map on
the $n^2$-dimensional vector space $M_n$. For a general $g$ it is hard to describe, but for a
\emph{diagonal} $t=\operatorname{diag}(t_1,\dots,t_n)$ it is as simple as possible. Using the two
formulas above (with $t$ and $t^{-1}=\operatorname{diag}(t_1^{-1},\dots,t_n^{-1})$ in place of $h$),
\[
  tE_{ij}t^{-1}=t_i\,t_j^{-1}\,E_{ij}.
\]
For example, with $t=\operatorname{diag}(t_1,t_2,t_3)$,
\[
  tE_{12}t^{-1}=
  \begin{pmatrix}t_1&&\\&t_2&\\&&t_3\end{pmatrix}
  \begin{pmatrix}0&1&0\\0&0&0\\0&0&0\end{pmatrix}
  \begin{pmatrix}t_1^{-1}&&\\&t_2^{-1}&\\&&t_3^{-1}\end{pmatrix}
  =\begin{pmatrix}0&t_1/t_2&0\\0&0&0\\0&0&0\end{pmatrix}=\frac{t_1}{t_2}\,E_{12}.
\]
So \emph{$\operatorname{Ad}_t$ is a diagonal linear map in the basis $\{E_{ij}\}$}, with eigenvalue
$t_i/t_j$ on $E_{ij}$. Better still, the same basis diagonalizes $\operatorname{Ad}_t$ for \emph{every}
diagonal $t$ at once. Writing the eigenvalue on $E_{ij}$ in position $(i,j)$, for $n=3$:
\[
  \operatorname{Ad}_t\ \text{has eigenvalues}\quad
  \begin{pmatrix}1&t_1/t_2&t_1/t_3\\ t_2/t_1&1&t_2/t_3\\ t_3/t_1&t_3/t_2&1\end{pmatrix}.
\]
The diagonal matrices themselves are fixed (eigenvalue $1$), since diagonal matrices commute.

\subsection*{Differentiating: the bracket appears}
Multiplicative eigenvalues like $t_1/t_2$ are awkward, so take logarithms. Put
\[
  t(s)=\operatorname{diag}(e^{sh_1},\dots,e^{sh_n}),
\]
a curve of diagonal matrices with $t(0)=I$ and velocity $t'(0)=h=\operatorname{diag}(h_1,\dots,h_n)$. Then
\[
  \operatorname{Ad}_{t(s)}E_{ij}=e^{s(h_i-h_j)}E_{ij}.
\]
Differentiate at $s=0$. On the right we get $(h_i-h_j)E_{ij}$. On the left, the product rule applied
to $t(s)Xt(s)^{-1}$, with $\frac{d}{ds}\big|_0t(s)^{-1}=-h$, gives for any matrix $X$
\[
  \frac{d}{ds}\Big|_{s=0}t(s)\,X\,t(s)^{-1}=hX-Xh .
\]
This is where the \emph{bracket} comes from. Define
\[
  \operatorname{ad}_h\colon M_n\to M_n,\qquad \operatorname{ad}_h(X)=[h,X]:=hX-Xh .
\]
Lower-case $\operatorname{ad}$ is the \emph{derivative} of upper-case $\operatorname{Ad}$: conjugation
seen infinitesimally. We did not invent the commutator; differentiating the symmetry produced it. (In
Section~\ref{LIE:sec:exp} we will integrate this back up: $\operatorname{Ad}_{e^X}=e^{\operatorname{ad}_X}$.)
Since $\operatorname{Ad}_{t(s)}$ is diagonal on each $E_{ij}$, so is its derivative:
\[
  \operatorname{ad}_h(E_{ij})=hE_{ij}-E_{ij}h=(h_i-h_j)\,E_{ij},
\]
which you can also read off the two $3\times3$ matrices above: $hE_{12}-E_{12}h=(a-b)E_{12}$.

\subsection*{The eigenvalues are the roots}
Now regard the eigenvalue as a function of $h$. On $E_{ij}$ it is the linear functional
$h\mapsto h_i-h_j$. If we identify functionals with vectors using the standard inner product, so that
$h\mapsto h_i$ corresponds to $e_i$, then the eigenvalue on $E_{ij}$ is the vector $e_i-e_j$. So:
\begin{center}
\emph{the nonzero eigenvalues of $\operatorname{ad}_h$, as $h$ runs over diagonal matrices, are exactly
the roots $e_i-e_j$ of $A_{n-1}$.}
\end{center}
For $n=3$, with $\alpha_1=e_1-e_2$ and $\alpha_2=e_2-e_3$, write the eigenvalue of $E_{ij}$ in position
$(i,j)$. The root system $A_2$ appears in the shape of a $3\times 3$ matrix, with positive roots above the
diagonal and negative roots below:
\[
  \begin{pmatrix}0&e_1-e_2&e_1-e_3\\ e_2-e_1&0&e_2-e_3\\ e_3-e_1&e_3-e_2&0\end{pmatrix}
  =\begin{pmatrix}0&\alpha_1&\alpha_1+\alpha_2\\ -\alpha_1&0&\alpha_2\\ -\alpha_1-\alpha_2&-\alpha_2&0\end{pmatrix}.
\]

Where is the hyperplane $\sum x_i=0$? Scalar matrices $h=cI$ commute with everything, so they tell us
nothing, and it is natural to discard them. That leaves the \emph{traceless} matrices
\[
  \mathfrak{sl}_n=\{X\in M_n:\operatorname{tr}X=0\},
\]
whose diagonal part
\[
  \mathfrak h=\{\operatorname{diag}(h_1,\dots,h_n): h_1+\dots+h_n=0\}
\]
is exactly the hyperplane $\sum x_i=0$ of Lecture~6, of dimension $n-1$. Every traceless matrix splits
into its diagonal part and its off-diagonal entries. For $n=3$:
\[
  \begin{pmatrix}a&x_{12}&x_{13}\\x_{21}&b&x_{23}\\x_{31}&x_{32}&c\end{pmatrix}
  =\underbrace{\operatorname{diag}(a,b,c)}_{\in\,\mathfrak h,\ a+b+c=0}
  +\,x_{12}E_{12}+x_{23}E_{23}+x_{13}E_{13}+x_{21}E_{21}+x_{32}E_{32}+x_{31}E_{31},
\]
that is,
\[
  \mathfrak{sl}_3=\mathfrak h\oplus\underbrace{\C E_{12}}_{\alpha_1}\oplus\underbrace{\C E_{23}}_{\alpha_2}
  \oplus\underbrace{\C E_{13}}_{\alpha_1+\alpha_2}\oplus\underbrace{\C E_{21}}_{-\alpha_1}
  \oplus\underbrace{\C E_{32}}_{-\alpha_2}\oplus\underbrace{\C E_{31}}_{-\alpha_1-\alpha_2},
  \qquad 8=2+6.
\]
In general we get the decomposition into eigenspaces of $\operatorname{ad}\mathfrak h$,
\begin{equation}\label{LIE:eq:rootdecomp}
  \mathfrak{sl}_n=\mathfrak h\ \oplus\ \bigoplus_{\alpha\in\Delta}\mathfrak g_\alpha,
  \qquad
  \mathfrak g_{e_i-e_j}=\C E_{ij},
\end{equation}
where $\Delta=\{e_i-e_j: i\neq j\}$ is the root system $A_{n-1}$ and $\mathfrak g_\alpha$ is the space on
which every $\operatorname{ad}_h$ acts by the scalar $\alpha(h)$. As a check on dimensions:
$n^2-1=(n-1)+n(n-1)$, which is (rank) $+$ (number of roots), using the count
$\abs{\Delta}=n(n-1)$ from Lecture~6.

\begin{example}[$\mathfrak{sl}_3$]\label{LIE:ex:sl3}
Here $\mathfrak h$ is two-dimensional and the six off-diagonal matrix units sit at the six roots of
$A_2$. Placing each $E_{ij}$ at its root gives the familiar hexagon:
\begin{center}
\begin{tikzpicture}[scale=1.55]
  \foreach \a/\lab/\pos in {0/E_{12}/right, 60/E_{13}/above right, 120/E_{23}/above left,
                            180/E_{21}/left, 240/E_{31}/below left, 300/E_{32}/below right}
    {\draw[-{Stealth}] (0,0)--(\a:1); \fill (\a:1) circle (0.03);
     \node[\pos,font=\small] at (\a:1) {$\lab$};}
  \fill (0,0) circle (0.045); \fill (0.09,0) circle (0.045);
  \node[below,font=\small] at (0.05,-0.08) {$\mathfrak h$};
  \draw[dashed,-{Stealth},addcol,thick] (120:1.05) to[bend left=35] node[above,font=\scriptsize]{$\operatorname{ad}_{E_{12}}$} (60:1.05);
  \node[font=\scriptsize,gray] at (2.6,0.55) {$\alpha_1=e_1-e_2$};
  \node[font=\scriptsize,gray] at (2.6,0.25) {$\alpha_2=e_2-e_3$};
  \node[font=\scriptsize,gray] at (2.6,-0.05) {$\alpha_1+\alpha_2=e_1-e_3$};
\end{tikzpicture}
\end{center}
The two dots at the centre stand for the two-dimensional zero eigenspace $\mathfrak h$. The dashed
arrow anticipates Section~\ref{LIE:sec:bracket}: bracketing with $E_{12}$ moves $E_{23}$ to $E_{13}$,
that is, it \emph{adds the root} $\alpha_1$.
\end{example}

The choices made in Lecture~4 also have matrix meanings. Take the direction $d=(n,n-1,\dots,1)$, so
that $\ip{e_i-e_j}{d}=j-i$. Then:
\begin{itemize}
\item the positive roots $e_i-e_j$ with $i<j$ are the matrix units \emph{above} the diagonal, and
  the negative roots are those below it;
\item the simple roots $\alpha_i=e_i-e_{i+1}$ are the entries just above the diagonal,
  $E_{12},E_{23},\dots,E_{n-1,n}$;
\item every positive root is a sum of simple roots: $e_i-e_j=\alpha_i+\alpha_{i+1}+\dots+\alpha_{j-1}$.
  On the matrix side this is the statement that $E_{ij}$ is an iterated product (indeed an iterated
  commutator) of superdiagonal units.
\end{itemize}
The Weyl group, generated by the reflections $s_{e_i-e_j}$, acts on $\mathfrak h$ by swapping $h_i$ and
$h_j$. It is the symmetric group $\Sigma_n$ permuting the diagonal entries, which is conjugation by
permutation matrices. The Weyl chamber containing $d$ is $\{h_1>h_2>\dots>h_n\}$, one of the $n!$
chambers of $\Conf(n,\R)$ from Lecture~8.

%======================================================================
\section{The bracket and the definition of a Lie algebra}
\label{LIE:sec:bracket}

In Section~\ref{LIE:sec:sln} the bracket appeared as the derivative of conjugation by diagonal
matrices. Nothing there needed $h$ to be diagonal: if $g(s)$ is any curve of invertible matrices with
$g(0)=I$ and velocity $g'(0)=X$, the same product-rule computation gives
\[
  \frac{d}{ds}\Big|_{s=0}g(s)\,Y\,g(s)^{-1}=XY-YX=[X,Y]=\operatorname{ad}_X(Y).
\]
So every matrix $X$ has its linear map $\operatorname{ad}_X\colon M_n\to M_n$, the infinitesimal
conjugation in the direction $X$. Diagonal $X$ gave diagonal $\operatorname{ad}_X$; for a matrix unit,
for instance, $\operatorname{ad}_{E_{12}}$ is nilpotent and moves matrix units around, as in
$\operatorname{ad}_{E_{12}}(E_{23})=E_{12}E_{23}-E_{23}E_{12}=E_{13}$ by \eqref{LIE:eq:units}.

The operation $[X,Y]=XY-YX$ is the \emph{commutator}. It has three properties that are easy to
check:
\begin{enumerate}[label=(\roman*)]
\item it is bilinear;
\item it is antisymmetric, $[X,Y]=-[Y,X]$;
\item it satisfies the \emph{Jacobi identity}
  \[
    [X,[Y,Z]]=[[X,Y],Z]+[Y,[X,Z]],
  \]
  which says that $\operatorname{ad}_X=[X,\cdot\,]$ obeys the product rule (it is a derivation of the
  bracket). Expanding both sides into products of three matrices proves it.
\end{enumerate}

\begin{definition}\label{LIE:def:liealg}
A \emph{Lie algebra} over a field $k$ is a vector space $\mathfrak g$ with a bilinear map
$[\cdot,\cdot]\colon\mathfrak g\times\mathfrak g\to\mathfrak g$ that is antisymmetric and satisfies the
Jacobi identity. A \emph{subalgebra} is a subspace closed under the bracket; an \emph{ideal} is a
subspace $I$ with $[\mathfrak g,I]\subseteq I$. A Lie algebra is \emph{simple} if it is not abelian and
has no ideals other than $0$ and $\mathfrak g$. A \emph{representation} of $\mathfrak g$ on a vector space
$V$ is a linear map $\rho\colon\mathfrak g\to\operatorname{End}(V)$ with
$\rho([X,Y])=\rho(X)\rho(Y)-\rho(Y)\rho(X)$.
\end{definition}

\begin{example}\leavevmode
\begin{enumerate}[nosep]
\item $\mathfrak{gl}_n=M_n$ with the commutator.
\item $\mathfrak{sl}_n$ is a subalgebra, since $\operatorname{tr}(XY-YX)=0$ for all $X,Y$. It is simple for
  $n\ge2$.
\item $\mathfrak{so}_n=\{X\in M_n: X^T=-X\}$ is a subalgebra: if $X,Y$ are antisymmetric, then
  $[X,Y]^T=Y^TX^T-X^TY^T=YX-XY=-[X,Y]$.
\item $\R^3$ with the cross product is a Lie algebra (Jacobi is the identity
  $a\times(b\times c)=(a\times b)\times c+b\times(a\times c)$). Sending $\omega\in\R^3$ to the matrix
  $X_\omega$ with $X_\omega v=\omega\times v$ gives an isomorphism $(\R^3,\times)\cong\mathfrak{so}_3(\R)$.
\item By the Jacobi identity, $X\mapsto\operatorname{ad}_X$ is a representation of $\mathfrak g$ on itself,
  the \emph{adjoint representation}. Section~\ref{LIE:sec:sln} was a study of the adjoint
  representation of $\mathfrak{sl}_n$, restricted to $\mathfrak h$.
\end{enumerate}
\end{example}

The Jacobi identity also explains the dashed arrow in Example~\ref{LIE:ex:sl3}. Suppose
$[h,x]=\alpha(h)x$ and $[h,y]=\beta(h)y$ for all $h\in\mathfrak h$. Then
\[
  [h,[x,y]]=[[h,x],y]+[x,[h,y]]=\big(\alpha(h)+\beta(h)\big)[x,y].
\]

\begin{lemma}\label{LIE:lem:add}
$[\mathfrak g_\alpha,\mathfrak g_\beta]\subseteq\mathfrak g_{\alpha+\beta}$, where
$\mathfrak g_{\alpha+\beta}=0$ if $\alpha+\beta$ is neither a root nor $0$, and $\mathfrak g_0=\mathfrak h$.
\end{lemma}

So the bracket moves around the root diagram by \emph{adding roots}. In $\mathfrak{sl}_3$,
$[E_{12},E_{23}]=E_{13}$ is the statement $\alpha_1+\alpha_2=e_1-e_3$, while $[E_{12},E_{13}]=0$
because $2e_1-e_2-e_3$ is not a root.

%======================================================================
\section{The atom: \texorpdfstring{$\mathfrak{sl}_2$}{sl(2)}}
\label{LIE:sec:sl2}

Fix a root $\alpha=e_i-e_j$ of $A_{n-1}$ and put
\[
  e_\alpha=E_{ij},\qquad f_\alpha=E_{ji},\qquad h_\alpha=[e_\alpha,f_\alpha]=E_{ii}-E_{jj}\in\mathfrak h.
\]
(The middle equality is \eqref{LIE:eq:units}: $E_{ij}E_{ji}-E_{ji}E_{ij}=E_{ii}-E_{jj}$.) The smallest
case is $n=2$, where
\[
  e=\begin{pmatrix}0&1\\0&0\end{pmatrix},\quad
  f=\begin{pmatrix}0&0\\1&0\end{pmatrix},\quad
  h=\begin{pmatrix}1&0\\0&-1\end{pmatrix}
\]
are a basis of $\mathfrak{sl}_2$. Multiplying out,
\[
  he=\begin{pmatrix}0&1\\0&0\end{pmatrix},\quad eh=\begin{pmatrix}0&-1\\0&0\end{pmatrix},\quad
  ef=\begin{pmatrix}1&0\\0&0\end{pmatrix},\quad fe=\begin{pmatrix}0&0\\0&1\end{pmatrix},
\]
so $[h,e]=2e$, $[e,f]=h$, and similarly $[h,f]=-2f$. The same computation with
\eqref{LIE:eq:units} works for any root:
\begin{equation}\label{LIE:eq:sl2rel}
  [h_\alpha,e_\alpha]=2e_\alpha,\qquad [h_\alpha,f_\alpha]=-2f_\alpha,\qquad [e_\alpha,f_\alpha]=h_\alpha .
\end{equation}
(The first is $\operatorname{ad}_{h_\alpha}(E_{ij})=(h_i-h_j)E_{ij}$ with $h_i=1$, $h_j=-1$.) So every root $\alpha$ comes with a copy
$\mathfrak s_\alpha=\operatorname{span}(e_\alpha,f_\alpha,h_\alpha)\cong\mathfrak{sl}_2$ sitting inside $\mathfrak{sl}_n$.

Now look at how $h_\alpha$ acts on the other root spaces. For a root $\beta=e_k-e_l$,
\[
  \beta(h_\alpha)=(\delta_{ki}-\delta_{kj})-(\delta_{li}-\delta_{lj})=\ip{\beta}{\alpha}
  =\frac{2\ip{\beta}{\alpha}}{\ip{\alpha}{\alpha}}=n_{\beta\alpha},
\]
since $\ip{\alpha}{\alpha}=2$. So \emph{the Cartan integer $n_{\beta\alpha}$ is the eigenvalue of
$h_\alpha$ on $\mathfrak g_\beta$}. ($h_\alpha$ is called the \emph{coroot} of $\alpha$.) To see why these
eigenvalues have to be integers, we need one theorem about $\mathfrak{sl}_2$.

\begin{theorem}[Representations of $\mathfrak{sl}_2$]\label{LIE:thm:sl2}
Let $V$ be a finite-dimensional complex representation of $\mathfrak{sl}_2$. Then $h$ acts
diagonalizably on $V$, its eigenvalues are integers, and for every $k\in\Z$ the eigenspaces satisfy
$\dim V_k=\dim V_{-k}$. If $V$ is irreducible of dimension $m+1$, then the eigenvalues of $h$ are
$m,m-2,\dots,-m$, each with multiplicity one.
\end{theorem}

\begin{proof}[Proof for irreducible $V$]
The key observation is that $e$ raises $h$-eigenvalues by $2$ and $f$ lowers them by $2$. If
$hv=\lambda v$, then by \eqref{LIE:eq:sl2rel}
\[
  h(ev)=e(hv)+[h,e]v=(\lambda+2)\,ev,\qquad h(fv)=(\lambda-2)\,fv .
\]
Take any eigenvector of $h$ and apply $e$ repeatedly. The vectors obtained are eigenvectors for
distinct eigenvalues while they are nonzero, and $V$ is finite-dimensional, so eventually we reach a
\emph{highest weight vector}: $v\neq0$ with $hv=\lambda v$ and $ev=0$. Put $v_k=f^kv$. Then
$hv_k=(\lambda-2k)v_k$, and induction on $k$ using $ef=fe+h$ gives
\[
  e\,v_k=k(\lambda-k+1)\,v_{k-1}.
\]
Let $m$ be the largest index with $v_m\neq0$. Then $0=e\,v_{m+1}=(m+1)(\lambda-m)\,v_m$, so
$\lambda=m$. In particular $\lambda$ is a non-negative integer. The span of $v_0,\dots,v_m$ is closed
under $e$, $f$ and $h$, so by irreducibility it is all of $V$. The eigenvalues are therefore
$m,m-2,\dots,-m$, as claimed.
\end{proof}

\begin{figure}[ht]
\centering
\begin{tikzpicture}[xscale=1.7]
  \foreach \x/\l in {0/{-m},1/{-m+2},2/{\cdots},3/{m-2},4/{m}} {\fill (\x,0) circle (0.05); \node[font=\small] at (\x,-0.62) {$\l$};}
  \foreach \x in {0,1,2,3} {\draw[-{Stealth},addcol] (\x+0.1,0.1) to[bend left=40] (\x+0.9,0.1);
                            \draw[-{Stealth},flagcol] (\x+0.9,-0.1) to[bend left=40] (\x+0.1,-0.1);}
  \node[addcol,font=\small] at (2,0.7) {$e$ raises by $2$};
  \node[flagcol,font=\small] at (2,-1.05) {$f$ lowers by $2$};
\end{tikzpicture}
\caption{The $(m+1)$-dimensional irreducible representation of $\mathfrak{sl}_2$: a string of
$h$-eigenvalues, symmetric about $0$.}
\end{figure}

For a general finite-dimensional $V$ one uses \emph{Weyl's complete reducibility theorem}: $V$ is a
direct sum of irreducible representations (see Humphreys, \emph{Introduction to Lie Algebras and
Representation Theory}, \S6.3, or Hall, \emph{Lie Groups, Lie Algebras, and Representations}, Chs.~4 and~10).
The statements about eigenvalues then follow from the irreducible case.

\begin{example}[An $\mathfrak{sl}_2$ inside $\mathfrak{sl}_3$]\label{LIE:ex:string}
Take $\alpha=\alpha_1=e_1-e_2$ in $\mathfrak{sl}_3$, so $e_\alpha=E_{12}$, $f_\alpha=E_{21}$ and
$h_\alpha=\operatorname{diag}(1,-1,0)$. Let it act by brackets on the two matrix units $E_{23}$ (root
$\alpha_2$) and $E_{13}$ (root $\alpha_1+\alpha_2$). By \eqref{LIE:eq:units},
\[
  [E_{12},E_{23}]=E_{13},\quad [E_{12},E_{13}]=0,\quad [E_{21},E_{13}]=E_{23},\quad [E_{21},E_{23}]=0,
\]
and $h_\alpha$ acts on $E_{23}$ by $h_2-h_3=-1$ and on $E_{13}$ by $h_1-h_3=+1$. So
$\operatorname{span}(E_{23},E_{13})$ is the two-dimensional irreducible representation of $\mathfrak{sl}_2$
($m=1$): a string of eigenvalues $-1,+1$, with $E_{13}$ the highest weight vector. The eigenvalue $-1$
on $E_{23}$ is the Cartan integer $n_{\alpha_2\alpha_1}=2\ip{\alpha_2}{\alpha_1}/\ip{\alpha_1}{\alpha_1}=-1$,
and the string $\alpha_2,\ \alpha_2+\alpha_1$ runs from $\alpha_2$ to its mirror image
$s_{\alpha_1}(\alpha_2)=\alpha_2+\alpha_1$. That is the reflection axiom happening in front of us.
\end{example}

%======================================================================
\section{Why the root system axioms hold}
\label{LIE:sec:axioms}

Now apply Theorem~\ref{LIE:thm:sl2} to the copy $\mathfrak s_\alpha\cong\mathfrak{sl}_2$, acting by
brackets. Fix roots $\beta\neq\pm\alpha$ and form the \emph{$\alpha$-string through $\beta$},
\[
  V=\bigoplus_{k\in\Z}\mathfrak g_{\beta+k\alpha}.
\]
By Lemma~\ref{LIE:lem:add}, bracketing with $e_\alpha$ or $f_\alpha$ shifts $k$ by $\pm1$, so $V$ is a
representation of $\mathfrak s_\alpha$. On the summand $\mathfrak g_{\beta+k\alpha}$, $h_\alpha$ acts by the
scalar
\[
  (\beta+k\alpha)(h_\alpha)=n_{\beta\alpha}+2k .
\]

\begin{itemize}
\item \textbf{Axiom 4 (crystallographic).} $n_{\beta\alpha}$ is an eigenvalue of $h_\alpha$ on $V$
  (take $k=0$), so by Theorem~\ref{LIE:thm:sl2} it is an integer.
\item \textbf{Axiom 3 (reflections).} The eigenvalues are symmetric about $0$, so $-n_{\beta\alpha}$ is
  also an eigenvalue: $n_{\beta\alpha}+2k=-n_{\beta\alpha}$ for some $k$ with
  $\mathfrak g_{\beta+k\alpha}\neq0$. That is, $k=-n_{\beta\alpha}$, and
  \[
    \beta-n_{\beta\alpha}\,\alpha=s_\alpha(\beta)
  \]
  is a root. The reflection $s_\alpha$ maps $\Delta$ to itself.
\item \textbf{Strings have no gaps.} Each irreducible piece of $V$ has eigenvalues forming an unbroken
  string with step $2$. In fact $V$ is irreducible here, since the root spaces turn out to be
  one-dimensional. So the roots $\beta+k\alpha$ form an unbroken string $\beta-p\alpha,\dots,\beta+q\alpha$
  with $p-q=n_{\beta\alpha}$. This is the ``$\alpha-\beta$'' phenomenon behind Lemma~\ref{L04a:lem:diff} of Lecture~4.
\end{itemize}

For $\mathfrak{sl}_n$ all of this can of course be checked directly. The point is that the argument used
only three ingredients: a commutative subalgebra $\mathfrak h$ acting diagonalizably, the copies of
$\mathfrak{sl}_2$ attached to the roots, and finite-dimensionality. These are exactly what a
\emph{complex semisimple Lie algebra} provides. In such an algebra a maximal commutative subalgebra of
diagonalizable elements (a \emph{Cartan subalgebra}) always exists, and its nonzero eigenvalues form a
reduced, crystallographic root system. (Reducedness, axiom~2, takes one more argument; see Humphreys
\S8.4.) \emph{The root system axioms are the shadow of $\mathfrak{sl}_2$ representation theory.}

%======================================================================
\section{The other classical families}
\label{LIE:sec:classical}

The families $B_n$, $C_n$, $D_n$ from Lecture~6 come from the Lie algebras that preserve a bilinear
form. Let $J$ be an invertible $N\times N$ matrix and put
\[
  \mathfrak g_J=\{X\in M_N: X^TJ+JX=0\}.
\]
This is a subalgebra (the same computation as for $\mathfrak{so}_n$), and in Section~\ref{LIE:sec:exp}
we will see that it consists of the \emph{infinitesimal symmetries} of the form $\ip{v}{w}_J=v^TJw$.
For $J=I$ this is $\mathfrak{so}_N$.

The trick is to choose $J$ so that a large family of \emph{diagonal} matrices lies in $\mathfrak g_J$.
Write $i'=N+1-i$ and let $J$ be the antidiagonal matrix with $J_{i,i'}=1$. For $X\in M_N$ one computes
$(JX)_{ij}=X_{i'j}$ and $(X^TJ)_{ij}=X_{j'i}$, so
\[
  X\in\mathfrak g_J\iff X_{ij}=-X_{j'i'}\ \text{for all } i,j .
\]
(Over $\C$ any two nondegenerate symmetric forms are equivalent, so $\mathfrak g_J\cong\mathfrak{so}_N(\C)$.)
\begin{itemize}
\item The diagonal matrices in $\mathfrak g_J$ are
  $h=\operatorname{diag}(h_1,\dots,h_n,[0],-h_n,\dots,-h_1)$, where the middle $0$ is present only when
  $N=2n+1$ is odd. These form a commutative subalgebra $\mathfrak h$ of dimension $n$.
\item The matrices $E_{ij}-E_{j'i'}$, for pairs with $i\neq j$ and $j\neq i'$, span the rest of
  $\mathfrak g_J$ (for $j=i'$ the expression is $0$: the antidiagonal of $X$ vanishes). Each is
  an eigenvector of $\operatorname{ad}_h$ with eigenvalue $\tilde h_i-\tilde h_j$, where $\tilde h$ is the
  diagonal of $h$. (The two matrix units give the same eigenvalue, because $\tilde h_{k'}=-\tilde h_k$.)
\item Reading off the eigenvalues: $i,j\le n$ gives $\pm(e_i-e_j)$; $i\le n<j$ gives
  $\pm(e_i+e_k)$ with $j=k'$ and $k\neq i$ (the case $k=i$ is excluded since then $j=i'$); and when $N$
  is odd the middle index gives $\pm e_i$.
\end{itemize}
So $\mathfrak{so}_{2n}$ has root system $\{\pm e_i\pm e_j\}=D_n$ and $\mathfrak{so}_{2n+1}$ has
$\{\pm e_i\pm e_j\}\cup\{\pm e_i\}=B_n$. If instead
$J=\left(\begin{smallmatrix}0&K\\-K&0\end{smallmatrix}\right)$, with $K$ the $n\times n$ antidiagonal
matrix of ones, the form is \emph{symplectic}. Now signs enter: with $\varepsilon_i=1$ for $i\le n$
and $\varepsilon_i=-1$ for $i>n$, the condition becomes $X_{ij}=-\varepsilon_i\varepsilon_jX_{j'i'}$, the
diagonal part is the same $\mathfrak h$ as before, and the off-diagonal basis is
$E_{ij}-\varepsilon_i\varepsilon_jE_{j'i'}$. The same bookkeeping gives the symplectic algebra
$\mathfrak{sp}_{2n}$, whose root system is $\{\pm e_i\pm e_j\}\cup\{\pm 2e_i\}=C_n$: now the case $k=i$ is
allowed, since for $j=i'$ the basis element is $2E_{ii'}\neq0$, and it produces the long roots $\pm2e_i$.

\begin{center}
\renewcommand{\arraystretch}{1.25}
\begin{tabular}{llllll}
\toprule
type & Lie algebra & condition & roots & rank $+\abs{\Delta}$ & $\dim\mathfrak g$\\
\midrule
$A_n$ & $\mathfrak{sl}_{n+1}$ & $\operatorname{tr}X=0$ & $e_i-e_j$ & $n+n(n+1)$ & $(n+1)^2-1$\\
$B_n$ & $\mathfrak{so}_{2n+1}$ & $X^TJ+JX=0$ & $\pm e_i\pm e_j,\ \pm e_i$ & $n+2n^2$ & $n(2n+1)$\\
$C_n$ & $\mathfrak{sp}_{2n}$ & $X^TJ+JX=0$ ($J$ skew) & $\pm e_i\pm e_j,\ \pm2e_i$ & $n+2n^2$ & $n(2n+1)$\\
$D_n$ & $\mathfrak{so}_{2n}$ & $X^TJ+JX=0$ & $\pm e_i\pm e_j$ & $n+2n(n-1)$ & $n(2n-1)$\\
\bottomrule
\end{tabular}
\end{center}
In every row, $\dim\mathfrak g=\operatorname{rank}+\abs{\Delta}$, as it must be by the root decomposition
\eqref{LIE:eq:rootdecomp}. The root counts are the ones computed in Lecture~6.

\begin{remark}[The ``repetitions'' in the classification]
Theorem~\ref{L04a:thm:classification} restricted the ranges of $n$ ($B_n$ for $n\ge2$, $C_n$ for $n\ge3$, $D_n$ for $n\ge4$) to
avoid repeating diagrams. Each coincidence of small diagrams is an isomorphism of Lie algebras:
$B_1=C_1=A_1$ gives $\mathfrak{so}_3\cong\mathfrak{sp}_2=\mathfrak{sl}_2$; $B_2=C_2$ gives
$\mathfrak{so}_5\cong\mathfrak{sp}_4$; $D_2=A_1\times A_1$ gives $\mathfrak{so}_4\cong\mathfrak{sl}_2\times\mathfrak{sl}_2$; and
$D_3=A_3$ gives $\mathfrak{so}_6\cong\mathfrak{sl}_4$.
\end{remark}

%======================================================================
\section{The classification of simple Lie algebras}
\label{LIE:sec:classification}

Everything so far points to a dictionary between Dynkin diagrams and Lie algebras. The result is
due to Killing (1888--90) and \'E.~Cartan (thesis, 1894); the diagrams themselves are due to
Dynkin (1946).

\begin{theorem}[Killing--Cartan classification]\label{LIE:thm:KC}
Let $\mathfrak g$ be a complex semisimple Lie algebra.
\begin{enumerate}[nosep]
\item $\mathfrak g$ has a Cartan subalgebra $\mathfrak h$, unique up to automorphisms of $\mathfrak g$, and the
  nonzero eigenvalues of $\operatorname{ad}\mathfrak h$ form a root system $\Delta(\mathfrak g)$ in (the real
  span of) $\mathfrak h^*$, with one-dimensional root spaces.
\item $\mathfrak g$ is simple if and only if $\Delta(\mathfrak g)$ is irreducible.
\item $\mathfrak g\cong\mathfrak g'$ if and only if $\Delta(\mathfrak g)\cong\Delta(\mathfrak g')$.
\item (Existence; Serre's theorem) Every root system arises this way.
\end{enumerate}
Consequently the complex simple Lie algebras correspond exactly to the connected Dynkin diagrams
$A_n$ ($n\ge1$), $B_n$ ($n\ge2$), $C_n$ ($n\ge3$), $D_n$ ($n\ge4$), $E_6$, $E_7$, $E_8$, $F_4$, $G_2$.
\end{theorem}

Serre's theorem even writes the algebra down from the diagram. Let $\alpha_1,\dots,\alpha_r$ be the
simple roots and $a_{ij}=n_{\alpha_j\alpha_i}$ the \emph{Cartan matrix}. Then $\mathfrak g$ is generated by
$e_i,f_i,h_i$ ($1\le i\le r$), which are exactly the $\mathfrak{sl}_2$-triples attached to the simple
roots, subject to the relations
\[
  [h_i,h_j]=0,\quad [h_i,e_j]=a_{ij}e_j,\quad [h_i,f_j]=-a_{ij}f_j,\quad [e_i,f_j]=\delta_{ij}h_i,
\]
\[
  (\operatorname{ad}e_i)^{1-a_{ij}}(e_j)=0,\qquad (\operatorname{ad}f_i)^{1-a_{ij}}(f_j)=0\qquad(i\neq j).
\]
The last two relations say that an $\alpha_i$-string through $\alpha_j$ has exactly the length the
Cartan integer predicts. For $\mathfrak{sl}_n$ this recovers $e_i=E_{i,i+1}$, $f_i=E_{i+1,i}$ and
$h_i=E_{ii}-E_{i+1,i+1}$.

The dimension formula $\dim\mathfrak g=\operatorname{rank}+\abs{\Delta}$ and the root counts ($12$ for
$G_2$ from Lecture~4; $48$, $126$, $240$ from Lecture~6; and $72$ for $E_6$, which one checks by counting
the roots of $E_8$ orthogonal to an $A_2$) give the dimensions of the exceptional algebras:
\begin{center}
\renewcommand{\arraystretch}{1.2}
\begin{tabular}{lccccc}
\toprule
 & $G_2$ & $F_4$ & $E_6$ & $E_7$ & $E_8$\\
\midrule
rank & 2 & 4 & 6 & 7 & 8\\
$\abs{\Delta}$ & 12 & 48 & 72 & 126 & 240\\
$\dim\mathfrak g$ & 14 & 52 & 78 & 133 & 248\\
\bottomrule
\end{tabular}
\end{center}
In particular, the $240$ minimal vectors of the $E_8$ lattice, the kissing configuration in
dimension $8$, \emph{literally} label the root spaces of the largest exceptional Lie algebra,
$\mathfrak e_8=\mathfrak h\oplus\bigoplus_{\alpha\in E_8}\mathfrak g_\alpha$, $248=8+240$. The smallest one,
$\mathfrak g_2$, is the algebra of derivations of the octonions (\'E.~Cartan).

%======================================================================
\section{From the algebra to the group: the exponential map}
\label{LIE:sec:exp}

The name ``Lie algebra'' comes from Sophus Lie's idea that a group of transformations is determined,
near the identity, by its \emph{infinitesimal} transformations. For matrices this can be made
completely explicit.

\begin{definition}
For $X\in M_n$ (real or complex) the \emph{matrix exponential} is
\[
  \exp X=e^X=\sum_{k=0}^\infty\frac{X^k}{k!}.
\]
The series converges absolutely, since $\norm{X^k}\le\norm{X}^k$ in the operator norm.
\end{definition}

\begin{proposition}\label{LIE:prop:exp}
For $X,Y\in M_n$ and $g\in\GL_n$:
\begin{enumerate}[nosep]
\item if $XY=YX$, then $e^{X+Y}=e^Xe^Y$; in particular $e^{(s+t)X}=e^{sX}e^{tX}$ and $(e^X)^{-1}=e^{-X}$;
\item $\frac{d}{dt}e^{tX}=Xe^{tX}=e^{tX}X$;
\item $e^{gXg^{-1}}=g\,e^Xg^{-1}$ and $(e^X)^T=e^{X^T}$;
\item $\det e^X=e^{\operatorname{tr}X}$;
\item $\exp$ maps a neighbourhood of $0$ diffeomorphically onto a neighbourhood of $I$, with inverse
  $\log(I+A)=\sum_{k\ge1}(-1)^{k+1}A^k/k$ for $\norm{A}<1$.
\end{enumerate}
\end{proposition}

\begin{proof}[Sketch]
(1) For commuting $X$ and $Y$ the binomial theorem applies to $(X+Y)^k$, and the Cauchy product of the
two series gives the result. (2) Differentiate the series term by term. (3) Apply the identities
$(gXg^{-1})^k=gX^kg^{-1}$ and $(X^k)^T=(X^T)^k$ to each term. (4) Over $\C$, write $X=gTg^{-1}$ with
$T$ upper triangular (Schur). Then $e^T$ is upper triangular with diagonal entries $e^{t_{ii}}$, so
$\det e^X=\det e^T=e^{\sum t_{ii}}=e^{\operatorname{tr}X}$. (5) The derivative of $\exp$ at $0$ is the
identity, so the inverse function theorem applies; the series for $\log$ inverts $\exp$ near $I$ as a
formal identity of power series.
\end{proof}

By (1), $t\mapsto e^{tX}$ is a homomorphism $(\R,+)\to\GL_n$: a \emph{one-parameter subgroup}. This is
the matrix version of the one-parameter groups of diffeomorphisms of Lecture~9. The linear vector field $v\mapsto Xv$ on $\R^n$ generates the
one-parameter group of diffeomorphisms $\varphi_t(v)=e^{tX}v$, because $\frac{d}{dt}e^{tX}v=X(e^{tX}v)$
by (2). So $X$ is the ``velocity'' of the subgroup at the identity.

This also closes the loop on the names $\operatorname{Ad}$ and $\operatorname{ad}$ from
Section~\ref{LIE:sec:sln}.

\begin{proposition}\label{LIE:prop:Adad}
For all $X,Y\in M_n$, \ $e^{X}Ye^{-X}=e^{\operatorname{ad}_X}(Y)=Y+[X,Y]+\tfrac1{2!}[X,[X,Y]]+\cdots$,
that is, $\operatorname{Ad}_{e^X}=e^{\operatorname{ad}_X}$.
\end{proposition}

\begin{proof}
Both $A(t)=\operatorname{Ad}_{e^{tX}}$ and $B(t)=e^{t\operatorname{ad}_X}$ are curves of linear maps on $M_n$
with value $\mathrm{id}$ at $t=0$. Both satisfy $\frac{d}{dt}C(t)=C(t)\circ\operatorname{ad}_X$: for $A$,
$\frac{d}{dt}e^{tX}Ye^{-tX}=e^{tX}(XY-YX)e^{-tX}$ (because $X$ commutes with $e^{tX}$), and for $B$ this
is property (2) of Proposition~\ref{LIE:prop:exp}, applied to the linear map $\operatorname{ad}_X$ of the
$n^2$-dimensional space $M_n$. By uniqueness of solutions of linear ODEs, $A=B$; now put $t=1$.
\end{proof}

For diagonal $X=h$ this is the formula we started from: $\operatorname{ad}_h$ has eigenvalue $h_i-h_j$
on $E_{ij}$, so $e^{\operatorname{ad}_h}$ has eigenvalue $e^{h_i-h_j}=t_i/t_j$, where $t=e^h$.

\begin{figure}[ht]
\centering
\begin{tikzpicture}[scale=1.4]
  \draw[thick] (0,0) circle (1);
  \draw[addcol,thick] (1,-1.35)--(1,1.35) node[above,font=\small]{$T_I\SO(2)=\mathfrak{so}_2\cong\R$};
  \fill (1,0) circle (0.03) node[right,font=\small]{$I$};
  \foreach \t in {0.6,1.2} {\fill[addcol] (1,\t) circle (0.03); \fill ({cos(\t r)},{sin(\t r)}) circle (0.03);
     \draw[-{Stealth},gray,dashed] (0.97,\t) to[bend right=20] ({0.97*cos(\t r)},{0.97*sin(\t r)});}
  \node[font=\small] at (0,0) {$\SO(2)\cong\Sph^1$};
  \node[font=\scriptsize,gray] at (2.45,0.45) {$\theta\mapsto\begin{pmatrix}\cos\theta&-\sin\theta\\ \sin\theta&\cos\theta\end{pmatrix}$};
\end{tikzpicture}
\caption{The exponential map wraps the tangent line $\mathfrak{so}_2$ around the circle $\SO(2)$.
Its kernel is $2\pi\Z$, a lattice.}
\label{LIE:fig:so2}
\end{figure}

\subsection*{Matrix groups and their Lie algebras}

\begin{definition}
A \emph{matrix group} is a closed subgroup $G\subseteq\GL_n(\R)$ or $\GL_n(\C)$. Its \emph{Lie algebra} is
\[
  \operatorname{Lie}(G)=\{X\in M_n: e^{tX}\in G\text{ for all }t\in\R\}.
\]
\end{definition}

\begin{theorem}[von Neumann 1929; Cartan 1930]\label{LIE:thm:closed}
Let $G$ be a matrix group. Then $\operatorname{Lie}(G)$ is a real vector space closed under the commutator,
so it is a Lie algebra. Moreover $G$ is a smooth submanifold of $M_n$ (a \emph{Lie group}) with tangent
space $T_IG=\operatorname{Lie}(G)$, and $\exp$ maps a neighbourhood of $0$ in $\operatorname{Lie}(G)$
diffeomorphically onto a neighbourhood of $I$ in $G$.
\end{theorem}

We will not prove that $G$ is a manifold (see Hall, Ch.~3, or Stillwell, \emph{Naive Lie Theory}, Ch.~7), but it is worth seeing where the bracket comes from. Closure under addition follows from the
\emph{Lie product formula} $e^{X+Y}=\lim_{k\to\infty}\big(e^{X/k}e^{Y/k}\big)^k$. For the bracket, take
$X,Y\in\operatorname{Lie}(G)$. For each $t$ the matrix $e^{tX}Ye^{-tX}$ lies in $\operatorname{Lie}(G)$, because
$\exp\!\big(s\,e^{tX}Ye^{-tX}\big)=e^{tX}e^{sY}e^{-tX}\in G$. Since $\operatorname{Lie}(G)$ is a closed
subspace, it also contains the derivative
\[
  \frac{d}{dt}\Big|_{t=0}e^{tX}Ye^{-tX}=XY-YX=[X,Y].
\]
Equivalently, expanding the group commutator gives
\[
  e^{sX}e^{tY}e^{-sX}e^{-tY}=I+st\,[X,Y]+(\text{terms of total degree}\ge3\text{ in }s,t).
\]
\emph{The bracket is the infinitesimal shadow of the failure of the group to commute.}

\begin{example}[The classical groups recover the classical algebras]\leavevmode
\begin{itemize}[nosep]
\item $\operatorname{Lie}(\GL_n)=M_n=\mathfrak{gl}_n$.
\item $\operatorname{Lie}(\SL_n)=\mathfrak{sl}_n$. If $\operatorname{tr}X=0$ then $\det e^{tX}=1$ by
  Proposition~\ref{LIE:prop:exp}(4); conversely, $\det e^{tX}=e^{t\operatorname{tr}X}=1$ for all $t$ forces
  $\operatorname{tr}X=0$.
\item $\operatorname{Lie}(G_J)=\mathfrak g_J$ for the isometry group $G_J=\{g: g^TJg=J\}$ of the form $J$. If
  $X^TJ+JX=0$, then $X^T=-JXJ^{-1}$, so $(e^{X})^TJe^{X}=Je^{-X}J^{-1}Je^{X}=J$; applying this to $tX$ shows $\mathfrak g_J\subseteq
  \operatorname{Lie}(G_J)$. Conversely, differentiating $(e^{tX})^TJe^{tX}=J$ at $t=0$ gives $X^TJ+JX=0$. Thus $\mathfrak{so}_N$
  belongs to $\Orth(N)$ and $\SO(N)$, and $\mathfrak{sp}_{2n}$ to $\mathrm{Sp}_{2n}$. This is why $\mathfrak g_J$
  deserves the name ``infinitesimal isometries''.
\item The unitary group $\mathrm U(n)=\{g\in\GL_n(\C): g^*g=I\}$ has Lie algebra
  $\mathfrak u(n)=\{X: X^*=-X\}$ (skew-Hermitian matrices), and $\mathrm{SU}(n)$ has
  $\mathfrak{su}(n)=\mathfrak u(n)\cap\mathfrak{sl}_n$. Note that $\mathfrak{su}(n)\otimes\C=\mathfrak{sl}_n(\C)$: the
  compact group $\mathrm{SU}(n)$ is a \emph{real form} of $\mathfrak{sl}_n(\C)$.
\end{itemize}
\end{example}

\begin{example}[$\SO(2)$ and $\SO(3)$]\label{LIE:ex:so3}
For $X=\theta\left(\begin{smallmatrix}0&-1\\1&0\end{smallmatrix}\right)$ we have $X^2=-\theta^2I$, and
splitting the series into even and odd terms gives
\[
  e^X=\cos\theta\,I+\sin\theta\begin{pmatrix}0&-1\\1&0\end{pmatrix},
\]
which is rotation by $\theta$ (Figure~\ref{LIE:fig:so2}). In $\SO(3)$, write $X_\omega v=\omega\times v$
with $\theta=\abs{\omega}$. Using $X_\omega^3=-\theta^2X_\omega$, the same splitting gives
\emph{Rodrigues' formula}
\[
  e^{X_\omega}=I+\frac{\sin\theta}{\theta}X_\omega+\frac{1-\cos\theta}{\theta^2}X_\omega^2,
\]
which is rotation by angle $\theta$ about the axis $\omega$. Every rotation has an axis, so $\exp$ is
onto $\SO(3)$. It maps the closed ball $\abs{\omega}\le\pi$ onto $\SO(3)$, identifying only antipodal
points of the boundary sphere, so $\SO(3)\cong\R P^3$.
\end{example}

\begin{remark}[Same algebra, different groups]
The group $\mathrm{SU}(2)$ of unit quaternions is a $3$-sphere. It acts on the purely imaginary
quaternions $\cong\R^3$ by conjugation, and this gives a two-to-one homomorphism
$\mathrm{SU}(2)\to\SO(3)$ with kernel $\{\pm I\}$. The two groups have isomorphic Lie algebras,
$\mathfrak{su}(2)\cong\mathfrak{so}(3)\cong(\R^3,\times)$, but they are not isomorphic: one is simply connected,
the other has $\pi_1=\Z/2$. The Lie algebra determines a connected group only up to coverings. The
missing global information turns out to be a \emph{lattice}, as the next section shows.
\end{remark}

%======================================================================
\section{Tori, Weyl groups, and the lattice inside the group}
\label{LIE:sec:torus}

Return to $A_{n-1}$, now in the compact group $\mathrm{SU}(n)$. The diagonal matrices in $\mathrm{SU}(n)$ form
the \emph{maximal torus}
\[
  T=\big\{\operatorname{diag}(e^{i\theta_1},\dots,e^{i\theta_n}):\ \theta_1+\dots+\theta_n\equiv0\pmod{2\pi}\big\},
\]
with Lie algebra
\[
  \mathfrak t=\{\operatorname{diag}(i\theta_1,\dots,i\theta_n):\ \theta_1+\dots+\theta_n=0\}
  \cong\{\theta\in\R^n:\textstyle\sum\theta_j=0\},
\]
which is again the hyperplane of Lecture~6. Since $\mathfrak t$ is commutative, $\exp\colon\mathfrak t\to T$ is a
surjective homomorphism, and its kernel consists of the $\theta$ with every $\theta_j\in2\pi\Z$ and
$\sum\theta_j=0$:
\[
  \ker\big(\exp|_{\mathfrak t}\big)=2\pi\cdot\{x\in\Z^n:\textstyle\sum x_j=0\}=2\pi\cdot A_{n-1},
\]
the $A_{n-1}$ root lattice from Lecture~6. So
\[
  T\;\cong\;\mathfrak t\,/\,2\pi A_{n-1}:
\]
\emph{the maximal torus of $\mathrm{SU}(n)$ is Euclidean space modulo the $A_{n-1}$ root lattice.} For
$\mathrm{SU}(3)$ it is the plane modulo the hexagonal lattice, the lattice of the densest disk packing
from Lecture~1.

\paragraph{The Weyl group.} Permutation matrices (with one sign changed if necessary, to make the
determinant $1$) normalize $T$ and permute the $\theta_j$. The resulting group
\[
  W=N(T)/T\cong\Sigma_n
\]
is the Weyl group of $A_{n-1}$, generated by the reflections $s_{e_i-e_j}$, which swap $\theta_i$ and
$\theta_j$. By the spectral theorem every element of $\mathrm{SU}(n)$ is conjugate to an element of $T$,
uniquely up to the action of $W$, so
\[
  \{\text{conjugacy classes of }\mathrm{SU}(n)\}\;\longleftrightarrow\;T/W
  \;\longleftrightarrow\;\Big\{\begin{array}{@{}c@{}}\text{unordered $n$-tuples of points}\\
  \text{on }\Sph^1\text{ with product }1\end{array}\Big\}.
\]

\begin{remark}[Back to configuration spaces]
An element of $\mathrm{SU}(n)$ is \emph{regular} when its eigenvalues are distinct, that is, when its
$n$ eigenvalues form a configuration of distinct points on the circle. Choose real eigenvalue angles with
$\sum\theta_j=0$, ordered so that $\theta_1\ge\theta_2\ge\dots\ge\theta_n\ge\theta_1-2\pi$ (this is
always possible, in exactly one way), and record the gaps
\[
  g_k=\theta_k-\theta_{k+1}\ (1\le k<n),\qquad g_n=2\pi-(\theta_1-\theta_n).
\]
The gaps are non-negative and sum to $2\pi$, and together with $\sum\theta_j=0$ they determine the
$\theta_j$. So $T/W$ is a closed simplex of dimension $n-1$, the \emph{Weyl alcove}. Its interior, the
regular classes, is identified through the gaps with the reduced configuration space of $n$ points on a
circle in a fixed cyclic order (one of the open simplices of gaps in Lecture~7). For $n=3$ it is the same open triangle of gap coordinates $(x,y,1-x-y)$, up to the scale
factor $2\pi$. Its boundary walls, where two eigenvalues collide, are the walls $\theta_i=\theta_j$ of
$\Conf(n,\R)$ from Lecture~8, wrapped around the circle.
\end{remark}

All of this generalizes.

\begin{theorem}[Maximal tori; Weyl, Cartan]\label{LIE:thm:torus}
Let $G$ be a compact connected Lie group. Then $G$ has a maximal torus $T$, unique up to conjugacy,
and every element of $G$ is conjugate into $T$. The Weyl group $W=N(T)/T$ is finite and is the Weyl
group of the root system of $\mathfrak g\otimes\C$. The kernel of $\exp\colon\mathfrak t\to T$ is a lattice
containing $2\pi$ times the coroot lattice (the coroots $h_\alpha$ placed in $\mathfrak t$ as $ih_\alpha$,
as in the $\mathrm{SU}(n)$ example), with equality exactly when $G$ is simply connected; in general the
quotient is $\pi_1(G)$. The
compact simply connected simple Lie groups correspond to the connected Dynkin diagrams:
\[
  \mathrm{SU}(n+1)\ (A_n),\quad \mathrm{Spin}(2n+1)\ (B_n),\quad \mathrm{Sp}(n)\ (C_n),\quad
  \mathrm{Spin}(2n)\ (D_n),\quad G_2,\ F_4,\ E_6,\ E_7,\ E_8 .
\]
\end{theorem}

For example, $\mathrm{SU}(2)$ has $T=\{\operatorname{diag}(e^{i\theta},e^{-i\theta})\}$ and kernel
$\theta\in2\pi\Z$, while $\SO(3)=\mathrm{SU}(2)/\{\pm I\}$ has the finer kernel $\theta\in\pi\Z$. The covering
$\mathrm{SU}(2)\to\SO(3)$ of the previous section is recorded by an index-two inclusion of lattices.

%======================================================================
\section{What else a Dynkin diagram knows}
\label{LIE:sec:readoff}

A Dynkin diagram is a tiny picture: a few dots, some single, double or triple edges, and a few arrows.
It is remarkable how much can be read off from it. We collect the main items, with the
$\mathfrak{sl}_n$ case as the running check.

\paragraph{1. The Cartan matrix.} Number the nodes $1,\dots,r$. The Cartan matrix
$a_{ij}=n_{\alpha_j\alpha_i}=2\ip{\alpha_j}{\alpha_i}/\ip{\alpha_i}{\alpha_i}$ has $2$ on the diagonal,
$0$ for non-adjacent nodes, and $-1,-1$ for a single edge. For a double (triple) edge the entry in the
row of the \emph{short} root is $-2$ ($-3$) and the other is $-1$. For example, with $\alpha_1$ long and
$\alpha_2$ short (for $B_2$: $\alpha_1=e_1-e_2$, $\alpha_2=e_2$, so $a_{21}=2\ip{\alpha_1}{\alpha_2}/\ip{\alpha_2}{\alpha_2}=-2$),
\[
  A_2:\ \begin{pmatrix}2&-1\\-1&2\end{pmatrix},\qquad
  B_2:\ \begin{pmatrix}2&-1\\-2&2\end{pmatrix},\qquad
  G_2:\ \begin{pmatrix}2&-1\\-3&2\end{pmatrix}.
\]
This is the convention that makes the Serre relations read $[h_i,e_j]=a_{ij}e_j$ (it is Kac's);
Humphreys and Bourbaki use the transpose, so their $B_2$ matrix has the $-2$ in the first row. From
the Cartan matrix the whole root system, and by Serre's theorem the whole Lie algebra, can be
rebuilt (Section~\ref{LIE:sec:classification}).

\paragraph{2. The Weyl group, by generators and relations.} The Coxeter graph gives a presentation
\[
  W=\big\langle s_1,\dots,s_r\ \big|\ s_i^2=1,\ (s_is_j)^{m_{ij}}=1\big\rangle,\qquad
  m_{ij}=2,3,4,6\ \text{for}\ 0,1,2,3\ \text{edges}.
\]
For $A_{n-1}$ this is the classical presentation of $\Sigma_n$ by adjacent transpositions:
$s_is_j=s_js_i$ for $\abs{i-j}\ge2$ and $s_is_{i+1}s_i=s_{i+1}s_is_{i+1}$. Dropping the relations
$s_i^2=1$ gives Artin's \emph{braid group} (Interlude~D, Section~\ref{TOP:sec:braids}).

\paragraph{3. Counting: roots, Coxeter number, exponents, $\abs{W}$.} Let $c=s_1s_2\cdots s_r$ be the
product of all the simple reflections (a \emph{Coxeter element}). Its order $h$ is the \emph{Coxeter
number}, and its eigenvalues are $e^{2\pi im_j/h}$ for integers $1=m_1\le\dots\le m_r=h-1$, the
\emph{exponents}. Then
\[
  \abs{\Delta}=rh,\qquad \#\{\text{positive roots}\}=m_1+\dots+m_r,\qquad
  \abs{W}=\prod_{j=1}^r(m_j+1).
\]
The numbers $d_j=m_j+1$ are the degrees of the basic $W$-invariant polynomials (Chevalley). For
$A_{n-1}$ the invariants of $\Sigma_n$ are the elementary symmetric polynomials of degrees $2,\dots,n$
(degree $1$ vanishes on the hyperplane), so $\abs{W}=2\cdot3\cdots n=n!$, and $h=n$. Moreover $c$
permutes the roots in $r$ orbits of size $h$. Projecting onto the plane where $c$ acts as a rotation
by $2\pi/h$ (the \emph{Coxeter plane}), the roots land on $r$ concentric regular $h$-gons. For $E_8$
these are the famous eight rings of $30$ points.

\paragraph{4. The determinant: the density of the root lattice.} For the simply laced diagrams
($A$, $D$, $E$), normalize all roots to have $\ip{\alpha}{\alpha}=2$. Then the Gram matrix of the simple
roots \emph{is} the Cartan matrix, so the root lattice has covolume $\sqrt{\det C}$. Its shortest
vectors are the roots, of length $\sqrt2$, so the lattice packing has balls of radius $1/\sqrt2$ and
\[
  \dens(\text{root lattice})=\frac{\vol\Ball^r\big(1/\sqrt2\big)}{\sqrt{\det C}}.
\]
The packing density of Lectures~1, 4 and 6 \emph{can be read off the diagram}: $\det C=3,4,4,3,2,1$ for
$A_2,A_3,D_4,E_6,E_7,E_8$ gives
\[
  \frac{\pi}{\sqrt{12}},\quad \frac{\pi}{\sqrt{18}},\quad \frac{\pi^2}{16},\quad
  \frac{\pi^3}{48\sqrt3},\quad \frac{\pi^3}{105},\quad \frac{\pi^4}{384}.
\]
The same determinant has a Lie-theoretic meaning. It is the index of the root lattice in the weight
lattice, which equals the order of the centre of the simply connected compact group, or equivalently
$\abs{\pi_1}$ of the adjoint group. For $\mathrm{SU}(n)$ ($A_{n-1}$, $\det C=n$) the centre is the $n$-th
roots of unity times $I$. For $E_8$ ($\det C=1$) the root lattice is unimodular and the group is
simultaneously simply connected and adjoint.

\paragraph{5. Symmetries of the diagram: outer automorphisms.} Automorphisms of the diagram give the
automorphisms of $\mathfrak g$ that are not conjugations. $A_n$ ($n\ge2$), $D_n$ ($n\ge5$) and $E_6$
have a flip; $D_4$ has the symmetry group $\Sigma_3$ permuting its three outer nodes, whose elements of
order three are called \emph{triality}. For $\mathfrak{sl}_n$ the flip is, up to an inner automorphism,
$X\mapsto-X^T$.

\paragraph{6. Folding.} Fold a diagram along a symmetry, identifying the nodes in each orbit, and you
get a non-simply-laced diagram. The corresponding Lie algebra is the subalgebra fixed by the symmetry:
\[
  A_{2n-1}\rightsquigarrow C_n,\qquad D_{n+1}\rightsquigarrow B_n,\qquad E_6\rightsquigarrow F_4,\qquad
  D_4\rightsquigarrow G_2\ \text{(by triality)}.
\]
Here the symmetry $\sigma$ has to be lifted to $\mathfrak g$ by permuting the Serre generators,
$e_i\mapsto e_{\sigma(i)}$, $f_i\mapsto f_{\sigma(i)}$. For $\mathfrak{sl}_{2n}$ a lift (up to signs, which
can be absorbed into the generators) is $X\mapsto-JX^TJ^{-1}$ with the skew $J$ of
Section~\ref{LIE:sec:classical}, and its fixed points are exactly $\mathfrak{sp}_{2n}$; the plain flip
$X\mapsto-X^T$ differs from it by an inner automorphism and fixes $\mathfrak{so}_{2n}$ instead. Which
node becomes short is visible on the diagram: averaging the roots in an orbit of two orthogonal simple
roots halves the squared length. So in $A_{2n-1}$ the one fixed (middle) node stays long, giving $C_n$,
while in $D_{n+1}$ the two swapped end nodes become one short node, giving $B_n$.
So the multiple edges remember a hidden symmetry of a simply laced diagram. Folding with the golden
ratio instead leads to the non-crystallographic $E_8\rightsquigarrow H_4$, $D_6\rightsquigarrow H_3$,
$A_4\rightsquigarrow H_2$ and to quasicrystals (Interlude~C).

\paragraph{7. The extended diagram and McKay's correspondence.} Adding the negative of the highest
root as an extra node gives the \emph{extended} (affine) Dynkin diagram, the diagram of the affine
Weyl group that tiles $\mathfrak t$ by alcoves (Section~\ref{LIE:sec:torus}). For the simply laced types
something unexpected happens. The finite subgroups $\Gamma\subset\mathrm{SU}(2)$ (cyclic, binary dihedral,
binary tetrahedral, octahedral, icosahedral) correspond exactly to the extended diagrams
$\tilde A_{n-1},\tilde D_n,\tilde E_6,\tilde E_7,\tilde E_8$ (McKay, 1980). The nodes are the
irreducible representations of $\Gamma$, and the edges record tensoring with the two-dimensional
representation. So the binary icosahedral group, the $120$ unit quaternions that are the vertices of
the $600$-cell and the double cover of the rotation group of the icosahedron of Interlude~C, \emph{is}
$\tilde E_8$. The same ADE list, now with the ordinary (unextended) diagrams, classifies the simple
surface singularities $\C^2/\Gamma$, whose minimal resolutions have the Dynkin diagram as dual graph
(du~Val, 1934; Arnold, 1972), and the quivers of finite representation type (Gabriel, 1972).

\paragraph{8. Representations: label the nodes.} The irreducible finite-dimensional representations
of $\mathfrak g$ correspond to labellings of the nodes by non-negative integers $(n_1,\dots,n_r)$, the
coordinates of the highest weight in the basis of fundamental weights. Their dimensions are given by
Weyl's dimension formula. For $\mathfrak{sl}_2$ the label $m$ gives dimension $m+1$, as in
Theorem~\ref{LIE:thm:sl2}. For $\mathfrak{sl}_3$ the labels $(a,b)$ give
\[
  \dim V_{(a,b)}=\tfrac12(a+1)(b+1)(a+b+2),
\]
so $(1,0)$ is $\C^3$ (dimension $3$) and $(1,1)$ is the adjoint representation (dimension $8$).

\begin{center}
\renewcommand{\arraystretch}{1.2}
\small
\begin{tabular}{lcccccc}
\toprule
type & $\abs{\Delta}$ & $h$ & exponents & $\abs{W}$ & $\det C$ & $\dim\mathfrak g$\\
\midrule
$A_n$ & $n(n+1)$ & $n+1$ & $1,2,\dots,n$ & $(n+1)!$ & $n+1$ & $n(n+2)$\\
$B_n$, $C_n$ & $2n^2$ & $2n$ & $1,3,\dots,2n-1$ & $2^nn!$ & $2$ & $n(2n+1)$\\
$D_n$ & $2n(n-1)$ & $2n-2$ & $1,3,\dots,2n-3;\ n-1$ & $2^{n-1}n!$ & $4$ & $n(2n-1)$\\
$E_6$ & $72$ & $12$ & $1,4,5,7,8,11$ & $51\,840$ & $3$ & $78$\\
$E_7$ & $126$ & $18$ & $1,5,7,9,11,13,17$ & $2\,903\,040$ & $2$ & $133$\\
$E_8$ & $240$ & $30$ & $1,7,11,13,17,19,23,29$ & $696\,729\,600$ & $1$ & $248$\\
$F_4$ & $48$ & $12$ & $1,5,7,11$ & $1\,152$ & $1$ & $52$\\
$G_2$ & $12$ & $6$ & $1,5$ & $12$ & $1$ & $14$\\
\bottomrule
\end{tabular}
\end{center}
(References: Humphreys, \emph{Reflection Groups and Coxeter Groups}, \S\S3.7--3.9 and 3.16--3.20; Bourbaki,
\emph{Lie Groups and Lie Algebras}, Ch.~VI, Plates; McKay, \emph{Graphs, singularities, and finite
groups}, Proc.\ Sympos.\ Pure Math.\ 37 (1980).)

%======================================================================
\section{Summary: a dictionary}
\label{LIE:sec:dictionary}

\begin{center}
\renewcommand{\arraystretch}{1.25}
\begin{tabular}{p{0.42\textwidth}p{0.50\textwidth}}
\toprule
Packings and lattices (Lectures 4--8) & Lie theory\\
\midrule
root system $\Delta$ & nonzero eigenvalues of $\operatorname{ad}\mathfrak h$ on $\mathfrak g$\\
Cartan integer $n_{\beta\alpha}\in\Z$ & eigenvalue of the coroot $h_\alpha$; integral by $\mathfrak{sl}_2$ theory\\
reflection $s_\alpha$, Weyl group $W$ & symmetry of $\alpha$-strings; $W=N(T)/T$\\
choice of $d$, positive roots, Weyl chamber & upper triangular matrices (a Borel subalgebra)\\
simple roots, Dynkin diagram & generators $e_i,f_i,h_i$ and Serre relations\\
rank $+\abs{\Delta}$ & $\dim\mathfrak g$\\
$\det$ of the Cartan matrix; density of the root lattice & order of the centre of the simply connected group\\
symmetries and foldings of the diagram & outer automorphisms; fixed subalgebras\\
root lattice $\Z\Delta$ (coroot lattice, in general) & $\tfrac1{2\pi}\,$kernel of $\exp$ on the maximal torus (simply connected, compact)\\
$E_8$ kissing configuration ($240$ vectors) & root spaces of $\mathfrak e_8$, $248=8+240$\\
walls $x_i=x_j$ of $\Conf(N,\R)$ & eigenvalue collisions; non-regular elements of $\mathrm{SU}(N)$\\
one-parameter groups of diffeomorphisms (Lecture 9) & $t\mapsto e^{tX}$\\
\bottomrule
\end{tabular}
\end{center}

%======================================================================
\section*{Exercises}
\addcontentsline{toc}{section}{Exercises}

\begin{exercise}
Write out the root decomposition of $\mathfrak{sl}_3$ explicitly. Check that $[E_{12},E_{23}]=E_{13}$,
$[E_{13},E_{32}]=E_{12}$ and $[E_{12},E_{13}]=0$, and interpret each identity on the hexagon of
Example~\ref{LIE:ex:sl3}.
\end{exercise}

\begin{exercise}
Prove the Jacobi identity for the commutator, and show that $\omega\mapsto X_\omega$ is an isomorphism
$(\R^3,\times)\to\mathfrak{so}_3(\R)$.
\end{exercise}

\begin{exercise}
Expand $e^{sX}e^{tY}e^{-sX}e^{-tY}$ to second order and confirm that the result is $I+st[X,Y]$ plus
higher-order terms.
\end{exercise}

\begin{exercise}
Show that $\exp\colon\mathfrak{sl}_2(\R)\to\SL_2(\R)$ is not surjective: the matrix
$\left(\begin{smallmatrix}-1&1\\0&-1\end{smallmatrix}\right)$ is not in the image. \emph{Hint:} a real
traceless $2\times2$ matrix has eigenvalues $\pm\lambda$ or $\pm i\mu$. Compute the possible traces of
$e^X$, and decide when the trace can equal $-2$.
\end{exercise}

\begin{exercise}
Using the antidiagonal form $J$, find the roots of $\mathfrak{so}_6$ and check that they form a copy of
$A_3$ (twelve roots, simply laced, rank $3$). Do the same for $\mathfrak{so}_4$ and $A_1\times A_1$, and for
$\mathfrak{sp}_4$ and $\mathfrak{so}_5$ (show that $C_2$ and $B_2$ differ by a rotation of $45^\circ$ and a scaling).
\end{exercise}

\begin{exercise}
Prove Rodrigues' formula, and show that $\exp$ maps the closed ball of radius $\pi$ in $\mathfrak{so}_3$
onto $\SO(3)$, identifying exactly the antipodal points of the boundary.
\end{exercise}

\begin{exercise}
Draw the kernel of $\exp$ on the maximal torus of $\mathrm{SU}(3)$ inside the plane
$\theta_1+\theta_2+\theta_3=0$, together with the Weyl alcove. Check that the kernel is a hexagonal
lattice and that the alcove is an equilateral triangle. How many alcoves fit into a fundamental domain
of the lattice?
\end{exercise}

\begin{exercise}
Using only the Dynkin diagrams, compute $\det C$ for $D_n$ and $E_7$ and recover the densities of the
lattices $D_5$ and $E_7$. Compare with the laminated lattices of Lecture~6
(Section~\ref{L04b:sec:laminated}).
\end{exercise}

\begin{exercise}
Check $\abs{W(E_8)}=696\,729\,600$ from the exponents, and verify that the exponents of $A_{n-1}$ give
$\abs{\Sigma_n}=n!$. Find a Coxeter element of $\Sigma_4$ and show that it has order $h=4$.
\end{exercise}

\section*{Further reading}
\addcontentsline{toc}{section}{Further reading}
\begin{itemize}[nosep]
\item J.~Stillwell, \emph{Naive Lie Theory}, Springer UTM, 2008. Matrix groups first, with minimal prerequisites.
\item B.~C.~Hall, \emph{Lie Groups, Lie Algebras, and Representations}, 2nd ed., Springer GTM 222, 2015.
  Matrix groups, the exponential map, and a full treatment of the classification.
\item J.~E.~Humphreys, \emph{Introduction to Lie Algebras and Representation Theory}, Springer GTM 9, 1972.
  The algebraic route from $\mathfrak{sl}_2$ to root systems and Serre's theorem.
\item W.~Fulton and J.~Harris, \emph{Representation Theory: A First Course}, Springer GTM 129, 1991.
  Lectures 11--13 work through $\mathfrak{sl}_2$ and $\mathfrak{sl}_3$ by hand, and Lecture 21 does the classification.
\item J.~H.~Conway and N.~J.~A.~Sloane, \emph{Sphere Packings, Lattices and Groups}, Ch.~4 and 8, for the
  lattice side of $E_8$.
\end{itemize}
% ===== end chapters/Lie.tex =====
\endgroup
\begingroup\renewcommand{\thechapter}{B}\renewcommand{\theHchapter}{interludeB}
% ===== begin chapters/NT.tex =====
% MAT670 -- Interlude B (end of Part II): lattices, quadratic forms and the geometry of numbers.
% Entirely new material (not in the handwritten notes). Labels prefixed NT:.

\chapter[Interlude B: Number theory]{Interlude B: Lattices, quadratic forms and the geometry of numbers}
\label{NT:ch}

\begin{addedblock}
This interlude is new: it is not in the handwritten notes. It starts from the title of Blichfeldt's
1929 paper quoted in Lecture~2, \emph{``The minimum value of quadratic forms, and the closest packing of
spheres''}, and follows the first half of that title into number theory. To keep it readable it is set
in black; all of it should be read as added material.
\end{addedblock}

Blichfeldt's title puts two problems side by side, and in Lectures~1--6 we only looked at the second
one. This interlude is about the first. We shall see that
\begin{itemize}
\item a lattice \emph{is} a positive definite quadratic form, so the lattice packing problem is the
  question of how large the minimum of a form of given determinant can be (Hermite's problem);
\item counting lattice points in balls is the classical problem of counting representations of integers
  by forms, and the leading term of the count is the covolume, which is what density measures;
\item Blichfeldt's averaging trick from Lecture~2 has a famous cousin, Blichfeldt's lemma of 1914, which
  gives Minkowski's convex body theorem, and through it Fermat's and Lagrange's theorems on sums of squares;
\item the generating function of the lattice point counts, the \emph{theta series}, is a modular form.
  Poisson summation turns this into strong restrictions on very good lattices, and the same Poisson
  summation is the engine of the Cohn--Elkies bound behind Viazovska's theorem.
\end{itemize}
Throughout, $\Lat\subset\R^n$ is a full-rank lattice, $\covol(\Lat)=\abs{\det\Lat}$ as in Lecture~4,
$\lambda_1(\Lat)$ is the length of a shortest nonzero vector, $V_n=\vol(\Ball^n)$, and
\[
  \dens(\Lat)=\frac{V_n\,(\lambda_1/2)^n}{\covol(\Lat)}
\]
is the density of the associated packing (Lecture~4). To avoid a clash with the centre density
$\cent$, the theta series will be written $\Theta_\Lat$.

%======================================================================
\section{Lattices are quadratic forms}
\label{NT:sec:forms}

Let $B$ be a basis matrix of $\Lat$, with the basis vectors $b_1,\dots,b_n$ as its columns. A lattice
vector is $Bx$ with $x\in\Z^n$, and its squared length is
\[
  \abs{Bx}^2=x^TB^TBx=x^TGx=:Q(x),\qquad G=B^TB=\big(\ip{b_i}{b_j}\big)_{i,j}.
\]
The \emph{Gram matrix} $G$ is symmetric and positive definite, and $\det G=\covol(\Lat)^2$. Conversely,
every positive definite symmetric $G$ factors as $G=B^TB$ (Cholesky), and $B$ is unique up to
$B\mapsto UB$ with $U\in\Orth(n)$. Changing the basis replaces $B$ by $BK$ with $K\in\GL_n(\Z)$
(Lecture~4), hence $G$ by $K^TGK$. So:
\[
  \frac{\{\text{lattices in }\R^n\}}{\text{rotations}}
  \;\;\longleftrightarrow\;\;
  \frac{\{\text{positive definite }n\times n\text{ symmetric matrices}\}}{G\sim K^TGK,\ K\in\GL_n(\Z)} .
\]
Under this dictionary
\[
  \lambda_1(\Lat)^2=\min_{x\in\Z^n\setminus0}Q(x)=:\min Q,\qquad \covol(\Lat)=\sqrt{\det G},
\]
and the number of lattice vectors of squared length $m$ is the number of \emph{representations} of $m$
by $Q$. Integral lattices (Lecture~4) are the forms with $G\in M_n(\Z)$. For example $\Z^n$ is the sum
of squares $x_1^2+\dots+x_n^2$, and the hexagonal lattice with basis $(1,0)$,
$(\tfrac12,\tfrac{\sqrt3}2)$ is
\[
  Q_{\mathrm{hex}}(x,y)=x^2+xy+y^2,\qquad G=\begin{pmatrix}1&\tfrac12\\ \tfrac12&1\end{pmatrix},
\]
while the root lattice $A_2$ (Lecture~6, minimal norm $2$) is $2Q_{\mathrm{hex}}$.

\begin{definition}[Hermite's constant]\label{NT:def:hermite}
For a positive definite form put $\gamma(Q)=\min Q/(\det G)^{1/n}$; this is invariant under scaling and
$\GL_n(\Z)$. \emph{Hermite's constant} is $\gamma_n=\sup_Q\gamma(Q)$, the supremum over all positive
definite $n$-ary forms. Equivalently, $\gamma_n=\sup_\Lat\lambda_1(\Lat)^2/\covol(\Lat)^{2/n}$.
\end{definition}

Comparing with the formula for $\dens(\Lat)$ gives
\begin{equation}\label{NT:eq:hermite-density}
  \dens(\Lat)=\frac{V_n}{2^n}\,\gamma(\Lat)^{n/2},
  \qquad\text{so}\qquad
  \sup_\Lat\dens(\Lat)=\frac{V_n}{2^n}\,\gamma_n^{n/2}.
\end{equation}
\emph{Finding the densest lattice packing is the same problem as computing Hermite's constant.} The
known values are below. They are attained exactly by the lattices of Lectures~4--6, and the reader
can recompute each from minimal norm $2$ and the determinants $\det G=n+1$ for $A_n$, $4$ for $D_n$, and
$3,2,1$ for $E_6,E_7,E_8$. For the Leech lattice the minimal norm is $4$ and $\det G=1$.
\begin{center}
\renewcommand{\arraystretch}{1.2}
\begin{tabular}{lccccccccc}
\toprule
$n$ & 1 & 2 & 3 & 4 & 5 & 6 & 7 & 8 & 24\\
\midrule
$\gamma_n^n$ & $1$ & $\tfrac43$ & $2$ & $4$ & $8$ & $\tfrac{64}3$ & $64$ & $2^8$ & $4^{24}$\\
optimal lattice & $\Z$ & $A_2$ & $A_3$ & $D_4$ & $D_5$ & $E_6$ & $E_7$ & $E_8$ & $\Lambda_{24}$\\
density & $1$ & $.9069$ & $.7405$ & $.6169$ & $.4653$ & $.3729$ & $.2953$ & $.2537$ & $.00193$\\
\bottomrule
\end{tabular}
\end{center}
(Lagrange for $n=2$; Gauss for $n=3$; Korkine--Zolotarev for $n=4,5$; Blichfeldt for $n=6,7,8$;
Cohn--Kumar for $n=24$. See the Gallery for references.)

\subsection*{Binary forms: Lagrange's reduction}

For $n=2$ write $Q(x,y)=ax^2+bxy+cy^2$, so $G=\left(\begin{smallmatrix}a&b/2\\b/2&c\end{smallmatrix}\right)$
and $\det G=ac-b^2/4$. Following Lagrange (1773) and Gauss (\emph{Disquisitiones}, 1801), call $Q$
\emph{reduced} if
\[
  \abs b\le a\le c .
\]

\begin{proposition}[Lagrange reduction]\label{NT:prop:reduction}
Every positive definite binary form is equivalent under $\SL_2(\Z)$ to a reduced form, and if $Q$ is
reduced then $\min Q=a$.
\end{proposition}

\begin{proof}
Let $v_1$ be a shortest nonzero vector of the lattice $\Lat$ of $Q$. It is primitive, since
$v_1=kw$ with $k\ge2$ would make $w$ shorter. So $v_1$ extends to a basis $(v_1,w)$, and every other
completion of $v_1$ to a basis is $\pm w+kv_1$, $k\in\Z$. Among the vectors $w+kv_1$ choose a shortest
one, $v_2$, and replace $v_2$ by $-v_2$ if needed to make the determinant positive. In the basis
$(v_1,v_2)$ we have $a=\abs{v_1}^2=\min Q$, $c=\abs{v_2}^2\ge a$, and $b=2\ip{v_1}{v_2}$. Minimality
of $v_2$ gives $c\le\abs{v_2\pm v_1}^2=c\pm b+a$, that is, $\abs b\le a$. For the second statement:
if $\abs b\le a\le c$ and $\abs x\ge\abs y>0$, then
$Q(x,y)\ge a x^2-a\abs{x}\abs{y}+cy^2\ge cy^2\ge a$; the case $\abs y\ge\abs x>0$ is symmetric
(use $\abs b\le c$), and $Q(x,0)=ax^2\ge a$, $Q(0,y)=cy^2\ge a$.
\end{proof}

\begin{theorem}[Lagrange; $\gamma_2=2/\sqrt3$]\label{NT:thm:gamma2}
For every positive definite binary form, $\min Q\le\frac{2}{\sqrt3}\sqrt{\det G}$, with equality only
for forms equivalent to a multiple of $x^2+xy+y^2$. Consequently every lattice packing of disks has
density at most $\pi/\sqrt{12}$, with equality only for the hexagonal lattice.
\end{theorem}

\begin{proof}
Take $Q$ reduced. Then $\det G=ac-\tfrac{b^2}4\ge a^2-\tfrac{a^2}4=\tfrac34a^2$, so
$\min Q=a\le\sqrt{4/3}\sqrt{\det G}$. Equality forces $c=a$ and $\abs b=a$, i.e.\
$Q=a(x^2\pm xy+y^2)$, and $x^2-xy+y^2$ is equivalent to $x^2+xy+y^2$ via $y\mapsto-y$. The density
statement is \eqref{NT:eq:hermite-density}: $\frac{\pi}{4}\cdot\frac{2}{\sqrt3}=\frac{\pi}{2\sqrt3}$.
\end{proof}

This is Theorem~\ref{L1:thm:lattice2} of Lecture~1 again. In fact the proof there, which stacked the
rows $\Z v$ at height $h$ and slid the next row so that one of its points sits over the gap, is
exactly the choice of $v_1$ and $v_2$ above. That proof generalizes to all $n$ as Hermite's inequality
(Exercise~\ref{NT:ex:hermite}).

\subsection*{The picture: the modular figure}

Rotate and scale so that $v_1=1\in\C$, and write $v_2=\tau$ with $\operatorname{Im}\tau>0$. Then
\[
  Q(x,y)=\abs{v_1}^2\abs{x+y\tau}^2=a\big(x^2+2\operatorname{Re}\tau\,xy+\abs\tau^2y^2\big),
\]
so $b/a=2\operatorname{Re}\tau$ and $c/a=\abs\tau^2$. The reduction conditions become
\[
  \abs{\operatorname{Re}\tau}\le\tfrac12,\qquad \abs\tau\ge1 ,
\]
which describe the classical fundamental domain $\mathcal F$ of $\SL_2(\Z)$ acting on the upper half
plane $\mathbb H$ by M\"obius transformations. A change of basis in $\SL_2(\Z)$ acts on $\tau=v_2/v_1$
exactly by $\tau\mapsto\frac{a\tau+b}{c\tau+d}$. So the space of planar lattices up to rotation and
scaling is $\SL_2(\Z)\backslash\mathbb H$, and $\mathcal F$ is a map of it (Figure~\ref{NT:fig:modular}).
Allowing $\GL_2(\Z)$ also identifies $\tau$ with $-\bar\tau$, which folds $\mathcal F$ in half.

In these coordinates the density is very simple. The covolume is $\abs{v_1}^2\operatorname{Im}\tau$ and,
for $\tau\in\mathcal F$, the minimum is $\abs{v_1}^2$, so
\[
  \dens(\tau)=\frac{\pi\abs{v_1}^2/4}{\abs{v_1}^2\operatorname{Im}\tau}=\frac{\pi}{4\operatorname{Im}\tau}
  \qquad(\tau\in\mathcal F).
\]
The level sets of the density are horizontal lines. The densest lattice is the \emph{lowest point} of
$\mathcal F$, the corner $\rho=e^{\pi i/3}$ with $\operatorname{Im}\rho=\frac{\sqrt3}{2}$: the hexagonal
lattice. The square lattice is $\tau=i$, with $\dens=\pi/4$. Lattices high up in the cusp are long and
thin and have small density; this is Mahler's criterion (Lecture~4) in its simplest form.

\begin{figure}[ht]
\centering
\begin{tikzpicture}[scale=2.3]
  \fill[addcol!12] (-0.5,3.0) -- (-0.5,0.866) arc (120:60:1) -- (0.5,3.0) -- cycle;
  \draw[->] (-1.7,0) -- (1.8,0) node[right]{$\operatorname{Re}\tau$};
  \draw[->] (0,-0.05) -- (0,3.15);
  \foreach \c in {-1,1} \draw[gray,dashed] (\c,0) ++(0:1) arc (0:180:1);
  \draw[gray,dashed] (1,0) arc (0:180:1);
  \foreach \c in {-1.5,1.5} \draw[gray,dashed] (\c,0) -- (\c,3.0);
  \draw[thick] (-0.5,3.0) -- (-0.5,0.866) arc (120:60:1) -- (0.5,3.0);
  \draw[flagcol,thin] (-1.6,0.866) -- (1.7,0.866) node[below right,font=\scriptsize]{$\dens=\frac{\pi}{\sqrt{12}}$};
  \draw[flagcol,thin] (-1.6,1) -- (1.7,1) node[above right,font=\scriptsize]{$\dens=\frac{\pi}{4}$};
  \draw[flagcol,thin] (-1.6,2) -- (1.7,2) node[right,font=\scriptsize]{$\dens=\frac{\pi}{8}$};
  \fill (0,1) circle (0.025) node[above left,font=\small]{$i$};
  \fill (0.5,0.866) circle (0.025) node[below right,font=\small]{$\rho$};
  \fill (-0.5,0.866) circle (0.025) node[below left,font=\small]{$\rho^2$};
  \foreach \x/\l in {-1/{-1},-0.5/{-\frac12},0.5/{\frac12},1/{1}} \node[below,font=\scriptsize] at (\x,0) {$\l$};
  \node[font=\small] at (0,2.6) {$\mathcal F$};
  \node[font=\scriptsize,gray,align=center] at (1.2,2.75) {long thin lattices\\(cusp)};
  % little lattice pictures
  \begin{scope}[shift={(-1.35,1.55)},scale=0.13]
    \foreach \i in {-2,...,2} \foreach \j in {-1,...,1} \fill ({\i+0.5*\j},{0.866*\j}) circle (0.18);
  \end{scope}
  \draw[gray,-{Stealth}] (-1.1,1.4) -- (0.42,0.9);
  \begin{scope}[shift={(1.25,1.55)},scale=0.13]
    \foreach \i in {-2,...,2} \foreach \j in {-1,...,1} \fill (\i,\j) circle (0.18);
  \end{scope}
  \draw[gray,-{Stealth}] (1.0,1.45) -- (0.06,1.03);
\end{tikzpicture}
\caption{The reduced domain $\mathcal F=\{\abs{\operatorname{Re}\tau}\le\frac12,\ \abs\tau\ge1\}$ for
binary forms, with some of its $\SL_2(\Z)$-translates dashed. Each point is a planar lattice with basis
$1,\tau$, up to rotation and scaling. The density $\pi/(4\operatorname{Im}\tau)$ is constant on
horizontal lines and is largest at the corners $\rho,\rho^2$, the hexagonal lattice.}
\label{NT:fig:modular}
\end{figure}

\begin{remark}[Class numbers]\label{NT:rem:classno}
For integral forms ($a,b,c\in\Z$) the discriminant $D=b^2-4ac<0$ is an invariant, and reduced forms
satisfy $3a^2\le -D$. So there are only finitely many reduced forms of a given discriminant. The number of
\emph{primitive} ones ($\gcd(a,b,c)=1$), with the usual boundary convention ($b\ge0$ if $\abs b=a$ or $a=c$),
is the \emph{class number} $h(D)$. For instance, $h(-4)=1$ (only $x^2+y^2$) and $h(-20)=2$ ($x^2+5y^2$ and $2x^2+2xy+3y^2$). Gauss
asked which $D$ have $h(D)=1$; see the Gallery.
\end{remark}

%======================================================================
\section{Counting lattice points: the Gauss circle problem}
\label{NT:sec:circle}

The density of a lattice packing is a count in disguise: by Exercise~\ref{L1:ex:periodic-density} of
Lecture~1, the number of lattice points in a large ball is
\[
  N_\Lat(r)=\#\{x\in\Lat:\abs x\le r\}\;\sim\;\frac{V_nr^n}{\covol(\Lat)},
  \qquad\text{so}\qquad
  \dens(\Lat)=\lim_{r\to\infty}\frac{N_\Lat(r)\,V_n(\lambda_1/2)^n}{V_nr^n}.
\]
How good is the approximation? For $\Lat=\Z^2$ this is \emph{Gauss's circle problem}. Write
\[
  N(r)=\#\{(x,y)\in\Z^2:x^2+y^2\le r^2\}=\pi r^2+E(r).
\]

\begin{proposition}[Gauss]\label{NT:prop:gauss}
$\abs{E(r)}\le\pi\big(\sqrt2\,r+\tfrac12\big)$ for $r\ge1$; in particular $E(r)=O(r)$.
\end{proposition}

\begin{proof}
Attach to each lattice point $p$ the unit square $p+[-\frac12,\frac12]^2$; these squares tile the plane
and each has diameter $\sqrt2$. The squares of the $N(r)$ points in the closed disk of radius $r$ lie
in the disk of radius $r+\frac{\sqrt2}2$. Conversely, every point of the disk of radius
$r-\frac{\sqrt2}2$ lies in a square whose centre is in the disk of radius $r$. Comparing areas,
\[
  \pi\big(r-\tfrac{\sqrt2}{2}\big)^2\le N(r)\le\pi\big(r+\tfrac{\sqrt2}{2}\big)^2,
\]
and $\pi(r\pm\frac{\sqrt2}2)^2-\pi r^2=\pm\pi\sqrt2\,r+\frac\pi2$.
\end{proof}

\begin{figure}[ht]
\centering
\begin{tikzpicture}[scale=0.42]
  \def\rr{4.3}
  \foreach \i in {-5,...,5} \foreach \j in {-5,...,5} {
    \pgfmathtruncatemacro{\inside}{(\i*\i+\j*\j<=18.49)?1:0}
    \ifnum\inside=1
      \fill[addcol!15] (\i-0.5,\j-0.5) rectangle (\i+0.5,\j+0.5);
      \draw[addcol!50,very thin] (\i-0.5,\j-0.5) rectangle (\i+0.5,\j+0.5);
      \fill (\i,\j) circle (2.6pt);
    \else
      \fill[gray!60] (\i,\j) circle (1.6pt);
    \fi
  }
  \draw[thick] (0,0) circle (\rr);
  \draw[flagcol,dashed] (0,0) circle ({\rr+0.7071});
  \draw[flagcol,dashed] (0,0) circle ({\rr-0.7071});
  \draw[-{Stealth}] (0,0) -- (30:\rr) node[pos=0.6,above,font=\small]{$r$};
  \node[font=\small,flagcol] at (6.9,3.6) {$r\pm\tfrac{\sqrt2}{2}$};
\end{tikzpicture}
\caption{Gauss's argument for $r=4.3$: the $N(r)=61$ unit squares centred at the lattice points in the
disk are trapped between the disks of radii $r\pm\frac{\sqrt2}{2}$, so $\abs{N(r)-\pi r^2}=O(r)$.}
\label{NT:fig:gauss}
\end{figure}

The same argument, with a fundamental parallelepiped in place of the square, gives
$N_\Lat(r)=V_nr^n/\covol(\Lat)+O(r^{n-1})$ for every lattice and every dimension, and also for any
ellipsoid $\{Q(x)\le r^2\}$ in place of the ball. So the error term is a ``boundary layer'', as in the
density remarks of Lecture~1. The question is how much the lattice points conspire along the boundary.

\paragraph{Representation numbers.} Let $r_2(m)=\#\{(x,y)\in\Z^2:x^2+y^2=m\}$. Then
$N(\sqrt X)=1+\sum_{1\le m\le X}r_2(m)$, so Proposition~\ref{NT:prop:gauss} says that \emph{the average
value of $r_2(m)$ is $\pi$}. At the same time $r_2(m)$ is given exactly by a formula of Jacobi (1829):
with $\chi(d)=0,1,-1$ for $d$ even, $d\equiv1$, $d\equiv3\pmod4$,
\begin{equation}\label{NT:eq:r2}
  r_2(m)=4\sum_{d\mid m}\chi(d)=4\big(d_1(m)-d_3(m)\big),
\end{equation}
where $d_i(m)$ counts the divisors of $m$ that are $\equiv i\pmod 4$. (For example $r_2(5)=8$ and
$r_2(3)=0$; it was checked numerically for all $m<2000$ in preparing these notes.) Combining the
geometric count with the arithmetic formula gives a lattice point proof of Leibniz's series.

\begin{proposition}[Lattice points prove Leibniz]\label{NT:prop:leibniz}
$1-\frac13+\frac15-\frac17+\dots=\frac\pi4$.
\end{proposition}

\begin{proof}
Put $L=\sum_{d\ge1}\chi(d)/d$ and $S(t)=\sum_{d\le t}\chi(d)\in\{0,1\}$. By \eqref{NT:eq:r2},
$N(\sqrt X)-1=4\sum_{dm\le X}\chi(d)$, a sum over the lattice points under the hyperbola $dm=X$. With
$y=\sqrt X$, Dirichlet's hyperbola method splits it as
\[
  \sum_{dm\le X}\chi(d)=\sum_{d\le y}\chi(d)\Big\lfloor\frac Xd\Big\rfloor+\sum_{m\le y}S\Big(\frac Xm\Big)-S(y)\lfloor y\rfloor .
\]
The last two terms are $O(y)$. In the first, $\lfloor X/d\rfloor=X/d+O(1)$, and the alternating series
gives $\sum_{d\le y}\chi(d)/d=L+O(1/y)$. Hence $N(\sqrt X)=4LX+O(\sqrt X)$. Proposition~\ref{NT:prop:gauss}
gives $N(\sqrt X)=\pi X+O(\sqrt X)$. Divide by $X$ and let $X\to\infty$.
\end{proof}

\paragraph{How big is $E(r)$?} Gauss's $O(r)$ uses nothing about circles except that they are curves
of finite length. Curvature helps: Sierpi\'nski (1906) proved $E(r)=O(r^{2/3})$, and the exponent has
been lowered slowly since. The conjecture, usually attributed to Hardy, is that
\[
  E(r)=O_\varepsilon\big(r^{1/2+\varepsilon}\big)\qquad\text{for every }\varepsilon>0,
\]
and this would be optimal: Hardy and, independently, Landau showed in 1915 that $E(r)$ is not
$o(r^{1/2})$, and Hardy (1916) that it is \emph{not} $o(r^{1/2}(\log r)^{1/4})$, so the
exponent $\frac12$ cannot be lowered. Sharper logarithmic factors are due to Hafner (1981) and
Soundararajan (2003). The exponent $\frac12$ is also the right one on average: Cram\'er (Math.\ Z.\ 15,
1922) showed
that $\frac{1}{R}\int_0^RE(r)^2\,dr$ grows like a constant times $R$, so a typical $\abs{E(r)}$ is of
size $r^{1/2}$. The current upper bounds, as exponents of $r$, are
\begin{center}
\renewcommand{\arraystretch}{1.15}
\begin{tabular}{lll}
\toprule
Gauss (c.~1800) & $1$ & elementary\\
Sierpi\'nski (1906) & $2/3\approx0.6667$ & \\
Huxley (2003) & $131/208\approx0.62981$ & Proc.\ London Math.\ Soc.\ 87\\
Bourgain--Watt (2017) & $517/824\approx0.62743$ & arXiv:1709.04340 (claimed)\\
Li--Yang (2023) & $\approx0.628966$ & arXiv:2308.14859\\
conjecture & $1/2+\varepsilon$ & \\
\bottomrule
\end{tabular}
\end{center}
Huxley's exponent is the best one published in a refereed journal: as of 2026 both later papers are
unpublished arXiv preprints. Both build on the Bombieri--Iwaniec method, and Bourgain--Watt add
Bourgain's decoupling theory. The Bourgain--Watt preprint claims the smaller exponent, but Li and Yang
describe Huxley's bound as the best known before their work and present theirs as its first improvement.
So the record is not settled, and should be quoted with care.

\begin{remark}[Why a packing person should care]
The circle problem is the cleanest model of a phenomenon that appears all through these notes: the bulk
term of a count is a volume, and hence a density, while the correction is a boundary term whose size
depends on fine arithmetic. In Lecture~2 we removed the boundary term by letting a box grow. Here it is
the main object of study.
\end{remark}

%======================================================================
\section{Blichfeldt again: Minkowski's convex body theorem}
\label{NT:sec:minkowski}

In Lecture~2 Blichfeldt's method compared a pointwise bound with an average: if
$\sum_if(\varphi_i^{-1}x)\le1$ everywhere, then its average, which is proportional to the density, is at
most $1$. Fifteen years earlier Blichfeldt had used the same comparison the other way round.

\begin{lemma}[Blichfeldt 1914]\label{NT:lem:blichfeldt}
Let $S\subset\R^n$ be measurable with $\vol(S)>\covol(\Lat)$. Then $S$ contains two distinct points
$s,s'$ with $s-s'\in\Lat$.
\end{lemma}

\begin{proof}
Consider the periodization $F(x)=\sum_{\lambda\in\Lat}\chi_S(x+\lambda)$, a $\Lat$-periodic function.
Integrating over a fundamental domain $P$,
\[
  \int_PF(x)\,dx=\sum_{\lambda\in\Lat}\vol\big(S\cap(P+\lambda)\big)=\vol(S)>\vol(P).
\]
So $F$ cannot be $\le1$ everywhere: there is an $x$ with $F(x)\ge2$, that is, two different
$\lambda,\lambda'\in\Lat$ with $s=x+\lambda$ and $s'=x+\lambda'$ in $S$. Then $s-s'=\lambda-\lambda'\in\Lat\setminus0$.
\end{proof}

\begin{figure}[ht]
\centering
\begin{tikzpicture}[scale=1.25]
  \def\blob{plot[smooth cycle,tension=0.8] coordinates {(0.25,0.35) (1.3,0.15) (2.35,0.55) (2.6,1.45) (1.75,2.05) (0.8,1.8) (0.15,1.2)}}
  \foreach \i in {0,...,3} \draw[gray!60] (\i,0) -- (\i,2.3);
  \foreach \j in {0,...,2} \draw[gray!60] (-0.1,\j) -- (3.1,\j);
  \foreach \i in {0,...,3} \foreach \j in {0,...,2} \fill[gray] (\i,\j) circle (0.025);
  \fill[addcol!25] \blob; \draw[addcol,thick] \blob;
  \node[addcol,font=\small] at (1.4,1.05) {$S$};
  \draw[-{Stealth},thick] (3.35,1) -- (4.15,1) node[midway,above,font=\scriptsize]{mod $\Lat$};
  \begin{scope}[xshift=4.5cm,yshift=0cm,scale=2]
    \foreach \i in {0,1,2} \foreach \j in {0,1} {
      \begin{scope}[shift={(-\i,-\j)}]
        \clip (\i,\j) rectangle (\i+1,\j+1);
        \fill[addcol,opacity=0.28] \blob;
        \draw[addcol] \blob;
      \end{scope}}
    \draw[thick] (0,0) rectangle (1,1);
    \node[below,font=\small] at (0.5,0) {$P=[0,1)^2$};
  \end{scope}
\end{tikzpicture}
\caption{Blichfeldt's lemma for $\Lat=\Z^2$. The pieces of $S$ in the different unit squares are
translated into one fundamental square $P$. Their total area is $\vol(S)>1=\vol(P)$, so two pieces
overlap (the darker regions). A point of overlap comes from two points of $S$ that differ by a lattice
vector.}
\label{NT:fig:blichfeldt}
\end{figure}

Notice the parallel. In Lecture~2 the periodized function was bounded by $1$ (packing condition), and
the conclusion was that its average, and so the density, is at most $1$. Here the average exceeds $1$, so
the function is somewhere at least $2$: the translates $S+\lambda$ are \emph{not} a packing. Both are
pigeonhole arguments for the periodization of a function.

\begin{theorem}[Minkowski's convex body theorem, 1889/1896]\label{NT:thm:minkowski}
Let $K\subset\R^n$ be convex and symmetric ($K=-K$). If $\vol(K)>2^n\covol(\Lat)$, then $K$ contains a
nonzero point of $\Lat$. If $K$ is also compact, $\vol(K)\ge2^n\covol(\Lat)$ suffices.
\end{theorem}

\begin{proof}
Apply Lemma~\ref{NT:lem:blichfeldt} to $S=\frac12K$, of volume $2^{-n}\vol(K)>\covol(\Lat)$. This
gives $s\ne s'$ in $\frac12K$ with $s-s'\in\Lat$. By symmetry $-s'\in\frac12K$, and by convexity
$\frac12(2s)+\frac12(-2s')=s-s'\in K$. For compact $K$ apply the first part to $(1+\varepsilon)K$ for
each $\varepsilon>0$. Each of these contains a nonzero lattice point, and they all lie in the bounded set
$2K$, which contains only finitely many lattice points. So one nonzero point lies in $(1+\varepsilon)K$
for arbitrarily small $\varepsilon$, hence in $K$.
\end{proof}

\begin{remark}[Minkowski and the trivial bound]
Applied to the ball $K=\Ball^n(\lambda_1)$, Minkowski gives $V_n\lambda_1^n\le2^n\covol(\Lat)$, which
is $\dens(\Lat)\le1$, the trivial bound that the balls of radius $\lambda_1/2$ are disjoint. (So the two
estimates quoted in Lecture~4 are actually the same inequality.) For sphere packing Minkowski's theorem
is therefore weak. Its power lies in two other directions: in arithmetic, where the lattice is
cut out by congruences, and in the bodies $K$ other than balls. Here are the two classic examples.
\end{remark}

\begin{theorem}[Fermat]\label{NT:thm:fermat}
Every prime $p\equiv1\pmod4$ is a sum of two squares.
\end{theorem}

\begin{proof}
Since $p\equiv1\pmod4$, $-1$ is a square modulo $p$; for instance $u=\big(\frac{p-1}2\big)!$ satisfies
$u^2\equiv-1$ by Wilson's theorem. Let
\[
  \Lat=\{(x,y)\in\Z^2:\ y\equiv ux\pmod p\},
\]
a sublattice of index $p$ in $\Z^2$, so $\covol(\Lat)=p$. The open disk of radius $\sqrt{2p}$ has area
$2\pi p>4p$, so by Theorem~\ref{NT:thm:minkowski} it contains some $(x,y)\in\Lat\setminus0$. Then
$0<x^2+y^2<2p$ and $x^2+y^2\equiv x^2(1+u^2)\equiv0\pmod p$. Hence $x^2+y^2=p$.
\end{proof}

\begin{figure}[ht]
\centering
\begin{tikzpicture}[scale=0.4]
  \foreach \x in {-6,...,6} \foreach \y in {-6,...,6} \fill[gray!40] (\x,\y) circle (1.5pt);
  \foreach \x in {-6,...,6} \foreach \y in {-6,...,6} {
    \pgfmathtruncatemacro{\t}{mod(\y-5*\x+130,13)}
    \ifnum\t=0 \fill (\x,\y) circle (4pt); \fi }
  \draw[thick,addcol] (0,0) circle ({sqrt(26)});
  \draw[-{Stealth},flagcol,thick] (0,0) -- (2,-3) node[below right,font=\scriptsize]{$(2,-3)$};
  \draw[-{Stealth},flagcol,thick] (0,0) -- (3,2) node[right,font=\scriptsize]{$(3,2)$};
  \node[addcol,font=\small] at (5.6,4.9) {$\abs v=\sqrt{26}$};
\end{tikzpicture}
\caption{Fermat by Minkowski for $p=13$, $u=5$ ($5^2\equiv-1$). The lattice
$\{y\equiv5x\bmod13\}$ (black) has covolume $13$. The disk of area $26\pi>4\cdot13$ must contain
lattice vectors, and they have $x^2+y^2=13$. Here the lattice is a square lattice of side $\sqrt{13}$,
with reduced basis $(3,2)$, $(2,-3)$.}
\label{NT:fig:fermat}
\end{figure}

\begin{theorem}[Lagrange's four square theorem, 1770]\label{NT:thm:lagrange4}
Every positive integer is a sum of four squares.
\end{theorem}

\begin{proof}
By Euler's identity, which expresses $(a_1^2+\dots+a_4^2)(b_1^2+\dots+b_4^2)$ as a sum of four squares
(it is the multiplicativity of the quaternion norm), it suffices to treat primes, and $2=1^2+1^2$. Let
$p$ be an odd prime. The sets $\{a^2\bmod p\}$ and $\{-1-b^2\bmod p\}$ each have $\frac{p+1}2$
elements, so they meet: $a^2+b^2+1\equiv0\pmod p$ for some $a,b$. Let
\[
  \Lat=\{x\in\Z^4:\ x_3\equiv ax_1+bx_2,\ \ x_4\equiv bx_1-ax_2\pmod p\},
\]
of index $p^2$ in $\Z^4$. For $x\in\Lat$,
\[
  x_1^2+x_2^2+x_3^2+x_4^2\equiv(x_1^2+x_2^2)(1+a^2+b^2)\equiv0\pmod p .
\]
The open ball of radius $\sqrt{2p}$ in $\R^4$ has volume $\frac{\pi^2}{2}(2p)^2=2\pi^2p^2>16p^2$, so
it contains a nonzero $x\in\Lat$, and then $0<\abs x^2<2p$ and $p\mid\abs x^2$ force $\abs x^2=p$.
\end{proof}

Three squares are more delicate. There is no product identity for sums of three squares, and the known
geometric proof (Ankeny, 1957) needs Dirichlet's theorem on primes in arithmetic progressions in
addition to Minkowski's theorem.

\begin{theorem}[Legendre 1798; Gauss 1801]\label{NT:thm:legendre}
A positive integer is a sum of three squares if and only if it is not of the form $4^a(8b+7)$.
\end{theorem}

The ``only if'' half is a computation mod $8$. The ``if'' half needs the theory of ternary forms (or
Dirichlet's theorem on primes in progressions). Which other forms represent \emph{every} positive integer?
Ramanujan (1917) listed $55$ diagonal quaternary forms $ax^2+by^2+cz^2+dw^2$ with this property. One of
them, $x^2+2y^2+5z^2+5w^2$, fails to represent $15$, and the other $54$ are exactly the universal ones. The definitive answer is surprisingly finite; see the $15$ and $290$ theorems in
the Gallery.

%======================================================================
\section{Theta series and modular forms}
\label{NT:sec:theta}

All the representation numbers of a lattice can be packaged into one function.

\begin{definition}\label{NT:def:theta}
The \emph{theta series} of $\Lat$ is
\[
  \Theta_\Lat(\tau)=\sum_{x\in\Lat}e^{\pi i\tau\abs x^2}=\sum_{t\ge0}N_t\,e^{\pi i\tau t},
  \qquad \tau\in\mathbb H,
\]
where $N_t$ is the number of vectors of squared length $t$. The series converges absolutely and locally
uniformly on $\mathbb H$, because $\abs{e^{\pi i\tau\abs x^2}}=e^{-\pi\operatorname{Im}\tau\abs x^2}$.
\end{definition}

For example $\Theta_{\Z}(\tau)=\sum_{k\in\Z}q^{k^2}$ with $q=e^{\pi i\tau}$ is Jacobi's theta function,
and $\Theta_{\Lat\oplus\Lat'}=\Theta_\Lat\Theta_{\Lat'}$, so
$\Theta_{\Z^2}=\Theta_\Z^2=\sum_mr_2(m)q^m$. Formula \eqref{NT:eq:r2} and Jacobi's four square formula
\begin{equation}\label{NT:eq:r4}
  r_4(m)=8\sum_{d\mid m,\ 4\nmid d}d
\end{equation}
are identities between $q$-series (both were checked numerically for small $m$). The tool that makes
theta series powerful is the following.

\begin{theorem}[Poisson summation over a lattice]\label{NT:thm:poisson}
Let $f\colon\R^n\to\C$ be continuous and satisfy $\abs{f(x)},\abs{\hat f(x)}\le C(1+\abs x)^{-n-\eta}$ for some $C,\eta>0$,
where $\hat f(y)=\int_{\R^n}f(x)e^{-2\pi i\ip xy}dx$. Then for every $v\in\R^n$,
\[
  \sum_{x\in\Lat}f(x+v)=\frac{1}{\covol(\Lat)}\sum_{y\in\Lat^*}\hat f(y)\,e^{2\pi i\ip yv},
  \qquad \Lat^*=\{y\in\R^n:\ip xy\in\Z\ \text{for all }x\in\Lat\}.
\]
Both sides converge absolutely.
\end{theorem}

\begin{proof}[Sketch]
The left side $F(v)$ is a continuous $\Lat$-periodic function, so it is a function on the torus
$\R^n/\Lat$, whose characters are $v\mapsto e^{2\pi i\ip yv}$, $y\in\Lat^*$. Its Fourier coefficients are
\[
  \frac{1}{\covol(\Lat)}\int_{\R^n/\Lat}F(v)e^{-2\pi i\ip yv}dv
  =\frac{1}{\covol(\Lat)}\int_{\R^n}f(v)e^{-2\pi i\ip yv}dv=\frac{\hat f(y)}{\covol(\Lat)},
\]
by unfolding the sum against the integral. These coefficients are absolutely summable by the decay of
$\hat f$, so the Fourier series converges uniformly to $F$.
\end{proof}

The \emph{dual lattice} $\Lat^*$ has $\covol(\Lat^*)=1/\covol(\Lat)$ and Gram matrix $G^{-1}$. $\Lat$ is
integral iff $\Lat\subseteq\Lat^*$, and \emph{unimodular} iff $\Lat=\Lat^*$ (equivalently, it is integral with
$\covol=1$). It is \emph{even} if all $\abs x^2$ are even. $\Z^n$ is unimodular but odd; $E_8$ is even
unimodular: it is integral and has determinant $1$ (Lecture~6).

Apply Theorem~\ref{NT:thm:poisson} to a Gaussian. For $\tau\in\mathbb H$ the function
$f(x)=e^{\pi i\tau\abs x^2}$ has $\hat f(y)=(\tau/i)^{-n/2}e^{\pi i(-1/\tau)\abs y^2}$ (principal branch;
for $\tau=it$ this is the familiar fact that $e^{-\pi t\abs x^2}$ has Fourier transform
$t^{-n/2}e^{-\pi\abs y^2/t}$, and both sides are holomorphic in $\tau$). With $v=0$ this gives $\Theta_\Lat(\tau)=\covol(\Lat)^{-1}(i/\tau)^{n/2}\,\Theta_{\Lat^*}(-1/\tau)$; now replace $\tau$ by $-1/\tau$.

\begin{corollary}[Theta transformation formula]\label{NT:cor:thetatrans}
$\displaystyle \Theta_\Lat\Big(-\frac1\tau\Big)=\Big(\frac\tau i\Big)^{n/2}\frac{1}{\covol(\Lat)}\,\Theta_{\Lat^*}(\tau).$
\end{corollary}

So the theta series of a lattice controls that of its dual, through the substitution $\tau\mapsto-1/\tau$.
For an even unimodular lattice we get two symmetries:
\begin{equation}\label{NT:eq:modular}
  \Theta_\Lat(\tau+1)=\Theta_\Lat(\tau),\qquad
  \Theta_\Lat\Big(-\frac1\tau\Big)=\Big(\frac\tau i\Big)^{n/2}\Theta_\Lat(\tau).
\end{equation}
The maps $T\colon\tau\mapsto\tau+1$ and $S\colon\tau\mapsto-1/\tau$ generate $\SL_2(\Z)$, the group
whose fundamental domain we drew in Figure~\ref{NT:fig:modular}. The fact that \emph{lattices in the
plane} are parametrized by $\SL_2(\Z)\backslash\mathbb H$ and the fact that \emph{theta series of lattices
in any dimension} transform under $\SL_2(\Z)$ are two separate appearances of the same group.

\begin{theorem}\label{NT:thm:8n}
If $\R^n$ contains an even unimodular lattice, then $8\mid n$.
\end{theorem}

\begin{proof}
Let $f=\Theta_\Lat$. Since $\Theta_{\Lat\oplus\Lat}=f^2$ and $\Lat\oplus\Lat$ is again even unimodular,
we may pass to $\Lat\oplus\Lat$ or to $\Lat^{\oplus4}$. So if $8\nmid n$ we may assume $n\equiv4\pmod 8$.
Put $k=n/2\equiv2\pmod4$. Then $(\tau/i)^{k}=i^{-k}\tau^k=-\tau^k$, and \eqref{NT:eq:modular} reads
$f(\tau+1)=f(\tau)$, $f(-1/\tau)=-\tau^kf(\tau)$. For $\gamma=\left(\begin{smallmatrix}a&b\\c&d\end{smallmatrix}\right)\in\SL_2(\Z)$
define $(f|\gamma)(\tau)=(c\tau+d)^{-k}f(\gamma\tau)$. The chain rule for $\gamma\tau$ shows
that $(f|\gamma)|\gamma'=f|(\gamma\gamma')$, and $f|(-I)=f$ since $k$ is even. With
$T=\left(\begin{smallmatrix}1&1\\0&1\end{smallmatrix}\right)$ and $S=\left(\begin{smallmatrix}0&-1\\1&0\end{smallmatrix}\right)$
we have $f|T=f$ and $f|S=-f$. A direct multiplication gives $(ST)^3=-I$. Therefore
\[
  f=f|(-I)=f|STSTST=(-1)^3f=-f,
\]
so $f\equiv0$. This contradicts $f(\tau)\to1$ as $\operatorname{Im}\tau\to\infty$ (only $x=0$ survives).
\end{proof}

Conversely, even unimodular lattices exist in every dimension $8k$: $E_8^{\oplus k}$, or
$D_{8k}^+=D_{8k}\cup(D_{8k}+\frac12\mathbf 1)$ (Exercise~\ref{NT:ex:D16}; $E_8=D_8^+$ by Lecture~6).

\subsection*{The theta series of \texorpdfstring{$E_8$}{E8}}

For $8\mid n$ the factor $(\tau/i)^{n/2}$ equals $\tau^{n/2}$, and \eqref{NT:eq:modular} says that
$\Theta_\Lat$ is a \emph{modular form of weight $n/2$} for $\SL_2(\Z)$: a holomorphic function on
$\mathbb H$, bounded as $\operatorname{Im}\tau\to\infty$, with $f(\gamma\tau)=(c\tau+d)^{n/2}f(\tau)$.
Write $\mathsf q=e^{2\pi i\tau}$ (so $\mathsf q=q^2$). The basic theorem on such forms (see Serre,
\emph{A Course in Arithmetic}, Ch.~VII) is that the space $M_k$ of modular forms of weight $k$ is
finite-dimensional, with
\[
  M_4=\C E_4,\qquad M_8=\C E_4^2,\qquad M_{12}=\C E_4^3\oplus\C\Delta,
\]
where $E_4=1+240\sum_{m\ge1}\sigma_3(m)\mathsf q^m$ is the Eisenstein series
($\sigma_3(m)=\sum_{d\mid m}d^3$) and $\Delta=\mathsf q\prod_{m\ge1}(1-\mathsf q^m)^{24}$. Since
$\Theta_{E_8}=\sum_mN_{2m}\mathsf q^m$ lies in $M_4$ and has constant term $1$,
\begin{equation}\label{NT:eq:E8theta}
  \Theta_{E_8}=E_4,\qquad\text{i.e.}\qquad
  \#\{x\in E_8:\abs x^2=2m\}=240\,\sigma_3(m).
\end{equation}
For $m=1$ this is $240$, the root count from Lecture~6 and the dimension~$8$ kissing number. For $m=2$
it predicts $240\cdot9=2160$ vectors of norm $4$, and for $m=3$, $240\cdot28=6720$. A brute-force count
of $E_8$ vectors confirms $240,\,2160,\,6720$. \emph{A single modular form knows every shell of $E_8$.}

\subsection*{Hearing the shape of a drum, and Niemeier's lattices}

The flat torus $\R^n/\Lat$ has Laplace eigenfunctions $e^{2\pi i\ip yx}$, $y\in\Lat^*$, with eigenvalues
$4\pi^2\abs y^2$. So its spectrum, with multiplicities, is recorded by $\Theta_{\Lat^*}$.

\begin{theorem}[Milnor 1964]\label{NT:thm:milnor}
The $16$-dimensional flat tori $\R^{16}/(E_8\oplus E_8)$ and $\R^{16}/D_{16}^+$ have the same Laplace
spectrum but are not isometric.
\end{theorem}

\begin{proof}
Both lattices are even unimodular, so both theta series lie in $M_8=\C E_4^2$ and have constant term $1$.
So they are equal, and the lattices are self-dual, which makes the spectra equal. (Numerically, both have
$480$ vectors of norm $2$ and $61920$ of norm $4$.) An isometry of the tori would lift to an isometry of
$\R^{16}$ carrying one lattice onto the other. But the $480$ norm-$2$ vectors of $E_8\oplus E_8$ span it,
while those of $D_{16}^+$ (the vectors $\pm e_i\pm e_j$) span only the index-$2$ sublattice $D_{16}$.
\end{proof}

This was the first counterexample to the question ``can one hear the shape of a drum?'', two years
before Kac named the question. In dimension $24$, $M_{12}$ is two-dimensional, so $\Theta_\Lat$ is
determined by its constant term and by the number of roots. Taking zero roots gives
\[
  \Theta_{\Lambda_{24}}=E_4^3-720\Delta=1+196560\,\mathsf q^2+16773120\,\mathsf q^3+\cdots,
\]
and $196560$ is the kissing number of the Leech lattice. Niemeier (1973) classified the even unimodular
lattices in dimension $24$. There are exactly $24$ of them, determined by their root systems (each one of
rank $24$ or empty), and the Leech lattice is the unique one without roots (Conway 1969).

%======================================================================
\section{Poisson summation as the engine: LP bounds}
\label{NT:sec:LP}

Lecture~3 mentioned that the solutions in dimensions $8$ and $24$ use ``linear programming duality'' and a
``magic'' auxiliary function. With Poisson summation in hand, the argument for lattices takes only a few lines.

\begin{definition}
Call $f\colon\R^n\to\R$ \emph{admissible} if $f$ and $\hat f$ satisfy the decay hypothesis of
Theorem~\ref{NT:thm:poisson}, and $f$ is even (so $\hat f$ is real).
\end{definition}

\begin{theorem}[Cohn--Elkies bound for lattices]\label{NT:thm:CElattice}
Let $f$ be admissible with $f(0)=\hat f(0)>0$, and suppose
\begin{enumerate}[label=(\roman*),nosep]
\item $f(x)\le0$ for $\abs x\ge r$, and
\item $\hat f(y)\ge0$ for all $y$.
\end{enumerate}
Then every lattice packing in $\R^n$ has density at most $V_n(r/2)^n$.
\end{theorem}

\begin{proof}
Scale the lattice so that $\lambda_1(\Lat)=r$. Apply Poisson summation with $v=0$:
\[
  f(0)\;\ge\;f(0)+\sum_{x\in\Lat\setminus0}f(x)\;=\;\frac{1}{\covol(\Lat)}\sum_{y\in\Lat^*}\hat f(y)
  \;\ge\;\frac{\hat f(0)}{\covol(\Lat)}.
\]
The first inequality is (i), since every $x\ne0$ has $\abs x\ge r$, and the last is (ii), dropping
all $y\ne0$. Hence $\covol(\Lat)\ge\hat f(0)/f(0)=1$, and $\dens(\Lat)=V_n(r/2)^n/\covol(\Lat)\le V_n(r/2)^n$.
\end{proof}

The full theorem drops the lattice assumption at almost no cost.

\begin{theorem}[Cohn--Elkies 2003]\label{NT:thm:CE}
Under the same hypotheses, every sphere packing in $\R^n$ has density at most $V_n(r/2)^n$.
\end{theorem}

\begin{proof}
By Exercise~\ref{L1:ex:periodic-approx} of Lecture~1 it suffices to treat periodic packings, with
centres $\bigcup_{j=1}^N(\Lat+v_j)$, the $v_j$ pairwise incongruent modulo $\Lat$, scaled to have
minimal distance $r$. Sum Poisson summation over all pairs
$(j,k)$ with $v=v_j-v_k$:
\[
  \sum_{j,k=1}^N\sum_{x\in\Lat}f(x+v_j-v_k)
  =\frac{1}{\covol(\Lat)}\sum_{y\in\Lat^*}\hat f(y)\,\Big|\sum_{j=1}^Ne^{2\pi i\ip y{v_j}}\Big|^2 .
\]
On the left, $x+v_j-v_k$ is the difference of two centres. It is $0$ exactly when $j=k$ and $x=0$, and
otherwise it has length at least $r$. So the left side is at most $Nf(0)$. On the right every term is
$\ge0$, and the $y=0$ term is $N^2\hat f(0)/\covol(\Lat)$. So $N/\covol(\Lat)\le f(0)/\hat f(0)=1$. The
density is $(N/\covol(\Lat))V_n(r/2)^n$.
\end{proof}

\begin{remark}[Blichfeldt's gauge versus Cohn--Elkies]\label{NT:rem:blichfeldt-vs-CE}
Recall the proof that $f_0=1-\frac12\abs x^2$ is a gauge (Lemma~\ref{L2:lem:f0}). It rested on the identity
\[
  \sum_{i<j}\abs{x_i-x_j}^2=m\sum_i\abs{x_i}^2-\Big|\sum_ix_i\Big|^2,
\]
and the inequality came from \emph{throwing away a square}, $\abs{\sum_ix_i}^2\ge0$. The proof of
Theorem~\ref{NT:thm:CE} throws away the squares $\hat f(y)\abs{\sum_je^{2\pi i\ip y{v_j}}}^2\ge0$, one for
each frequency $y\ne0$. Blichfeldt uses one quadratic ``moment'', while Cohn--Elkies use all the
exponentials. There is also a difference in what the auxiliary function must satisfy:
\begin{itemize}[nosep]
\item a Blichfeldt gauge is a condition in physical space only: $f\ge0$ and
  $\sum_if(\varphi_i^{-1}x)\le1$ for every packing;
\item a Cohn--Elkies function has a sign condition in \emph{both} spaces: $f\le0$ beyond $r$, and
  $\hat f\ge0$ everywhere.
\end{itemize}
The uncertainty principle makes the second kind of condition hard to satisfy, and this is why the
optimal functions are rare and remarkable.
\end{remark}

\paragraph{When is the bound sharp?} Suppose $\Lat$ is unimodular, $\lambda_1(\Lat)=r$, and $f$ gives
equality. Then both inequalities in the proof of Theorem~\ref{NT:thm:CElattice} are equalities, so
\[
  f(x)=0\ \text{ for all }x\in\Lat\setminus0,\qquad \hat f(y)=0\ \text{ for all }y\in\Lat^*\setminus0.
\]
For $E_8$ ($r=\sqrt2$) this says that $f$ and $\hat f$ vanish on all spheres $\abs x=\sqrt{2m}$, $m\ge1$. Since
$f\le0$ beyond $\sqrt2$ and $\hat f\ge0$ everywhere, these zeros must have even order (they are double roots),
except that $f$ may change sign, with a simple root, at $\sqrt2$ itself.
Cohn and Elkies found numerically that such functions seem to exist in dimensions $8$ and $24$ (and
$2$), but could not construct them.

\paragraph{Viazovska's idea.} Look for $f$ radial and built out of Gaussians:
\[
  f(x)=\int_{\Gamma}\phi(z)\,e^{\pi iz\abs x^2}\,dz
\]
along a path $\Gamma$ in $\mathbb H$. In dimension $8$ the Fourier transform of $e^{\pi iz\abs x^2}$ is
$z^{-4}e^{\pi i(-1/z)\abs y^2}$ (the computation before Corollary~\ref{NT:cor:thetatrans}). Substituting
$w=-1/z$ shows that $\hat f$ is given by the same kind of integral with weight $\phi(-1/w)\,w^{2}$ in place of
$\phi(w)$, over the image path. So the Fourier transform acts on the \emph{weight} $\phi$ through
$z\mapsto-1/z$, and $f$ is a Fourier eigenfunction when $\phi$ transforms like a modular form of weight $-2$
(up to the change of path). Viazovska (2016) built two eigenfunctions of the Fourier transform, with
eigenvalues $\pm1$: one from a weakly holomorphic \emph{quasimodular} form for $\SL_2(\Z)$ (a rational
expression in $E_2,E_4,E_6$ and $\Delta$), and one from a weakly holomorphic modular form for the subgroup
$\Gamma(2)$ (built from Jacobi theta functions). A factor $\sin^2(\pi\abs x^2/2)$ supplies the
double roots at $\abs x^2\in2\Z$. A suitable linear combination of the two is the magic function, and it proves that $E_8$ is the densest
packing in $\R^8$. A week later, Cohn, Kumar, Miller, Radchenko and Viazovska did the same in dimension
$24$ for the Leech lattice. The theta series of Section~\ref{NT:sec:theta} and the magic functions use
the same symmetry $\tau\mapsto-1/\tau$, once applied to a lattice and once to an auxiliary function.

\paragraph{Universal optimality.} The same circle of ideas gives a much stronger statement (Gallery,
Theorem~\ref{NT:gal:universal}): $E_8$ and Leech minimize energy for every completely monotone function of
squared distance, Gaussians $e^{-\alpha\abs x^2}$ included. For Gaussian potentials, Poisson summation
turns the energy of a lattice into the energy of its dual. The proof needs a \emph{Fourier interpolation
formula}: a radial Schwartz function on $\R^8$ is determined by the values of $f$ and $\hat f$ (and
their radial derivatives) on the spheres $\abs x=\sqrt{2m}$.

%======================================================================
\section{Averaging over all lattices: Siegel and Minkowski--Hlawka}
\label{NT:sec:siegel}

So far we have looked at special lattices. At the opposite extreme, take a \emph{random} lattice. Let
$X_n=\SL_n(\R)/\SL_n(\Z)$ be the space of lattices of covolume $1$ in $\R^n$. It carries a unique
$\SL_n(\R)$-invariant probability measure $\mu$ (Minkowski, Siegel). For $n=2$, after quotienting by
rotations, $X_2$ becomes the modular surface of Figure~\ref{NT:fig:modular}, and $\mu$ becomes the
hyperbolic area $\frac{3}{\pi}\frac{dx\,dy}{y^2}$ restricted to $\mathcal F$.

\begin{theorem}[Siegel's mean value theorem, 1945]\label{NT:thm:siegel}
For $n\ge2$ and every integrable $f\colon\R^n\to\R$,
\[
  \int_{X_n}\sum_{x\in\Lat\setminus0}f(x)\,d\mu(\Lat)=\int_{\R^n}f(x)\,dx,
  \qquad
  \int_{X_n}\sum_{x\in\Lat\ \mathrm{primitive}}f(x)\,d\mu(\Lat)=\frac{1}{\zeta(n)}\int_{\R^n}f(x)\,dx .
\]
\end{theorem}

The two formulas are equivalent, because every nonzero vector is uniquely $kx$ with $k\ge1$ and $x$
primitive, and $\sum_kk^{-n}=\zeta(n)$. In words: \emph{on average, a lattice looks like Lebesgue measure.}
From this one gets a lower bound in a few lines.

\begin{corollary}[Minkowski--Hlawka]\label{NT:cor:MH}
For $n\ge2$, $\sup_\Lat\dens(\Lat)\ge\dfrac{2\zeta(n)}{2^n}$.
\end{corollary}

\begin{proof}
Let $B$ be a ball centred at $0$ with $\vol(B)<2\zeta(n)$, and apply the primitive form of
Theorem~\ref{NT:thm:siegel} to $f=\chi_B$. The average number of primitive vectors in $B$ is
$\vol(B)/\zeta(n)<2$. Primitive vectors come in pairs $\pm x$, so some $\Lat\in X_n$ has no primitive
vector in $B$, and hence no nonzero vector at all (if $kx\in B$ with $k\ge1$, then $x\in B$). Balls of
half the radius centred at $\Lat$ form a packing of density $\vol(B)/2^n$. Let
$\vol(B)\uparrow2\zeta(n)$.
\end{proof}

This is the ``$2\zeta(d)2^{-d}$'' of the remark after Proposition~\ref{L1:prop:saturated-bound} in
Lecture~1, and it is still within a polynomial factor of the best known lower bounds (for lattices,
Klartag's 2025 bound $c\,n^2\,2^{-n}$). Upper bounds are exponentially larger
(Kabatiansky--Levenshtein, Lecture~3). It is still unknown whether the best lattices in high dimension are
exponentially better than random ones.

Siegel's theorem treats $X_n$ as a probability space. Interlude~C treats it as a dynamical system: the
diagonal flow on $X_n$, Mahler's criterion as a statement about the cusp, and equidistribution results
that sharpen Siegel's average.

%======================================================================
\section{A dictionary}
\label{NT:sec:dictionary}

\begin{center}
\renewcommand{\arraystretch}{1.25}
\begin{tabular}{p{0.42\textwidth}p{0.50\textwidth}}
\toprule
Packings and lattices & Number theory\\
\midrule
lattice up to rotation & positive definite form up to $\GL_n(\Z)$\\
$\lambda_1^2$, $\covol^2$ & $\min Q$, $\det G$\\
densest lattice packing & Hermite's constant $\gamma_n$\\
planar lattices, Figure~\ref{NT:fig:modular} & $\SL_2(\Z)\backslash\mathbb H$, reduced binary forms\\
number of points in a ball & $\sum_{m\le r^2}r_Q(m)$; Gauss circle problem\\
kissing number of $\Lat$ & number of representations of $\min Q$\\
Blichfeldt's gauge (Lecture 2) & Blichfeldt's lemma, Minkowski's theorem\\
shells of $E_8$, Leech & coefficients of $E_4$, $E_4^3-720\Delta$\\
dual lattice $\Lat^*$ & $G^{-1}$; $\tau\mapsto-1/\tau$\\
LP bound, magic function & Poisson summation, (quasi)modular forms\\
random lattice & Siegel's mean value theorem\\
\bottomrule
\end{tabular}
\end{center}

%======================================================================
\section{Gallery of theorems}
\label{NT:sec:gallery}

Classical results first, then recent ones. Each comes with a sentence on why it belongs in this course.

\begin{theorem}[Hermite's inequality; Hermite 1850]\label{NT:gal:hermite}
$\gamma_n\le(4/3)^{(n-1)/2}$ for all $n\ge1$.
\end{theorem}
Published in Hermite's letters to Jacobi, J.\ reine angew.\ Math.\ 40 (1850). For $n=2$ it is
Theorem~\ref{NT:thm:gamma2}. The proof (Exercise~\ref{NT:ex:hermite}) is the row-stacking argument of
Lecture~1, run by induction on the dimension.

\begin{theorem}[The known Hermite constants]\label{NT:gal:hermiteknown}
The values of $\gamma_n$ in the table of Section~\ref{NT:sec:forms} are correct, and in each of these
dimensions the optimal lattice is unique up to similarity.
\end{theorem}
Lagrange (1773) and Gauss (1831, in his review of Seeber) for $n=2,3$; Korkine and Zolotarev, Math.\
Ann.\ 6 (1873) and 11 (1877), for $n=4,5$; Blichfeldt, Math.\ Z.\ 39 (1935), for $n=6,7,8$ (uniqueness
there is due to Vet\v{c}inkin, 1980); H.~Cohn and A.~Kumar, \emph{Optimality and uniqueness of the Leech
lattice among lattices}, Ann.\ of Math.\ 170 (2009), for $n=24$. The case $n=24$ uses the Cohn--Elkies
bound of Section~\ref{NT:sec:LP} together with an argument special to lattices. These are exactly the
lattices of Lectures~4--6.

\begin{theorem}[Voronoi 1908]\label{NT:gal:voronoi}
A positive definite form $Q$ is a local maximum of $\gamma(Q)$ (an \emph{extreme} form) if and only if it is
\emph{perfect} (it is the only form with its minimum value and its set of minimal vectors; equivalently,
the matrices $xx^T$, $x$ minimal, span the space of symmetric matrices) and \emph{eutactic} ($G^{-1}$ is a
combination of the matrices $xx^T$, $x$ minimal, with strictly positive coefficients). Up to equivalence and scaling there are finitely many perfect forms in each dimension.
\end{theorem}
G.~Voronoi, J.\ reine angew.\ Math.\ 133 (1908). The numbers of perfect forms in dimensions $2,\dots,8$
are $1,1,2,3,7,33,10916$ (dimension $7$: Jaquet-Chiffelle, Ann.\ Inst.\ Fourier 43 (1993); dimension $8$:
Dutour Sikiri\'c, Schürmann and Vallentin, Electron.\ Res.\ Announc.\ AMS 13 (2007)). This is a finite
algorithm for the lattice packing problem in each fixed dimension. Eutaxy is the lattice analogue of the
force balanced contact graphs of Lecture~10: the minimal vectors are the contacts, and the positive coefficients are
the stresses.

\begin{theorem}[Class number one; Heegner 1952, Baker 1966, Stark 1967]\label{NT:gal:classno}
The imaginary quadratic fields $\Q(\sqrt{-d})$ with class number $1$ are exactly those with
$d=1,2,3,7,11,19,43,67,163$.
\end{theorem}
Gauss conjectured the list. Heegner's proof (Math.\ Z.\ 56, 1952) was not accepted at first. Baker
(Mathematika 13, 1966) and Stark (Michigan Math.\ J.\ 14, 1967) gave independent proofs, and Heegner's
argument was later found to be essentially correct. In the language of Remark~\ref{NT:rem:classno}: for the
corresponding discriminants $D=-3,-4,-7,-8,-11,-19,-43,-67,-163$ there is exactly one reduced primitive form.
For $D=-3$ and $D=-4$ it is the hexagonal and the square lattice. For $D=-163$ it is $x^2+xy+41y^2$, and
Euler's observation that $k^2+k+41$ is prime for $0\le k\le39$ is a shadow of $h(-163)=1$ (Rabinowitsch, 1913).

\begin{theorem}[The $15$ theorem; Conway--Schneeberger 1993, Bhargava 2000]\label{NT:gal:15}
A positive definite quadratic form with integer Gram matrix represents every positive integer if and only if it
represents $1,2,3,5,6,7,10,14,15$.
\end{theorem}
Conway and Schneeberger did not publish their proof. M.~Bhargava, \emph{On the Conway--Schneeberger
fifteen theorem}, in \emph{Quadratic Forms and Their Applications} (Dublin 1999), Contemp.\ Math.\ 272,
AMS (2000), 27--37, gave a short proof, built by ``escalating'' lattices: one adds a vector whose
norm is the smallest integer not yet represented. This is the arithmetic version of the laminated
lattices of Lecture~6. None of the nine numbers can be omitted.

\begin{theorem}[The $290$ theorem; Bhargava--Hanke]\label{NT:gal:290}
A positive definite integer-valued quadratic form represents every positive integer if and only if it
represents
the $29$ numbers 1, 2, 3, 5, 6, 7, 10, 13, 14, 15, 17, 19, 21, 22, 23, 26, 29, 30, 31, 34, 35, 37, 42, 58, 93,
110, 145, 203, 290.
\end{theorem}
Conjectured by Conway. M.~Bhargava and J.~Hanke, \emph{Universal quadratic forms and the 290-theorem},
preprint (2011). It is widely cited, but as of 2026 it has not appeared in a journal.
``Integer-valued'' allows odd off-diagonal coefficients, such as $x^2+xy+y^2$; that is why the list is longer.

\begin{theorem}[Hardy 1916; Huxley 2003]\label{NT:gal:circle}
In the circle problem, $E(r)$ is not $o\big(r^{1/2}(\log r)^{1/4}\big)$. On the other hand,
$E(r)=O_\varepsilon(r^{131/208+\varepsilon})$ for every $\varepsilon>0$.
\end{theorem}
The lower bound is in G.~H.~Hardy, \emph{On Dirichlet's divisor problem}, Proc.\ London Math.\ Soc.\ (2)
15 (1916). The upper bound is in M.~N.~Huxley, \emph{Exponential sums and lattice points III}, Proc.\
London Math.\ Soc.\ (3) 87 (2003), 591--609. The logarithmic factors in the lower bound were sharpened
by Hafner (Invent.\ Math.\ 63, 1981) and by Soundararajan (Int.\ Math.\ Res.\ Not.\ 2003, no.~36). For
the arXiv improvements of the upper bound, see Section~\ref{NT:sec:circle}. This is the prototype of a boundary-layer error in a density count.

\begin{theorem}[Milnor 1964; Niemeier 1973]\label{NT:gal:niemeier}
There are non-isometric isospectral flat tori in dimension $16$. There are exactly $24$ even unimodular
lattices in $\R^{24}$, and the Leech lattice is the only one with no vectors of norm $2$.
\end{theorem}
J.~Milnor, \emph{Eigenvalues of the Laplace operator on certain manifolds}, Proc.\ Nat.\ Acad.\ Sci.\
USA 51 (1964), 542; H.-V.~Niemeier, J.\ Number Theory 5 (1973), 142--178; J.~H.~Conway, Invent.\ Math.\ 7
(1969), 137--142, for the characterization of Leech. Both are consequences of the theta series being a
modular form (Section~\ref{NT:sec:theta}).

\begin{theorem}[Cohn--Elkies 2003]\label{NT:gal:CE}
Theorem~\ref{NT:thm:CE} holds.
\end{theorem}
H.~Cohn and N.~Elkies, \emph{New upper bounds on sphere packings I}, Ann.\ of Math.\ 157 (2003), 689--714.
Their numerical bounds were the best known in many dimensions, and in dimensions $8$ and $24$ they came
extremely close to the densities of $E_8$ and Leech. This is the bridge from Blichfeldt's method to the solution in dimensions $8$ and $24$.

\begin{conjecture}[Cohn--Elkies]\label{NT:gal:CE2}
The linear programming bound of Theorem~\ref{NT:thm:CE} is sharp in dimension $2$, that is, its optimum
equals $\pi/\sqrt{12}$.
\end{conjecture}
Numerical optimization supports it, but no exact auxiliary function is known. Thue's
theorem is known by other means (Lecture~1), so what is open is whether \emph{this method} reproves it.

\begin{theorem}[Viazovska 2016; Cohn--Kumar--Miller--Radchenko--Viazovska 2016]\label{NT:gal:viazovska}
The $E_8$ lattice packing is the densest sphere packing in $\R^8$, and the Leech lattice packing is the
densest in $\R^{24}$.
\end{theorem}
M.~Viazovska, \emph{The sphere packing problem in dimension 8}, Ann.\ of Math.\ 185 (2017), 991--1015;
H.~Cohn, A.~Kumar, S.~D.~Miller, D.~Radchenko and M.~Viazovska, \emph{The sphere packing problem in dimension
24}, Ann.\ of Math.\ 185 (2017), 1017--1033. In both cases the magic function is built from (quasi)modular
forms, as sketched in Section~\ref{NT:sec:LP}.

\begin{theorem}[Universal optimality; Cohn--Kumar--Miller--Radchenko--Viazovska 2022]\label{NT:gal:universal}
Let $p(r)=g(r^2)$ with $g\colon(0,\infty)\to\R$ completely monotone. Among all point configurations in $\R^8$
(resp.\ $\R^{24}$) of density one point per unit volume, $E_8$ (resp.\ the Leech lattice) minimizes the
lower $p$-energy.
\end{theorem}
H.~Cohn, A.~Kumar, S.~D.~Miller, D.~Radchenko and M.~Viazovska, \emph{Universal optimality of the $E_8$ and
Leech lattices and interpolation formulas}, Ann.\ of Math.\ 196 (2022), no.~3, 983--1082. Sphere packing is
a limiting case (very steep potentials). This connects the lattice problem to the Boltzmann weights and
energy landscapes of Lectures~7--11: $E_8$ and Leech are ground states for a whole class of interactions,
and not only optimal hard spheres.

\begin{theorem}[Siegel 1945; Minkowski--Hlawka 1943]\label{NT:gal:siegel}
Theorem~\ref{NT:thm:siegel} and Corollary~\ref{NT:cor:MH} hold.
\end{theorem}
C.~L.~Siegel, \emph{A mean value theorem in geometry of numbers}, Ann.\ of Math.\ 46 (1945), 340--347;
E.~Hlawka, \emph{Zur Geometrie der Zahlen}, Math.\ Z.\ 49 (1943/44), 285--312. Minkowski had stated the bound
without a complete proof. The lattice analogue of the saturation bound of Lecture~1.

\begin{theorem}[Klartag 2025]\label{NT:gal:klartag}
There is a constant $c>0$ such that for all large $n$ there is a lattice packing in $\R^n$ of density at least
$c\,n^2\,2^{-n}$.
\end{theorem}
B.~Klartag, \emph{Lattice packing of spheres in high dimensions using a stochastically evolving ellipsoid},
Invent.\ Math.\ (2026), doi:10.1007/s00222-026-01412-w; arXiv:2504.05042 (2025). It is the first lower bound, for lattices or for general packings, that beats
$n\,2^{-n}$ by more than a logarithmic factor. The proof grows an ellipsoid that contains no nonzero point
of a lattice, by a random process. Compare the Minkowski--Hlawka bound
$2\zeta(n)2^{-n}$ above and the table in Lecture~3.

%======================================================================
\section*{Exercises}
\addcontentsline{toc}{section}{Exercises}

\begin{exercise}\label{NT:ex:reduce}
Reduce the form $26x^2+130xy+169y^2$ (the Gram matrix of the basis $(1,5)$, $(0,13)$ of the lattice in
Figure~\ref{NT:fig:fermat}). Check that you get $13(x^2+y^2)$, and read off a representation $13=a^2+b^2$.
Show that $h(-20)=2$ by listing the reduced forms of discriminant $-20$.
\end{exercise}

\begin{exercise}\label{NT:ex:hermite}
(Hermite's inequality.) Let $v_1$ be a shortest vector of $\Lat\subset\R^n$, and let $\Lat'$ be the
orthogonal projection of $\Lat$ onto $v_1^\perp$.
\begin{enumerate}[nosep]
\item Show that $\Lat'$ is a lattice with $\covol(\Lat')=\covol(\Lat)/\abs{v_1}$.
\item Show that every $w'\in\Lat'$ lifts to $w\in\Lat$ with $\abs{w}^2\le\abs{w'}^2+\frac14\abs{v_1}^2$, and
  deduce $\frac34\lambda_1(\Lat)^2\le\lambda_1(\Lat')^2$.
\item Conclude that $\gamma_n\le\big(\frac43\gamma_{n-1}\big)^{(n-1)/n}$, and hence $\gamma_n\le(4/3)^{(n-1)/2}$.
  Compare with the proof of Theorem~\ref{L1:thm:lattice2}.
\end{enumerate}
\end{exercise}

\begin{exercise}
Show that $N_\Lat(r)=V_nr^n/\covol(\Lat)+O(r^{n-1})$ for every lattice $\Lat\subset\R^n$, by Gauss's
argument with a fundamental parallelepiped. Then use it to show that the average of $r_2(m)$ for
$m\le X$ tends to $\pi$, and that the average of $r_4(m)$ over $m\le X$ is asymptotic to $\frac{\pi^2}{2}X$.
\end{exercise}

\begin{exercise}
Use Minkowski's theorem to show that every prime $p\equiv1$ or $3\pmod8$ can be written as $x^2+2y^2$.
(Use that $-2$ is a square mod $p$ for these primes, and the ellipse $x^2+2y^2<2p$, of area $\sqrt2\,\pi p>4p$.)
\end{exercise}

\begin{exercise}
Prove Corollary~\ref{NT:cor:thetatrans} for $\Lat=\Z$ directly, by applying Poisson summation to
$e^{-\pi tx^2}$, $t>0$. Deduce that $t^{1/4}\,\Theta_\Z(it)$ is invariant under $t\mapsto1/t$, and use this to
show that $\Theta_\Z(it)\sim t^{-1/2}$ as $t\to0^+$. Interpret this limit as a lattice point count.
\end{exercise}

\begin{exercise}\label{NT:ex:D16}
Show that $D_n^+=D_n\cup(D_n+\frac12\mathbf 1)$ is a lattice if and only if $n$ is even, that it is
unimodular when it is integral, and that it is even unimodular if and only if $8\mid n$. Count the vectors of
norm $2$ and $4$ in $D_{16}^+$ and compare with $E_8\oplus E_8$.
\end{exercise}

\begin{exercise}
Using $\Theta_{E_8}=E_4$, show that the number of vectors of norm $2m$ in $E_8$ is divisible by $240$ for every $m\ge1$.
Then verify by hand that $E_8$ has $2160$ vectors of norm $4$. (They are $\pm2e_i$, the vectors
$\pm e_i\pm e_j\pm e_k\pm e_l$, and the half-integral vectors with one coordinate $\pm\frac32$.)
\end{exercise}

\begin{exercise}
(The LP bound in dimension one.) Let $f(x)=\max(0,1-\abs x)$. Compute $\hat f(y)=\big(\frac{\sin\pi y}{\pi y}\big)^2$,
check that $f$ satisfies the hypotheses of Theorem~\ref{NT:thm:CE} with $r=1$, and conclude that the bound is
sharp in dimension $1$. Which lattice packings achieve it, and how does this show up in the equality analysis of
Section~\ref{NT:sec:LP}?
\end{exercise}

\begin{exercise}
Using Siegel's theorem, compute the average over $X_n$ of the number of nonzero lattice vectors in a ball
of volume $V$ centred at $0$. Deduce that some unimodular lattice has $\lambda_1\ge(2/V_n)^{1/n}$, which is
asymptotically $\sqrt{n/(2\pi e)}$. Compare with the upper bound $\lambda_1\le2V_n^{-1/n}$ coming from
Minkowski's theorem, and translate both into bounds on $\gamma_n$.
\end{exercise}

\section*{Further reading}
\addcontentsline{toc}{section}{Further reading}
\begin{itemize}[nosep]
\item J.-P.~Serre, \emph{A Course in Arithmetic}, Springer GTM 7, 1973. Chapter~V: quadratic forms over $\Z$
  (including $8\mid n$ for even unimodular lattices); Chapter~VII: modular forms and theta series.
\item J.~H.~Conway and N.~J.~A.~Sloane, \emph{Sphere Packings, Lattices and Groups}, 3rd ed., Springer, 1999.
  Ch.~2 (theta series), Ch.~16--18 (Niemeier lattices, Leech).
\item J.~W.~S.~Cassels, \emph{An Introduction to the Geometry of Numbers}, Springer, 1959. Blichfeldt,
  Minkowski, Siegel and Hlawka.
\item J.~Martinet, \emph{Perfect Lattices in Euclidean Spaces}, Springer, 2003. Hermite constants and Voronoi's theory.
\item J.~H.~Conway (with F.~Fung), \emph{The Sensual (Quadratic) Form}, MAA, 1997. Binary forms
  and the topograph, and the 15 theorem.
\item H.~Cohn, \emph{A conceptual breakthrough in sphere packing}, Notices AMS 64 (2017), 102--115. The LP
  bound and Viazovska's construction, for non-specialists.
\item M.~N.~Huxley, \emph{Area, Lattice Points, and Exponential Sums}, Oxford, 1996. For the circle problem.
\end{itemize}
% ===== end chapters/NT.tex =====
\endgroup
\begingroup\renewcommand{\thechapter}{C}\renewcommand{\theHchapter}{interludeC}
% ===== begin chapters/ERG.tex =====
% MAT670 -- Interlude C: the torus, Diophantine approximation and aperiodic order.
% Entirely new material (not in the handwritten notes). Labels prefixed ERG:.
% Placed after Interlude B (number theory). Numerical claims checked with python3.

\tikzset{
  ERGv/.style={circle,draw,fill=white,inner sep=0pt,minimum size=6pt},
  ERGw/.style={circle,draw,fill=addcol!25,inner sep=0pt,minimum size=6pt},
  ERGproj/.style={gray,densely dotted},
  ERGthick/.style={draw=black,line width=0.3pt,fill=addcol!35},
  ERGthin/.style={draw=black,line width=0.3pt,fill=flagcol!30},
}

\chapter[Interlude C: Torus and aperiodic order]{Interlude C: The torus, Diophantine approximation and aperiodic order}
\label{ERG:ch}

\begin{addedblock}
This interlude is new: it is not in the handwritten notes. It starts from the torus $\R^n/\Lat$ of
Lecture~4 and asks what a single straight line does on it. The answer leads through Diophantine
approximation, the points on a circle of Lecture~7, the Fibonacci words of Lecture~11, and finally
to quasicrystals, which are obtained by \emph{folding} the Dynkin diagrams of Lectures~4--6. To keep
it readable it is set in black; all of it should be read as added material.
\end{addedblock}

Several objects in these notes have quietly been tori. Lecture~4 attached to a lattice $\Lat\subset\R^n$
the torus $\R^n/\Lat$, whose volume is $\abs{\det\Lat}$. Interlude~A found the maximal torus of
$\mathrm{SU}(n)$ to be $\mathfrak t/2\pi A_{n-1}$ (Section~\ref{LIE:sec:torus}). And the circle
$\Sph^1=\R/\Z$ carrying the points of Lecture~7 is the one-dimensional torus. The simplest dynamical
question we can ask about a torus is:
\begin{center}
\emph{draw a straight line on $\R^2/\Z^2$. Where does it go?}
\end{center}
If the slope is rational, the line closes up. If the slope is irrational, the line never closes and
\emph{shreds} the torus: it comes arbitrarily close to every point. How fast it does so is the subject
of Diophantine approximation, and the answer singles out the golden ratio
\[
  \varphi=\tau=\frac{1+\sqrt5}2=1.618\ldots,\qquad \varphi^2=\varphi+1 .
\]
(Both letters are standard. We write $\varphi$ in the number-theoretic sections and $\tau$ in the
Coxeter group sections, following each literature; $\tau'=-1/\tau=\frac{1-\sqrt5}{2}$ is its Galois
conjugate. Here $\tau$ is \emph{not} the kissing number.) Cutting the line out of the torus in a
different way produces the Fibonacci chain, the simplest quasicrystal. Its higher-dimensional
relatives come from the lattices $A_4$, $D_6$ and $E_8$ of Lecture~6.

Throughout, $\{x\}=x-\lfloor x\rfloor$ is the fractional part, $\norm{x}_{\Z}=\min_{p\in\Z}\abs{x-p}$
is the distance to the nearest integer, and $e(x)=e^{2\pi ix}$.

%======================================================================
\section{Shredding the torus}
\label{ERG:sec:torus}

Let $\ell_\alpha$ be the image in $\T^2=\R^2/\Z^2$ of the line $y=\alpha x$. Each time the line crosses
the vertical circle $\{x=0\}$ (that is, at $x=k\in\Z$) its height is $k\alpha$ modulo $1$. So the
\emph{first return map} to this circle is the rotation
\[
  R_\alpha\colon\T^1\to\T^1,\qquad t\mapsto t+\alpha \pmod 1,
\]
and $\ell_\alpha$ is the suspension of $R_\alpha$: it is dense in $\T^2$ exactly when the orbit
$\{k\alpha\}$ is dense in $\T^1$.

\begin{figure}[ht]
\centering
\begin{tikzpicture}[scale=2.5]
  \begin{scope}
    \begin{scope}\clip (0,0) rectangle (1,1);
      \foreach \k in {0,...,4} {\pgfmathsetmacro{\y}{\k*0.4-floor(\k*0.4+0.0001)}
        \draw[addcol,thick] (0,\y)--++(1,0.4); \draw[addcol,thick] (0,\y-1)--++(1,0.4);}
    \end{scope}
    \draw[thick] (0,0) rectangle (1,1);
    \foreach \k in {0,...,4} {\pgfmathsetmacro{\y}{\k*0.4-floor(\k*0.4+0.0001)} \fill (0,\y) circle (0.025);}
    \node[below,font=\footnotesize] at (0.5,-0.05) {slope $\frac25$: a closed curve};
  \end{scope}
  \begin{scope}[shift={(1.8,0)}]
    \begin{scope}\clip (0,0) rectangle (1,1);
      \foreach \k in {0,...,29} {\pgfmathsetmacro{\y}{\k*0.41421356-floor(\k*0.41421356)}
        \draw[addcol] (0,\y)--++(1,0.41421356); \draw[addcol] (0,\y-1)--++(1,0.41421356);}
    \end{scope}
    \draw[thick] (0,0) rectangle (1,1);
    \foreach \k in {0,...,29} {\pgfmathsetmacro{\y}{\k*0.41421356-floor(\k*0.41421356)} \fill (0,\y) circle (0.018);}
    \node[below,font=\footnotesize] at (0.5,-0.05) {slope $\sqrt2-1$: $30$ windings};
  \end{scope}
  \foreach \s in {0,1.8} {
    \draw[-{Stealth[length=3pt]}] (\s+0.47,0)--(\s+0.53,0); \draw[-{Stealth[length=3pt]}] (\s+0.47,1)--(\s+0.53,1);
    \draw[-{Stealth[length=3pt]}] (\s,0.44)--(\s,0.5); \draw[-{Stealth[length=3pt]}] (\s,0.5)--(\s,0.56);
    \draw[-{Stealth[length=3pt]}] (\s+1,0.44)--(\s+1,0.5); \draw[-{Stealth[length=3pt]}] (\s+1,0.5)--(\s+1,0.56);}
\end{tikzpicture}
\caption{The line $y=\alpha x$ wrapped around the torus $\R^2/\Z^2$ (opposite edges identified as
marked). The dots on the left edge are the returns $\{k\alpha\}$ to the circle $x=0$. For rational
$\alpha=p/q$ the line closes after $q$ turns; for irrational $\alpha$ it shreds the torus.}
\label{ERG:fig:torus}
\end{figure}

If $\alpha=p/q$ in lowest terms, the line closes up after $q$ horizontal and $p$ vertical turns, a
closed geodesic of length $\sqrt{p^2+q^2}$. The irrational case is Kronecker's theorem, and the proof
is the pigeonhole principle.

\begin{proposition}[Dirichlet 1842; Kronecker 1884]\label{ERG:prop:dirichlet}
Let $\alpha\in\R$ and let $Q\ge1$ be an integer. There is an integer $q$ with $1\le q\le Q$ and
$\norm{q\alpha}_\Z<1/Q$. If $\alpha$ is irrational, the orbit $\{k\alpha\}_{k\ge0}$ is dense in
$[0,1)$, and so $\ell_\alpha$ is dense in $\T^2$.
\end{proposition}

\begin{proof}
The $Q+1$ numbers $\{k\alpha\}$, $0\le k\le Q$, lie in the $Q$ intervals $[i/Q,(i+1)/Q)$, so two of
them, say for $j<k$, lie in the same interval. Then $q=k-j$ satisfies $\norm{q\alpha}_\Z<1/Q$. If
$\alpha$ is irrational then $\delta=\norm{q\alpha}_\Z>0$, and the multiples of $q\alpha$ march around
the circle in steps of length $\delta<1/Q$ (all in the same direction), so every interval of length
$1/Q$ contains a point of the orbit. Since $Q$ was arbitrary, the orbit is dense.
\end{proof}

Density says nothing about \emph{how often} the orbit visits a given arc. The right notion is the one
behind the Boltzmann distribution of Lecture~7: time averages should equal space averages.

\begin{definition}
A sequence $(x_k)$ in $[0,1)$ is \emph{equidistributed} if for every interval $I\subseteq[0,1)$
\[
  \lim_{N\to\infty}\frac1N\#\{0\le k<N: x_k\in I\}=\abs{I}.
\]
\end{definition}

\begin{theorem}[Weyl's criterion; Weyl 1916]\label{ERG:thm:weylcrit}
$(x_k)$ is equidistributed if and only if
\[
  \frac1N\sum_{k<N}e(mx_k)\to0\qquad\text{for every integer } m\neq0.
\]
Equivalently, $\frac1N\sum_{k<N}f(x_k)\to\int_0^1f$ for every continuous $1$-periodic $f$.
\end{theorem}

\begin{proof}[Sketch]
If $(x_k)$ is equidistributed, then the averages of step functions converge, and continuous functions are
uniform limits of step functions. Conversely, the condition gives convergence for trigonometric
polynomials, which are uniformly dense in the continuous periodic functions (Fej\'er, or
Stone--Weierstrass). For an interval $I$, choose continuous $f_-\le\mathbf 1_I\le f_+$ with
$\int(f_+-f_-)<\varepsilon$ and squeeze.
\end{proof}

\begin{theorem}[Weyl 1916; also Bohl and Sierpi\'nski 1909--10]\label{ERG:thm:weyl}
If $\alpha$ is irrational, $(\{k\alpha\})_{k\ge0}$ is equidistributed.
\end{theorem}

\begin{proof}
For $m\neq0$ we have $m\alpha\notin\Z$, so $e(m\alpha)\neq1$ and the geometric series gives
\[
  \Big|\frac1N\sum_{k=0}^{N-1}e(mk\alpha)\Big|
  =\frac{\abs{e(mN\alpha)-1}}{N\abs{e(m\alpha)-1}}\le\frac{1}{N\abs{\sin\pi m\alpha}}\longrightarrow0 .
  \qedhere
\]
\end{proof}

The same Fourier argument gives a statement about measures rather than orbits.

\begin{proposition}[Unique ergodicity]\label{ERG:prop:ue}
If $\alpha$ is irrational, Lebesgue measure is the only Borel probability measure on $\T^1$ invariant
under $R_\alpha$.
\end{proposition}

\begin{proof}
Let $\mu$ be invariant and $\hat\mu(m)=\int e(-mt)\,d\mu(t)$. Invariance gives
$\hat\mu(m)=\int e(-m(t+\alpha))\,d\mu=e(-m\alpha)\hat\mu(m)$, so $\hat\mu(m)=0$ for $m\neq0$. A
measure on the circle is determined by its Fourier coefficients, and these are those of Lebesgue measure.
\end{proof}

For a general measure-preserving system there is no such uniformity, and the best one can say is
\emph{Birkhoff's ergodic theorem} (1931): if $T$ preserves a probability measure $\mu$ and
$f\in L^1(\mu)$, then $\frac1N\sum_{k<N}f(T^kx)$ converges for $\mu$-almost every $x$, and the limit is
$\int f\,d\mu$ when $T$ is \emph{ergodic} (every invariant set has measure $0$ or $1$). Unique
ergodicity upgrades ``almost every $x$'' to ``every $x$, uniformly''.

\subsection*{Higher dimensions, and the torus inside a Lie group}

Now let $\Lat\subset\R^n$ be a lattice and $\T_\Lat=\R^n/\Lat$ the torus of Lecture~4. Its
characters are $x\mapsto e(\ip{w}{x})$ with $w$ in the \emph{dual lattice}
$\Lat^*=\{w:\ip{w}{\lambda}\in\Z\ \forall\lambda\in\Lat\}$ (the same dual lattice that appears in the
Poisson summation formula of Interlude~B). Weyl's criterion holds on $\T_\Lat$ with these characters, and
the geometric series argument gives:

\begin{theorem}[Kronecker--Weyl]\label{ERG:thm:kroneckerweyl}
Let $v\in\R^n$.
\begin{enumerate}[nosep]
\item The sequence $kv\bmod\Lat$ ($k\in\N$) is equidistributed in $\T_\Lat$ if and only if
  $\ip{w}{v}\notin\Z$ for every $w\in\Lat^*\setminus\{0\}$.
\item The line $tv\bmod\Lat$ ($t\in\R$) is equidistributed (in particular dense) if and only if
  $\ip{w}{v}\neq0$ for every $w\in\Lat^*\setminus\{0\}$, i.e.\ $v$ lies in no rational hyperplane of
  $\Lat$.
\end{enumerate}
For $\Lat=\Z^n$, (1) says that $1,v_1,\dots,v_n$ are linearly independent over $\Q$.
\end{theorem}

If the condition fails for some $w$, the orbit stays in the closed subgroup $\{x:\ip{w}{x}\in\Z\}$,
a finite union of parallel subtori. So orbit closures
of linear flows on tori are always subtori. This is the abelian ancestor of Ratner's theorem in
Section~\ref{ERG:sec:lattices}.

\begin{remark}[Generic elements of a maximal torus]
In Interlude~A the maximal torus of a compact group is $T=\mathfrak t/\Lat_0$, with $\Lat_0$ the kernel
of $\exp$ (Section~\ref{LIE:sec:torus}). By Theorem~\ref{ERG:thm:kroneckerweyl}(1), an element
$t=\exp(H)$ generates a dense subgroup of $T$ if and only if $\chi(t)\neq1$ for every nontrivial
character $\chi$ of $T$. Such \emph{topological generators} form a dense set of full measure. They are
the standard tool for deducing from Theorem~\ref{LIE:thm:torus} (every element is conjugate into $T$)
that any two maximal tori are conjugate: conjugate a generator $t'$ of $T'$ into $T$, and its closure
$T'$ goes along with it.
\end{remark}

%======================================================================
\section{How close does the line come? Diophantine approximation}
\label{ERG:sec:dioph}

The line $\ell_\alpha$ passes within $\norm{q\alpha}_\Z$ of the lattice point $(q,p)$, measured
vertically. Dirichlet's bound $\norm{q\alpha}_\Z<1/q$ for infinitely many $q$ is a statement about
lattice points near a line, and it has a second proof from the geometry of numbers of Lecture~4.

\begin{proof}[Second proof of Dirichlet, via Minkowski]
The parallelogram $K_Q=\{(x,y):\abs{x}\le Q,\ \abs{x\alpha-y}\le1/Q\}$ is convex, symmetric and compact,
of area $2Q\cdot 2/Q=4$. By Minkowski's convex body theorem it contains a lattice point
$(q,p)\neq0$. If $q=0$ then $\abs{p}\le1/Q<1$ forces $p=0$; so $1\le\abs{q}\le Q$ and
$\abs{q\alpha-p}\le1/Q$.
\end{proof}

\subsection*{Continued fractions}

Write $\alpha=[a_0;a_1,a_2,\dots]=a_0+\cfrac{1}{a_1+\cfrac{1}{a_2+\dotsb}}$ with $a_i\in\Z$, $a_i\ge1$
for $i\ge1$, and let $p_n/q_n=[a_0;a_1,\dots,a_n]$ be the convergents, so
$q_{n+1}=a_{n+1}q_n+q_{n-1}$. The following facts are classical (see e.g.\ Khinchin's
\emph{Continued Fractions}, or Cassels' \emph{Introduction to Diophantine Approximation}, Ch.~I).
\begin{enumerate}[nosep,label=(\roman*)]
\item $\frac{1}{(a_{n+1}+2)}<q_n\,\norm{q_n\alpha}_\Z<\frac{1}{a_{n+1}}$ for $n\ge1$, and $q_n\alpha-p_n$
  alternates in sign.
\item (Best approximation, Lagrange) If $1\le q<q_{n+1}$, then $\norm{q\alpha}_\Z\ge\norm{q_n\alpha}_\Z$.
  So the record holders for $\norm{q\alpha}_\Z$ are exactly the $q_n$.
\item For $\varphi=[1;1,1,1,\dots]$ the convergents are $F_{n+1}/F_n$ (Fibonacci numbers,
  $F_1=F_2=1$), and $F_n\varphi-F_{n+1}=-(\tau')^{n}=(-1)^{n+1}\varphi^{-n}$.
\end{enumerate}
Fact (iii) is Binet's formula $\varphi^n=F_n\varphi+F_{n-1}$ applied to $\tau'$, which satisfies the
same quadratic equation. By (i), $\inf_q q\norm{q\alpha}_\Z>0$ if and only if the partial quotients $a_n$
are bounded; such $\alpha$ are called \emph{badly approximable}. The golden ratio has the smallest
possible partial quotients, and it turns out to be the worst approximable number of all.

\begin{theorem}[Hurwitz 1891]\label{ERG:thm:hurwitz}
For every irrational $\alpha$ there are infinitely many fractions $p/q$ with
$\abs{\alpha-p/q}<\frac1{\sqrt5\,q^2}$. The constant $\sqrt5$ cannot be increased: for
$\alpha=\varphi$ and any $A>\sqrt5$ only finitely many $p/q$ satisfy $\abs{\alpha-p/q}<\frac1{Aq^2}$.
\end{theorem}

\begin{proof}[Proof of the sharpness]
The existence part follows by showing that of any three consecutive convergents one satisfies the
inequality (Hardy--Wright, Thm.~195). For the optimality, note that
$(p-q\varphi)(p-q\tau')=p^2-pq-q^2$ is a nonzero integer for $(p,q)\neq(0,0)$, because $\varphi$ is
irrational. Hence
\[
  q^2\Big|\varphi-\frac pq\Big|=\frac{q\abs{p^2-pq-q^2}}{\abs{p-q\tau'}}\ \ge\ \frac{1}{\abs{p/q-\tau'}}
  \longrightarrow\frac1{\varphi-\tau'}=\frac1{\sqrt5}
\]
as $p/q\to\varphi$. So $\liminf q^2\abs{\varphi-p/q}\ge1/\sqrt5$, and the convergents
$F_{n+1}/F_n$ attain this value in the limit.
\end{proof}

The proof is really about the indefinite quadratic form $x^2-xy-y^2$, which takes no small nonzero
values on $\Z^2$. Markov turned this into a complete description of the top of the spectrum.

\begin{theorem}[Markov 1879--80]\label{ERG:thm:markov}
For a real indefinite binary quadratic form $f$ of discriminant $d>0$, put
$\mu(f)=\inf_{(x,y)\in\Z^2\setminus0}\abs{f(x,y)}/\sqrt d$. The values of $\mu(f)$ that exceed $\frac13$
are exactly
\[
  \frac{m}{\sqrt{9m^2-4}},\qquad m\in\{1,2,5,13,29,34,89,169,194,\dots\},
\]
where $m$ runs over the \emph{Markov numbers}, the integers occurring in positive solutions of
$x^2+y^2+z^2=3xyz$. The value for $m$ is attained only by forms equivalent to a multiple of an
explicit integral form $f_m$. For $m=1,2$ these are $x^2-xy-y^2$ ($\mu=1/\sqrt5$) and $x^2-2y^2$
($\mu=1/\sqrt8$).
\end{theorem}

The dual statement for numbers is that the \emph{Lagrange spectrum}
$\{\limsup_q(q\norm{q\alpha}_\Z)^{-1}\}$ begins $\sqrt5,\sqrt8,\sqrt{221}/5,\dots$ and accumulates
only at $3$. All Markov triples arise from $(1,1,1)$ by the moves $(x,y,z)\mapsto(x,y,3xy-z)$. Whether
each Markov number is the largest entry of only one triple (Frobenius 1913) is the \emph{Markov
uniqueness conjecture}, still open.

\subsection*{Badly approximable numbers are bounded orbits}

Here the torus meets the space of lattices $X_2=\SL_2(\R)/\SL_2(\Z)$ of Lecture~4. Put
\[
  \Lat_\alpha=\begin{pmatrix}1&0\\ \alpha&1\end{pmatrix}\Z^2=\{(q,\,q\alpha+p):p,q\in\Z\},
  \qquad g_t=\begin{pmatrix}e^{-t}&0\\0&e^{t}\end{pmatrix}.
\]
The flow $g_t$ squeezes the first coordinate and stretches the second. A lattice point with small
$q\abs{q\alpha+p}$ becomes, at the right time $t$, a short vector of $g_t\Lat_\alpha$.

\begin{proposition}[Dani correspondence, simplest case]\label{ERG:prop:dani}
$\alpha$ is badly approximable if and only if the forward orbit $\{g_t\Lat_\alpha:t\ge0\}$ is bounded
in $X_2$.
\end{proposition}

\begin{proof}
By Mahler's criterion (Lecture~4) a set of unimodular lattices is bounded exactly when the shortest
vector length $\lambda_1$ is bounded below on it. We use the sup norm, which changes $\lambda_1$ by at
most a factor $\sqrt2$. Suppose $\abs q\abs{q\alpha+p}\ge c$ for all $q\neq0$, where we may take $c\le1$.
A vector $v=(e^{-t}q,e^t(q\alpha+p))$ of $g_t\Lat_\alpha$ with $q=0$ has $\norm v_\infty=e^t\abs p\ge1$.
If $q\neq0$ then $\norm v_\infty^2\ge\abs q\abs{q\alpha+p}\ge c$. So $\lambda_1\ge\sqrt c$ along the
orbit. Conversely, suppose $\norm v_\infty\ge\varepsilon$ for all $t\ge0$ and $v\neq0$, where
$\varepsilon\le1$. Given $q\ge1$ and $p$, choose $t\ge0$ with $e^{-t}q=\varepsilon/2$. Then the first
coordinate of $v$ is less than $\varepsilon$, so $e^t\abs{q\alpha+p}\ge\varepsilon$, that is,
$q\abs{q\alpha+p}\ge\varepsilon^2/2$.
\end{proof}

This dictionary, due to Dani (1985) in all dimensions, turns questions about numbers into questions about
orbits on the space of lattices. We return to it in Section~\ref{ERG:sec:lattices}.

%======================================================================
\section{Three gaps}
\label{ERG:sec:threegap}

Return to the circle of Lecture~7 (Section~\ref{L07:subsec:circle}), which we take to have length $1$, and
place on it the first $N$ points of an irrational rotation,
\[
  x_k=\{k\alpha\},\qquad 0\le k\le N-1 .
\]
They cut the circle into $N$ arcs, a point $(g_1,\dots,g_N)$ of the open gap simplex of Lecture~7. One
would expect $N$ different gap lengths. There are at most three.

\begin{theorem}[Three gap theorem; conjectured by Steinhaus, proved by S\'os 1957/58, Sur\'anyi 1958, \'Swierczkowski 1959]\label{ERG:thm:threegap}
Let $\alpha$ be irrational and $N\ge2$. Let $u,v\in\{1,\dots,N-1\}$ be the indices of the points
closest to $0$ from the right and from the left, and put $a=\{u\alpha\}$, $b=1-\{v\alpha\}$. Then every
gap between neighbouring points has length $a$, $b$ or $a+b$. More precisely, the right neighbour of
$x_k$ is
\[
  x_{k+u}\ \ (\text{gap } a)\ \text{if } k<N-u,\qquad
  x_{k-v}\ \ (\text{gap } b)\ \text{if } k\ge v,\qquad
  x_{k+u-v}\ \ (\text{gap } a+b)\ \text{otherwise}.
\]
There are $N-u$ gaps of length $a$, $N-v$ of length $b$, and $u+v-N$ of length $a+b$.
\end{theorem}

\begin{proof}
All the indices named are in $[0,N-1]$, so each claimed neighbour is one of the points, at the stated
distance to the right of $x_k$ (because $x_{k+u}=x_k+a$ and $x_{k-v}=x_k+b$ modulo $1$). We must
show that no point $x_j$ lies strictly inside the arc. We use only the minimality of $a$ and $b$: no
$1\le m\le N-1$ has $\{m\alpha\}\in(0,a)$ or $\{m\alpha\}\in(1-b,1)$.

\emph{Gap $a$} ($k+u\le N-1$). If $x_j\in(x_k,x_k+a)$ with $j>k$, then $\{(j-k)\alpha\}\in(0,a)$,
contradicting minimality of $a$. If $j<k$, then $x_{k+u}=x_k+a$ lies in $(x_j,x_j+a)$, so
$m=k+u-j\in[1,N-1]$ has $\{m\alpha\}\in(0,a)$, again a contradiction.

\emph{Gap $b$} ($k\ge v$). If $x_j\in(x_k,x_k+b)$ with $j<k$, then $\{(j-k)\alpha\}\in(0,b)$, i.e.\
$m=k-j$ has $\{m\alpha\}\in(1-b,1)$. If $j>k$, then $x_j\in(x_{k-v}-b,x_{k-v})$, and $m=j-k+v$
(which lies in $[1,N-1]$) has $\{m\alpha\}\in(1-b,1)$. Both contradict minimality of $b$.

\emph{Gap $a+b$} ($N-u\le k\le v-1$). Suppose $x_j\in(x_k,x_k+a+b)$. If $j>k$, then
$\{(j-k)\alpha\}\in(0,a+b)$, and in fact in $(a,a+b)$ by minimality of $a$. (It cannot equal $a$,
since then $j=k+u\ge N$.) Then $m=k+u-j$ lies in $[1,N-1]$ and $\{-m\alpha\}=\{(j-k)\alpha\}-a\in(0,b)$, so
$\{m\alpha\}\in(1-b,1)$, which is impossible. If $j<k$, then $s=k-j\in[1,N-1]$ has
$1-\{s\alpha\}\in(b,a+b)$, and $m=v-s\in[1,N-1]$ has
$\{m\alpha\}=\{v\alpha\}-\{s\alpha\}=(1-\{s\alpha\})-b\in(0,a)$, which is impossible.

The three ranges of $k$ are disjoint: a $k$ in the first two would give $a=b$, hence
$(u+v)\alpha\in\Z$. Counting the $k$ in each range gives the multiplicities.
\end{proof}

\begin{figure}[ht]
\centering
\begin{tikzpicture}[scale=1.75]
  \draw[gray!40] (0,0) circle (1);
  % arcs: start x, gap, style
  \foreach \x/\g/\c in {0/0.0711/addcol, 0.0711/0.1716/black, 0.2426/0.0711/addcol, 0.3137/0.1005/flagcol,
     0.4142/0.0711/addcol, 0.4853/0.1716/black, 0.6569/0.0711/addcol, 0.7279/0.1005/flagcol,
     0.8284/0.0711/addcol, 0.8995/0.1005/flagcol}
    {\draw[\c,line width=1.8pt] ({90-360*\x}:1) arc[start angle={90-360*\x}, end angle={90-360*(\x+\g)}, radius=1];}
  \foreach \k/\x in {0/0, 5/0.0711, 3/0.2426, 8/0.3137, 1/0.4142, 6/0.4853, 4/0.6569, 9/0.7279, 2/0.8284, 7/0.8995}
    {\fill ({90-360*\x}:1) circle (0.035); \node[font=\scriptsize] at ({90-360*\x}:1.17) {$\k$};}
  \node[font=\small,align=left,anchor=west] at (1.45,0.35) {\textcolor{addcol}{$\blacksquare$} $a=\{5\alpha\}\approx0.0711$ \ ($\times5$)};
  \node[font=\small,align=left,anchor=west] at (1.45,0) {\textcolor{flagcol}{$\blacksquare$} $b=1-\{7\alpha\}\approx0.1005$ \ ($\times3$)};
  \node[font=\small,align=left,anchor=west] at (1.45,-0.35) {$\blacksquare$ $a+b\approx0.1716$ \ ($\times2$)};
\end{tikzpicture}
\caption{The points $\{k\alpha\}$, $0\le k\le9$, for $\alpha=\sqrt2-1$, drawn clockwise from the top.
Here $u=5$, $v=7$, $N=10$, and the multiplicities are $N-u=5$, $N-v=3$ and $u+v-N=2$.}
\label{ERG:fig:threegap}
\end{figure}

\begin{example}[Fibonacci: two gaps in ratio $\varphi$]\label{ERG:ex:fibgaps}
Take $\alpha=\varphi^{-1}=\varphi-1$ and $N=F_n$ with $n\ge4$. By facts (ii) and (iii) above, the
closest returns to $0$ among $1\le k<F_n$ occur at $k=F_{n-1}$ and $k=F_{n-2}$. (Since all partial
quotients equal $1$, there are no intermediate fractions, so the one-sided best approximations are also
convergents.) They are at distances
$\varphi^{-(n-1)}$ and $\varphi^{-(n-2)}$ on opposite sides (the convergents of $\varphi^{-1}$ are
$F_{k-1}/F_k$). So $\{u,v\}=\{F_{n-1},F_{n-2}\}$ and $u+v-N=0$: \emph{there are only two gap lengths,
and their ratio is $\varphi$}. There are $F_{n-2}$ short gaps and $F_{n-1}$ long ones, and indeed
$F_{n-2}\varphi^{-(n-1)}+F_{n-1}\varphi^{-(n-2)}=1$. Adding the next point splits one long gap into a
long and a short gap, in the ratio $\varphi:1$. This is the arrangement of leaves (``phyllotaxis'') that
a plant produces when it turns each new leaf by the golden angle $2\pi/\varphi^2\approx137.5^\circ$.
(Checked numerically for $N\le40$.)
\end{example}

\begin{proposition}[How well spread is an orbit?]\label{ERG:prop:spread}
Let $m_N(\alpha)$ be the smallest gap of $\{k\alpha\}_{0\le k<N}$, so that the injectivity radius of
Lecture~7 is $\rho=m_N/2$. The regular $N$-gon has $Nm_N=1$. For every irrational $\alpha$,
\[
  \liminf_{N\to\infty}N\,m_N(\alpha)=\liminf_{q\to\infty}q\,\norm{q\alpha}_\Z\le\frac1{\sqrt5},
\]
with equality for $\alpha=\varphi$. Moreover $\inf_NNm_N(\alpha)>0$ if and only if $\alpha$ is badly
approximable.
\end{proposition}

\begin{proof}
The smallest gap is the smallest circular distance between two of the points, so
$m_N=\min_{1\le k<N}\norm{k\alpha}_\Z$. Taking $N=q+1$ gives
$Nm_N\le(q+1)\norm{q\alpha}_\Z$. Conversely $m_N=\norm{k\alpha}_\Z$ for some $k<N$, so
$Nm_N\ge k\norm{k\alpha}_\Z$, and $k\to\infty$ as $N\to\infty$ because $m_N\to0$. This proves the
equality of the two $\liminf$s, and the same two inequalities give the last statement. The bound
$1/\sqrt5$ and the equality case are Hurwitz's Theorem~\ref{ERG:thm:hurwitz}.
\end{proof}

So for \emph{every} rotation, infinitely often the injectivity radius falls short of the optimal value
$1/(2N)$ of Lecture~7 by a factor of (almost) $\sqrt5$. The golden rotation does no worse than this
asymptotically. Rotations by numbers not equivalent to $\varphi$ do worse (by a factor $\sqrt8$ or
more), by Markov's Theorem~\ref{ERG:thm:markov}.

\begin{remark}[Gaps as a function on the space of lattices]
The vector of gaps lives in the gap simplex of Lecture~7. In Interlude~A the same simplex appeared as
the Weyl alcove $T/W$ of $\mathrm{SU}(N)$ (the eigenvalues of a unitary matrix are $N$ points on a
circle). Marklof and Str\"ombergsson (2017) observed that the gap lengths of $\{k\alpha\}_{k<N}$ are
values of a simple function on the space of lattices $X_2$, evaluated at a lattice obtained from
$\Lat_\alpha$ by the diagonal flow at time $t\approx\log N$. The three gap theorem then becomes an elementary
statement about two-dimensional lattices.
\end{remark}

%======================================================================
\section{Fibonacci points and discrepancy}
\label{ERG:sec:fib}

How uniform can $N$ points in the unit square be? Measure the error with the \emph{star discrepancy}
\[
  D^*(P)=\sup_{0\le x,y\le1}\Big|\frac{\#\big(P\cap[0,x)\times[0,y)\big)}{N}-xy\Big| .
\]
Random points have $D^*\approx N^{-1/2}$. The three gap structure suggests a better choice. The
\emph{Fibonacci lattice} with $N=F_m$ points is
\[
  P_m=\Big\{\Big(\frac{k}{F_m},\Big\{\frac{kF_{m-1}}{F_m}\Big\}\Big):0\le k<F_m\Big\},
\]
the points of the line of slope $F_{m-1}/F_m\approx\varphi^{-1}$ on the torus
(Figure~\ref{ERG:fig:torus}), sampled at the $F_m$ times it crosses $x=k/F_m$. It is a lattice in
$\T^2$, namely $\Z^2+\Z\,(1,F_{m-1})/F_m$.

\begin{figure}[ht]
\centering
\begin{tikzpicture}[scale=3.4]
  \draw (0,0) rectangle (1,1);
  \foreach \k in {0,...,54} {\pgfmathsetmacro{\y}{mod(34*\k,55)/55} \fill[addcol] ({\k/55},\y) circle (0.012);}
  \node[below,font=\small] at (0.5,0) {$F_{10}=55$ Fibonacci lattice points};
\end{tikzpicture}
\qquad
\begin{tikzpicture}[scale=1.55]
  \draw[gray!50] (0,0) circle (1);
  \foreach \k in {0,...,99} {
    \pgfmathsetmacro{\z}{1-(2*\k+1)/100}
    \pgfmathsetmacro{\r}{sqrt(1-\z*\z)}
    \pgfmathsetmacro{\th}{360*mod(\k*0.6180339887,1)}
    \pgfmathsetmacro{\c}{cos(\th)}
    \ifdim \c pt>0pt \fill[addcol] ({\r*sin(\th)},\z) circle (0.022); \fi}
  \node[below,font=\small] at (0,-1.05) {$100$ spherical Fibonacci points};
\end{tikzpicture}
\caption{Left: the Fibonacci lattice $P_{10}$. Right: the front half of the spherical Fibonacci point set
\eqref{ERG:eq:sphfib} with $N=100$, viewed from the side.}
\label{ERG:fig:fib}
\end{figure}

\begin{theorem}[Discrepancy in the square]\label{ERG:thm:discrepancy}\leavevmode
\begin{enumerate}[nosep]
\item (Schmidt 1972) There is $c>0$ such that every $N$-point set $P\subset[0,1]^2$ has
  $D^*(P)\ge c\,\frac{\log N}{N}$. (Roth 1954 had proved the weaker bound $c\sqrt{\log N}/N$. It is a bound for the
  $L^2$ average of the discrepancy, and in that form it is sharp, by Davenport 1956.)
\item The Fibonacci lattice has $D^*(P_m)=O\big(\frac{\log F_m}{F_m}\big)$. More generally, a
  two-dimensional lattice $\{(k/N,\{kg/N\})\}$ has $D^*=O(K\log N/N)$ if the partial quotients of
  $g/N$ are at most $K$ (see Niederreiter, \emph{Random Number Generation and Quasi-Monte Carlo
  Methods}, 1992, Ch.~5).
\end{enumerate}
\end{theorem}

So the Fibonacci lattice is optimal up to a constant, because all partial quotients of
$F_{m-1}/F_m=[0;1,1,\dots,1]$ equal $1$. The same mechanism appears in one dimension: for badly
approximable $\alpha$ the discrepancy of $\{k\alpha\}_{k<N}$ is $O(\log N/N)$ (see Kuipers--Niederreiter,
Ch.~2), while Schmidt showed that every infinite sequence in $[0,1)$ has discrepancy at least $c\log N/N$ for infinitely many $N$.

\paragraph{On the sphere.} Pull the Fibonacci lattice back to $\Sph^2$ by the area-preserving
cylindrical (Lambert) projection $(\theta,z)\mapsto(\sqrt{1-z^2}\cos\theta,\sqrt{1-z^2}\sin\theta,z)$.
Up to small variations this gives the \emph{spherical Fibonacci points}
\begin{equation}\label{ERG:eq:sphfib}
  z_k=1-\frac{2k+1}{N},\qquad \theta_k=2\pi k\varphi^{-1},\qquad 0\le k<N,
\end{equation}
a golden spiral from pole to pole. The relevant measure of uniformity is the \emph{spherical cap
discrepancy} $D_{\mathrm{cap}}$, the supremum over spherical caps $C$ of $\abs{\#(P\cap C)/N-\sigma(C)}$,
with $\sigma$ the normalized area. Beck (1984) showed that every $N$-point set has
$D_{\mathrm{cap}}\ge cN^{-3/4}$, and that some $N$-point sets have $D_{\mathrm{cap}}\le C N^{-3/4}\sqrt{\log N}$.
For the spherical Fibonacci lattice, Aistleitner, Brauchart and Dick (2012) proved
$D_{\mathrm{cap}}\le 44\sqrt{8/F_m}$, that is $O(N^{-1/2})$. Their numerical experiments suggest
that the true order is close to $N^{-3/4}$, but this is \emph{not} proved. Spherical Fibonacci points
are also a standard starting configuration for numerical searches in the Tammes and Thomson problems
of Lectures~2 and~11. This is an empirical observation; they are not optimal.

%======================================================================
\section{Sturmian words and the Fibonacci chain}
\label{ERG:sec:sturm}

Now \emph{code} the irrational line. Record, at each integer time $k$, whether the rotation
$t\mapsto t+\alpha$ wraps around:
\[
  s_k=\lfloor(k+1)\alpha+\rho\rfloor-\lfloor k\alpha+\rho\rfloor\in\{0,1\}\qquad(0<\alpha<1).
\]
Equivalently, $s_k=1$ if and only if $\{k\alpha+\rho\}\in[1-\alpha,1)$. Such a word is called a
\emph{mechanical} (or \emph{Sturmian}) word of slope $\alpha$. It is also the \emph{cutting sequence} of
the line $y=\alpha x+\rho$: the record of which horizontal grid lines it crosses between consecutive
vertical ones. The \emph{complexity} $p(n)$ of an infinite word is its number of distinct factors (blocks)
of length $n$.

\begin{theorem}[Morse--Hedlund 1938/1940; Coven--Hedlund 1973]\label{ERG:thm:morsehedlund}
An infinite word with $p(n)\le n$ for some $n$ is eventually periodic. The one-sided infinite words of
the smallest remaining complexity, $p(n)=n+1$ for all $n$, are exactly the mechanical words of
irrational slope. (For two-sided words this needs recurrence: $\cdots aaabaaa\cdots$ also has
$p(n)=n+1$.)
\end{theorem}

\begin{proof}[Sketch that a mechanical word has $p(n)=n+1$]
The factor $s_k\cdots s_{k+n-1}$ is determined by which of the arcs $[1-\alpha,1)-j\alpha$,
$0\le j<n$, contain $y=\{k\alpha+\rho\}$. As $y$ moves around the circle, the factor can change only when
$y$ crosses one of the endpoints of these arcs, which are the $n+1$ points $\{-j\alpha\}$, $0\le j\le n$.
These cut the circle into $n+1$ arcs, so $p(n)\le n+1$. The word is not eventually periodic (its frequency
of $1$'s is the irrational number $\alpha$), so $p(n)\ge n+1$ by the first half of the theorem
(Exercise~\ref{ERG:ex:mh}). Hence $p(n)=n+1$, and the $n+1$ arcs correspond bijectively to the factors.
\end{proof}

By equidistribution, the frequency of a factor is the length of its arc, and by the three gap theorem
the $n+1$ arcs have at most three different lengths. So \emph{the factors of length $n$ of a Sturmian
word occur with at most three different frequencies}.

\paragraph{The Fibonacci word.} Let $\sigma$ be the substitution $a\mapsto ab$, $b\mapsto a$. Then
$\sigma^n(a)$ is a prefix of $\sigma^{n+1}(a)$, and the limit is the \emph{Fibonacci word}
\[
  f=abaababaabaababaababaabaababaabaab\cdots,\qquad \abs{\sigma^n(a)}=F_{n+2}.
\]
The lengths obey $\abs{\sigma^{n+1}(a)}=\abs{\sigma^n(a)}+\abs{\sigma^{n-1}(a)}$, because
$\sigma^{n+1}(a)=\sigma^n(a)\,\sigma^{n-1}(a)$. The letter $b$ never appears twice in a row, since
every $b$ comes from an $a\mapsto ab$ and the next block begins with $a$. The frequency of $b$ is
$\lim F_n/F_{n+2}=\varphi^{-2}\approx0.382$. In fact $f$ is the mechanical word of slope
$\varphi^{-2}$ with $\rho=\varphi^{-2}$ and $b\leftrightarrow1$ (the word with $\rho=0$ is $af$): its $k$-th letter ($k\ge0$) is $b$ exactly when
$\lfloor(k+2)\varphi^{-2}\rfloor-\lfloor(k+1)\varphi^{-2}\rfloor=1$ (Lothaire,
\emph{Algebraic Combinatorics on Words}, Ch.~2; we checked the first $987$ letters numerically).

\begin{remark}[Back to the channel of Lecture~11]
With $b=1$ and $a=0$, the Fibonacci word is a binary word with no two adjacent $1$'s, which is exactly
the jamming condition for discs in the narrow channel of Section~\ref{L11:sec:channel}. It is a
jammed configuration that is neither periodic nor random. Its fraction of $1$'s is
$\varphi^{-2}\approx0.382$. This lies strictly between the typical value
$\frac12-\frac{\sqrt5}{10}\approx0.276$ of Exercise~\ref{L11:ex:typical} and the maximum $\frac12$. Its
length per contact is $2\varphi^{-2}+L\varphi^{-1}$. Every finite block of it occurs in it with bounded
gaps, so it is ``as ordered as an aperiodic jammed state can be''. The count $F_{\ell+2}$ of all such
words and the lengths $\abs{\sigma^n(a)}=F_{n+2}$ obey the same Fibonacci recursion.
\end{remark}

%======================================================================
\section{Cut and project}
\label{ERG:sec:cutproject}

Here is the Fibonacci chain as a geometric object. In $\R^2$ take the line $E_\parallel$ of slope
$\varphi^{-1}$ (direction $(\varphi,1)$), the perpendicular line $E_\perp$, and the strip
$S=E_\parallel+W$, where the \emph{window} $W$ is the projection of the unit square to $E_\perp$, shifted
generically. Select the lattice points of $\Z^2$ in the strip and project them to $E_\parallel$.

\begin{figure}[ht]
\centering
\begin{tikzpicture}[scale=0.78]
  \begin{scope}
    \clip (-0.6,-2.4) rectangle (10.4,5.6);
    \fill[addcol!10] (-0.6,{(-0.6-1.3)/1.618034}) -- (10.4,{(10.4-1.3)/1.618034}) -- (10.4,{(10.4+1.318034)/1.618034}) -- (-0.6,{(-0.6+1.318034)/1.618034}) -- cycle;
    \draw[addcol!60] (-0.6,{(-0.6-1.3)/1.618034}) -- (10.4,{(10.4-1.3)/1.618034});
    \draw[addcol!60] (-0.6,{(-0.6+1.318034)/1.618034}) -- (10.4,{(10.4+1.318034)/1.618034});
  \end{scope}
  \foreach \m in {0,...,8} \foreach \k in {0,...,5} \fill[gray!60] (\m,\k) circle (0.04);
  \draw[thick,addcol] (0.667,-2.188) -- (10.18,3.691) node[right,font=\small]{$E_\parallel$};
\draw[ERGproj] (0,0)--(1.163,-1.881); \fill[addcol] (1.163,-1.881) circle (0.06);
\draw[ERGproj] (1,0)--(1.886,-1.434); \fill[addcol] (1.886,-1.434) circle (0.06);
\draw[ERGproj] (1,1)--(2.334,-1.158); \fill[addcol] (2.334,-1.158) circle (0.06);
\draw[ERGproj] (2,1)--(3.057,-0.711); \fill[addcol] (3.057,-0.711) circle (0.06);
\draw[ERGproj] (2,2)--(3.504,-0.434); \fill[addcol] (3.504,-0.434) circle (0.06);
\draw[ERGproj] (3,2)--(4.228,0.013); \fill[addcol] (4.228,0.013) circle (0.06);
\draw[ERGproj] (4,2)--(4.952,0.460); \fill[addcol] (4.952,0.460) circle (0.06);
\draw[ERGproj] (4,3)--(5.399,0.737); \fill[addcol] (5.399,0.737) circle (0.06);
\draw[ERGproj] (5,3)--(6.122,1.184); \fill[addcol] (6.122,1.184) circle (0.06);
\draw[ERGproj] (6,3)--(6.846,1.631); \fill[addcol] (6.846,1.631) circle (0.06);
\draw[ERGproj] (6,4)--(7.293,1.907); \fill[addcol] (7.293,1.907) circle (0.06);
\draw[ERGproj] (7,4)--(8.017,2.355); \fill[addcol] (8.017,2.355) circle (0.06);
\draw[ERGproj] (7,5)--(8.464,2.631); \fill[addcol] (8.464,2.631) circle (0.06);
\draw[ERGproj] (8,5)--(9.188,3.078); \fill[addcol] (9.188,3.078) circle (0.06);
\foreach \x/\y/\l in {1.672/-1.896/L, 2.257/-1.534/S, 2.843/-1.172/L, 3.428/-0.811/S, 4.013/-0.449/L,
  4.737/-0.002/L, 5.322/0.360/S, 5.908/0.722/L, 6.631/1.169/L, 7.217/1.531/S, 7.802/1.893/L,
  8.388/2.255/S, 8.973/2.617/L} \node[font=\scriptsize] at (\x,\y) {$\l$};
\draw[thick] (0,0)--(1,0)--(1,1)--(2,1)--(2,2)--(4,2)--(4,3)--(6,3)--(6,4)--(7,4)--(7,5)--(8,5);
\foreach \p in {(0,0),(1,0),(1,1),(2,1),(2,2),(3,2),(4,2),(4,3),(5,3),(6,3),(6,4),(7,4),(7,5),(8,5)} \fill \p circle (0.07);
  \draw[-{Stealth}] (9.6,0.2) -- ++(-0.4,0.647) node[left,font=\small]{$E_\perp$};
  \node[font=\small,addcol,anchor=west] at (-0.5,5.3) {strip $S=E_\parallel+W$};
\end{tikzpicture}
\caption{The Fibonacci chain by cut and project. The lattice points of $\Z^2$ in the strip (black) form a
staircase. Projected orthogonally onto $E_\parallel$ they give long ($L$) and short ($S$) intervals,
of lengths in the ratio $\varphi:1$, in the order of the Fibonacci word.}
\label{ERG:fig:cutproject}
\end{figure}

\begin{proposition}\label{ERG:prop:fibchain}
Let $w(m,k)=\varphi k-m+s$ (this is $\sqrt{1+\varphi^2}$ times the $E_\perp$-coordinate, shifted by
$s$), and select the points with $w\in(-1,\varphi]$. Each selected point $x$ is followed in the strip by
exactly one of $x+(1,0)$ (if $w(x)>0$) or $x+(0,1)$ (if $w(x)\le0$). The steps project to intervals of
lengths $\varphi/\sqrt{1+\varphi^2}$ and $1/\sqrt{1+\varphi^2}$. In the normalized coordinate
$y=(w+1)/\varphi^2\in(0,1]$, every step acts as $y\mapsto y-\varphi^{-2}\pmod1$, and the step is short
exactly when $y\le\varphi^{-2}$. So the sequence of tiles is a mechanical word of slope $\varphi^{-2}$:
a Fibonacci word (for $s$ chosen suitably, exactly $f$).
\end{proposition}

\begin{proof}
If $w(x)>0$ then $w(x+(1,0))=w(x)-1\in(-1,\varphi-1]$ is in the window, while $w(x+(0,1))=w(x)+\varphi>\varphi$
is not. If $w(x)\le0$ the roles are reversed, since $w(x)-1\le-1$ and $w(x)+\varphi\in(\varphi-1,\varphi]$.
A step $(1,0)$ changes $y$ by $-1/\varphi^2$, and a step $(0,1)$ by $\varphi/\varphi^2=1-\varphi^{-2}$,
and $w>0$ means $y>\varphi^{-2}$. It remains to see that the staircase through one selected point
contains all of them. Every selected point lies on such a staircase (run the rule backwards as well), and
distinct staircases are disjoint. By Theorem~\ref{ERG:thm:weyl} a staircase has a fraction
$\varphi^{-2}$ of short steps, so it has $1/(\varphi^{-1}\cdot\frac{\varphi}{\sqrt{1+\varphi^2}}+
\varphi^{-2}\frac{1}{\sqrt{1+\varphi^2}})=\varphi^2/\sqrt{1+\varphi^2}$ points per unit length of
$E_\parallel$. This is already the density of all lattice points in the strip, which is (strip width)/(covolume) $=\varphi^2/\sqrt{1+\varphi^2}$. So there is only one staircase.
\end{proof}

This is the prototype of a \emph{model set}.

\begin{definition}
A \emph{cut and project scheme} is a lattice $\mathcal L\subset E_\parallel\times E_\perp$ such that the
projection $\pi_\parallel$ is injective on $\mathcal L$ and $\pi_\perp(\mathcal L)$ is dense in
$E_\perp$. For a \emph{window} $W\subset E_\perp$, relatively compact with nonempty interior, the
\emph{model set} is $\Lambda(W)=\{\pi_\parallel(x):x\in\mathcal L,\ \pi_\perp(x)\in W\}$. It is
\emph{regular} if $\partial W$ has measure zero.
\end{definition}

A model set is a Delone set (uniformly discrete and relatively dense), so it is a packing in the sense of
Lecture~1 and has a density. It is in general not periodic. It is a \emph{Meyer set}:
$\Lambda-\Lambda\subseteq\Lambda+F$ for a finite set $F$ (Meyer 1972). Its density is
$\vol(W)/\covol(\mathcal L)$, by Weyl equidistribution of $\pi_\perp(\mathcal L)$ in $W$. Most importantly
for physics, \emph{regular model sets have pure point diffraction}: the Fourier transform of the
autocorrelation of the Dirac comb $\sum_{x\in\Lambda}\delta_x$ is a sum of point masses, supported on the
projection of the dual lattice $\mathcal L^*$ (Hof 1995 for Euclidean internal spaces; Schlottmann 2000
in general). The proof is a Poisson summation formula (Interlude~B) for $\mathcal L$, cut down by the
window.

\paragraph{Penrose tilings and Shechtman's crystals.} De Bruijn (1981) showed that the rhombic Penrose
tilings (Penrose 1974) are dual to \emph{pentagrids}: superpositions of five families of parallel lines
at angles $2\pi j/5$. Equivalently, their vertex sets are obtained by cut and project from $\Z^5$, with
$E_\parallel$ the plane on which the cyclic permutation of coordinates acts as rotation by $2\pi/5$. Since
the vertices lie in finitely many layers along $(1,1,1,1,1)$, the minimal setting is the
four-dimensional root lattice $A_4$, with a window made of four pentagons in a two-dimensional internal
space (Baake--Kramer--Schlottmann--Zeidler 1990). Figure~\ref{ERG:fig:penrose} was drawn from a
pentagrid. In 1984 Shechtman, Blech, Gratias and Cahn reported a rapidly
solidified Al--Mn alloy whose electron diffraction pattern had sharp spots and icosahedral symmetry. Such
symmetry is impossible for a crystal: a rotation of order $5$ preserving a lattice in $\R^2$ or $\R^3$
would have integer trace, but its trace is $2\cos\frac{2\pi}5=\varphi^{-1}$, or $1+\varphi^{-1}$. The
discovery earned Shechtman the 2011 Nobel Prize in Chemistry. Model sets with icosahedral symmetry are
the standard mathematical models of such quasicrystals.

\begin{figure}[ht]
\centering
\begin{tikzpicture}[scale=0.52]
\begin{scope}
\clip (0,0) circle (4.9);
\filldraw[ERGthick] (-5.545,0.951)--(-4.545,0.951)--(-4.236,1.902)--(-5.236,1.902)--cycle;
\filldraw[ERGthick] (-3.427,-2.490)--(-2.427,-2.490)--(-2.118,-1.539)--(-3.118,-1.539)--cycle;
\filldraw[ERGthick] (-3.427,-0.588)--(-2.427,-0.588)--(-2.118,0.363)--(-3.118,0.363)--cycle;
\filldraw[ERGthick] (-3.427,2.490)--(-2.427,2.490)--(-2.118,3.441)--(-3.118,3.441)--cycle;
\filldraw[ERGthick] (-0.809,-5.568)--(0.191,-5.568)--(0.500,-4.617)--(-0.500,-4.617)--cycle;
\filldraw[ERGthick] (-1.618,-3.078)--(-0.618,-3.078)--(-0.309,-2.127)--(-1.309,-2.127)--cycle;
\filldraw[ERGthick] (-0.809,-0.588)--(0.191,-0.588)--(0.500,0.363)--(-0.500,0.363)--cycle;
\filldraw[ERGthick] (-1.618,1.902)--(-0.618,1.902)--(-0.309,2.853)--(-1.309,2.853)--cycle;
\filldraw[ERGthick] (-0.809,4.392)--(0.191,4.392)--(0.500,5.343)--(-0.500,5.343)--cycle;
\filldraw[ERGthick] (1.309,-4.029)--(2.309,-4.029)--(2.618,-3.078)--(1.618,-3.078)--cycle;
\filldraw[ERGthick] (1.309,-2.127)--(2.309,-2.127)--(2.618,-1.176)--(1.618,-1.176)--cycle;
\filldraw[ERGthick] (1.309,0.951)--(2.309,0.951)--(2.618,1.902)--(1.618,1.902)--cycle;
\filldraw[ERGthick] (2.118,3.441)--(3.118,3.441)--(3.427,4.392)--(2.427,4.392)--cycle;
\filldraw[ERGthick] (3.927,-2.127)--(4.927,-2.127)--(5.236,-1.176)--(4.236,-1.176)--cycle;
\filldraw[ERGthick] (3.927,0.951)--(4.927,0.951)--(5.236,1.902)--(4.236,1.902)--cycle;
\filldraw[ERGthin] (-2.618,-3.078)--(-1.618,-3.078)--(-2.427,-2.490)--(-3.427,-2.490)--cycle;
\filldraw[ERGthin] (-2.618,1.902)--(-1.618,1.902)--(-2.427,2.490)--(-3.427,2.490)--cycle;
\filldraw[ERGthin] (-0.809,-3.665)--(0.191,-3.665)--(-0.618,-3.078)--(-1.618,-3.078)--cycle;
\filldraw[ERGthin] (-0.500,0.363)--(0.500,0.363)--(-0.309,0.951)--(-1.309,0.951)--cycle;
\filldraw[ERGthin] (2.118,-4.617)--(3.118,-4.617)--(2.309,-4.029)--(1.309,-4.029)--cycle;
\filldraw[ERGthin] (2.118,0.363)--(3.118,0.363)--(2.309,0.951)--(1.309,0.951)--cycle;
\filldraw[ERGthin] (4.236,-1.176)--(5.236,-1.176)--(4.427,-0.588)--(3.427,-0.588)--cycle;
\filldraw[ERGthin] (-2.618,1.902)--(-1.618,1.902)--(-2.427,1.314)--(-3.427,1.314)--cycle;
\filldraw[ERGthin] (-2.618,-3.078)--(-1.618,-3.078)--(-2.427,-3.665)--(-3.427,-3.665)--cycle;
\filldraw[ERGthin] (-0.500,3.441)--(0.500,3.441)--(-0.309,2.853)--(-1.309,2.853)--cycle;
\filldraw[ERGthin] (-0.500,-1.539)--(0.500,-1.539)--(-0.309,-2.127)--(-1.309,-2.127)--cycle;
\filldraw[ERGthin] (2.118,3.441)--(3.118,3.441)--(2.309,2.853)--(1.309,2.853)--cycle;
\filldraw[ERGthin] (2.427,-0.588)--(3.427,-0.588)--(2.618,-1.176)--(1.618,-1.176)--cycle;
\filldraw[ERGthin] (4.236,-0.000)--(5.236,-0.000)--(4.427,-0.588)--(3.427,-0.588)--cycle;
\filldraw[ERGthick] (-3.427,4.392)--(-2.427,4.392)--(-2.118,3.441)--(-3.118,3.441)--cycle;
\filldraw[ERGthick] (-3.427,1.314)--(-2.427,1.314)--(-2.118,0.363)--(-3.118,0.363)--cycle;
\filldraw[ERGthick] (-3.427,-0.588)--(-2.427,-0.588)--(-2.118,-1.539)--(-3.118,-1.539)--cycle;
\filldraw[ERGthick] (-3.427,-3.665)--(-2.427,-3.665)--(-2.118,-4.617)--(-3.118,-4.617)--cycle;
\filldraw[ERGthick] (-0.809,4.392)--(0.191,4.392)--(0.500,3.441)--(-0.500,3.441)--cycle;
\filldraw[ERGthick] (-1.618,1.902)--(-0.618,1.902)--(-0.309,0.951)--(-1.309,0.951)--cycle;
\filldraw[ERGthick] (-0.809,-0.588)--(0.191,-0.588)--(0.500,-1.539)--(-0.500,-1.539)--cycle;
\filldraw[ERGthick] (-0.809,-3.665)--(0.191,-3.665)--(0.500,-4.617)--(-0.500,-4.617)--cycle;
\filldraw[ERGthick] (1.309,2.853)--(2.309,2.853)--(2.618,1.902)--(1.618,1.902)--cycle;
\filldraw[ERGthick] (2.118,0.363)--(3.118,0.363)--(3.427,-0.588)--(2.427,-0.588)--cycle;
\filldraw[ERGthick] (1.309,-2.127)--(2.309,-2.127)--(2.618,-3.078)--(1.618,-3.078)--cycle;
\filldraw[ERGthick] (3.927,2.853)--(4.927,2.853)--(5.236,1.902)--(4.236,1.902)--cycle;
\filldraw[ERGthick] (3.927,0.951)--(4.927,0.951)--(5.236,-0.000)--(4.236,-0.000)--cycle;
\filldraw[ERGthick] (3.118,-4.617)--(3.427,-3.665)--(2.618,-3.078)--(2.309,-4.029)--cycle;
\filldraw[ERGthick] (0.191,-3.665)--(0.500,-2.714)--(-0.309,-2.127)--(-0.618,-3.078)--cycle;
\filldraw[ERGthick] (-1.618,-3.078)--(-1.309,-2.127)--(-2.118,-1.539)--(-2.427,-2.490)--cycle;
\filldraw[ERGthick] (-4.545,-2.127)--(-4.236,-1.176)--(-5.045,-0.588)--(-5.354,-1.539)--cycle;
\filldraw[ERGthick] (3.927,-2.127)--(4.236,-1.176)--(3.427,-0.588)--(3.118,-1.539)--cycle;
\filldraw[ERGthick] (1.309,-2.127)--(1.618,-1.176)--(0.809,-0.588)--(0.500,-1.539)--cycle;
\filldraw[ERGthick] (-0.809,-0.588)--(-0.500,0.363)--(-1.309,0.951)--(-1.618,0.000)--cycle;
\filldraw[ERGthick] (-3.736,0.363)--(-3.427,1.314)--(-4.236,1.902)--(-4.545,0.951)--cycle;
\filldraw[ERGthick] (3.118,0.363)--(3.427,1.314)--(2.618,1.902)--(2.309,0.951)--cycle;
\filldraw[ERGthick] (1.000,1.902)--(1.309,2.853)--(0.500,3.441)--(0.191,2.490)--cycle;
\filldraw[ERGthick] (-1.618,1.902)--(-1.309,2.853)--(-2.118,3.441)--(-2.427,2.490)--cycle;
\filldraw[ERGthick] (3.927,2.853)--(4.236,3.804)--(3.427,4.392)--(3.118,3.441)--cycle;
\filldraw[ERGthick] (2.118,3.441)--(2.427,4.392)--(1.618,4.980)--(1.309,4.029)--cycle;
\filldraw[ERGthin] (1.000,-3.078)--(1.309,-2.127)--(0.500,-2.714)--(0.191,-3.665)--cycle;
\filldraw[ERGthin] (-3.736,-1.539)--(-3.427,-0.588)--(-4.236,-1.176)--(-4.545,-2.127)--cycle;
\filldraw[ERGthin] (3.118,-1.539)--(3.427,-0.588)--(2.618,-1.176)--(2.309,-2.127)--cycle;
\filldraw[ERGthin] (-1.618,0.000)--(-1.309,0.951)--(-2.118,0.363)--(-2.427,-0.588)--cycle;
\filldraw[ERGthin] (3.927,0.951)--(4.236,1.902)--(3.427,1.314)--(3.118,0.363)--cycle;
\filldraw[ERGthin] (0.191,2.490)--(0.500,3.441)--(-0.309,2.853)--(-0.618,1.902)--cycle;
\filldraw[ERGthin] (1.309,4.029)--(1.618,4.980)--(0.809,4.392)--(0.500,3.441)--cycle;
\filldraw[ERGthin] (-3.736,-1.539)--(-3.427,-0.588)--(-3.118,-1.539)--(-3.427,-2.490)--cycle;
\filldraw[ERGthin] (1.000,-3.078)--(1.309,-2.127)--(1.618,-3.078)--(1.309,-4.029)--cycle;
\filldraw[ERGthin] (-3.736,0.363)--(-3.427,1.314)--(-3.118,0.363)--(-3.427,-0.588)--cycle;
\filldraw[ERGthin] (0.191,-0.588)--(0.500,0.363)--(0.809,-0.588)--(0.500,-1.539)--cycle;
\filldraw[ERGthin] (-3.736,3.441)--(-3.427,4.392)--(-3.118,3.441)--(-3.427,2.490)--cycle;
\filldraw[ERGthin] (1.000,1.902)--(1.309,2.853)--(1.618,1.902)--(1.309,0.951)--cycle;
\filldraw[ERGthin] (0.191,4.392)--(0.500,5.343)--(0.809,4.392)--(0.500,3.441)--cycle;
\filldraw[ERGthick] (3.927,-2.127)--(3.118,-1.539)--(2.309,-2.127)--(3.118,-2.714)--cycle;
\filldraw[ERGthick] (2.118,-4.617)--(1.309,-4.029)--(0.500,-4.617)--(1.309,-5.204)--cycle;
\filldraw[ERGthick] (4.236,1.902)--(3.427,2.490)--(2.618,1.902)--(3.427,1.314)--cycle;
\filldraw[ERGthick] (2.118,0.363)--(1.309,0.951)--(0.500,0.363)--(1.309,-0.225)--cycle;
\filldraw[ERGthick] (1.309,-2.127)--(0.500,-1.539)--(-0.309,-2.127)--(0.500,-2.714)--cycle;
\filldraw[ERGthick] (-0.500,-4.617)--(-1.309,-4.029)--(-2.118,-4.617)--(-1.309,-5.204)--cycle;
\filldraw[ERGthick] (2.118,3.441)--(1.309,4.029)--(0.500,3.441)--(1.309,2.853)--cycle;
\filldraw[ERGthick] (1.000,1.902)--(0.191,2.490)--(-0.618,1.902)--(0.191,1.314)--cycle;
\filldraw[ERGthick] (-0.809,-0.588)--(-1.618,0.000)--(-2.427,-0.588)--(-1.618,-1.176)--cycle;
\filldraw[ERGthick] (-2.618,-3.078)--(-3.427,-2.490)--(-4.236,-3.078)--(-3.427,-3.665)--cycle;
\filldraw[ERGthick] (-0.809,4.392)--(-1.618,4.980)--(-2.427,4.392)--(-1.618,3.804)--cycle;
\filldraw[ERGthick] (-2.618,1.902)--(-3.427,2.490)--(-4.236,1.902)--(-3.427,1.314)--cycle;
\filldraw[ERGthick] (-3.427,-0.588)--(-4.236,0.000)--(-5.045,-0.588)--(-4.236,-1.176)--cycle;
\filldraw[ERGthin] (3.118,-2.714)--(2.309,-2.127)--(2.618,-3.078)--(3.427,-3.665)--cycle;
\filldraw[ERGthin] (3.927,2.853)--(3.118,3.441)--(3.427,2.490)--(4.236,1.902)--cycle;
\filldraw[ERGthin] (1.309,-0.225)--(0.500,0.363)--(0.809,-0.588)--(1.618,-1.176)--cycle;
\filldraw[ERGthin] (-0.809,-3.665)--(-1.618,-3.078)--(-1.309,-4.029)--(-0.500,-4.617)--cycle;
\filldraw[ERGthin] (0.191,1.314)--(-0.618,1.902)--(-0.309,0.951)--(0.500,0.363)--cycle;
\filldraw[ERGthin] (-1.618,-1.176)--(-2.427,-0.588)--(-2.118,-1.539)--(-1.309,-2.127)--cycle;
\filldraw[ERGthin] (-1.618,3.804)--(-2.427,4.392)--(-2.118,3.441)--(-1.309,2.853)--cycle;
\filldraw[ERGthin] (-3.736,0.363)--(-4.545,0.951)--(-4.236,0.000)--(-3.427,-0.588)--cycle;
\filldraw[ERGthick] (3.118,3.441)--(2.309,2.853)--(2.618,1.902)--(3.427,2.490)--cycle;
\filldraw[ERGthick] (3.927,0.951)--(3.118,0.363)--(3.427,-0.588)--(4.236,-0.000)--cycle;
\filldraw[ERGthick] (-0.809,4.392)--(-1.618,3.804)--(-1.309,2.853)--(-0.500,3.441)--cycle;
\filldraw[ERGthick] (1.000,1.902)--(0.191,1.314)--(0.500,0.363)--(1.309,0.951)--cycle;
\filldraw[ERGthick] (2.118,0.363)--(1.309,-0.225)--(1.618,-1.176)--(2.427,-0.588)--cycle;
\filldraw[ERGthick] (3.927,-2.127)--(3.118,-2.714)--(3.427,-3.665)--(4.236,-3.078)--cycle;
\filldraw[ERGthick] (-3.736,3.441)--(-4.545,2.853)--(-4.236,1.902)--(-3.427,2.490)--cycle;
\filldraw[ERGthick] (-1.618,1.902)--(-2.427,1.314)--(-2.118,0.363)--(-1.309,0.951)--cycle;
\filldraw[ERGthick] (-0.809,-0.588)--(-1.618,-1.176)--(-1.309,-2.127)--(-0.500,-1.539)--cycle;
\filldraw[ERGthick] (1.000,-3.078)--(0.191,-3.665)--(0.500,-4.617)--(1.309,-4.029)--cycle;
\filldraw[ERGthick] (-4.545,0.951)--(-5.354,0.363)--(-5.045,-0.588)--(-4.236,0.000)--cycle;
\filldraw[ERGthick] (-3.736,-1.539)--(-4.545,-2.127)--(-4.236,-3.078)--(-3.427,-2.490)--cycle;
\filldraw[ERGthick] (-1.618,-3.078)--(-2.427,-3.665)--(-2.118,-4.617)--(-1.309,-4.029)--cycle;
\end{scope}
\end{tikzpicture}
\caption{A patch of a rhombic Penrose tiling, computed from a generic pentagrid (de Bruijn's method).
Thick rhombi ($72^\circ$) blue, thin rhombi ($36^\circ$) orange.}
\label{ERG:fig:penrose}
\end{figure}

%======================================================================
\section{Folding Dynkin diagrams: quasicrystals from \texorpdfstring{$A_4$, $D_6$, $E_8$}{A4, D6, E8}}
\label{ERG:sec:fold}

Where do the icosahedral windows and lattices come from? Return to Lecture~4. Dropping the
crystallographic axiom~4) from Definition~\ref{L04a:def:rootsystem} allows more root systems. The
exercise after that definition asked for the decagon and the icosidodecahedron. Coxeter's classification
of finite reflection groups (1934) adds exactly three types to the list of Lecture~5: the dihedral groups
$I_2(m)$ ($m\neq2,3,4,6$), and two exceptional groups $H_3$ and $H_4$, with Coxeter graphs
\[
  H_2=I_2(5):\ \circ\overset{5}{\text{---}}\circ,\qquad
  H_3:\ \circ\text{---}\circ\overset{5}{\text{---}}\circ,\qquad
  H_4:\ \circ\text{---}\circ\text{---}\circ\overset{5}{\text{---}}\circ .
\]
These are the symmetry groups of the regular pentagon (order $10$), the icosahedron (order $120$) and the
$600$-cell (order $14400$). An edge labelled $5$ means an angle of $\pi-\pi/5$ between simple roots, so
$\cos\frac\pi5=\frac\tau2$ enters, and the Cartan numbers are irrational. These groups preserve no lattice
in their own dimension. But they do preserve a lattice of \emph{twice} the dimension, and that is where
the quasicrystals come from.

The construction is a \emph{folding} of a simply laced Dynkin diagram (Shcherbak 1988; Moody--Patera
1993). Pair up the $2r$ nodes of $\Gamma=A_4,D_6,E_8$ as $(\alpha_i,\beta_i)$, $1\le i\le r$ (so $r=2,3,4$), as in
Figure~\ref{ERG:fig:fold}. The rule is simple. Each plain edge $i$---$j$ of $H_r$ gives the two edges
$\alpha_i$---$\alpha_j$ and $\beta_i$---$\beta_j$. Each $5$-edge $i$---$j$ gives the three edges
$\alpha_i$---$\beta_j$, $\beta_i$---$\alpha_j$ and $\beta_i$---$\beta_j$. For $H_4$ this produces the
tree with arms of lengths $4,2,1$ from the node $\beta_3$, which is $E_8$. For $H_3$ it produces $D_6$, and
for $H_2$ the chain $\alpha_1$---$\beta_2$---$\beta_1$---$\alpha_2$, which is $A_4$.

\begin{figure}[ht]
\centering
\begin{tikzpicture}[scale=0.86]
  % E8 -> H4
  \begin{scope}
    \foreach \i in {1,...,4} {\node[ERGv] (a\i) at (\i-1,1) {}; \node[font=\scriptsize,above=2pt] at (\i-1,1) {$\alpha_{\i}$};
                              \node[ERGw] (b\i) at (\i-1,0) {}; \node[font=\scriptsize,below=2pt] at (\i-1,0) {$\beta_{\i}$};
                              \draw[gray,dashed] (a\i)--(b\i);}
    \draw (a1)--(a2)--(a3); \draw (b1)--(b2)--(b3)--(b4); \draw (a3)--(b4); \draw (a4)--(b3);
    \node[font=\small] at (1.5,-0.85) {$E_8$};
    \draw[-{Stealth},thick] (3.5,0.5)--(4.3,0.5);
    \foreach \i in {1,...,4} {\node[ERGv] (h\i) at (4.6+0.7*\i-0.7,0.5) {}; \node[font=\scriptsize,above=2pt] at (4.6+0.7*\i-0.7,0.5) {$a_{\i}$};}
    \draw (h1)--(h2)--(h3)--node[below,font=\scriptsize]{$5$}(h4);
    \node[font=\small] at (5.65,-0.85) {$H_4$};
  \end{scope}
  % D6 -> H3
  \begin{scope}[xshift=7.6cm]
    \foreach \i in {1,...,3} {\node[ERGv] (a\i) at (\i-1,1) {}; \node[font=\scriptsize,above=2pt] at (\i-1,1) {$\alpha_{\i}$};
                              \node[ERGw] (b\i) at (\i-1,0) {}; \node[font=\scriptsize,below=2pt] at (\i-1,0) {$\beta_{\i}$};
                              \draw[gray,dashed] (a\i)--(b\i);}
    \draw (a1)--(a2); \draw (b1)--(b2)--(b3); \draw (a2)--(b3); \draw (a3)--(b2);
    \node[font=\small] at (1,-0.85) {$D_6$};
    \draw[-{Stealth},thick] (2.5,0.5)--(3.2,0.5);
    \foreach \i in {1,...,3} {\node[ERGv] (h\i) at (3.5+0.7*\i-0.7,0.5) {}; \node[font=\scriptsize,above=2pt] at (3.5+0.7*\i-0.7,0.5) {$a_{\i}$};}
    \draw (h1)--(h2)--node[below,font=\scriptsize]{$5$}(h3);
    \node[font=\small] at (4.2,-0.85) {$H_3$};
  \end{scope}
  % A4 -> H2
  \begin{scope}[xshift=13.3cm]
    \foreach \i in {1,2} {\node[ERGv] (a\i) at (\i-1,1) {}; \node[font=\scriptsize,above=2pt] at (\i-1,1) {$\alpha_{\i}$};
                          \node[ERGw] (b\i) at (\i-1,0) {}; \node[font=\scriptsize,below=2pt] at (\i-1,0) {$\beta_{\i}$};
                          \draw[gray,dashed] (a\i)--(b\i);}
    \draw (b1)--(b2); \draw (a1)--(b2); \draw (a2)--(b1);
    \node[font=\small] at (0.5,-0.85) {$A_4$};
    \draw[-{Stealth},thick] (1.4,0.5)--(2.0,0.5);
    \node[ERGv] (h1) at (2.3,0.5) {}; \node[ERGv] (h2) at (3.0,0.5) {};
    \node[font=\scriptsize,above=2pt] at (2.3,0.5) {$a_1$}; \node[font=\scriptsize,above=2pt] at (3.0,0.5) {$a_2$};
    \draw (h1)--node[below,font=\scriptsize]{$5$}(h2);
    \node[font=\small] at (2.65,-0.85) {$H_2$};
  \end{scope}
\end{tikzpicture}
\caption{Folding $E_8\to H_4$, $D_6\to H_3$, $A_4\to H_2$. The dashed lines pair each node $\alpha_i$
with its partner $\beta_i$ (shaded). These are orthogonal, and under the projection
$\beta_i\mapsto\tau a_i$, $\alpha_i\mapsto a_i$. A plain edge of $H_r$ comes from two parallel edges, and a
$5$-edge from a crossing pair plus the bottom edge. The $E_8$ diagram here is the chain
$\alpha_1\alpha_2\alpha_3\beta_4\beta_3\beta_2\beta_1$ with $\alpha_4$ attached to $\beta_3$.}
\label{ERG:fig:fold}
\end{figure}

\begin{theorem}[Folding]\label{ERG:thm:fold}
Let $\Gamma\in\{A_4,D_6,E_8\}$ with simple roots labelled as in Figure~\ref{ERG:fig:fold}, $L$ its root
lattice in $V=\R^{2r}$, and define $M\colon L\to L$ on the simple roots by
\[
  M\alpha_i=\beta_i,\qquad M\beta_i=\alpha_i+\beta_i .
\]
\begin{enumerate}[nosep]
\item $M$ is self-adjoint and $M^2=M+1$. Hence $V=E_\tau\oplus E_{\tau'}$ is an orthogonal sum of the
  $\tau$- and $\tau'$-eigenspaces of $M$, each of dimension $r$, and $L$ is a module over
  $\Z[\tau]$ (with $\tau$ acting as $M$).
\item Let $P$ be the orthogonal projection onto $E_\tau$ and $a_i=P\alpha_i$. Then $P\beta_i=\tau a_i$, and
  $a_1,\dots,a_r$ are the simple roots of $H_r$, all of squared length $1-\frac1{\sqrt5}$.
\item The products $w_i=s_{\alpha_i}s_{\beta_i}$ preserve $E_\tau$ and act on it as the reflections
  $s_{a_i}$. They generate a subgroup of the Weyl group $W(\Gamma)$ isomorphic to $W(H_r)$.
\item $P$ is injective on $L$, and $P(L)=\Z[\tau]a_1+\dots+\Z[\tau]a_r$ is dense in $E_\tau$.
\end{enumerate}
\end{theorem}

\begin{proof}
(1) $M$ is self-adjoint if and only if $\ip{\beta_i}{\alpha_j}=\ip{\alpha_i}{\beta_j}$ and
$\ip{\beta_i}{\beta_j}=\ip{\alpha_i}{\alpha_j}+\ip{\alpha_i}{\beta_j}$ for all $i,j$. Both identities can
be read off the edge rule. Each side of the first is $-1$ exactly for a $5$-edge. For the second, take
$i=j$ ($2=2+0$), a plain edge ($-1=-1+0$), a $5$-edge ($-1=0-1$), and no edge ($0=0$). Clearly
$M^2\alpha_i=M\beta_i=\alpha_i+\beta_i=(M+1)\alpha_i$, and similarly for $\beta_i$. A self-adjoint
operator with $M^2=M+1$ has eigenvalues among $\tau,\tau'$ and orthogonal eigenspaces. Since $M$ has an
integer matrix, the two eigenvalues, being Galois conjugate, have equal multiplicity.

(2) $P=(M-\tau')/\sqrt5$ commutes with $M$, so $P\beta_i=PM\alpha_i=MP\alpha_i=\tau a_i$. Then
$\ip{a_i}{a_j}=\ip{P\alpha_i}{\alpha_j}=\big(\ip{\beta_i}{\alpha_j}-\tau'\ip{\alpha_i}{\alpha_j}\big)/\sqrt5$.
This gives $-2\tau'/\sqrt5=1-\frac1{\sqrt5}=:x$ for $i=j$, $\tau'/\sqrt5=-\frac x2$ for a plain edge, so
the angle is $2\pi/3$, and $-1/\sqrt5=-x\frac\tau2$ for a $5$-edge, so the angle is $\pi-\pi/5$. Non-adjacent
pairs give $0$.

(3) The roots $\alpha_i,\beta_i$ are orthogonal, so $w_i=-1$ on $U_i=\operatorname{span}(\alpha_i,\beta_i)$
and $w_i=1$ on $U_i^\perp$. The plane $U_i$ is $M$-invariant, hence so is $U_i^\perp$, so $w_i$ commutes
with $M$ and preserves $E_\tau$. On $E_\tau$ its $(-1)$-eigenspace is $U_i\cap E_\tau=\R a_i$. So
$w_i|_{E_\tau}=s_{a_i}$, and these generate $W(H_r)$. If some $w\in\langle w_i\rangle$ is trivial on
$E_\tau$, then $w-1$ is a rational matrix vanishing on $E_\tau$, hence also on its Galois conjugate
$E_{\tau'}$, so $w=1$.

(4) A nonzero $v\in L$ with $Pv=0$ would satisfy $Mv=\tau'v$, impossible for a rational vector $v$ and an
integer matrix $M$. The image is the $\Z[\tau]$-span of the $a_i$, which is dense because $\Z[\tau]$ is
dense in $\R$.
\end{proof}

A direct computation (we did it in python3) shows that $P$ maps the roots of $\Gamma$ onto
$\Delta(H_r)\cup\tau\Delta(H_r)$:
\[
  A_4:\ 20=10+10,\qquad D_6:\ 60=30+30,\qquad E_8:\ 240=120+120 .
\]
Here $\Delta(H_r)$ is the root system of $H_r$: the vertices of a decagon, of an icosidodecahedron, and
of a $600$-cell. The $120$ unit quaternions forming the $600$-cell are the \emph{icosians}; their
$\Z[\tau]$-span, the icosian ring, is a $\Z[\tau]$-order in the quaternions, and Theorem~\ref{ERG:thm:fold}(1) is a shadow of the fact that $E_8$ is the
icosian ring, with a suitably modified norm, viewed as a lattice over $\Z$ (Conway--Sloane, Ch.~8, \S2; Moody--Patera 1993). The planar
case is drawn in Figure~\ref{ERG:fig:a4}.

\begin{figure}[ht]
\centering
\begin{tikzpicture}[scale=1.25]
  \foreach \k in {0,...,9} {\draw[gray!40] (0,0)--({36*\k}:1.618);}
  \draw[addcol] (0:1) \foreach \k in {1,...,10} {-- ({36*\k}:1)};
  \draw[flagcol] (0:1.618) \foreach \k in {1,...,10} {-- ({36*\k}:1.618)};
  \foreach \k in {0,...,9} {\fill[addcol] ({36*\k}:1) circle (0.045); \fill[flagcol] ({36*\k}:1.618) circle (0.045);}
  \draw[-{Stealth},thick] (0,0)--(0:1) node[below right=-1pt,font=\scriptsize]{$a_1$};
  \draw[-{Stealth},thick] (0,0)--(144:1) node[below left=-1pt,font=\scriptsize]{$a_2$};
  \node[font=\scriptsize,flagcol] at (0:2.05) {$\tau a_1$};
  \node[font=\small,align=left,anchor=west] at (2.4,0.3) {$P(\alpha)$ for the $20$ roots of $A_4$:};
  \node[font=\small,align=left,anchor=west] at (2.4,-0.15) {\textcolor{addcol}{inner decagon} $=\Delta(H_2)$,};
  \node[font=\small,align=left,anchor=west] at (2.4,-0.6) {\textcolor{flagcol}{outer decagon} $=\tau\,\Delta(H_2)$.};
\end{tikzpicture}
\caption{The $20$ roots of $A_4$ projected to the plane $E_\tau$ (the Coxeter plane of $A_4$). The two
concentric decagons have radii in the ratio $\tau$.}
\label{ERG:fig:a4}
\end{figure}

\paragraph{Quasicrystals.} By Theorem~\ref{ERG:thm:fold}(4), $(L,E_\tau,E_{\tau'})$ is a cut and project
scheme whose symmetry group $W(H_r)$ acts on both factors, on $E_\tau$ as $H_r$ and on $E_{\tau'}$ as its
Galois twin. Choosing an $H_r$-invariant window in $E_{\tau'}$ gives model sets with decagonal ($A_4$),
icosahedral ($D_6$; face-centred icosahedral phases such as Al--Cu--Fe are modelled on it, while
primitive ones such as Shechtman's Al--Mn use $\Z^6$) or $H_4$ symmetry ($E_8$, the
four-dimensional quasicrystal of Elser--Sloane 1987). In this way the lattices of Lecture~6 \emph{are}
the hidden periodic structures behind the quasicrystals.

\paragraph{Affine extensions and viruses.} For a crystallographic diagram, adding the extra node of the
extended (affine) Dynkin diagram adds a translation, and the group becomes an infinite discrete affine
reflection group, as in Interlude~A's alcoves. For $H_r$, adding a translation produces a group whose
orbits are dense, because $\Z[\tau]$ is. Patera and Twarock (2002) constructed an affine extension of each of
$H_2,H_3,H_4$ (unique under their axioms), and used finitely many steps of it to generate finite
quasicrystal fragments.
Dechant, Boehm and Twarock (2013) showed that some of these extensions are exactly what one gets by
folding the \emph{affine} diagrams $\tilde A_4,\tilde D_6,\tilde E_8$ and projecting. Twarock's
motivation was biology. Many viral capsids are icosahedral. Caspar--Klug theory explains many of them,
but not, for instance, the papovaviruses, whose $72$ capsomers are all pentamers. Twarock (2004)
explained their architecture by tilings built from the icosahedral group. Keef and Twarock (2009) used
affine extensions of the icosahedral group to describe how material at different radii of a virus
(capsid, genome) is organised by one common icosahedral blueprint.

%======================================================================
\section{The space of lattices as a dynamical system}
\label{ERG:sec:lattices}

Section~\ref{ERG:sec:dioph} used the space $X_n=\SL_n(\R)/\SL_n(\Z)$ of unimodular lattices from
Lecture~4. It is not compact (Mahler's criterion characterizes its bounded sets), but it has finite
volume. Siegel's mean value theorem (Interlude~B) says that, for the normalized invariant probability
measure $\mu_n$ on $X_n$ and $f$ integrable on $\R^n$,
\[
  \int_{X_n}\sum_{v\in\Lat\setminus0}f(v)\,d\mu_n(\Lat)=\int_{\R^n}f .
\]
Take $f=\mathbf 1_B$ for a ball $B$ of radius $r$ and volume $V<2$. The average number of nonzero
lattice vectors in $B$ is $V<2$. Since nonzero lattice vectors come in pairs $\pm v$, \emph{some} lattice
has none. Its balls of radius $r/2$ form a packing of density $V2^{-n}$. This gives the Minkowski--Hlawka bound
$\dens\ge2^{1-n}$ for lattice packings, and counting only primitive vectors improves it to
$\zeta(n)2^{1-n}$. It is a non-constructive existence statement, like the saturation bound of
Lecture~1, but for lattices.

$\SL_n(\R)$ acts on $X_n$, and the dynamics of its subgroups encode arithmetic. The diagonal flow
$g_t$ of Proposition~\ref{ERG:prop:dani} is the geodesic flow on the modular surface. \emph{Unipotent}
subgroups behave much more rigidly.

\begin{theorem}[Margulis 1987, the Oppenheim conjecture]\label{ERG:thm:margulis}
Let $Q$ be a real nondegenerate indefinite quadratic form in $n\ge3$ variables that is not a multiple of a
form with rational coefficients. Then for every $\varepsilon>0$ there is $x\in\Z^n\setminus0$ with
$\abs{Q(x)}<\varepsilon$. (Dani and Margulis, \emph{Invent.\ Math.}~98 (1989), strengthened this to the
density of $Q(x)$, $x\in\Z^n$ primitive, in $\R$.)
\end{theorem}

Margulis's proof is a statement about orbits. It suffices to treat $n=3$. Write $Q=Q_0\circ g$ with
$Q_0$ a fixed form of signature $(2,1)$, and let $H=\SO(Q_0)^\circ\cong\SO(2,1)^\circ$, which is generated by
unipotent elements. Raghunathan observed that it suffices to show: \emph{every bounded $H$-orbit in $X_3$ is
compact}. Margulis proved this. A compact orbit $Hg\Z^3$ occurs only when $Q$ is a multiple of a rational
form. Otherwise the orbit is unbounded, so by Mahler's criterion some $hg\Z^3$ contains a very short vector
$hgx$, and then $\abs{Q(x)}=\abs{Q_0(hgx)}$ is small. Theorem~\ref{ERG:thm:margulis} fails for $n=2$, and the
reason is Section~\ref{ERG:sec:dioph}. For $Q(x,y)=(x-\sqrt2y)(x-\sqrt3y)$ one has
$\inf_{\Z^2\setminus0}\abs{Q}>0$ (Exercise~\ref{ERG:ex:oppenheim2}), because $\sqrt2$ and $\sqrt3$ are badly
approximable. Markov's Theorem~\ref{ERG:thm:markov} is a detailed study of exactly these binary forms.

\begin{theorem}[Ratner 1991]\label{ERG:thm:ratner}
Let $G$ be a Lie group, $\Gamma\subset G$ a lattice and $U\subseteq G$ a subgroup generated by
one-parameter unipotent subgroups. Then the closure of every orbit $Ux\Gamma$ in $G/\Gamma$ is
\emph{homogeneous}: it is an orbit $Fx\Gamma$ of a closed subgroup $F\supseteq U$, carrying a finite
$F$-invariant measure. Every ergodic $U$-invariant probability measure is of this algebraic form.
\end{theorem}

This is the non-abelian descendant of Theorem~\ref{ERG:thm:kroneckerweyl}, where orbit closures of linear
flows were subtori. For diagonal flows nothing like it holds in general (bounded orbits abound, as the
badly approximable numbers show). But for higher rank diagonal actions there is a rigidity conjecture of
Margulis, of which the Littlewood conjecture in the Gallery is a consequence.

%======================================================================
\section{Gallery of theorems}
\label{ERG:sec:gallery}

\begin{itemize}[leftmargin=1.2em,itemsep=0.45em]
\item \textbf{Weyl's polynomial theorem} (Weyl 1916, \emph{Math.\ Ann.}~77). If a real polynomial $p$
  has an irrational coefficient other than the constant term, then $\{p(k)\}$ is equidistributed. This is
  the first step beyond rotations; the proof uses Weyl differencing.
\item \textbf{Bounded remainder intervals} (Hecke 1922, Ostrowski; Kesten 1966, \emph{Acta Arith.}~12,
  193--212). For irrational $\alpha$, the discrepancy $\#\{k<N:\{k\alpha\}\in I\}-N\abs I$ stays bounded
  if and only if $\abs I\in\Z+\alpha\Z$. The window of the Fibonacci chain is such an interval, which is
  why its tile counts deviate from $N\varphi^{-1}$, $N\varphi^{-2}$ by $O(1)$.
\item \textbf{Khinchin's theorem} (Khinchin 1924, \emph{Math.\ Ann.}~92). Let $\psi\colon\N\to[0,\infty)$ be
  non-increasing. For almost every $\alpha$, $\norm{q\alpha}_\Z<\psi(q)$ has infinitely many solutions if
  $\sum\psi(q)=\infty$, and only finitely many if $\sum\psi(q)<\infty$. In particular almost every number
  is \emph{not} badly approximable, yet the badly approximable numbers have full Hausdorff dimension
  (Jarn\'ik 1928).
\item \textbf{The Duffin--Schaeffer conjecture} (Duffin--Schaeffer 1941; proved by Koukoulopoulos and
  Maynard, \emph{Ann.\ of Math.}~192 (2020), 251--307). For \emph{any} $\psi\colon\N\to[0,\infty)$, if
  $\sum_q\psi(q)\phi(q)/q=\infty$ (with $\phi$ Euler's function), then almost every $\alpha$ has infinitely
  many reduced fractions $a/q$ with $\abs{\alpha-a/q}\le\psi(q)/q$. This removes monotonicity from
  Khinchin's theorem. The proof is a delicate study of overlaps of the approximation sets, organised as a
  graph-theoretic compression argument.
\item \textbf{Littlewood's conjecture} (c.~1930, open). For all real $\alpha,\beta$,
  $\liminf_{q\to\infty}q\norm{q\alpha}_\Z\norm{q\beta}_\Z=0$. Einsiedler, Katok and Lindenstrauss
  (\emph{Ann.\ of Math.}~164 (2006), 513--560) classified the positive entropy invariant measures for the
  diagonal group on $\SL_k(\R)/\SL_k(\Z)$, and deduced that the set of exceptions $(\alpha,\beta)$ has
  Hausdorff dimension $0$. By the Dani correspondence, Littlewood concerns orbits of the $2$-parameter
  diagonal group on $X_3$.
\item \textbf{Three gaps and lattices} (Marklof--Str\"ombergsson, \emph{Amer.\ Math.\ Monthly}~124 (2017),
  741--745). A proof of Theorem~\ref{ERG:thm:threegap} in which the gaps are the values of a function on the
  space of $2$-dimensional lattices, as in the remark at the end of Section~\ref{ERG:sec:threegap}.
\item \textbf{Gaps of $\sqrt n$ mod 1} (Elkies--McMullen, \emph{Duke Math.\ J.}~123 (2004), 95--139). The
  points $\{\sqrt n\}$, $n\le N$, also cut the circle (Lecture~7). Their normalized gap distribution
  exists, is \emph{not} the exponential law of random points, and is given by an explicit piecewise analytic
  formula. The proof uses Ratner's theorem on a space of affine lattices.
\item \textbf{Optimal dispersion} (Breneis--Hinrichs, arXiv 2019; in \emph{Discrepancy Theory}, de Gruyter
  2020, 117--132). The largest empty axis-parallel box among $n\ge2$ points in the torus $\T^2$ has area at least $2/n$,
  and for $n$ a Fibonacci number the Fibonacci lattice attains exactly $2/n$. This is a packing problem
  for boxes in which the golden lattice is optimal.
\item \textbf{Spherical Fibonacci points} (Aistleitner--Brauchart--Dick, \emph{Discrete Comput.\ Geom.}~48
  (2012), 990--1024). The spherical cap discrepancy is $O(N^{-1/2})$, with an explicit constant, for the spherical
  Fibonacci lattice, and the same order for spherical digital nets. Compare Beck's lower bound $cN^{-3/4}$
  (\emph{Mathematika}~31, 1984).
\item \textbf{Pure point diffraction} (Hof, \emph{Comm.\ Math.\ Phys.}~169 (1995), 25--43; Schlottmann,
  in \emph{Directions in Mathematical Quasicrystals}, CRM Monogr.~13, AMS 2000, 143--159). Regular model
  sets have pure point diffraction. This is the precise sense in which a model set is ``crystalline
  without being periodic''.
\item \textbf{Gauss circle problem for Penrose tilings} (Haynes--Lutsko, arXiv:2512.21444, 2025). The number of
  vertices of a rhombic Penrose tiling (unit edges) in a disc of radius $R$ is
  $\pi C_PR^2+O(R^{2/3}(\log R)^{2/3})$, uniformly over all such tilings, with $C_P\approx1.231$. The proof
  combines the cut and project description with Poisson summation in $\Q(\zeta_5)$, a direct bridge to
  Interlude~B. A recent preprint.
\item \textbf{Undecidability of tiling} (Wang 1961; Berger, \emph{Mem.\ Amer.\ Math.\ Soc.}~66 (1966);
  Robinson, \emph{Invent.\ Math.}~12 (1971)). It is undecidable whether a finite set of Wang tiles (unit
  squares with coloured edges, matched edge to edge, no rotations) tiles the plane. The proof needs, and
  produces, an \emph{aperiodic} tile set: one that tiles, but only non-periodically. Berger used $20426$
  tiles, and Robinson used $6$ tiles (allowing rotations and reflections). The minimum is $11$ Wang tiles, on $4$ colours (Jeandel--Rao,
  \emph{Advances in Combinatorics} 2021).
\item \textbf{Penrose tilings} (Penrose 1974, \emph{Bull.\ Inst.\ Math.\ Appl.}~10; de Bruijn 1981,
  \emph{Indag.\ Math.}~43). Two rhombi with matching rules tile only aperiodically. Every finite patch
  recurs with bounded gaps, and the vertex sets are model sets (Section~\ref{ERG:sec:cutproject}).
\item \textbf{The einstein problem} (Smith--Myers--Kaplan--Goodman-Strauss, ``An aperiodic monotile'',
  \emph{Combinatorial Theory}~4(1) (2024); ``A chiral aperiodic monotile'', \emph{ibid.}~4(2) (2024)). The
  ``hat'', a polygon made of $8$ kites of the deltoidal trihexagonal tiling, tiles the plane but never
  periodically. Its tilings use the hat and its mirror image. The ``spectre'', a modification of a related
  tile, tiles only aperiodically and needs no mirror images. Earlier, Socolar and Taylor (\emph{J.\ Combin.\ Theory Ser.~A}~118 (2011), 2207--2231) had
  given an aperiodic hexagonal monotile whose matching rules are not enforced by shape alone.
\item \textbf{Keller's conjecture} (Keller 1930). In every tiling of $\R^n$ by unit cubes, some two cubes
  share a full facet. It is true for $n\le6$ (Perron 1940). It is false for $n\ge10$ (Lagarias--Shor,
  \emph{Bull.\ Amer.\ Math.\ Soc.}~27 (1992), 279--283) and for $n\ge8$ (Mackey, \emph{Discrete Comput.\
  Geom.}~28 (2002), 275--279). It is true for $n=7$ (Brakensiek--Heule--Mackey--Narv\'aez, IJCAR 2020;
  \emph{J.\ Automated Reasoning}~66 (2022), 277--300, by a SAT solver). By Haj\'os' 1942 theorem the
  \emph{lattice} case, Minkowski's conjecture, is true in all dimensions. By a reduction of Szab\'o (1986)
  one may assume that the tiling is periodic modulo $2\Z^n$, so that it is a configuration of cubes in the
  flat torus $\R^n/2\Z^n$.
\item \textbf{Periodic tiling conjecture} (Grünbaum--Shephard; Lagarias--Wang 1996). If a single finite set
  $F\subset\Z^d$ tiles $\Z^d$ by translations, does it also tile periodically? Yes for $d=2$ (Bhattacharya,
  \emph{Amer.\ J.\ Math.}~142 (2020), 255--266), hence tiling by one tile is decidable there. No in general: Greenfeld and
  Tao (\emph{Ann.\ of Math.}~200 (2024), 301--363) found a counterexample in sufficiently high dimension,
  encoding a ``Sudoku'' by a single tiling equation. So translational tilings by one tile can be forced to
  be aperiodic.
\item \textbf{Convex pentagons} (Reinhardt 1918 to Mann--McLoud-Mann--Von Derau, arXiv 2015,
  \emph{Geom.\ Dedicata}~194 (2018), 141--167). Fifteen families of
  convex pentagons are known to tile the plane, the fifteenth found by computer search in 2015. Rao (arXiv
  2017, computer-assisted) announced a proof that the list is complete; Hales independently reproduced
  the main part of the computation, but to our knowledge the proof has not appeared in a refereed journal. As a consequence, no convex polygon
  tiles the plane only non-periodically.
\item \textbf{Oppenheim and Ratner} (Margulis, \emph{C.\ R.\ Acad.\ Sci.\ Paris}~304 (1987); Ratner,
  \emph{Ann.\ of Math.}~134 (1991), 545--607, and \emph{Duke Math.\ J.}~63 (1991), 235--280). See
  Theorems~\ref{ERG:thm:margulis} and~\ref{ERG:thm:ratner}. They are the reason that counting lattice
  points in thin regions (Interlude~B) can be done by dynamics.
\item \textbf{Quasicrystals and icosians} (Elser--Sloane, \emph{J.\ Phys.\ A}~20 (1987), 6161; Moody--Patera,
  \emph{J.\ Phys.\ A}~26 (1993), 2829). These are the $E_8\to H_4$ projection and its algebraic form, via the
  icosian ring. Shcherbak (\emph{Russian Math.\ Surveys}~43 (1988), 149) found the foldings
  $A_4\to H_2$, $D_6\to H_3$, $E_8\to H_4$ in singularity theory. Affine versions: Patera--Twarock,
  \emph{J.\ Phys.\ A}~35 (2002), 1551--1574; Dechant--Boehm--Twarock, \emph{J.\ Phys.\ A}~45 (2012), 285202
  and \emph{J.\ Math.\ Phys.}~54 (2013), 093508. Viruses: Twarock, \emph{J.\ Theor.\ Biol.}~226 (2004),
  477--482; Keef--Twarock, \emph{J.\ Math.\ Biol.}~59 (2009), 287--313.
\item \textbf{Quasicrystals in the lab} (Shechtman--Blech--Gratias--Cahn, \emph{Phys.\ Rev.\ Lett.}~53 (1984),
  1951; Nobel Prize in Chemistry 2011). Sharp icosahedral diffraction in Al--Mn: the experimental
  counterpart of pure point diffraction of model sets.
\end{itemize}

%======================================================================
\section*{Exercises}
\addcontentsline{toc}{section}{Exercises}

\begin{exercise}
Let $\gcd(p,q)=1$. Show that the closed geodesic $\ell_{p/q}$ on $\R^2/\Z^2$ has length $\sqrt{p^2+q^2}$,
and that its complement is a single open annulus of width $1/\sqrt{p^2+q^2}$. (Its lifts to $\R^2$ are the
lines $qy-px\in\Z$.) So rational lines are spread evenly, but they leave a gap that shrinks only like
$1/\text{length}$.
\end{exercise}

\begin{exercise}
Prove Theorem~\ref{ERG:thm:kroneckerweyl} for $\Lat=\Z^n$ using Weyl's criterion on $\T^n$ (characters
$e(\ip{m}{x})$, $m\in\Z^n$). Then deduce the general case by a linear change of variables, and identify
the dual lattice.
\end{exercise}

\begin{exercise}
For $\alpha=\varphi^{-1}$ and $N=F_n$, list the points $\{k\alpha\}$ for $n=5$ ($N=5$) and $n=6$
($N=8$), and verify Example~\ref{ERG:ex:fibgaps} directly. What happens for $N=F_n+1$?
\end{exercise}

\begin{exercise}\label{ERG:ex:oppenheim2}
Show that $Q(x,y)=(x-\sqrt2y)(x-\sqrt3y)$ satisfies $\inf_{(x,y)\in\Z^2\setminus0}\abs{Q(x,y)}>0$.
\emph{Hint:} $\sqrt2=[1;2,2,\dots]$ and $\sqrt3=[1;1,2,1,2,\dots]$ have bounded partial quotients. If
$\abs{x-\sqrt2y}<\frac{\sqrt3-\sqrt2}{2}\abs y$, then the other factor is at least
$\frac{\sqrt3-\sqrt2}{2}\abs y$.
\end{exercise}

\begin{exercise}
Show that the Fibonacci word $f$ is \emph{balanced}: any two factors of the same length contain numbers of
$b$'s differing by at most one. Deduce again that $bb$ does not occur, and show that $aaa$ does not occur
either. Interpret both facts for the jammed channel of Lecture~11.
\end{exercise}

\begin{exercise}
Verify the edge rule of Section~\ref{ERG:sec:fold} for $H_3$: the resulting graph on the
six nodes $\alpha_i,\beta_i$ is $D_6$. Then compute the $2\times2$ Gram matrix
of $a_1,a_2$ in the case $A_4\to H_2$ and check that the angle is $4\pi/5$.
\end{exercise}

\begin{exercise}
Using the explicit $A_4$ roots $e_i-e_j\in\R^5$, find the plane on which the cyclic shift
$e_1\to e_2\to\dots\to e_5\to e_1$ acts as rotation by $2\pi/5$, project the $20$ roots to it, and recover
Figure~\ref{ERG:fig:a4}. (The cyclic shift is a Coxeter element of $W(A_4)=\Sigma_5$.)
\end{exercise}

\begin{exercise}
Using only Siegel's mean value theorem, show that for every $\varepsilon>0$ there is a lattice in $\R^n$
with packing density at least $(2-\varepsilon)2^{-n}$. Where exactly is the symmetry $v\leftrightarrow-v$
used?
\end{exercise}

\begin{exercise}\label{ERG:ex:mh}
Prove the first half of the Morse--Hedlund theorem. If $p(n)\le n$ for some $n$, then $p(m)=p(m+1)$ for
some $m\le n$. Deduce that every factor of length $m$ has a unique extension to the right, so the word is
eventually periodic.
\end{exercise}

\section*{Further reading}
\addcontentsline{toc}{section}{Further reading}
\begin{itemize}[nosep]
\item M.~Einsiedler and T.~Ward, \emph{Ergodic Theory with a view towards Number Theory}, Springer GTM 259,
  2011. Rotations, Weyl, the space of lattices, and the Dani correspondence.
\item J.~W.~S.~Cassels, \emph{An Introduction to Diophantine Approximation}, Cambridge, 1957; and
  M.~Aigner, \emph{Markov's Theorem and 100 Years of the Uniqueness Conjecture}, Springer, 2013.
\item L.~Kuipers and H.~Niederreiter, \emph{Uniform Distribution of Sequences}, Wiley, 1974. Weyl's
  criterion, discrepancy, $\{k\alpha\}$.
\item M.~Lothaire, \emph{Algebraic Combinatorics on Words}, Cambridge, 2002, Ch.~2 (Sturmian words).
\item M.~Baake and U.~Grimm, \emph{Aperiodic Order}, Vol.~1, Cambridge, 2013. Substitutions, cut and
  project, Penrose tilings, diffraction.
\item M.~Senechal, \emph{Quasicrystals and Geometry}, Cambridge, 1995; B.~Gr\"unbaum and G.~C.~Shephard,
  \emph{Tilings and Patterns}, Freeman, 1987.
\item J.~E.~Humphreys, \emph{Reflection Groups and Coxeter Groups}, Cambridge, 1990, for $H_3$, $H_4$ and
  Coxeter elements. Conway--Sloane, Ch.~8, for the icosians.
\item D.~W.~Morris, \emph{Ratner's Theorems on Unipotent Flows}, Chicago Lectures in Math., 2005.
\end{itemize}
% ===== end chapters/ERG.tex =====
\endgroup
\setcounter{chapter}{6}

\part{Configuration spaces and Morse theory}
% ===== begin chapters/L07.tex =====
% ======================================================================
% MAT670 -- Lectures 7 and 8 (source: MAT670_78, pp. 1--31)
% First-pass transcription: "clean + fill gaps".
% ======================================================================

\chapter{Lecture 7: A short introduction to configuration spaces}
\label{L07:ch:lecture7}

\section{Configuration spaces}
\label{L07:sec:confspaces}

A \emph{configuration space} is a topological space that records the possible
configurations of an object we wish to study. For example:
\begin{itemize}
  \item the $2^N$ possible results of $N$ coin flips
    \added{(a discrete space $\{0,1\}^N$ with $2^N$ points)};
  \item the tilings of a checkerboard by dominoes \added{(again a finite set)};
  \item the positions and orientations of a polytope in space
    \added{(for a polytope with no symmetries this is $\R^3\times\SO(3)$; in general
    one divides $\SO(3)$ by the rotational symmetry group of the polytope)};
  \item the configurations of a mechanical linkage;
  \item the positions (and velocities) of particles in a box.
\end{itemize}

This last example is the one most closely related to what we want to study, and to the
first few lectures: positions of points in space.

\begin{definition}[Configuration space of $N$ points]
\label{L07:def:conf}
Let $X$ be a topological space. The \emph{(classical) configuration space of $N$
distinct labelled points in $X$} is
\[
  C^N(X) = \Conf(N,X) = \prod_{i=1}^N X \setminus \Delta,
  \qquad
  \Delta = \{(x_i)_{i=1}^N : x_i = x_j \text{ for some } i\neq j\},
\]
\added{with the subspace topology from the product $X^N$.}
\orig{crossed out: ``statistical mechanics'', ``also study''}
\end{definition}

\begin{exercise}
\label{L07:ex:weyl}
Think about how this might be related to the Weyl group associated to the root system
$A_n$.
\added{(Hint: take $X=\R$ and $N=n+1$. The Weyl group of $A_n$ is the symmetric group
$\Sigma_{n+1}$, acting on $\R^{n+1}$ by permuting coordinates, and its reflecting
hyperplanes are exactly the hyperplanes $x_i=x_j$ that make up $\Delta$. This is taken up
again in Section~\ref{L07:sec:weyl}.)}
\end{exercise}

There are also \emph{unlabelled} versions, $\Conf(N,X)/\Sigma_N$, where the symmetric
group $\Sigma_N$ acts by permuting the indices. When $X$ is endowed with a metric we can
also consider \emph{reduced} versions, $\Conf(N,X)/\operatorname{Isom}(X)$.
\added{(One can of course do both: $\Conf(N,X)/(\Sigma_N\times\operatorname{Isom}(X))$.)}
So the unlabelled version picks a single representative class for each permutation of the
indices, and the reduced version divides out by the symmetries of the ambient space.

\begin{example}[$\Conf(N,\R)$]
\label{L07:ex:confR}
We have $\Conf(N,\R)=\R^N\setminus\Delta$.
\begin{itemize}
  \item \emph{Unlabelled:} each point of $\Conf(N,\R)/\Sigma_N$ has exactly one
    representative with $x_1<x_2<x_3<\dots<x_N$, so we can choose to work in this region
    of the space.
    \added{$\R^N\setminus\Delta$ has $N!$ connected components, one for each ordering of
    the coordinates; each is an open convex cone, hence contractible, and $\Sigma_N$
    permutes them simply transitively.}
  \item \emph{Reduced:} we can shift one of the $x_i$ to $0$, for example the smallest.
    Together with the ordering this gives $\{0=x_1<x_2<\dots<x_N\}$.
    \added{(Modulo translations, the ordered region becomes the open cone
    $\{0<x_2<\dots<x_N\}\subset\R^{N-1}$; if reflections $x\mapsto -x$ are divided out as well, one may in addition
    fix the order of the two ``outer'' gaps.)}
\end{itemize}
\end{example}

\begin{center}
\begin{tikzpicture}[scale=0.9]
  % N = 2
  \fill[blue!8] (0,0) -- (2.4,2.4) -- (0,2.4) -- cycle;
  \draw[->] (-1.2,0) -- (2.6,0) node[right] {$x_1$};
  \draw[->] (0,-1.2) -- (0,2.6) node[above] {$x_2$};
  \draw[dashed] (-1.2,-1.2) -- (2.4,2.4) node[above right] {$\Delta$};
  \draw[very thick,addcol] (0,0.05) -- (0,2.4);
  \fill[white] (0,0) circle (2pt); \draw (0,0) circle (2pt);
  \node[align=center] at (0.7,1.75) {\small $x_1<x_2$};
  \node[addcol,left] at (0,1.2) {\small $x_1=0$};
\end{tikzpicture}
\end{center}
\noindent\textit{Figure:} $\Conf(2,\R)=\R^2\setminus\Delta$. The shaded half-plane
$\{x_1<x_2\}$ is a copy of the unlabelled space; the thick ray $\{0=x_1<x_2\}$ is the
reduced unlabelled space.
\orig{p.\,3 also sketches $N=3$, labelled ``$x_1=0$'', ``$x_1<x_2$'', ``$0<x_2<x_3$''}

\begin{example}[$\Conf(N,\Sph^2)$]
\label{L07:ex:confS2}
Consider labelled distinct points on a sphere. The reduced space is the space of $N$
labelled points on $\Sph^2$ modulo \emph{ambient} rotations.
For $N=4$: two configurations that differ by a rotation of the sphere are the same point of
the reduced space; a configuration and its mirror image are in general \emph{not} related by
a rotation, but they may become identified in the unlabelled space after relabelling
(e.g.\ swapping the labels $2$ and $3$).
\orig{p.\,4: sketches of 4 labelled points on spheres}
\added{(If one divides by all of $\operatorname{Isom}(\Sph^2)=\Orth(3)$ rather than by
$\SO(3)$, mirror images are identified as well.)}
\end{example}

\section{Braid groups}
\label{L07:sec:braid}

These spaces are related to the Artin braid groups, which we may define as follows.

\begin{definition}[Braid groups]
\label{L07:def:braid}
The \emph{braid group} and the \emph{pure braid group} of $X$ (on $N$ strands) are
\[
  B_N(X) = \pi_1\big(\Conf(N,X)/\Sigma_N\big),
  \qquad
  P_N(X) = \pi_1\big(\Conf(N,X)\big).
\]
\added{For $X=\R^2$ these are Artin's braid group $B_N$ and pure braid group $P_N$.
Since $\Conf(N,X)\to\Conf(N,X)/\Sigma_N$ is a covering map with deck group $\Sigma_N$, for
connected $\Conf(N,X)$ there is a short exact sequence
$1\to P_N(X)\to B_N(X)\to\Sigma_N\to 1$.}
\end{definition}

To see why these are called braid groups, consider what a path is in $\Conf(N,\R^2)$.
Fix a base point $x_0=(x_0^1,\dots,x_0^N)$. A loop $\gamma\colon[0,1]\to\Conf(N,\R^2)$
consists of $N$ paths $\gamma_1,\dots,\gamma_N$ in the plane which never collide at any
time $t$. Drawing the graphs $\{(\gamma_i(t),t)\}\subset\R^2\times[0,1]$ gives $N$
disjoint strands running from the plane $t=0$ to the plane $t=1$: a braid.
\added{In the unlabelled space a loop may end at a permutation of the starting labels,
which is why $B_N$ surjects onto $\Sigma_N$; homotopic loops give braids that are isotopic
with endpoints fixed.}

\begin{center}
\begin{tikzpicture}[scale=0.8,yscale=1.1]
  % top and bottom planes
  \draw (-0.6,0.15) -- (2.6,0.15) -- (3.0,0.65) -- (-0.2,0.65) -- cycle;
  \draw (-0.6,-3.15) -- (2.6,-3.15) -- (3.0,-2.65) -- (-0.2,-2.65) -- cycle;
  \node[left] at (-0.6,0.4) {\small $x_0$};
  \node[left] at (-0.6,-2.9) {\small $x_0$};
  % level 0: sigma_1 (strand from pos 0 over)
  \draw[thick] (1,0) .. controls (1,-0.5) and (0,-0.5) .. (0,-1);
  \draw[white,line width=5pt] (0,0) .. controls (0,-0.5) and (1,-0.5) .. (1,-1);
  \draw[thick] (0,0) .. controls (0,-0.5) and (1,-0.5) .. (1,-1);
  \draw[thick] (2,0) -- (2,-1);
  % level 1: sigma_2^{-1}
  \draw[thick] (1,-1) .. controls (1,-1.5) and (2,-1.5) .. (2,-2);
  \draw[white,line width=5pt] (2,-1) .. controls (2,-1.5) and (1,-1.5) .. (1,-2);
  \draw[thick] (2,-1) .. controls (2,-1.5) and (1,-1.5) .. (1,-2);
  \draw[thick] (0,-1) -- (0,-2);
  % level 2: sigma_1
  \draw[thick] (1,-2) .. controls (1,-2.5) and (0,-2.5) .. (0,-3);
  \draw[white,line width=5pt] (0,-2) .. controls (0,-2.5) and (1,-2.5) .. (1,-3);
  \draw[thick] (0,-2) .. controls (0,-2.5) and (1,-2.5) .. (1,-3);
  \draw[thick] (2,-2) -- (2,-3);
  \foreach \x/\l in {0/1,1/2,2/3} {
    \fill (\x,0) circle (1.6pt); \node[above] at (\x,0.05) {\scriptsize $\l$};
    \fill (\x,-3) circle (1.6pt); \node[below] at (\x,-3.05) {\scriptsize $\l$};
  }
  \draw[->] (3.6,0) -- (3.6,-3) node[midway,right] {\small $t$};
\end{tikzpicture}
\end{center}
\noindent\textit{Figure:} a loop in $\Conf(3,\R^2)/\Sigma_3$ drawn as a braid on three
strands. \orig{p.\,5: second sketch (strand graphs) illegible, omitted}

\begin{remark}
\added{For $X=\R$ the components of $\Conf(N,\R)$ are contractible
(Example~\ref{L07:ex:confR}), so all braid groups of $\R$ are trivial. For $X=\R^n$ with
$n\geq 3$ the diagonal has codimension $n\geq3$, so $\Conf(N,\R^n)$ is simply connected:
$P_N(\R^n)=1$ and $B_N(\R^n)\cong\Sigma_N$. The plane is the interesting case.}
\end{remark}

\section{Functions on configuration spaces}
\label{L07:sec:functions}

To configurations we would like to associate real-valued functions, based on what they
model. For example:
\begin{itemize}
  \item an electrical potential, or more generally a potential given as a sum of pair
    interactions \added{$E(x_1,\dots,x_N)=\sum_{i<j}\phi(x_i,x_j)$};
  \item a hard shell (hard sphere) potential
    \added{($\phi=+\infty$ if two spheres overlap, $0$ otherwise)};
\end{itemize}
which can describe different energy levels and can also constrain the system.

We may ask about
\begin{itemize}
  \item minima and maxima (local and global);
  \item critical points;
  \item generic configurations.
\end{itemize}
This breaks down into a couple of categories:
\begin{itemize}
  \item small $N$: by hand;
  \item large $N$: by statistical mechanics, or by averaging.
\end{itemize}
The classic example for large $N$ is the expected value of $N$ coin flips with values in
$\{0,1\}$, which leads to the central limit theorem.

\subsection{The Boltzmann distribution}
\label{L07:subsec:boltzmann}

So, for example, consider $N$ particles, each in one of $K$ states, with energy $E_i$
associated to the state $i\in\{1,\dots,K\}$. Let $N_i$ be the number of particles in state
$i$. The system is closed, so the total energy is constant, and
\[
  E=\sum_{i=1}^K E_i N_i, \qquad N=\sum_{i=1}^K N_i .
\]
The number of ways to place the $N$ particles according to the partition
$\{N_1,\dots,N_K\}$ is the multinomial coefficient
\[
  W=\frac{N!}{N_1!\cdots N_K!}.
\]
For large $N$ the optimal partition dominates, so we want to determine the partition that
maximises $W$ \added{subject to the two constraints above}. One can derive the
\emph{Boltzmann measure}
\[
  \frac{N_i}{N}=\frac{e^{-\beta E_i}}{\sum_{j=1}^K e^{-\beta E_j}},
  \qquad\text{i.e.}\qquad
  \Pr(\text{state } i)=\frac{e^{-\beta E_i}}{Z},
\]
where $Z=\sum_{j=1}^K e^{-\beta E_j}$ is the \emph{partition function} and
$\beta=1/(k_B T)$, with $k_B$ \added{Boltzmann's} constant and $T$ the temperature.
\orig{derivation deferred to ``next time'' in the notes}

\begin{addedblock}
Sketch of the derivation. By Stirling's formula,
$\ln W\approx N\ln N-\sum_i N_i\ln N_i$. Maximising this subject to $\sum_i N_i=N$ and
$\sum_i E_iN_i=E$ with Lagrange multipliers $\alpha,\beta$ gives
$-\ln N_i-1-\alpha-\beta E_i=0$, i.e.\ $N_i\propto e^{-\beta E_i}$; normalising gives the
formula above. The multiplier $\beta$ is fixed by the energy constraint, and physics
identifies it with $1/(k_BT)$.
\end{addedblock}

\begin{remark}
Note:
\begin{itemize}
  \item as $T\to\infty$ (so $\beta\to0$), $\Pr(i)\to 1/K$: all states become equally likely;
  \item as $T\to 0$ (so $\beta\to\infty$), $\Pr$ concentrates on the states of minimal
    energy \added{(uniformly on the ground states if there are several)}.
\end{itemize}
\end{remark}

\subsection{A gas in a box}
\label{L07:subsec:gas}

But let us consider another (sketchy) physics analysis: the distribution of a gas in a box,
in the simplest non-interacting setting. Take $2N$ points, i.e.\ a configuration in
\[
  \Conf\big(2N,[-1,1]^3\big),
\]
\added{chosen uniformly at random.} Let $P_m$ be the probability that exactly $N+m$
particles have positive first coordinate.

There are $2^{2N}$ choices for the signs of the first coordinates, all equally likely, and
$\binom{2N}{N+m}$ of them have exactly $N+m$ particles with first coordinate $>0$. (This is
really the coin-flip analysis.) So
\[
  P_m=2^{-2N}\,\frac{(2N)!}{(N+m)!\,(N-m)!}.
\]

We use a very bad version of Stirling's approximation,
\[
  n!\sim\Big(\frac ne\Big)^n\sqrt{2\pi n}\;\leadsto\;n!\approx\Big(\frac ne\Big)^n .
\]
Then
\begin{align*}
  P_m&\approx 2^{-2N}\Big(\frac{2N}{e}\Big)^{2N}\Big/
     \Big[\Big(\frac{N+m}{e}\Big)^{N+m}\Big(\frac{N-m}{e}\Big)^{N-m}\Big]
   =\frac{N^{2N}}{(N+m)^{N+m}(N-m)^{N-m}}\\
  &=\Big(1+\frac mN\Big)^{-(N+m)}\Big(1-\frac mN\Big)^{-(N-m)}
   =\Big(1-\frac{m^2}{N^2}\Big)^{-N}\Big(1+\frac mN\Big)^{-m}\Big(1-\frac mN\Big)^{m}\\
  &\overset{m\ll N}{\approx}\big(e^{-m^2/N^2}\big)^{-N}\big(e^{m/N}\big)^{-m}\big(e^{-m/N}\big)^{m}
   = e^{-m^2/N}.
\end{align*}
Since the crude Stirling formula loses the factors $\sqrt{2\pi n}$, we only keep the shape
and write
\[
  P_m\approx P_0\,e^{-m^2/N}
\]
(a crude central limit theorem). Normalising,
\[
  1=\sum_{m=-N}^{N}P_m\approx\int_{-\infty}^{\infty}P_0\,e^{-m^2/N}\,dm=P_0\sqrt{\pi N}
  \quad\Longrightarrow\quad
  P_0=\frac{1}{\sqrt{\pi N}},\qquad P_m\approx\frac{1}{\sqrt{\pi N}}\,e^{-m^2/N}.
\]
Comparing with $e^{-m^2/(2\sigma^2)}$ we get $\sigma^2=N/2$, so $\sigma=\sqrt{N/2}$.
\added{This agrees with the exact variance $2N\cdot\frac12\cdot\frac12=N/2$ of the
binomial distribution, and $P_0=1/\sqrt{\pi N}$ is exactly what the full Stirling formula
gives for $\binom{2N}{N}4^{-N}$.}

So the number of points with positive first coordinate is approximately Gaussian, centred
at the expected value $N$ with standard deviation $\sqrt{N/2}$. For example, for large $N$,
about $95\%$ of configurations are within $2\sigma=2\sqrt{N/2}$ of the expected value.
\added{The relative fluctuation $\sigma/N=1/\sqrt{2N}$ tends to $0$.}

\begin{remark}
\added{The approximations $\ln(1\pm x)\approx\pm x-x^2/2$ used above each drop cubic
terms of size $m^3/N^2$, but these cancel between the factors $(1+\frac mN)^{-m}$ and
$(1-\frac mN)^{m}$ (the exact exponent is even in $m$). The first neglected term is
$-m^4/(6N^3)$, so the Gaussian approximation is accurate for $m\ll N^{3/4}$; this covers
the typical range $m=O(\sqrt N)$.}\fixed{error order was stated as $m^3/N^2$}
\end{remark}

So if we pick a point in high dimensions, i.e.\ a configuration in
$[-1,1]^{3\cdot 2N}$\fixed{notes: $[-1,1]^{3N}$; $2N$ points in $\R^3$ give $6N$ coordinates}
(and the same argument works for the other variables), we find roughly half of the
coordinates positive and half negative, just like the coin flips. This makes sense.
We can also talk about energy, etc., and generic configurations in this way.

But we may be interested in smaller configurations, like nearest neighbours, where
Stirling's approximation is bad.
\unclear{reading of ``smaller configurations, like nearest neighbours'' uncertain}
How do we analyse such things?

\subsection{Points on the circle}
\label{L07:subsec:circle}

Let us look at points on $\Sph^1$, \added{which we take to have length $1$}. Consider
$\operatorname{Reduced}(3,\Sph^1)$, three points on the circle modulo rotation. Use the
rotation to fix one point at $0$; going around the circle, the three gaps between
consecutive points are $x$, $y$ and $1-x-y$.
\fixed{notes write $1-x+y$}
So the reduced space is described by
\[
  \{(x,y): x>0,\ y>0,\ x+y<1\},
\]
an open triangle, whose boundary edges $x=0$, $y=0$, $x+y=1$ are the collisions.
\added{(For labelled points there are two such triangles, one for each cyclic order of the
labels $1,2,3$; compare Example~\ref{L07:ex:confS1} below.)}
In general, for $N$ points we have a simplex
\added{$\{(g_1,\dots,g_N): g_i>0,\ \sum g_i=1\}$ of gaps, an open $(N-1)$-simplex
(one for each of the $(N-1)!$ cyclic orders of the labels)}.

We can attach the \emph{injectivity radius} function
\[
  \rho=\min_{i\neq j}\tfrac12\,\operatorname{geodesic}(x_i,x_j).
\]
\added{Since the smallest gap is at most $1/N\le 1/2$, $\rho$ is half the smallest gap:
$\rho=\frac12\min(x,y,1-x-y)$ for $N=3$. It has a unique maximum $\rho=1/6$ at the
equilateral configuration $x=y=1/3$; its level sets are triangles, and it fails to be
smooth along the three segments from the centre to the vertices. In general the maximum is
$1/(2N)$, attained at the regular $N$-gon.}

\begin{center}
\begin{tikzpicture}[scale=3.6]
  \draw[dashed] (0,0) -- (1,0) -- (0,1) -- cycle;
  \draw[thin,gray] (1/3,1/3) -- (0,0) (1/3,1/3) -- (1,0) (1/3,1/3) -- (0,1);
  \draw[addcol] (1/6,1/6) -- (2/3,1/6) -- (1/6,2/3) -- cycle;
  \fill (1/3,1/3) circle (0.5pt);
  \draw[gray,thin] (1/3,1/3) -- (0.62,0.62); \node[right] at (0.62,0.64) {\scriptsize $x=y=\tfrac13$};
  \node[below] at (0.5,0) {\scriptsize $y=0$};
  \node[left] at (0,0.5) {\scriptsize $x=0$};
  \node[above right] at (0.72,0.3) {\scriptsize $x+y=1$};
  \node[addcol] at (0.47,0.21) {\tiny $\rho=\frac1{12}$};
\end{tikzpicture}
\end{center}
\noindent\textit{Figure:} $\operatorname{Reduced}(3,\Sph^1)$ as an open triangle
(dashed boundary = collisions), with a level set of $\rho$ and the lines where $\rho$ is
not smooth.

\begin{remark}
\orig{p.\,13: closing remark largely illegible}
\unclear{remark ends ``\dots it's a (box--?)\dots''}
The notes end Lecture~7 with a remark about computing with the injectivity radius.
\added{The point to keep in mind for the next lecture is that $\rho$ is a minimum of smooth
functions, hence only piecewise smooth, so smooth Morse theory does not apply to it
directly.}
\end{remark}

% ======================================================================

\chapter{Lecture 8: (Smooth) Morse theory}
\label{L07:ch:lecture8}

\section{Configuration spaces on a line and Weyl chambers}
\label{L07:sec:weyl}

We started to talk about the spaces
\[
  \Conf(N,X)=\prod_{i=1}^N X\setminus\Delta,
  \qquad
  \Delta=\{(x_i)_{i=1}^N: x_i=x_j \text{ for some } i\neq j\}.
\]

\begin{example}[$\Conf(3,\Sph^1)$]
\label{L07:ex:confS1}
Fixing the first point at $0$ by a rotation, the other two points lie in
$\Sph^1\setminus\{0\}\cong(0,1)$ and must be distinct, so
\[
  \Conf(3,\Sph^1)\cong \Sph^1\times\big((0,1)^2\setminus\{x_2=x_3\}\big),
\]
\added{and the reduced space is the open square minus its diagonal: two open triangles, one
for each cyclic order (cf.\ Section~\ref{L07:subsec:circle}).}
\fixed{notes: ``$\Conf(3,\Sph^1)\approx$ square''; this is the reduced space}
\end{example}

\begin{center}
\begin{tikzpicture}[scale=1.6]
  \draw[dashed] (0,0) rectangle (1,1);
  \draw[dashed] (0,0) -- (1,1);
  \node[below] at (0.5,0) {\scriptsize $x_2\in(0,1)$};
  \node[left] at (0,0.5) {\scriptsize $x_3$};
\end{tikzpicture}
\end{center}

\begin{example}[$\Conf(N,\R)$ and the $A_{N-1}$ Weyl chambers]
Write $n=N$. Consider the affine subspaces
\[
  V=V_c=\Big\{(x_i)_{i=1}^n:\ \sum_i x_i=c\Big\}\subset\R^n .
\]
\added{Translation along $(1,\dots,1)$ splits off, $\Conf(n,\R)\cong\R\times(V_c\setminus\Delta)$.}
In these subspaces the hyperplanes $\{x_i=x_j\}\cap V$ are the boundaries of the Weyl
chambers in our construction of the $A_{n-1}$ root system in $V\subset\R^n$.
\added{Recall $A_{n-1}=\{e_i-e_j: i\neq j\}$, with $n(n-1)$ roots, simple roots
$e_i-e_{i+1}$, and Weyl group $\Sigma_n$ generated by the reflections in the hyperplanes
$x_i=x_j$ (i.e.\ swapping two coordinates). The connected components of
$V\setminus\Delta$ are exactly the $n!$ open Weyl chambers.}

For $n=3$\fixed{notes: ``$n=2$''; the example is $A_2$, $n=3$} we had the simple roots
\[
  (1,-1,0),\qquad (0,1,-1).
\]
So the root system is given by the six vectors
\[
  \pm(1,-1,0),\qquad \pm(0,1,-1),\qquad \pm(1,0,-1),
\]
generated by swapping coordinates, i.e.\ by reflecting through the hyperplanes $x_i=x_j$,
$i\neq j$. To draw them in the plane $V_3=\{x+y+z=3\}$ we translate by $(1,1,1)$.
\end{example}

\begin{center}
\begin{tikzpicture}[scale=1.25]
  \coordinate (A) at (-1.732,-1);
  \coordinate (B) at (1.732,-1);
  \coordinate (C) at (0,2);
  \draw (A) -- (B) -- (C) -- cycle;
  \node[below left] at (A) {\scriptsize $(3,0,0)$};
  \node[below right] at (B) {\scriptsize $(0,3,0)$};
  \node[above] at (C) {\scriptsize $(0,0,3)$};
  \coordinate (O) at (barycentric cs:A=1,B=1,C=1);
  % walls (medians)
  \draw[dashed] (C) -- (barycentric cs:A=1,B=1);
  \draw[dashed] (A) -- (barycentric cs:B=1,C=1);
  \draw[dashed] (B) -- (barycentric cs:A=1,C=1);
  \node[below] at (barycentric cs:A=1,B=1) {\scriptsize $x=y$};
  \node[above right] at (barycentric cs:B=1,C=1) {\scriptsize $y=z$};
  \node[above left] at (barycentric cs:A=1,C=1) {\scriptsize $x=z$};
  % roots + (1,1,1)
  \foreach \a/\b/\c in {2/0/1, 0/2/1, 1/2/0, 2/1/0, 1/0/2, 0/1/2} {
    \draw[->,addcol,thick] (O) -- (barycentric cs:A=\a,B=\b,C=\c);
    \fill[addcol] (barycentric cs:A=\a,B=\b,C=\c) circle (1.2pt);
  }
  \fill (O) circle (1.2pt);
  \node[below right] at ($(O)+(0.05,-0.02)$) {\tiny $(1,1,1)$};
  \node[addcol,left] at ($(barycentric cs:A=2,B=0,C=1)+(-0.05,0)$) {\tiny $(2,0,1)$};
\end{tikzpicture}
\end{center}
\noindent\textit{Figure:} the plane $x+y+z=3$. The dashed lines are the walls
$x_i=x_j$; they cut the plane into the six Weyl chambers (the components of
$V_3\setminus\Delta$). The arrows are the six roots of $A_2$, translated by $(1,1,1)$; they
end at the trisection points of the edges, e.g.\ $(1,1,1)+(1,-1,0)=(2,0,1)$, and each
root is perpendicular to one wall.

\subsection{Manifolds}
\label{L07:subsec:manifolds}

We will consider spaces $M$ that are Hausdorff, second countable topological spaces: the
topology separates points, and there is a countable basis for the topology. (Basically
this is to eliminate pathological examples.)

\begin{definition}[Homeomorphism]
A \emph{homeomorphism} $\varphi\colon M\to N$ is a continuous bijection with continuous
inverse. Equivalently: $\varphi$ is a bijection, and $U\subset M$ is open if and only if
$\varphi(U)\subset N$ is open.
\end{definition}

\begin{definition}[Coordinate chart]
A \emph{coordinate chart} for a space $M$ is a homeomorphism from a coordinate
neighbourhood $U$ \added{(an open subset of $M$)} onto an open subset of $\R^n$,
\[
  \varphi_U\colon U\subset M\longrightarrow V\subset\R^n .
\]
\end{definition}

\begin{center}
\begin{tikzpicture}[scale=0.85]
  \draw (0,0) ellipse (2.2 and 1.2); \node at (1.6,0.6) {$M$};
  \draw[fill=gray!15] (-0.5,0) circle (0.55); \node at (-0.5,0) {\small $U$};
  \draw[->] (-0.5,-1.0) -- (-0.5,-2.2) node[midway,left] {$\varphi_U$};
  \draw (-2.4,-3.3) -- (1.6,-3.3) -- (2.2,-2.5) -- (-1.8,-2.5) -- cycle;
  \node at (1.6,-2.8) {\small $\R^n$};
  \draw[fill=gray!15] (-0.3,-2.9) ellipse (0.6 and 0.28); \node at (-0.3,-2.9) {\small $V$};
\end{tikzpicture}
\end{center}

A smooth bijection with smooth inverse is a \emph{diffeomorphism}. So in the smooth
setting, coordinate changes should be diffeomorphisms.

A \emph{manifold} is a space $M$ with a collection of charts $\{\varphi_{U_i}\}$ (an
\emph{atlas}), indexed by a cover $\{U_i\}$ of $M$, together with \emph{transition maps}
that patch the charts together:
\[
  \tau_{ij}=\tau_{U_iU_j}=\varphi_{U_j}\circ\varphi_{U_i}^{-1}\colon
  \varphi_{U_i}(U_i\cap U_j)\to\varphi_{U_j}(U_i\cap U_j).
\]

\begin{center}
\begin{tikzpicture}[scale=0.85]
  \draw (0,0) ellipse (2.4 and 1.2);
  \draw (-0.55,0) circle (0.7); \draw (0.45,0) circle (0.7);
  \begin{scope} \clip (-0.55,0) circle (0.7); \fill[gray!35] (0.45,0) circle (0.7); \end{scope}
  \node at (-0.95,0) {\small $U_i$}; \node at (0.85,0) {\small $U_j$};
  \draw (-2,-3) circle (0.6); \node at (-2.2,-3) {\small $V_i$};
  \draw (2,-3) circle (0.6); \node at (2.2,-3) {\small $V_j$};
  \draw[->] (-0.9,-0.9) -- (-1.8,-2.3) node[midway,left] {$\varphi_{U_i}$};
  \draw[->] (0.8,-0.9) -- (1.8,-2.3) node[midway,right] {$\varphi_{U_j}$};
  \draw[->] (-1.3,-2.75) to[bend left=20] node[midway,above] {\small $\tau_{ij}$} (1.3,-2.75);
  \draw[->] (1.3,-3.25) to[bend left=20] node[midway,below] {\small $\tau_{ji}=\tau_{ij}^{-1}$} (-1.3,-3.25);
\end{tikzpicture}
\end{center}

\begin{definition}[Smooth structure]
\label{L07:def:smoothstructure}
A \emph{smooth differentiable structure} on $M$ is a collection of (smooth) coordinate
charts $\varphi_i\colon U_i\to V_i\subset\R^n$, $U_i\subset M$, such that
\begin{enumerate}
  \item $M=\bigcup_i U_i$;
  \item every transition map $\tau_{ij}=\varphi_j\circ\varphi_i^{-1}$ is smooth
    \added{(on its domain $\varphi_i(U_i\cap U_j)$)};
  \item it is a \emph{maximal} atlas: if a chart is compatible with all $\varphi_i$ as
    above, then it is in the collection.
\end{enumerate}
A topological space with a smooth structure is a \emph{smooth manifold} (of dimension $n$).
\end{definition}

Basically we can ignore maximality for computations, as maximal atlases are ``unique'':
\added{every atlas satisfying (1) and (2) is contained in exactly one maximal atlas, so it
suffices to write down any atlas.}

\section{Examples; smooth maps}
\label{L07:sec:examples}

\begin{example}
\label{L07:ex:manifolds}
\leavevmode
\begin{itemize}
  \item $\R^n$: choose the identity map as the chart \added{(one chart, $U=\R^n$)}.
  \item $\C^n\to\R^{2n}$, using $(\operatorname{Re},\operatorname{Im})$ on each $U_i$.
  \item $\Sph^n\to\R^n$: stereographic projection on $\Sph^n\setminus\{N\}$ and on
    $\Sph^n\setminus\{S\}$ \added{(north and south poles; the transition map is the
    inversion $y\mapsto y/\abs{y}^2$ on $\R^n\setminus\{0\}$, which is smooth)}.
  \item On tori, by products of the $\Sph^1$ charts.
  \item On open subsets of manifolds, by carrying over the subspace topology
    \added{and restricting the charts}.
\end{itemize}
\added{In particular $\Conf(N,X)$ is an open subset of $X^N$, so it is a smooth manifold of
dimension $N\dim X$ whenever $X$ is.}
\end{example}

We can define \emph{smooth maps} between manifolds.

\begin{definition}[Smooth map]
\label{L07:def:smoothmap}
A map $f\colon M\to N$ is \emph{smooth} if for each $p\in M$ and some charts $\varphi$ on
$U\ni p$ for $M$ and $\psi$ on $V\ni f(p)$ for $N$ (and then for all such charts),
\added{with $f(U)\subset V$,} the induced map
\[
  \begin{array}{ccc}
    U & \xrightarrow{\ f\ } & V\\[2pt]
    {\scriptstyle\varphi}\downarrow & & \downarrow{\scriptstyle\psi}\\[2pt]
    \varphi(U) & \xrightarrow{\ \psi f\varphi^{-1}\ } & \psi(V)
  \end{array}
\]
$\psi\circ f\circ\varphi^{-1}$ is smooth on its domain.
\added{(``Some $\Rightarrow$ all'' because the transition maps are smooth.)}
\end{definition}
\orig{p.\,12 (= scan p.\,25) repeats this diagram}

Similarly for maps between manifolds of other kinds. So we can define diffeomorphisms
(homeomorphisms) of manifolds: a (smooth) bijection $M\to N$ with (smooth) inverse,
written $M\cong N$.
If we change the underlying conditions on the maps \added{(e.g.\ continuous, $C^k$, or
piecewise linear transition maps)} we can still define coordinate charts, atlases, etc.
\orig{crossed out: ``Sadly, not all the maps we want to work with will be smooth''}

\section{Homotopy type}
\label{L07:sec:homotopy}

\orig{a first ``8.3 Morse theory'' is crossed out here}
We will want to talk about \emph{homotopy type} as well, since we may be dealing with
things that are not diffeomorphic but are topologically ``similar''.

\begin{definition}[Homotopy]
\label{L07:def:homotopy}
Two maps $f,g\colon M\to N$ are \emph{homotopic}, written $f\simeq g$, if there exists a
continuous function
\[
  H\colon M\times I\to N,\qquad I=[0,1],
\]
such that $H(x,0)=f(x)$ and $H(x,1)=g(x)$ for all $x\in M$.
\end{definition}

\begin{definition}[Homotopy equivalence]
\label{L07:def:hequiv}
Two (smooth) spaces $M$ and $N$ are \emph{homotopy equivalent} (have the same homotopy
type) if there are maps
\[
  f\colon M\to N,\qquad g\colon N\to M
  \qquad\text{such that}\qquad
  f\circ g\simeq\mathrm{id}_N,\quad g\circ f\simeq\mathrm{id}_M .
\]
\orig{letters $X,Y$ corrected to $M,N$ in blue on the page}
\end{definition}

The idea: the spaces can be continuously (smoothly) deformed into each other.

\begin{example}[$D^2\simeq\{0\}$]
\label{L07:ex:disk}
Let $f\colon D^2\to\{0\}$ be the radial retraction to $0$ (the constant map) and
$g\colon\{0\}\to D^2$ the inclusion. In polar coordinates put
\[
  h\big(t,(r,\theta)\big):=\big((1-t)r,\theta\big).
\]
Then $h(0,\cdot)=\mathrm{id}_{D^2}$ and $h(1,\cdot)=g\circ f$ (the map $D^2\to 0$), so
$g\circ f\simeq\mathrm{id}_{D^2}$; and $f\circ g=\mathrm{id}_{\{0\}}$.
\added{($h$ is continuous also at $r=0$, since $(1-t)\cdot 0=0$ regardless of $\theta$.)}
\end{example}

\begin{definition}[Relative homotopy; deformation retract]
\label{L07:def:relhom}
A \emph{homotopy relative to a subspace} is a homotopy which fixes the points of the
subspace: for $f,g\colon M\to N$ and $K\subset M$, $f$ and $g$ are \emph{homotopic
relative to $K$} if there is a homotopy $h\colon M\times I\to N$ between $f$ and $g$ such
that
\[
  h(k,t)=f(k)=g(k)\qquad\text{for all } k\in K,\ t\in I .
\]
If $f=\mathrm{id}_M$ and $g$ is a \emph{retraction}, i.e.\ a continuous map $M\to K$ with
$g|_K=\mathrm{id}$ \added{(followed by the inclusion $K\hookrightarrow M$)}, then this
homotopy defines a \emph{strong deformation retraction} of $M$ onto $K$.
\end{definition}
See the last example: $h$ is a strong deformation retraction of $D^2$ onto $\{0\}$.

\section{An aside about homotopy groups}
\label{L07:sec:homotopygroups}

We consider homotopy classes of maps to pointed spaces
\added{(with homotopies relative to the base point)}. Consider
$f\colon \Sph^1\to(X,x_0)$, or equivalently paths $f\colon I\to X$ with
$f(0)=f(1)=x_0$. These have a group structure under concatenation,
\[
  (f\cdot g)(t)=\begin{cases} f(2t), & t\in[0,\tfrac12],\\ g(2t-1), & t\in[\tfrac12,1].\end{cases}
\]
This is the fundamental group $\pi_1(X,x_0)$.

In general one can consider maps $f\colon\Sph^n\to(X,x_0)$, viewed as maps
$f\colon I^n\to X$ with $f(\partial I^n)=x_0$, with concatenation defined appropriately
\added{(in the first coordinate)}; this gives $\pi_n(X,x_0)$. If $X$ is path connected we
may ignore the base point $x_0$ \added{(up to isomorphism)}.

\begin{example}
Draw $(D^2,x_0)\setminus\{0\}$.
\orig{example left blank on the page}
\added{The punctured disk deformation retracts (radially) onto the circle of radius
$\abs{x_0}$, so $\pi_1(D^2\setminus\{0\},x_0)\cong\pi_1(\Sph^1)\cong\Z$, generated by
the loop going once around the puncture. Compare $\pi_1(D^2)=0$ from
Example~\ref{L07:ex:disk}: removing a point changes the homotopy type.}
\end{example}

We will not deal with this too much, but note that this is a nice invariant for spaces:
\added{homotopy equivalent spaces have isomorphic homotopy groups.} This is also the
source of the braid groups of Section~\ref{L07:sec:braid}.

\section{Morse theory}
\label{L07:sec:morse}
\orig{numbered ``8.2'' in the notes}

Morse theory is a fairly general way to study manifolds and functions on them. The idea is
that the structure of the sets of critical points of functions and the topological
structure of the manifold are related. (We will follow Milnor and others
\added{(see J.~Milnor, \emph{Morse Theory}, Annals of Math.\ Studies 51, Part~I)}.)

\begin{example}[The classic picture]
\label{L07:ex:torus}
Consider $f\colon\T^2\to\R$ given by
\[
  \T^2\hookrightarrow\R^3\xrightarrow{\ \mathrm{proj}_z\ }\R,
\]
the height function $h=\mathrm{proj}_z(x,y,z)=z$ on a torus standing upright, normalised
so that it takes values in $[0,1]$. It has four critical points: $p$ (minimum, height $0$),
$q$ and $r$ (the bottom and top of the hole) and $s$ (maximum, height $1$).
\end{example}

\begin{center}
\begin{tikzpicture}[scale=0.95]
  \draw[thick] (0,0) ellipse (1.25 and 1.9);
  \draw[thick] (-0.1,-0.75) .. controls (-0.38,-0.3) and (-0.38,0.3) .. (-0.1,0.75);
  \draw[thick] (-0.17,-0.6) .. controls (0.12,-0.25) and (0.12,0.25) .. (-0.17,0.6);
  \foreach \y/\n in {-1.9/p,-0.68/q,0.68/r,1.9/s} {
    \draw[dashed,gray] (-1.6,\y) -- (3.2,\y);
  }
  \fill (0,-1.9) circle (1.5pt) node[below left] {$p$};
  \fill (-0.13,-0.68) circle (1.5pt) node[left] {\small $q$};
  \fill (-0.13,0.68) circle (1.5pt) node[left] {\small $r$};
  \fill (0,1.9) circle (1.5pt) node[above left] {$s$};
  \node[right] at (3.2,-1.9) {\small $0$};
  \node[right] at (3.2,1.9) {\small $1$};
  \draw[->] (2.4,-2.3) -- (2.4,2.3) node[above] {\small $h=z$};
\end{tikzpicture}
\end{center}
\noindent\textit{Figure:} the torus standing on the plane $z=0$, seen from the side, with
the height function and its critical points.

Consider the \emph{sublevel sets}
\[
  M_a:=\{x\in M: f(x)<a\}.
\]
\added{(Milnor uses $M^a=\{f\le a\}$; for $a$ not a critical value the two have the same
homotopy type.)}
\begin{itemize}
  \item $a<0=f(p)$: $M_a=\emptyset$.
  \item $f(p)<a<f(q)$: $M_a$ is a $2$-cell, i.e.\ a $2$-ball (a disk).
  \item $f(q)<a<f(r)$: $M_a$ is a cylinder.
  \item $f(r)<a<f(s)$: $M_a$ is a torus minus a disk.
  \item $f(s)<a$: $M_a$ is the whole torus.
\end{itemize}

Up to homotopy, these are handle or \emph{cell attachments}.

\begin{definition}[Attaching a cell]
\label{L07:def:cell}
Let $e^k$ be a closed $k$-dimensional ball, with $\partial e^k=\Sph^{k-1}$, let $Y$ be a
topological space and $g\colon\Sph^{k-1}\to Y$ continuous. Then $Y\cup_g e^k$, the result
of \emph{attaching a $k$-cell to $Y$ by $g$}, is defined as
\[
  Y\cup_g e^k=\big(Y\sqcup e^k\big)/\sim,
  \qquad x\sim g(x)\ \text{ for } x\in\Sph^{k-1}.
\]
\end{definition}

\orig{scan pp.\,28--29 are in reverse order (circled 16, 15)}
The sublevel sets of the torus, up to homotopy, are:
\begin{itemize}
  \item a disk $\simeq$ a point ($0$-cell);
  \item the cylinder $\simeq$ a disk with a $1$-handle attached (a ``U'' shape)
    $\simeq$ a point with a $1$-cell attached, i.e.\ a circle;
  \item the torus minus a disk $\simeq$ a disk with two $1$-handles
    $\simeq$ a point with two $1$-cells attached, i.e.\ a figure eight
    $\Sph^1\vee\Sph^1$;
  \item the torus $\simeq$ the figure eight with a $2$-cell attached
    \added{along the loop $aba^{-1}b^{-1}$}.
\end{itemize}
\orig{p.\,15: sketches of the handle pictures; one equivalence marked ``?''}

\begin{center}
\begin{tikzpicture}[scale=0.85]
  % point
  \fill (0,0) circle (2pt);
  \node at (0,-1.3) {\small $\{p\}$};
  % circle
  \begin{scope}[xshift=2.4cm]
    \draw[thick] (0,0.5) circle (0.5); \fill (0,0) circle (2pt);
    \node at (0,-1.3) {\small $e^0\cup e^1$};
  \end{scope}
  % figure eight
  \begin{scope}[xshift=5.2cm]
    \draw[thick] (-0.5,0) circle (0.5); \draw[thick] (0.5,0) circle (0.5);
    \fill (0,0) circle (2pt);
    \node at (0,-1.3) {\small $e^0\cup e^1\cup e^1$};
  \end{scope}
  % torus as square
  \begin{scope}[xshift=8.4cm,yshift=-0.6cm]
    \fill[gray!15] (0,0) rectangle (1.2,1.2);
    \draw[thick,->] (0,0) -- (0.65,0); \draw[thick] (0.6,0) -- (1.2,0);
    \draw[thick,->] (0,1.2) -- (0.65,1.2); \draw[thick] (0.6,1.2) -- (1.2,1.2);
    \draw[thick,->] (0,0) -- (0,0.65); \draw[thick] (0,0.6) -- (0,1.2);
    \draw[thick,->] (1.2,0) -- (1.2,0.65); \draw[thick] (1.2,0.6) -- (1.2,1.2);
    \node[below] at (0.6,0) {\scriptsize $a$}; \node[above] at (0.6,1.2) {\scriptsize $a$};
    \node[left] at (0,0.6) {\scriptsize $b$}; \node[right] at (1.2,0.6) {\scriptsize $b$};
    \node at (0.6,-0.7) {\small $\cdots\cup e^2=\T^2$};
  \end{scope}
  \node at (1.2,0) {$\leadsto$};
  \node at (3.75,0) {$\leadsto$};
  \node at (7.0,0) {$\leadsto$};
\end{tikzpicture}
\end{center}
\noindent\textit{Figure:} homotopy types of the sublevel sets $M_a$ as $a$ passes
$p$, $q$, $r$, $s$.

The homotopy type (and the diffeomorphism type, outside of a neighbourhood) changes
exactly at the critical points of the function $f$: those points where the gradient is
zero, $\nabla f=0$.

If we look at the Hessian \added{$\Hess f_x=\big(\partial^2 f/\partial u_i\partial u_j(x)\big)$
in local coordinates $u$}, which here is non-degenerate at each critical point, we see
\added{(writing the signature as (\#positive, \#negative) eigenvalues)}:
\[
  \operatorname{sig}(\Hess f_p)=(2,0),\qquad
  \operatorname{sig}(\Hess f_q)=\operatorname{sig}(\Hess f_r)=(1,1),\qquad
  \operatorname{sig}(\Hess f_s)=(0,2).
\]
We call the negative part the \emph{index}, i.e.\ the number of decreasing directions at the
critical point, and we see that it corresponds to the dimension of the attached cell:
index $0$ at $p$ ($0$-cell), index $1$ at $q$ and $r$ ($1$-cells), index $2$ at $s$
($2$-cell).

This is why we want to start with smooth functions.

\begin{definition}[Morse function]
\label{L07:def:morse}
We will call a smooth function $f\colon M\to\R$ \emph{Morse} if it has no degenerate
critical points\fixed{notes say ``no nondegenerate critical pts''}, i.e.\ \added{at every
critical point $x$ (where $df_x=0$) the Hessian $\Hess f_x$ is a non-singular matrix}.
\end{definition}

We will formalise this next time.

\begin{addedblock}
For reference, the two basic theorems (Milnor, \emph{Morse Theory}, Thms.~3.1 and 3.2):
let $f\colon M\to\R$ be smooth and write $M^a=f^{-1}(-\infty,a]$.
(A) If $a<b$, $f^{-1}[a,b]$ is compact and contains no critical points, then $M^a$ is
diffeomorphic to $M^b$, and $M^a$ is a deformation retract of $M^b$.
(B) If $x$ is a non-degenerate critical point of index $\lambda$ with $f(x)=c$, and
$f^{-1}[c-\varepsilon,c+\varepsilon]$ is compact and contains no other critical point, then
for small $\varepsilon>0$, $M^{c+\varepsilon}$ has the homotopy type of
$M^{c-\varepsilon}\cup e^{\lambda}$.
\end{addedblock}
% ===== end chapters/L07.tex =====
% ===== begin chapters/L09.tex =====
% ============================================================
% MAT670 -- Lectures 9 and 10 (source: MAT670_9X, pp. 1--29)
%   Lecture 9:  pp. 1--15
%   Lecture 10: pp. 16--29
% ============================================================
\chapter{Lecture 9: More on functions and their critical points}
\label{L9:chap}

Last time we defined a smooth function $f\colon M\to\R$ on a manifold $M$ to be
\emph{Morse} if it has no degenerate critical points.%
\fixed{notes say ``no non-degenerate critical points''}
\added{That is, every critical point of $f$ is non-degenerate.}

\section{What does this mean?}
\label{L9:sec:meaning}

To make sense of this definition we need to understand criticality, so we introduce
the tangent space $T_pM$ at a point $p\in M$.

\begin{center}
\begin{tikzpicture}[scale=0.8]
  \draw[thick] (0,0) ellipse (2 and 1.3);
  \node at (1.5,-0.6) {$M$};
  \draw[fill=blue!8,opacity=0.8] (-1.6,-0.1) -- (0.6,-0.5) -- (1.6,0.7) -- (-0.6,1.1) -- cycle;
  \fill (0,0.3) circle (1.5pt) node[below left] {$p$};
  \node at (1.9,1.0) {$T_pM$};
\end{tikzpicture}
\end{center}

We consider velocities of curves
\[
  \gamma\colon[-1,1]\to M,\qquad \gamma(0)=p,
\]
and, with respect to a chart $\varphi\colon U\to\R^n$ with $p\in U\subset M$, we ask
that $\varphi\circ\gamma\colon[-1,1]\to\R^n$ be differentiable at $0$.
The classes of such curves, where two curves are identified when their tangents
\added{$(\varphi\circ\gamma)'(0)$} at $0$ agree, define the \emph{tangent space} $T_pM$ at $p$.
\added{(This identification does not depend on the chart, since the transition maps
are diffeomorphisms; the tangent vector represented by $\gamma$ is written $\gamma'(0)$.)}

Given a smooth map $f\colon M\to N$, there is an induced \emph{linear} map
\[
  f_*\colon T_pM\to T_{f(p)}N,\qquad f_*\bigl(\gamma'(0)\bigr)=(f\circ\gamma)'(0),
\]
\fixed{notes: target $T_pN$}
called the \emph{differential} or \emph{push-forward} of $f$.

\begin{definition}
\label{L9:def:critical}
Let $f\colon M\to\R$ be smooth. A point $p\in M$ is a \emph{critical point} of $f$ if
\[
  f_*\colon T_pM\to T_{f(p)}\R
\]
is the zero map. The value $f(p)$ is then called a \emph{critical value}.
\end{definition}

In local coordinates $(x_1,\dots,x_n)$ \added{about $p$} this is written
\[
  \frac{\partial f}{\partial x_i}\Big|_p = 0\qquad\text{for all } i=1,\dots,n .
\]
Also, there is a bilinear form \added{on $T_pM$, well defined at a critical point $p$,} which can be
written similarly in coordinates as
\[
  H_f(p)=\Bigl(\frac{\partial^2 f}{\partial x_i\,\partial x_j}\Big|_p\Bigr)_{i,j},
  \qquad\text{also written } f_{**}(p).
\]
This is the \emph{Hessian} of $f$ at $p$.

\begin{definition}
\label{L9:def:nondeg}
A critical point $p$ of $f$ is \emph{non-degenerate} if the Hessian $H_f(p)$ is
non-degenerate, i.e.\ its matrix is non-singular. The \emph{index} of $f$ at $p$ is the
number of negative eigenvalues of $H_f(p)$ \added{(the dimension of a maximal subspace of
$T_pM$ on which $H_f(p)$ is negative definite)}.
\end{definition}

The notion of index is well defined in this case.
\added{Indeed, under a change of coordinates the Hessian at a critical point changes by
$H\mapsto J^{\mathsf T}HJ$ with $J$ the (invertible) Jacobian, and by Sylvester's law of inertia the number
of negative eigenvalues is unchanged.}

\section{Why do we care about non-degenerate critical points?}
\label{L9:sec:morselemma}
\orig{numbered 6.2--6.4 in the notes; read as 9.2--9.4}

\begin{lemma}[Morse Lemma]
\label{L9:lem:morse}
If $p$ is a non-degenerate critical point of a smooth function $f$, there is a local
coordinate system $(x_1,\dots,x_n)$ about $p$ \added{with $p$ at the origin} such that
\[
  f = \added{f(p)}\ - \sum_{i=1}^{\lambda} x_i^2 + \sum_{j=\lambda+1}^{n} x_j^2 ,
\]
where $\lambda$ is the index of $f$ at $p$.
\end{lemma}

\begin{proof}[Proof (sketch)]
We may shift so that $f(p)=0$ \added{and $p=0$ in some chart}. By the fundamental theorem of
calculus and the integral form of Taylor's theorem,
\[
  f(x_1,\dots,x_n)=\int_0^1 \frac{d}{dt}f(tx)\,dt
  =\int_0^1\sum_{i=1}^n \frac{\partial f}{\partial x_i}(tx_1,\dots,tx_n)\,x_i\,dt ,
\]
\added{so $f=\sum_i x_i g_i$ with $g_i$ smooth and $g_i(0)=\frac{\partial f}{\partial x_i}(0)=0$.
Applying the same argument to each $g_i$,}
and in fact
\[
  f(x_1,\dots,x_n)=\sum_{i,j} h_{ij}(x_1,\dots,x_n)\,x_ix_j
\]
for some smooth functions $h_{ij}$. We can symmetrize ($h_{ij}=h_{ji}$, replacing $h_{ij}$
by $\tfrac12(h_{ij}+h_{ji})$) and compute
\[
  h_{ij}(0)=\frac12\,\frac{\partial^2 f}{\partial x_i\,\partial x_j}(0)
\]
\fixed{notes: $\partial f$ for $\partial^2 f$}
in some coordinate system for $T_pM$.
This can be inductively diagonalized in a neighbourhood of $p$, for example via the
spectral theorem plus a continuity argument.

\emph{Example in dimension 2.} Write
\[
  f(x,y)=x^2h_{11}+2xy\,h_{12}+y^2h_{22},
  \qquad \frac{\partial^2 f}{\partial x^2}\Big|_{(0,0)}=2h_{11}(0,0),\ \text{etc.}
\]
By a \added{linear} coordinate change we may assume $h_{11}(0,0)\neq0$; by continuity there is
a neighbourhood on which $h_{11}\neq0$. Then we can define a local change of coordinates
\[
  \bar x=\sqrt{\abs{h_{11}}}\Bigl(x+\frac{h_{12}}{h_{11}}\,y\Bigr)
  \quad\Longrightarrow\quad
  \bar x^2=\abs{h_{11}}\Bigl(x^2+2\frac{h_{12}}{h_{11}}xy+\frac{h_{12}^2}{h_{11}^2}y^2\Bigr).
\]
In the case $h_{11}>0$,
\[
  \bar x^2=h_{11}x^2+2h_{12}xy+\frac{h_{12}^2}{h_{11}}y^2 ,
\]
and since
\[
  f=\underbrace{x^2h_{11}+2xy\,h_{12}}_{=\ \bar x^2-\frac{h_{12}^2}{h_{11}}y^2}+\,y^2h_{22},
  \qquad\text{we get}\qquad
  f=\bar x^2+\Bigl(h_{22}-\frac{h_{12}^2}{h_{11}}\Bigr)y^2 .
\]
\added{Non-degeneracy at $0$ means $h_{11}h_{22}-h_{12}^2\neq0$ there, so the coefficient of $y^2$
is non-zero near $0$, and $\bar y=\sqrt{\abs{h_{22}-h_{12}^2/h_{11}}}\;y$ finishes the job.
(The Jacobian of $(x,y)\mapsto(\bar x,\bar y)$ at $0$ is invertible, so this is a coordinate
change by the inverse function theorem.)}
The proof proceeds the same way for the other cases and for higher dimensions.
\begin{addedblock}
In general, after a linear change we have $h_{11}(0)\ne0$; put
$\bar x_1=\sqrt{\abs{h_{11}}}\,\bigl(x_1+\sum_{j>1}x_jh_{1j}/h_{11}\bigr)$, so that
$f=\pm\bar x_1^2+\sum_{i,j\ge2}x_ix_j\,h'_{ij}$ with $h'_{ij}$ smooth and symmetric and
$(h'_{ij}(0))$ non-singular; induct on $n$. (See Milnor, \emph{Morse Theory}, Lemma 2.2.)
\end{addedblock}
\end{proof}

\begin{corollary}
\label{L9:cor:isolated}
Non-degenerate critical points are isolated.
\end{corollary}
\begin{addedblock}
Indeed, in Morse coordinates $df=(-2x_1,\dots,-2x_\lambda,2x_{\lambda+1},\dots,2x_n)$ vanishes
only at the origin.
\end{addedblock}

\section{One-parameter groups of diffeomorphisms}
\label{L9:sec:oneparam}

\begin{definition}
\label{L9:def:oneparam}
A \emph{one-parameter group of diffeomorphisms} of a manifold $M$ is a smooth map
\[
  \Phi\colon\R\times M\to M
\]
such that
\begin{enumerate}[nosep]
  \item for each $t\in\R$, the map $\varphi_t\colon M\to M$ defined as
        $\varphi_t(p)=\Phi(t,p)$ is a diffeomorphism of $M$ onto $M$, and
  \item for all $t,s\in\R$, $\varphi_{t+s}=\varphi_t\circ\varphi_s$.
\end{enumerate}
\end{definition}
\orig{$\Phi$ vs.\ $\varphi_t$ distinction added in blue on the page}

This defines a vector field $X$ on $M$: for a smooth $f\colon M\to\R$, $X$ is characterized by
\[
  X_q(f)=\lim_{h\to0}\frac{f(\varphi_h(q))-f(q)}{h}.
\]
So this is an assignment of a direction to each point $q$ of $M$, and $X_q$ takes a
directional derivative of $f$. We say that $X$ \emph{generates} the group $\varphi$.

\begin{lemma}
\label{L9:lem:flow}
A smooth vector field on $M$ which vanishes outside of a compact set $K\subset M$ generates
a unique one-parameter group of diffeomorphisms of $M$.
\end{lemma}

\begin{proof}
Consider a smooth curve $\gamma\colon t\mapsto\gamma(t)\in M$ with velocity vector
$\frac{d\gamma}{dt}\in T_{\gamma(t)}M$, characterized by
\[
  \frac{d\gamma}{dt}f=\lim_{h\to0}\frac{f(\gamma(t+h))-f(\gamma(t))}{h}.
\]
Let $\varphi$ be a one-parameter group of diffeomorphisms generated by a vector field $X$ as in
the lemma. Then for fixed $q$, the curve $t\mapsto\varphi_t(q)$ satisfies
\begin{equation}\label{L9:eq:ode}
  \frac{d\varphi_t(q)}{dt}=X_{\varphi_t(q)},\qquad \varphi_0(q)=q ,
\end{equation}
since
\[
  \frac{d\varphi_t(q)}{dt}(f)=\lim_{h\to0}\frac{f(\varphi_{t+h}(q))-f(\varphi_t(q))}{h}
  =\lim_{h\to0}\frac{f\bigl(\varphi_h(\varphi_t(q))\bigr)-f(\varphi_t(q))}{h}
  =X_{\varphi_t(q)}(f).
\]
\added{This already gives uniqueness: any group generated by $X$ consists of solutions of
\eqref{L9:eq:ode}, which are unique.}

\emph{Existence.} By existence and uniqueness of solutions of ODEs (Picard--Lindel\"of, applied in
local coordinates), the system \eqref{L9:eq:ode} has a unique solution that depends smoothly on
the initial conditions. Hence for every $q\in M$ there is a neighbourhood $U$ of $q$ and
an $\varepsilon>0$ such that \eqref{L9:eq:ode} has a unique smooth solution for initial points in
$U$ and $\abs{t}<\varepsilon$.

Since $K$ is compact, \added{finitely many such neighbourhoods cover $K$, so} there is a minimal
$\varepsilon>0$ for such a cover of $K$. For $q\notin K$, let $\varphi_t(q)=q$
\added{(this solves \eqref{L9:eq:ode} since $X$ vanishes there)}. Hence there is a unique smooth
solution $\varphi_t(q)$, $\abs{t}<\varepsilon$, for all $q\in M$.
Note that
\[
  \varphi_{t+s}=\varphi_t\circ\varphi_s\qquad\text{for }\abs{t},\abs{s},\abs{t+s}<\varepsilon
\]
\added{(both sides solve \eqref{L9:eq:ode} in $t$ with the same initial value $\varphi_s(q)$)}.
So $\varphi_t$ is a diffeomorphism for $\abs{t}<\varepsilon$ \added{with inverse $\varphi_{-t}$}.

But we can iterate this process for $t\ge\varepsilon$:
\unclear{``by LNP\dots''; step size $\varepsilon$ vs.\ $\varepsilon/2$ unclear}
writing $t=k\,\tfrac{\varepsilon}{2}+r$ with $k=\lfloor 2t/\varepsilon\rfloor$ and
$0\le r<\tfrac{\varepsilon}{2}$, \added{define
$\varphi_t=\varphi_{\varepsilon/2}\circ\cdots\circ\varphi_{\varepsilon/2}\circ\varphi_r$ ($k$ factors),
and $\varphi_{-t}=(\varphi_t)^{-1}$. One checks that this is well defined, smooth, and satisfies the group law.}
\end{proof}

\section{The diffeomorphism theorem}
\label{L9:sec:diffeo}
\orig{title ``Homotopy Theorem'' crossed out, replaced by ``Diffeomorphism''}

For $f\colon M\to\R$ we have the \emph{sublevel sets}
\[
  M^a:=f^{-1}(-\infty,a]=\{p\in M: f(p)\le a\}.
\]
\orig{notes mix $M_a$ and $M^a$}

\begin{theorem}
\label{L9:thm:diffeo}
Let $f\colon M\to\R$ be a smooth function, let $a<b$, and suppose $f^{-1}[a,b]$ is compact and
contains no critical points of $f$. Then $M^a$ is diffeomorphic to $M^b$.
\added{Moreover, $M^a$ is a deformation retract of $M^b$.}
\end{theorem}

\emph{Idea.} Push $f^{-1}(b)$ down to $f^{-1}(a)$ along $-\nabla f$.

\begin{center}
\begin{tikzpicture}[scale=0.9]
  % lower cap region M^a
  \draw[thick] (-2,0) arc (180:360:2 and 1.3);
  \draw[thick] (-2,0) -- (-2,2.5);
  \draw[thick] (2,0) -- (2,2.5);
  \draw[thick] (0,2.5) ellipse (2 and 0.4);
  \draw (-2,0) arc (180:360:2 and 0.4);
  \draw[dashed] (2,0) arc (0:180:2 and 0.4);
  \draw[->,decorate,decoration={snake,amplitude=1pt,segment length=6pt}] (-0.6,2.3) -- (-0.6,-0.2);
  \draw[->,decorate,decoration={snake,amplitude=1pt,segment length=6pt}] (0.6,2.3) -- (0.6,-0.2);
  \node[left] at (-2.1,2.5) {$f^{-1}(b)$};
  \node[left] at (-2.1,0) {$f^{-1}(a)$};
  \node at (0,1.2) {\small $-\nabla f$};
  \draw[|-|] (3,-1.3) -- (3,0) node[midway,right] {$M^a$};
  \draw[|-|] (4,-1.3) -- (4,2.5) node[midway,right] {$M^b$};
\end{tikzpicture}
\end{center}

\begin{proof}
$M$ admits a \added{Riemannian} metric; we need it to define $\langle X,Y\rangle$, the inner product
of tangent vectors. So we can define the gradient $\nabla f$, which is the vector field
characterized by
\[
  \langle X,\nabla f\rangle = Xf
\]
for any vector field $X$; that is, $\langle X,\nabla f\rangle$ is the directional derivative of $f$
along $X$. Note that $\nabla f$ vanishes exactly at the critical points of $f$.
If $\gamma\colon\R\to M$ is a curve with velocity vector $\frac{d\gamma}{dt}$, then
\fixed{notes: $\gamma\colon\R\to\R$}
\[
  \Bigl\langle\frac{d\gamma}{dt},\nabla f\Bigr\rangle=\frac{d(f\circ\gamma)}{dt}.
\]

Define a smooth function $\rho\colon M\to\R$ that is
\[
  \rho=\frac{1}{\langle\nabla f,\nabla f\rangle}\quad\text{on } f^{-1}[a,b],
\]
and $0$ outside of a compact neighbourhood of $f^{-1}[a,b]$
\added{(chosen to contain no critical points, which is possible since the critical set is closed
and $f^{-1}[a,b]$ is compact)}.
Then the vector field $X$ defined as
\[
  X_q=\rho(q)\,(\nabla f)_q
\]
satisfies the hypotheses of Lemma~\ref{L9:lem:flow} (it is smooth and vanishes outside a compact
set), so there is a one-parameter group of diffeomorphisms $\varphi_t\colon M\to M$ generated by $X$.

For fixed $q$, consider $t\mapsto f(\varphi_t(q))$. If $\varphi_t(q)\in f^{-1}[a,b]$, then
\[
  \frac{d f(\varphi_t(q))}{dt}=\Bigl\langle\frac{d\varphi_t(q)}{dt},\nabla f\Bigr\rangle
  =\langle X,\nabla f\rangle = 1 .
\]
So $t\mapsto f(\varphi_t(q))$ is linear with derivative $1$ while $f(\varphi_t(q))$ lies in $[a,b]$.
Therefore $\varphi_{b-a}\colon M\to M$ is a diffeomorphism which carries $M^a$
\added{onto} $M^b$.

\begin{center}
\begin{tikzpicture}[scale=0.8]
  \draw[thick] (-1.5,0) -- (-1.5,3); \draw[thick] (1.5,0) -- (1.5,3);
  \draw[thick] (0,3) ellipse (1.5 and 0.3);
  \draw (0,1.5) ellipse (1.5 and 0.3);
  \draw[thick] (-1.5,0) arc (180:360:1.5 and 0.3);
  \draw[dashed] (1.5,0) arc (0:180:1.5 and 0.3);
  \draw[-{Latex}] (0,-0.3) -- (0,1.2); \draw[-{Latex}] (0,1.2) -- (0,2.7);
  \fill (0,-0.3) circle (1.5pt);
  \node[right] at (1.6,3) {$f=b$}; \node[right] at (1.6,0) {$f=a$};
  \node[left] at (-1.6,1.5) {$\varphi_t$};
  \node at (4.2,1.5) {\small $\varphi_{b-a}\colon M^a\xrightarrow{\ \cong\ }M^b$};
\end{tikzpicture}
\end{center}
\orig{p.\,15: sketch of this flow}

Furthermore, define $r_t\colon M^b\to M^b$, $0\le t\le 1$, by
\fixed{notes: $r_t\colon M_b\to M_a$}
\[
  r_t(q)=\begin{cases}
    q & \text{if } f(q)\le a,\\
    \varphi_{t(a-f(q))}(q) & \text{if } a\le f(q)\le b.
  \end{cases}
\]
Then $r_0=\mathrm{id}$, and $r_1$ is a retraction from $M^b$ to $M^a$; \added{$r_t$ is continuous
since the two formulas agree on $f^{-1}(a)$.} Hence $r_t$ is a deformation retraction
\added{of $M^b$ onto $M^a$.}
\added{(Cf.\ Milnor, \emph{Morse Theory}, Theorem 3.1.)}
\end{proof}

% ============================================================
\chapter{Lecture 10: Configuration spaces and Morse-like results}
\label{L10:chap}
\orig{p.\,16, boxed: ``need to be rewritten, out of order / missing pieces''}

\section{Recap and the handle attachment theorem}
\label{L10:sec:recap}

For smooth functions we have the diffeomorphism result (Theorem~\ref{L9:thm:diffeo}) that says:
if $f\colon M\to\R$, $a<b$, $f^{-1}[a,b]$ is compact, and there are no critical values in $[a,b]$,
then $M^a$ is diffeomorphic to $M^b$.

There is also a handle attachment theorem (not proved here) for smooth $f\colon M\to\R$:

\begin{theorem}[Handle attachment]
\label{L10:thm:handle}
Let $p$ be a non-degenerate critical point of $f$ with index $\lambda$ and $f(p)=c$.
Suppose there is $\varepsilon>0$ such that $f^{-1}[c-\varepsilon,c+\varepsilon]$ is compact and
contains no critical points other than $p$. Then $M^{c+\varepsilon}$ has the homotopy type of
$M^{c-\varepsilon}\cup_g e^\lambda$, \added{that is, $M^{c-\varepsilon}$ with a $\lambda$-cell attached
along a map $g\colon\partial e^\lambda\to M^{c-\varepsilon}$}.
\end{theorem}
\added{(Milnor, \emph{Morse Theory}, Theorem 3.2.)}

\subsection*{Morse-type inequalities}
Let $c_\lambda$ be the number of critical points of index $\lambda$ \added{of a Morse function on a
compact manifold $M$}. The homotopy type changes by attaching various cells of index $\lambda$ as we
``climb'' $M$ via the Morse function. Consequently
\[
  \sum_\lambda(-1)^\lambda c_\lambda=\chi(M),
\]
and in general we have the constraints
\[
  c_\lambda\ \ge\ b_\lambda(M),
\]
where $b_\lambda(M)$ is the $\lambda$-th Betti number of $M$ (the rank of the $\lambda$-th homology group).
\added{(These are the weak Morse inequalities.)}
The key point for us is that the Betti numbers can be computed for many spaces we like.

\section{Betti numbers for the configuration space of points on \texorpdfstring{$\Sph^2$}{S2}}
\label{L10:sec:betti}

Let $B\Conf(N)$ denote the configuration space of $N$ points on $\Sph^2$
\added{modulo rotations, $B\Conf(N)=\Conf(N,\Sph^2)/\SO(3)$.}
\unclear{meaning of ``$B$'' not on page; quotient by $\SO(3)$ inferred}
Via stereographic projection \added{from the first point}, for $N\ge3$ we have a
\added{diffeomorphism}
\[
  B\Conf(N)\ \xrightarrow{\ \cong\ }\ \Conf(N-1,\R^2)/\SO(2).
\]
\begin{addedblock}
Indeed, $\SO(3)$ acts transitively on $\Sph^2$, so we may rotate $u_1$ to the north pole; the
remaining freedom is the stabilizer $\SO(2)$ of the pole, and stereographic projection from the
pole identifies $\Sph^2\setminus\{u_1\}$ with $\R^2$, $\SO(2)$-equivariantly (rotations about the
origin). For $N\ge3$ the actions are free.
\end{addedblock}
That allows us to compute Betti numbers from known Betti numbers \added{via the
fibration $\SO(2)\to\Conf(N-1,\R^2)\to\Conf(N-1,\R^2)/\SO(2)$}: the Poincar\'e polynomial is
\[
  P\bigl(B\Conf(N)\bigr)=\frac{P\bigl(\Conf(N-1,\R^2)\bigr)}{P(\SO(2))},
\]
where
\[
  P\bigl(\Conf(N-1,\R^2)\bigr)=(1+t)(1+2t)\cdots\bigl(1+(N-2)t\bigr),\qquad P(\SO(2))=1+t .
\]
\begin{addedblock}
Here $P(X)(t)=\sum_i b_i(X)\,t^i$. The formula for $\Conf(n,\R^2)$ is due to Arnold (1969). The
division is justified because the $\SO(2)$-bundle is trivial: the map
$u\mapsto (u_2-u_1)/\abs{u_2-u_1}\in \Sph^1$ is $\SO(2)$-equivariant, so
$\Conf(N-1,\R^2)\cong \SO(2)\times\bigl(\Conf(N-1,\R^2)/\SO(2)\bigr)$, and the K\"unneth formula
applies. Hence
\[
  P\bigl(B\Conf(N)\bigr)=(1+2t)(1+3t)\cdots\bigl(1+(N-2)t\bigr).
\]
\end{addedblock}

The picture behind the formula for $\Conf(n,\R^2)$ is that of punctured planes, and the induction
up \added{in $n$: forgetting the last point gives a fibration
$\Conf(n,\R^2)\to\Conf(n-1,\R^2)$ whose fibre is the plane with $n-1$ punctures, with Poincar\'e
polynomial $1+(n-1)t$ (Fadell--Neuwirth)}.

So the Betti numbers $b_\lambda$ of $B\Conf(N)$ are:
\begin{center}
\begin{tabular}{c|cccc}
\toprule
 & $\lambda=0$ & $\lambda=1$ & $\lambda=2$ & $\lambda=3$ \\
\midrule
$N=3$ & $1$ & $0$ & $0$ & $0$\\
$N=4$ & $1$ & $2$ & $0$ & $0$\\
$N=5$ & $1$ & $5$ & $6$ & $0$\\
\added{$N=6$} & \added{$1$} & \added{$9$} & \added{$26$} & \added{$24$}\\
$\vdots$ & $1$ & & & \\
\bottomrule
\end{tabular}
\end{center}
In general $b_0=1$, and the top Betti number is $b_{N-3}=(N-2)!$.

One can use this idea to describe extensions and induction arguments for configuration spaces in
general; for example the $K(\pi,1)$ type of configuration spaces comes from the fact that such spaces
really change nicely \added{(the Fadell--Neuwirth fibrations: inductively, $\Conf(n,\R^2)$ is an
iterated fibration with aspherical fibres, so it is aspherical, with fundamental group the pure
braid group $P_n$)}.
\orig{some crossed-out words here}

\section{Min-type functions}
\label{L10:sec:mintype}
\orig{numbered 10.2 again in the notes}

\begin{definition}
\label{L10:def:mintype}
A \emph{min-type function} is given by a parametrized family of functions
\[
  f(x,p)\colon X\times M\to\R ,
\]
where $X$ is a parameter space, compact\dots; for our purposes, finite and discrete. We require
$f(x,\cdot)$ to be smooth (or smooth enough) for each parameter $x$. The min-type function is then
\[
  F(p)=\added{\min_{x\in X}f(x,p).}
\]
\end{definition}

Our function, the injectivity radius
\[
  \rho\colon\Conf(N)\to\R,\qquad \rho(u)=\min_{\substack{i\ne j\\ u_i,u_j\in u}} d(u_i,u_j),
\]
is of this type \added{(with $X$ the set of pairs $\{i,j\}$ and $f(\{i,j\},u)=d(u_i,u_j)$, the
geodesic distance on $\Sph^2$)}.
\added{(Here ``injectivity radius'' means the minimal distance; the largest radius of disjoint
spherical caps centred at the $u_i$ is $\rho(u)/2$.)}

In general, topological regularity is characterized by the non-emptiness of an intersection of
gradient cones of this type.
\begin{addedblock}
Precisely: for $p\in M$ let $X(p)=\{x\in X: f(x,p)=F(p)\}$ be the active parameters, and let
\[
  C_p=\{v\in T_pM:\ df_x|_p(v)>0\ \text{for all } x\in X(p)\}
\]
be the cone of directions that increase every active function, hence increase $F$ to first order.
Where $C_p\ne\emptyset$, a gradient-like vector field can be built locally, and the flow argument of
Theorem~\ref{L9:thm:diffeo} goes through for superlevel sets $\{F\ge c\}$
(cf.\ Baryshnikov--Bubenik--Kahle, \emph{Min-type Morse theory for configuration spaces of hard
spheres}, IMRN 2014).
\end{addedblock}

\begin{corollary}
\label{L10:cor:topchange}
The topology \added{of the superlevel sets of $F$} only changes at critical points $x$, which only
occur when the gradient cone is empty, i.e.\ when
\[
  \sum_{x} w_x\,df_x\big|_u=0\qquad\text{for some } w_x\ge0 \added{\text{ not all zero, } x\in X(u)}.
\]
\end{corollary}
\added{The equivalence of ``$C_u=\emptyset$'' and the existence of such weights is Gordan's theorem
of the alternative (a finite-dimensional Farkas/separation lemma).}

From our picture for the injectivity radius, this is when increasing some collection
\added{of active distances simultaneously} is not possible.
\unclear{p.\,21 right column only partly legible}

\begin{center}
\begin{tikzpicture}[scale=1]
  % regular: three points in a box, can all move apart
  \begin{scope}
    \draw[rounded corners] (-1.5,-1.1) rectangle (1.5,1.5);
    \coordinate (A) at (90:0.8); \coordinate (B) at (210:0.8); \coordinate (C) at (330:0.8);
    \draw (A) -- (B) -- (C) -- cycle;
    \foreach \P/\a in {A/90,B/210,C/330} { \fill (\P) circle (1.5pt); \draw[->,blue] (\P) -- ++(\a:0.5); }
    \node[below,align=center] at (0,-1.15) {\small regular:\\ \small all can move apart};
  \end{scope}
  % critical: three gradient directions at 120 degrees
  \begin{scope}[xshift=5.5cm,yshift=0.2cm]
    \fill (0,0) circle (1.5pt);
    \foreach \a in {90,210,330} \draw[->,thick] (0,0) -- (\a:1);
    \draw[->,dashed,gray] (0,0) -- (20:0.9) node[right] {$v$?};
    \node[below,align=center] at (0,-1.35) {\small critical: no $v$ in\\ \small the positive (dual) cone};
  \end{scope}
\end{tikzpicture}
\end{center}

\subsection{The contact graph and balanced configurations}
\orig{p.\,22 left column crossed out}

Let $u\in\Conf(N)$ with $\rho(u)\ge\theta$ \added{(we will take $\rho(u)=\theta$)}. There is a
\emph{contact graph}:
\begin{itemize}[nosep]
  \item vertices: the points $u_i\in u$;
  \item edges: $u_iu_j$ whenever $d(u_i,u_j)=\theta$.
\end{itemize}
Edges may be assigned positive weights $w_{ij}$, and if the resulting weighted graph is
\emph{force balanced}, this is equivalent to the cone being empty.
\begin{addedblock}
Write $e_{ij}\in T_{u_i}\Sph^2$ for the unit tangent vector at $u_i$ pointing along the geodesic
towards $u_j$. The differential of $d(u_i,u_j)$ has component $-e_{ij}$ at $u_i$ and $-e_{ji}$ at
$u_j$. So the condition of Corollary~\ref{L10:cor:topchange} reads: there are weights $w_{ij}\ge0$,
not all zero, on the edges of the contact graph with
\[
  \sum_{j\sim i} w_{ij}\,e_{ij}=0\qquad\text{at every vertex } u_i ,
\]
i.e.\ the ``forces'' along the edges balance at each point. The edges with $w_{ij}>0$ form the
\emph{balanced (stressed) subgraph}.
\end{addedblock}

So the only obstacles to topological regularity are exactly the force balanced configurations of
points.

\emph{Examples\dots}
\orig{left blank on p.\,23}

\begin{definition}
\label{L10:def:conftheta}
\added{$\Conf(N,\theta)=\{u\in\Conf(N,\Sph^2):\rho(u)\ge\theta\}$, and $\Conf(N,0):=\Conf(N,\Sph^2)$.}
\end{definition}

\begin{theorem}
\label{L10:thm:smalltheta}
Suppose $N\ge3$.
\unclear{``$N\ge3$'' or ``$N>3$''}
Then for $0<\theta<\pi/N$, $\Conf(N,\theta)$ is a strong deformation retract of, and hence homotopy
equivalent to, $\Conf(N,0)$.
\end{theorem}

\begin{proof}[Proof (sketch)]
First, there are no balanced configurations \added{with $\rho(u)=\theta'$ for $0<\theta'<\pi/N$};
then push each point apart \added{using the flow of a gradient-like vector field for $\rho$}\dots
\begin{addedblock}
\emph{No balanced configurations.} (i) A balanced subgraph is not contained in any open hemisphere:
if all its vertices lie in $\{y:\langle y,e\rangle>0\}$, take a vertex $v$ of the subgraph minimizing
$\langle v,e\rangle$. For each neighbour $u$, the vector $e_{vu}$ is a positive multiple of
$u-\langle u,v\rangle v$, whose $e$-component is $\langle u,e\rangle-\langle u,v\rangle\langle v,e\rangle
\ge\langle v,e\rangle(1-\langle u,v\rangle)>0$ (using $\langle u,e\rangle\ge\langle v,e\rangle$ and $\langle u,v\rangle<1$). So every force at $v$ has a positive $e$-component,
and they cannot sum to zero.
(ii) Each connected component of the balanced subgraph is itself balanced. If a component has $k$
vertices, a spanning tree has $k-1$ edges of length $\theta'$, total length $L=(k-1)\theta'$. A tree
has a centre $c$ within distance $L/2$ (along the tree) of all its points, so all vertices lie in the
spherical ball of radius $L/2$ about $c$. By (i), $L/2\ge\pi/2$, so $(k-1)\theta'\ge\pi$, so
$\theta'\ge\pi/(k-1)\ge\pi/(N-1)>\pi/N$.
(So this crude argument even gives the theorem for $\theta<\pi/(N-1)$. The sharp bound is
$\theta<2\pi/N$: a balanced graph has every valence $\ge2$, hence at least as many edges as
vertices, and Kusner--Kusner--Lagarias--Shlosman (Thm.~4.14, cited below) show its total length is at least
$2\pi$, with equality only for the $N$ points equally spaced on a great circle. Compare the
statement at the start of Lecture~11.)

\emph{Deformation.} So $\rho$ has no critical values in $(0,\theta]$, and the min-type version of
Theorem~\ref{L9:thm:diffeo} lets us push $\{\rho\ge\eta\}$ onto $\{\rho\ge\theta\}$ for every
$0<\eta<\theta$; some care is needed near $\rho=0$, where $\Conf(N,0)$ is not compact. (Cf.\
Kusner--Kusner--Lagarias--Shlosman, \emph{Configuration spaces of equal spheres touching a given
sphere: the twelve spheres problem}, arXiv:1611.10297 (2016); in \emph{New Trends in Intuitive
Geometry}, Bolyai Soc.\ Math.\ Stud.\ 27, Springer, 2018.)
\end{addedblock}
\end{proof}

For $N=4$ \dots, $\Conf(N,\theta)$ for small $\theta$ is therefore path connected.
\unclear{``$N=4$'' or ``$N>4$''}
\added{(Indeed $\Conf(N,\Sph^2)$ is connected, since the diagonals removed have real codimension~2.)}

\begin{example}[$N=4$]
\label{L10:ex:N4}
\orig{left blank in the notes}
\added{Balanced subgraphs need every vertex of valence $\ge2$, and a vertex of valence $2$ must have
its two edges in opposite directions, so it lies on a great circle through its two neighbours. One
finds: the square on a great circle ($\theta=\pi/2$), and the regular tetrahedron
($\theta=\arccos(-1/3)\approx109.47^\circ$, the maximum of $\rho$). An equilateral triangle on a great
circle ($\theta=2\pi/3$) plus a free fourth point is \emph{not} realizable, since every point of
$\Sph^2$ is within $\pi/2$ of one of the three. This is the editor's check; compare the question marks on p.\,29.}
\end{example}

\subsection{Outlook}
\orig{pp.\,24--25: headings only}
\begin{itemize}[nosep]
  \item Strata for Tammes maxima: table\dots
  \added{(The Tammes problem asks for $\max\rho$ over $\Conf(N,\Sph^2)$; it is solved for
  $N\le14$ and $N=24$, with $N=13,14$ due to Musin--Tarasov, 2012/2015.)}
  \item An algorithm for searching \added{for critical configurations} by evolution of force
  configurations\dots
  \item Local maxima; rings\dots
  \item Rigidity?, critical conditions and weights\dots
  \item Candidate graph sorting\dots
\end{itemize}

\section{Topological regularity revisited}
\label{L10:sec:revisit}
\orig{pp.\,26--29 seem to re-explain \S\ref{L10:sec:mintype}; possibly a later lecture}

What is this? This is a force balanced \added{configuration}.
\unclear{p.\,26 has only these two lines}

That is, for the injectivity radius, \added{the gradient cone} is the intersection of the convex
cones of motions that increase the distance between each \added{active} pair of points.

\emph{Picture: three points.} The configuration space of $3$ points on $\Sph^2$ is $6$-dimensional.
Fixing one point leaves a $4$-dimensional space; fixing also the rotation angle $\psi$ about it
leaves a $3$-dimensional space, which is \added{locally} parametrized by the distances
$d_{12},d_{13},d_{23}$. The cone for $d_{12}$ is the upper half space $\{d_{12}\text{ increasing}\}$,
etc.
\orig{angle called $\theta$ in the notes}

\begin{center}
\begin{tikzpicture}[scale=1]
  \begin{scope}
    \draw (0,0) circle (1);
    \fill (-0.3,0.15) circle (1.2pt); \fill (0.35,-0.05) circle (1.2pt); \fill (0.1,0.5) circle (1.2pt);
    \node at (0,-1.3) {\small 3 points on $\Sph^2$};
  \end{scope}
  \begin{scope}[xshift=3.5cm]
    \draw[dotted] (0,-1.3) -- (0,1.3);
    \coordinate (P3) at (0,-0.6); \coordinate (P1) at (-0.8,0.6); \coordinate (P2) at (0.8,0.6);
    \draw (P1) -- node[above]{\small $d_{12}$} (P2) -- node[right]{\small $d_{23}$} (P3)
          -- node[left]{\small $d_{13}$} cycle;
    \fill (P1) circle (1.3pt); \fill (P2) circle (1.3pt); \fill (P3) circle (2pt);
    \node[below right] at (P3) {\small fixed};
    \draw (0,-0.2) arc (90:55:0.4); \node at (0.22,-0.05) {\small $\psi$};
  \end{scope}
  \begin{scope}[xshift=7.5cm]
    \draw[->] (-1.3,0) -- (1.3,0) node[right]{\small $d_{13}$};
    \draw[->] (0,-1.1) -- (0,1.3) node[above]{\small $d_{12}$};
    \draw[->] (-0.9,-0.9) -- (0.9,0.9) node[above right]{\small $d_{23}$};
    \node at (0,-1.4) {\small coordinates near $u$};
  \end{scope}
\end{tikzpicture}
\end{center}
\added{When the three active differentials are linearly independent, the three half-spaces meet in
an open octant, so the cone is non-empty and the value is regular.}

So:
\[
  \text{$\theta$ is a topologically regular value}\iff
  \begin{minipage}[c]{0.5\linewidth}
  neighbourhoods of the value have superlevel sets that deformation retract to each other.
  \end{minipage}
\]
\added{(For $\theta'<\theta$ near $\theta$, $\{\rho\ge\theta\}\hookrightarrow\{\rho\ge\theta'\}$ is
a deformation retract.)} So critical values are where the topology changes\dots

For topologically regular values, this deformation retraction is defined by the flow of some vector
field defined by $f$, the gradient flow as in the last lecture (Theorem~\ref{L9:thm:diffeo}).
So the topology can only change at an obstruction to gradient flow.
For a min-type function, this obstruction is given by a convex cone, namely the intersection
of all variational cones that increase $f$ at the parameters $x\in X$.

\begin{exercise}
\label{L10:ex:classify}
Classify all distinct realizable balanced configurations of $N$ points on a sphere.
\end{exercise}

First, list the graphs. For $3$ points, the candidate contact graphs are:
\begin{center}
\begin{tikzpicture}[scale=0.7,every node/.style={circle,fill,inner sep=1.2pt}]
  \begin{scope}
    \node (a) at (0,0) {}; \node (b) at (1,0) {}; \node (c) at (2,0) {};
  \end{scope}
  \begin{scope}[xshift=3.5cm]
    \node (a) at (0,0) {}; \node (b) at (1,0) {}; \node (c) at (2,0) {}; \draw (a)--(b);
  \end{scope}
  \begin{scope}[xshift=7cm]
    \node (a) at (0,0) {}; \node (b) at (1,0) {}; \node (c) at (2,0) {}; \draw (a)--(b)--(c);
  \end{scope}
  \begin{scope}[xshift=10.5cm]
    \node (a) at (0,0) {}; \node (b) at (1,0) {}; \node (c) at (2,0) {}; \draw (a)--(b)--(c);
    \draw (a) to[bend right=40] (c);
  \end{scope}
\end{tikzpicture}
\end{center}
For example, to be force balanced every vertex needs valence $\ge2$
\unclear{reads ``valence $=2$''}
\added{(and a vertex of valence $2$ needs its two edges opposite, i.e.\ along a great circle). For
$N=3$ only the triangle survives, realized as three points at mutual distance $2\pi/3$ on a great
circle.}

Some graphs on $4$ points:
\begin{center}
\begin{tikzpicture}[scale=0.7,every node/.style={circle,fill,inner sep=1.2pt}]
  \begin{scope}
    \foreach \i in {0,1,2,3} \node (p\i) at (\i*0.7,0) {};
  \end{scope}
  \begin{scope}[xshift=3.5cm]
    \foreach \i in {0,1,2,3} \node (q\i) at (\i*0.7,0) {};
    \draw (q0)--(q1)--(q2)--(q3); \draw (q0) to[bend right=40] (q3);
  \end{scope}
  \begin{scope}[xshift=7cm]
    \node (a) at (0,0) {}; \node (b) at (1,0) {}; \node (c) at (0.5,0.85) {}; \node (d) at (1.7,0.4) {};
    \draw (a)--(b)--(c)--(a);
  \end{scope}
  \begin{scope}[xshift=10cm]
    \node (t) at (0.7,0.8) {}; \node (l) at (0,0) {}; \node (r) at (1.4,0) {}; \node (bt) at (0.7,-0.8) {};
    \draw (t)--(l)--(bt)--(r)--(t); \draw (l)--(r);
  \end{scope}
  \begin{scope}[xshift=13cm]
    \node (k1) at (0,0) {}; \node (k2) at (1.4,0) {}; \node (k3) at (0.2,1.0) {}; \node (k4) at (0.7,0.35) {};
    \draw (k1)--(k2)--(k3)--(k1); \draw (k4)--(k1); \draw (k4)--(k2); \draw (k4)--(k3);
  \end{scope}
\end{tikzpicture}
\end{center}
\orig{``?'' by the triangle${}+{}$point and the rhombus; ``opposite'', ``great circles''. One more sketch, apparently $K_4$, drawn at right.}
For the rhombus with a diagonal, the two valence-2 vertices need opposite edges, which forces great
circles; whether such configurations are realizable is a question of geometry\dots
\added{(It is not: the two valence-3 vertices would have to be at distance $\theta$ and at distance
$2\theta$ along a great circle, forcing $\theta=2\pi/3$, and then the two valence-2 vertices would
coincide.)}
\orig{p.\,29 ends with a sketch of two great circles}
% ===== end chapters/L09.tex =====
% ===== begin chapters/L11.tex =====
% MAT670 -- Lecture XI (scans MAT670_XI, pp. 1--14)
\chapter{Lecture 11: Critical points, jamming and entropy}
\label{L11:ch}

\section{Critical points and jamming}
\label{L11:sec:critical}

\added{Throughout, a configuration is an ordered $N$-tuple $X=(x_1,\dots,x_N)$ of points on the sphere $\Sph^2$, and $d$ denotes the geodesic (angular) distance. Let $f(X)=\min_{i<j}d(x_i,x_j)$ and write $\Conf(N;\theta)=\{X: f(X)\ge\theta\}$. A configuration in $\Conf(N;\theta)$ is the same thing as a packing of $N$ spherical caps of angular radius $\rho=\theta/2$.}

Last time we stated the following theorem, which characterizes the regular values of $f$ in terms of the geometry of the contact graph.

\begin{definition}[Contact graph]\label{L11:def:contact}
Let $X\in\Conf(N;\theta)$ with $f(X)=\theta$, where $\theta=2\rho$. The \emph{contact graph} of $X$ is the set of edges on $\Sph^2$ \added{(geodesic arcs)} joining those points $x_i,x_j$ that are at distance exactly $\theta$.
\end{definition}
\orig{crossed out: ``geodesic skeleton of'', ``struts''}

\begin{addedblock}
For an edge $\{i,j\}$ of the contact graph let $u_{ij}\in T_{x_i}\Sph^2$ be the unit tangent vector at $x_i$ pointing along the geodesic towards $x_j$ (well defined since $0<\theta<\pi$). A tangent vector $v=(v_1,\dots,v_N)$, $v_i\in T_{x_i}\Sph^2$, changes the length of the edge at rate
\[
\frac{d}{dt}\Big|_{t=0} d(x_i+tv_i,\,x_j+tv_j) \;=\; -\ip{u_{ij}}{v_i}-\ip{u_{ji}}{v_j}.
\]
We say $v$ lies in the \emph{tangent cone} (of increasing directions) at $X$ if this rate is strictly positive for every contact edge, and that $\theta$ is a \emph{regular value} of $f$ if every $X$ with $f(X)=\theta$ admits such a $v$; otherwise $\theta$ is a \emph{critical value} and the offending $X$ are \emph{critical configurations}.
\end{addedblock}

\begin{theorem}\label{L11:thm:forcebalance}
A configuration $X$ with $f(X)=\theta$ has no element in its tangent cone if and only if there exist weights $w_{ij}=w_{ji}\ge0$ on the contact edges, not all zero, such that the contact graph is \emph{force balanced}:
\[
\sum_{j:\,\{i,j\}\text{ contact}} w_{ij}\,u_{ij} \;=\; 0 \qquad\text{for every vertex } i.
\]
In particular, $\theta$ is a regular value if and only if no configuration at level $\theta$ admits such a balanced (non-negative, not all zero) system of weights.
\end{theorem}
\orig{``(non-negative, not all 0)'' replaces crossed-out ``(positive)''}

\begin{addedblock}
\emph{Proof sketch.} Let $A$ be the linear map $v\mapsto\big(-\ip{u_{ij}}{v_i}-\ip{u_{ji}}{v_j}\big)_{\{i,j\}}$ from $\bigoplus_i T_{x_i}\Sph^2$ to $\R^{E}$, $E$ the set of contact edges. By Gordan's alternative (a form of Farkas' lemma), either there is $v$ with $Av>0$ componentwise, or there is $w\in\R^E$, $w\ge0$, $w\ne0$ with $A^{T}w=0$. Computing $A^Tw$ vertex by vertex gives exactly $-\sum_j w_{ij}u_{ij}=0$. Together with the deformation theorem from last lecture (if $[a,b]$ contains no critical values, then $\Conf(N;b)$ is a deformation retract of $\Conf(N;a)$), this is the basic tool for computing the topology of $\Conf(N;\theta)$ as $\theta$ increases.
\end{addedblock}

The simplest application is to points on a circle.

\begin{theorem}\label{L11:thm:circle}
\added{For $N$ points on the circle $\Sph^1$,} $\Conf(N;0)\simeq\Conf(N;2\pi/N)$.
\end{theorem}
\unclear{Blotted symbol before $2\pi/N$; read here as the circle case. Possibly ``$<$'': the $\Sph^2$ theorem, see remark.}
\begin{addedblock}
\emph{Alternative reading.} The blot may be ``$\theta<$'', i.e.\ the statement for $\Sph^2$:
$\Conf(N;\theta)$ is a strong deformation retract of $\Conf(N;0)$ for $0\le\theta<2\pi/N$. This is
Theorem~4.14 of Kusner--Kusner--Lagarias--Shlosman (cited in \S\ref{L11:sec:examples}). It sharpens
the bound $\pi/N$ of Lecture~10, since a balanced contact graph has total length at least $2\pi$,
and the smallest critical value is attained only by $N$ equally spaced points on a great circle. The
proof sketch below covers only the circle.
\end{addedblock}

\begin{addedblock}
\emph{Proof sketch.} Here $\Conf(N;0)$ means the usual ordered configuration space (distinct points). On $\Sph^1$ the vectors $u_{ij}$ at $x_i$ are $\pm$ the unit tangent, so force balance at a vertex forces it to have one contact on each side with equal weights. Following the weights along a chain of contacts, a chain with an endpoint is unbalanced at that endpoint; hence the contact graph of a critical configuration is the whole $N$-cycle, i.e.\ $X$ is a regular $N$-gon and $\theta=2\pi/N$. Thus there are no critical values in $(0,2\pi/N)$, so $\Conf(N;\theta)\simeq\Conf(N;0)$ for $\theta<2\pi/N$. At $\theta=2\pi/N$ itself the space consists of the labelled regular $N$-gons, a disjoint union of $(N-1)!$ circles (one for each cyclic order of the labels); and $\Conf(N;0)$ has the same homotopy type, each of its $(N-1)!$ components being a circle times an open simplex.
\end{addedblock}

This also gives an algorithm for exploring critical and maximal configurations: \emph{$\rho$-evolution}. Starting from any configuration, increase $\rho$ by moving the points in a direction of the tangent cone; update the contact graph whenever new contacts form, and continue until no increasing direction exists, i.e.\ until the contact graph becomes force balanced.
\unclear{Page ends with ``update\ldots'' and ``conflicting\ldots'' (or ``conf-flip''?); not expanded.}
\added{The configuration reached is then a critical configuration, generically a local maximum of $f$ (a \emph{jammed} configuration).}

\section{Examples with pictures}
\label{L11:sec:examples}

\orig{p.~2: ``$N=3$, $N=4$: full picture'' -- pictures were presumably drawn on the board.}
For $N=3$ and $N=4$ one can give the full picture of the critical points.
\begin{addedblock}
For $N=3$: a vertex with a single contact cannot be balanced, so every vertex of a critical contact graph has degree $2$, and the two tangent vectors at it must be opposite; thus the three points lie on a great circle, equally spaced. The only critical value is the maximum $\theta=2\pi/3$, and $\Conf(3;\theta)\simeq\Conf(3;0)$ for $\theta<2\pi/3$.
For $N=4$: the maximum is the regular tetrahedron, $\theta=\arccos(-1/3)\approx109.47^\circ$. The square inscribed in a great circle ($\theta=\pi/2$) is also critical (each vertex has two opposite contacts), but it is not a local maximum: pushing alternate vertices slightly north and south increases all four contact distances to second order, deforming the square towards the tetrahedron.
\end{addedblock}

For larger $N$:
\begin{itemize}
\item For various $N$ the maximum is the solution of the \emph{Tammes problem}. \added{(The Tammes problem is solved for $N\le 14$ and $N=24$; e.g.\ $N=13$ and $N=14$ by Musin--Tarasov, $N=24$ by Robinson.)}
\item $N=10$: local maxima? \unclear{``$N=10$ -- local max?'': whether there are non-global local maxima, presumably.}
\item $N=12$: local maxima? Higher critical points. The jitterbug?
\end{itemize}
\begin{addedblock}
For $N=12$ the global maximum is the icosahedron ($\theta=\arctan 2\approx63.43^\circ$). The cuboctahedron (the kissing configuration of the FCC lattice, $\theta=60^\circ$) is a critical configuration that is not a local maximum; Buckminster Fuller's ``jitterbug'' motion deforms the cuboctahedron into the icosahedron. (See R.~Kusner, W.~Kusner, J.~Lagarias, S.~Shlosman, \emph{Configuration spaces of equal spheres touching a given sphere: the twelve spheres problem}, arXiv:1611.10297 (2016); in \emph{New Trends in Intuitive Geometry}, Bolyai Soc.\ Math.\ Stud.\ 27, Springer, 2018. That paper shows the cuboctahedral (FCC) configuration is critical. That it is not a local maximum follows from the jitterbug: along it the triangle edges lengthen from $60^\circ$ to $\arctan 2$, while the square diagonals shrink from $90^\circ$ to $\arctan 2$.)
\end{addedblock}

\section{Local maxima and Edwards entropy}
\label{L11:sec:edwards}

\begin{definition}[Edwards' ansatz]\label{L11:def:edwards}
All jammed states of equal volume have equal probability.
\end{definition}
\orig{``Edw.'' crossed out and replaced by ``Ansatz''}
\begin{addedblock}
This is the hypothesis of S.~F.~Edwards and R.~B.~S.~Oakeshott (\emph{Theory of powders}, Physica A, 1989). The \emph{Edwards entropy} at volume $V$ is the logarithm of the number (or measure) of jammed states of volume $V$; it plays the role that the Boltzmann entropy plays for the microcanonical ensemble.
\end{addedblock}

This is related to the evolution of the $\rho$-flow: \added{if configurations are chosen at random and then evolved by $\rho$-evolution, the probability of a given jammed state is the measure of its basin of attraction, and Edwards' ansatz says these basins are (for equal volumes) equally large.} It is also related to the model \ldots\ \unclear{``Also the model\ldots''; crossed out: ``ex.\ annealed''} It is not always true\,? \added{Numerical studies of small disc packings have found that jammed states are in general \emph{not} equally likely, and that the frequencies depend on the preparation protocol.}
\added{\emph{Update:} later work (Martiniani, Schrenk, Ramola, Chakraborty, Frenkel, \emph{Nature Physics}, 2017) gave numerical evidence that the ansatz does hold for soft-sphere packings exactly at the jamming threshold.}

\subsection{Boltzmann measure: the discrete setting}
\label{L11:sec:boltzmann}

Consider $N$ particles, each in one of $k$ states. Each state has an energy $E_i$, $i\in\{1,\dots,k\}$. Let $n_i$ be the number of particles in state $i$. The system is closed, so
\[
E=\sum_{i=1}^k E_i\,n_i,\qquad N=\sum_{i=1}^k n_i
\]
\orig{he writes both $n_i$ and $N_i$; unified to $n_i$}
are fixed. The number of ways to place the particles with occupation numbers $\{n_1,\dots,n_k\}$ is the multinomial coefficient
\[
W=\frac{N!}{n_1!\,n_2!\cdots n_k!}.
\]
By the law of large numbers, the optimal partition (the one maximizing $W$) dominates for $N$ large.
\added{(More precisely: $W$, as a function of the frequencies $n_i/N$, is $e^{N H+o(N)}$ with $H=-\sum (n_i/N)\ln(n_i/N)$, so frequencies whose entropy is smaller than the maximum by a fixed amount are exponentially rare.)}
So we determine the partition that maximizes $W$; since $\log$ is monotone increasing we may replace $W$ by $\log W$. By Stirling's formula,
\[
\log W=(N\ln N-N)-\sum_{i=1}^k\,(n_i\ln n_i-n_i)\;+\;\text{error}\quad\added{(\text{error }=O(k\ln N))}.
\]
We maximize subject to the two constraints with Lagrange multipliers $\alpha,\beta$:
\[
f=\ln W+\alpha\Big(N-\sum_{i=1}^k n_i\Big)+\beta\Big(E-\sum_{i=1}^k E_i n_i\Big).
\]
Since $\frac{\partial}{\partial n}(n\ln n-n)=\ln n$,
\[
0=\frac{\partial f}{\partial n_i}=-\ln n_i-(\alpha+\beta E_i),
\qquad\text{so}\qquad n_i=e^{-\alpha-\beta E_i}.
\]
\added{Summing over $i$ gives $N=e^{-\alpha}Z$, which eliminates $\alpha$:}
\[
\frac{n_i}{N}=\frac{e^{-\beta E_i}}{\sum_{j=1}^k e^{-\beta E_j}},\qquad
Z=\sum_{j=1}^k e^{-\beta E_j}\quad\added{\text{(the \emph{partition function})}}.
\]

\begin{definition}\label{L11:def:boltzmann}
Here $\beta=\dfrac1T$ \ (or $\dfrac{1}{k_BT}$ \added{with Boltzmann's constant $k_B$}) is the inverse temperature, and
\[
F=-T\ln Z
\]
is the \emph{free energy}. The \emph{Boltzmann measure} is
\[
\Pr(\text{state } i)=\frac{e^{-\beta E_i}}{Z}.
\]
\end{definition}

The mean energy per particle is
\[
\bar E=\frac EN=\sum_{i=1}^k E_i\frac{n_i}{N}
=\sum_{i=1}^k\frac{E_ie^{-\beta E_i}}{Z}
=\frac{-\frac{\partial}{\partial\beta}Z}{Z}
=-\frac{\partial}{\partial\beta}\ln Z.
\]
\fixed{page has $e^{\beta E_i}$ and final ``$-\partial_\beta Z$''; signs and $\ln$ corrected}

\begin{lemma}\label{L11:lem:concave}
$-\ln Z$ is concave \added{as a function of $\beta$}. Equivalently, $\bar E=-\frac{\partial}{\partial\beta}\ln Z$ is monotone \added{(non-increasing in $\beta$, i.e.\ non-decreasing in $T$)}.
\end{lemma}
\begin{proof}[Proof idea]
We must show $\dfrac{\partial^2}{\partial\beta^2}\ln Z\ge0$. Why? Write it as the variance of $E$, which is $\ge0$.
\begin{addedblock}
Indeed, with $\langle g\rangle=\sum_i g(E_i)e^{-\beta E_i}/Z$,
\[
\frac{\partial^2}{\partial\beta^2}\ln Z=\frac{Z''}{Z}-\Big(\frac{Z'}{Z}\Big)^2=\langle E^2\rangle-\langle E\rangle^2=\operatorname{Var}(E)\ge0,
\]
with equality only if all $E_i$ are equal.
\end{addedblock}
\end{proof}
\fixed{page has ``$>0$''; it is $\ge 0$}

Thus $-\dfrac1\beta\ln Z$ is the free energy.
\fixed{page: ``$\frac1\beta\ln Z$''; sign restored as in $F=-T\ln Z$}
As $T\to\infty$ (i.e.\ $\beta\to0$), $\Pr(i)\to\dfrac1k$: all states become equally likely. \added{As $T\to0$ the measure concentrates on the states of minimal energy.}

\begin{remark}
\added{For hard particles the energy of a configuration is $0$ if the particles do not overlap and $+\infty$ if they do. The Boltzmann weight $e^{-\beta E}$ is then the indicator of the admissible configurations, so the Boltzmann measure is normalized volume on the configuration space, at every temperature. This is why volumes of configuration spaces appear as partition functions below.}
\end{remark}

\subsection{The hard disc model in an interval}
\label{L11:sec:hardrods}

Place hard discs \added{of diameter~$1$ (that is, hard rods of length~$1$)} in the interval $[0,N]$, and let $a_1,a_2,a_3,\dots$ be the gaps \added{in front of the first, second, third, \ldots\ disc}.

\begin{center}
\begin{tikzpicture}[scale=0.9]
\draw[thick] (0,0) -- (8,0);
\draw (0,-0.15)--(0,0.15) node[below=6pt] {$0$};
\draw (8,-0.15)--(8,0.15) node[below=6pt] {$N$};
\foreach \x in {1.2,3.6,5.3} { \fill[gray!50] (\x,-0.12) rectangle (\x+1,0.12); \draw (\x,-0.12) rectangle (\x+1,0.12); }
\draw[<->] (0,0.35) -- node[above] {$a_1$} (1.2,0.35);
\draw[<->] (2.2,0.35) -- node[above] {$a_2$} (3.6,0.35);
\draw[<->] (4.6,0.35) -- node[above] {$a_3$} (5.3,0.35);
\end{tikzpicture}
\end{center}

With three discs the configuration space is
\[
\Omega_{N,3}=\{(a_1,a_2,a_3):\ 0\le a_1,a_2,a_3,\ \ a_1+a_2+a_3\le N-3\},
\]
a simplex \orig{p.~7: sketch of the simplex with axes $a_1,a_2,a_3$} with
\[
\vol\Omega_{N,3}=\frac{(N-3)^3}{3!}.
\]
In general, \added{for $k\le N$ discs,}
\[
\vol(\Omega_{N,k})=\frac{(N-k)^k}{k!}.
\]
\added{(The gaps describe the discs up to relabelling; for labelled discs multiply by $k!$.)}

Now give each particle an energy $E_0$ and define the \emph{fugacity}
\[
z=e^{-\beta E_0}.
\]
\added{(In the usual notation $z=e^{\beta\mu}$ with chemical potential $\mu$; here $\mu=-E_0$.)} The \added{grand canonical} partition function is
\[
Z_N(z)=\sum_{k=0}^N z^k\,\vol(\Omega_{N,k})=\sum_{k=0}^N\frac{(N-k)^k}{k!}\,z^k.
\]
\added{This is the classical \emph{Tonks gas} (L.~Tonks, 1936), one of the few exactly solvable models of hard particles.}

\begin{exercise}\label{L11:ex:tonks}
\added{Using Stirling's formula on the largest term, show that $\frac1N\ln Z_N(z)\to P$, where $P>0$ solves $z=Pe^{P}$, and that the dominant number of discs satisfies $k/N\to P/(1+P)$. ($P$ is $\beta$ times the pressure; note that the density $k/N$ increases from $0$ to $1$ as $z$ goes from $0$ to $\infty$.)}
\end{exercise}

\section{A quasi-one-dimensional model with open ends}
\label{L11:sec:channel}

\orig{crossed out: ``Model''}
We study discs in a narrow channel, where a combinatorial analysis is possible. Take discs of radius $r=1$ in a channel whose width $w$ satisfies
\[
2<w<2+\sqrt3 .
\]
\begin{addedblock}
The centres then lie in a strip of width $h=w-2\in(0,\sqrt3)$. Discs cannot pass each other ($h<2$). If two consecutive discs touch while lying against opposite walls, their centres are horizontally $L=\sqrt{4-h^2}$ apart, and $1<L<2$; the condition $h<\sqrt3$ is exactly $L>1$, which guarantees that each disc can touch only its two neighbours in the row (discs $i$ and $i+2$ are always more than $2$ apart). Consecutive discs touching along the same wall are horizontally $2$ apart.
\end{addedblock}

\begin{figure}[ht]
\centering
\begin{tikzpicture}[scale=0.33]
% channel width 3.2, h = 1.2, L = 1.6
\draw[thick] (-1.5,0) -- (15.1,0); \draw[thick] (-1.5,3.2) -- (15.1,3.2);
\foreach \x/\y in {0/1,1.6/2.2,3.2/1,5.2/1,6.8/2.2,8.4/1,10.4/1,12.0/2.2,13.6/1} { \draw (\x,\y) circle (1); }
\node[right] at (15.3,1.6) {\small jammed: locally maximal density};
\end{tikzpicture}

\vspace{2ex}
\begin{tikzpicture}[scale=0.33]
\draw[thick] (-1.5,0) -- (8.7,0); \draw[thick] (-1.5,3.2) -- (8.7,3.2);
\foreach \x/\y in {0/1,1.6/2.2,3.6/2.2,5.6/2.2,7.2/1} { \draw (\x,\y) circle (1); }
\draw[->,thick] (3.6,2.2) -- (3.6,0.6);
\node[right] at (9,1.6) {\small critical saddle};
\draw[<->] (19,0) -- node[right] {\small $w$} (19,3.2);
\end{tikzpicture}
\caption{Discs in a narrow channel ($h=1.2$, $L=1.6$). Top: a jammed configuration, contact word $00100100$. Bottom: three consecutive discs on the same wall (word $011\,0$); the middle one is force balanced but can be pushed into the gap, so this is a critical point but not a local maximum.}
\label{L11:fig:channel}
\end{figure}
\orig{p.~9: three sketches; third shows a narrow channel with discs in a staircase}

Let the length of the configuration be $t$. \added{(Discs are confined only by the walls; the ends are open, so ``jammed'' means that $t$ is locally minimal for the given number of discs.)}

\subsection{\texorpdfstring{Combinatorics: how many maxima for given $t$?}{Combinatorics: how many maxima for given t?}}

The length is determined by the number of adjacent pairs of discs at opposite walls \added{versus at the same wall}. Record each contact between consecutive discs by a letter:

\begin{center}
\begin{tikzpicture}[scale=0.4]
\draw (0,0) circle (1); \draw (2,0) circle (1); \draw[->] (0,0)--(2,0);
\node at (1,-1.9) {$1$: same wall};
\draw (7,0) circle (1); \draw (8.6,-1.2) circle (1); \draw[->] (7,0)--(8.6,-1.2);
\node at (7.8,-3.1) {$0$: opposite walls};
\end{tikzpicture}
\end{center}

\added{Two consecutive $1$'s are impossible in a jammed configuration: the middle disc of three discs along one wall can move away from the wall (Figure~\ref{L11:fig:channel}, bottom).} So counting jammed configurations amounts to counting binary words with no two adjacent $1$'s. This is equivalent to a domino tiling problem, or a ``stars and bars'' count: such words are sequences formed of the blocks
\[
0=\square,\qquad 01=\square\!\square\ (\text{a domino}).
\]
\added{(Precisely: prefixing a $0$, words of length $\ell$ with no $11$ correspond bijectively to tilings of a $1\times(\ell+1)$ strip by squares and dominoes.)}

This satisfies the Fibonacci recursion \added{(we count configurations up to the up/down reflection of the channel)}: a tiling of length $t$ is a tiling of length $t-1$ followed by a square, or a tiling of length $t-2$ followed by a domino; there are $1,2,\dots$ tilings of lengths $1,2,\dots$
\unclear{``(not up/down symmetric)'' -- reading as: words count configurations modulo reflection.}
\orig{crossed out: ``\# words ending in $0$ of length $\ell$''}
\added{Hence the number of tilings of length $m$ is the Fibonacci number $F_{m+1}$, and the number of binary words of length $\ell$ with no $11$ is $F_{\ell+2}$ ($F_1=F_2=1$).}

\begin{center}
\begin{tikzpicture}[scale=0.45]
\draw (0,0) rectangle (6,1); \draw (6,0) rectangle (7,1); \node[left] at (-0.2,0.5) {$t-1$};
\draw (0,-2) rectangle (5,-1); \draw (5,-2) rectangle (7,-1); \node[left] at (-0.2,-1.5) {$t-2$};
\end{tikzpicture}
\end{center}

Given exactly $k$ ones in a word of length $\ell$, we are inserting $k$ ones into a sequence of $\ell-k$ zeros,
\[
{}_{\wedge}0_{\wedge}0_{\wedge}0_{\wedge}\cdots{}_{\wedge}0_{\wedge},
\]
at most one into each of the $\ell-k+1$ slots. So for each $k$ there are $\dbinom{\ell-k+1}{k}$ choices of order. \added{(Summing over $k$ gives $\sum_k\binom{\ell-k+1}{k}=F_{\ell+2}$.)}

So if the length of a ``$1$'' is $2$ and that of a ``$0$'' is $L$, the configuration with word of length $\ell$ \added{(i.e.\ $\ell+1$ discs)} containing $k$ ones has length
\[
t=2k+(\ell-k)L+2,
\]
the $+2$ coming from the two ends.

\begin{remark}
\unclear{Check: with open ends, a $1$ in the first or last slot seems not locally length-minimizing.}
\added{With truly open ends, an end disc lying against the same wall as its neighbour can roll around that neighbour towards the opposite wall, shortening $t$. If such end configurations are excluded, the count becomes $\binom{\ell-k-1}{k}$ instead of $\binom{\ell-k+1}{k}$; the asymptotics below are unchanged.}
\end{remark}

We can also compute the maximum of $\binom{n-k+1}{k}$ over $k$ (with $n=\ell$): as $n\to\infty$ the maximizing $k$ satisfies
\[
\frac kn\longrightarrow\frac5{10}-\frac{\sqrt5}{10}=\frac12-\frac{\sqrt5}{10}\approx0.276 .
\]
\begin{addedblock}
Indeed, with $x=k/n$, Stirling gives $\frac1n\ln\binom{n-k+1}{k}\to g(x)=(1-x)\ln(1-x)-x\ln x-(1-2x)\ln(1-2x)$, and $g'(x)=0$ becomes $(1-2x)^2=x(1-x)$, i.e.\ $5x^2-5x+1=0$, whose root in $(0,\tfrac12)$ is $x=\frac{5-\sqrt5}{10}$. At this point $g(x)=\ln\varphi$, $\varphi=\frac{1+\sqrt5}2$, which matches the growth $F_{n+2}\sim\varphi^{n}$: the maximal term carries the whole exponential growth.
\end{addedblock}

And there is a ``unique'' densest packing: \added{the zigzag with all contacts at opposite walls (the word $00\cdots0$, length $\ell L+2$), unique up to reflection. At the other extreme the least dense jammed configuration alternates pairs on the two walls (word $0101\cdots$, or $1010\cdots$ with ends allowed).}
\unclear{Second sketch annotated ``$\leftarrow$ \ldots\ last \ldots\ as well'': read as the least dense case.}

\begin{center}
\begin{tikzpicture}[scale=0.28]
\draw (-1.3,0)--(14.1,0); \draw (-1.3,3.2)--(14.1,3.2);
\foreach \i in {0,...,8} { \pgfmathsetmacro{\y}{mod(\i,2)==0 ? 1 : 2.2} \draw (1.6*\i,\y) circle (1); }
\node[right] at (14.5,1.6) {\small densest: $0000\cdots$};
\end{tikzpicture}

\vspace{1ex}
\begin{tikzpicture}[scale=0.28]
\draw (-1.3,0)--(15.7,0); \draw (-1.3,3.2)--(15.7,3.2);
\foreach \x/\y in {0/1,2/1,3.6/2.2,5.6/2.2,7.2/1,9.2/1,10.8/2.2,12.8/2.2,14.4/1} { \draw (\x,\y) circle (1); }
\node[right] at (16.1,1.6) {\small least dense: $10101010$};
\end{tikzpicture}
\end{center}

\begin{exercise}\label{L11:ex:typical}
Using the law of large numbers and \added{Stirling's} approximation, show that \added{if all jammed configurations of $n+1$ discs are equally likely (Edwards' ansatz without the volume constraint),} a typical configuration has a fraction
\[
\frac12-\frac{\sqrt5}{10}
\]
of $1$'s. \added{Deduce that its length per disc is close to $2x+(1-x)L$ with $x=\frac12-\frac{\sqrt5}{10}$.} Compare to coin flips, where the fraction would be $\tfrac12$.
\end{exercise}
\orig{``Ex.\ law of large \ldots\ \& approx.\ \ldots''; partly illegible}

\subsection{Outlook}
\begin{itemize}
\item Probabilities of critical points. \unclear{First word unclear (``prune''? ``prior''?).} Flow to jammed configurations\ldots\ The resulting distribution is probably concentrated on the maxima, since critical saddles should have probability near $0$. \added{(A saddle attracts only a lower-dimensional set of starting configurations under the flow.)}
\item $\Longrightarrow$ Question: compute these values!
\item Also\ldots\ for large $N$ this is hard, so we study low-order cases.
\end{itemize}
\unclear{last line read as ``so study low order cases''}
% ===== end chapters/L11.tex =====
\begingroup\renewcommand{\thechapter}{D}\renewcommand{\theHchapter}{interludeD}
% ===== begin chapters/TOP.tex =====
% MAT670 -- Interlude D (after Lecture 11): min-max, configuration spaces and braids.
% Entirely new material (not in the handwritten notes). Labels prefixed TOP:.

\chapter[Interlude D: Min-max and braids]{Interlude D: Min-max, configuration spaces and braids}
\label{TOP:ch}

\begin{addedblock}
This interlude is new: it is not in the handwritten notes. It takes three threads from
Lectures~7--11 (the deformation theorem of Lecture~9, the max-min distance function and its balanced
contact graphs in Lectures~10--11, and the braid groups of Lecture~7) and follows each one into a
neighbouring field. To keep it readable it is set in black; all of it should be read as added material.
\end{addedblock}

Lectures~8--11 used Morse theory in two ways. For smooth functions, the Morse inequalities count critical
points of \emph{non-degenerate} functions by Betti numbers. For the min-type function
\[
  \rho(u)=\min_{i<j}d(u_i,u_j)\qquad\text{on }\Conf(N,\Sph^2),
\]
we only had a deformation theorem: the topology of $\{\rho\ge\theta\}$ changes only at balanced
configurations. This leaves some natural questions open.
\begin{itemize}
\item The deformation theorem alone already forces critical points to exist, without any
  non-degeneracy. How many, and why?
\item What do the Betti numbers of $\Conf(N,\Sph^2)/\SO(3)$ computed in Lecture~10 say about energies on
  that space, and why are they products like $(1+2t)(1+3t)\cdots$?
\item Lecture~7 defined braid groups as fundamental groups and Lecture~8 found the Weyl chambers of
  $A_{N-1}$ inside $\Conf(N,\R)$. Are these two facts related?
\end{itemize}
The answers lead to \emph{min-max theory} (Section~\ref{TOP:sec:minmax}), to Morse theory on
\emph{configuration spaces of hard particles} (Sections~\ref{TOP:sec:mintype}--\ref{TOP:sec:conf}),
and to \emph{braid and Artin groups} (Section~\ref{TOP:sec:braids}). Section~\ref{TOP:sec:gallery}
collects theorems from these areas.

Throughout, $M$ is a connected smooth manifold with a Riemannian metric and $f\colon M\to\R$ is
smooth. As in Lecture~9, $M^a=f^{-1}(-\infty,a]$. We say that $f$ has \emph{compact sublevel sets} if
every $M^a$ is compact. This holds when $M$ is compact, and also for energies such as
$\sum_{i<j}\abs{u_i-u_j}^{-1}$ on $\Conf(N,\Sph^2)$, which blow up at collisions.

%======================================================================
\section{From Morse theory to min-max}
\label{TOP:sec:minmax}

\subsection{The deformation lemma and the minimax principle}

The proof of Theorem~\ref{L9:thm:diffeo} gives slightly more than its statement. We record the
version we need.

\begin{lemma}[Deformation lemma]\label{TOP:lem:deform}
Let $f$ have compact sublevel sets, let $a<b$ and $\delta>0$, and suppose that $f^{-1}[a-\delta,b+\delta]$
contains no critical point. Then there are diffeomorphisms $\eta_t\colon M\to M$, $0\le t\le1$,
depending smoothly on $t$, with $\eta_0=\mathrm{id}$, $\eta_1(M^b)\subseteq M^a$, and $\eta_t(x)=x$
for all $t$ whenever $f(x)\le a-\delta$ or $f(x)\ge b+\delta$.
\end{lemma}

\begin{proof}
Choose a smooth $\chi\colon\R\to[0,1]$ with $\chi=1$ on $[a,b]$ and support in $(a-\delta,b+\delta)$, and
put
\[
  X=\chi(f)\,\frac{\nabla f}{\abs{\nabla f}^2}.
\]
This is smooth, because $\nabla f\neq0$ on the support of $\chi(f)$, and it vanishes outside the compact set
$f^{-1}[a-\delta,b+\delta]$. By Lemma~\ref{L9:lem:flow} it generates a one-parameter group $\varphi_t$. Along
a flow line, $\frac{d}{dt}f(\varphi_t(x))=df(X)=\chi(f)\in[0,1]$, with value $1$ while $f\in[a,b]$. Set
$\eta_t=\varphi_{-t(b-a)}$. Running the flow backwards for time $b-a$, a point with $f(x)\le b$ either
starts below $a$ and stays there ($f$ is non-increasing), or loses height at unit rate until it reaches
$a$. So $\eta_1(M^b)\subseteq M^a$. Points where $X=0$ are fixed by the flow, which gives the last claim.
\end{proof}

The lemma says that \emph{if there is no critical value near $c$, then everything up to level $c+\varepsilon$
can be pushed below $c-\varepsilon$.} Turning this around produces critical points.

\begin{theorem}[Minimax principle]\label{TOP:thm:minimax}
Let $f$ have compact sublevel sets. Let $Z\subseteq M$ be a set (possibly empty) and let $\mathcal F$ be a
family of non-empty subsets of $M$ such that $\eta_1(A)\in\mathcal F$ whenever $A\in\mathcal F$ and
$(\eta_t)$ is an isotopy of diffeomorphisms of $M$ fixing $Z$ pointwise. Put
\[
  c=\inf_{A\in\mathcal F}\ \sup_{x\in A}f(x).
\]
If $c<\infty$ and $\sup_Zf<c$, then $c$ is a critical value of $f$.
\end{theorem}

\begin{proof}
Since $f$ attains its minimum, $c>-\infty$. The critical points in $f^{-1}[c-1,c+1]$ form a compact
set, so their values form a compact set. If $c$ is not critical, there is therefore an $\varepsilon>0$ such
that $[c-3\varepsilon,c+3\varepsilon]$ contains no critical value; shrinking $\varepsilon$, we may also assume
$\sup_Zf<c-2\varepsilon$. Apply Lemma~\ref{TOP:lem:deform} with $a=c-\varepsilon$, $b=c+\varepsilon$,
$\delta=\varepsilon$. The isotopy fixes $\{f\le c-2\varepsilon\}\supseteq Z$. Pick $A\in\mathcal F$ with
$\sup_Af\le c+\varepsilon$. Then $\eta_1(A)\in\mathcal F$ and $\eta_1(A)\subseteq M^{c-\varepsilon}$, so
$\sup_{\eta_1(A)}f\le c-\varepsilon<c$, contradicting the definition of $c$.
\end{proof}

\begin{example}\label{TOP:ex:families}\leavevmode
\begin{enumerate}[nosep]
\item $\mathcal F=\{\{x\}:x\in M\}$ gives $c=\min f$. If $M$ is compact, $\mathcal F=\{M\}$ gives $c=\max f$.
\item \emph{Homology classes.} For $0\neq\alpha\in H_k(M)$ let $\mathcal F_\alpha$ be the family of compact
  $A\subseteq M$ such that $\alpha$ lies in the image of $H_k(A)\to H_k(M)$. A diffeomorphism isotopic to
  the identity acts trivially on homology, so $\mathcal F_\alpha$ is invariant. The critical value
  $c_\alpha$ is ``the lowest level at which $\alpha$ can be realized''. For the height function on the
  torus of Example~\ref{L07:ex:torus}, the classes in $H_1$ give the two saddle values $f(q)$ and $f(r)$.
\item \emph{Mountain pass.} For $p,q\in M$, take $Z=\{p,q\}$ and $\mathcal F=\{\gamma([0,1]):\gamma
  \text{ a path from }p\text{ to }q\}$.
\end{enumerate}
\end{example}

\begin{corollary}[Mountain pass]\label{TOP:cor:mp}
Let $f$ have compact sublevel sets and let $p\neq q$ be strict local minima of $f$. Then $f$ has a third
critical point, at the level
\[
  c=\inf_{\gamma\colon p\leadsto q}\ \max_{t\in[0,1]}f(\gamma(t))\ >\ \max\big(f(p),f(q)\big).
\]
\end{corollary}

\begin{proof}
Say $f(p)\le f(q)$. Since $q$ is a strict local minimum, a small geodesic sphere $S$ around $q$ (not
containing $p$ inside) satisfies $\min_Sf>f(q)$. Every path from $p$ to $q$ crosses $S$, so
$c\ge\min_Sf>f(q)\ge f(p)$. Also $c<\infty$ because $M$ is connected. Theorem~\ref{TOP:thm:minimax}
with Example~\ref{TOP:ex:families}(3) applies.
\end{proof}

\begin{figure}[ht]
\centering
\begin{tikzpicture}[scale=1.15]
  \draw[gray!70] plot[domain=0:360,samples=120,smooth] ({2.6*cos(\x)/(1+sin(\x)^2)},{2.6*sin(\x)*cos(\x)/(1+sin(\x)^2)});
  \draw[gray!50] (-1.55,0) ellipse (0.55 and 0.3); \draw[gray!50] (1.55,0) ellipse (0.55 and 0.3);
  \draw[gray!50] (-1.55,0) ellipse (0.95 and 0.5); \draw[gray!50] (1.55,0) ellipse (0.95 and 0.5);
  \draw[gray!50] (0,0) ellipse (3.2 and 1.25);
  \fill (-1.55,0) circle (1.5pt) node[below=2pt,font=\small] {$p$};
  \fill (1.55,0) circle (1.5pt) node[below=2pt,font=\small] {$q$};
  \draw[addcol,very thick] (-1.55,0) .. controls (-0.8,0.25) and (-0.3,0) .. (0,0) .. controls (0.3,0) and (0.8,-0.25) .. (1.55,0);
  \fill[flagcol] (0,0) circle (1.8pt) node[above=3pt,font=\small,flagcol] {$s$};
  \draw[thick,dashed] (-1.55,0) .. controls (-1.4,1.3) and (1.4,1.3) .. (1.55,0);
  \fill (0,0.98) circle (1.3pt) node[above,font=\scriptsize] {$\max f\circ\gamma$};
  \node[font=\scriptsize,gray] at (2.85,1.15) {level sets of $f$};
  % profile
  \begin{scope}[xshift=5.2cm,yshift=-0.6cm,scale=0.9]
    \draw[->] (0,0)--(2.9,0) node[right,font=\scriptsize]{$t$};
    \draw[->] (0,0)--(0,2.1) node[above,font=\scriptsize]{$f(\gamma(t))$};
    \draw[addcol,very thick] (0.1,0.3) .. controls (0.7,0.4) and (1.0,1.1) .. (1.4,1.1) .. controls (1.8,1.1) and (2.1,0.55) .. (2.7,0.5);
    \draw[thick,dashed] (0.1,0.3) .. controls (0.6,1.9) and (2.2,1.9) .. (2.7,0.5);
    \draw[dotted] (0,1.1)--(1.4,1.1); \node[left,font=\scriptsize] at (0,1.1) {$c$};
  \end{scope}
\end{tikzpicture}
\caption{The mountain pass. Every path from $p$ to $q$ has a highest point; the dashed path climbs
higher than necessary. The minimax level $c$ is attained by paths through the saddle $s$, where the
level set through $s$ is a figure eight.}
\label{TOP:fig:mp}
\end{figure}

\subsection{Lusternik--Schnirelmann category}

The minimax principle gives one critical value per invariant family. To get many critical points we need
many families, and a way to tell them apart. Lusternik and Schnirelmann (around 1930) measured the
families by how hard it is to contract them.

\begin{definition}\label{TOP:def:cat}
For $A\subseteq M$, the \emph{category of $A$ in $M$}, $\operatorname{cat}_M(A)$, is the least $m$ such that
$A\subseteq U_1\cup\dots\cup U_m$ with each $U_i$ open in $M$ and each inclusion $U_i\hookrightarrow M$
homotopic to a constant map. The \emph{Lusternik--Schnirelmann category} of $M$ is
$\operatorname{cat}(M)=\operatorname{cat}_M(M)$.
\end{definition}

So $\operatorname{cat}(\text{point})=1$ and $\operatorname{cat}(\Sph^n)=2$ (two hemispheres). This is the
unnormalized convention; many papers subtract $1$. The following properties are immediate.
Category is monotone in $A$ and subadditive, $\operatorname{cat}_M(A\cup B)\le\operatorname{cat}_M(A)+
\operatorname{cat}_M(B)$. A homeomorphism of $M$ homotopic to the identity preserves it. A finite set has
category $1$: a union of small disjoint balls can be contracted inside the connected manifold $M$.

\begin{theorem}[Cup-length bound]\label{TOP:thm:cuplength}
If there are classes $u_1,\dots,u_k\in H^*(M)$ of positive degree (any coefficients) with
$u_1\smile\dots\smile u_k\neq0$, then $\operatorname{cat}(M)\ge k+1$.
\end{theorem}

\begin{proof}
Suppose $M=U_1\cup\dots\cup U_k$ with each $U_i\hookrightarrow M$ null-homotopic. Then the restriction
of $u_i$ to $U_i$ factors through the cohomology of a point, which vanishes in positive degree, so
$u_i|_{U_i}=0$. By the long exact sequence of the pair $(M,U_i)$, $u_i$ is the image of some
$\tilde u_i\in H^*(M,U_i)$. The relative cup product (defined for open sets) gives
\[
  \tilde u_1\smile\dots\smile\tilde u_k\in H^*(M,U_1\cup\dots\cup U_k)=H^*(M,M)=0,
\]
and it maps to $u_1\smile\dots\smile u_k$. Hence the product is $0$, a contradiction.
\end{proof}

To count critical points we also need a version of the deformation lemma that is allowed to stop near the
critical points at level $c$.

\begin{lemma}[Deformation with a hole]\label{TOP:lem:hole}
Let $f$ have compact sublevel sets, let $c\in\R$, let $K_c$ be the (compact) set of critical points
with value $c$, and let $U\supseteq K_c$ be open. Then there are $\varepsilon>0$ and a diffeomorphism
$\eta$ of $M$, isotopic to the identity, with $\eta\big(M^{c+\varepsilon}\setminus U\big)\subseteq M^{c-\varepsilon}$.
\end{lemma}

\begin{proof}
Choose open sets $K_c\subseteq V\subseteq\overline V\subseteq U_1\subseteq\overline U_1\subseteq U$ with
$\overline U_1$ compact.

\emph{Step 1: the gradient is bounded below away from $V$ near level $c$.} There are $\varepsilon_0,\sigma>0$
with $\abs{\nabla f}\ge\sigma$ on $f^{-1}[c-\varepsilon_0,c+\varepsilon_0]\setminus V$. Otherwise there are
$x_n\notin V$ with $f(x_n)\to c$ and $\nabla f(x_n)\to0$. They lie in the compact set
$f^{-1}[c-1,c+1]$, so a subsequence converges to some $x\notin V$ with $\nabla f(x)=0$ and $f(x)=c$. Then
$x\in K_c\subseteq V$, a contradiction.

\emph{Step 2: the vector field.} Let $\psi\colon M\to[0,1]$ be smooth with $\psi=0$ on $V$ and $\psi=1$
outside $U_1$, and let $\chi$ be a cutoff with $\chi=1$ on $[c-\varepsilon,c+\varepsilon]$ and support in
$(c-\varepsilon_0,c+\varepsilon_0)$, for an $\varepsilon\le\varepsilon_0/2$ fixed below. Put
$X=\chi(f)\psi\nabla f/\abs{\nabla f}^2$. By Step~1 it is smooth, compactly supported, and
$\abs X\le1/\sigma$. Along its backward flow $f$ is non-increasing, and
$\frac{d}{dt}f=-\chi(f)\psi$.

\emph{Step 3: flow lines from outside $U$ do not reach $U_1$ in time $2\varepsilon$.} Let
$d=\dist(M\setminus U,\overline U_1)>0$. A flow line has speed at most $1/\sigma$, so in time $T$ it travels
at most $T/\sigma$. Choose $\varepsilon$ with $2\varepsilon<d\sigma$. A flow line starting at
$x\notin U$ then stays outside $U_1$, where $\psi=1$, for $0\le t\le2\varepsilon$. If moreover $f(x)\le
c+\varepsilon$, then $f$ drops at unit rate while it lies in $[c-\varepsilon,c+\varepsilon]$, so after time
$2\varepsilon$ we have $f\le c-\varepsilon$. Take $\eta$ to be the backward flow for time $2\varepsilon$.
\end{proof}

\begin{theorem}[Lusternik--Schnirelmann]\label{TOP:thm:LS}
Every smooth function on a compact connected manifold $M$ has at least $\operatorname{cat}(M)$ critical
points.
\end{theorem}

\begin{proof}
Let $m=\operatorname{cat}(M)$ and, for $1\le k\le m$, let $\mathcal F_k$ be the family of closed
$A\subseteq M$ with $\operatorname{cat}_M(A)\ge k$. These families are invariant under diffeomorphisms
isotopic to the identity and nested, $\mathcal F_1\supseteq\dots\supseteq\mathcal F_m\ni M$. So the values
$c_k=\inf_{A\in\mathcal F_k}\max_Af$ satisfy $c_1\le c_2\le\dots\le c_m\le\max f$, and each is critical by
Theorem~\ref{TOP:thm:minimax}. If they are distinct we have $m$ critical points. Suppose instead that
$c_k=c_{k+1}=c$. We claim that $K_c$ is infinite. If not, let $U$ be a union of small disjoint balls around
the points of $K_c$, so that $\operatorname{cat}_M(\overline U)=1$. Take $\varepsilon,\eta$ from
Lemma~\ref{TOP:lem:hole} and $A\in\mathcal F_{k+1}$ with $\max_Af\le c+\varepsilon$. The closed set $B=A\setminus
U$ has $\operatorname{cat}_M(B)\ge\operatorname{cat}_M(A)-\operatorname{cat}_M(\overline U)\ge k$, so
$\eta(B)\in\mathcal F_k$. But $\eta(B)\subseteq M^{c-\varepsilon}$, which gives $c_k\le c-\varepsilon$, a
contradiction.
\end{proof}

The argument also shows that if two of the minimax values coincide, then the critical set at that level
is not just infinite but has category at least $2$ in $M$.

\begin{example}[Three critical points on the torus]\label{TOP:ex:torus}
On $\T^2=\R^2/\Z^2$ the de Rham classes $a=[dx]$ and $b=[dy]$ in $H^1(\T^2;\R)$ have
$a\smile b=[dx\wedge dy]\neq0$. By Theorem~\ref{TOP:thm:cuplength}, $\operatorname{cat}(\T^2)\ge3$, so \emph{every}
smooth function on $\T^2$ has at least three critical points. A \emph{Morse} function needs at least
$b_0+b_1+b_2=4$ (Lecture~10), as for the height function of Lecture~8. The bound $3$ is sharp. Consider
\[
  f(x,y)=\sin(\pi x)\,\sin(\pi y)\,\sin(\pi(x+y)),
\]
which is well defined on $\R^2/\Z^2$ because each shift $x\mapsto x+1$ or $y\mapsto y+1$ changes the
sign of exactly two factors. Using $\sin u\cos v+\cos u\sin v=\sin(u+v)$,
\[
  \partial_xf=\pi\sin(\pi y)\sin(\pi(2x+y)),\qquad \partial_yf=\pi\sin(\pi x)\sin(\pi(x+2y)).
\]
If $y\equiv0$, then $\partial_yf=\pi\sin^2(\pi x)$ vanishes only for $x\equiv0$, and likewise if $x\equiv0$.
Otherwise $2x+y\equiv x+2y\equiv0\pmod1$, so $x\equiv y$ and $3x\equiv0$. The critical points are therefore
exactly
\[
  (0,0)\ (f=0),\qquad \big(\tfrac13,\tfrac13\big)\ \big(f=\tfrac{3\sqrt3}8,\ \text{max}\big),\qquad
  \big(\tfrac23,\tfrac23\big)\ \big(f=-\tfrac{3\sqrt3}8,\ \text{min}\big).
\]
Near the origin $f\approx\pi^3xy(x+y)$, a \emph{monkey saddle}: a degenerate critical point at which the
zero set consists of three lines through the point (Figure~\ref{TOP:fig:torus}).
\end{example}

\begin{figure}[ht]
\centering
\begin{tikzpicture}[scale=3]
  \fill[addcol!12] (0,0)--(1,0)--(0,1)--cycle;
  \fill[flagcol!12] (1,0)--(1,1)--(0,1)--cycle;
  \draw[thick] (0,0) rectangle (1,1);
  \draw[thick] (1,0)--(0,1);
  \fill (1/3,1/3) circle (0.6pt) node[below right,font=\small] {max};
  \fill (2/3,2/3) circle (0.6pt) node[above left,font=\small] {min};
  \node[addcol,font=\small] at (0.2,0.55) {$f>0$};
  \node[flagcol,font=\small] at (0.8,0.45) {$f<0$};
  \foreach \p in {(0,0),(1,0),(0,1),(1,1)} \fill[flagcol] \p circle (0.7pt);
  \node[below left,font=\scriptsize] at (0,0) {$(0,0)$};
  \node[font=\scriptsize,align=left,right] at (1.08,0.5) {all four corners are one point:\\ the monkey saddle; six sectors\\ of alternating sign meet there};
  \node[below,font=\scriptsize] at (0.5,-0.04) {$x$}; \node[left,font=\scriptsize] at (-0.03,0.5) {$y$};
\end{tikzpicture}
\caption{The function $\sin(\pi x)\sin(\pi y)\sin(\pi(x+y))$ on the torus $\R^2/\Z^2$. Its zero set is the
three closed curves $x\equiv0$, $y\equiv0$, $x+y\equiv0$, which cut the torus into two triangles. Each
triangle contains one extremum, at its centroid.}
\label{TOP:fig:torus}
\end{figure}

\begin{remark}[Back to Interlude~A]
Up to the constant $(2i)^3$, $f$ is the \emph{Weyl denominator}
$\prod_{j<k}\big(e^{i(\theta_j-\theta_k)/2}-e^{-i(\theta_j-\theta_k)/2}\big)$ of $\mathrm{SU}(3)$, written
in the coordinates $x=(\theta_1-\theta_2)/2\pi$, $y=(\theta_2-\theta_3)/2\pi$ on the maximal torus of
Section~\ref{LIE:sec:torus} (modulo its centre). Its zero set is the set of non-regular elements, where two
eigenvalues collide, that is, the walls of $\Conf(3,\R)$ wrapped around the circle. Its two extrema are
the regular elements whose eigenvalues are \emph{equally spaced} on the circle, in the two cyclic orders.
So maximizing $\abs f$ is the circle version of maximizing the injectivity radius $\rho$ in Lecture~7.
\end{remark}

\subsection{Closed geodesics and minimal surfaces}
\label{TOP:sec:geodesics}

Min-max was born in geometry, with the space of closed curves in place of $M$ and length in place
of $f$. Birkhoff (Trans.\ AMS, 1917) swept a Riemannian $2$-sphere by a one-parameter family of closed
curves, like the circles of latitude, and took
\[
  W=\inf_{\text{sweepouts}}\ \max_t\ \operatorname{length}(\gamma_t).
\]
Lemma~\ref{TOP:lem:deform} has an infinite-dimensional analogue, a curve-shortening process applied to the
whole family. If no closed geodesic had length $W$, this process would push every sweepout below level
$W$. So \emph{every Riemannian metric on $\Sph^2$ has a closed geodesic}, of length $W$. Lusternik and
Schnirelmann (1929) used three linked families to obtain \emph{three simple closed geodesics}. Their
deformation step had a gap, which was closed only much later; the proof now accepted uses Grayson's
curve-shortening flow (Ann.\ of Math., 1989). The number three is the category of $\R P^2$. On the round
sphere the simple closed geodesics are the great circles, parametrized by $\R P^2$, and the $\Z/2$
cup-length of $\R P^2$ is $2$.

Almgren and Pitts replaced curves by hypersurfaces and length by area. Their min-max theory (Pitts' book,
1981), with the regularity theory of Schoen and Simon (1981), produces a smooth closed embedded minimal
hypersurface in every closed Riemannian manifold of dimension between $3$ and $7$. In the last fifteen
years this machinery has solved two famous problems; see Theorems~\ref{TOP:thm:willmore}
and~\ref{TOP:thm:yau} in the gallery. The Willmore conjecture is a statement about tori, and Marques and Neves
prove it by min-max over a \emph{five}-parameter family. They show that the minimal surface it produces
has Morse index at most $5$ and is not a great sphere. By a theorem of Urbano (Proc.\ AMS, 1990), such a
minimal surface in $\Sph^3$ is a Clifford torus, of area $2\pi^2$. This is Lusternik--Schnirelmann in spirit: a topologically non-trivial family forces a
critical point, and an index count identifies it.

%======================================================================
\section{Min-type functions and Morse inequalities for configurations}
\label{TOP:sec:mintype}

\subsection{Minimax for the injectivity radius}

Now return to $\rho$ on $\Conf(N,\Sph^2)$, or on $B\Conf(N)=\Conf(N,\Sph^2)/\SO(3)$ as in Lecture~10. The
roles of sub- and superlevel sets are reversed: we study $\Conf(N,\theta)=\{\rho\ge\theta\}$, which shrinks
as $\theta$ grows. By Theorem~\ref{L11:thm:forcebalance}, the critical configurations are those whose
contact graph carries a balanced system of non-negative weights. The deformation theorem quoted in the
proof sketch of Lecture~11 states that if $[\theta_1,\theta_2]$ contains no critical value, then
$\Conf(N,\theta_2)$ is a deformation retract of $\Conf(N,\theta_1)$. (See Baryshnikov--Bubenik--Kahle for
the general min-type framework, and Kusner--Kusner--Lagarias--Shlosman for the sphere.) Given this, the proof
of Theorem~\ref{TOP:thm:minimax} applies to $\rho$ with all inequalities reversed:
\begin{center}
\emph{if $\mathcal F$ is a deformation-invariant family and
$\theta_{\mathcal F}=\sup_{A\in\mathcal F}\min_A\rho>0$, then $\theta_{\mathcal F}$ is a critical value,
realized by a balanced configuration.}
\end{center}
One ingredient should be checked: the critical values of $\rho$ in $[\theta_0,\pi]$, $\theta_0>0$, form a
closed set. In a limit of balanced configurations, contact edges can only be gained, and the normalized
weights have a convergent subsequence, so the limit is balanced.

The extreme case is $\mathcal F=\{\text{points}\}$. Then $\theta_{\mathcal F}=\max\rho$ is the solution of
the \emph{Tammes problem}, the top critical value of $\rho$. Beyond it $\Conf(N,\theta)$ is empty.

\begin{figure}[ht]
\centering
\begin{tikzpicture}[scale=1.25]
  % top view
  \begin{scope}
    \draw[thick] (0,0) circle (1.3);
    \foreach \a/\l in {90/1,210/2,330/3} {
      \draw[addcol,thick] (0,0)--(\a:1.3);
      \fill (\a:1.3) circle (1.6pt);
      \node[font=\scriptsize] at (\a:1.55) {$e_\l$};
      \draw[-{Stealth},flagcol,thick] (0,0)--(\a:0.55);
    }
    \fill (0,0) circle (1.8pt); \node[font=\scriptsize] at (0.32,0.16) {$n$};
    \node[font=\scriptsize,align=center] at (0,-1.85) {from above: the forces at the\\ pole $n$ cancel ($120^\circ$ apart)};
  \end{scope}
  % side view
  \begin{scope}[xshift=4.2cm]
    \draw[thick] (0,0) circle (1.3);
    \draw[gray] (-1.3,0)--(1.3,0);
    \draw[addcol,thick] (0,1.3)--(0,-1.3);
    \draw[addcol,dashed] (0,1.3) arc[start angle=90,end angle=-90,x radius=1.126,y radius=1.3];
    \draw[addcol,dashed] (0,1.3) arc[start angle=90,end angle=270,x radius=1.126,y radius=1.3];
    \fill (0,1.3) circle (1.8pt) node[above,font=\scriptsize] {$n$};
    \fill (0,-1.3) circle (1.8pt) node[below,font=\scriptsize] {$s$};
    \fill (0,0) circle (1.6pt); \node[font=\scriptsize] at (0.22,0.15) {$e_1$};
    \fill[gray] (1.126,0) circle (1.3pt); \fill[gray] (-1.126,0) circle (1.3pt);
    \draw[-{Stealth},flagcol,thick] (0,0)--(0,0.55);
    \draw[-{Stealth},flagcol,thick] (0,0)--(0,-0.55);
    \node[font=\scriptsize,align=center] at (0,-1.95) {from the side: the forces at an\\ equatorial point cancel};
  \end{scope}
\end{tikzpicture}
\caption{A balanced configuration for $N=5$: the two poles and three points on the equator. The
contact graph (blue) has the six pole-to-equator edges, all of length $\pi/2$, with equal weights. This is
a Tammes configuration, $\max\rho=\pi/2$. It is \emph{not} isolated: the equatorial points may slide
freely as long as their gaps stay at least $\pi/2$. Critical sets of min-type functions can be large.}
\label{TOP:fig:N5}
\end{figure}

\begin{example}[A min-max value computed exactly]\label{TOP:ex:N4pass}
Take $N=4$, labelled. The regular tetrahedra (the maxima of $\rho$, $\rho=\arccos(-\tfrac13)$) come
in two classes modulo $\SO(3)$, mirror images of each other. The signed volume
$\det(u_2-u_1,u_3-u_1,u_4-u_1)$ is rotation invariant and has opposite signs on the two classes, so every
path from one to the other passes through a configuration whose four points lie in a plane, hence on a
circle. On a circle, the closest pair of four points is at angular distance at most $\pi/2$, with equality
only for the square on a great circle (Exercise~\ref{TOP:exc:N4}). On the other hand, the path
\[
\begin{gathered}
  u(s)=\big((\cos s,0,\sin s),\,(-\cos s,0,\sin s),\,(0,\cos s,-\sin s),\,(0,-\cos s,-\sin s)\big),\\
  \abs{s}\le\arcsin\tfrac1{\sqrt3},
\end{gathered}
\]
runs from one tetrahedron ($s<0$) through the square ($s=0$) to the other ($s>0$). Along it,
$\rho=\min\big(\arccos(-\sin^2s),\,\pi-2\abs s\big)\ge\pi/2$. Hence the mountain-pass value between the two
tetrahedra is exactly $\pi/2$, and the pass is the balanced square of Lecture~11. This is the
``pushing alternate vertices'' deformation of Lecture~11 read as a min-max statement.
\end{example}

\subsection{What the Betti numbers of Lecture 10 force}

Lecture~10 computed the Poincar\'e polynomial
\[
  P\big(B\Conf(N)\big)(t)=\prod_{k=2}^{N-2}(1+kt).
\]
Setting $t=1$ and $t=-1$:
\[
  \sum_\lambda b_\lambda\big(B\Conf(N)\big)=\prod_{k=2}^{N-2}(k+1)=\frac{(N-1)!}2,\quad
  \chi\big(B\Conf(N)\big)=\prod_{k=2}^{N-2}(1-k)=(-1)^{N-3}(N-3)!.
\]
For $N=6$ this is $1+9+26+24=60=5!/2$, matching the table in Lecture~10, and $\chi=1-9+26-24=-6$.

For $N\ge3$, $\SO(3)$ acts freely on $\Conf(N,\Sph^2)$ (a non-trivial rotation fixes only two points of the
sphere), so $B\Conf(N)$ is a smooth manifold of dimension $2N-3$. Let $E$ be a smooth $\SO(3)$-invariant
function on $\Conf(N,\Sph^2)$ that tends to $+\infty$ at collisions, for example the Riesz energy
$E_s=\sum_{i<j}\abs{u_i-u_j}^{-s}$ ($s>0$). Its sublevel sets are compact, and the Morse theory of
Lectures~8--10 applies to the induced function on $B\Conf(N)$ (Milnor, \emph{Morse Theory}, Thm.~3.5).

\begin{proposition}\label{TOP:prop:count}
If $E$ induces a Morse function on $B\Conf(N)$, then $E$ has at least $(N-1)!/2$ critical $\SO(3)$-orbits
of labelled configurations, at least $b_\lambda$ of them of index $\lambda$. If there are finitely many,
then $\sum_\lambda(-1)^\lambda c_\lambda=(-1)^{N-3}(N-3)!$.
\end{proposition}

\begin{example}[$N=4$]\label{TOP:ex:N4morse}
Here $b_0=1$, $b_1=2$ and $\chi=-1$. For the Coulomb energy $E_1$, the regular tetrahedron and the square on a
great circle are critical by symmetry. In $B\Conf(4)$ they give $4!/12=2$ tetrahedral classes and
$4!/8=3$ square classes (the orders $12$ and $8$ are their rotation groups). A numerical computation of the
Hessian (editor's check) shows that both are non-degenerate modulo rotations, with index $0$ and $1$. The
count $2\cdot(+1)+3\cdot(-1)=-1=\chi$ is consistent, and $c_0=2\ge1$, $c_1=3\ge2$. The squares are
mountain passes between the tetrahedra, exactly as for $\rho$ in Example~\ref{TOP:ex:N4pass}.
\end{example}

\begin{remark}[Labelled versus unlabelled]
These bounds count \emph{labelled} configurations. The symmetric group $\Sigma_N$ acts on $B\Conf(N)$, and
an unlabelled critical configuration whose rotation group is $G$ accounts for $N!/\abs G$ labelled classes.
So the bound of Proposition~\ref{TOP:prop:count} only says $\sum 1/\abs G\ge 1/(2N)$, summed over the
unlabelled critical configurations. Topology forces essentially \emph{nothing} about how many genuinely
different equilibria an energy has. The plane shows the same effect in cohomology. $\Conf(N,\C)$ has total
Betti number $N!$, but by a theorem of Arnold (1970) the rational cohomology of the \emph{unordered} space
$\Conf(N,\C)/\Sigma_N$ is that of a circle. Forgetting the labels kills almost all of the rational topology.
Numerical experiments, by contrast, suggest that the number of distinct local minima of energies such as
$E_1$ grows quickly with $N$. That growth is a fact about geometry, not about topology.
\end{remark}

\begin{remark}[Without non-degeneracy]
The proof of Theorem~\ref{TOP:thm:LS} goes through for any function with compact sublevel sets. If such a
function has finitely many critical points, then a large sublevel set is a deformation retract of the whole
space, and that set has full category. Over $\Q$ the cohomology of $\Conf(N-1,\C)$ is generated in degree
$1$ and has top degree $N-2$ (Section~\ref{TOP:sec:conf}), so its cup-length is $N-2$. Since
$\Conf(N-1,\C)/\SO(2)\cong B\Conf(N)$ and the circle bundle is trivial (Lecture~10), K\"unneth gives
cup-length $N-3$ for $B\Conf(N)$. Hence \emph{every} smooth proper $\SO(3)$-invariant energy, degenerate or
not, has at least $N-2$ critical orbits. This is far weaker than $(N-1)!/2$: degenerate critical points can
absorb many non-degenerate ones, as the monkey saddle of Example~\ref{TOP:ex:torus} absorbs two saddles.
\end{remark}

%======================================================================
\section{Configuration spaces of the plane, and of thick particles}
\label{TOP:sec:conf}

\subsection{Arnold's computation}

Lecture~10 quoted the Poincar\'e polynomial of $\Conf(n,\R^2)=\Conf(n,\C)$. Here is the proof.
For $i\ne j$ let
\[
  \omega_{ij}=\frac1{2\pi}\,d\arg(z_i-z_j),
\]
a closed $1$-form on $\Conf(n,\C)$ with $\omega_{ij}=\omega_{ji}$. Its integral over a loop is the number of
times $z_i$ winds around $z_j$.

\begin{theorem}[Fadell--Neuwirth 1962; Arnold 1969]\label{TOP:thm:arnold}
\begin{enumerate}[nosep]
\item Forgetting the last point, the map $\pi\colon\Conf(n,\C)\to\Conf(n-1,\C)$ is a locally trivial
  fibration with fibre $\C\setminus\{n-1\text{ points}\}$, and it has a section.
\item $\Conf(n,\C)$ is aspherical, i.e.\ $\pi_k=0$ for $k\ge2$. Its fundamental group, the pure braid group
  $P_n$, is an iterated semidirect product $P_n\cong F_{n-1}\rtimes P_{n-1}$ of free groups.
\item $H^*(\Conf(n,\C);\Z)$ is torsion-free with Poincar\'e polynomial
  $\prod_{k=1}^{n-1}(1+kt)$. It is generated by the classes $[\omega_{ij}]$, and all relations follow from
  $\omega_{ij}^2=0$ and the \emph{Arnold relations} $\omega_{ij}\omega_{jk}+\omega_{jk}\omega_{ki}+\omega_{ki}\omega_{ij}=0$.
\end{enumerate}
\end{theorem}

\begin{proof}[Proof of (2) and of the Poincar\'e polynomial, given (1)]
(2) The fibre $F$ is homotopy equivalent to a wedge of $n-1$ circles, so it is aspherical. Inductively the base is
aspherical, and the base case $\Conf(1,\C)=\C$ is contractible. The long exact sequence
$\pi_k(F)\to\pi_k(\Conf(n,\C))\to\pi_k(\Conf(n-1,\C))$ then gives $\pi_k(\Conf(n,\C))=0$ for $k\ge2$, and in
degree one it gives $1\to F_{n-1}\to P_n\to P_{n-1}\to1$. The sequence splits by the section
$(z_1,\dots,z_{n-1})\mapsto(z_1,\dots,z_{n-1},1+\sum\abs{z_i})$.

(3) The fibre over $(z_1,\dots,z_{n-1})$ is the plane punctured at the $z_i$. Its $H^1$ is free with
basis dual to small loops around the punctures. The global classes $[\omega_{in}]$, $1\le i\le n-1$,
restrict to exactly this dual basis in every fibre: going once around $z_i$ changes $\arg(z_n-z_i)$ by $2\pi$
and leaves $\arg(z_n-z_j)$ unchanged. So the Leray--Hirsch theorem applies, $H^*(\Conf(n,\C))\cong
H^*(\Conf(n-1,\C))\otimes H^*(F)$ as modules, and $P_n(t)=P_{n-1}(t)\,(1+(n-1)t)$.
\end{proof}

For $n=3$ the change of variables $(z_1,z_2,z_3)\mapsto\big(z_1,\ z_2-z_1,\ \frac{z_3-z_1}{z_2-z_1}\big)$
gives $\Conf(3,\C)\cong\C\times\C^*\times(\C\setminus\{0,1\})$, with Poincar\'e polynomial $(1+t)(1+2t)$.
Dividing by the affine group $z\mapsto az+b$ in this way is the plane version of Lecture~10's
reduction of $\Conf(N,\Sph^2)$ by $\SO(3)$.

\subsection{Thick particles}

For points, $\Conf(n,\C)$ has the clean topology above. The physical systems of Lecture~11 are made of
\emph{hard discs}, and there the topology depends on the radius. Fix a region $R$ and let
$\Conf(n,r;R)$ be the space of $n$ non-overlapping discs of radius $r$ in $R$. This is a superlevel
set of a min-type ``tautological'' function: the minimum of half of all pairwise distances and of all
distances from centres to $\partial R$. Its critical points are again balanced configurations, now with
forces from the walls as well. This is the min-type Morse theory of Baryshnikov, Bubenik and Kahle
(Theorem~\ref{TOP:thm:BBK}). Some landmarks:
\begin{itemize}
\item For small $r$ nothing happens: $\Conf(n,r;R)\simeq\Conf(n,R)$. In a box with shortest side $L$, the
  first critical radius is $L/2n$, at the configuration of $n$ discs in a row across the box.
\item Carlsson, Gorham, Kahle and Mason (Phys.\ Rev.\ E 85, 2012) computed the homology of small hard-disc systems
  numerically and found topology changing already for very few discs, at the balanced configurations.
\item In an infinite strip, Alpert, Kahle and MacPherson (J.\ Appl.\ Comput.\ Topology,
  2021) found three ``homological phases'', named after gas, liquid and solid
  (Theorem~\ref{TOP:thm:AKM}). In their units (discs of diameter $1$), the channel of Lecture~11 has width
between $1$ and $1+\sqrt3/2$, just above the trivial case of width $1$.
\item The thick-disc spaces are not just the point spaces with holes filled in.
  Ramsey's \emph{thick braids} (Theorem~\ref{TOP:thm:thick}) are loops of discs that are non-trivial even
  though the corresponding loops of points can be contracted.
\end{itemize}

\subsection{Robots, motion planning and algorithms}

Configuration spaces are also the state spaces of machines. Ghrist (in a 2001 volume dedicated to
J.~Birman) studied $n$ robots moving on a network of tracks, the configuration space of $n$ points on a
graph. He showed that it is aspherical, with homological dimension bounded by the number of
\emph{essential} vertices (valence at least $3$); Abrams' thesis (2000) gave a cubical model.

Farber (Discrete Comput.\ Geom., 2003) asked how discontinuous a motion planner must be. A
\emph{motion planner} on $U\subseteq X\times X$ is a continuous map $s$ assigning to each $(x,y)\in U$ a path from $x$
to $y$. The \emph{topological complexity} $\operatorname{TC}(X)$ is the least number of open sets covering
$X\times X$ that carry motion planners. This is a Lusternik--Schnirelmann-type invariant, with the same
unnormalized convention as in Definition~\ref{TOP:def:cat}. If $\operatorname{TC}(X)=1$, then
$s(x_0,\cdot)$ is a contraction of $X$. So a non-contractible space, such as $\Conf(n,\R^2)$ for $n\ge2$,
needs at least two planning rules (Exercise~\ref{TOP:exc:TC}). For $n$ robots
in $\R^m$ the answer is known exactly (Theorem~\ref{TOP:thm:farber}).

Smale (J.\ Complexity, 1987) turned the same idea on numerical analysis. A monic polynomial of degree $d$
with distinct roots is the same thing as a point of $\Conf(d,\C)/\Sigma_d$, the unordered set of its roots.
An algorithm that outputs approximations to the roots must branch (``if \dots\ then \dots'') at least as often as the
\emph{Schwarz genus} of the covering $\Conf(d,\C)\to\Conf(d,\C)/\Sigma_d$ requires. Using the
cohomology of braid groups, Smale showed that for sufficiently accurate approximations the number of
branchings exceeds $(\log_2d)^{2/3}$. Vassiliev
(Funct.\ Anal.\ Appl., 1988; see also his book \emph{Complements of Discriminants of Smooth Maps}, AMS, 1992)
computed the genus exactly when $d$ is a prime power: it equals $d$. So for these degrees the lower bound
becomes linear in $d$. The genus is always at most $d$, and Arone (Contemp.\ Math.\ 399, 2006) showed
that it is strictly smaller than $d$ when $d$ is neither a prime power nor twice a prime power.

%======================================================================
\section{Braids}
\label{TOP:sec:braids}

\subsection{Artin's presentation}

Recall from Lecture~7 that $B_N=\pi_1\big(\Conf(N,\C)/\Sigma_N\big)$ and $P_N=\pi_1(\Conf(N,\C))$. Put the base
point at $1,2,\dots,N$ on the real line. Let $\sigma_i$ be the loop in which the points at $i$ and $i+1$
exchange places by a half turn, counterclockwise, with the others fixed. In a braid diagram (time
running downwards, seen from a fixed side) we draw this as strand $i$ crossing over strand $i+1$.

\begin{theorem}[Artin 1925, 1947]\label{TOP:thm:artin}
$B_N$ is generated by $\sigma_1,\dots,\sigma_{N-1}$ with the relations
\[
  \sigma_i\sigma_{i+1}\sigma_i=\sigma_{i+1}\sigma_i\sigma_{i+1}\quad(1\le i\le N-2),\qquad
  \sigma_i\sigma_j=\sigma_j\sigma_i\quad(\abs{i-j}\ge2).
\]
The map $B_N\to\Sigma_N$ sending $\sigma_i$ to the transposition $(i\ i{+}1)$ has kernel $P_N$.
\end{theorem}

\begin{figure}[ht]
\centering
\begin{tikzpicture}[scale=0.62,yscale=0.9]
  % sigma_i on N strands
  \begin{scope}
    \foreach \x in {0,2,3.5} \draw[thick] (\x,0)--(\x,-2);
    \draw[thick] (2,0) .. controls (2,-1) and (1,-1) .. (1,-2);
    \draw[white,line width=5pt] (1,0) .. controls (1,-1) and (2,-1) .. (2,-2);
    \draw[thick] (1,0) .. controls (1,-1) and (2,-1) .. (2,-2);
    \foreach \x in {0,1,2,3.5} {\fill (\x,0) circle (2pt); \fill (\x,-2) circle (2pt);}
    \node[font=\scriptsize] at (0,0.45) {$1$}; \node[font=\scriptsize] at (1,0.45) {$i$};
    \node[font=\scriptsize] at (2,0.45) {$i{+}1$}; \node[font=\scriptsize] at (3.5,0.45) {$N$};
    \node[font=\scriptsize] at (0.5,-1) {$\cdots$}; \node[font=\scriptsize] at (2.75,-1) {$\cdots$};
    \node at (1.75,-2.8) {$\sigma_i$};
  \end{scope}
  % braid relation: left side sigma1 sigma2 sigma1
  \begin{scope}[xshift=6.3cm]
    % level 1: sigma1
    \draw[thick] (1,0) .. controls (1,-0.5) and (0,-0.5) .. (0,-1);
    \draw[white,line width=5pt] (0,0) .. controls (0,-0.5) and (1,-0.5) .. (1,-1);
    \draw[thick] (0,0) .. controls (0,-0.5) and (1,-0.5) .. (1,-1);
    \draw[thick] (2,0)--(2,-1);
    % level 2: sigma2
    \draw[thick] (2,-1) .. controls (2,-1.5) and (1,-1.5) .. (1,-2);
    \draw[white,line width=5pt] (1,-1) .. controls (1,-1.5) and (2,-1.5) .. (2,-2);
    \draw[thick] (1,-1) .. controls (1,-1.5) and (2,-1.5) .. (2,-2);
    \draw[thick] (0,-1)--(0,-2);
    % level 3: sigma1
    \draw[thick] (1,-2) .. controls (1,-2.5) and (0,-2.5) .. (0,-3);
    \draw[white,line width=5pt] (0,-2) .. controls (0,-2.5) and (1,-2.5) .. (1,-3);
    \draw[thick] (0,-2) .. controls (0,-2.5) and (1,-2.5) .. (1,-3);
    \draw[thick] (2,-2)--(2,-3);
    \foreach \x in {0,1,2} {\fill (\x,0) circle (2pt); \fill (\x,-3) circle (2pt);}
    \node at (1,-3.8) {$\sigma_1\sigma_2\sigma_1$};
  \end{scope}
  \node at (9.6,-1.5) {$=$};
  % right side sigma2 sigma1 sigma2
  \begin{scope}[xshift=10.6cm]
    \draw[thick] (2,0) .. controls (2,-0.5) and (1,-0.5) .. (1,-1);
    \draw[white,line width=5pt] (1,0) .. controls (1,-0.5) and (2,-0.5) .. (2,-1);
    \draw[thick] (1,0) .. controls (1,-0.5) and (2,-0.5) .. (2,-1);
    \draw[thick] (0,0)--(0,-1);
    \draw[thick] (1,-1) .. controls (1,-1.5) and (0,-1.5) .. (0,-2);
    \draw[white,line width=5pt] (0,-1) .. controls (0,-1.5) and (1,-1.5) .. (1,-2);
    \draw[thick] (0,-1) .. controls (0,-1.5) and (1,-1.5) .. (1,-2);
    \draw[thick] (2,-1)--(2,-2);
    \draw[thick] (2,-2) .. controls (2,-2.5) and (1,-2.5) .. (1,-3);
    \draw[white,line width=5pt] (1,-2) .. controls (1,-2.5) and (2,-2.5) .. (2,-3);
    \draw[thick] (1,-2) .. controls (1,-2.5) and (2,-2.5) .. (2,-3);
    \draw[thick] (0,-2)--(0,-3);
    \foreach \x in {0,1,2} {\fill (\x,0) circle (2pt); \fill (\x,-3) circle (2pt);}
    \node at (1,-3.8) {$\sigma_2\sigma_1\sigma_2$};
  \end{scope}
\end{tikzpicture}
\caption{Left: the Artin generator $\sigma_i$. Right: the braid relation. In both braids the strand
starting on the left passes over the other two, and the strands starting in the middle and on the right
cross once. The two diagrams differ by sliding the middle crossing past the strand that lies over it.}
\label{TOP:fig:braid}
\end{figure}

The far commutation relation says that exchanges in disjoint regions of the plane can be performed in
either order. This uses the fact that \emph{points take up no room}, a point we return to in
Theorem~\ref{TOP:thm:thick}. A first gem: the braid relation forces all generators to become equal in
the abelianization, and a winding number shows that nothing more collapses, so $B_N$ is infinite.

\begin{proposition}\label{TOP:prop:abel}
For $N\ge2$ the abelianization of $B_N$ is $\Z$, generated by the image of any $\sigma_i$. The
abelianization map is induced by the discriminant
$\Delta(z)=\prod_{i<j}(z_i-z_j)^2\colon\Conf(N,\C)/\Sigma_N\to\C^*$, and it sends every $\sigma_i$ to $1$.
\end{proposition}

\begin{proof}
In an abelian quotient, $\sigma_i\sigma_{i+1}\sigma_i=\sigma_{i+1}\sigma_i\sigma_{i+1}$ becomes
$\sigma_i=\sigma_{i+1}$, so the abelianization is cyclic, generated by $\sigma_1$. Conversely,
$\Delta$ is symmetric, so it is defined on unordered configurations, and it takes values in $\C^*$. Along
$\sigma_i$, the argument of $z_i-z_{i+1}$ changes by $\pi$ while every other factor of $\Delta$ returns to
itself (the factors $(z_i-z_k)$ and $(z_{i+1}-z_k)$ are exchanged), so $\Delta\circ\sigma_i$ winds once around $0$.
Hence $\Delta_*\colon B_N\to\pi_1(\C^*)=\Z$ is onto with $\sigma_i\mapsto1$, and the abelianization is $\Z$.
\end{proof}

\begin{example}[$N=2$]\label{TOP:ex:N2}
$\Conf(2,\C)\cong\C\times\C^*$ via $(z_1,z_2)\mapsto\big(\tfrac{z_1+z_2}2,w=z_1-z_2\big)$, and $\Sigma_2$
acts by $w\mapsto-w$. Squaring identifies $\C^*/\pm1$ with $\C^*$, so
$\Conf(2,\C)/\Sigma_2\cong\C\times\C^*$, $B_2\cong\Z$ generated by $\sigma_1$, and $P_2=2\Z$ generated by the
full twist $\sigma_1^2$ (Figure~\ref{TOP:fig:N2}).
\end{example}

\begin{figure}[ht]
\centering
\begin{tikzpicture}[scale=1.05]
  \begin{scope}
    \draw[->,gray] (-1.6,0)--(1.7,0); \draw[->,gray] (0,-1.4)--(0,1.5);
    \draw[thick,fill=white] (0,0) circle (2pt);
    \fill (1,0) circle (1.6pt) node[below right,font=\scriptsize] {$w$};
    \fill (-1,0) circle (1.6pt) node[below left,font=\scriptsize] {$-w$};
    \draw[addcol,very thick,-{Stealth}] (1,0) arc (0:180:1);
    \node[addcol,font=\scriptsize] at (0.95,1.0) {$\sigma_1$};
    \node[font=\small] at (0,-1.75) {$w=z_1-z_2\in\C^*$};
  \end{scope}
  \draw[-{Stealth},thick] (2.2,0)--(3.4,0) node[midway,above,font=\scriptsize] {$w\mapsto w^2$};
  \begin{scope}[xshift=5.4cm]
    \draw[->,gray] (-1.6,0)--(1.7,0); \draw[->,gray] (0,-1.4)--(0,1.5);
    \draw[thick,fill=white] (0,0) circle (2pt);
    \fill (1,0) circle (1.6pt) node[below right,font=\scriptsize] {$w^2$};
    \draw[addcol,very thick,-{Stealth}] (1,0) arc (0:355:1);
    \node[font=\small] at (0,-1.75) {$\Conf(2,\C)/\Sigma_2\simeq\C^*$};
  \end{scope}
\end{tikzpicture}
\caption{Two points in the plane. The half turn exchanging $z_1$ and $z_2$ is not a loop in
$\Conf(2,\C)$, but it becomes one after the points are unlabelled. Under the identification
$w\mapsto w^2$ it is a generator of $\pi_1(\C^*)$. This is $\sigma_1$, and the discriminant of $(z-z_1)(z-z_2)$ is $w^2$.}
\label{TOP:fig:N2}
\end{figure}

\subsection{Braids from reflection groups}

Now combine Lecture~7 with Lecture~8 and Interlude~A. The space $\Conf(N,\C)$ is the complement in $\C^N$ of the
hyperplanes $z_i=z_j$. These are the \emph{complexified} reflecting hyperplanes of the Weyl group
$W(A_{N-1})=\Sigma_N$. Their real parts cut $\R^N$ into the $N!$ contractible Weyl chambers of Lecture~8. Over $\C$
a hyperplane has real codimension $2$, so the complement is connected, but now loops can wind around the
walls. The Weyl group acts freely on the complement, and
\[
  \pi_1\Big(\big(\C^N\setminus\textstyle\bigcup_{i<j}\{z_i=z_j\}\big)\big/W\Big)=B_N .
\]
The generator $\sigma_i$ is a half turn around the wall $z_i=z_{i+1}$ of the fundamental chamber
$x_1<\dots<x_N$, that is, around the wall of the simple reflection $s_i$. Compare the two presentations:
\[
  W:\ \ s_i^2=1,\ \ (s_is_j)^{m_{ij}}=1;\qquad\qquad
  B_N:\ \ \underbrace{\sigma_i\sigma_j\sigma_i\cdots}_{m_{ij}}=\underbrace{\sigma_j\sigma_i\sigma_j\cdots}_{m_{ij}},
\]
with $m_{ij}=3$ for adjacent nodes of the Dynkin diagram $A_{N-1}$ and $m_{ij}=2$ otherwise. The braid
group is the Coxeter group with the relations $s_i^2=1$ deleted. This works for every diagram of
Lectures~4--5.

\begin{definition}
For a Coxeter matrix $(m_{ij})$, the \emph{Artin group} (or Artin--Tits group) is
$A=\langle\sigma_1,\dots,\sigma_r\mid\underbrace{\sigma_i\sigma_j\sigma_i\cdots}_{m_{ij}}=
\underbrace{\sigma_j\sigma_i\sigma_j\cdots}_{m_{ij}}\ (i\neq j)\rangle$. Adding $\sigma_i^2=1$ recovers the
Coxeter group.
\end{definition}

\begin{theorem}[Brieskorn 1971; Deligne 1972]\label{TOP:thm:BD}
Let $W$ be a finite real reflection group acting on $V=\R^r$, and let $Y_W$ be the complement in $V\otimes\C$
of the complexified reflecting hyperplanes. Then $W$ acts freely on $Y_W$ and $\pi_1(Y_W/W)$ is the Artin
group of the Coxeter diagram of $W$ (Brieskorn). Moreover $Y_W$, and hence $Y_W/W$, is aspherical (Deligne).
\end{theorem}

So every finite Coxeter diagram, including the $E_8$ of the densest packing in dimension $8$, has a braid group. For
$W(B_n)$ the quotient $Y_W/W$ is the space of $n$ unordered points in $\C^*$, and its Artin group is the braid
group of an annulus (Exercise~\ref{TOP:exc:Bn}). For \emph{infinite} Coxeter groups the analogous statement
is the \emph{$K(\pi,1)$ conjecture}, still open in general. For the affine groups, whose fundamental
alcoves appeared in Section~\ref{LIE:sec:torus} as configurations of points on a circle, it was proved by
Paolini and Salvetti (Theorem~\ref{TOP:thm:PS}).

\begin{remark}[Further directions]\leavevmode
\begin{itemize}[nosep]
\item \emph{Linearity.} Braid groups have faithful finite-dimensional representations, by Bigelow and by
  Krammer (Theorem~\ref{TOP:thm:linear}).
\item \emph{Anyons.} For identical quantum particles, exchanges are loops in the \emph{unlabelled}
  configuration space. In $\R^3$ this space has fundamental group $\Sigma_N$, which allows only bosons and fermions. In the plane
  the group is $B_N$, and a one-dimensional representation $\sigma_i\mapsto e^{i\vartheta}$ gives particles with
  any exchange phase. This was observed by Leinaas and Myrheim (Nuovo Cimento, 1977), and Wilczek named the
  particles \emph{anyons} (1982). Non-abelian representations of $B_N$ are the basis of proposals for
  topological quantum computation (Kitaev, 2003; Freedman--Kitaev--Larsen--Wang, Bull.\ AMS, 2003). The
  difference between labelled and unlabelled configurations that we met in Section~\ref{TOP:sec:mintype}
  here has physical consequences.
\end{itemize}
\end{remark}

%======================================================================
\section{Gallery of theorems}
\label{TOP:sec:gallery}

Each statement is followed by a sentence or two on why it belongs in this course.

\subsection*{Min-max and critical points}

\begin{theorem}[Lusternik--Fet 1951; Birkhoff 1917 for $\Sph^2$]
Every closed Riemannian manifold carries a non-constant closed geodesic.
\end{theorem}
\noindent The proof is min-max over a non-trivial homotopy class of maps into the free loop space, which
is the minimax principle of Theorem~\ref{TOP:thm:minimax} in infinite dimensions.

\begin{theorem}[Lusternik--Schnirelmann 1929; Grayson, Ann.\ of Math.\ 1989]
Every Riemannian metric on $\Sph^2$ has at least three simple closed geodesics.
\end{theorem}
\noindent This is the theorem behind Theorem~\ref{TOP:thm:LS}, with $\operatorname{cat}(\R P^2)=3$.
Franks (1992) and Bangert (1993) showed that there are always infinitely many closed geodesics, not
necessarily simple.

\begin{theorem}[Conley--Zehnder, Invent.\ Math.\ 1983]
Every Hamiltonian diffeomorphism of the standard symplectic torus $\T^{2n}$ has at least $2n+1$ fixed points,
and at least $2^{2n}$ if all fixed points are non-degenerate.
\end{theorem}
\noindent This is Arnold's conjecture for tori. The two numbers are exactly
$\operatorname{cat}(\T^{2n})$ and $\sum_\lambda b_\lambda(\T^{2n})$, the degenerate and the Morse counts of
Example~\ref{TOP:ex:torus}.

\begin{theorem}[Marques--Neves, Ann.\ of Math.\ 179 (2014)]\label{TOP:thm:willmore}
Every immersed torus $\Sigma$ in $\R^3$ satisfies
\[
  \mathcal W(\Sigma)=\int_\Sigma H^2\,dA\ \ge\ 2\pi^2,\qquad H=\tfrac12(\kappa_1+\kappa_2).
\]
More precisely, every embedded closed surface of genus $g\ge1$ in $\Sph^3$ has
$\int_\Sigma(1+H^2)\,dA\ge2\pi^2$, with equality if and only if $\Sigma$ is the Clifford torus up to
conformal transformations of $\Sph^3$.
\end{theorem}
\noindent This is the Willmore conjecture (1965). In $\R^3$ the extremal tori are, up to conformal
maps, the torus of revolution with radii in ratio $\sqrt2:1$. The proof is a min-max argument over a
five-parameter family (Section~\ref{TOP:sec:geodesics}).

\begin{theorem}[Marques--Neves, Invent.\ Math.\ 2017; Song, Ann.\ of Math.\ 2023]\label{TOP:thm:yau}
Every closed Riemannian manifold of dimension $3\le n+1\le7$ contains infinitely many smooth, closed,
embedded minimal hypersurfaces. Marques and Neves proved this under the assumption of positive Ricci curvature;
Song proved it in general.
\end{theorem}
\noindent This was conjectured by Yau (1982). The proofs use min-max over families of arbitrarily many
parameters, the analogue of the classes $c_1\le c_2\le\cdots$ in the proof of Theorem~\ref{TOP:thm:LS}.

\subsection*{Configuration spaces of hard particles}

\begin{theorem}[Baryshnikov--Bubenik--Kahle, IMRN 2014]\label{TOP:thm:BBK}
For hard balls of radius $r$ in a rectangular box in $\R^d$ whose shortest side has length $L$, the
critical points of the tautological function are the mechanically balanced configurations. There are no
critical values in $(0,L/2n)$, so $\Conf(n,r)\simeq\Conf(n)$ for $r<L/2n$, and the topology changes just above
$L/2n$.
\end{theorem}
\noindent This is the box version of Theorem~\ref{L10:thm:smalltheta} and of the threshold $2\pi/N$ on the sphere
quoted in Lecture~11. The critical configuration is a row of balls spanning the box, the analogue of
$N$ points equally spaced on a great circle.

\begin{theorem}[Alpert--Kahle--MacPherson, J.\ Appl.\ Comput.\ Topology 5 (2021)]\label{TOP:thm:AKM}
Let $C(n,w)$ be the space of $n$ non-overlapping discs of diameter $1$ in an infinite strip of integer width $w\ge2$, and fix
$j$. If $j\le w-2$, then $H_j(C(n,w))\cong H_j(\Conf(n,\R^2))$ and $b_j$ grows like $n^{2j}$. If $j\ge w-1$,
then $b_j$ grows exponentially in $n$. More precisely, writing $j=q(w-1)+r$ with $0\le r<w-1$,
$b_j\asymp(q+1)^n\,n^{qw+2r}$.
\end{theorem}
\noindent Exponential growth of Betti numbers is a topological counterpart of the exponential count of
jammed states in the channel of Lecture~11, where Fibonacci numbers appeared.

\begin{theorem}[P.~Ramsey, arXiv:2311.16776, 2023]\label{TOP:thm:thick}
For four discs of radius $r$ in the unit disc with $\tfrac14<r\le\tfrac13$, there is a non-trivial
element of $\pi_1$ of the configuration space that becomes trivial when each disc is replaced by its centre.
It is represented by the commutator $\sigma_1^{-1}\sigma_3^{-1}\sigma_1\sigma_3$. For $n\ge5$ discs and
$\tfrac1n<r\le\tfrac1{n-1}$, the configuration space contains a non-contractible $(n-3)$-sphere.
\end{theorem}
\noindent In $B_4$ the relation $\sigma_1\sigma_3=\sigma_3\sigma_1$ makes this commutator trivial. Thick discs have
no room to perform the two exchanges independently, so far commutation \emph{fails} for them. This is a preprint (P.~Ramsey, \emph{Thick braids and other
non-trivial homotopy in configuration spaces of hard discs}).

\begin{theorem}[Farber--Grant, Proc.\ AMS 2009]\label{TOP:thm:farber}
For $n,m\ge2$, $\operatorname{TC}(\Conf(n,\R^m))=2n-1$ if $m$ is odd and $2n-2$ if $m$ is even.
\end{theorem}
\noindent Farber and Yuzvinsky had proved the cases $m=2$ and $m$ odd. The lower bound is a cup-length bound in the spirit of
Theorem~\ref{TOP:thm:cuplength}, using products of ``zero-divisors'' $u\otimes1-1\otimes u$ in
$H^*(\Conf(n,\R^m)\times\Conf(n,\R^m))$.

\subsection*{Braid and Artin groups}

\begin{theorem}[Paolini--Salvetti, Invent.\ Math.\ 224 (2021)]\label{TOP:thm:PS}
The $K(\pi,1)$ conjecture holds for affine Artin groups: for every affine Coxeter group, the quotient of the
complexified complement of the reflection arrangement by the group is a classifying space for the
affine Artin group.
\end{theorem}
\noindent This extends Theorem~\ref{TOP:thm:BD} from the finite Weyl groups of Lectures~4--6 to their affine
versions, whose alcoves are the simplices of Interlude~A.

\begin{theorem}[Bigelow, J.\ AMS 2001; Krammer, Ann.\ of Math.\ 2002]\label{TOP:thm:linear}
The Lawrence--Krammer representation $B_n\to\GL_{n(n-1)/2}\big(\Z[q^{\pm1},t^{\pm1}]\big)$ is faithful for all
$n$. In particular braid groups are linear.
\end{theorem}
\noindent Bigelow's proof is topological: it uses the configuration space of \emph{two} points in the punctured disc.

\subsection*{Rigidity and optimality of packings}

\begin{theorem}[Sch\"utte--van der Waerden, Math.\ Ann.\ 1953; Musin, Ann.\ of Math.\ 2008]
The kissing number is $\kiss_3=12$ and $\kiss_4=24$.
\end{theorem}
\noindent The first settles the Newton--Gregory problem. In the language of these notes it says that
$\max\rho$ on $\Conf(13,\Sph^2)$ is less than $\pi/3$. Musin's proof modifies Delsarte's linear programming method.

\begin{theorem}[Tammes problem: Musin--Tarasov 2012 and 2015; Robinson 1961]
The Tammes problem is solved for $N=13$ (Musin--Tarasov, Discrete Comput.\ Geom.\ 48 (2012); optimal angle
$\approx57.1367^\circ$, unique up to isometry), $N=14$ (Musin--Tarasov, Experiment.\ Math.\ 24 (2015);
$\approx55.6706^\circ$) and $N=24$ (Robinson, Math.\ Ann.\ 1961; the snub cube).
\end{theorem}
\noindent The proofs for $13$ and $14$ enumerate the possible \emph{irreducible contact graphs} by computer.
These are the balanced graphs of Lecture~11, so each proof is the statement that one particular
critical configuration of $\rho$ is the top one. Together with the classical cases ($N\le12$, by Fejes
T\'oth, Sch\"utte--van der Waerden, Danzer and B\"or\"oczky), this covers $N\le14$ and $N=24$.

%======================================================================
\section*{Exercises}
\addcontentsline{toc}{section}{Exercises}

\begin{exercise}
For $f=\sin(\pi x)\sin(\pi y)\sin(\pi(x+y))$ on $\R^2/\Z^2$, show that the index of $\nabla f$ at the origin
is $-2$ (expand $f$ to third order) and check the Poincar\'e--Hopf formula $1+1-2=\chi(\T^2)$. Why can a
function on $\T^2$ with three critical points never be Morse?
\end{exercise}

\begin{exercise}
Using $H^*(\R P^n;\Z/2)=\Z/2[a]/(a^{n+1})$, show that $\operatorname{cat}(\R P^n)\ge n+1$. Deduce that every
smooth function on $\SO(3)\cong\R P^3$, for example a smooth function of the orientation of a rigid
polytope (Lecture~7), has at least four critical points. How many does a Morse function need?
\end{exercise}

\begin{exercise}\label{TOP:exc:N4}
Complete Example~\ref{TOP:ex:N4pass}. (a) Show that if four distinct points lie on a circle in $\Sph^2$, then
two of them are at angular distance at most $\pi/2$, with equality only for a square on a great circle.
(b) Verify the formula for $\rho(u(s))$ and that it is at least $\pi/2$.
\end{exercise}

\begin{exercise}
Compute $P(B\Conf(7))$ and check that the Betti numbers sum to $360$ and that $\chi=24$. If a Morse energy on
$B\Conf(7)$ has $c_0$ local minima, show that it has at least $c_0+13$ critical points of index $1$. (Use the
strong Morse inequality $c_1-c_0\ge b_1-b_0$.)
\end{exercise}

\begin{exercise}
Show that $\Delta=(\sigma_1\sigma_2\sigma_1)^2$ commutes with $\sigma_1$ and $\sigma_2$ in $B_3$, and that it
is the full twist of all three strands. Show that $(\sigma_1\sigma_2)^3=\Delta$. Under the abelianization of
Proposition~\ref{TOP:prop:abel}, what is the image of $\Delta$?
\end{exercise}

\begin{exercise}\label{TOP:exc:Bn}
The Weyl group $W(B_n)=(\Z/2)^n\rtimes\Sigma_n$ has reflecting hyperplanes $x_i=\pm x_j$ and $x_i=0$. Show
that the map $(z_1,\dots,z_n)\mapsto\{z_1^2,\dots,z_n^2\}$ identifies $Y_{W(B_n)}/W(B_n)$ with the space of
$n$ unordered distinct points in $\C^*$. Write down the Artin relations for the diagram $B_n$, with
$m=4$ on the double edge, and interpret the extra generator as a strand going around the puncture.
\end{exercise}

\begin{exercise}\label{TOP:exc:TC}
(a) Show that $\operatorname{TC}(X)=1$ if and only if $X$ is contractible. (b) Show that
$\operatorname{TC}(\Sph^1)=2$ by writing down two explicit motion planners. (c) Two robots move
without collision on a circular track. Identify the configuration space and compute its $\operatorname{TC}$.
\end{exercise}

\begin{exercise}
Show that $C(n,1)$, the space of $n$ discs of diameter $1$ in a strip of width $1$, has exactly $n!$ connected
components and that each is contractible. Compare with $\Conf(n,\R)$ in Lecture~7. (This is the case
$w=1$, excluded from Theorem~\ref{TOP:thm:AKM}; Alpert, Kahle and MacPherson record $b_0=n!$ and $b_j=0$
for $j\ge1$.)
\end{exercise}

\section*{Further reading}
\addcontentsline{toc}{section}{Further reading}
\begin{itemize}[nosep]
\item J.~Milnor, \emph{Morse Theory}, Annals of Math.\ Studies 51, 1963. Part~I is the basis of
  Lectures~8--10 and of Lemma~\ref{TOP:lem:deform}.
\item O.~Cornea, G.~Lupton, J.~Oprea, D.~Tanr\'e, \emph{Lusternik--Schnirelmann Category}, AMS Math.\ Surveys and
  Monographs 103, 2003. Category, cup-length, and critical point theory.
\item M.~Struwe, \emph{Variational Methods}, Springer. Min-max and the mountain pass in infinite dimensions.
\item F.~C.~Marques, A.~Neves, \emph{The Willmore conjecture}, arXiv:1409.7664. A survey of the proof.
\item Y.~Baryshnikov, P.~Bubenik, M.~Kahle, \emph{Min-type Morse theory for configuration spaces of hard
  spheres}, Int.\ Math.\ Res.\ Not.\ 2014(9), 2577--2592.
\item C.~Kassel, V.~Turaev, \emph{Braid Groups}, Springer GTM 247, 2008.
\item M.~Farber, \emph{Invitation to Topological Robotics}, EMS Zurich Lectures in Advanced Mathematics, 2008.
\item R.~Ghrist, \emph{Elementary Applied Topology}, 2014. Configuration spaces, robots, and Morse theory
  for applications.
\end{itemize}
% ===== end chapters/TOP.tex =====
\endgroup
\begingroup\renewcommand{\thechapter}{E}\renewcommand{\theHchapter}{interludeE}
% ===== begin chapters/SM.tex =====
% MAT670 -- Interlude E: random packings, entropy and information.
% Entirely new material (not in the handwritten notes). Labels prefixed SM:.
% Includes the coordinator's section chapters/tetra.tex (labels TET:) verbatim.

\chapter[Interlude E: Entropy and information]{Interlude E: Random packings, entropy and information}
\label{SM:ch}

\begin{addedblock}
This interlude is new: it is not in the handwritten notes. It picks up three threads that the course
left loose: the coin flips and the gas in a box of Lecture~7, the Edwards entropy and the Fibonacci
words of Lecture~11, and the project topics ``error correcting codes and information theory'' and
``entropy of jammed systems'' listed at the end of Lecture~6. To keep it readable it is set in black;
all of it should be read as added material.
\end{addedblock}

Most of this course was about the \emph{best} configurations: the densest packing, the largest
kissing configuration, the maximum of the injectivity radius. Lecture~7 already suggested a second
point of view. For large $N$ one does not ask for the best configuration but for the \emph{typical}
one, and one answers with Stirling's formula. Lecture~11 then did exactly this for jammed discs in a
channel, and found the golden ratio.

In this interlude we follow that second point of view: hard spheres as statistical mechanics with no
energy at all (Section~\ref{SM:sec:hard}); random and jammed packings (Section~\ref{SM:sec:random}); why
dense liquids do not simply crystallize (Sections~\ref{SM:sec:frustration}--\ref{SM:sec:curved}); Edwards'
statistical mechanics of the Lecture~11 channel, solved completely (Section~\ref{SM:sec:edwards}); and
information theory (Section~\ref{SM:sec:info}), where the channel's maximal Edwards entropy turns out to
be a Shannon capacity and $E_8$ and Leech come from perfect codes.

Throughout, $\ln$ is the natural logarithm and $\log_2$ the logarithm to base $2$. Boltzmann's constant
is set to $1$.

%======================================================================
\section{Hard spheres as statistical mechanics}
\label{SM:sec:hard}

Recall the remark at the end of \S\ref{L11:sec:boltzmann}: for hard particles the energy of a
configuration is $0$ if no two particles overlap and $+\infty$ if two do, so the Boltzmann weight
$e^{-\beta E}$ is the indicator function of the admissible configurations, at every temperature.
For $N$ balls of diameter $1$ with centres in a region $\Omega\subset\R^d$ of volume $V$, the
admissible configurations form the \emph{hard-sphere configuration space}
\[
  \Conf_{\ge1}(N,\Omega)=\{(x_1,\dots,x_N)\in\Omega^N:\ \abs{x_i-x_j}\ge1\ \text{for all }i\ne j\},
\]
a closed subset of the configuration space $\Conf(N,\Omega)$ of Lecture~7
(Definition~\ref{L07:def:conf}). On the sphere it is the space $\Conf(N,\theta)$ of
Definition~\ref{L10:def:conftheta}, whose topology we computed with min-type Morse theory. The
(configurational) \emph{partition function} is its volume,
\[
  Z_N(\Omega)=\frac1{N!}\vol\Conf_{\ge1}(N,\Omega),
\]
the factor $1/N!$ accounting for the indistinguishability of the particles. The \emph{entropy} is
$S=\ln Z_N$, and since the energy is identically zero the free energy is $F=-TS$. Temperature only
sets a scale. The \emph{pressure} is
\[
  \beta P=\frac{\partial S}{\partial V},
\]
the rate at which the configuration space grows when the box grows. (Change $\Omega$ by a homothety
to make this a derivative in one variable.) So for hard spheres \emph{everything is entropy}. This
is worth remembering in Section~\ref{SM:sec:edwards}, where Edwards runs the same dictionary in the
opposite direction and treats volume as if it were energy.

\subsection*{The Kepler problem is about where the free energy is finite}

Let $\dens_d$ be the maximal packing density of balls in $\R^d$, and let
\[
  \rho_{\mathrm{cp}}=\frac{\dens_d}{\vol\Ball^d(\tfrac12)}
\]
be the corresponding maximal \emph{number density} of centres for balls of diameter~$1$. In
dimension~$3$, Hales' theorem gives $\dens_3=\pi/\sqrt{18}$ and so $\rho_{\mathrm{cp}}=\sqrt2$, the
number density of the face-centred cubic lattice with nearest-neighbour distance~$1$.

\begin{proposition}\label{SM:prop:kepler}
Let $Q_\ell=[0,\ell]^d$, and let $N=N(\ell)$ with $N/\ell^d\to\rho$ as $\ell\to\infty$.
\begin{enumerate}[nosep]
\item If $N>\rho_{\mathrm{cp}}(\ell+1)^d$, then $Z_N(Q_\ell)=0$. In particular $Z_N(Q_\ell)=0$ for all
  large $\ell$ when $\rho>\rho_{\mathrm{cp}}$.
\item If $\rho<\rho_{\mathrm{cp}}$, then $\liminf_{\ell\to\infty}\frac1N\ln Z_N(Q_\ell)>-\infty$.
\end{enumerate}
\end{proposition}

\begin{proof}
(1) Given $N$ admissible centres in $Q_\ell$, the balls of radius $\frac12$ around them lie in the
cube $[-\frac12,\ell+\frac12]^d$ of side $\ell+1$. Repeat this cube periodically. The result is a
periodic packing of density $N\vol\Ball^d(\frac12)/(\ell+1)^d$, which is at most $\dens_d$.

(2) Choose $\rho<\rho''<\rho_{\mathrm{cp}}$. There is a periodic packing of balls of diameter $1$
with some number density $\rho'>\rho''$, since periodic packings come arbitrarily close to the optimal density
(Exercise~\ref{L1:ex:periodic-approx} of Lecture~1). Scale it by $\lambda=(\rho'/\rho'')^{1/d}>1$. The
centres now have number density $\rho''$ and are at mutual distance at least $\lambda$, so each centre may
be moved independently anywhere in a ball of radius $r=(\lambda-1)/2$ without creating an overlap.
The cube $Q_\ell$ contains $\rho''\ell^d-O(\ell^{d-1})$ of these wiggle balls, which is more than $N$ for
large $\ell$; keep $N$ of them. Assigning the $N$ labelled particles to the $N$ wiggle balls in any of
the $N!$ ways gives disjoint pieces of $\Conf_{\ge1}(N,Q_\ell)$, so
$Z_N\ge\frac1{N!}N!\,(\vol\Ball^d(r))^N$ and $\frac1N\ln Z_N\ge\ln\vol\Ball^d(r)$.
\end{proof}

The limit $s(\rho)=\lim\frac1N\ln Z_N$ exists for $\rho<\rho_{\mathrm{cp}}$ (subadditivity; see Ruelle's
book), so the free energy is finite at every density below $\rho_{\mathrm{cp}}$ and infinite at every
density above it. At $\rho_{\mathrm{cp}}$ itself one expects $s=-\infty$ and a divergent pressure, as for the
Tonks gas below, but the proposition does not decide this. Thus the theorems of Hales in $\R^3$, and of Viazovska and
Cohn--Kumar--Miller--Radchenko--Viazovska in $\R^8$ and $\R^{24}$, determine the density at which the
free energy of the hard-sphere gas becomes infinite.

\subsection*{One dimension: the Tonks gas}

On a line everything can be computed. Lecture~11 (\S\ref{L11:sec:hardrods}) found that the gaps
between $n$ rods of length $1$ in an interval of length $\ell$ form a simplex of volume
$(\ell-n)^n/n!$. Ordering the labels gives a factor $n!$, so
\[
  Z_n(\ell)=\frac{1}{n!}\cdot n!\cdot\frac{(\ell-n)^n}{n!}=\frac{(\ell-n)^n}{n!}.
\]

\begin{proposition}[Tonks 1936]\label{SM:prop:tonks}
The hard-rod gas has equation of state
\[
  \beta P=\frac{\partial}{\partial\ell}\ln Z_n(\ell)=\frac{n}{\ell-n}=\frac{\rho}{1-\rho},\qquad\rho=\frac n\ell,
\]
and entropy per particle $\frac1n\ln Z_n\to1+\ln\frac{1-\rho}{\rho}$ as $n,\ell\to\infty$ with
$n/\ell\to\rho\in(0,1)$.
\end{proposition}

The first formula is a one-line computation, and the second is Stirling's formula. Compare
Exercise~\ref{L11:ex:tonks}, where the same system was treated in the grand canonical ensemble: there
the density was $\rho=P/(1+P)$, which is the same equation of state solved for $P$. The pressure is a
real-analytic function of $\rho\in[0,1)$ and blows up at $\rho=1=\rho_{\mathrm{cp}}$, as
Proposition~\ref{SM:prop:kepler} predicts. There is no phase transition. That is not an accident of the
model.

\begin{theorem}[van Hove 1950]\label{SM:thm:vanhove}
Consider particles on a line interacting by a pair potential that has a hard core ($\phi(r)=+\infty$ for
$r<a$) and finite range ($\phi(r)=0$ for $r>R$), and is bounded on $[a,R]$. Then the thermodynamic
limit of the pressure is a real-analytic function of the temperature and the density: one-dimensional
systems of this kind have no phase transitions.
\end{theorem}

The proof uses a \emph{transfer operator}: by the hard core and the finite range, the partition function
is an iterated integral $\langle f,\mathcal T^ng\rangle$ of a positive integral operator, whose simple top
eigenvalue (Perron--Frobenius--Jentzsch) depends analytically on the parameters. Section~\ref{SM:sec:edwards}
has the same mechanism with a $2\times2$ matrix. (Cuesta and S\'anchez, J.~Stat.~Phys.\ 2004, discuss
which hypotheses are needed, and give one-dimensional models with transitions when they fail.)

\subsection*{Two and three dimensions: transitions seen only on the computer}

In $\R^3$ the hard-sphere gas is believed to \emph{freeze}. In 1957 Alder and Wainwright (molecular
dynamics) and Wood and Jacobson (Monte Carlo) published, back to back in J.~Chem.\ Phys.~27, the
first computer evidence of a fluid--solid transition. In the accepted picture (Hoover and Ree 1968) the
fluid is stable up to packing fraction about $0.494$, the close-packed crystal from about $0.545$, and
the two coexist in between. All of this is \emph{numerical}: that hard spheres crystallize, or have any
phase transition at all, is open in every dimension $d\ge2$. In $\R^2$ long-range positional order is
impossible (Mermin 1968 for smooth potentials; Richthammer 2016 for hard discs, see the gallery), and
what can survive is \emph{orientational} bond order, the hexatic phase. Bernard and Krauth (Phys.\ Rev.\
Lett.\ 107, 2011) found numerically that hard discs melt in two steps: a first-order liquid--hexatic
transition and a continuous hexatic--solid transition. Rigorous results are mostly at zero temperature,
where the problem becomes an optimization problem of the kind studied in this course (Theil's theorem in
the gallery). The first continuum model with finite-range interactions and a \emph{proved} phase
transition was the two-species Widom--Rowlinson model (Ruelle, Phys.\ Rev.\ Lett.\ 27, 1971), with hard
cores only between unlike particles.

%======================================================================
\section{Random packings}
\label{SM:sec:random}

\subsection{R\'enyi's parking problem}

Lecture~1 (Proposition~\ref{L1:prop:saturated-bound}) gave the easiest lower bound for packing
densities: a \emph{saturated} packing, one to which no ball can be added, has density at least
$2^{-d}$. The most natural random way of producing a saturated packing is to throw balls at random,
one at a time, and keep each one that does not overlap those already placed, until there is no room
left. This is \emph{random sequential adsorption} (RSA). In dimension one it is R\'enyi's \emph{parking
problem}: cars of length $1$ park at uniformly random free positions along a kerb of length $x$ until no
gap of length $\ge1$ remains.

Let $M(x)$ be the expected number of cars that park. Clearly $M(x)=0$ for $x<1$, and $M(x)=1$ for
$1\le x<2$. For $x\ge1$, condition on the position $t\in[0,x-1]$ of the first car, which is uniform. It
leaves two independent kerbs, of lengths $t$ and $x-1-t$:
\begin{equation}\label{SM:eq:renyi}
  M(x)=1+\frac1{x-1}\int_0^{x-1}\big(M(t)+M(x-1-t)\big)\,dt=1+\frac2{x-1}\int_0^{x-1}M(t)\,dt .
\end{equation}

\begin{theorem}[R\'enyi 1958]\label{SM:thm:renyi}
$M(x)/x\to m$ as $x\to\infty$, where
\[
  m=\int_0^\infty\exp\Big(-2\int_0^t\frac{1-e^{-u}}{u}\,du\Big)\,dt=0.7475979202\ldots
\]
\end{theorem}

\begin{proof}
Since the cars are disjoint subintervals of $[0,x]$, $0\le M(x)\le x$, so all Laplace transforms below
converge for $s>0$. Put $N(y)=M(y+1)$ for $y\ge0$ and multiply \eqref{SM:eq:renyi} by $y=x-1$:
\[
  yN(y)=y+2\int_0^yM(t)\,dt .
\]
Take Laplace transforms $\hat f(s)=\int_0^\infty e^{-sy}f(y)\,dy$. The transform of $yN(y)$ is
$-\hat N'(s)$, that of $\int_0^yM$ is $\hat M(s)/s$, and $\hat M(s)=e^{-s}\hat N(s)$ because $M$ vanishes on
$[0,1)$. So
\[
  \hat N'(s)+\frac{2e^{-s}}{s}\hat N(s)=-\frac1{s^2}.
\]
Solutions of the homogeneous equation are constant multiples of $\exp\big(2\int_s^\infty e^{-u}u^{-1}du\big)$,
which tends to a nonzero constant as $s\to\infty$, whereas $\hat N(s)\le\int_0^\infty e^{-sy}(y+1)\,dy\to0$. Hence
$\hat N$ is the unique decaying solution,
\[
  \hat N(s)=\int_s^\infty\frac{1}{t^2}\exp\Big(2\int_s^t\frac{e^{-u}}u\,du\Big)\,dt
  =\frac1{s^2}\int_s^\infty\exp\Big(-2\int_s^t\frac{1-e^{-u}}u\,du\Big)\,dt,
\]
where we wrote $e^{-u}=1-(1-e^{-u})$ and used $\exp(2\int_s^t du/u)=t^2/s^2$. As $s\to0^+$ the last
integral tends to $m$ by dominated convergence (for $t\ge1$ the integrand is at most
$\exp(-2\int_1^t\frac{1-e^{-u}}u du)=O(t^{-2})$). Hence $\hat M(s)=e^{-s}\hat N(s)\sim m/s^2$.

Now $U(x)=\int_0^xM$ is nondecreasing and $\int_0^\infty e^{-sx}dU(x)=\hat M(s)\sim m\,s^{-2}$, so Karamata's
Tauberian theorem (Feller, \emph{An Introduction to Probability Theory}, Vol.~II, \S XIII.5) gives
$U(x)\sim mx^2/2$. Finally \eqref{SM:eq:renyi} itself says $M(x)=1+2U(x-1)/(x-1)$, so $M(x)\sim mx$.
\end{proof}

\paragraph{Numerical check.} Quadrature gives $m=0.74759792025$ (R\'enyi, Publ.\ Math.\ Inst.\ Hung.\
Acad.\ Sci.\ 3, 1958). Twenty simulated kerbs of length $10^5$ gave
$0.74762\pm0.00012$. Solving \eqref{SM:eq:renyi} on a grid shows that $M(x)-mx$ settles down to
$m-1\approx-0.2524$ already for $x\approx10$. The density $m\approx0.748$ lies well above the saturation
bound $\frac12$ and well below~$1$.

\begin{figure}[ht]
\centering
\begin{tikzpicture}[scale=0.95]
  % 1D parking: first car splits the kerb
  \draw[thick] (0,0)--(7,0);
  \draw (0,-0.12)--(0,0.12) node[below=7pt,font=\small]{$0$};
  \draw (7,-0.12)--(7,0.12) node[below=7pt,font=\small]{$x$};
  \fill[flagcol!60] (2.6,0.05) rectangle (3.6,0.35);
  \draw (2.6,0.05) rectangle (3.6,0.35);
  \node[font=\scriptsize] at (3.1,0.6) {first car};
  \draw[<->] (0,-0.45)--node[below,font=\small]{$t$}(2.6,-0.45);
  \draw[<->] (3.6,-0.45)--node[below,font=\small]{$x-1-t$}(7,-0.45);
  \foreach \a in {0.3,1.45,4.1,5.2} {\fill[addcol!30] (\a,0.05) rectangle (\a+1,0.35); \draw (\a,0.05) rectangle (\a+1,0.35);}
  \node[font=\small,align=center] at (3.5,-1.6) {two independent smaller\\ parking problems};
  % 2D RSA picture
  \begin{scope}[xshift=8.0cm,yshift=-1.9cm,scale=0.5]
    \clip (0,0) rectangle (9,9);
    \foreach \p in {(5.63,8.07),(6.98,2.03),(2.70,7.86),(0.05,7.39),(9.05,7.39),(7.17,4.21),(2.73,2.51),(2.29,4.01),(4.54,4.98),(1.94,1.44),(5.51,0.40),(5.51,9.39),(0.32,4.63),(9.32,4.63),(4.20,8.26),(5.66,4.63),(4.47,2.23),(0.11,1.73),(9.11,1.73),(3.33,0.03),(3.33,9.03),(5.76,6.68),(3.25,5.38),(0.53,3.49),(6.09,1.36),(6.05,2.70),(7.87,5.96),(1.19,7.61),(5.13,3.39),(2.90,6.76),(7.86,3.10),(6.89,8.18),(3.77,7.34),(0.13,5.66),(9.13,5.66),(1.61,3.12),(6.68,5.23),(7.56,1.01),(0.27,0.18),(0.27,9.18),(9.27,0.18),(9.27,9.18),(1.69,5.10),(1.49,6.10),(7.76,7.17),(1.45,0.24),(1.45,9.24),(4.11,5.92),(3.40,3.78),(8.00,8.25),(-0.45,3.90),(8.55,3.90),(6.75,6.37),(6.25,3.79),(3.53,1.38),(4.72,6.71),(4.60,1.03),(5.32,5.65),(6.81,0.18),(6.81,9.19),(8.28,4.90),(8.03,2.01),(0.82,1.02),(0.95,2.32),(-0.14,2.70),(8.86,2.70),(2.48,0.57),(1.28,4.15),(8.21,0.25),(8.21,9.25),(0.60,6.55)} {\filldraw[fill=addcol!20,draw=addcol] \p circle (0.5);}
    \draw[thick] (0,0) rectangle (9,9);
  \end{scope}
\end{tikzpicture}
\caption{Left: the recursion behind R\'enyi's integral equation \eqref{SM:eq:renyi}. Right: a saturated
random sequential adsorption of $57$ discs of diameter $1$ in a periodic $9\times9$ square (computed for
this figure). Every point of the square is within distance $1$ of a centre, yet no two discs touch.}
\label{SM:fig:parking}
\end{figure}

In the plane, RSA of discs saturates at packing fraction about $0.547$ (Feder 1980). This is a
numerical value; Zhang and Torquato (Phys.\ Rev.\ E 88, 2013) give $0.5470735(28)$, and in $\R^3$
$0.3841307(21)$. (A quick simulation for this interlude, ten saturated periodic configurations of about
$1740$ discs each, gave $0.5468\pm0.0009$.) No closed form is known in any dimension $d\ge2$.

Notice what RSA does \emph{not} produce. In a saturated RSA configuration, with probability one no two
discs touch: the contact graph of Lectures~10--11 is empty. RSA packings are saturated but not
jammed; a small shake would let every disc move. To get mechanically stable random packings one has to
compress.

\subsection{Random close packing and the maximally random jammed state}

Pour ball bearings into a container and shake it: the packing fraction settles near $0.64$. This
\emph{random close packing} (RCP) value goes back to the experiments of Bernal and of Scott (both in
\emph{Nature}, 1960); Scott and Kilgour (1969) measured about $0.637$. It is well below the crystal value
$0.7405$, and remarkably reproducible.

Torquato, Truskett and Debenedetti (\emph{Is random close packing of spheres well defined?}, Phys.\
Rev.\ Lett.\ 84, 2000) argued that it is not: by compressing faster or slower one trades a little
disorder for a little density, with no canonical place to stop. They proposed instead the
\emph{maximally random jammed} (MRJ) state, the jammed packing minimizing a chosen order parameter,
again at about $0.64$ (numerical). This needs a definition of ``jammed''; Torquato and Stillinger
(J.~Phys.\ Chem.~B 105, 2001) give a hierarchy. For a packing in a container with rigid walls:
\begin{itemize}[nosep]
\item \emph{locally jammed}: no single ball can be moved while all the others are held fixed;
\item \emph{collectively jammed}: locally jammed, and no subset of balls can be moved
  collectively, the walls being fixed;
\item \emph{strictly jammed}: collectively jammed, also under small volume-nonincreasing
  deformations of the walls (for periodic packings: of the lattice).
\end{itemize}
For instance, with suitable walls, the simple cubic packing is collectively but not strictly jammed (shearing the box lets it
slide), FCC is strictly jammed, and the honeycomb disc packing is locally but not collectively jammed.
(In a torus there are no walls to stop a row, and the square and simple cubic packings are not even
collectively jammed; see Figure~\ref{SM:fig:maxwell}.)
To first order, a ball with the others held fixed is stuck exactly when the directions to its contacts
do not all lie in a closed hemisphere. In particular a locally jammed ball has at least $d+1$ contacts.

\subsection{Jamming, force balance and Maxwell counting}

Collective jamming is the version of Lectures~10--11 that sees all balls at once. Take $N$ balls with
centres $x_i$ and radii $r_i$ in a flat torus $\R^d/\Lambda$ (periodic boundary conditions), with contact
set $E=\{\{i,j\}:\abs{x_i-x_j}=r_i+r_j\}$, and let $u_{ij}=(x_j-x_i)/\abs{x_j-x_i}$. A velocity
$v=(v_1,\dots,v_N)\in(\R^d)^N$ changes the contact distances at the rates
\[
  (Av)_{ij}=\ip{u_{ij}}{v_j-v_i},\qquad\{i,j\}\in E.
\]
The linear map $A\colon(\R^d)^N\to\R^E$ is the \emph{rigidity matrix} of the contact graph. To first order a
motion keeps the balls disjoint if and only if $Av\ge0$ componentwise. Translations $v_i=c$ always
satisfy $Av=0$; call them \emph{trivial}.

\begin{definition}\label{SM:def:jam}
The packing is \emph{first-order jammed} if every $v$ with $Av\ge0$ is trivial.
\end{definition}

For sphere packings, first-order jamming is equivalent to collective jamming (with the lattice fixed):
a sphere packing is rigid if and only if it is infinitesimally rigid. This is a theorem of Connelly (1988)
for packings in a container, extended to periodic packings by Donev, Torquato, Stillinger and Connelly
(J.~Comput.\ Phys.\ 197, 2004). It is special to spheres; it fails for ellipsoids.

\begin{proposition}\label{SM:prop:jam}
A packing is first-order jammed if and only if
\begin{enumerate}[nosep]
\item $\ker A$ consists of the trivial motions (the contact graph is infinitesimally rigid as a bar
  framework), and
\item there are \emph{strictly positive} weights $w_{ij}>0$ on all contacts with
  $\sum_jw_{ij}u_{ij}=0$ at every ball $i$ (a positive self-stress: the forces balance at every ball).
\end{enumerate}
\end{proposition}

\begin{proof}
If (1) and (2) hold and $Av\ge0$, then
\[
  0=\sum_i\Big\langle v_i,\sum_jw_{ij}u_{ij}\Big\rangle=-\sum_{\{i,j\}\in E}w_{ij}\ip{u_{ij}}{v_j-v_i}=-\ip{w}{Av},
\]
a sum of non-positive terms, so $Av=0$, and $v$ is trivial by (1). Conversely, (1) is necessary because
$Av=0$ implies $Av\ge0$. For (2), Stiemke's alternative (a relative of Gordan's alternative used in the
proof of Theorem~\ref{L11:thm:forcebalance}) says that either there is $v$ with $Av\ge0$ and $Av\ne0$, or
there is $w>0$ with $A^Tw=0$. In a first-order jammed packing the first option is excluded, and $A^Tw=0$
is exactly the balance condition.
\end{proof}

So the jammed packings are the force-balanced contact graphs of Lecture~10 with two upgrades: the
weights are positive on \emph{every} contact, and the graph is rigid. (Compare
Corollary~\ref{L10:cor:topchange}, where critical configurations of the min-distance function only
needed weights $\ge0$, not all zero.) The rigidity condition already forces a lot of contacts.

\begin{corollary}[Maxwell counting]\label{SM:cor:maxwell}
A first-order jammed packing of $N$ balls in a $d$-dimensional torus has at least $dN-d+1$ contacts, so the
mean number of contacts per ball satisfies
\[
  z=\frac{2\abs E}{N}\ \ge\ 2d-\frac{2(d-1)}{N}.
\]
\end{corollary}

\begin{proof}
By (1), $\ker A$ has dimension $d$, so $\operatorname{rank}A=dN-d$. By (2), $A^T$ has a nonzero kernel, so
$\abs E\ge\operatorname{rank}A+1=dN-d+1$.
\end{proof}

This is Maxwell's 1864 count for frames. Conversely, for generic radii each contact is one equation
$\abs{x_i-x_j}=r_i+r_j$ on $dN-d$ free coordinates plus one free parameter, the common scale of the radii
at which the packing jams, and generic equations should not be overdetermined.
Together the counts give $z=2d$: jammed frictionless packings should be \emph{isostatic}. The upper bound
is a theorem in the plane (Connelly--Gortler--Solomonides--Yampolskaya, gallery) and open for $d\ge3$, but
it holds very precisely in simulations. O'Hern, Silbert, Liu and Nagel (Phys.\ Rev.~E 68, 2003), compressing
soft repulsive spheres to the threshold ``point~J'' where they first jam, found $z=3.98$ in $2$d and $5.98$ in
$3$d (both $\pm0.02$, rattlers removed) and a $3$d monodisperse threshold $\phi^*=0.639\pm0.001$,
essentially the RCP value.

\begin{figure}[ht]
\centering
\begin{tikzpicture}[scale=0.55]
  \foreach \i in {0,...,5} \foreach \j in {0,1,3} {\draw (\i,\j) circle (0.5);}
  \foreach \i in {0,...,5} {\filldraw[fill=addcol!20,draw=addcol] (\i+0.35,2) circle (0.5);}
  \draw[-{Stealth},thick,flagcol] (5.6,2)--(6.6,2);
  \draw[dashed,gray] (-0.6,-0.6) rectangle (5.6,3.6);
  \node[right,font=\small,align=left] at (7.0,1.5) {square lattice in a torus: $z=4=2d$,\\
    but the shaded row slides; to first order\\ $Av=0$ for this motion, so it is not jammed};
\end{tikzpicture}
\caption{Maxwell's count is necessary, not sufficient. In the square packing every disc has four
contacts, yet a whole row can slide along itself: its vertical contacts open only to second order.}
\label{SM:fig:maxwell}
\end{figure}


%======================================================================
\section{Geometric frustration}
\label{SM:sec:frustration}

Why does a dense liquid not simply crystallize? In 1952 F.~C.~Frank (\emph{Supercooling of liquids},
Proc.\ R.\ Soc.\ Lond.\ A 215) pointed out that the question has a \emph{local} answer. Put twelve atoms
around a thirteenth. There are three natural candidates: the cuboctahedron (the FCC neighbourhood), its
twisted relative from HCP, and the icosahedron. Lecture~11 met the first and the last, which the jitterbug
connects. For
Lennard-Jones atoms Frank computed that the icosahedral cluster has the lowest energy of the three: it
has $30$ near-neighbour pairs among the outer twelve, against $24$. But icosahedral symmetry is
incompatible with a lattice, so a liquid of locally icosahedral clusters must rearrange substantially
before it can crystallize, and it can be supercooled. This is the first instance of \emph{geometric
frustration}: the locally preferred structure cannot be extended to fill space.

For hard spheres the same tension appears as a statement about densities. The densest \emph{local}
arrangement around one ball beats the densest global one. Fejes T\'oth conjectured, and Hales and
McLaughlin proved (J.~Amer.\ Math.\ Soc.\ 23, 2010), that in any packing of unit balls in $\R^3$ every
Voronoi cell (Lecture~1, Definition~\ref{L1:def:voronoi}) has volume at least that of the regular
dodecahedron of inradius $1$. That cell would give density
\[
  \frac{\vol\Ball^3}{\vol(\text{dodecahedron of inradius }1)}\approx0.7547\ >\ \frac{\pi}{\sqrt{18}}\approx0.7405 .
\]
The dodecahedral cell is the Voronoi cell of a ball surrounded by an icosahedron of twelve touching balls,
and regular dodecahedra do not tile space. The local optimum is not the global one.

Frank's icosahedron is often described as twenty regular tetrahedra, the densest arrangement of four
balls, glued around the centre. That description is slightly wrong, and the error is the heart of the
matter.

% ===== begin chapters/tetra.tex =====
% Section for Interlude E: why 20 regular tetrahedra do not make a regular icosahedron.
% Written to follow W. Kusner's argument. Labels prefixed TET:.

\section{Twenty tetrahedra do not make an icosahedron}
\label{TET:sec}

Frank's icosahedral cluster (and the jitterbug of Lecture~11) rests on a small geometric fact: a regular
icosahedron is \emph{not} made of $20$ regular tetrahedra. Cone the icosahedron from its centre and
you get $20$ congruent tetrahedra, but they are slightly flattened. The edges from the centre are a
little \emph{shorter} than the edges of the icosahedron:
\[
  \frac{\text{circumradius}}{\text{edge}}=\sin\frac{2\pi}{5}\approx0.951<1 .
\]
In terms of balls: twelve unit balls placed at the vertices of an icosahedron around a central unit
ball can all touch the central ball, but then they do not touch each other. This slack is why the
twelve outer balls can roll around (the kissing configuration in $\R^3$ is not rigid), and it is the
root of the Newton--Gregory dispute over whether a thirteenth ball fits.

The usual proof is a computation, for instance with the dihedral angle $\arccos\frac13$ of the regular
tetrahedron. The following argument needs no computation. It uses only the plane fact that \emph{six}
equal circles fit exactly around a seventh, together with the triangle inequality on the sphere.

\begin{proposition}\label{TET:prop}
A regular icosahedron is not the union of $20$ regular tetrahedra sharing a vertex at its centre.
\end{proposition}

\paragraph{The argument.} Suppose it were, with edge length $a$. Put balls of radius $a/2$ at the centre
and at the $12$ vertices. Then the central ball touches all twelve, and each vertex ball touches its
five neighbours. Project the twelve contact points radially onto a sphere about the centre, and
measure distances by angle. Neighbouring vertices are then exactly $60^\circ$ apart, because the centre
and two neighbours form an equilateral triangle.

Stand the icosahedron on a vertex: a north pole $N$, the antipodal vertex $S$, and two
\emph{tropical rings} of five vertices each. The upper ring consists of the neighbours of $N$, the
lower ring of the neighbours of $S$, and the rings are staggered by $36^\circ$ of longitude. The equatorial
band between the rings is made of ten triangles (Figure~\ref{TET:fig:band}, left). Every vertex of the
upper ring is $60^\circ$ from $N$, so the upper ring lies at latitude $30^\circ$\,N. Likewise the lower ring
lies at latitude $30^\circ$\,S.

Now \emph{rotate the lower ring} by $36^\circ$ about the axis $NS$, keeping its latitude. Each lower vertex
$\ell$ now sits directly below an upper vertex $u$, and the ten equatorial triangles have become five
equatorial quadrilaterals, the ``squares'' (Figure~\ref{TET:fig:band}, right). Each of the five
meridians through $N$ and $S$ now carries four of our points,
\[
  N,\ u,\ \ell,\ S\qquad\text{at angles}\qquad 0^\circ,\ 60^\circ,\ 120^\circ,\ 180^\circ
\]
along the meridian, since $u$ is $60^\circ$ from $N$, $\ell$ is $60^\circ$ from $S$, and $N$ and $S$ are
$180^\circ$ apart. The meridian is a great circle, and $60^\circ$ steps are exactly the spacing of a
\emph{ring of six} balls around the central one, as in the hexagonal packing
(Figure~\ref{TET:fig:meridian}). So after the rotation, $u$ and $\ell$ touch: the five meridians have
become five half-rings of six, one great circle each.

But before the rotation $u$ and $\ell$ also touched, since they were neighbours in the icosahedron, and
in between we only slid $\ell$ sideways along its latitude circle. The closest point of a latitude
circle to $u$ is the one on $u$'s meridian, so sliding $\ell$ off the meridian moves it strictly
\emph{farther} from $u$. In other words, the deformation from ten equatorial triangles to five
equatorial squares, made without changing any latitude, shortened the distance from $u$ to $\ell$. Yet
it is $60^\circ$ at both ends. That is a contradiction. \qed

\medskip
In one line, the same argument is the triangle inequality on the sphere:
\[
  d(u,\ell)\ \ge\ d(N,\ell)-d(N,u)\ =\ (180^\circ-60^\circ)-60^\circ\ =\ 60^\circ,
\]
with equality only if $u$ lies on the arc from $N$ to $\ell$, that is, on $\ell$'s meridian. Each upper
vertex $u$ has two lower neighbours, and they cannot both lie on $u$'s meridian, so for at least one of
them $d(u,\ell)>60^\circ$. (In the icosahedron both are off the meridian, by $\pm36^\circ$.) The ring of six
is the picture of the equality case.

\begin{figure}[ht]
\centering
\begin{tikzpicture}[scale=0.95]
  % left: antiprism band (10 triangles), right: prism band (5 squares)
  \begin{scope}
    \foreach \k in {0,...,5} {\coordinate (U\k) at (\k,1); }
    \foreach \k in {0,...,4} {\coordinate (L\k) at (\k+0.5,0); }
    \coordinate (Lm) at (-0.5,0);
    \draw[thick] (U0)--(U5); \draw[thick] (Lm)--(L4);
    \foreach \k [evaluate=\k as \j using int(\k+1)] in {0,...,4} {\draw (U\k)--(L\k)--(U\j);}
    \draw (U0)--(Lm);
    \foreach \k in {0,...,5} \fill (U\k) circle (0.06);
    \foreach \k in {0,...,4} \fill[flagcol] (L\k) circle (0.06);
    \node[font=\small] at (2.5,1.55) {upper ring, latitude $30^\circ$\,N};
    \node[font=\small] at (2.5,-0.55) {lower ring, latitude $30^\circ$\,S};
    \node[font=\small,gray] at (2.5,-1.1) {ten triangles (icosahedron)};
  \end{scope}
  \draw[-{Stealth},thick] (5.8,0.5)--(6.9,0.5) node[midway,above,font=\scriptsize]{rotate};
  \node[font=\scriptsize] at (6.35,0.2) {lower ring};
  \begin{scope}[xshift=7.4cm]
    \foreach \k in {0,...,5} {\coordinate (U\k) at (\k,1); \coordinate (L\k) at (\k,0);}
    \draw[thick] (U0)--(U5); \draw[thick] (L0)--(L5);
    \foreach \k in {0,...,5} \draw (U\k)--(L\k);
    \foreach \k in {0,...,5} {\fill (U\k) circle (0.06); \fill[flagcol] (L\k) circle (0.06);}
    \node[font=\small,gray] at (2.5,-1.1) {five squares, same latitudes};
  \end{scope}
\end{tikzpicture}
\caption{The equatorial band, unrolled (the left and right ends are identified; longitudes in units of
$72^\circ$). Rotating the lower ring by $36^\circ$ turns the ten equatorial triangles into five
quadrilaterals without changing any latitude.}
\label{TET:fig:band}
\end{figure}

\begin{figure}[ht]
\centering
\begin{tikzpicture}[scale=0.62]
  \fill[gray!15] (0,0) circle (1); \draw (0,0) circle (1); \node[font=\scriptsize] at (0,0) {$O$};
  \foreach \a/\lab in {90/N,30/u,-30/\ell,-90/S}
    {\filldraw[fill=addcol!15,draw=addcol,thick] (\a:2) circle (1); \node[font=\small] at (\a:2) {$\lab$};}
  \foreach \a in {150,210} {\draw[dashed,gray] (\a:2) circle (1);}
  \draw[gray,thin] (0,0) circle (2);
  \foreach \a in {90,30,-30,-90} \draw[gray,thin] (0,0)--(\a:2);
  \node[font=\scriptsize] at (60:0.65) {$60^\circ$};
  \node[font=\scriptsize] at (0:0.7) {$60^\circ$};
  \node[font=\scriptsize] at (-60:0.65) {$60^\circ$};
\end{tikzpicture}
\caption{After the rotation, each meridian is half of a ring of six balls around the central ball: $N$,
$u$, $\ell$, $S$ are consecutive at $60^\circ$ steps. The dashed balls complete the ring.}
\label{TET:fig:meridian}
\end{figure}

\begin{remark}[What the argument shows, quantitatively]
In the hypothetical configuration with rings at latitudes $\pm30^\circ$, an upper and a lower vertex
whose longitudes differ by $36^\circ$ are at angular distance $\arccos\big(\tfrac34\cos36^\circ-\tfrac14\big)
\approx69.1^\circ$, not $60^\circ$. In the true icosahedron the rings sit lower, at latitudes
$\pm\arctan\frac12\approx\pm26.6^\circ$, every edge subtends $\arctan2\approx63.4^\circ$ at the centre, and
the gaps between neighbouring outer balls are what let the twelve balls roll.
\end{remark}

\begin{remark}[Two related classical facts]
(1) Regular tetrahedra do not even fill the space around an \emph{edge}: the dihedral angle is
$\arccos\frac13\approx70.53^\circ$, and five such angles leave a gap of about $7.36^\circ$. (Aristotle
claimed that regular tetrahedra tile space; see Lagarias--Zong, \emph{Mysteries in packing regular
tetrahedra}, Notices AMS 59 (2012).) The same angle appears here: five points at $60^\circ$ from $N$ and
$60^\circ$ from each other would need angles of $\arccos\frac13$ at $N$, and five of these do not close up.
(2) In curved space the frustration disappears. In the $3$-sphere, $600$ regular tetrahedra (with
the right curvature) fit together exactly, five around each edge, forming the regular polytope
$\{3,3,5\}$ with $120$ vertices, each surrounded by a regular icosahedron. This is the ``ideal''
icosahedral order of Frank's liquids; the flat-space packing problem has to compromise.
\end{remark}

\begin{remark}[Source]
\unclear{Wording of the attribution is a placeholder for the author to edit.}
This argument came out of a conversation between the author and T.~C.~Hales.
\end{remark}
% ===== end chapters/tetra.tex =====

%======================================================================
\section{Ideal order in curved space}
\label{SM:sec:curved}

The previous section measured the frustration: five regular tetrahedra around an edge leave a gap of
about $7.36^\circ$. As noted there, the gap closes in positive curvature. The dihedral angle of a regular
spherical tetrahedron grows with its size from the flat value $\arccos\frac13\approx70.53^\circ$, and for
edge $36^\circ$ on the unit $\Sph^3$ it is exactly $72^\circ=360^\circ/5$. Six hundred such tetrahedra are the cells of the
regular polytope $\{3,3,5\}$ (Coxeter, \emph{Regular Polytopes}), whose $120$ vertices are at angular
distance at least $36^\circ$, each with $12$ nearest neighbours forming a regular icosahedron. Balls of
angular radius $18^\circ$ at these points realize Frank's local structure \emph{everywhere}. By a theorem of
Cohn and Kumar (gallery) these $120$ points even minimize every completely monotonic pair potential on
$\Sph^3$: the ideal order is a zero-temperature ground state, though in curved space.

This is the starting point of the curved-space approach to liquids and glasses of Sadoc, Mosseri and
others, and of Nelson (\emph{Order, frustration, and defects in liquids and glasses}, Phys.\ Rev.~B 28,
1983): a dense flat liquid is the ideal $\{3,3,5\}$ structure ``flattened'', with the curvature
concentrated on line defects (disclinations) along which six rather than five tetrahedra share an edge.
See Sadoc and Mosseri, \emph{Geometrical Frustration} (Cambridge, 1999).

%======================================================================
\section{Edwards entropy}
\label{SM:sec:edwards}

\subsection{The ansatz and its tests}

A pile of sand is not at thermal equilibrium; $k_BT$ is negligible next to the energy needed to lift a
grain. Yet tapped or sheared granular packings settle into reproducible states, and Edwards and
Oakeshott (\emph{Theory of powders}, Physica~A 157, 1989) proposed a statistical mechanics for them in
which the \emph{volume} plays the role of energy. This is Definition~\ref{L11:def:edwards}: all jammed
(mechanically stable) states of the same volume are equally probable. Writing $\Omega(V)$ for the number
(or a suitable measure) of jammed states of volume $V$, the dictionary is
\begin{center}
\renewcommand{\arraystretch}{1.2}
\begin{tabular}{ll}
\toprule
thermal statistical mechanics & Edwards' granular ensemble\\
\midrule
energy $E$ & volume $V$\\
all states of energy $E$ & jammed states of volume $V$\\
entropy $S(E)=\ln\Omega(E)$ & Edwards entropy $S(V)=\ln\Omega(V)$\\
temperature $T=\partial E/\partial S$ & compactivity $X=\partial V/\partial S$\\
Boltzmann weight $e^{-E/T}$ & weight $e^{-V/X}$\\
\bottomrule
\end{tabular}
\end{center}
The compactivity measures how much more compact the packing could be. As $X\to0^+$ the weight
concentrates on the densest jammed states, and as $X\to\infty$ on the most numerous ones. Compare
Section~\ref{SM:sec:hard}: for the thermal hard-sphere gas $\partial S/\partial V=\beta P$, while for
Edwards $\partial S/\partial V=1/X$.

In the continuum, jammed packings are isolated points of configuration space, so ``equal
probability'' needs a measure, and a real protocol samples a jammed state with probability equal to the
volume of its \emph{basin of attraction}, as in the $\rho$-evolution of Lecture~11. Edwards' ansatz then
says that basins of equal-volume states are equally large; as noted in Lecture~11, small disc packings
violate this. The most careful test so far (Martiniani, Schrenk, Ramola, Chakraborty and Frenkel,
\emph{Nature Physics} 13, 2017) computed basin volumes by thermodynamic integration for $N=64$
polydisperse discs in $2$d (hard cores with soft repulsive shells), with $N=32$ to $128$ for finite-size
scaling. Above jamming the basin volumes are broadly distributed and the ansatz fails; the distribution
narrows towards the unjamming point, extrapolated to $\phi^*=0.841(3)$, where all packings become
equally probable (numerically). Baule, Morone, Herrmann and Makse (Rev.\ Mod.\ Phys.\ 90, 2018) review
the theory and its tests.

Where Edwards' ensemble can be written down \emph{exactly}, its predictions can be checked against the
ansatz-free dynamics. The channel of Lecture~11 is the simplest such model, and we now solve it.

\subsection{The channel model solved}
\label{SM:sec:channel}

Recall the setting of \S\ref{L11:sec:channel}. Discs of radius $1$ lie in a channel of width $w=2+h$
with $0<h<\sqrt3$, so that each disc touches at most its two neighbours. Two consecutive discs in
contact at opposite walls are horizontally $L=\sqrt{4-h^2}\in(1,2)$ apart, and two at the same wall are
$2$ apart. A jammed configuration of $n=\ell+1$ discs is recorded by its contact word
$\sigma\in\{0,1\}^\ell$ with no two adjacent $1$'s ($1$ for same wall, $0$ for opposite walls). A word
with $k$ ones has length
\[
  t(\sigma)=2k+(\ell-k)L+2=\ell L+(2-L)k+2,
\]
and the number of such words is
\[
  \Omega(\ell,k)=\binom{\ell-k+1}{k},
\]
or $\binom{\ell-k-1}{k}$ if the end discs are restricted, as flagged in Lecture~11.
We take the volume to be the channel length, $V=t$ (multiply by $w$ for the area). Since $t$ is an affine
function of $k$, fixing the volume is the same as fixing the number $k$ of defects. Write
$x=k/\ell\in[0,\frac12]$ for the defect density, and use $\eta$ for the packing fraction (the letter
$\varphi$ is reserved for the golden ratio).

\begin{proposition}[Edwards entropy of the channel]\label{SM:prop:entropy}
As $\ell\to\infty$ with $k/\ell\to x\in[0,\frac12]$,
\[
  \frac1\ell\ln\Omega(\ell,k)\ \longrightarrow\ s(x)=(1-x)\ln(1-x)-x\ln x-(1-2x)\ln(1-2x),
\]
for both counts. The function $s$ is strictly concave on $[0,\frac12]$, vanishes at both ends, and has
its maximum $s(x^*)=\ln\varphi$ at $x^*=\frac{5-\sqrt5}{10}$. The volume per contact is
$v(x)=L+(2-L)x$ and the packing fraction is
\[
  \eta(x)=\frac{\pi}{w\,v(x)}=\frac{\pi}{w\,(L+(2-L)x)} .
\]
\end{proposition}

\begin{proof}
Stirling's formula in the form $\ln\binom ab=a\,H\!\big(\tfrac ba\big)+O(\ln a)$, with
$H(p)=-p\ln p-(1-p)\ln(1-p)$, applied with $a=\ell(1-x)+O(1)$ and $b=\ell x$, gives the limit
$(1-x)H\big(\frac x{1-x}\big)$. Expanding, with $1-\frac x{1-x}=\frac{1-2x}{1-x}$, this is $s(x)$. The
shift $O(1)$ in $a$ does not change the limit, so the two counts agree to leading order. Next,
\[
  s'(x)=\ln\frac{(1-2x)^2}{x(1-x)},\qquad s''(x)=\frac1{1-x}-\frac1x-\frac4{1-2x}<0,
\]
since $\frac1{1-x}\le2<\frac4{1-2x}$ on $[0,\frac12)$. The critical point solves $(1-2x)^2=x(1-x)$, that is
$5x^2-5x+1=0$, as in Lecture~11. At $x^*$ one has $1-x^*=\frac{5+\sqrt5}{10}=\frac{\varphi}{\sqrt5}$,
$x^*=\frac{1}{\sqrt5\,\varphi}$ and $1-2x^*=\frac{1}{\sqrt5}$, so that
$s(x^*)=\frac\varphi{\sqrt5}\ln\frac\varphi{\sqrt5}-\frac1{\sqrt5\varphi}\ln\frac1{\sqrt5\varphi}-\frac1{\sqrt5}\ln\frac1{\sqrt5}$.
Since $\frac\varphi{\sqrt5}-\frac1{\sqrt5\varphi}-\frac1{\sqrt5}=0$, the terms in $\ln\sqrt5$ cancel and what is
left is $\big(\frac\varphi{\sqrt5}+\frac1{\sqrt5\varphi}\big)\ln\varphi=\ln\varphi$, using
$\varphi+\varphi^{-1}=\sqrt5$. The formula for $v$ is $t/\ell$ up to $O(1/\ell)$, and $\eta$ is the area
$\pi n$ of the discs divided by the area $wt$ of the channel, in the limit.
\end{proof}

The densest jammed state is the zigzag, $x=0$, with $\eta_{\max}=\pi/(wL)$; the loosest is
$x=\frac12$, with $\eta_{\min}=2\pi/(w(L+2))$. The zigzag word $00\cdots0$ is unique, the loosest word
$1010\cdots1$ is unique for odd $\ell$ (for even $\ell$ there are $\frac\ell2+1$ words with $\frac\ell2$
ones), and both have Edwards entropy zero. For $h=1.2$, the channel of Figure~\ref{L11:fig:channel}, one has $L=1.6$, $w=3.2$ and
\[
  \eta_{\max}\approx0.6136,\qquad \eta(x^*)\approx0.5739,\qquad \eta_{\min}\approx0.5454 .
\]

\paragraph{Compactivity and the equation of state.} Changing $k$ by one changes $V$ by $2-L$ and $\ln\Omega$
by $s'(x)+o(1)$, so
\begin{equation}\label{SM:eq:compactivity}
  \frac1X=\frac{\partial S}{\partial V}=\frac{s'(x)}{2-L}=\frac1{2-L}\ln\frac{(1-2x)^2}{x(1-x)} .
\end{equation}

\begin{proposition}[Equation of state]\label{SM:prop:eos}
For $X\in(-\infty,0)\cup(0,\infty]$, the defect density at compactivity $X$ is
\[
  x(X)=\frac12\Big(1-\frac1{\sqrt{1+4e^{-(2-L)/X}}}\Big),
\]
and $\eta(X)=\pi/\big(w(L+(2-L)x(X))\big)$. In particular:
\begin{itemize}[nosep]
\item as $X\to0^+$, $x\sim e^{-(2-L)/X}\to0$ and $\eta\to\eta_{\max}$ (densest);
\item $X=\pm\infty$ corresponds to $x=x^*$, the most numerous states;
\item the states looser than typical, $x^*<x<\frac12$, have \emph{negative} compactivity, and as $X\to0^-$,
  $x\to\frac12$ and $\eta\to\eta_{\min}$ (loosest).
\end{itemize}
\end{proposition}

\begin{proof}
Put $y=e^{(2-L)/X}$. By \eqref{SM:eq:compactivity}, $(1-2x)^2=y\,x(1-x)$, that is
$(4+y)x^2-(4+y)x+1=0$, whose root in $[0,\frac12]$ is $x=\frac12\big(1-\sqrt{y/(y+4)}\big)$, which is the
formula. The limits follow by inspection: $y\to\infty$, $y=1$, and $y\to0$ respectively. The sign of
$X$ is the sign of $s'(x)$, which changes at $x^*$.
\end{proof}

The small-$X$ behaviour $x\approx e^{-(2-L)/X}$ is an Arrhenius law: a defect costs ``volume'' $2-L$,
just as a thermal excitation costs energy. The negative branch is the analogue of negative temperatures
for spin systems, whose energy is bounded above. Here the jammed volume is bounded above, and past the
entropy maximum adding volume \emph{decreases} the number of states. For $h=1.2$, for example,
$X=0.25,1,-1,-0.25$ give $x\approx0.128,0.239,0.311,0.390$ and $\eta\approx0.595,0.579,0.569,0.559$.

\begin{figure}[ht]
\centering
\begin{tikzpicture}[scale=1.0]
  \draw[->] (0,0)--(7.6,0) node[below right,font=\small]{$\eta$};
  \draw[->] (0,0)--(0,4.5) node[left,font=\small]{$s$};
  \foreach \e/\x in {0.55/0.95,0.56/1.9,0.57/2.85,0.58/3.8,0.59/4.75,0.60/5.7,0.61/6.65}
    {\draw (\x,0)--(\x,-0.08) node[below,font=\scriptsize]{$\e$};}
  \foreach \v in {0.1,0.2,0.3,0.4} {\draw (0,8*\v)--(-0.08,8*\v) node[left,font=\scriptsize]{$\v$};}
  \draw[dashed,gray] (0,3.8497)--(3.224,3.8497) node[pos=0,left,font=\scriptsize,black]{$\ln\varphi$};
  \draw[dashed,gray] (3.224,0)--(3.224,3.8497);
  \draw[very thick,addcol] plot[smooth] coordinates {(6.991,0.000) (6.955,0.140) (6.918,0.252) (6.882,0.353) (6.846,0.447) (6.701,0.781) (6.557,1.071) (6.414,1.330) (6.272,1.567) (6.130,1.785) (5.989,1.987) (5.848,2.174) (5.709,2.349) (5.570,2.512) (5.431,2.663) (5.293,2.804) (5.156,2.935) (5.020,3.057) (4.884,3.169) (4.749,3.272) (4.615,3.367) (4.481,3.452) (4.348,3.530) (4.215,3.599) (4.084,3.660) (3.952,3.712) (3.822,3.756) (3.692,3.792) (3.562,3.819) (3.434,3.838) (3.305,3.848) (3.178,3.849) (3.051,3.841) (2.924,3.824) (2.799,3.798) (2.673,3.761) (2.549,3.715) (2.425,3.657) (2.301,3.589) (2.178,3.509) (2.056,3.416) (1.934,3.310) (1.813,3.191) (1.692,3.055) (1.572,2.903) (1.452,2.733) (1.333,2.542) (1.215,2.328) (1.097,2.086) (0.979,1.812) (0.862,1.497) (0.746,1.128) (0.688,0.915) (0.659,0.799) (0.630,0.675) (0.601,0.541) (0.572,0.393) (0.543,0.224) (0.514,0.000)};
  \fill (6.991,0) circle (0.06) node[above right,font=\scriptsize,align=left]{zigzag\\ $X\to0^+$};
  \fill (0.514,0) circle (0.06) node[right,xshift=3pt,yshift=6pt,font=\scriptsize,align=left]{$0101\cdots$, $X\to0^-$};
  \fill (3.224,3.8497) circle (0.06) node[above,font=\scriptsize]{$x^*=\frac{5-\sqrt5}{10}$, $X=\pm\infty$};
  \node[font=\small,flagcol] at (5.6,3.6) {$X>0$};
  \node[font=\small,flagcol] at (1.0,3.6) {$X<0$};
\end{tikzpicture}
\caption{Edwards entropy per contact of the channel model as a function of the packing fraction, for
$h=1.2$ ($L=1.6$, $w=3.2$); computed from Proposition~\ref{SM:prop:entropy}. Since $V$ decreases as $\eta$
increases, the compactivity $X=\partial V/\partial S$ is positive on the dense side of the maximum and
negative on the loose side.}
\label{SM:fig:entropy}
\end{figure}

\paragraph{The canonical ensemble is a transfer matrix.} Instead of fixing $V$, weight every jammed
state by $e^{-V/X}$:
\[
  Z_\ell(X)=\sum_\sigma e^{-t(\sigma)/X}=e^{-2/X}\sum_\sigma a^{\#0(\sigma)}\,b^{\#1(\sigma)},\qquad a=e^{-L/X},\ b=e^{-2/X},
\]
the sum running over words of length $\ell$ with no $11$. Keep track of the last letter. If
$(c_\ell(0),c_\ell(1))$ are the weighted sums over words ending in $0$ and in $1$, then
$(c_1(0),c_1(1))=(a,b)$ and
\[
  \big(c_{\ell+1}(0),c_{\ell+1}(1)\big)=\big(c_\ell(0),c_\ell(1)\big)\,T,\qquad
  T=\begin{pmatrix}a&b\\a&0\end{pmatrix},
\]
since after a $0$ we may write either letter and after a $1$ only a $0$. So we have
$Z_\ell(X)=e^{-2/X}(a,b)\,T^{\ell-1}(1,1)^T$. Since $T\ge0$ and $T^2>0$, Perron--Frobenius gives
\[
  \frac1\ell\ln Z_\ell(X)\longrightarrow\ln\lambda(X),\qquad \lambda(X)=\frac{a+\sqrt{a^2+4ab}}2 ,
\]
the larger root of $\lambda^2-a\lambda-ab=0$. This is analytic in $X\ne0$: as in van Hove's theorem, the
one-dimensional channel has no Edwards phase transition.

\begin{proposition}[Equivalence of ensembles]\label{SM:prop:canon}
The mean defect density in the canonical ensemble at compactivity $X$,
$\lim_\ell\langle k\rangle/\ell=b\,\partial_b\ln\lambda$ (with $a$ held fixed), equals the microcanonical
$x(X)$ of Proposition~\ref{SM:prop:eos}. Moreover
\[
  \ln\lambda(X)=\max_{0\le x\le1/2}\Big(s(x)-\frac{v(x)}X\Big),
\]
so $-X\ln\lambda$ is the Legendre transform of the entropy (an ``Edwards free energy'' per contact).
\end{proposition}

\begin{proof}
Write $\lambda=a\,g(u)$ with $u=b/a=e^{-(2-L)/X}$ and $g(u)=\frac12(1+r)$, $r=\sqrt{1+4u}$. Then
\[
  b\,\partial_b\ln\lambda=u\,\frac{g'(u)}{g(u)}=\frac{u\cdot 2/r}{1+r}=\frac{r^2-1}{2r(1+r)}=\frac12\Big(1-\frac1r\Big),
\]
and $1/r=\big(1+4e^{-(2-L)/X}\big)^{-1/2}$, as in Proposition~\ref{SM:prop:eos}. For the second statement,
group the words by $k$: $Z_\ell=\sum_k\Omega(\ell,k)e^{-(\ell L+(2-L)k+2)/X}$ has at most $\ell+1$ terms, so
$\frac1\ell\ln Z_\ell$ converges to the maximum over $x$ of $s(x)-v(x)/X$ (Laplace's method). Since $s$ is strictly concave, the
maximizer is unique and is characterized by $s'(x)=(2-L)/X$, which is again \eqref{SM:eq:compactivity}.
\end{proof}

At $X=\infty$ we have $a=b=1$, and $T$ becomes the Fibonacci matrix $\left(\begin{smallmatrix}1&1\\1&0\end{smallmatrix}\right)$ with
$\lambda=\varphi$. The identity $\ln\lambda(\infty)=\max s=\ln\varphi$ is Lecture~11's observation that ``the
maximal term carries the whole exponential growth''. Exercise~\ref{L11:ex:typical} asked for the
typical defect fraction $x^*$; we now see it as the $X=\infty$ point of a whole equation of state. Section~\ref{SM:sec:info}
identifies the measure behind it.

\subsection{Comparison with Bowles and Ashwin}

The same model was analysed independently, and in the same spirit, by Bowles and Ashwin (\emph{Edwards
entropy and compactivity in a model of granular matter}, Phys.\ Rev.~E 83, 2011). They use discs of
diameter $\sigma$ between walls a distance $H$ apart, with $H/\sigma<1+\sqrt{3/4}$, so that each disc
touches only its neighbours and one wall. With $\sigma=2$ and $H=w$ this is exactly our condition
$h<\sqrt3$. Their bond volumes $v_0=H\sqrt{(2-H)H}$ (cross-channel, $\sigma=1$) and $v_1=H$
(same wall) are our $L$ and $2$, rescaled by $\frac12$ and multiplied by the width. They count
$\Omega_J=\binom{N-M}{M}$ jammed states with $M$ defects among $N$ bonds, which differs from Lecture~11's
count only at the ends, and their entropy per particle and compactivity are exactly our $s(x)$ and
\eqref{SM:eq:compactivity}, in their units. They note that $X\to\infty$ at
$\theta=\frac12-\frac{\sqrt5}{10}$ and they restrict attention to $X>0$, remarking that denser-defect
packings would need negative compactivities. They also show that all the jammed packings are isostatic,
and that two subsystems exchanging defects equilibrate at equal compactivity. Ashwin and Bowles had
earlier (\emph{Complete jamming landscape of confined hard discs}, Phys.\ Rev.\ Lett.\ 102, 2009) treated the wider channel $1+\sqrt{3/4}<H/\sigma<2$ with a transfer
matrix on $32$ local tiles.

The open-ends caveat of Lecture~11 is precisely the boundary term on which the two counts differ; it
affects none of the limits above.

%======================================================================
\section{Information theory}
\label{SM:sec:info}

\subsection{Entropy and maximum entropy}

Lecture~7 counted coin flips and Lecture~11 maximized a multinomial coefficient. In both, the number of
sequences of length $N$ with letter frequencies $p$ is $e^{NH(p)+O(\ln N)}$, $H(p)=-\sum_ip_i\ln p_i$.
Shannon (1948) took $H$ as the measure of \emph{information}: such a source can be encoded with
$H(p)/\ln2$ bits per letter and no fewer. Jaynes (\emph{Information theory and statistical mechanics},
Phys.\ Rev.~106, 1957) turned the Lagrange multiplier computation of \S\ref{L11:sec:boltzmann} into a
principle: \emph{the least biased distribution compatible with what one knows has maximal entropy.}

\begin{proposition}[Maximum entropy gives Boltzmann]\label{SM:prop:jaynes}
Let $E_1,\dots,E_k$ be energies and $q_i=e^{-\beta E_i}/Z$ a Boltzmann distribution with mean energy
$\bar E=\sum q_iE_i$. Every probability vector $p$ with $\sum p_iE_i=\bar E$ satisfies $H(p)\le H(q)$, with
equality only for $p=q$.
\end{proposition}

\begin{proof}
The \emph{relative entropy} $D(p\Vert q)=\sum p_i\ln(p_i/q_i)$ is non-negative and vanishes only for $p=q$,
by Jensen's inequality for $\ln$. Since $-\ln q_i=\beta E_i+\ln Z$,
\[
  0\le D(p\Vert q)=-H(p)+\beta\sum_ip_iE_i+\ln Z=-H(p)+\beta\bar E+\ln Z=-H(p)+H(q).
\]
\end{proof}

Edwards' ansatz is the same principle with volume for energy: the flat measure on jammed states of
volume $V$ is the maximum entropy measure given $V$, and the canonical weight $e^{-V/X}$ is the maximum
entropy measure given the mean volume.

\subsection{The golden-mean channel}

Shannon's 1948 paper also treated \emph{constrained} channels, where not every sequence of symbols may be
sent. Its first example is the telegraph, whose symbols have different durations and may not follow each
other arbitrarily. The capacity of such a channel is the exponential growth rate of the number of
admissible messages. The simplest constraint is ``no two consecutive $1$'s'', the $(1,\infty)$
run-length-limited constraint of magnetic and optical recording. In symbolic dynamics the admissible
bi-infinite sequences form the \emph{golden-mean shift}. These are exactly the contact words of the
channel model.

\begin{figure}[ht]
\centering
\begin{tikzpicture}[scale=1.0,>={Stealth}]
  \node[circle,draw,thick,minimum size=0.8cm] (A) at (0,0) {$0$};
  \node[circle,draw,thick,minimum size=0.8cm] (B) at (2.6,0) {$1$};
  \draw[->,thick] (A) to[bend left=30] node[above,font=\small]{$1$} (B);
  \draw[->,thick] (B) to[bend left=30] node[below,font=\small]{$0$} (A);
  \draw[->,thick] (A) to[out=150,in=210,looseness=6] node[left,font=\small]{$0$} (A);
  \node[font=\small] at (6.2,0.3) {$T=\begin{pmatrix}1&1\\1&0\end{pmatrix}$,\quad $\lambda=\varphi$};
  \node[font=\small] at (6.2,-0.5) {Parry: $P(0\to1)=\varphi^{-2}$, $P(0\to0)=\varphi^{-1}$};
\end{tikzpicture}
\caption{The golden-mean shift. Admissible words are paths in this graph, where the state is the last
letter written. Its adjacency matrix is the Fibonacci matrix. With the weights $a=e^{-L/X}$, $b=e^{-2/X}$ on
the edges labelled $0$ and $1$, it becomes the Edwards transfer matrix of the channel.}
\label{SM:fig:golden}
\end{figure}

\begin{theorem}[Shannon 1948; Parry 1964]\label{SM:thm:golden}
\begin{enumerate}[nosep]
\item The number of binary words of length $n$ with no two consecutive $1$'s is $F_{n+2}$, and the
  capacity of the constraint is
  \[
    C=\lim_{n\to\infty}\frac1n\log_2F_{n+2}=\log_2\varphi\approx0.6942\ \text{bits per symbol}.
  \]
\item Among stationary Markov chains supported on admissible sequences, the entropy rate is maximized
  uniquely by the chain with $P(0\to1)=\varphi^{-2}$, $P(0\to0)=\varphi^{-1}$, $P(1\to0)=1$ (the
  \emph{Parry measure}). Its entropy rate is $\ln\varphi$ nats $=C$, and its frequency of $1$'s is
  $\frac{5-\sqrt5}{10}$.
\end{enumerate}
\end{theorem}

\begin{proof}
(1) The count is Lecture~11's Fibonacci argument. Equivalently it is $(1,1)T^{n-1}(1,1)^T$ with $T$ the
Fibonacci matrix, whose eigenvalues are $\varphi$ and $-\varphi^{-1}$, so the count is $c\varphi^n+O(\varphi^{-n})$.

(2) Such a chain is determined by $p=P(0\to1)$, since after a $1$ a $0$ is forced. Its stationary law
is $\pi_0=\frac1{1+p}$, $\pi_1=\frac p{1+p}$, and its entropy rate is $f(p)=\pi_0H(p)=H(p)/(1+p)$. Setting
$f'(p)=0$ gives $(1+p)\ln\frac{1-p}{p}=H(p)$, which simplifies to $2\ln(1-p)=\ln p$, that is
$(1-p)^2=p$, so $p=\frac{3-\sqrt5}2=\varphi^{-2}$ and $1-p=\varphi^{-1}$. Then
$H(p)=2p\ln\varphi+(1-p)\ln\varphi=(1+p)\ln\varphi$, so $f(p)=\ln\varphi$, and
$\pi_1=\frac{\varphi^{-2}}{1+\varphi^{-2}}=\frac1{\varphi^2+1}=\frac{5-\sqrt5}{10}$. Uniqueness: $f$ is
continuous on $[0,1]$, vanishes at $p=0$ and at $p=1$, and has only one critical point in $(0,1)$.
\end{proof}

(The variational principle of symbolic dynamics upgrades (2) to all shift-invariant measures: the Parry
measure is the unique measure of maximal entropy, and its entropy is the topological entropy
$\ln\varphi$. See Lind and Marcus in the Further reading.)

\begin{corollary}\label{SM:cor:capacity}
The maximal Edwards entropy per contact of the channel model equals the Shannon capacity of the
golden-mean constraint, $\ln\varphi$ nats $=\log_2\varphi$ bits. Edwards' flat measure on jammed
channel packings (the case $X=\infty$) is, in the limit, the Parry measure. More generally, the
canonical Edwards ensemble at compactivity $X$ is the Markov measure that maximizes entropy minus mean volume divided by $X$, for the weights $a,b$ of
Figure~\ref{SM:fig:golden}: an equilibrium (Gibbs) measure in the sense of dynamics (Exercise~\ref{SM:ex:markov}).
\end{corollary}

So the typical jammed state of Lecture~11, with its fraction $\frac{5-\sqrt5}{10}$ of same-wall contacts, is
the typical message of an optimally used golden-mean channel.

\begin{remark}[Capacity per unit length]
Shannon's telegraph asks for capacity per unit \emph{time} when symbols have durations $t_i$; for free
concatenation it is $\ln X_0$ with $\sum_iX_0^{-t_i}=1$. A no-$11$ word is a free concatenation of the blocks
$0$ and $10$, of channel lengths $L$ and $2+L$, so jammed configurations (any number of discs) of length
at most $t$ number $e^{C't+o(t)}$ with $e^{-C'L}+e^{-C'(L+2)}=1$. This is $\lambda(X)=1$ with $X=1/C'$, and by
Proposition~\ref{SM:prop:canon} $C'=\max_xs(x)/v(x)$. For $h=1.2$: $C'\approx0.2817$ at $x\approx0.266$, so
$X\approx3.55$.
\end{remark}

\subsection{Noise spheres: Shannon's packing argument}

Codes enter sphere packing through noise. A sender transmits a vector $x\in\R^n$ (one real number per
use of the channel) subject to a power constraint $\abs x^2\le nP$. The receiver sees $y=x+z$, where $z$ has
independent $N(0,N)$ coordinates. The capacity of this \emph{additive white Gaussian noise} channel is
the largest rate $R$ (bits per coordinate) such that $2^{nR}$ codewords can be chosen and decoded with
error probability tending to $0$ as $n\to\infty$.

\begin{theorem}[Shannon 1948, 1949]\label{SM:thm:awgn}
The capacity of the AWGN channel with signal power $P$ and noise power $N$ is
\[
  C=\tfrac12\log_2\Big(1+\frac PN\Big)\ \text{bits per coordinate}.
\]
\end{theorem}


\paragraph{Why packing?} The heuristic is the gas in a box of Lecture~7 again. By the law of large
numbers $\abs z^2=nN(1+O(n^{-1/2}))$: in high dimension the noise is concentrated on a thin shell of
radius $\sqrt{nN}$ around the transmitted codeword, just as the particle counts were concentrated near
their mean. It is also nearly orthogonal to $x$, so $\abs y^2\approx\abs x^2+\abs z^2\le n(P+N)$. Reliable
decoding requires the noise balls of radius $\sqrt{nN}$ around different codewords to be essentially
disjoint, and they all lie in a ball of radius about $\sqrt{n(P+N)}$. Comparing volumes gives at most
$\big(\sqrt{n(P+N)}/\sqrt{nN}\big)^n=(1+P/N)^{n/2}$ codewords, the bound in the theorem. Shannon's 1949 paper (\emph{Communication in the presence of noise}, Proc.\
IRE 37) turned this picture into a proof; achievability uses codewords chosen \emph{at random}, the
analogue of the non-constructive lower bounds of Lecture~1.

Demanding \emph{zero} error against every noise vector of length less than $\rho$ gives a genuine packing
problem. For codewords of equal power, on the sphere of radius $\sqrt{nP}$, this means angular separation
$\theta$ with $\sin(\theta/2)\ge\rho/\sqrt{nP}$, and the best such codes are counted by $M(n,\theta)$ of
Definition~\ref{L3:def:M}. So the Kabatiansky--Levenshtein bound of \S\ref{L3:sec:KL} is a coding theorem,
published as one (Problems of Information Transmission, 1978), and through Corollary~\ref{L3:cor:KL} it is
also the best known exponential upper bound on packing density.

\subsection{Perfect codes, Golay and Leech}

For binary codes (Definition~\ref{L04b:def:code}) the Hamming distance replaces the Euclidean one, and the
noise ball is a Hamming ball.

\begin{proposition}[Hamming's sphere packing bound]\label{SM:prop:hamming}
A binary code $C\subseteq\{0,1\}^n$ with minimal distance $d(C)\ge2e+1$ satisfies
\[
  \abs C\sum_{i=0}^e\binom ni\le2^n .
\]
\end{proposition}

\begin{proof}
The Hamming balls of radius $e$ around the codewords are disjoint, by the triangle inequality, and each
contains $\sum_{i\le e}\binom ni$ words.
\end{proof}

A code attaining equality is \emph{perfect}: its Hamming balls tile the cube, so every received word lies
within distance $e$ of exactly one codeword. The trivial perfect codes are a single word, the whole cube, and the
binary repetition codes of odd length. The nontrivial examples are rare and beautiful.
\begin{itemize}[nosep]
\item The \emph{Hamming codes} (1950): $n=2^r-1$, $e=1$, $\abs C=2^{n-r}$, and indeed
  $2^{n-r}(1+n)=2^n$.
\item The binary \emph{Golay code} (1949): $n=23$, $e=3$, $\abs C=2^{12}$, and
  \[
    2^{12}\Big(1+23+\binom{23}2+\binom{23}3\Big)=2^{12}(1+23+253+1771)=2^{12}\cdot2^{11}=2^{23}.
  \]
\item The ternary Golay code: $n=11$, $e=2$, $3^6$ words over $\mathbb F_3$, and
  $3^6(1+11\cdot2+\binom{11}2\cdot4)=3^6\cdot243=3^{11}$.
\end{itemize}
In fact these are all of them (Theorem~\ref{SM:thm:perfect} in the gallery).

Adding a parity bit turns the Hamming code of length $7$ into the extended Hamming code $[8,4,4]$, and
Construction~A (Definition~\ref{L04b:def:constrA}) turns that into $E_8$, as we saw in Lecture~6. Adding a
parity bit to the Golay code gives the extended Golay code $\mathcal G_{24}$, a $[24,12,8]$ code whose
$759$ words of weight $8$ (the \emph{octads}) form a Steiner system $S(5,8,24)$. Leech (1964, 1967) built
the Leech lattice from it. In the form given by Conway and Sloane (Ch.~4, \S11), $\sqrt8\,\Lambda_{24}$ is the set
of $x\in\Z^{24}$ such that, for some $m\in\{0,1\}$, all $x_i\equiv m\pmod2$, each set
$\{i:x_i\equiv\nu\pmod4\}$ ($\nu\in\Z/4$) is a word of $\mathcal G_{24}$, and $\sum_ix_i\equiv4m\pmod8$.
The minimal vectors have norm $32$ in these coordinates, and there are
\[
  \underbrace{\tbinom{24}2\cdot4}_{(\pm4,\pm4,0^{22})}+\underbrace{759\cdot2^7}_{(\pm2^8,0^{16})\text{ on an octad}}+\underbrace{24\cdot2^{12}}_{(\mp3,\pm1^{23})}
  =1104+97152+98304=196560
\]
of them, the kissing number in dimension $24$ quoted in Lecture~3. So $E_8$ and the Leech lattice, the two lattices now known to be optimal packings beyond dimension $3$,
come from the perfect Hamming code of length $7$ and the perfect Golay code of length $23$ by adding a
parity bit.

\subsection{Lattice codes}

Construction~A also runs the other way: a lattice $\Lambda\subset\R^n$ is a code for the Gaussian channel,
with infinitely many codewords. To respect the power constraint one keeps only the points of $\Lambda$ in a
bounded region. A good choice of region, which keeps the code a group, is the Voronoi cell of a coarser lattice $\Lambda'\subset\Lambda$,
which gives the \emph{nested lattice codes} $\Lambda/\Lambda'$. Erez and Zamir (IEEE Trans.\ Inform.\ Theory 50,
2004) proved that, with a random dither and a scaling at the receiver, nested lattice codes with
lattice decoding achieve the capacity $\frac12\log_2(1+P/N)$. Decoding is then
the nearest-point (Voronoi cell) problem of Lecture~1.

%======================================================================
\section{Gallery of theorems}
\label{SM:sec:gallery}

Results stated in the text (R\'enyi, Tonks and van Hove, Shannon, Kabatiansky--Levenshtein) are not
repeated here.

\begin{theorem}[Hales, Ann.\ of Math.\ 162, 2005; Viazovska, Ann.\ of Math.\ 185, 2017; Cohn--Kumar--Miller--Radchenko--Viazovska, Ann.\ of Math.\ 185, 2017]
The optimal sphere packing densities in dimensions $3$, $8$, $24$ are $\pi/\sqrt{18}$, $\pi^4/384$ and
$\pi^{12}/12!$. Equivalently, the hard-sphere gas of unit diameter in $\R^3$, $\R^8$, $\R^{24}$ has finite
free energy at number densities below, and infinite free energy above, $\sqrt2$, $16$ and $2^{24}$ respectively.
\end{theorem}
\noindent By Proposition~\ref{SM:prop:kepler}, $\rho_{\mathrm{cp}}$ is the optimal density divided by
$\vol\Ball^d(\frac12)$; it is the number density of FCC, $E_8$ and Leech, scaled to minimal distance~$1$.

\begin{theorem}[Hales--McLaughlin, J.~Amer.\ Math.\ Soc.\ 23, 2010]
In a packing of unit balls in $\R^3$, every Voronoi cell has volume at least that of a regular dodecahedron
of inradius $1$.
\end{theorem}
\noindent Fejes T\'oth's dodecahedral conjecture: the locally best density $\approx0.7547$ exceeds the global
optimum $0.7405$. This is geometric frustration as a theorem.

\begin{theorem}[Theil, Comm.\ Math.\ Phys.\ 262, 2006]
For a class of short-range, Lennard-Jones-like pair potentials in the plane, the ground-state energy per
particle of $N$ particles converges as $N\to\infty$ to the energy per particle of the triangular lattice:
pieces of the triangular lattice are asymptotically optimal.
\end{theorem}
\noindent Crystallization at zero temperature, the problem whose finite-temperature version for hard
spheres is open (\S\ref{SM:sec:hard}). Heitmann--Radin (1980) had done ``sticky discs''.

\begin{theorem}[Richthammer, Comm.\ Math.\ Phys.\ 345, 2016]
In the hard disc model in a $2n\times2n$ box with arbitrary boundary condition, the mean square
displacement of a particle near the centre is at least $c\log n$.
\end{theorem}
\noindent Hard discs have no long-range positional order: the most a two-dimensional ``crystal'' can
have is the hexatic orientational order seen numerically by Bernard and Krauth.

\begin{theorem}[Connelly--Gortler--Solomonides--Yampolskaya, Discrete Comput.\ Geom.\ 2020]
A collectively jammed packing of $n$ discs in a flat torus $\R^2/\Lambda$ whose radii and lattice are
generic has exactly $2n-1$ contacts, and its contact graph is isostatic.
\end{theorem}
\noindent The isostatic conjecture of Section~\ref{SM:sec:random} as a theorem in the plane (the lattice is
held fixed, and must itself be generic). Corollary~\ref{SM:cor:maxwell} gives the lower bound $2n-1$; the
content is the matching upper bound, $2n-2$ positions modulo translation plus the common scale of the
radii. The authors leave dimension $\ge3$ open.

\begin{theorem}[Cohn--Kumar, J.~Amer.\ Math.\ Soc.\ 20, 2007]
A configuration on $\Sph^{n-1}$ that has $m$ distinct distances and is a spherical $(2m-1)$-design
(``sharp'') minimizes the energy $\sum_{i\ne j}f(\abs{x_i-x_j}^2)$ among configurations of the same size,
for every completely monotonic $f$. The same holds for the $120$ vertices of the $600$-cell.
\end{theorem}
\noindent Examples: the icosahedron, the $240$ roots of $E_8$, the $196560$ minimal vectors of the Leech
lattice. These are zero-temperature ground states for a whole family of interactions at once, including
Frank's ideal $\{3,3,5\}$ order.

\begin{theorem}[Cohn--Kumar--Miller--Radchenko--Viazovska, Ann.\ of Math.\ 196, 2022]
The $E_8$ and Leech lattices are universally optimal among point configurations in $\R^8$ and $\R^{24}$:
they minimize energy for every potential that is a completely monotonic function of squared distance,
for example Gaussians $e^{-\alpha r^2}$ and inverse powers.
\end{theorem}
\noindent So the zero-temperature Gaussian core model is solved in dimensions $8$ and $24$ at every
density. The sphere packing theorems are recovered in the limit $\alpha\to\infty$.

\begin{theorem}[Tiet\"av\"ainen, SIAM J.~Appl.\ Math.\ 24, 1973; van Lint 1971; Zinov'ev--Leont'ev 1973]\label{SM:thm:perfect}
Every nontrivial perfect code over a finite field has the same length, number of words and minimal
distance as a Hamming code or one of the two Golay codes.
\end{theorem}
\noindent Tilings of the Hamming cube by Hamming balls are classified by their parameters (codes with
the Hamming parameters need not be linear). The Golay code's tiling is the
discrete shadow of the Leech lattice.

\begin{theorem}[Erez--Zamir, IEEE Trans.\ Inform.\ Theory 50, 2004]
Nested lattice codes with dithering, MMSE scaling and lattice decoding achieve the capacity
$\frac12\log_2(1+P/N)$ of the Gaussian channel.
\end{theorem}
\noindent Structured codes, decoded by finding the nearest lattice point, reach Shannon's capacity.

%======================================================================
\section*{Exercises}
\addcontentsline{toc}{section}{Exercises}


\begin{exercise}
(a) Solve \eqref{SM:eq:renyi} on $[2,3]$: show $M(x)=1+\frac{2(x-2)}{x-1}$ there. (b) Show that every saturated
packing of unit intervals in $[0,x]$ has at least $(x-1)/2$ intervals, and deduce $m\ge\frac12$ without
computing anything.
\end{exercise}

\begin{exercise}
(a) Show that the honeycomb disc packing is locally jammed but not collectively jammed: find a collective
motion. (b) Verify the claim of Figure~\ref{SM:fig:maxwell}.
\end{exercise}


\begin{exercise}\label{SM:ex:markov}
Let $r=(r_0,r_1)^T$ be a positive right eigenvector of $T=\left(\begin{smallmatrix}a&b\\a&0\end{smallmatrix}\right)$
for $\lambda=\lambda(X)$. Show that $P(i\to j)=T_{ij}r_j/(\lambda r_i)$ is a stochastic matrix, that its
stationary chain has frequency of $1$'s equal to $x(X)$, and that for $X=\infty$ it is the Parry measure
of Theorem~\ref{SM:thm:golden}.
\end{exercise}

\begin{exercise}
Compute the variance of the number of defects $k$ in the canonical ensemble to leading order in $\ell$, by
differentiating $\ln\lambda$ twice in $\ln b$. Show that it is $\ell\big/\!\big(-s''(x)\big)$, $x=x(X)$ (compare
Lemma~\ref{L11:lem:concave}).
\end{exercise}

\begin{exercise}
(a) Check that $1+90+\binom{90}2=2^{12}$, so the sphere packing bound allows a perfect binary code of
length $90$ with $2^{78}$ words and $e=2$. (b) Show that no such code exists. \emph{Hint:} assume
$0\in C$, so every word of weight $3$ lies at distance $\le2$ from a unique codeword, necessarily of
weight $5$. Count the weight-$3$ words that have $1$'s in the first two positions and are covered by
codewords of weight $5$ with $1$'s in the first two positions: each such codeword covers exactly three of them.
\end{exercise}

\begin{exercise}
(a) Let $z\in\R^n$ have independent $N(0,N)$ coordinates. Show that
$\abs z^2/n\to N$ in probability, and compare with the gas in a box of
\S\ref{L07:subsec:gas}. (b) Show that a fixed $x$ with $\abs x^2=nP$ satisfies $\ip xz/n\to0$ in
probability, and deduce $\abs{x+z}^2/n\to P+N$. (c) If disjoint balls of radius $\sqrt{n(N-\varepsilon)}$
lie in a ball of radius $\sqrt{n(P+N+\varepsilon)}$, bound their number by $2^{n(C+\delta(\varepsilon))}$.
\end{exercise}


\section*{Further reading}
\addcontentsline{toc}{section}{Further reading}
\begin{itemize}[nosep]
\item D.~Ruelle, \emph{Statistical Mechanics: Rigorous Results}, Benjamin, 1969.
\item S.~Torquato and F.~H.~Stillinger, \emph{Jammed hard-particle packings: From Kepler to Bernal and
  beyond}, Rev.\ Mod.\ Phys.\ 82 (2010).
\item A.~Baule, F.~Morone, H.~J.~Herrmann and H.~A.~Makse, \emph{Edwards statistical mechanics for jammed
  granular matter}, Rev.\ Mod.\ Phys.\ 90 (2018).
\item T.~M.~Cover and J.~A.~Thomas, \emph{Elements of Information Theory}, Wiley.
\item D.~Lind and B.~Marcus, \emph{An Introduction to Symbolic Dynamics and Coding}, Cambridge, 1995.
\item J.~H.~Conway and N.~J.~A.~Sloane, \emph{Sphere Packings, Lattices and Groups}, Springer, Chs.~1--5.
\end{itemize}
% ===== end chapters/SM.tex =====
\endgroup
\begingroup\renewcommand{\thechapter}{F}\renewcommand{\theHchapter}{interludeF}
% ===== begin chapters/COMP.tex =====
% MAT670 -- Interlude F (placed last, after Interlude E): computation and certified proof.
% Entirely new material (not in the handwritten notes). Labels prefixed COMP:.
% Scripts referred to in the text live in /home/claude/mat670/code/.

\chapter[Interlude F: Computation and proof]{Interlude F: Computation and certified proof}
\label{COMP:ch}

\begin{addedblock}
This interlude is new: it is not in the handwritten notes. It collects the computational ideas that
sit behind several results quoted in the course (the Kepler conjecture, the Tammes problem for $13$
and $14$ points, the kissing numbers in dimensions $4$, $8$ and $24$) and tries them out on small
problems that the notes already contain. To keep it readable it is set in black; all of it should be
read as added material. The short Python programs mentioned below are in the directory
\texttt{code/} next to these notes.
\end{addedblock}

Look again at the proof of Thue's theorem in Lecture~1. It begins with an infinite object, a
saturated packing of the plane, and it ends with a statement about a single triangle
(Lemma~\ref{L1:lem:triangle-density}, in the form its proof actually uses): a triangle with sides at
least $2$ and circumradius less than $2$ has area at least $\sqrt3$. Everything in between is bookkeeping: a Delaunay triangulation cuts the
packing into pieces, a weighted average reduces the density of the whole to the density of the pieces,
and the boundary terms disappear in the limit. The Fejes T\'oth inequality of Lecture~2 has the same
shape. A triangulation of the sphere is cut into $2n-4$ spherical triangles, Euler's formula counts
them, and one inequality about spherical triangles finishes the job.

This is the shape of almost every sharp result in the subject.
\begin{enumerate}[label=(\arabic*)]
\item \emph{Reduce} a continuum of configurations to finitely many combinatorial \emph{cases}, indexed by
  a geometric object: a triangulation, a contact graph, a Delaunay star, a perfect form.
\item \emph{Prune} the cases by cheap necessary conditions: angle sums, areas, Euler's formula, planarity.
\item \emph{Settle} each surviving case by an inequality (it cannot occur, or it is worse than the
  candidate) or by solving equations (it is the candidate).
\item \emph{Certify} every step, so that the proof can be checked by somebody who does not trust the
  person, or the program, that found it.
\end{enumerate}
In Lecture~1 there was one case and the inequality took two lines. In Hales's proof of the Kepler
conjecture there are thousands of cases and tens of thousands of inequalities. This interlude is about
what happens in between: how cases are indexed and enumerated (Section~\ref{COMP:sec:index}), how
searches are organized (Section~\ref{COMP:sec:bb}), which pieces of linear algebra and optimization do
the work (Sections~\ref{COMP:sec:matrix}--\ref{COMP:sec:kkt}), what algebra can and cannot decide
(Section~\ref{COMP:sec:poly}), and how a computer can produce a \emph{proof} rather than a number
(Section~\ref{COMP:sec:verify}). In Section~\ref{COMP:sec:thue} we put the pieces together and prove
Thue's triangle lemma by a certified computer search.

%======================================================================
\section{The shape of a computer-assisted proof}
\label{COMP:sec:shape}

Here are the problems from the course, arranged by the four steps above.
\begin{center}
\renewcommand{\arraystretch}{1.25}
\small
\begin{tabular}{p{0.19\textwidth}p{0.23\textwidth}p{0.27\textwidth}p{0.18\textwidth}}
\toprule
problem & cases indexed by & each case settled by & reference\\
\midrule
disks in $\R^2$ & Delaunay triangles & area $\ge\sqrt3$ & Lecture~1\\
$N$ caps on $\Sph^2$ & triangulations, Euler & spherical area inequality & Lecture~2\\
critical points of $\rho$ & balanced contact graphs & force balance (an LP) & Lectures~10--11\\
Tammes, $N{=}13,14$ & irreducible graphs & LP and inequalities & Musin--Tarasov\\
kissing, $n=8,24$ & --- (one case) & one polynomial & Section~\ref{COMP:sec:matrix}\\
Kepler, $\R^3$ & tame plane graphs & nonlinear inequalities, LPs & Hales\\
lattices in $\R^n$ & perfect forms & Voronoi's algorithm & Lect.~6, Int.~B\\
\bottomrule
\end{tabular}
\end{center}
Two remarks before we start. First, the line between ``by hand'' and ``by computer'' is not where one
might think. The kissing number $240$ in dimension $8$ is proved by a single polynomial, which a
computer found but which anyone can check with pencil and paper (Example~\ref{COMP:ex:240}). On the
other hand, the enumeration of all tame graphs for the Kepler conjecture cannot be checked by hand at
all. Second, what matters in the end is not whether a computer was used but whether every step
produced a \emph{certificate}: a finite object (a list of graphs, a polynomial with rational
coefficients, a collection of boxes, a vector of weights) together with a simple procedure that checks
it. Most of this interlude is about which certificates are available and how to produce them.

%======================================================================
\section{Indexing cases by geometric objects}
\label{COMP:sec:index}

\subsection{Contact graphs}
Lectures~10 and~11 attached to a configuration $X=(x_1,\dots,x_N)$ on $\Sph^2$ with minimal distance
$\theta=\rho(X)$ its \emph{contact graph} (Definition~\ref{L11:def:contact}), with edges the geodesic arcs of
length $\theta$. Two simple lemmas turn this graph into a good index for case analysis.

\begin{lemma}[Contact graphs are planar]\label{COMP:lem:planar}
Let $0<\theta<\pi$ and let $X\subset\Sph^2$ have all pairwise distances $\ge\theta$. Then two contact
edges (arcs of length $\theta$ between points of $X$) can meet only at a common endpoint.
\end{lemma}

\begin{proof}
Two minor arcs from a common endpoint $a$ lie on great circles through $a$, which meet again only at
$-a$, so they meet only at $a$. Suppose the arcs $[a,b]$ and $[c,d]$ have four distinct endpoints and
cross at a point $p$. By the triangle inequality,
\[
  d(a,c)+d(b,d)\le d(a,p)+d(p,c)+d(b,p)+d(p,d)=d(a,b)+d(c,d)=2\theta .
\]
Both terms on the left are $\ge\theta$, so both equal $\theta$ and both triangle inequalities are
equalities. Since $d(a,c)=\theta<\pi$, equality in $d(a,c)\le d(a,p)+d(p,c)$ means that $p$ lies on the
minor arc from $a$ to $c$. That arc and $[a,b]$ leave $a$ in the same direction (towards $p$) and both
have length $\theta$, so $c=b$, a contradiction.
\end{proof}

\begin{lemma}[Angle lemma]\label{COMP:lem:angle}
Let $u,v,w\in\Sph^2$ with $d(u,v)=d(v,w)=\theta\in(0,\pi)$ and $d(u,w)\ge\theta$. Then the angle
$\varphi$ at $v$ between the arcs $vu$ and $vw$ satisfies $\varphi\ge\alpha(\theta)$, where
\[
  \cos\alpha(\theta)=\frac{\cos\theta}{1+\cos\theta}
\]
is the angle of the equilateral spherical triangle of side $\theta$. Consequently a point of $X$ with
$k\ge2$ contacts has $k\,\alpha(\theta)\le2\pi$.
\end{lemma}

\begin{proof}
By the spherical law of cosines, $\cos d(u,w)=\cos^2\theta+\sin^2\theta\cos\varphi$. The hypothesis
$\cos d(u,w)\le\cos\theta$ gives
$\cos\varphi\le(\cos\theta-\cos^2\theta)/\sin^2\theta=\cos\theta/(1+\cos\theta)$. The $k$ angles between
consecutive contact arcs at a point add up to $2\pi$.
\end{proof}

For $\theta=\pi/3$ this gives $\cos\alpha=\frac13$: \emph{at most five contacts}, because
$5\arccos\frac13<2\pi<6\arccos\frac13$. This is the same number $\arccos\frac13\approx70.53^\circ$ as the
dihedral angle of the regular tetrahedron in Section~\ref{TET:sec}, and for the same reason. For
$\theta=\pi/2$ we get $\alpha=\pi/2$ and at most four contacts.

\subsection{A worked example: the critical configurations of five points}
Exercise~\ref{L10:ex:classify} of Lecture~10 asked for the classification of balanced configurations,
and the notes did $N=3$ and $N=4$ by hand. Let us do $N=5$ the way a computer would.

Let $X$ be a critical configuration of $\rho$ on $\Conf(5,\Sph^2)$, with balancing weights $w\ge0$, not
all zero, as in Theorem~\ref{L11:thm:forcebalance}. Let $S$ be the \emph{stressed graph}: the edges with
$w_e>0$, together with all five points (points without stressed edges are isolated vertices of $S$;
they feel no force). We look for all graphs that can occur as $S$, using four rules.
\begin{enumerate}[label=\textbf{R\arabic*},leftmargin=*]
\item $S$ is planar. (Lemma~\ref{COMP:lem:planar}, since $\theta\le\max\rho<\pi$.)
\item No vertex has degree $1$. (A single positive force cannot balance.)
\item No vertex of degree $2$ lies in a triangle of $S$. (If $v$ has exactly two stressed edges, the two
  forces are opposite, so the neighbours $u,w$ lie on a great circle on either side of $v$ and
  $d(u,w)=\min(2\theta,2\pi-2\theta)$. If also $d(u,w)=\theta$, then $\theta=2\pi/3$. But four unit vectors
  with pairwise inner products $\le-\frac12$ would have $\abs{x_1+\dots+x_4}^2\le4-12\cdot\frac12<0$.)
\item If two vertices $a,b$ have three common neighbours, then $a,b$ are not adjacent, and their
  common neighbours span no triangle. (The circles of radius $\theta$ about $a$ and $b$ meet in at most two
  points unless they coincide, which forces $b=-a$ and $\theta=\pi/2$. Then $d(a,b)=\pi\ne\theta$, and the
  common neighbours lie on the equator between $a$ and $b$, where three points cannot be pairwise
  $\pi/2$ apart.)
\end{enumerate}
The program \texttt{code/contact5.py} runs through the $34$ graphs on five vertices (from the
\emph{Atlas of Graphs} built into the \texttt{networkx} library). One is empty; R1 removes $K_5$; R2 removes $18$;
R3 removes $7$ and R4 removes one more. Six graphs survive (Figure~\ref{COMP:fig:N5}).

\begin{figure}[ht]
\centering
\begin{tikzpicture}[scale=0.95,
  v/.style={circle,fill,inner sep=1.15pt}, f/.style={circle,draw,fill=white,inner sep=1.15pt},
  lab/.style={font=\scriptsize,align=center}]
  % row 1
  \begin{scope}[shift={(0,0)}]
    \node[v] (a) at (0,0.6) {}; \node[v] (b) at (-0.55,-0.35) {}; \node[v] (c) at (0.55,-0.35) {};
    \draw (a)--(b)--(c)--(a); \node[f] at (-0.7,0.55) {}; \node[f] at (0.7,0.55) {};
    \node[lab] at (0,-0.95) {$C_3$\\ R3};
  \end{scope}
  \begin{scope}[shift={(2.1,0)}]
    \node[v] (a) at (-0.45,0.45) {}; \node[v] (b) at (0.45,0.45) {}; \node[v] (c) at (0.45,-0.45) {}; \node[v] (d) at (-0.45,-0.45) {};
    \draw[addcol,thick] (a)--(b)--(c)--(d)--(a); \node[f] at (0,0) {};
    \node[lab] at (0,-0.95) {$C_4$\\ \textbf{realized}, $\pi/2$};
  \end{scope}
  \begin{scope}[shift={(4.2,0)}]
    \node[v] (a) at (0,0.6) {}; \node[v] (b) at (-0.45,0) {}; \node[v] (c) at (0.45,0) {}; \node[v] (d) at (0,-0.6) {};
    \draw (a)--(b)--(d)--(c)--(a) (b)--(c); \node[f] at (0.75,0.55) {};
    \node[lab] at (0,-0.95) {diamond\\ R3};
  \end{scope}
  \begin{scope}[shift={(6.3,0)}]
    \foreach \k in {0,...,4} \node[v] (p\k) at (90+72*\k:0.6) {};
    \draw[addcol,thick] (p0)--(p1)--(p2)--(p3)--(p4)--(p0);
    \node[lab] at (0,-0.95) {$C_5$\\ \textbf{realized}, $2\pi/5$};
  \end{scope}
  \begin{scope}[shift={(8.4,0)}]
    \node[v] (a) at (0,0.62) {}; \node[v] (b) at (-0.58,-0.38) {}; \node[v] (c) at (0.58,-0.38) {}; \node[v] (d) at (0,0) {};
    \draw (a)--(b)--(c)--(a) (d)--(a) (d)--(b) (d)--(c); \node[f] at (0.72,0.55) {};
    \node[lab] at (0,-0.95) {$K_4$\\ free point};
  \end{scope}
  \begin{scope}[shift={(10.5,0)}]
    \node[v] (o) at (0,0) {}; \node[v] (a) at (-0.6,0.42) {}; \node[v] (b) at (-0.6,-0.42) {};
    \node[v] (c) at (0.6,0.42) {}; \node[v] (d) at (0.6,-0.42) {};
    \draw (o)--(a)--(b)--(o)--(c)--(d)--(o);
    \node[lab] at (0,-0.95) {bowtie\\ R3};
  \end{scope}
  \begin{scope}[shift={(12.6,0)}]
    \node[v] (a) at (-0.4,-0.5) {}; \node[v] (b) at (0.4,-0.5) {}; \node[v] (c) at (0.4,0.12) {}; \node[v] (d) at (-0.4,0.12) {};
    \node[v] (e) at (0,0.62) {};
    \draw (a)--(b)--(c)--(d)--(a) (c)--(e)--(d);
    \node[lab] at (0,-0.95) {house\\ R3};
  \end{scope}
  % row 2
  \begin{scope}[shift={(0,-2.5)}]
    \node[v] (n) at (0,0.6) {}; \node[v] (s) at (0,-0.6) {};
    \node[v] (a) at (-0.6,0) {}; \node[v] (b) at (0,0) {}; \node[v] (c) at (0.6,0) {};
    \draw[addcol,thick] (n)--(a)--(s) (n)--(b)--(s) (n)--(c)--(s);
    \node[lab] at (0,-0.95) {$K_{2,3}$\\ \textbf{realized}, $\pi/2$};
  \end{scope}
  \begin{scope}[shift={(2.1,-2.5)}]
    \node[v] (n) at (0,0.6) {}; \node[v] (s) at (0,-0.6) {};
    \node[v] (a) at (-0.6,0) {}; \node[v] (b) at (0.28,0) {}; \node[v] (c) at (0.72,0) {};
    \draw (n)--(a)--(s) (n)--(b)--(s) (n)--(c)--(s) (n)--(s);
    \node[lab] at (0,-0.95) {$K_{1,1,3}$\\ R3};
  \end{scope}
  \begin{scope}[shift={(4.2,-2.5)}]
    \node[v] (t) at (0,0.6) {}; \node[v] (a) at (-0.72,-0.3) {}; \node[v] (b) at (-0.25,-0.45) {};
    \node[v] (c) at (0.25,-0.45) {}; \node[v] (d) at (0.72,-0.3) {};
    \draw (a)--(b)--(c)--(d) (t)--(a) (t)--(b) (t)--(c) (t)--(d);
    \node[lab] at (0,-0.95) {gem\\ R3};
  \end{scope}
  \begin{scope}[shift={(6.3,-2.5)}]
    \node[v] (n) at (0,0.6) {}; \node[v] (s) at (0,-0.6) {};
    \node[v] (a) at (-0.6,0) {}; \node[v] (b) at (0,0) {}; \node[v] (c) at (0.6,0) {};
    \draw (n)--(a)--(s) (n)--(b)--(s) (n)--(c)--(s) (a)--(b);
    \node[lab] at (0,-0.95) {$K_{2,3}+e$\\ weights};
  \end{scope}
  \begin{scope}[shift={(8.4,-2.5)}]
    \node[v] (a) at (0,0.62) {}; \node[v] (b) at (-0.55,-0.3) {}; \node[v] (c) at (0.55,-0.3) {}; \node[v] (d) at (0,0.02) {};
    \node[v] (e) at (0,-0.68) {};
    \draw (a)--(b)--(c)--(a) (d)--(a) (d)--(b) (d)--(c) (e)--(b) (e)--(c);
    \node[lab] at (0,-0.95) {$K_4$ + ear\\ R3};
  \end{scope}
  \begin{scope}[shift={(10.5,-2.5)}]
    \node[v] (a) at (-0.45,0.45) {}; \node[v] (b) at (0.45,0.45) {}; \node[v] (c) at (0.45,-0.45) {}; \node[v] (d) at (-0.45,-0.45) {};
    \node[v] (o) at (0,0) {};
    \draw (a)--(b)--(c)--(d)--(a) (o)--(a) (o)--(b) (o)--(c) (o)--(d);
    \node[lab] at (0,-0.95) {$W_4$\\ weights};
  \end{scope}
  \begin{scope}[shift={(12.6,-2.5)}]
    \node[v] (a) at (-0.62,-0.4) {}; \node[v] (b) at (0.62,-0.4) {}; \node[v] (c) at (0,0.62) {};
    \node[v] (o) at (0,0) {}; \node[v] (e) at (0,-0.75) {};
    \draw (a)--(b)--(c)--(a) (o)--(a) (o)--(b) (o)--(c) (e)--(a) (e)--(b);
    \draw (e) .. controls (1.15,-0.8) and (0.95,0.6) .. (c);
    \node[lab] at (0,-0.95) {$K_5-e$\\ R4};
  \end{scope}
\end{tikzpicture}
\caption{The $14$ graphs on five vertices that pass R1 and R2 (hollow: a point without stressed
contacts). R3 and R4 remove eight of them. Of the six that remain, three are removed by a one-line
geometric argument and three (blue) are realized by critical configurations.}
\label{COMP:fig:N5}
\end{figure}

Now the last, geometric, step, which here is easy enough to do by hand.
\begin{itemize}
\item $K_4$ plus a point: four pairwise equidistant points form a regular tetrahedron, so
  $x_1+x_2+x_3+x_4=0$. For any fifth point $y$ the four numbers $\ip y{x_i}$ add up to $0$, so one of them is
  $\ge0$ and $d(y,x_i)\le\pi/2<\theta$.
\item $W_4$: by R4 applied to the two pairs of opposite rim vertices, $\theta=\pi/2$, the rim is a square on a
  great circle and the hub is its pole. At a rim vertex the two rim forces are opposite and the force
  towards the hub is perpendicular to both, so its weight must vanish.
\item $K_{2,3}+e$: the same argument at an endpoint of the extra edge (north and south forces opposite, the
  extra force along the equator).
\item The survivors are realized: $C_5$ by the regular pentagon on a great circle ($\theta=2\pi/5$); $C_4$
  by a square on a great circle with the fifth point at a pole ($\theta=\pi/2$); $K_{2,3}$ by the two
  poles and three points on the equator with all gaps in $[\pi/2,\pi)$, as in Figure~\ref{TOP:fig:N5} of
  Interlude~D. For $C_5$ and $C_4$ equal weights balance. For $K_{2,3}$ put the same weight $w_k$ on both
  edges at the $k$-th equatorial point (so the equatorial points balance); since all gaps are less than
  $\pi$, the three directions at a pole are not in a closed half-plane, so there are $w_1,w_2,w_3>0$
  balancing the poles. (Equal weights work only for gaps $2\pi/3$.)
\end{itemize}

\begin{proposition}\label{COMP:prop:N5}
The critical values of $\rho$ on $\Conf(5,\Sph^2)$ are $2\pi/5$ and $\pi/2$. At $2\pi/5$ the critical
configurations are the regular pentagons on great circles; at $\pi/2$ they are the configurations with
two antipodal points and three points on the equator between them, with all gaps $\ge\pi/2$.
\end{proposition}

Together with Lecture~11 this is the complete list for $N\le5$. The lowest critical value is $2\pi/N$,
as predicted by the Kusner--Kusner--Lagarias--Shlosman theorem quoted there. For $N=6$ the same rules
leave many more graphs, and the geometric step needs actual computation: one has to solve, for each
graph, the equations ``all edges have length $\theta$, all non-edges are longer, the forces balance''
(Section~\ref{COMP:sec:poly}). That is the general pattern.

\subsection{The same idea at full scale}
\paragraph{Tammes for $N=13,14$.} The Tammes problem asks for $\max\rho$. Musin and Tarasov call a
contact graph \emph{irreducible} (or jammed) if no single point can be shifted so as to increase its
distance to all its neighbours and no ``Danzer flip'' is possible, and they show that for $N>6$ the
contact graph of an optimal configuration is irreducible. The
contact graph of an irreducible configuration is planar, every vertex has degree $0$, $3$, $4$ or $5$
(the upper bound is Lemma~\ref{COMP:lem:angle}), every face is a convex spherical polygon with sides $\theta$ and at most
$\lfloor2\pi/\theta\rfloor$ vertices, and isolated vertices lie in faces with at least six vertices.
These are the analogues of R1--R4. They enumerate the planar graphs with these properties by computer,
discard most of them by area and angle inequalities (in the spirit of the proof of the Fejes T\'oth
inequality: the face areas add up to $4\pi$, and each face has a minimal area depending on its
shape), and treat the few survivors by optimization. This gives the optimal $13$- and $14$-point
configurations (Section~\ref{COMP:sec:gallery}). Generating planar graphs with prescribed degrees and face
sizes is itself a well-developed subject; the standard tool is \texttt{plantri} (Brinkmann and McKay).

\paragraph{Kepler.} In Hales's proof the cases are indexed by \emph{tame plane graphs}. Around each ball
one draws, on a sphere about its centre, the planar graph whose vertices are (the radial projections of)
the neighbouring centres within a fixed distance ($2.51$ in the original proof, $2.52$ in the revised and formal
one) and whose edges join neighbours that are close to each other. The local density (the \emph{score} of a
\emph{decomposition star}, later replaced by \emph{Marchal cells}) is then bounded case by case.
``Tame'' is a list of combinatorial conditions of the R1--R4 type (degree bounds, face sizes, a weight
inequality), designed so that a non-tame graph is excluded by an easy argument. Hales's original archive
had $5128$ tame graphs. When Nipkow, Bauer and Schultz re-did the enumeration inside the proof assistant
Isabelle (2006), they confirmed that the archive was complete and found it redundant, reducing it to
$2771$ graphs. The revised proof uses a larger archive of $18\,762$ graphs.

\paragraph{Lattices.} For lattice packings the index is a \emph{perfect form}
(Theorem~\ref{NT:gal:voronoi} of Interlude~B). Voronoi's algorithm walks from one perfect form to its
neighbours across the facets of a polyhedral cone, so the enumeration is a graph search. Its cost is
dominated by symmetry: one must decide whether a new form is equivalent to one already found. In
dimension $8$ there are $10\,916$ perfect forms (Dutour Sikiri\'c, Sch\"urmann and Vallentin, 2007).

%======================================================================
\section{Branch and bound, and meet in the middle}
\label{COMP:sec:bb}

\subsection{Branch and bound}
Suppose we want to show that a function $F$ on a compact set $K$ satisfies $F\ge0$, or to find
$\max_KF$. \emph{Branch and bound} keeps a list of pieces of $K$ (boxes, cells, partial assignments).
For each piece $B$ it computes a cheap \emph{bound}: a number $\underline F(B)\le\inf_BF$ (for the
inequality) or $\overline F(B)\ge\sup_BF$ (for the maximum). If the bound settles the piece, the
piece is discarded; otherwise it is \emph{branched}, that is, split into smaller pieces. Two
properties make it work: the bounds must be \emph{valid} (never wrong), and they must become
\emph{sharp} as the pieces shrink.

\paragraph{Max-min problems.} The min-type function $\rho$ of Lecture~10 is well suited to this. Let
$B_1,\dots,B_N\subset\Sph^2$ be small cells and suppose $x_i\in B_i$. Then
\[
  \rho(x_1,\dots,x_N)\ \le\ \overline\rho(B_1,\dots,B_N):=\min_{i<j}\ \max_{y\in B_i,\,z\in B_j}d(y,z),
\]
and the right side is a minimum of $\binom N2$ numbers that are easy to compute. If
$\overline\rho<\rho(X_0)$ for a known configuration $X_0$, no configuration with $x_i\in B_i$ can be
optimal. The difficulty is the number of pieces: with cells of size $\varepsilon$ there are about
$(4\pi/\varepsilon^2)^N$ choices, and symmetry multiplies the work by $\abs{\SO(3)\times\Sigma_N}$-many
copies of the same configuration. For the analogous problem of $n$ equal circles in a square,
Mark\'ot and Csendes made this approach work with interval arithmetic and a clever elimination of
regions (the ``method of active areas''), proving the optimal packings for $n=28,29,30$
(about $50$, $50$ and $20$ hours of computing in 2005); Mark\'ot later did $n=31,32,33$ (2021). For
the Tammes problem with $13$ points, the graph enumeration of Section~\ref{COMP:sec:index} is used
instead: fixing the combinatorics first removes both the symmetry and most of the dimension.

\paragraph{Branch and bound on integer points.} The same idea finds short vectors in a lattice. Let $G$ be
a positive definite Gram matrix and $C>0$; we want all $x\in\Z^n$ with $x^TGx\le C$. Completing squares
(the $LDL^T$ decomposition, in exact rational arithmetic) writes
\[
  x^TGx=\sum_{i=1}^nq_{ii}\Big(x_i+\sum_{j>i}q_{ij}x_j\Big)^2,\qquad q_{ii}>0 .
\]
Choose $x_n$ first: the $i=n$ term alone must be $\le C$, which leaves finitely many integers. Given
$x_n$, the remaining budget bounds $x_{n-1}$, and so on. This is the algorithm of Fincke and Pohst
(1985), a branch and bound whose pieces are the partial vectors $(x_k,\dots,x_n)$. Run on the Cartan
matrix of $E_8$ (\texttt{code/fincke\_pohst.py}), it visits $739$ nodes and finds the $240$ roots; with
$C=4$ it visits $5477$ nodes and finds the $2400$ nonzero vectors of norm at most $4$. The size of the
tree depends heavily on the basis: a basis with long, nearly parallel vectors gives a very unbalanced
tree. The LLL algorithm of Lenstra, Lenstra and Lov\'asz (1982) computes, in polynomial time, a reduced
basis whose first vector is within a factor $2^{(n-1)/2}$ of the shortest, and in practice LLL followed by
Fincke--Pohst enumeration is how minimal vectors are found. Voronoi's algorithm needs exactly this at
every step: the minimal vectors of each new form.

\subsection{Meet in the middle}
The idea of \emph{meet in the middle} is to split a search over pairs $(u,v)$ into two half-searches and
match the halves through a table. The classic example is \emph{subset sum}: given $n$ integers, is there
a subset with sum $T$? Split the integers into two halves, list the $2^{n/2}$ subset sums of each half, sort
one list, and look up $T-s$ for every $s$ in the other. This takes about $2^{n/2}$ steps instead of $2^n$
(Horowitz and Sahni, 1974), and Schroeppel and Shamir (1981) reduced the memory to about $2^{n/4}$.

Here is the same idea counting vectors of $E_8$, as defined in Lecture~6
(Definition~\ref{L04b:def:E8}). Write $x=(u,v)$ with $u,v\in\R^4$. Then $x\in E_8$ exactly when $u,v$ are
both integral or both in $(\Z+\frac12)^4$, and the coordinate sum is even. So it suffices to tabulate the
half-vectors $u$ by the pair $\big(\abs u^2,\ \text{sum of $u$ modulo }2\big)$, separately for the two
cosets, and to multiply the counts of matching classes. The program \texttt{code/e8\_mitm.py} lists only
$625$ half-vectors and recovers
\[
  240,\quad 2160,\quad 6720,\quad 17520\qquad\text{vectors of norm }2,4,6,8,
\]
which is $240\,\sigma_3(m)$ as predicted by the theta series (equation~\eqref{NT:eq:E8theta} in
Interlude~B), whereas the full list has $26\,641$ vectors. In lattice language: $E_8$ is a union of four cosets of
$D_4\oplus D_4$, and the theta series of an orthogonal sum is the product of the theta series. This
``gluing'' is how Conway and Sloane describe many lattices.

To be honest about where meet in the middle is used: as an \emph{algorithm} it is central in coding
theory and cryptography. Stern's method (1989) for finding low-weight codewords splits an information set
in two halves and looks for collisions, and Howgrave-Graham (2007) combined lattice reduction with a
meet-in-the-middle search to attack the NTRU cryptosystem. In the proofs about packings quoted in this
course it appears mainly in the gluing form above, as a way to count or list vectors without listing the
whole space.

%======================================================================
\section{Matrix techniques}
\label{COMP:sec:matrix}

\subsection{Gram matrices}
A configuration of $N$ unit vectors is determined, up to an orthogonal transformation, by its Gram
matrix $G_{ij}=\ip{x_i}{x_j}$.

\begin{proposition}\label{COMP:prop:gram}
A real symmetric $N\times N$ matrix $G$ is the Gram matrix of $N$ vectors in $\R^d$ if and only if $G$ is
positive semidefinite of rank at most $d$. The vectors are then unique up to $\Orth(d)$.
\end{proposition}

\begin{proof}
If $G=V^TV$ with $V$ a $d\times N$ matrix, then $y^TGy=\abs{Vy}^2\ge0$ and $\operatorname{rank}G\le d$.
Conversely, diagonalize $G=Q\Lambda Q^T$ with $\Lambda\ge0$ having at most $d$ nonzero entries, and take
$V$ to be the nonzero rows of $\Lambda^{1/2}Q^T$. If $V^TV=W^TW$, the map $Vy\mapsto Wy$ is a well defined
isometry between the column spaces, and it extends to an element of $\Orth(d)$.
\end{proof}

So ``is there a configuration of $N$ points on $\Sph^{d-1}$ with all angles at least $\theta$?'' becomes ``is
there a positive semidefinite matrix with ones on the diagonal, off-diagonal entries $\le\cos\theta$, and
rank at most $d$?'' Without the rank condition this is a \emph{semidefinite program}. The rank condition
is what makes the problem hard, and dropping it is what makes bounds computable. Sometimes, however, the
sign pattern alone settles everything.

\begin{proposition}[Rankin 1955]\label{COMP:prop:rankin}
Let $x_1,\dots,x_N\in\R^d$ be nonzero.
\begin{enumerate}[label=(\alph*),nosep]
\item If $\ip{x_i}{x_j}<0$ for all $i\ne j$, then $N\le d+1$.
\item If $\ip{x_i}{x_j}\le0$ for all $i\ne j$, then $N\le2d$.
\end{enumerate}
\end{proposition}

\begin{proof}
\emph{Claim:} vectors with pairwise non-positive inner products that lie in an open half-space
$\{\ip ye>0\}$ are linearly independent. Indeed, a linear relation can be written
$u=\sum_{i\in P}a_ix_i=\sum_{j\in Q}b_jx_j$ with $a,b\ge0$ and $P\cap Q=\emptyset$. Then
$\abs u^2=\sum a_ib_j\ip{x_i}{x_j}\le0$, so $u=0$, and $0=\ip ue=\sum a_i\ip{x_i}e$ forces $a=0$; similarly
$b=0$. For (a), apply the claim to $x_1,\dots,x_{N-1}$ and $e=-x_N$. For (b), choose $e$ with
$\ip{x_i}e\ne0$ for all $i$ and apply the claim to the vectors on each side.
\end{proof}

Since $d(x,y)\ge\pi/2$ on $\Sph^2$ means $\ip xy\le0$, this settles two cases of the Tammes problem at
once and bounds a third: five or six points cannot all be more than $\pi/2$ apart (by (a), $N\le4$), and
the triangular bipyramid and the octahedron show that $\pi/2$ is attained for $N=5,6$; seven points cannot
all be $\pi/2$ apart (by (b), $N\le6$), so the Tammes distance for $N=7$ is less than $\pi/2$ (its actual
value, $\approx77.87^\circ$, needs much more work).

\subsection{Positive definite kernels and the linear programming bound}
The proof of Proposition~\ref{COMP:prop:rankin} used one positive semidefinite matrix. The linear
programming bound uses infinitely many at once. For $S^{n-1}\subset\R^n$ let $G_k$ be the Gegenbauer
polynomial $C_k^{(n-2)/2}$, normalized by $G_k(1)=1$. By the addition theorem for spherical harmonics,
$G_k(\ip xy)$ is a positive multiple of $\sum_mY_m(x)Y_m(y)$, summed over an orthonormal basis of the
harmonic polynomials of degree $k$. Hence for every finite $C\subset S^{n-1}$
\begin{equation}\label{COMP:eq:posdef}
  \sum_{x,y\in C}G_k(\ip xy)\ \ge\ 0,
\end{equation}
that is, the matrix $\big(G_k(\ip xy)\big)_{x,y\in C}$ is positive semidefinite (Schoenberg, 1942).

\begin{theorem}[Delsarte--Goethals--Seidel 1977; Kabatiansky--Levenshtein 1978]\label{COMP:thm:LP}
Let $f(t)=\sum_{k=0}^Dc_kG_k(t)$ with $c_0>0$ and $c_k\ge0$ for all $k$, and suppose $f(t)\le0$ for
$-1\le t\le\cos\theta$. Then every $C\subset S^{n-1}$ with all angles at least $\theta$ has
$\abs C\le f(1)/c_0$.
\end{theorem}

\begin{proof}
Sum $f(\ip xy)$ over all pairs. The off-diagonal terms are $\le0$, so the sum is at most $\abs Cf(1)$. By
\eqref{COMP:eq:posdef} every $k\ge1$ contributes $\ge0$ and $k=0$ contributes $c_0\abs C^2$. Hence
$c_0\abs C^2\le\abs Cf(1)$.
\end{proof}

This is the third time in these notes that a bound comes from \emph{throwing away squares}: Blichfeldt
dropped $\abs{\sum x_i}^2$ (Lemma~\ref{L2:lem:f0}), Cohn and Elkies dropped one square for each
frequency (Remark~\ref{NT:rem:blichfeldt-vs-CE}), and here each $G_k$ hides a sum of squares.

\begin{example}[The kissing number in dimension $8$]\label{COMP:ex:240}
Take $n=8$, $\theta=\pi/3$, and the polynomial of Odlyzko--Sloane and Levenshtein,
\[
  f(t)=(t+1)\Big(t+\frac12\Big)^2t^2\Big(t-\frac12\Big).
\]
Clearly $f\le0$ on $[-1,\frac12]$. Expanding in the $G_k$ for $n=8$ (exact rational arithmetic; a few lines of
\texttt{sympy}) gives
\[
  (c_0,\dots,c_6)=\Big(\frac3{320},\ \frac3{40},\ \frac{15}{64},\ \frac{39}{80},\ \frac{399}{640},\ \frac9{16},\
  \frac{33}{128}\Big),
\]
all positive, and $f(1)/c_0=\frac94\cdot\frac{320}3=240$. So $\kiss_8\le240$, and $E_8$ gives equality.
The same computation with
$f(t)=(t+1)(t+\frac12)^2(t+\frac14)^2t^2(t-\frac14)^2(t-\frac12)$ (degree $10$) in dimension $24$ gives
$196\,560$, with all eleven coefficients $c_0,\dots,c_{10}$ positive. Note the roots: they are the inner products $-1,-\frac12,0,\frac12$
occurring in the $E_8$ kissing configuration, and $-1,\pm\frac12,\pm\frac14,0$ in the Leech one. Equality in
the proof forces $f(\ip xy)=0$ for $x\ne y$, which is how one guesses $f$, and how Bannai and Sloane
(1981) proved that both kissing configurations are unique.
\end{example}

The polynomial was found by solving a linear program (maximize $c_0$ subject to $f(1)=1$ and to
sampled constraints $f(t_\ell)\le0$), but the \emph{certificate} is the polynomial: rational coefficients,
checked exactly. In dimensions $3$ and $4$ the LP bound is not sharp (it gives $13$ and $25$). Musin
proved $\kiss_4=24$ (2008) by a modification of the method. Replacing pairs by triples of points gives
\emph{semidefinite} programming bounds (Bachoc and Vallentin, 2008), which improved the upper bounds in
dimensions $5,6,7,9,10$; Leijenhorst and de Laat (2024) improved them in dimensions $11$ to $23$. For
packings in $\R^n$ the continuous analogue is the Cohn--Elkies bound (Theorem~\ref{NT:thm:CE}), and its
three-point refinement (Cohn, de Laat and Salmon, 2022) improves the density bounds in dimensions $4$--$7$
and $9$--$16$.

\subsection{Rigidity matrices and stresses}
The other matrix in this subject is the \emph{rigidity matrix} $A$ of a contact graph: the linear map that
sends a motion $v=(v_1,\dots,v_N)$ to the rates of change of the contact distances (Lecture~11, and
Interlude~E, Proposition~\ref{SM:prop:jam}). Its kernel is the space of infinitesimal motions, and the
kernel of $A^T$ is the space of \emph{stresses}, systems of edge weights that balance at every vertex.
Theorem~\ref{L11:thm:forcebalance} says that a configuration is critical when $\ker A^T$ meets the
non-negative orthant. Two computational remarks. First, if the configuration has algebraic coordinates,
ranks and kernels can be computed exactly over a number field, and a positive stress, once found
numerically, can be checked exactly. Second, deciding whether $\ker A^T$ contains a non-negative nonzero
vector is a linear program: maximize nothing subject to $A^Tw=0$, $w\ge0$, $\sum w_e=1$. Its feasibility
is exactly Gordan's alternative from Lecture~10. The next section explains why this is no accident.

%======================================================================
\section{KKT: forces are Lagrange multipliers}
\label{COMP:sec:kkt}

The Tammes problem $\max\rho$ is a nonsmooth optimization problem, since $\rho$ is a minimum of smooth
functions. The standard trick is to add a variable. Write $s$ for the cosine of the minimal distance and
consider
\begin{equation}\label{COMP:eq:tammes}
  \text{minimize } s\quad\text{subject to}\quad \ip{x_i}{x_j}\le s\ \ (i<j),\qquad \abs{x_i}^2=1\ \ (1\le i\le N).
\end{equation}
Now everything is smooth (indeed polynomial), and a configuration $X$ is a local maximum of $\rho$ if and
only if $(X,\cos\rho(X))$ is a local minimum of \eqref{COMP:eq:tammes}. The basic necessary condition for
smooth constrained problems is due to F.~John (1948) for inequality constraints, and to Mangasarian and
Fromovitz (1967) with equality constraints.

\begin{theorem}[Fritz John conditions]
Let $p$ be a local minimum of $f$ subject to $g_k\le0$ and $h_l=0$, all $C^1$. Then there are
$\lambda_0\ge0$, $\lambda_k\ge0$ and $\mu_l\in\R$, not all zero, with
\[
  \lambda_0\nabla f(p)+\sum_k\lambda_k\nabla g_k(p)+\sum_l\mu_l\nabla h_l(p)=0,\qquad \lambda_kg_k(p)=0\ \text{for all }k.
\]
If $\lambda_0>0$ one may take $\lambda_0=1$; these are the \emph{Karush--Kuhn--Tucker (KKT)} conditions.
\end{theorem}

Usually one needs a \emph{constraint qualification} (for instance, linear independence of the gradients of
the active constraints) to guarantee $\lambda_0>0$. For the Tammes problem it comes for free, and the
multipliers turn out to be the forces of Lecture~11.

\begin{proposition}\label{COMP:prop:kkt}
A configuration $X$ with $\rho(X)=\theta\in(0,\pi)$ satisfies the KKT conditions for
\eqref{COMP:eq:tammes} (with $s=\cos\theta$) if and only if it is critical in the sense of
Theorem~\ref{L11:thm:forcebalance}. The multipliers of the constraints $\ip{x_i}{x_j}\le s$ are balancing
weights on the contact graph. In particular every local maximum of $\rho$ is critical.
\end{proposition}

\begin{proof}
With $g_{ij}=\ip{x_i}{x_j}-s$ and $h_i=\abs{x_i}^2-1$, the derivative of the Lagrangian in $s$ and in $x_i$
gives
\[
  \lambda_0=\sum_{i<j}\lambda_{ij},\qquad \sum_{j\ne i}\lambda_{ij}x_j+2\mu_ix_i=0\quad(1\le i\le N),
\]
with $\lambda_{ij}=\lambda_{ji}$, and $\lambda_{ij}=0$ unless $\ip{x_i}{x_j}=s$, that is, unless $ij$ is a
contact. For a contact, $x_j=\cos\theta\,x_i+\sin\theta\,u_{ij}$, with $u_{ij}$ the unit tangent vector at
$x_i$ pointing to $x_j$. The tangential part of the second equation is therefore
$\sin\theta\sum_j\lambda_{ij}u_{ij}=0$: the $\lambda_{ij}$ balance the contact graph. If all
$\lambda_{ij}$ vanished, then $\lambda_0=0$ and $2\mu_ix_i=0$, so all multipliers would vanish; hence the
weights are not all zero and $\lambda_0=\sum\lambda_{ij}>0$. Conversely, balancing weights $w_{ij}$, scaled
so that $\sum w_{ij}=1$, give KKT multipliers $\lambda_{ij}=w_{ij}$, $\lambda_0=1$ and
$\mu_i=-\frac12\cos\theta\sum_jw_{ij}$ (the normal part of the second equation). The last sentence follows
from the Fritz John theorem.
\end{proof}

So the dictionary is: \emph{multipliers are forces}, \emph{complementary slackness says that only contacts
carry force}, and \emph{Gordan's alternative is LP duality}. The weights are found by a linear program,
and the forces at each vertex close up into a polygon (Figure~\ref{COMP:fig:forces}); that polygon is the
first step of the Maxwell--Cremona correspondence between stresses and reciprocal diagrams.

\begin{figure}[ht]
\centering
\begin{tikzpicture}[scale=1.5]
  \begin{scope}
    \fill (0,0) circle (1.6pt) node[above left,font=\small] {$x_i$};
    \foreach \a/\w/\l in {80/0.684/1,190/1/2,330/1/3} {
      \draw[gray] (0,0)--(\a:1.35); \fill (\a:1.35) circle (1.3pt) node[font=\scriptsize,shift={(\a:0.22)}] {$x_{\l}$};
      \draw[-{Stealth},flagcol,very thick] (0,0)--(\a:\w);}
    \draw[gray,dashed] (0,0)--(265:1.35); \fill (265:1.35) circle (1.3pt) node[font=\scriptsize,below] {$x_4$};
    \node[font=\scriptsize,align=center] at (0,-1.85) {weights $0.68,\ 1,\ 1$ on the solid contacts;\\ the dashed contact carries no force};
  \end{scope}
  \begin{scope}[xshift=3.6cm,yshift=-0.2cm]
    \coordinate (P0) at (0,0); \coordinate (P1) at (0.119,0.674); \coordinate (P2) at (-0.866,0.5);
    \draw[-{Stealth},flagcol,very thick] (P0)--(P1); \draw[-{Stealth},flagcol,very thick] (P1)--(P2);
    \draw[-{Stealth},flagcol,very thick] (P2)--(P0);
    \fill (P0) circle (1pt);
    \node[font=\scriptsize,align=center] at (-0.3,-1.65) {the three forces, head to tail,\\ close up: $\sum_jw_{ij}u_{ij}=0$};
  \end{scope}
\end{tikzpicture}
\caption{Force balance at one vertex of a critical configuration, seen in the tangent plane at $x_i$. The
multipliers of Proposition~\ref{COMP:prop:kkt} are the lengths of the arrows. A fourth contact with zero
weight is allowed: the weights are $\ge0$, not necessarily $>0$.}
\label{COMP:fig:forces}
\end{figure}

\begin{remark}[Necessary, not sufficient]
KKT points need not be maxima: the square of Lecture~11 is critical but not a local maximum of $\rho$. A
sufficient condition comes from Interlude~E. If $X$ is \emph{first-order jammed} (every nontrivial
infinitesimal motion strictly shortens some contact, which by Proposition~\ref{SM:prop:jam} means a positive
stress on every contact plus infinitesimal rigidity modulo rotations), then a compactness argument on the
unit sphere of motions shows that $X$ is a strict local maximum of $\rho$ up to rotation. The
$\rho$-evolution algorithm of Lecture~11 is, in this language, an \emph{active set method}: it moves in a
direction that increases all active constraints until a new constraint becomes active, and it stops at a
KKT point.
\end{remark}

%======================================================================
\section{Polynomial methods}
\label{COMP:sec:poly}

\subsection{Equations for balanced configurations}
Once a contact graph is fixed, a critical configuration is a solution of polynomial equations:
$\abs{x_i}^2=1$, $\ip{x_i}{x_j}=c$ on the edges, and the balance equations
$\sum_jw_{ij}(x_j-c\,x_i)=0$, modulo rotations, together with the inequalities $\ip{x_i}{x_j}<c$ off the
edges and $w\ge0$. Polynomial algebra solves the equations exactly.

\begin{example}[The square antiprism, $N=8$]
The Tammes configuration for $8$ points (Sch\"utte and van der Waerden, 1951) is a square antiprism with
all $16$ edges equal. Put the two squares at heights $\pm h$ on circles of radius $r$, rotated by
$45^\circ$ against each other. A side of a square has squared length $2r^2$ and a slanted edge has
$(2-\sqrt2)r^2+4h^2$, and $\cos\theta=1-r^2$. With $s=\sqrt2$ and $r_2=r^2$, the equations are
\[
  r_2+h^2-1=0,\qquad s^2-2=0,\qquad s\,r_2-4h^2=0,\qquad c-1+r_2=0 .
\]
A Gr\"obner basis in the lexicographic order $r_2>h>s>c$ (Buchberger's algorithm, one line in
\texttt{sympy}) is
\[
  \{\,r_2+c-1,\ \ h^2-c,\ \ 2s-7c-1,\ \ 7c^2+2c-1\,\}.
\]
It is triangular, like the result of Gaussian elimination, and its last element is a polynomial in $c$
alone. So $c=\cos\theta$ satisfies $7c^2+2c-1=0$, that is $\cos\theta=(2\sqrt2-1)/7$ and
$\theta\approx74.8585^\circ$.
\end{example}

For larger systems one needs to decide how many real roots a univariate polynomial has in an
interval. \emph{Sturm's theorem} (1829) does it exactly: with $p_0=p$, $p_1=p'$ and
$p_{k+1}=-\operatorname{rem}(p_{k-1},p_k)$, the number of roots in $(a,b]$ is the drop in the number of sign
changes of $(p_0,p_1,\dots)$ from $a$ to $b$. For $p=7c^2+2c-1$ the sequence is $(7c^2+2c-1,\,14c+2,\,\frac87)$,
with signs $(-,+,+)$ at $0$ and $(+,+,+)$ at $1$: exactly one root in $(0,1]$. Descartes' rule of signs
and its refinements (Vincent; Collins and Akritas, 1976) give faster \emph{real root isolation}.
\emph{Resultants} eliminate one variable at a time, and \emph{homotopy continuation} (as in the programs
Bertini and PHCpack) follows the solutions of an easy system along a path to all isolated complex solutions
of the given one, numerically. Numerical solutions can be certified afterwards: by Smale's
\emph{$\alpha$-theory}, if $\alpha(f,x)=\beta(f,x)\gamma(f,x)<(13-3\sqrt{17})/4\approx0.1577$, where
$\beta=\norm{Df(x)^{-1}f(x)}$ is the Newton step and
$\gamma=\sup_{k\ge2}\norm{Df(x)^{-1}D^kf(x)/k!}^{1/(k-1)}$, then Newton's method started at $x$ converges
quadratically to a true zero. The program \texttt{alphaCertified} (Hauenstein and Sottile, 2012) checks
these conditions in exact or interval arithmetic.

\subsection{Tarski--Seidenberg: everything is decidable}
A \emph{first-order formula over $\R$} is built from polynomial equations and inequalities with rational
coefficients, using $\wedge,\vee,\neg$ and quantifiers $\exists x,\forall x$ over real numbers. For example,
``there are $13$ points on $\Sph^2$ with pairwise distances at least $\pi/3$'' is
\[
  \exists x_1,\dots,x_{13}\in\R^3:\ \bigwedge_i\abs{x_i}^2=1\ \wedge\ \bigwedge_{i<j}\ip{x_i}{x_j}\le\tfrac12 .
\]

\begin{theorem}[Tarski 1948/1951; Seidenberg 1954]
Every first-order formula over $\R$ is equivalent to a quantifier-free one, and the equivalent formula can
be computed. In particular the truth of every first-order sentence over $\R$ is decidable.
\end{theorem}

So $\kiss_3\le12$ could, \emph{in principle}, be decided by running an algorithm. There is also a
nice theoretical consequence. Apply quantifier elimination to the formula in one free variable
$c$: ``there are $N$ unit vectors with $\ip{x_i}{x_j}\le c$''. The result defines a finite union of
points and intervals with real algebraic endpoints, and its least element is the cosine of the Tammes
distance. \emph{Hence every Tammes value $\cos d_N$ is an algebraic number}, as are the optimal
values of all the finite configuration problems of the course, without anyone having to find them.

Why not just do it? Collins's \emph{cylindrical algebraic decomposition} (1975), the first practical
algorithm, has running time roughly $(sd)^{2^{O(k)}}$ for $s$ polynomials of degree $d$ in $k$ variables,
\emph{doubly} exponential in $k$, and Davenport and Heintz (1988) showed that quantifier elimination
can be inherently doubly exponential. For a single block of existential quantifiers there are better
algorithms (Grigoriev--Vorobjov, Renegar, Basu--Pollack--Roy) with running time of the order
$s^{k+1}d^{O(k)}$, singly exponential. For the $13$-sphere sentence above, $k=39$ and $s=91$; that is far
beyond any computer. In practice, cylindrical algebraic decomposition handles problems with a handful of
variables. Tarski's theorem decides everything in principle and nothing big in practice. The art of
the subject is to reduce a big problem to many small ones, each with few variables, and then to treat
those by the methods of the next section.

%======================================================================
\section{Verification}
\label{COMP:sec:verify}

\subsection{Interval arithmetic}
\emph{Interval arithmetic} (R.~E.~Moore, 1966) computes with closed intervals $X=[\underline x,\overline x]$
instead of numbers:
\[
  X+Y=[\underline x+\underline y,\ \overline x+\overline y],\qquad
  X\cdot Y=[\min P,\max P],\quad P=\{\underline x\underline y,\underline x\overline y,\overline x\underline y,\overline x\overline y\},
\]
and so on. On a computer each lower endpoint is rounded down and each upper endpoint up (the IEEE~754
standard provides these \emph{directed rounding} modes), or one uses exact rational endpoints, as we will.
By induction on the length of an expression one gets the \emph{fundamental theorem}: if an expression $F$ is
evaluated on intervals, the result contains $F(x)$ for every choice of $x$ in the intervals. So a single
evaluation on a box gives a \emph{proved} enclosure of the range of $F$ on the box.

The price is overestimation, the \emph{dependency problem}: each occurrence of a variable is treated as
independent. For $f(x)=x-x^2$ on $X=[0,1]$, the natural evaluation $X-X\cdot X=[0,1]-[0,1]=[-1,1]$, while the
true range is $[0,\frac14]$. Rewriting $f=\frac14-(x-\frac12)^2$ gives the exact range. The \emph{centred
form} $f(m)+f'(X)(X-m)$, here $\frac14+[-1,1]\cdot[-\frac12,\frac12]=[-\frac14,\frac34]$, is still not exact,
but its overestimation is $O(w^2)$ in the width $w$ of the box, against $O(w)$ for the natural evaluation.
This is why bounds become sharp as boxes shrink, which is what branch and bound needs.

\paragraph{Certified zeros.} Interval arithmetic also proves that solutions exist.

\begin{proposition}[Interval Newton]\label{COMP:prop:newton}
Let $f\in C^1[a,b]$, $X=[a,b]$, $m\in X$, and suppose that the interval $f'(X)$ (any enclosure of $f'$ on
$X$) does not contain $0$. Put $N(X)=m-f(m)/f'(X)$. If $N(X)\subseteq X$, then $f$ has exactly one zero in $X$,
and it lies in $N(X)$.
\end{proposition}

\begin{proof}
Since $f'\ne0$ on $X$, $f$ is strictly monotone, so there is at most one zero. Say $f'>0$ and $f(m)\ge0$ (the
other cases are symmetric). By the mean value theorem $f(a)=f(m)-f'(\xi)(m-a)$ for some $\xi\in X$, and
$m-f(m)/f'(\xi)\in N(X)\subseteq X$ gives $f(m)/f'(\xi)\le m-a$, so $f(a)\le0$. By the intermediate value
theorem there is a zero $x^*$ in $[a,m]$. Finally $x^*=m-f(m)/f'(\eta)$ for some $\eta$ between $x^*$ and $m$,
so $x^*\in N(X)$.
\end{proof}

For the antiprism, take $f(c)=7c^2+2c-1$, $X=[0.26,0.27]$ and $m=0.265$. Then $f(m)=0.021575$,
$f'(X)=[5.64,5.78]$, and $N(X)=[0.26117\ldots,0.26126\ldots]\subset X$, so the Tammes cosine for $8$ points
lies in $N(X)$. One more step shrinks the interval to width $4\cdot10^{-9}$: Newton's quadratic convergence,
now with a proof attached. In several variables the same role is played by the \emph{Krawczyk operator}
$K(X)=m-Yf(m)+(I-Yf'(X))(X-m)$, with $Y$ an approximate inverse of $f'(m)$: if $K(X)$ lies in the interior of
$X$, then $f$ has a unique zero in $X$ (Krawczyk 1969; Moore 1977). This is how a numerically computed
balanced configuration is turned into a theorem.

\subsection{Dual numbers and jets}
To use centred forms or interval Newton one needs enclosures of derivatives on boxes. They come almost for
free from \emph{dual numbers}: compute in $\R[\varepsilon]/(\varepsilon^2)$. For a polynomial,
\[
  f(a+b\varepsilon)=f(a)+f'(a)\,b\,\varepsilon,
\]
because $\varepsilon^2=0$ kills everything else. So evaluating $f$ at $x+\varepsilon$ with the ordinary rules
$(a+a'\varepsilon)(b+b'\varepsilon)=ab+(ab'+a'b)\varepsilon$ produces $f(x)$ and $f'(x)$ together. This is
\emph{forward-mode automatic differentiation}. With $k$ variables one carries a gradient
$(a,a_1',\dots,a_k')$ instead; with $\R[\varepsilon]/(\varepsilon^{r+1})$ one carries the Taylor coefficients
up to order $r$, a \emph{jet}. Now let the entries be intervals. Evaluating on a box $B$ then gives, in one
pass, enclosures of $f(B)$ and of $\nabla f(B)$, by the fundamental theorem. We use exactly this in
Section~\ref{COMP:sec:thue}.

Going one step further, a \emph{Taylor model} (Berz and Makino, 1998) represents a function on a box by a
polynomial of degree $r$ together with an interval bounding the remainder. Arithmetic on Taylor models
keeps the dependency in the polynomial part, so the overestimation is $O(w^{r+1})$. The nonlinear inequalities
of the Kepler conjecture were proved with second-order versions of this idea (``Taylor intervals''), and
Solovyev and Hales (2013) implemented them inside the proof assistant HOL Light.

\subsection{Sums of squares}
A polynomial $p$ that is a sum of squares of polynomials is visibly nonnegative, and the squares are a
certificate. Write $p=m^TQm$ with $m$ the vector of monomials of degree at most $\frac12\deg p$. Then $p$ is
a sum of squares if and only if this holds for some positive semidefinite $Q$, since $Q=L^TL$ gives
$p=\abs{Lm}^2$. The conditions on $Q$ are linear equations (matching coefficients) plus $Q\succeq0$:
\emph{finding a sum-of-squares certificate is a semidefinite program} (Parrilo, 2000/2003; Lasserre, 2001).

Not every nonnegative polynomial is a sum of squares (Hilbert, 1888). Motzkin's
$x^4y^2+x^2y^4-3x^2y^2+1$ is nonnegative by the inequality between arithmetic and geometric means applied
to $x^4y^2,x^2y^4,1$, but it is not a sum of squares. To prove $p\ge0$ only on a set
$K=\{g_1\ge0,\dots,g_r\ge0\}$ one looks for certificates of the form
\[
  p=\sigma_0+\sum_i\sigma_ig_i\qquad(\text{Putinar, 1993}),\qquad\text{or}\qquad
  p=\sum_{I\subseteq\{1,\dots,r\}}\sigma_I\prod_{i\in I}g_i\qquad(\text{Schm\"udgen, 1991}),
\]
with sums of squares $\sigma$. Schm\"udgen showed that a polynomial strictly positive on a compact $K$ always
has the second kind of certificate, and Putinar that, under an extra hypothesis on $K$ (Archimedean), it has the
first. The Positivstellensatz of Krivine (1964) and Stengle (1974) gives certificates for the
\emph{emptiness} of a semialgebraic set, and so for nonnegativity in general, at the price of denominators.
Compare: the $\sigma_i$ are polynomial versions of the KKT multipliers of Section~\ref{COMP:sec:kkt}, and the
multiplier $\lambda(z)$ of Section~\ref{COMP:sec:thue} is an example.

The semidefinite program is solved numerically, so the matrix $Q$ it returns is only approximately positive
semidefinite and only approximately satisfies the linear equations. Peyrl and Parrilo (2008) showed how to
round it to an \emph{exact} rational certificate when $p$ is strictly inside the cone: project the rounded $Q$
onto the affine space of exact solutions and check $Q\succeq0$ exactly, for example by an exact Cholesky
factorization. The same rounding idea turns the numerical solutions of the LP and SDP bounds of
Section~\ref{COMP:sec:matrix} into rigorous bounds.

\subsection{Formal proof}
A certificate still has to be checked by \emph{some} program, and programs have bugs. A \emph{proof
assistant} (HOL Light, Isabelle, Coq, Lean) reduces this risk to a small trusted kernel that checks every
logical step.
\begin{itemize}
\item \emph{Kepler.} Hales's proof was published in 2005 (with a joint paper of Hales and Ferguson). The
  refereeing could not certify all of the computer code, and Hales started the Flyspeck project to produce a
  formal proof. It was completed on 10 August 2014, in HOL Light and
  Isabelle. The formal proof checks more than $23\,000$ nonlinear inequalities (after subdividing domains),
  which took about $5000$ processor hours on the Microsoft Azure cloud, together with $43\,078$ linear
  programs and the tame graph enumeration of Nipkow, Bauer and Schultz.
\item \emph{Dimension $8$.} Viazovska's proof (Interlude~B) relies on two inequalities between
  modular forms. In her paper they were checked with computer help: the functions were replaced by truncated
  expansions at the two ends of the interval, the truncation errors were bounded rigorously, and the
  inequalities for the truncations were checked by a finite calculation. Romik (2023) later gave a proof
  without computer calculations. A Lean formalization of the whole dimension-$8$ theorem, begun by Hariharan
  and Viazovska in 2024, was completed in February 2026; its final stages were produced by the
  autoformalization system Gauss of Math, Inc.\ (Hariharan, Birkbeck, Lee, Ma, Mehta, Poiroux and
  Viazovska, arXiv:2604.23468). According to EPFL's announcement (March 2026) the dimension-$24$ case followed.
\item \emph{Outside this course.} The four colour theorem (Appel and Haken, 1976--77; a simpler proof by
  Robertson, Sanders, Seymour and Thomas, 1997) was formalized in Coq by Gonthier (2005, published 2008).
  Tucker's proof (1999, 2002) that the Lorenz equations have a strange attractor, Smale's 14th problem,
  is a classic of interval analysis: it combines rigorous integration of the flow with a normal form near
  the origin.
\end{itemize}

%======================================================================
\section{A certified proof of Thue's triangle lemma}
\label{COMP:sec:thue}

We now prove the inequality of Lemma~\ref{L1:lem:triangle-density} by computer, using branch and bound, exact interval
arithmetic, dual numbers and a KKT multiplier. Of course Lecture~1 proves it in two lines. The point is to
watch every piece of the machinery work on an example where we know the answer.

\paragraph{The polynomial form.} Let a triangle have sides $a,b,c\ge2$ and circumradius $R\le2$, and let
$c$ be the longest side. Put $x=a^2$, $y=b^2$, $z=c^2$. Heron's formula says
\[
  16\operatorname{Area}^2=H(x,y,z):=2xy+2yz+2zx-x^2-y^2-z^2,
\]
and $R=abc/(4\operatorname{Area})$ gives $R^2=xyz/H$, so $R\le2$ means $g:=4H-xyz\ge0$. Since $c\le2R\le4$,
the domain is
\[
  D=\{4\le x\le z,\ 4\le y\le z,\ z\le16,\ g\ge0\},
\]
and the inequality of Lemma~\ref{L1:lem:triangle-density} (area $\ge\sqrt3$, i.e.\ density
$\le\pi/\sqrt{12}$) follows from $H\ge48$ on $D$. (We allow $R=2$; this only makes the claim stronger.) All transcendental functions are gone: this is
the first rule of thumb.

\paragraph{Where it is tight.} $H=48$ at the equilateral triangle $(4,4,4)$, a corner of $D$. It is also
$48$ at $(4,4,12)$, the triangle with sides $2,2,2\sqrt3$ and largest angle $2\pi/3$, which has $R=2$
exactly. (It is excluded by the strict inequality $R<2$ in Lecture~1, but it lies in the closure, so the
inequality on $D$ is sharp at two points.) Tight points are the enemy of branch and bound: near them no bound
of the form ``$H\ge48$ on this box'' can be proved with any slack, and a naive search never ends.

\paragraph{The KKT multiplier.} At $(4,4,12)$ the constraint $g\ge0$ is active. Solving the KKT
equations $\nabla H=\mu_x e_x+\mu_ye_y+\lambda\nabla g$ there gives $\nabla H=(24,24,-8)$,
$\nabla g=(48,48,-48)$, hence $\lambda=\frac16$ and $\mu_x=\mu_y=16$. At $(4,4,4)$ the constraint $g$ is not
active ($g=128$) and its multiplier is $0$. On the edge $x=y=4$ that joins the two tight points,
\[
  H-48=(z-4)(12-z),\qquad g=4z(12-z),\qquad\text{so}\qquad H-48=\frac{z-4}{4z}\,g .
\]
This suggests the variable multiplier $\lambda(z)=\frac{z-4}{4z}$, which is $0$ at $z=4$ and $\frac16$ at
$z=12$. A two-line computation (checked symbolically in the script) gives the identity
\begin{equation}\label{COMP:eq:cert}
  H-48=\frac{P}{z}+\frac{z-4}{4z}\,g,\qquad
  P=\frac{xy-16}4\,z^2+\big(16-(x-8)(y-8)\big)z-4(x-y)^2 .
\end{equation}
We write $P=Az^2+Bz+C$ with $A=\frac{xy-16}4$, $B=16-(x-8)(y-8)$, $C=-4(x-y)^2$. Since $\lambda(z)\ge0$ and
$g\ge0$ on $D$, it suffices to prove
\begin{equation}\label{COMP:eq:Pclaim}
  P(x,y,z)\ge0\qquad\text{on the polyhedron } \{4\le x\le z,\ 4\le y\le z\le16\}.
\end{equation}
The curved constraint has disappeared. The price: $P$ vanishes on the whole edge $x=y=4$.

\paragraph{The search.} We prove \eqref{COMP:eq:Pclaim} by branch and bound over boxes, starting from
$[4,16]^3$ and bisecting the longest side. A box is \emph{settled} if one of three tests succeeds:
\begin{enumerate}[label=(\roman*),nosep]
\item it misses the region ($z<x$ or $z<y$ everywhere on the box);
\item \emph{value test:} the interval evaluation of $P$ on the box has lower end $\ge0$;
\item \emph{derivative test:} interval enclosures of $\partial P/\partial x$ and $\partial P/\partial y$
  (from dual numbers) have lower ends $\ge0$, so $P$ is smallest on the face $x=\underline x$, $y=\underline y$,
  and the interval evaluation of $P$ on that face has lower end $\ge0$.
\end{enumerate}
Test (iii) is what handles the tight edge: on it, $A=B=C=0$ exactly, so the evaluation on the face gives
the interval $[0,0]$. All arithmetic is in exact rationals (Python's
\texttt{fractions.Fraction}), so there is no rounding at all. Here is the heart of
\texttt{code/thue\_bb.py}, slightly abridged:
{\footnotesize
\begin{verbatim}
from fractions import Fraction as Fr
class I:                                   # closed interval, rational endpoints
    def __init__(s, lo, hi=None): s.lo, s.hi = Fr(lo), Fr(lo if hi is None else hi)
    def __add__(s, t): t = lift(t); return I(s.lo + t.lo, s.hi + t.hi)
    def __neg__(s): return I(-s.hi, -s.lo)
    def __mul__(s, t):
        t = lift(t); p = [s.lo*t.lo, s.lo*t.hi, s.hi*t.lo, s.hi*t.hi]
        return I(min(p), max(p))
    def sq(s):                             # a square is >= 0
        lo, hi = sorted([abs(s.lo), abs(s.hi)])
        return I(0 if s.lo <= 0 <= s.hi else lo*lo, hi*hi)
class D:                                   # dual number: value and gradient, all intervals
    def __init__(s, v, g): s.v, s.g = v, g
    def __add__(s, t): t = liftD(t); return D(s.v + t.v, [a + b for a, b in zip(s.g, t.g)])
    def __mul__(s, t):                     # Leibniz rule
        t = liftD(t); return D(s.v*t.v, [s.v*b + t.v*a for a, b in zip(s.g, t.g)])
    def sq(s): return s*s
def P(x, y, z):
    A = Fr(1, 4)*(x*y - 16); B = 16 - (x - 8)*(y - 8); C = -4*(x - y).sq()
    return A*z.sq() + B*z + C
def grad(X, Y, Z):
    e = lambda k: [I(1) if j == k else I(0) for j in range(3)]
    return P(D(X, e(0)), D(Y, e(1)), D(Z, e(2))).g
def check(X, Y, Z):
    if Z.hi < X.lo or Z.hi < Y.lo: return True             # (i)   outside the region
    if P(X, Y, Z).lo >= 0: return True                      # (ii)  value test
    gx, gy, _ = grad(X, Y, Z)                               # (iii) derivative test
    return gx.lo >= 0 and gy.lo >= 0 and P(I(X.lo), I(Y.lo), Z).lo >= 0
def bb(box):                    # returns only if P >= 0 is proved on the box
    if check(*box): return
    k = max(range(3), key=lambda j: box[j].hi - box[j].lo)
    m = (box[k].lo + box[k].hi) / 2
    for half in (I(box[k].lo, m), I(m, box[k].hi)):
        b = list(box); b[k] = half; bb(tuple(b))
bb((I(4, 16), I(4, 16), I(4, 16)))
\end{verbatim}
}
\noindent (Here \texttt{lift} and \texttt{liftD} turn numbers into point intervals and constant dual
numbers; subtraction and the reversed operations are defined in the same way.)

\paragraph{The result.} The program terminates in under a second. It examines $53$ boxes; the $27$ leaves
are settled by test (i) in $9$ cases, by test (ii) in $15$ and by test (iii) in $3$; the deepest box is at
depth $8$. Since every leaf passed a valid test and the leaves cover $[4,16]^3$, this proves
\eqref{COMP:eq:Pclaim}, hence $H\ge48$ on $D$, hence the inequality of
Lemma~\ref{L1:lem:triangle-density}. (The equality case is not addressed by the search.)

Without test (iii) the search never ends. With the depth capped at $3k$, the number of unsettled boxes
is $8,16,32,64,128,256$ for $k=3,\dots,8$, all of them strung along the edge $x=y=4$: each halving of the
box size doubles the work. Figure~\ref{COMP:fig:bb} shows both runs on the slice $z=6$.

\colorlet{COMPcolval}{addcol!25}\colorlet{COMPcolmono}{green!45!black!35}
\colorlet{COMPcolout}{gray!18}\colorlet{COMPcolopen}{flagcol!70}
\begin{figure}[ht]
\centering
\begin{tikzpicture}
  % panel A: certified run, slice z=6, coordinates (x,y) in [4,16]^2, 0.36 cm per unit
  \begin{scope}[x=0.36cm,y=0.36cm,shift={(-4,-4)}]
    \foreach \t/\a/\b/\c/\d in {
      mono/4/5.5/4/7, val/5.5/7/4/7, val/4/5.5/7/8.5, out/4/5.5/8.5/10, val/5.5/7/7/10, val/7/8.5/4/7,
      out/8.5/10/4/7, val/7/10/7/10, out/4/7/10/13, out/4/7/13/16, val/7/10/10/13, out/7/10/13/16,
      out/10/13/4/7, val/10/13/7/10, out/13/16/4/10, val/10/13/10/13, out/10/13/13/16, out/13/16/10/16} {
      \draw[thin,fill=COMPcol\t] (\a,\c) rectangle (\b,\d);
    }
    \draw[thick,dashed] (6,4)--(6,6)--(4,6);
    \fill[black] (4,4) circle (2.2pt);
    \draw[thick] (4,4) rectangle (16,16);
    \node[font=\scriptsize,below] at (4,4) {$4$}; \node[font=\scriptsize,below] at (16,4) {$16$};
    \node[font=\scriptsize,left] at (4,16) {$16$};
    \node[font=\small,below] at (10,3.3) {$x$}; \node[font=\small,left] at (3.2,10) {$y$};
    \node[font=\small,align=center] at (10,18) {certified run (all tests)};
  \end{scope}
  % panel B: no derivative test, depth cap 18, zoom on [4,7]^2, 1.45 cm per unit
  \begin{scope}[x=1.45cm,y=1.45cm,xshift=6.6cm,shift={(-4,-4)}]
    \clip (3.98,3.98) rectangle (7.02,7.02);
    \foreach \t/\a/\b/\c/\d in {
      open/4.0/4.1875/4.0/4.1875, val/4.0/4.1875/4.1875/4.375, val/4.1875/4.375/4.0/4.375,
      val/4.0/4.375/4.375/4.75, val/4.375/4.75/4.0/4.75, val/4.0/4.75/4.75/5.5, val/4.75/5.5/4.0/5.5,
      val/4.0/5.5/5.5/7.0, val/5.5/7.0/4.0/7.0, val/5.5/7.0/7.0/10.0, val/4.0/5.5/7.0/8.5} {
      \draw[thin,fill=COMPcol\t] (\a,\c) rectangle (\b,\d);
    }
    \draw[thick,dashed] (6,4)--(6,6)--(4,6);
    \fill[black] (4,4) circle (2.2pt);
  \end{scope}
  \begin{scope}[x=1.45cm,y=1.45cm,xshift=6.6cm,shift={(-4,-4)}]
    \draw[thick] (4,4) rectangle (7,7);
    \node[font=\scriptsize,below] at (4,4) {$4$}; \node[font=\scriptsize,below] at (7,4) {$7$};
    \node[font=\scriptsize,left] at (4,7) {$7$};
    \node[font=\small,below] at (5.5,3.83) {$x$}; \node[font=\small,left] at (3.8,5.5) {$y$};
    \node[font=\small,align=center] at (5.5,7.45) {without (iii), depth $18$ (zoom)};
  \end{scope}
\end{tikzpicture}
\\[4pt]
{\small
\tikz\draw[fill=COMPcolval] (0,0) rectangle (0.3,0.2); value test \quad
\tikz\draw[fill=COMPcolmono] (0,0) rectangle (0.3,0.2); derivative test \quad
\tikz\draw[fill=COMPcolout] (0,0) rectangle (0.3,0.2); outside the region \quad
\tikz\draw[fill=COMPcolopen] (0,0) rectangle (0.3,0.2); unsettled}
\caption{Branch and bound for \eqref{COMP:eq:Pclaim}, shown on the slice $z=6$ (the boxes are
three-dimensional; we draw their intersections with this plane). The dashed square is the region
$x,y\le z=6$, and the dot is the tight edge $x=y=4$. Left: the certified run, with $18$ boxes on this
slice. Right: without the derivative test the boxes pile up geometrically towards the tight edge, and the
corner box is never settled.}
\label{COMP:fig:bb}
\end{figure}

\begin{remark}[What was learned]
(1) The tight points dictated the method: a strict inequality would have been proved by test (ii)
alone, and the multiplier would not have been needed. Every sharp inequality in Hales's proof needed a
local argument near its equality case of exactly this kind. (2) The multiplier $\lambda(z)$ is a
Positivstellensatz-type certificate in the sense of Section~\ref{COMP:sec:verify}: \eqref{COMP:eq:cert} reads
$z(H-48)=P+\frac{z-4}4\,g$, and only the nonnegativity of $P$ on a polyhedron is left to the computer. A
semidefinite program could look for a sum-of-squares certificate for $P$ too (Exercise~\ref{COMP:exc:sos}).
(3) The derivative test is the monotonicity argument of Lecture~1 in disguise: there the area
$\frac12ac\sin B$ is increased by lengthening $a$ and $c$.
\end{remark}

%======================================================================
\section{Gallery of theorems}
\label{COMP:sec:gallery}

The theorems below were proved with essential help from computers, or are certified by finite objects
that were found by computer. Each comes with a sentence on how it fits the course.

\begin{theorem}[Hales, Ann.\ of Math.\ 162 (2005); Hales--Ferguson, Discrete Comput.\ Geom.\ 36 (2006)]
No packing of congruent balls in $\R^3$ has density greater than $\pi/\sqrt{18}$.
\end{theorem}
\noindent The formal proof is T.~Hales et al., \emph{A formal proof of the Kepler conjecture}, Forum Math.\ Pi 5
(2017), e2. Cases are indexed by tame plane graphs (Section~\ref{COMP:sec:index}) and settled by linear programs
and interval arithmetic; the whole of this interlude is a miniature of it.

\begin{theorem}[Hales--McLaughlin, J.\ Amer.\ Math.\ Soc.\ 23 (2010)]
Every Voronoi cell of a packing of unit balls in $\R^3$ has volume at least that of the regular dodecahedron
of inradius $1$.
\end{theorem}
\noindent Fejes T\'oth's dodecahedral conjecture (Interlude~E, Section~\ref{SM:sec:frustration}), proved by the
same computer-assisted methodology as the Kepler conjecture.

\begin{theorem}[Musin--Tarasov, Discrete Comput.\ Geom.\ 48 (2012); Experiment.\ Math.\ 24 (2015)]
The optimal Tammes configurations for $N=13$ and $N=14$ are unique up to isometry, with minimal angular
distances $\approx57.1367^\circ$ and $\approx55.6706^\circ$.
\end{theorem}
\noindent The proofs enumerate irreducible contact graphs: the full-scale version of the five-point example of
Section~\ref{COMP:sec:index}. The case $N=13$ is the ``strong thirteen spheres problem'': how large can $13$ equal balls touching a unit ball
be? Since $57.1367^\circ<60^\circ$, it contains a new proof of $\kiss_3=12$.

\begin{theorem}[Odlyzko--Sloane, J.\ Combin.\ Theory Ser.~A 26 (1979); Levenshtein, Soviet Math.\ Dokl.\ 20 (1979)]
$\kiss_8=240$ and $\kiss_{24}=196\,560$.
\end{theorem}
\noindent Example~\ref{COMP:ex:240}. The computer found the polynomials; the proof is checkable by hand.
Bannai and Sloane (Canad.\ J.\ Math.\ 33, 1981) proved uniqueness of the configurations.

\begin{theorem}[Musin, Ann.\ of Math.\ 168 (2008)]
$\kiss_4=24$.
\end{theorem}
\noindent The plain LP bound gives only $25$. Musin modifies Delsarte's method so that the auxiliary polynomial
may be positive on part of the range, and controls that part by a separate geometric argument.

\begin{theorem}[Bachoc--Vallentin, J.\ Amer.\ Math.\ Soc.\ 21 (2008)]
$\kiss_5\le45$, $\kiss_6\le78$, $\kiss_7\le135$, $\kiss_9\le366$, $\kiss_{10}\le567$.
\end{theorem}
\noindent Three-point semidefinite bounds, a matrix version of Theorem~\ref{COMP:thm:LP}. N.~Leijenhorst and
D.~de Laat (Math.\ Program.\ Comput., 2024) improved the bounds in dimensions $11$--$23$, and
H.~Cohn, D.~de Laat and A.~Salmon (\emph{Three-point bounds for sphere packing}, arXiv:2206.15373) gave the
analogous improvement of the Cohn--Elkies bound for packings in dimensions $4$--$7$ and $9$--$16$.

\begin{theorem}[Mark\'ot--Csendes, SIAM J.\ Optim.\ 16 (2005); Mark\'ot, J.\ Global Optim.\ 81 (2021)]
The best known packings of $n$ equal circles in a square are optimal for $n=28,\dots,33$.
\end{theorem}
\noindent Interval branch and bound in configuration space, the method of Section~\ref{COMP:sec:bb} with a
great deal of engineering.

\begin{theorem}[Brakensiek--Heule--Mackey--Narv\'aez, IJCAR 2020; J.\ Automat.\ Reason.\ 66 (2022)]
Every tiling of $\R^7$ by unit cubes contains two cubes sharing a full face (Keller's conjecture holds in
dimension $7$).
\end{theorem}
\noindent Keller's conjecture was known for $n\le6$ (Perron) and false for $n\ge8$ (Lagarias--Shor for $n\ge10$, Mackey
for $n=8$). The case $n=7$ was reduced to the non-existence of a clique of size $128$ in certain graphs and settled by
a SAT solver; the proofs of unsatisfiability (about $200$ gigabytes) were checked by an independent, formally
verified checker. A packing problem for cubes, settled by a certificate of non-existence.

\begin{theorem}[Appel--Haken 1977; Robertson--Sanders--Seymour--Thomas, J.\ Combin.\ Theory Ser.~B 70 (1997); Gonthier, Notices AMS 55 (2008)]
Every planar graph is $4$-colourable.
\end{theorem}
\noindent The model case of ``cases indexed by planar graphs'': an unavoidable set of configurations ($633$ in the
1997 proof), each shown reducible by computer. Gonthier's Coq proof removed the doubts about the code.

\begin{theorem}[Tucker, C.~R.\ Acad.\ Sci.\ Paris 328 (1999); Found.\ Comput.\ Math.\ 2 (2002)]
The Lorenz equations, with the classical parameters, have a robust strange attractor.
\end{theorem}
\noindent Smale's 14th problem, solved by interval arithmetic: an example far from geometry of the methods of
Section~\ref{COMP:sec:verify}.

\begin{theorem}[Lower bounds by construction, 2025]
$\kiss_{11}\ge593$ (AlphaEvolve, Google DeepMind, arXiv:2506.13131); $\kiss_{10}\ge510$, $\kiss_{14}\ge1932$
(M.~Ganzhinov, \emph{Highly symmetric lines}, Linear Algebra Appl.\ 722 (2025); arXiv:2207.08266).
\end{theorem}
\noindent A lower bound is certified by a list of vectors and an exact check of all pairwise inner products.
Finding the list is the hard part; checking it is easy. This is the opposite situation to the upper bounds above.

%======================================================================
\section*{Exercises}
\addcontentsline{toc}{section}{Exercises}

\begin{exercise}
Using Lemma~\ref{COMP:lem:angle}, show that in a configuration with $\rho\ge\pi/3$ every point has at most five
contacts, and that the bound is still five for the Tammes configuration of $13$ points
($\theta\approx57.14^\circ$). Show that the bound drops to four exactly when $\theta>\arctan2$, the edge of the
regular icosahedron.
\end{exercise}

\begin{exercise}
Modify the program \texttt{code/contact5.py} to run the rules R1--R4 on the $156$ graphs on six vertices.
Which graphs survive?
Find the critical configurations whose stressed graphs are the octahedron, the hexagon $C_6$, and $C_5$ plus an
isolated point. Are there others?
\end{exercise}

\begin{exercise}
Prove the equality case of Proposition~\ref{COMP:prop:rankin}(b): if $2d$ nonzero vectors in $\R^d$ have pairwise
non-positive inner products, they are $\pm$ an orthogonal basis up to positive scaling. Deduce that the octahedron
is the unique optimal configuration for $N=6$, while for $N=5$ the optimum is not unique.
\end{exercise}

\begin{exercise}
Verify Example~\ref{COMP:ex:240}: expand $f$ in Gegenbauer polynomials for $n=8$ and compute $f(1)/c_0$. Then
use the same polynomial in dimension $3$. Show that its coefficients are still non-negative, and that it only
proves $\kiss_3\le1890/127\approx14.9$. Why should one expect a polynomial with roots at $-\frac12,0,\frac12$ to
be good in dimension $8$ but not in dimension $3$?
\end{exercise}

\begin{exercise}\label{COMP:exc:kkt}
For $N=4$, compute the KKT multipliers of Proposition~\ref{COMP:prop:kkt} at the regular tetrahedron and at the
square on a great circle. Show that the square is not first-order jammed: find a nontrivial infinitesimal motion $v$ with $Av=0$, and
check that along the corresponding curve all four contact distances increase to second order. Relate this to
Example~\ref{TOP:ex:N4pass} of Interlude~D.
\end{exercise}

\begin{exercise}
(a) Compute the natural interval extension of $f(x)=x^2-2x$ on $[0,2]$, and of the equivalent forms $x(x-2)$ and
$(x-1)^2-1$. (b) Use Proposition~\ref{COMP:prop:newton} with exact rational arithmetic to enclose $\sqrt2$ in an
interval of width less than $10^{-20}$, starting from $[1,2]$.
\end{exercise}

\begin{exercise}\label{COMP:exc:sos}
(a) Check that the three coefficients $A,B,C$ of $P$ vanish on the edge $x=y=4$.
Compute $\partial P/\partial x$ and $\partial P/\partial y$, and show by interval arithmetic that both are
$\ge0$ on the box $[4,5.5]\times[4,7]\times[4,7]$, one of the three boxes settled by test (iii). (b) Using an SDP solver,
look for a certificate $P=\sigma_0+\sum_i\sigma_ig_i$ on the polyhedron of \eqref{COMP:eq:Pclaim} with
$g\in\{x-4,\,y-4,\,z-x,\,z-y,\,16-z\}$, and round it to exact rationals.
\end{exercise}

\begin{exercise}
Write the statement ``there are $5$ points on $\Sph^2$ with pairwise distances greater than $\pi/2$'' as a
first-order sentence. How many variables and polynomials does it have? Using
Proposition~\ref{COMP:prop:rankin}, decide it by hand. Why does the existence of a decision procedure imply that
$\cos d_N$ is algebraic for every $N$?
\end{exercise}

\begin{exercise}
Modify \texttt{code/e8\_mitm.py} to count the vectors of norm $2$ and $4$ in $D_4$ (split $2+2$) and in $E_7$
(take $E_7$ to be the vectors of $E_8$ orthogonal to the root $e_1-e_2$, and keep the first two coordinates in the
same half). You should find $24,24$ and $126,756$.
\end{exercise}

%======================================================================
\section*{Further reading}
\addcontentsline{toc}{section}{Further reading}
\begin{itemize}[nosep]
\item R.~E.~Moore, R.~B.~Kearfott, M.~J.~Cloud, \emph{Introduction to Interval Analysis}, SIAM, 2009. Interval
  arithmetic, interval Newton and Krawczyk, with code.
\item W.~Tucker, \emph{Validated Numerics: A Short Introduction to Rigorous Computations}, Princeton University
  Press, 2011. Includes automatic differentiation and the Lorenz attractor.
\item G.~Blekherman, P.~A.~Parrilo, R.~R.~Thomas (eds.), \emph{Semidefinite Optimization and Convex Algebraic
  Geometry}, MOS--SIAM Series on Optimization, 2013. Sums of squares, Positivstellens\"atze, SDP.
\item S.~Basu, R.~Pollack, M.-F.~Roy, \emph{Algorithms in Real Algebraic Geometry}, 2nd ed., Springer, 2006.
  Sturm sequences, cylindrical algebraic decomposition, complexity.
\item D.~Cox, J.~Little, D.~O'Shea, \emph{Ideals, Varieties, and Algorithms}, Springer UTM. Gr\"obner bases.
\item J.~H.~Conway, N.~J.~A.~Sloane, \emph{Sphere Packings, Lattices and Groups}, Springer; Ch.~9 (bounds for
  codes and sphere packings), Ch.~13 (Odlyzko--Sloane, bounds on kissing numbers) and Ch.~14 (Bannai--Sloane,
  uniqueness).
\item T.~C.~Hales et al., \emph{A formal proof of the Kepler conjecture}, Forum Math.\ Pi 5 (2017), e2; and
  T.~C.~Hales, \emph{Dense Sphere Packings: A Blueprint for Formal Proofs}, LMS Lecture Note Series 400, Cambridge,
  2012.
\item J.~C.~Lagarias (ed.), \emph{The Kepler Conjecture: The Hales--Ferguson Proof}, Springer, 2011.
\item H.~Cohn, \emph{Packing, coding, and ground states}, PCMI lecture notes, arXiv:1603.05202. LP and SDP bounds
  from the point of view of these notes.
\end{itemize}
% ===== end chapters/COMP.tex =====
\endgroup
\begingroup\renewcommand{\thechapter}{G}\renewcommand{\theHchapter}{interludeG}
% ===== begin chapters/PHYS.tex =====
% MAT670 -- Interlude G (placed after Interlude F): packings in the laboratory.
% Entirely new material (not in the handwritten notes). Labels prefixed PHYS:.
% Numerical claims checked with python3 (cell areas, Coxeter's count, Bragg spacings, Brillouin-zone
% distances, the free-volume model, Letokhov's radius, the n=6 sticky-cluster partition functions).

\newcommand{\PHYStag}[1]{\textsf{\scriptsize[#1]}}
\tikzset{
  PHYSth/.style={fill=black,draw=black},
  PHYSball/.style={fill=addcol!25,draw=addcol,fill opacity=0.45},
}
\newcommand{\PHYSmth}[1]{\fill[black] #1 ++(-2.2pt,-2.2pt) rectangle ++(4.4pt,4.4pt);}
\newcommand{\PHYSmsim}[1]{\filldraw[draw=black,fill=white,thick] #1 circle (2.3pt);}
\newcommand{\PHYSmexp}[1]{\fill[flagcol] #1 ++(0pt,3.2pt) -- ++(-2.9pt,-5.2pt) -- ++(5.8pt,0pt) -- cycle;}
\newcommand{\PHYSmcon}[1]{\fill[addcol] #1 ++(0pt,3pt) -- ++(3pt,-3pt) -- ++(-3pt,-3pt) -- ++(-3pt,3pt) -- cycle;}

\chapter[Interlude G: Packings in the lab]{Interlude G: Packings in the laboratory}
\label{PHYS:ch}

\begin{addedblock}
This interlude is new: it is not in the handwritten notes. It follows the aside at the end of
Lecture~6 (``ellipsoid packings; polygon packings; honeycomb and bubble problems'') and the physical
remarks scattered through Lectures~7--11 out of the mathematics and into the laboratory: squashed peas,
foams, colloids, opals, lasers made of sand, sticky spheres and dice. To keep it readable it is set in
black; all of it should be read as added material.
\end{addedblock}

Everything in these notes so far has been about ideal balls: exactly round, exactly hard, exactly
equal. Physicists have been squeezing, shaking and photographing real ones for three centuries, and
many of the objects of the course were met on a laboratory bench before they were met in a proof. The
Voronoi cell was a squashed pea, the face-centred cubic lattice was an opal, and the contact graph was
a pattern of paint marks on ball bearings.

In physics one has to be careful about the status of a statement, more careful than in the rest of the
course. Throughout, claims are tagged \PHYStag{theorem} (proved), \PHYStag{simulation} (seen on a
computer), \PHYStag{experiment} (measured) or \PHYStag{model} (a calculation in an idealized model).
Most of the famous numbers of the subject (the random close packing fraction $0.64$, the freezing
point $0.494$ of hard spheres, the jamming exponents) are of the second and third kind; Figure~\ref{PHYS:fig:strip} sorts them. Interlude~E
discussed the statistical mechanics behind several of them (hard spheres, random close packing, Maxwell
counting, frustration, Edwards' entropy), and we refer to it rather than repeat it.

%======================================================================
\section{Squashed peas, lead shot and soap}
\label{PHYS:sec:squash}

\subsection*{The first Hales}
Lecture~1 introduced the Voronoi cell (Definition~\ref{L1:def:voronoi}) as a bookkeeping device: it
assigns to each ball the part of space that is closer to it than to any other ball. There is a much
more physical way to see Voronoi cells. Take soft balls, and press them together until there is no
space left. If every ball keeps its centre and grows at the same rate, each one ends up filling exactly
its Voronoi cell.

The first famous squashing experiment was made by a Hales, although not the one who proved the
Kepler conjecture. Stephen Hales, a clergyman and physiologist, reported in his \emph{Vegetable
Staticks} (1727):
\begin{quote}\small
I compressed several fresh parcels of Pease in the same Pot, with a force equal to 1600, 800, and 400
pounds; \dots\ being thereby formed into pretty regular Dodecahedrons.
\end{quote}
(The peas swelled with water inside an iron pot, under a weighted lid.) What could ``pretty regular
dodecahedrons'' have been? Not regular pentagonal dodecahedra, as a short argument shows.

\begin{proposition}\label{PHYS:prop:dodec}
Congruent regular dodecahedra do not tile $\R^3$.
\end{proposition}

\begin{proof}
The dihedral angle of the regular dodecahedron is $\delta=\arccos(-1/\sqrt5)\approx116.57^\circ$.
Suppose the dodecahedra tile, and let $e$ be an edge of one tile $T$. The edges and vertices of the
other tiles meet the open segment $e$ in finitely many points or in subsegments of $e$, so we may pick
a point $p\in e$ that is not a vertex of any tile and at which no edge of another tile crosses $e$
transversally. Let $T'\ne T$ be a tile containing $p$. If $p$ lies in the relative interior of a face
of $T'$, that face must contain the direction of $e$: otherwise the face plane would cross $e$ at $p$
and the half-space bounded by it would overlap the interior of $T$ near $p$. If $p$ lies on an edge of
$T'$, then that edge runs along $e$. So in a small disc perpendicular to $e$ at $p$, the tiles cut out
sectors of angle $\delta$ (wedges along $e$) and $\pi$ (half-planes), with disjoint interiors, filling
the disc. Hence $k\delta+m\pi=2\pi$ for integers $k\ge1$, $m\ge0$. Since $\pi/2<\delta<\pi$, the case
$m=1$ is impossible, and so is $m=2$. For $m=0$ we would need $k\delta=2\pi$, but
$3\delta\approx349.7^\circ$ and $4\delta\approx466.3^\circ$.
\end{proof}

So three regular dodecahedra around an edge leave a gap of about $10.3^\circ$. This is the
dodecahedral counterpart of the $7.36^\circ$ gap left by five regular tetrahedra around an edge
(Section~\ref{TET:sec}), and the two facts are dual to each other: a regular dodecahedron is the
Voronoi cell of a ball surrounded by a regular icosahedron of twelve neighbours, whose Delaunay cells
would be the twenty slightly flattened tetrahedra of Interlude~E. By the theorem of Hales and
McLaughlin quoted there, no Voronoi cell of a unit ball packing is smaller than this dodecahedron, and
by Proposition~\ref{PHYS:prop:dodec} not all of them can be equal to it. The peas must have been
something else.

\subsection*{Kepler's pomegranate}
A century earlier, Kepler had already given the right answer. In \emph{Strena seu de nive sexangula}
(1611), the little book on snowflakes that also contains his packing conjecture, he describes the
seeds of a pomegranate, squeezed together as the fruit grows, and the rhombic shape they take: the
rhombic dodecahedron. Here is the computation behind it.

\begin{proposition}[The Voronoi cell of FCC]\label{PHYS:prop:rd}
Let $D_3=\{x\in\Z^3:x_1+x_2+x_3\ \text{even}\}$ be the FCC lattice, with minimal distance $\sqrt2$. Its
Voronoi cell at $0$ is
\[
  V=\{x\in\R^3:\ \abs{x_i}+\abs{x_j}\le1\ \text{for all }i\ne j\}.
\]
This is the rhombic dodecahedron: the cube $[-\frac12,\frac12]^3$ with a square pyramid of height
$\frac12$ erected on each face. It has $12$ rhombic faces, one for each minimal vector $\pm e_i\pm e_j$, $6$
vertices of degree four at $(\pm1,0,0),(0,\pm1,0),(0,0,\pm1)$, and $8$ vertices of degree three at
$(\pm\frac12,\pm\frac12,\pm\frac12)$. Its volume is $2=\covol(D_3)$.
\end{proposition}

\begin{proof}
The Voronoi cell is $\{x:\ip xv\le\frac12\abs v^2\ \forall v\in D_3\setminus0\}$. For the $12$ minimal
vectors $v=\pm e_i\pm e_j$ the condition reads $\pm x_i\pm x_j\le1$, so $V$ contains the Voronoi cell. For
the converse we show that every other $v\in D_3$ gives a redundant condition. Let $x\in V$ and put
$a_i=\abs{x_i}$; then $a_i+a_j\le1$, so each $a_i\le1$ and at most one $a_i$ exceeds $\frac12$. With
$m_i=\abs{v_i}\in\Z_{\ge0}$,
\[
  \tfrac12\abs v^2-\ip xv\ \ge\ \sum_i\Big(\tfrac12m_i^2-a_im_i\Big).
\]
A term with $m_i=0$ is $0$, and a term with $m_i\ge2$ is at least $\frac12m_i(m_i-2)\ge0$. A term with
$m_i=1$ is $\frac12-a_i$, which is negative for at most one index $i_0$. If that happens, $v_{i_0}$ is odd,
and since $v_1+v_2+v_3$ is even, some other coordinate $v_j$ is odd too. If $m_j=1$, the two terms add up
to $1-a_{i_0}-a_j\ge0$; if $m_j\ge3$, the $j$-term is at least $\frac32$ and the $i_0$-term at least
$-\frac12$. Either way the sum is $\ge0$. So $V$ is the Voronoi cell. (For every simply laced root lattice
the same holds, with the roots in place of the $\pm e_i\pm e_j$; see Lemma~\ref{DYN:lem:relevant} of
Interlude~H.)

On the cube $\max\abs{x_i}\le\frac12$ all the inequalities hold. Over the face $x_3=\frac12$ of the cube,
the inequalities $\abs{x_1}+x_3\le1$ and $\abs{x_2}+x_3\le1$ cut out the pyramid with apex $(0,0,1)$, and
similarly for the other faces. The volume is $1+6\cdot\frac13\cdot1\cdot\frac12=2$. Each face
$\{\pm x_i\pm x_j=1\}\cap V$ is a rhombus with two degree-three and two degree-four vertices, for instance
$(1,0,0),(\frac12,\frac12,\pm\frac12),(0,1,0)$ on the face $x_1+x_2=1$.
\end{proof}

The two kinds of vertices are the two kinds of holes of FCC (Interlude~H, Section~\ref{DYN:sec:fcc}): the
degree-four vertices, at distance $1$ from the centre, are octahedral holes, and the degree-three ones,
at distance $\frac{\sqrt3}{2}$, are tetrahedral holes. The rhombic faces have diagonals $\sqrt2$ and $1$.
In the hexagonal close packing (HCP) the layers above and below each ball are mirror images of each
other, and the cell becomes the \emph{trapezo-rhombic dodecahedron}: the rhombic dodecahedron cut across a
three-fold axis, with one half turned by $60^\circ$ (six rhombi and six trapezoids, the same volume). Both
cells have twelve faces, one for each of the twelve kissing neighbours. Whatever Stephen Hales saw, the
experiment below suggests that most of his peas were less regular than either.

\begin{figure}[ht]
\centering
\begin{minipage}{0.47\textwidth}
\centering
\begin{tikzpicture}[x={(1.55cm,0cm)},y={(0.62cm,0.42cm)},z={(0cm,1.55cm)},line join=round]
  \foreach \a in {0.5,-0.5} \foreach \b in {0.5,-0.5} {
    \draw[dashed,gray] (-0.5,\a,\b)--(0.5,\a,\b);
    \draw[dashed,gray] (\a,-0.5,\b)--(\a,0.5,\b);
    \draw[dashed,gray] (\a,\b,-0.5)--(\a,\b,0.5);}
  \foreach \h in {0.5,-0.5} \foreach \a in {0.5,-0.5} \foreach \b in {0.5,-0.5} {
    \draw[thick] ({2*\h},0,0)--(\h,\a,\b);
    \draw[thick] (0,{2*\h},0)--(\a,\h,\b);
    \draw[thick] (0,0,{2*\h})--(\a,\b,\h);}
  \fill[flagcol] (0,0,1) circle (2.2pt) node[above,font=\scriptsize,black] {$(0,0,1)$: octahedral hole};
  \fill[addcol] (0.5,-0.5,0.5) circle (2.2pt) node[right,font=\scriptsize,black,xshift=2pt] {$\tfrac12(1,-1,1)$: tetrahedral hole};
  \fill (0,0,0) circle (1.2pt);
\end{tikzpicture}
\end{minipage}\hfill
\begin{minipage}{0.5\textwidth}
\centering\small
\renewcommand{\arraystretch}{1.15}
\begin{tabular}{lcc}
\toprule
cell & faces & $A/V^{2/3}$\\
\midrule
ball (does not tile) & -- & $4.836$\\
rhombic dodecahedron (FCC) & $12$ & $5.345$\\
truncated octahedron (BCC) & $14$ & $5.315$\\
Kelvin (curved faces) & $14$ & $5.306$\\
Weaire--Phelan (average) & $13.5$ & $5.288$\\
cube & $6$ & $6$\\
\bottomrule
\end{tabular}
\end{minipage}
\caption{Left: the Voronoi cell of FCC is a cube (dashed) with a pyramid on each face. Right: surface area
per cell for cells of unit volume. The flat polyhedra were computed for this table; the Kelvin and
Weaire--Phelan values are numerical minimizations (Surface Evolver), as quoted by Weaire and
Phelan.}
\label{PHYS:fig:rd}
\end{figure}

\subsection*{Lead shot: what disorder does}
Hales' peas presumably went into the pot anyhow. In 1939 the botanist J.~W.~Marvin repeated the experiment with
lead shot, compressed in a cylinder \PHYStag{experiment}. Shot that had been carefully stacked in a
close-packed arrangement became rhombic dodecahedra. Shot that had simply been poured in became
irregular polyhedra with on average $14.17$ faces, not twelve. So disorder costs faces: a randomly packed ball has
more Voronoi neighbours than a crystalline one, even though it has fewer \emph{touching} neighbours.

Bernal turned this into a programme around 1960. To model the structure of liquids he poured steel
ball bearings into rubber balloons, kneaded them, ran black paint into the pile and let it set, and
then took the pile apart ball by ball, recording which pairs had been in contact by the paint marks
(Bernal, \emph{Nature} 183, 1959; Bernal and Mason, \emph{Nature} 188, 1960). Bernal and Mason counted
about $6.4$ contacts per ball \PHYStag{experiment}, close to the isostatic value $2d=6$ of
Corollary~\ref{SM:cor:maxwell}. Bernal also classified the holes of a random packing: besides the tetrahedra and octahedra of FCC there
are three larger ``canonical holes'' bounded by equilateral triangles (\emph{Proc.\ R.\ Soc.\ A} 280, 1964). The density of these packings, about $0.64$, is the
random close packing value of Interlude~E (Section~\ref{SM:sec:random}).

\subsection*{Soap: the cell that minimizes its area}
Squeeze the peas harder, or replace them by bubbles, and the cells choose their own shapes: a dry
foam minimizes its total film area subject to fixed cell volumes. Plateau (\emph{Statique
exp\'erimentale et th\'eorique des liquides}, 1873) found the rules empirically: films meet in threes, at
$120^\circ$, along curves (the Plateau borders), and the borders meet in fours at points, at the
tetrahedral angle $\arccos(-\frac13)\approx109.47^\circ$. This is the angle between two lines from the
centre of a regular tetrahedron to its vertices, the Tammes angle for $N=4$ points of
Lecture~11 (Section~\ref{L11:sec:examples}). Jean Taylor proved a century later that these are the only
singularities a soap film can have (gallery, Theorem~\ref{PHYS:thm:taylor}).

Which cell has the least area? Kelvin (1887) proposed a slightly curved truncated octahedron, the
Voronoi cell of the body-centred cubic lattice, with $8$ hexagonal and $6$ square faces. For a century
it was the best known. In 1994 Weaire and Phelan found a better foam by computer \PHYStag{simulation}:
a periodic structure with two kinds of cells of equal volume, pentagonal dodecahedra and $14$-hedra
with $12$ pentagons and $2$ hexagons, in the ratio $1:3$. Its area is about $0.3\%$ smaller
(Figure~\ref{PHYS:fig:rd}). The Weaire--Phelan foam is the Voronoi diagram, relaxed, of the A15
($\beta$-tungsten) structure, one of the \emph{Frank--Kasper} phases of metallurgy, whose Delaunay cells
are all nearly regular tetrahedra. So the best known foam is Frank's frustrated tetrahedral packing of
Interlude~E seen from the dual side. It took until 2012 to make it in the laboratory: monodisperse
bubbles tend to crystallize into Kelvin's structure next to flat walls, and Gabbrielli, Meagher, Weaire, Brakke
and Hutzler (\emph{Phil.\ Mag.\ Lett.} 92, 2012) obtained Weaire--Phelan foam in a container whose walls
were shaped to fit it \PHYStag{experiment}. Whether Weaire--Phelan is optimal is open, and so is
Kelvin's problem itself.

Matzke (\emph{Amer.\ J.\ Botany} 33, 1946) blew a foam of about a thousand equal bubbles one at a time and
examined the $600$ in its interior. He found on average $13.70$ faces per bubble and not a single Kelvin cell
\PHYStag{experiment}. Euler's formula gives a feeling for where such numbers come from.

\begin{proposition}[Coxeter's statistical honeycomb]\label{PHYS:prop:coxeter}
Let $P$ be a convex polyhedron with exactly three faces at each vertex, and let $\bar n$ be the average
number of sides of its faces. Then the number of faces is
\[
  F=\frac{12}{6-\bar n}.
\]
If every face of a foam had $\bar n=2\pi/\arccos\frac13\approx5.104$ sides, the number of regular
tetrahedra that fit around an edge, every cell would have $F\approx13.40$ faces.
\end{proposition}

\begin{proof}
With $V$ vertices and $E$ edges, $3V=2E$ (three edges at each vertex, two vertices on each edge) and $\bar
nF=2E$. Euler's formula $V-E+F=2$ becomes $F-\frac E3=2$, that is $F(1-\frac{\bar n}6)=2$. For the second
statement: by Plateau's rules every vertex of a dry foam is simple, and each face of a cell is dual to
an edge of the Delaunay complex of the cell centres. The number of sides of the face is the number of
Delaunay cells around that edge. If those cells were regular tetrahedra, the dihedral angle
$\arccos\frac13$ would have to go $\bar n$ times around the edge.
\end{proof}

Coxeter (\emph{Illinois J.\ Math.} 2, 1958) called such idealized averages ``statistical honeycombs''. The
formula is a dictionary: $\bar n=5$ gives the dodecahedron ($F=12$), Kelvin's cell has $\bar
n=\frac{8\cdot6+6\cdot4}{14}=\frac{36}7$ ($F=14$), Weaire--Phelan has $F=13.5$, Matzke's bubbles $13.70$,
and Marvin's shot $14.17$. (Voronoi cells of completely random points have on average
$2+\frac{48\pi^2}{35}\approx15.54$ faces, a theorem of Meijering, \emph{Philips Res.\ Rep.} 8, 1953.) Five tetrahedra around every edge would give
$\bar n=5$ exactly, and then all the cells would be the regular dodecahedra of
Proposition~\ref{PHYS:prop:dodec}. This is impossible in flat space, but it is exactly what happens in the
curved space of Section~\ref{SM:sec:curved}. The Voronoi cells of the $120$ vertices of the $600$-cell in
$\Sph^3$ are the $120$ regular dodecahedra of the $120$-cell, spherical dodecahedra whose dihedral angle is exactly
$120^\circ$, three around each edge. Stephen Hales' ``regular dodecahedrons'' are
what the peas would have become in a pot shaped like the $3$-sphere.

In the plane everything is settled: the hexagonal honeycomb is the least-perimeter way to divide the
plane into equal areas (Hales, gallery, Theorem~\ref{PHYS:thm:honeycomb}). Bubble rafts on water were
used by Bragg and Nye (\emph{Proc.\ R.\ Soc.\ A} 190, 1947) as models of two-dimensional crystals, complete with
dislocations and grain boundaries.

%======================================================================
\section{Colloids: atoms you can watch}
\label{PHYS:sec:colloids}

\subsection*{The hard-sphere gas, made visible}
Interlude~E (Section~\ref{SM:sec:hard}) treated hard spheres as a system with no energy at all, in which
everything is entropy. Atoms are too soft for this, but colloids, particles between about
$100$ nanometres and a few micrometres across, suspended in a liquid, come remarkably close. The standard
hard-sphere colloid is a sphere of poly(methyl methacrylate) carrying a thin brush of polymer that
stops the spheres from sticking, suspended in a liquid of nearly the same refractive index, which
removes most of the van der Waals attraction and lets light through.

Colloids are big enough to see in an optical microscope, and slow. By the Stokes--Einstein relation a sphere of radius $a$ in a liquid of viscosity $\eta$ diffuses
with $D=k_BT/(6\pi\eta a)$. For $a=1\,\mu$m in water at room temperature,
$D\approx2\times10^{-13}\,\mathrm m^2/\mathrm s$, and the time to diffuse its own radius is $a^2/D\approx5$
seconds. An atom in a liquid does the same in picoseconds. So colloids are atoms that can be filmed, and since
the only energy scale is $k_BT$, the phase behaviour depends only on the packing fraction $\varphi$.

Pusey and van Megen (\emph{Nature} 320, 1986) prepared hard-sphere suspensions at a range of
concentrations and watched them for days \PHYStag{experiment}. At low $\varphi$ the samples stayed fluid;
in a window of concentrations they separated into a fluid and a crystal, with the crystal fraction
growing across the window; above it they crystallized throughout; and above about $\varphi\approx0.58$ they
did not crystallize at all, but stayed amorphous, a colloidal \emph{glass} (Pusey and van Megen,
\emph{Phys.\ Rev.\ Lett.} 59, 1987). This is the sequence predicted by the computer for hard spheres, with coexistence between
$0.494$ and $0.545$ (Interlude~E) \PHYStag{simulation}. (The volume fractions were in fact
calibrated by identifying the onset of freezing with $0.494$, so the experiment tests the shape of the
phase diagram rather than that number; on this scale melting came out near $0.535$. Polydispersity, a few
percent in real colloids, also shifts the phase boundaries, and in simulations a size spread above about
$7\%$ suppresses crystallization altogether.) Whether hard spheres crystallize at all is still not a theorem in any dimension
$d\ge2$, as Interlude~E explained. In microgravity, on the Space Shuttle, dense samples that had stayed glassy on
earth for more than a year crystallized completely within two weeks (Zhu et al., \emph{Nature} 387,
1997), so part of the colloidal ``glass'' is gravity's doing.

\subsection*{Why order wins: a free-volume calculation}
That a system with no energy should \emph{order} when compressed seems paradoxical, since ordering
sounds like a loss of entropy. The resolution is that the relevant entropy is the room each ball has to
rattle. The following idealized model, the \emph{cell model}, makes this quantitative.

\begin{proposition}[Cell model \PHYStag{model}]\label{PHYS:prop:cell}
Suppose each of $N$ hard balls of diameter $1$ in $\R^d$ moves independently in a cage formed by its
neighbours, frozen at the sites of a lattice packing scaled from its close-packed density $\rho_{\mathrm{cp}}$
to density $\rho=N/V$, and suppose the volume of each cage is $c(a-1)^d$, where
$a=(\rho_{\mathrm{cp}}/\rho)^{1/d}$ is the nearest-neighbour distance. Then the entropy per ball is
$s=d\ln(a-1)+\ln c$ and the equation of state is
\[
  \frac{\beta P}{\rho}=\frac1{1-(\rho/\rho_{\mathrm{cp}})^{1/d}} .
\]
\end{proposition}

\begin{proof}
The partition function is $Z=(c(a-1)^d)^N$, so $S=\ln Z=Nd\ln(a-1)+N\ln c$. The volume is
$V=Na^d/\rho_{\mathrm{cp}}$, so $dV/da=Nda^{d-1}/\rho_{\mathrm{cp}}$. With $\beta P=\partial S/\partial V$ as in
Section~\ref{SM:sec:hard},
\[
  \beta P=\frac{\partial S/\partial a}{\partial V/\partial a}=\frac{Nd/(a-1)}{Nda^{d-1}/\rho_{\mathrm{cp}}}
  =\frac{\rho_{\mathrm{cp}}}{a^{d-1}(a-1)},
\]
and dividing by $\rho=\rho_{\mathrm{cp}}/a^d$ gives $\beta P/\rho=a/(a-1)=1/(1-a^{-1})$.
\end{proof}

In dimension one this is \emph{exactly} the Tonks gas of Proposition~\ref{SM:prop:tonks}: with
$\rho_{\mathrm{cp}}=1$ it reads $\beta P/\rho=1/(1-\rho)$. In one dimension the cell picture is not a
model at all, since the gaps between consecutive rods really are free lengths. In three dimensions the
cell model is crude, but not bad: for FCC at the melting density $\varphi=0.545$ it gives
$\beta P/\rho\approx10.3$, against about $11$ in simulations.

Now compare a crystal with an amorphous packing at the same density. The amorphous packing jams at
$\varphi_J\approx0.64$ instead of $\pi/\sqrt{18}\approx0.7405$, so in the same model its cages have linear size
$(\varphi_J/\varphi)^{1/3}-1$ instead of $(0.7405/\varphi)^{1/3}-1$. At $\varphi=0.55$ the ratio of the two cage sizes
is $2.01$: each crystalline ball has twice the room to move in every direction. The vibrational entropy
gained by crystallizing is therefore about $3\ln2\approx2.1\,k_B$ per ball. Against this the disordered
state has its \emph{configurational} entropy, the logarithm of the number of distinct disordered
arrangements, which wins at low density; roughly speaking, freezing is where the two balance. Order is the high-entropy state of a crowded system.

\subsection*{FCC or HCP?}
Every stacking of hexagonal layers (each layer in the hollows of the one below, the Barlow packings)
has density $\pi/\sqrt{18}$, so Hales' theorem does not choose between them. Nor does counting: with $L$
layers there are $3\cdot2^{L-1}$ stacking sequences, but each layer has $m$ balls, so the stacking entropy
per ball is $(\ln3+(L-1)\ln2)/(Lm)\to0$ as $m\to\infty$. The choice is made by the vibrational entropy,
which sees the difference between FCC and HCP only at the third shell of neighbours. Bolhuis, Frenkel,
Mau and Huse (\emph{Nature} 388, 1997) computed that FCC is favoured, by about $10^{-3}k_BT$ per particle
near melting \PHYStag{simulation}; Mau and Huse (\emph{Phys.\ Rev.~E} 59, 1999) found $0.00115(4)\,k_B$ per sphere at close
packing, and that FCC has the highest entropy of all stackings. The difference is so small that colloidal
crystals grown in the laboratory often stack \emph{randomly} (Pusey, van Megen, Bartlett, Ackerson,
Rarity and Underwood, \emph{Phys.\ Rev.\ Lett.} 63, 1989) \PHYStag{experiment}.

\paragraph{Looking inside.} Confocal microscopy locates every particle of a three-dimensional sample, in
glasses and in growing crystal nuclei (Gasser et al., \emph{Science} 292, 2001). In two dimensions, colloidal hard spheres confined to a plane
show the melting scenario that Bernard and Krauth found in simulations (Section~\ref{SM:sec:hard}): a
first-order liquid--hexatic transition followed by a continuous hexatic--solid transition (Thorneywork,
Abbott, Aarts and Dullens, \emph{Phys.\ Rev.\ Lett.} 118, 2017) \PHYStag{experiment}.

%======================================================================
\section{Opals and photonic crystals}
\label{PHYS:sec:opal}

\subsection*{A gemstone made of lattice planes}
Precious opal flashes pure colours that change as the stone is turned. In 1964 J.~V.~Sanders and his
colleagues at CSIRO looked at opal in the electron microscope and found a regular array of silica
spheres, a few hundred nanometres across, stacked in hexagonal layers (Jones, Sanders and Segnit,
\emph{Nature} 204, 1964; Sanders, \emph{Nature} 204, 1964). The colours are Bragg diffraction of light by the
layers \PHYStag{experiment}. To see which layers and which colours, we need the planes of a lattice, and
those are counted by the dual lattice $\Lat^*=\{y:\ip xy\in\Z\ \forall x\in\Lat\}$ of
Interlude~B (Theorem~\ref{NT:thm:poisson}).

\begin{proposition}[Lattice planes]\label{PHYS:prop:planes}
Let $\Lat\subset\R^3$ be a lattice and $w\in\Lat^*$ a primitive vector (not a multiple $kw'$ of
some $w'\in\Lat^*$ with $k\ge2$). Then $\Lat$ is the disjoint union of its intersections with the parallel planes
$\{x:\ip xw=k\}$, $k\in\Z$, each of which is nonempty, and these planes are $1/\abs w$ apart. Conversely, every equally spaced family of parallel planes that covers $\Lat$, each plane containing a
lattice point, arises this way, so the largest spacing between lattice planes is
$1/\min\{\abs w:w\in\Lat^*\setminus0\}$.
\end{proposition}

\begin{proof}
The values $\ip xw$, $x\in\Lat$, form a subgroup $g\Z$ of $\Z$ with $g\ge1$. Then $\ip x{w/g}\in\Z$ for all
$x\in\Lat$, so $w/g\in\Lat^*$, and primitivity forces $g=1$. So every integer $k$ is a value, every lattice
point lies on one of the planes, and the planes $\ip xw=k$ are $1/\abs w$ apart. Conversely, suppose $\Lat$ lies on the planes $\ip xu\in s\Z$, for a unit vector
$u$ and some $s>0$, and each of these planes contains a lattice point. Then $w=u/s$ has $\ip xw\in\Z$ for all
$x\in\Lat$, so $w\in\Lat^*$, and every integer is a value, so $w$ is not a multiple $kw'$ with $k\ge2$.
\end{proof}

\begin{example}[The colour of an opal]\label{PHYS:ex:opal}
For $\Lat=D_3$ the dual is $D_3^*=\Z^3\cup(\Z^3+\frac12(1,1,1))$, the body-centred cubic lattice, whose
shortest vectors are the eight vectors $\frac12(\pm1,\pm1,\pm1)$ of length $\frac{\sqrt3}2$. So the widest-spaced
planes of FCC are perpendicular to the body diagonals, at distance $2/\sqrt3$: these are the hexagonal
close-packed layers. Rescaling from minimal distance $\sqrt2$ to sphere diameter $D$, the layer spacing is
\[
  d=\frac{2}{\sqrt3}\cdot\frac{D}{\sqrt2}=D\sqrt{\tfrac23},
\]
which is also the height of a regular tetrahedron of edge $D$. The tetrahedron shows that the spacing
is the same for every Barlow stacking, so the colour does not care about stacking faults. Bragg's law for
light inside a medium of effective refractive index $n_{\mathrm{eff}}$ reflects the wavelength
\[
  \lambda=2n_{\mathrm{eff}}\,d\sin\theta,
\]
where $\theta$ is the glancing angle. At normal incidence, with $n_{\mathrm{eff}}\approx1.45$ for silica, $\lambda\approx2.37\,D$.
The visible range $400$--$700$\,nm corresponds to $D$ between about $170$ and $300$\,nm. Turning the stone
changes $\theta$, and the colour slides towards the blue.
\end{example}


\subsection*{Photonic band gaps}
In 1987 Yablonovitch (\emph{Phys.\ Rev.\ Lett.} 58, 2059), who wanted to switch off spontaneous emission, and
John (\emph{Phys.\ Rev.\ Lett.} 58, 2486), who wanted to localize light, independently proposed to do with
photons what a crystal does with electrons. In a medium whose refractive index is periodic with lattice
$\Lat$, Maxwell's equations separate into Bloch waves, and the frequencies form bands over the Brillouin
zone, the Voronoi cell of $\Lat^*$ (up to the factor $2\pi$). A \emph{complete photonic band gap} is an
interval of frequencies in which light cannot propagate in any direction.

The Brillouin zone of an opal is the Voronoi cell of BCC: the truncated octahedron, that is,
Kelvin's cell from Section~\ref{PHYS:sec:squash}. A gap opens on each face of the zone, at the frequency where
the Bragg condition for the corresponding planes holds, and for a complete gap the gaps on all faces must
overlap. With $D_3^*$ as above, the centres of the hexagonal faces (the points the physicists call $L$) are
at distance $\frac{\sqrt3}4\approx0.433$, the centres of the square faces ($X$) at $\frac12$, and the
vertices ($W$) at $\frac{\sqrt5}4\approx0.559$. The free-photon frequencies at $W$ and at $L$ differ by the factor
$\sqrt{5/3}\approx1.29$, which is hard to bridge. In fact Ho, Chan and Soukoulis (\emph{Phys.\ Rev.\ Lett.} 65,
1990) found that dielectric spheres on an FCC lattice have no complete gap in the lowest bands, because
of a band degeneracy at $W$, while spheres on the \emph{diamond} structure have one, already at an
index contrast of about $2$ \PHYStag{simulation}.

The diamond structure is $D_3^+=D_3\cup(D_3+\frac12(1,1,1))$: FCC together with one of its two classes of
tetrahedral holes (Interlude~H, Section~\ref{DYN:sec:fcc}). This is the $n=3$ case of the construction
$D_n^+$ of Lecture~6 (Section~\ref{L04b:sec:rootdense}), which is a lattice only for even $n$
(Exercise~\ref{NT:ex:D16}). As a packing of balls it is very loose: the nearest distance drops to
$\frac{\sqrt3}2$, there is one ball per unit volume, and the density is
$\frac43\pi(\frac{\sqrt3}4)^3=\frac{\pi\sqrt3}{16}\approx0.340$, less than half that of FCC. The best crystal
for light is among the worst for packing. Weevils found it first: the iridescent scales of the beetle
\emph{Lamprocyphus augustus} contain a diamond-structured photonic crystal of chitin (Galusha et al.,
\emph{Phys.\ Rev.~E} 77, 2008) \PHYStag{experiment}.

Opals can be made to work after all by inverting them. Let colloidal spheres crystallize into a synthetic
opal, fill the gaps with a material of high refractive index, and dissolve the spheres. The result, an
\emph{inverse opal} of air spheres, has a complete gap between its eighth and ninth bands once the index
contrast exceeds about $2.8$ (Busch and John, \emph{Phys.\ Rev.~E} 58, 1998) \PHYStag{simulation}.
Silicon, with index about $3.5$, qualifies, and silicon inverse opals with a complete gap near
$1.5\,\mu$m, the wavelength of optical fibres, were made by Blanco et al.\ (\emph{Nature} 405, 2000; a gap of
roughly $5\%$ of its centre frequency) and Vlasov, Bo, Sturm and Norris (\emph{Nature} 414, 2001) \PHYStag{experiment}.
For a theorem guaranteeing gaps in a model of high-contrast media, see the gallery
(Theorem~\ref{PHYS:thm:FK}).

%======================================================================
\section{From crystal to glass: random lasers and granular lasers}
\label{PHYS:sec:laser}

Take the same colloidal spheres and let them settle quickly, before they can crystallize, or add salt so
that they stick on contact. The result is a \emph{photonic glass} (Garc\'ia, Sapienza, Blanco and L\'opez,
\emph{Adv.\ Mater.} 19, 2007): a random packing of monodisperse spheres, the optical counterpart of Bernal's
ball bearings. Light no longer diffracts; it scatters many times and performs a random walk, with steps
of the order of the \emph{transport mean free path} $\ell$. Over distances much larger than $\ell$, the
energy density $W$ of the light diffuses, $\partial_tW=D\Delta W$ with $D=v\ell/3$ ($v$ the speed of light in
the medium), the continuum limit of the coin flips in the gas in a box of Lecture~7
(Section~\ref{L07:subsec:gas}).

Now soak the packing in a laser dye and pump it. Letokhov (\emph{Sov.\ Phys.\ JETP} 26, 1968) predicted
that a scattering medium with gain will lase without any mirrors, once it is large enough: the
scattering keeps the light inside long enough to be amplified. The threshold is a pretty piece of
spectral theory.

\begin{proposition}[Letokhov's critical radius \PHYStag{model}]\label{PHYS:prop:letokhov}
Let $W(r,t)$ solve $\partial_tW=D\Delta W+\frac v{\ell_g}W$ in a ball of radius $R$ in $\R^3$, with $W=0$ on the
boundary, where $D=v\ell/3$ and $\ell_g$ is the gain length (the distance over which the intensity grows
by a factor $e$). Then every solution decays to $0$ if $R<R_{\mathrm{cr}}$ and generic solutions grow
exponentially if $R>R_{\mathrm{cr}}$, where
\[
  R_{\mathrm{cr}}=\pi\sqrt{\frac{\ell\,\ell_g}3}.
\]
\end{proposition}

\begin{proof}
Expand in the Dirichlet eigenfunctions $-\Delta\phi_k=\lambda_k\phi_k$ of the ball:
\[
  W=\sum_kc_k\,e^{(v/\ell_g-D\lambda_k)t}\phi_k .
\]
Everything decays if and only if $v/\ell_g<D\lambda_1$. The first
eigenvalue is simple, with a positive eigenfunction, which is therefore invariant under rotations. For
radial $\phi$, the function $u=r\phi$ satisfies $u''=-\lambda u$ with $u(0)=u(R)=0$, so $u=\sin(\pi r/R)$ and
$\lambda_1=\pi^2/R^2$. The threshold $v/\ell_g=\frac{v\ell}3\cdot\frac{\pi^2}{R^2}$ gives the formula.
\end{proof}

The threshold volume scales like $(\ell\ell_g)^{3/2}$: stronger scattering means a smaller random laser.
And $\ell$ depends on the packing. For independent scatterers, $\ell\approx1/(n\sigma_{\mathrm{tr}})$, with
number density $n=6\varphi/(\pi D^3)$ and transport cross-section $\sigma_{\mathrm{tr}}$ of order $D^2$ for grains
much larger than the wavelength, so $\ell$ is of order $D/\varphi$ and $R_{\mathrm{cr}}\propto\varphi^{-1/2}$: compacting a
granular gain medium lowers its threshold. (Near jamming, positional correlations between the grains
reduce scattering, and the independent-scatterer estimate has to be corrected by the structure factor of
the packing.)

Lawandy, Balachandran, Gomes and Sauvain (\emph{Nature} 368, 1994) saw laser action in a dye solution
with titania particles \PHYStag{experiment}, and Cao and coworkers (\emph{Phys.\ Rev.\ Lett.} 82, 1999) saw
sharp lasing modes in zinc oxide powder, where the feedback comes from closed loops of multiply scattered
light. In a photonic glass the Mie resonances of the equal spheres can be used to select the lasing
wavelength (Gottardo et al., \emph{Nature Photonics} 2, 2008). Wiersma and Cavalieri (\emph{Nature} 414, 2001) built a random laser out
of compressed glass powder whose pores were filled with a dye-doped liquid crystal. The alignment of the
liquid crystal, and with it the strength of the scattering, depends on temperature, so the laser can be
tuned by heating it, and it switches off at the phase transition of the liquid crystal near $42^\circ$C
\PHYStag{experiment}. This was already a granular laser, tuned through the scattering of its grains.

\paragraph{Granular lasers.} Then the packing itself was made to move. Folli, Puglisi, Leuzzi and Conti
(\emph{Shaken granular lasers}, \emph{Phys.\ Rev.\ Lett.} 108, 2012) immersed about $1500$ millimetre-sized
spheres in a dye solution, put the container on a vibrating plate and pumped it with a pulsed laser
\PHYStag{experiment}. The random laser emission depends on whether and how the grains move: at strong
shaking a second emission peak appeared, attributed to transient cavities formed by a few grains. In
a sequel (Folli, Ghofraniha, Puglisi, Leuzzi and Conti, \emph{Sci.\ Rep.} 3, 2013) with $1$\,mm glass
beads in a rhodamine solution, the laser light distinguished a solid-like regime, in which the grains keep
their neighbours and the spectra are steady, from a fluidized regime in which the density fluctuations
switch the system above and below threshold, so that the spectra alternate between fluorescence and lasing.
The authors call the result a ``granular'' laser with a tunable emission spectrum: the configuration of
the beads, and whether the packing is solid-like or fluidized, is part of the optical cavity.

%======================================================================
\section{Sticky spheres}
\label{PHYS:sec:sticky}

\subsection*{From hard to sticky}
Add small polymer coils to a colloidal suspension and the spheres attract each other: when two spheres
come closer than the size of a coil, the coils are excluded from the gap between them, and the osmotic
pressure of the coils pushes the spheres together (Asakura and Oosawa, \emph{J.\ Chem.\ Phys.} 22, 1954).
The range of this \emph{depletion} attraction is the size of the coil, which can be a few percent of the
sphere diameter. In the limit of infinitely deep and narrow attraction, Baxter's \emph{sticky hard spheres}
(\emph{J.\ Chem.\ Phys.} 49, 1968), the energy of a configuration is minus the number of contacts: it is
read off the contact graph of Lectures~10--11 (Definition~\ref{L11:def:contact}), now for balls in $\R^3$.

So the ground states of $n$ sticky spheres are the packings of $n$ unit balls with the most contacts.
This \emph{contact number problem} is a cousin of the kissing problem, and it is hard. In the plane it is
solved: $n$ unit discs have at most $\lfloor3n-\sqrt{12n-3}\rfloor$ contacts (Harborth, gallery), attained by
hexagonal patches. In $\R^3$ only bounds are known (Bezdek and Reid, gallery). For small $n$ the answers
come from enumeration.

\subsection*{Rigidity and counting}
A cluster that is held together only by its contacts should be rigid. The rigidity matrix $A$ of
Interlude~E (Section~\ref{SM:sec:random}; see also Interlude~F, Section~\ref{COMP:sec:matrix}) has a row
$(\dots,-u_{ij},\dots,u_{ij},\dots)$ for each contact. Now there is no box: a free cluster can be
translated and rotated.

\begin{proposition}[Maxwell count for clusters \PHYStag{theorem}]\label{PHYS:prop:cluster}
Let $n\ge3$ unit balls in $\R^3$ have centres not all on a line. If the contact framework is
infinitesimally rigid (every infinitesimal motion preserving all contact distances to first order is a
rigid motion), then it has at least $3n-6$ contacts, and if moreover $n\ge4$, every ball has at least
three contacts.
\end{proposition}

\begin{proof}
The infinitesimal rigid motions $v_i=b+\omega\times x_i$ form a $6$-dimensional subspace of $\ker A$ when the
centres are not collinear (a motion with $b+\omega\times x_i=0$ for three non-collinear $x_i$ is zero).
Infinitesimal rigidity says this is all of $\ker A$, so $\operatorname{rank}A=3n-6$, and $A$ needs at least
that many rows. For the second claim, suppose ball $i$ has at most two contacts. Pick $w\ne0$ orthogonal
to its contact directions, and put $v_i=w$ and $v_l=0$ for $l\ne i$; this lies in $\ker A$. A rigid motion
vanishing at three non-collinear centres is zero, so if this $v$ is a rigid motion, the other $n-1\ge3$
centres lie on a line $L$, and $x_i\notin L$. Then the other balls form a straight chain along $L$, and an
end ball of the chain touches only its neighbour in the chain and possibly ball $i$. The centres other
than this end ball include $x_i$ and two points of $L$, so they are not collinear, and the same
construction at the end ball gives an element of $\ker A$ that is not a rigid motion.
\end{proof}

\noindent For clusters that are rigid but not infinitesimally rigid the second claim can fail: a ball
squeezed exactly between two balls at distance $2$ of an otherwise rigid cluster is pinned by its two
contacts.

Clusters with exactly $3n-6$ contacts that are rigid are called \emph{isostatic} (or minimally rigid),
those with more \emph{hyperstatic}, and rigid ones with fewer \emph{hypostatic}. A rigid cluster that is
not infinitesimally rigid is \emph{singular}. Arkus, Manoharan and Brenner (\emph{Phys.\ Rev.\ Lett.} 103, 2009;
\emph{SIAM J.\ Discrete Math.} 25, 2011) enumerated the candidate contact graphs (adjacency matrices with
at least $3n-6$ edges and minimum degree $3$, the two conditions suggested by Proposition~\ref{PHYS:prop:cluster}), and
solved for the distances with geometric rules. Hoy, Harwayne-Gidansky and O'Hern (\emph{Phys.\ Rev.~E} 85,
2012) went to $n=11$, and Holmes-Cerfon (\emph{SIAM Review} 58, 2016) to $n=14$ by a different method,
walking from one rigid cluster to the next along the one-parameter families obtained by breaking a
contact. Her counts:
\begin{center}
\small
\renewcommand{\arraystretch}{1.12}
\begin{tabular}{lccccccccc}
\toprule
$n$ & $6$ & $7$ & $8$ & $9$ & $10$ & $11$ & $12$ & $13$ & $14$\\
\midrule
rigid clusters & $2$ & $5$ & $13$ & $52$ & $263$ & $1659$ & $11980$ & $98529$ & $895478$\\
of which hypostatic & $0$ & $0$ & $0$ & $0$ & $1$ & $20$ & $159$ & $1308$ & $11245$\\
of which hyperstatic & $0$ & $0$ & $0$ & $0$ & $3$ & $21$ & $183$ & $1411$ & $11241$\\
most contacts & $12$ & $15$ & $18$ & $21$ & $25$ & $29$ & $33$ & $36$ & $40$\\
\bottomrule
\end{tabular}
\end{center}
These are numerical results \PHYStag{simulation}, believed complete apart possibly from a few highly
singular clusters, but not proved complete. The first singular cluster appears at $n=9$, and the first
clusters with more than $3n-6$ contacts at $n=10$: there the ground states stop being isostatic and become
fragments of the hexagonal close packing. Even the statement that $3n-6$ is the maximal number of
contacts for $6\le n\le9$ is, strictly speaking, a conjecture, conditional on the completeness of these
lists (Bezdek and Khan, arXiv:1601.00145).

\begin{figure}[ht]
\centering
\begin{tikzpicture}[scale=1.75]
  \begin{scope}
    \coordinate (O0) at (0.637,0.126); \coordinate (O1) at (-0.637,-0.126);
    \coordinate (O2) at (-0.308,0.260); \coordinate (O3) at (0.308,-0.260);
    \coordinate (O4) at (0.000,0.645); \coordinate (O5) at (0.000,-0.645);
    \foreach \i in {2,5,0,1,4,3} \filldraw[PHYSball] (O\i) circle (0.5);
    \foreach \a/\b in {0/2,0/3,0/4,0/5,1/2,1/3,1/4,1/5,2/4,2/5,3/4,3/5} \draw[thick] (O\a)--(O\b);
    \foreach \i in {0,...,5} \fill (O\i) circle (0.025);
    \node[font=\small,align=center] at (0,-1.3) {octahedron\\ $12$ contacts, $\sigma=24$};
  \end{scope}
  \begin{scope}[xshift=3.1cm]
    \coordinate (P0) at (0.661,-0.123); \coordinate (P1) at (-0.331,-0.254);
    \coordinate (P2) at (0.052,0.670); \coordinate (P3) at (0.127,0.098);
    \coordinate (P4) at (0.253,-0.856); \coordinate (P5) at (-0.762,0.466);
    \foreach \i in {3,5,4,2,0,1} \filldraw[PHYSball] (P\i) circle (0.5);
    \foreach \a/\b in {0/1,0/2,0/3,0/4,1/2,1/3,1/4,1/5,2/3,2/5,3/4,3/5} \draw[thick] (P\a)--(P\b);
    \foreach \i in {0,...,5} \fill (P\i) circle (0.025);
    \node[font=\small,align=center] at (0,-1.3) {polytetrahedron\\ $12$ contacts, $\sigma=2$};
  \end{scope}
\end{tikzpicture}
\caption{The two rigid clusters of six sticky spheres, with their contact graphs. The octahedron is a
piece of FCC (six balls around an octahedral hole). The polytetrahedron is three regular tetrahedra
glued face to face; the view is along the normal of one face of the middle tetrahedron. Both have
$3n-6=12$ contacts, so they have the same energy, but in the sticky-limit model the polytetrahedron is $24$
times as likely (Example~\ref{PHYS:ex:six}).}
\label{PHYS:fig:six}
\end{figure}

\subsection*{Six spheres: entropy chooses the less symmetric cluster}
For $n=6$ there are two ground states, both with $12$ contacts (Figure~\ref{PHYS:fig:six}). They have the same
energy, so which one forms is decided by entropy, and the entropy has three parts. Holmes-Cerfon,
Gortler and Brenner (\emph{Proc.\ Natl.\ Acad.\ Sci.\ USA} 110, 2013) showed that in the sticky limit the
partition function of a rigid, infinitesimally rigid cluster with $m$ contacts is, up to factors common
to all clusters of the same size,
\begin{equation}\label{PHYS:eq:sticky}
  z=\kappa^m\cdot\frac{\sqrt{\det I}}{\sigma\,\sqrt{\det(AA^T)}}\qquad(m=3n-6),
\end{equation}
where $\kappa$ is a ``stickiness'' parameter that depends on the depth and range of the attraction, $I$ is the
moment of inertia tensor of the centres (unit masses), $\sigma$ is the order of the rotation group of the
cluster, and $A$ is the rigidity matrix. The three factors are rotational entropy (a cluster with more
symmetry has fewer distinguishable orientations, and a more spread-out cluster has more), and vibrational
entropy: $\det(AA^T)$ is the product of the nonzero eigenvalues of the stiffness matrix $A^TA$, and
softer clusters wiggle more.

\begin{example}[$n=6$ \PHYStag{model}]\label{PHYS:ex:six}
For the octahedron, with centres $\pm\frac1{\sqrt2}e_i$, the rotation group is that of the cube,
$\sigma=24$, and $I=2\cdot\mathrm{Id}$, so $\det I=8$. The $12\times12$ matrix $AA^T$ has entries
$\ip{u_e}{u_{e'}}$ (with the signs of the shared endpoints) and $\det(AA^T)=864$. For the polytetrahedron,
with centres
\[
  (0,0,0),\ (1,0,0),\ (\tfrac12,\tfrac{\sqrt3}2,0),\ (\tfrac12,\tfrac{\sqrt3}6,\tfrac{\sqrt6}3),\
  (\tfrac12,-\tfrac{7\sqrt3}{18},\tfrac{2\sqrt6}9),\ (\tfrac43,\tfrac{4\sqrt3}9,\tfrac{2\sqrt6}9),
\]
only a half-turn survives, $\sigma=2$, and an exact computation gives $\det I=\frac{261569}{26244}\approx9.967$ and
$\det(AA^T)=\frac{261569}{972}\approx269.1$. Hence
\[
  \frac{z_{\mathrm{poly}}}{z_{\mathrm{oct}}}
  =\frac{24}2\cdot\sqrt{\frac{\det I_{\mathrm{poly}}}{\det I_{\mathrm{oct}}}\cdot\frac{\det(AA^T)_{\mathrm{oct}}}{\det(AA^T)_{\mathrm{poly}}}}
  =12\cdot\sqrt{\frac{972}{26244}\cdot\frac{864}8}=12\cdot\sqrt{\frac{108}{27}}=24 .
\]
Symmetry contributes the factor $12$, and the moments of inertia and the vibrations together exactly $2$
(about $1.12$ and $1.79$). So the model predicts that the octahedron forms with probability $\frac1{25}=4\%$.
\end{example}

Meng, Arkus, Brenner and Manoharan (\emph{Science} 327, 2010) assembled clusters of colloidal spheres with a
short-range attraction in microscopic wells and counted them \PHYStag{experiment}: the polytetrahedron
was observed more than $20$ times as often as the octahedron, so the octahedron made up only about
$4$--$5\%$ of the six-particle clusters, in good agreement with the $4\%$ (a ratio of $24$) of
Example~\ref{PHYS:ex:six}. For clusters of up to ten particles, highly symmetric clusters were strongly
suppressed by rotational entropy. So the cluster that is a piece of
the crystal loses, and the polytetrahedral one wins, as in Frank's picture of liquids (Section~\ref{SM:sec:frustration}),
though here the cause is rotational entropy rather than energy. In sticky systems the crystal comes back at
$n\ge10$, through the hyperstatic clusters, which gain a factor $\kappa$ for every extra contact. Kallus
and Holmes-Cerfon (\emph{Phys.\ Rev.~E} 95, 2017) extended \eqref{PHYS:eq:sticky} to singular clusters,
whose zero-frequency modes make the harmonic formula diverge, computed the partition functions of all
known rigid clusters with $n\le19$, and found that the hyperstatic ones, which look like fragments of a
close-packed lattice, dominate the landscape: this is how crystalline order emerges from sticky spheres,
already at $n=10$.

\paragraph{The configuration space is stratified.} The space of $n$ sticky spheres modulo rigid motions is
a union of strata, one for each contact graph $G$, of dimension $3n-6-\abs{E(G)}$ when $G$ is regular. Rigid
clusters are the $0$-dimensional strata, and breaking one contact of an isostatic cluster gives a
$1$-dimensional stratum, a path to other clusters. This is the configuration space of Lectures~7--10 (and
Interlude~D, Section~\ref{TOP:sec:conf}) cut up by contact graphs instead of by levels of the min-distance
function; see Holmes-Cerfon's review (\emph{Annu.\ Rev.\ Condens.\ Matter Phys.} 8, 2017).

%======================================================================
\section{Minimal rigidity and random jamming}
\label{PHYS:sec:jam}

\subsection*{Near the jamming point}
Interlude~E showed that a first-order jammed packing of frictionless spheres has at least $2d$ contacts
per ball on average, and that jammed packings found in simulations have exactly this many: they are
isostatic at the jamming threshold $\varphi_J\approx0.64$ (point~J of O'Hern, Silbert, Liu and Nagel). What
happens just above $\varphi_J$? Make the spheres slightly soft, with a repulsion proportional to the overlap
(harmonic spheres), and compress. Then (Figure~\ref{PHYS:fig:jam}) \PHYStag{simulation}
\[
  z-2d\ \propto\ (\varphi-\varphi_J)^{1/2},\qquad P\ \propto\ \varphi-\varphi_J ,
\]
found for a bubble model of foams by Durian (\emph{Phys.\ Rev.\ Lett.} 75, 1995) and in detail by O'Hern et
al.\ (\emph{Phys.\ Rev.~E} 68, 2003). In the laboratory, Majmudar, Sperl, Luding and Behringer
(\emph{Phys.\ Rev.\ Lett.} 98, 2007) measured the contact numbers of photoelastic discs, which light up under
stress, and found a jump at $\varphi_J$ followed by growth with exponent about $0.5$, and pressure growing with
exponent about $1.1$ \PHYStag{experiment}. The square root, rather than the linear growth one might expect, says
that new contacts are cheap: the isostatic packing is marginal.

There is a pretty counting argument behind the marginality, due to Wyart, Nagel and Witten
(\emph{Europhys.\ Lett.} 72, 2005) \PHYStag{model}. Cut a cube of side $L$ (in units of the ball diameter)
out of a packing with $z=2d+\Delta z$ and free its boundary. The contacts crossing the boundary, about
$L^{d-1}$ of them, are lost, and in an isostatic packing each lost contact releases one floppy mode. The
interior has about $\frac12\Delta z\,L^d$ redundant contacts to block them. So floppy modes exist in boxes
of size up to
\[
  \ell^*\sim\frac1{\Delta z}\sim(\varphi-\varphi_J)^{-1/2},
\]
a length that diverges at the jamming point. Extended soft modes of this size explain the excess of
low-frequency vibrations (the ``boson peak'') of jammed packings.

\begin{figure}[ht]
\centering
\begin{tikzpicture}[x=55cm,y=0.48cm]
  \draw[->] (0.595,0)--(0.705,0) node[below right,font=\small]{$\varphi$};
  \draw[->] (0.6,0)--(0.6,10.5) node[left,font=\small]{$z$};
  \foreach \v in {4,6} {\draw (0.6,\v)--(0.598,\v) node[left,font=\scriptsize]{$\v$};}
  \draw (0.64,0)--(0.64,-0.25) node[below,font=\scriptsize]{$\varphi_J$};
  \draw[very thick,addcol] (0.6,0)--(0.64,0);
  \draw[very thick,addcol] plot[domain=0.64:0.70,samples=60] (\x,{6+14*sqrt(\x-0.64)});
  \fill[addcol] (0.64,6) circle (1.6pt);
  \draw[addcol,dashed] (0.64,0)--(0.64,6);
  \draw[thick,flagcol,dashed] (0.64,0)--(0.70,6.6) node[right,font=\scriptsize,flagcol]{$P\propto\varphi-\varphi_J$};
  \node[font=\scriptsize,addcol,anchor=west] at (0.70,9.45) {$z-6\propto(\varphi-\varphi_J)^{1/2}$};
  \draw[gray,densely dotted] (0.6,4)--(0.635,4);
  \node[font=\scriptsize,gray,anchor=west] at (0.6015,4.55) {friction: $z\ge4$};
  \node[font=\scriptsize,anchor=west] at (0.603,8.8) {$d=3$, schematic};
\end{tikzpicture}
\caption{Schematic behaviour of soft frictionless spheres near the jamming point in $d=3$: below $\varphi_J$
there are no contacts (after rattlers are removed); at $\varphi_J$ the mean contact number jumps to the
isostatic value $2d$; above it the excess grows like a square root, while the pressure of harmonic spheres
grows linearly. Frictional packings can be stable with fewer contacts.}
\label{PHYS:fig:jam}
\end{figure}

\subsection*{Friction, and the loosest random packing}
Real grains have friction, and friction lowers the number of contacts needed. The counting is the same as
in Corollary~\ref{SM:cor:maxwell}, but with forces instead of motions \PHYStag{model}. A frictional contact
carries a force \emph{vector} with $d$ components (normal and tangential), so a packing with mean
coordination $z$ has $\frac{zN}2\cdot d$ unknown force components. Each ball must satisfy force balance
($d$ equations) and torque balance ($\binom d2$ equations), $N\frac{d(d+1)}2$ in all. For generic packings
to carry a load we need at least as many unknowns as equations, $z\ge d+1$. Without friction a contact
carries only a normal force and there is no torque on a sphere, so $\frac{zN}2\ge dN$, which is $z\ge2d$
again. On the other hand each contact is one equation on the positions, whether or not there is friction,
so for generic sizes $\frac{zN}2\le dN$, as in Interlude~E. Hence
\[
  d+1\ \le\ z\ \le\ 2d,
\]
with $z=d+1$ approached for very rough grains and $z=2d$ for frictionless ones: $4\le z\le6$ in three
dimensions.

Fewer contacts allow looser packings. Onoda and Liniger (\emph{Phys.\ Rev.\ Lett.} 64, 1990) let glass
spheres settle slowly through liquids of nearly the same density and extrapolated to zero effective
gravity, finding a \emph{random loose packing} fraction $0.555\pm0.005$ \PHYStag{experiment}. Song, Wang and Makse
(\emph{Nature} 453, 2008) proposed a mean-field theory of jammed packings in which the packing fraction is
tied to the mean geometric coordination by $\varphi=z/(z+2\sqrt3)$ \PHYStag{model}. It gives
$\varphi=6/(6+2\sqrt3)\approx0.634$ for $z=6$ (frictionless, random close packing) and
$4/(4+2\sqrt3)\approx0.536$ for $z=4$ (infinite friction, random loose packing), and a phase diagram of
jammed matter between these two lines. With adhesion the constraint count drops further, and simulations of
small sticky grains give much looser stable packings; Liu, Li, Baule and Makse (\emph{Soft Matter} 11, 2015)
conjecture a limiting loosest packing at $z=2$ with $\varphi=1/2^3$, which their simulations approach but
do not reach \PHYStag{model}.

\subsection*{How packings are made}
Slow sedimentation gives loose packings; vibration, shearing and tapping compact them. Knight, Fandrich, Lau, Jaeger and Nagel (\emph{Phys.\ Rev.~E}
51, 1995) tapped glass beads in a tall narrow tube and found that the density approaches its final value
only logarithmically in the number of taps $t$ \PHYStag{experiment},
\[
  \varphi(t)=\varphi_\infty-\frac{\Delta\varphi_\infty}{1+B\ln(1+t/\tau)},
\]
with $\varphi_\infty$ between about $0.637$ and $0.647$ in the fits of Nowak, Knight, Ben-Naim, Jaeger and Nagel
(\emph{Phys.\ Rev.~E} 57, 1998). So tapping creeps towards random close packing and essentially stops there. Once the beads are annealed,
cycling the tapping intensity up and down retraces a reversible curve of density against intensity
(Nowak et al.), a standard testing ground for Edwards' statistical mechanics of powders
(Section~\ref{SM:sec:edwards}).

%======================================================================
\section{Dice, candies and Ulam's conjecture}
\label{PHYS:sec:shapes}

\subsection*{Counting for other shapes}
Lecture~6 ended with an aside on packing ellipsoids and polygons (Section~\ref{L04b:sec:aside}). In the
laboratory, non-spherical particles are the rule. The counting of Section~\ref{PHYS:sec:jam} goes through
for frictionless convex particles with $f$ degrees of freedom each: an isostatic packing has $2f$
constraints per particle on average. Spheres have $f=3$; spheroids (ellipsoids of revolution) $f=5$, since
rotation about the axis does nothing; general ellipsoids and tetrahedra $f=6$.

Donev, Cisse, Sachs, Variano, Stillinger, Connelly, Torquato and Chaikin (\emph{Science} 303, 2004) poured
M\&M candies, which are oblate spheroids, into a flask and measured a random packing fraction of
$0.685\pm0.01$, well above the $0.64$ of spheres \PHYStag{experiment}. Their simulations gave about $0.71$ for
spheroids and up to $0.735$ for general ellipsoids with aspect ratio about $1.3$ \PHYStag{simulation}.
The packing fraction has a cusp-like minimum at the sphere: any small deformation helps. The M\&Ms had
$9.82$ contacts on average, close to $2f=10$, but jammed general ellipsoids in the simulations had at most
about $11.4$, short of $2f=12$. Such \emph{hypostatic} packings are possible because the equivalence of jamming and first-order jamming,
which Interlude~E noted is special to spheres, fails for ellipsoids: the curvature of the surfaces can
block motions at second order, so a jammed packing of ellipsoids need not satisfy the first-order count.

\subsection*{Tetrahedral dice}
Aristotle believed that regular tetrahedra fill space, and Section~\ref{TET:sec} explained why they do not
even fill the space around an edge. Jaoshvili, Esakia, Porrati and Chaikin (\emph{Phys.\ Rev.\ Lett.} 104,
2010) poured tetrahedral dice, of the kind used in role-playing games, into containers and
measured, by volume and by magnetic resonance imaging of the interior, a random packing fraction of
$0.76\pm0.02$ \PHYStag{experiment}: denser than random spheres (by far) and even denser than the densest
sphere packing. A die touches only about $6.3$ neighbours on average, far below $2f=12$. But contacts between
polyhedra come in kinds: a face-to-face contact blocks one translation and two rotations, so counts as three
constraints, an edge-to-face contact as two, and a vertex-to-face or edge-to-edge contact as one. With these
weights the dice had about $12$ constraints per particle, isostatic again.

The densest packings of regular tetrahedra were a race between 2006 and 2010 \PHYStag{simulation}. The densest
lattice packing, of density $\frac{18}{49}\approx0.367$ (Hoylman 1970), is far from optimal; Conway and
Torquato (2006) found $0.72$, Chen (2008) $0.7786$, and then Haji-Akbari, Engel, Keys, Zheng, Petschek, Palffy-Muhoray and
Glotzer (\emph{Nature} 462, 2009), compressing tetrahedra on the computer, found that they
spontaneously form a \emph{dodecagonal quasicrystal}, a cousin of the cut-and-project sets of Interlude~C
(Section~\ref{ERG:sec:cutproject}), whose periodic approximant reaches $0.8503$. Within months the record
passed to explicit periodic packings built from pairs of tetrahedra glued along a face: $\frac{100}{117}\approx0.8547$
(Kallus, Elser and Gravel), $\frac{12250}{14319}\approx0.8555$ (Torquato and Jiao) and
$\frac{4000}{4671}\approx0.8563$ (Chen, Engel and Glotzer, gallery), the densest known. These constructions are exact,
but nobody knows whether $0.8563$ is optimal; the best upper bound is still astronomically close to $1$
(Gravel, Elser and Kallus, gallery). Lagarias and Zong (\emph{Notices AMS} 59, 2012) tell the story.

\subsection*{Is the ball the worst?}
Every one of these shapes packs more densely than the ball. Ulam suspected that this is always so: among
convex bodies in $\R^3$, the ball has the smallest optimal packing density (reported by Martin Gardner,
1995). This is \emph{Ulam's packing conjecture}, and it is open. Numerical work has ruled out many
candidates, including the regular tetrahedron. The best result is local: Kallus (gallery,
Theorem~\ref{PHYS:thm:kallus}) proved that every origin-symmetric convex body sufficiently close to the
ball, but not a ball, packs more densely than the ball, so the ball is a \emph{local pessimum}. In the plane
the disc is not the worst: the regular heptagon is conjectured to be, and among centrally symmetric convex
bodies Reinhardt (1934) conjectured that the worst is a smoothed octagon, of density
$\frac{8-4\sqrt2-\ln2}{2\sqrt2-1}\approx0.9024<\frac{\pi}{\sqrt{12}}\approx0.9069$. For centrally symmetric
bodies in the plane the densest packing is a lattice packing (the aside of Lecture~6), so Reinhardt's
problem is a problem about lattices; it is open too.

\begin{figure}[ht]
\centering
\begin{tikzpicture}
\begin{scope}[x=25cm,xshift=-8.25cm]
  \draw[->] (0.33,0)--(0.885,0) node[right,font=\small]{$\varphi$};
  \foreach \t in {0.35,0.40,0.45,0.50,0.55,0.60,0.65,0.70,0.75,0.80,0.85}
    {\draw (\t,0)--(\t,-0.08) node[below,font=\scriptsize]{$\t$};}
  % spheres row
  \fill[gray!15] (0.33,1.25) rectangle (0.494,1.55);
  \fill[addcol!35] (0.494,1.25) rectangle (0.545,1.55);
  \fill[addcol!15] (0.545,1.25) rectangle (0.7405,1.55);
  \node[font=\scriptsize] at (0.45,1.4) {fluid};
  \node[font=\scriptsize] at (0.68,1.4) {crystal};
  \node[font=\scriptsize,anchor=east] at (0.33,1.4) {spheres};
  \PHYSmsim{(0.384,1.0)} \node[font=\scriptsize,below] at (0.384,0.92) {RSA};
  \PHYSmsim{(0.494,1.0)} \PHYSmsim{(0.545,1.0)}
  \node[font=\scriptsize,above] at (0.5195,1.55) {coexistence};
  \PHYSmexp{(0.555,0.62)} \node[font=\scriptsize,below] at (0.555,0.55) {RLP};
  \PHYSmexp{(0.58,1.0)} \node[font=\scriptsize,below] at (0.58,0.92) {glass};
  \PHYSmexp{(0.64,1.0)} \PHYSmsim{(0.64,0.62)} \node[font=\scriptsize,below] at (0.64,0.55) {RCP, J};
  \PHYSmth{(0.7405,1.0)} \node[font=\scriptsize,below] at (0.7405,0.92) {FCC};
  % non-spherical row
  \node[font=\scriptsize,anchor=east] at (0.33,2.35) {other};
  \draw[gray!50] (0.33,2.35)--(0.885,2.35);
  \PHYSmcon{(0.340,2.35)} \node[font=\scriptsize,above] at (0.340,2.45) {diamond};
  \PHYSmth{(0.3673,2.35)} \node[font=\scriptsize,below] at (0.3673,2.25) {tetra.\ lattice};
  \PHYSmexp{(0.685,2.35)} \node[font=\scriptsize,above] at (0.685,2.45) {M\&M};
  \PHYSmsim{(0.735,2.35)} \node[font=\scriptsize,below] at (0.735,2.25) {ellipsoids};
  \PHYSmexp{(0.76,2.35)} \node[font=\scriptsize,above] at (0.765,2.45) {dice};
  \PHYSmsim{(0.8503,2.35)} \PHYSmcon{(0.8563,2.35)} \node[font=\scriptsize,above] at (0.853,2.45) {tetrahedra};
  % legend
  \PHYSmth{(0.36,3.2)} \node[font=\scriptsize,anchor=west] at (0.365,3.2) {theorem};
  \PHYSmcon{(0.46,3.2)} \node[font=\scriptsize,anchor=west] at (0.465,3.2) {exact construction};
  \PHYSmsim{(0.60,3.2)} \node[font=\scriptsize,anchor=west] at (0.605,3.2) {simulation};
  \PHYSmexp{(0.72,3.2)} \node[font=\scriptsize,anchor=west] at (0.725,3.2) {experiment};
\end{scope}
\end{tikzpicture}
\caption{Packing fractions of this interlude and of Interlude~E, by status. Spheres: random sequential
adsorption ($0.384$), the fluid--crystal coexistence region ($0.494$--$0.545$), random loose packing ($0.555$), the
colloidal glass (about $0.58$), random close packing and point~J (about $0.64$), and the Kepler bound
$\pi/\sqrt{18}$. Other shapes: the diamond packing of balls ($\pi\sqrt3/16$), the densest lattice packing of
regular tetrahedra ($18/49$), M\&M candies, simulated random ellipsoids, tetrahedral dice, the tetrahedral
quasicrystal approximant ($0.8503$) and the densest known tetrahedron packing ($4000/4671$).}
\label{PHYS:fig:strip}
\end{figure}

%======================================================================
\section{Gallery of theorems}
\label{PHYS:sec:gallery}

Results quoted in Interludes~E and~F (Hales, Hales--McLaughlin, Connelly et al., Richthammer, Theil) are not
repeated.

\begin{theorem}[J.~E.~Taylor, Ann.\ of Math.\ 103 (1976)]\label{PHYS:thm:taylor}
Every soap-film-like surface in $\R^3$, in the sense of Almgren's $(\mathbf M,\varepsilon,\delta)$-minimal sets, is
locally either a smooth minimal surface, or three smooth sheets meeting at $120^\circ$ along a curve, or
four such curves meeting at a point at angles $\arccos(-\frac13)$, with six sheets, in the shape of the cone over
the edges of a regular tetrahedron.
\end{theorem}
\noindent Plateau's experimental laws (Section~\ref{PHYS:sec:squash}) as a theorem. The tetrahedral angle is the
Tammes angle for $N=4$.

\begin{theorem}[T.~C.~Hales, Discrete Comput.\ Geom.\ 25 (2001)]\label{PHYS:thm:honeycomb}
Any partition of the plane into regions of equal area has perimeter at least that of the regular hexagonal
honeycomb (in the sense of an averaged limit over large discs).
\end{theorem}
\noindent The two-dimensional Kelvin problem is solved; the three-dimensional one is open.

\begin{question}[Kelvin's problem, 1887]
Which partition of $\R^3$ into cells of unit volume has the least area per cell? Is it the Weaire--Phelan
foam?
\end{question}
\noindent Kelvin's own candidate was beaten by $0.3\%$ in 1994, numerically but with a margin far beyond the
numerical error. That Weaire--Phelan is optimal is open.

\begin{theorem}[H.~Harborth, Elem.\ Math.\ 29 (1974)]\label{PHYS:thm:harborth}
The maximal number of touching pairs among $n$ non-overlapping unit discs in the plane is
$\lfloor3n-\sqrt{12n-3}\rfloor$.
\end{theorem}
\noindent The ground states of $n$ sticky discs. For $n=7$ (a hexagon around a disc) the formula gives
$21-9=12$.

\begin{theorem}[K.~Bezdek and S.~Reid, J.\ Geom.\ 104 (2013)]\label{PHYS:thm:bezdek}
Every packing of $n>2$ unit balls in $\R^3$ has fewer than $6n-0.926\,n^{2/3}$ touching pairs.
\end{theorem}
\noindent A lower bound for the sticky-sphere energy, with a surface term of the right order $n^{2/3}$. Exact values
are known only for $n\le5$ (Section~\ref{PHYS:sec:sticky}).

\begin{theorem}[E.~R.~Chen, M.~Engel and S.~C.~Glotzer, Discrete Comput.\ Geom.\ 44 (2010)]\label{PHYS:thm:CEG}
There is a periodic packing of regular tetrahedra of density $\frac{4000}{4671}\approx0.856347$, with four
tetrahedra (two face-sharing dimers) per fundamental domain. It is optimal within a three-parameter family of
dimer packings.
\end{theorem}
\noindent Found by numerical compression, then made exact; its global optimality is open. (Lattice packings,
whose tetrahedra all point the same way, reach only $\frac{18}{49}$: Hoylman, \emph{Bull.\ Amer.\ Math.\ Soc.}
76, 1970.)

\begin{theorem}[S.~Gravel, V.~Elser and Y.~Kallus, Discrete Comput.\ Geom.\ 46 (2011)]\label{PHYS:thm:GEK}
The packing density of regular tetrahedra is at most $1-2.6\cdots\times10^{-25}$.
\end{theorem}
\noindent A quantitative version of Section~\ref{TET:sec}; the gap between $0.856$ and $1-10^{-25}$ shows how
hard upper bounds are.

\begin{theorem}[Y.~Kallus, Adv.\ Math.\ 264 (2014)]\label{PHYS:thm:kallus}
There is $\varepsilon>0$ such that every origin-symmetric convex body $K\subset\R^3$ that is not a ball, and whose
asphericity is less than $\varepsilon$, has packing density greater than $\pi/\sqrt{18}$. In dimensions $4$ to $8$ and
$24$ there are origin-symmetric bodies of arbitrarily small asphericity whose lattice packing density is
smaller than that of the ball.
\end{theorem}
\noindent The best evidence for Ulam's conjecture. In dimensions $4$--$8$ and $24$ the best lattices for balls
($D_4,\dots,E_8$ of Lectures~4--6, and Leech) are so good that deforming the ball can make things worse.

\begin{theorem}[A.~Figotin and P.~Kuchment, SIAM J.\ Appl.\ Math.\ 56 (1996)]\label{PHYS:thm:FK}
Consider the periodic medium in $\R^d$, $d=2,3$, made of cubes of air, $[\delta,1]^d+\Z^d$, separated by walls of
width $\delta$ of a material with dielectric constant $\varepsilon>1$, and the operator $\nabla\cdot\varepsilon(x)^{-1}\nabla$. For any
$N$, if $\delta$ is small enough and $\varepsilon$ large enough (suitably coupled), the spectrum in $[0,N]$ lies in small
intervals around the Neumann eigenvalues of the cube. In particular the spectrum has gaps.
\end{theorem}
\noindent Band gaps for high contrast, as a theorem (for acoustic waves, or $2$-dimensional photonics); the
Maxwell case of inverse opals rests on computation (Section~\ref{PHYS:sec:opal}).

%======================================================================
\section*{Exercises}
\addcontentsline{toc}{section}{Exercises}

\begin{exercise}
(a) Show that the dihedral angle of the regular dodecahedron is $\arccos(-1/\sqrt5)$. (b) Show that the
truncated octahedron is the Voronoi cell of the body-centred cubic lattice $\Z^3\cup(\Z^3+\frac12(1,1,1))$, and
verify the values $5.345$ and $5.315$ of $A/V^{2/3}$ in Figure~\ref{PHYS:fig:rd} for it and for the rhombic
dodecahedron.
\end{exercise}

\begin{exercise}
(a) Use Proposition~\ref{PHYS:prop:coxeter} to find the average number of sides of a face for Kelvin's cell, for
the Weaire--Phelan cells, and for Matzke's bubbles. (b) The rhombic dodecahedron is not simple. Find the
analogue of the formula $F=12/(6-\bar n)$ for polyhedra in which every vertex has degree $3$ or $4$, and check it
on the rhombic dodecahedron. (c) Show that a convex polyhedron with three faces at each vertex and only
pentagonal and hexagonal faces has exactly $12$ pentagons.
\end{exercise}

\begin{exercise}
In the cell model of Proposition~\ref{PHYS:prop:cell}, at $\varphi=0.60$, compare the cage sizes of FCC and of an amorphous packing jamming at $0.64$, and
estimate the vibrational entropy difference per ball.
\end{exercise}

\begin{exercise}
(a) Show that the number of Barlow stackings of $L$ labelled layers is $3\cdot2^{L-1}$. (b) Show that in every
Barlow stacking each ball has $12$ touching neighbours, and that its Voronoi cell is either the rhombic or the
trapezo-rhombic dodecahedron. (c) Using Proposition~\ref{PHYS:prop:planes}, find the second largest spacing of
lattice planes in FCC, and the corresponding Bragg wavelength of an opal at normal incidence.
\end{exercise}

\begin{exercise}
(a) Redo Proposition~\ref{PHYS:prop:letokhov} for a cube of side $L$ (Dirichlet conditions on its boundary),
and show that the threshold is $L=\pi\sqrt{\ell\ell_g}$.
(b) With $\ell=1/(n\sigma)$ and $n=6\varphi/(\pi D^3)$, how does the critical number of grains in a ball of
radius $R_{\mathrm{cr}}$ depend on $\varphi$?
\end{exercise}

\begin{exercise}
(a) Check that the six points listed in Example~\ref{PHYS:ex:six} form three regular tetrahedra glued face to face,
and that they have exactly $12$ contacts. (b) Verify $\det I=8$ and $\det(AA^T)=864$ for the octahedron.
(\emph{Hint:} $AA^T=2\,\mathrm{Id}+B$, where $B$ records the angles between contacts sharing a ball.) (c) For $n=5$
show that a rigid cluster must be the triangular bipyramid, with $9$ contacts.
\end{exercise}

\begin{exercise}
(a) Carry out the frictional counting of Section~\ref{PHYS:sec:jam} for discs in the plane: show that
$3\le z\le4$. (b) For frictionless tetrahedra, explain why the
isostatic count is $12$ constraints per particle, and check it against the contact weights of
Section~\ref{PHYS:sec:shapes}.
\end{exercise}

\begin{exercise}
Check Harborth's formula for the hexagonal patches $n=3k^2-3k+1$: show that such a patch has $3(3k^2-5k+2)$
contacts and that this equals $\lfloor3n-\sqrt{12n-3}\rfloor$.
\end{exercise}

\section*{Further reading}
\addcontentsline{toc}{section}{Further reading}
\begin{itemize}[nosep]
\item H.~S.~M.~Coxeter, \emph{Close-packing and froth}, Illinois J.\ Math.\ 2 (1958). The statistical honeycombs of
  Proposition~\ref{PHYS:prop:coxeter}.
\item D.~Weaire (ed.), \emph{The Kelvin Problem: Foam Structures of Minimal Surface Area}, Taylor and Francis, 1996;
  D.~Weaire and S.~Hutzler, \emph{The Physics of Foams}, Oxford, 1999.
\item P.~N.~Pusey et al., \emph{Hard spheres: crystallization and glass formation}, Phil.\ Trans.\ R.\ Soc.\ A 367
  (2009), arXiv:0908.1663.
\item J.~D.~Joannopoulos, S.~G.~Johnson, J.~N.~Winn and R.~D.~Meade, \emph{Photonic Crystals: Molding the Flow of
  Light}, 2nd ed., Princeton, 2008.
\item D.~S.~Wiersma, \emph{The physics and applications of random lasers}, Nature Physics 4 (2008).
\item M.~Holmes-Cerfon, \emph{Sticky-sphere clusters}, Annu.\ Rev.\ Condens.\ Matter Phys.\ 8 (2017),
  arXiv:1709.05138; and \emph{Enumerating rigid sphere packings}, SIAM Review 58 (2016).
\item A.~J.~Liu and S.~R.~Nagel, \emph{Jamming is not just cool any more}, Nature 396 (1998).
\item J.~C.~Lagarias and C.~Zong, \emph{Mysteries in packing regular tetrahedra}, Notices AMS 59 (2012).
\end{itemize}
% ===== end chapters/PHYS.tex =====
\endgroup

\part*{Coda}\addcontentsline{toc}{part}{Coda}
\begingroup\renewcommand{\thechapter}{H}\renewcommand{\theHchapter}{interludeH}
% ===== begin chapters/DYN.tex =====
% MAT670 -- Interlude H: Dynkin diagrams everywhere.
% Entirely new material (not in the handwritten notes). Labels prefixed DYN:.
% Placed at the very end of the book. All numerical claims checked with python3
% (Gram determinants, alcove vertices, hole counts, cell volumes, group orders).

\tikzset{
  DYNv/.style={circle,draw,fill=white,inner sep=0pt,minimum size=6.5pt},
  DYNa/.style={circle,draw,fill=black,inner sep=0pt,minimum size=6.5pt},
  DYNx/.style={circle,draw=flagcol,thick,fill=flagcol!35,inner sep=0pt,minimum size=7.5pt},
  DYNm/.style={font=\scriptsize,addcol},
  DYNl/.style={font=\scriptsize},
}

\chapter[Interlude H: Dynkin diagrams everywhere]{Interlude H: Dynkin diagrams everywhere}
\label{DYN:ch}

\begin{addedblock}
This interlude is new: it is not in the handwritten notes. It collects the places where the Dynkin
diagrams of Lectures~4--6 turn up again, in these notes and outside them, and asks in each place what
the diagram lets us read off. To keep it readable it is set in black; all of it should be read as
added material.
\end{addedblock}

Lecture~5 drew Dynkin diagrams in order to \emph{classify} root systems. Once the list was on the
board, the diagrams had done their job. But they kept coming back. Interlude~A found them inside
matrix algebras, Interlude~C folded them into quasicrystals, Interlude~D turned them into braid groups.
This last interlude takes the diagrams as the main characters and asks a few concrete questions about
objects we already know.
\begin{itemize}[nosep]
\item The inequality $\frac1p+\frac1q+\frac1r>1$ that ended the classification in Lecture~5: is it the
  same inequality that lists the Platonic solids?
\item Interlude~E proved that the icosahedron is not made of regular tetrahedra. Can the diagram of
  the icosahedron see that?
\item The FCC lattice has octahedral and tetrahedral holes, a cubic and a tetrahedral description,
  and the cuboctahedron as its kissing configuration. Which of these are visible on the diagram
  \tikz[baseline=-3pt]{\node[DYNv] (a) at (0,0) {}; \node[DYNv] (b) at (0.45,0) {}; \node[DYNv] (c) at (0.9,0) {};
  \draw (a)--(b)--(c);}\,?
\item Lecture~6 built $E_8$ by filling the deep holes of $D_8$. Can we \emph{see} the holes of $E_8$
  itself?
\end{itemize}
The answers are all yes, and most of them come from one extra node: the extended diagram of
Interlude~A (Section~\ref{LIE:sec:readoff}, item~7). At the end we tour the other places where the
ADE list appears (singularities, quivers, conformal field theory, subfactors, cubic surfaces), and
record in each case what the diagram says there.

Throughout, simply laced root systems are normalized with roots of norm $\ip\alpha\alpha=2$, so that
the Cartan matrix $C$ is the Gram matrix of the simple roots. We use the labelling of $E_8$ from
Lecture~6 (Section~\ref{L04b:sec:E8}): $\alpha_1,\dots,\alpha_6$ is a chain, $\alpha_7=e_6+e_7$ is
attached to $\alpha_5$, and $\alpha_8=-\frac12(1,\dots,1)$ is attached to $\alpha_7$.

%======================================================================
\section{One inequality, three times}
\label{DYN:sec:pqr}

\subsection*{The branch point, as a determinant}
Recall the last step of the classification (Section~\ref{L04b:sec:coxeter-cont}). A diagram with one
branch point has a centre $x$ and three arms of $p-1$, $q-1$, $r-1$ nodes. Call it $T_{p,q,r}$
(Figure~\ref{DYN:fig:Tpqr}); it has $p+q+r-2$ nodes. The lecture projected $x$ onto the span of the
arms and found that the diagram is admissible exactly when
$\frac1p+\frac1q+\frac1r>1$. The same computation, done with the Cartan matrix, gives an exact formula.

\begin{proposition}\label{DYN:prop:detT}
The Cartan matrix of $T_{p,q,r}$ (all edges simple, roots of norm $2$) has determinant
\[
  \det C(T_{p,q,r})=pqr\Big(\frac1p+\frac1q+\frac1r-1\Big).
\]
In particular $C(T_{p,q,r})$ is positive definite if and only if $\frac1p+\frac1q+\frac1r>1$.
\end{proposition}

\begin{proof}
Order the nodes with the centre first. Then
$C=\left(\begin{smallmatrix}2&b^T\\ b&C'\end{smallmatrix}\right)$, where
$C'=C(A_{p-1})\oplus C(A_{q-1})\oplus C(A_{r-1})$ is the Cartan matrix of the three arms and $b$ has an
entry $-1$ at the three arm nodes next to the centre. The Schur complement formula gives
\[
  \det C=\det C'\cdot\big(2-b^TC'^{-1}b\big).
\]
Two facts about the chain $A_k$ finish the job: $\det C(A_k)=k+1$ (expand along the last row:
$d_k=2d_{k-1}-d_{k-2}$), and the corner entry of $C(A_k)^{-1}$ is
$d_{k-1}/d_k=k/(k+1)$. Hence
\[
  \det C=pqr\Big(2-\frac{p-1}p-\frac{q-1}q-\frac{r-1}r\Big)=pqr\Big(\frac1p+\frac1q+\frac1r-1\Big).
\]
Since $C'$ is positive definite, $C$ is positive definite exactly when the Schur complement is
positive.
\end{proof}

The three terms $\frac{p-1}p$ are, up to the factor $2$ from our normalization, the squared
projections $\frac12(1-\frac1p)$ of the lecture. So the proof of Lecture~5 \emph{was} a Schur complement.
The formula is more than a test of positivity, though. Plugging in:
{\small\setlength{\arraycolsep}{4pt}\[
  \begin{array}{l|cccccccc}
  (p,q,r) & (1,q,r) & (2,2,n) & (2,3,3) & (2,3,4) & (2,3,5) & \begin{smallmatrix}(3,3,3),(2,4,4),\\(2,3,6)\end{smallmatrix} & (2,3,7)\\ \hline
  T_{p,q,r} & A_{q+r-1} & D_{n+2} & E_6 & E_7 & E_8 & \tilde E_6,\ \tilde E_7,\ \tilde E_8 & E_{10}\\
  \det C & q+r & 4 & 3 & 2 & 1 & 0 & -1
  \end{array}
\]}
These are the determinants $\det C$ that Interlude~A used to read off the packing densities of the
root lattices, now all from one formula. On the $E$-series $E_n=T_{2,3,n-3}$ the formula is linear:
\begin{equation}\label{DYN:eq:Edet}
  \det C(E_n)=6(n-3)\Big(\frac1{n-3}-\frac16\Big)=9-n .
\end{equation}
At $n=9$ the determinant vanishes. $E_9=T_{2,3,6}$ is the extended diagram $\tilde E_8$, the
diagram of the affine Weyl group that tiles $\R^8$ by alcoves (Section~\ref{LIE:sec:torus}). At
$n=10$ the form becomes indefinite, of signature $(9,1)$. The determinant $-1$ says the $E_{10}$ lattice is
unimodular; it is the even unimodular Lorentzian lattice $\mathrm{II}_{9,1}$.

\begin{figure}[ht]
\centering
\begin{tikzpicture}[scale=0.78]
  \begin{scope}
    \node[DYNv] (x) at (0,0) {}; \node[DYNl,below=2pt] at (0,0) {$x$};
    \node[DYNv] (u1) at (-1,0) {}; \node[DYNv] (u2) at (-2.6,0) {};
    \node[DYNv] (v1) at (1,0) {}; \node[DYNv] (v2) at (2.6,0) {};
    \node[DYNv] (w1) at (0,1) {}; \node[DYNv] (w2) at (0,2.4) {};
    \draw (u1)--(x)--(v1); \draw (x)--(w1);
    \draw[dotted,thick] (u1)--(u2); \draw[dotted,thick] (v1)--(v2); \draw[dotted,thick] (w1)--(w2);
    \draw[gray] (-2.7,-0.3)--(-2.7,-0.4)--(-0.9,-0.4)--(-0.9,-0.3); \node[DYNl,below] at (-1.8,-0.4) {$p-1$};
    \draw[gray] (0.9,-0.3)--(0.9,-0.4)--(2.7,-0.4)--(2.7,-0.3); \node[DYNl,below] at (1.8,-0.4) {$q-1$};
    \draw[gray] (-0.3,0.9)--(-0.4,0.9)--(-0.4,2.5)--(-0.3,2.5); \node[DYNl,left] at (-0.4,1.7) {$r-1$};
    \node at (0,-1.4) {$T_{p,q,r}$};
  \end{scope}
  % E6~ = T333
  \begin{scope}[xshift=5.2cm]
    \foreach \i in {-2,...,2} \node[DYNv] (a\i) at (0.7*\i,0) {};
    \node[DYNv] (b1) at (0,0.7) {}; \node[DYNa] (b2) at (0,1.4) {};
    \foreach \i [evaluate=\i as \j using int(\i+1)] in {-2,...,1} \draw (a\i)--(a\j);
    \draw (a0)--(b1)--(b2);
    \node at (0,-1.4) {$\tilde E_6=T_{3,3,3}$};
  \end{scope}
  % E7~ = T244
  \begin{scope}[xshift=9.3cm]
    \foreach \i in {-3,...,3} \node[DYNv] (a\i) at (0.6*\i,0) {};
    \node[DYNa] (a-3) at (-1.8,0) {};
    \node[DYNv] (b1) at (0,0.7) {};
    \foreach \i [evaluate=\i as \j using int(\i+1)] in {-3,...,2} \draw (a\i)--(a\j);
    \draw (a0)--(b1);
    \node at (0,-1.4) {$\tilde E_7=T_{2,4,4}$};
  \end{scope}
  % E8~ = T236
  \begin{scope}[xshift=14.2cm]
    \foreach \i in {-5,...,2} \node[DYNv] (a\i) at (0.55*\i,0) {};
    \node[DYNa] (a-5) at (-2.75,0) {};
    \node[DYNv] (b1) at (0,0.7) {};
    \foreach \i [evaluate=\i as \j using int(\i+1)] in {-5,...,1} \draw (a\i)--(a\j);
    \draw (a0)--(b1);
    \node at (-0.6,-1.4) {$\tilde E_8=T_{2,3,6}$};
  \end{scope}
\end{tikzpicture}
\caption{The diagram $T_{p,q,r}$, with arms of $p-1$, $q-1$, $r-1$ nodes, and the three cases with
$\frac1p+\frac1q+\frac1r=1$. These are the extended diagrams $\tilde E_6,\tilde E_7,\tilde E_8$; the
black node is the extra (affine) node.}
\label{DYN:fig:Tpqr}
\end{figure}

\subsection*{Triangles on the sphere}
Now forget Dynkin diagrams and draw a triangle on the unit sphere with angles $\frac\pi p$, $\frac\pi q$,
$\frac\pi r$. By Gauss--Bonnet its area is the angle excess,
\[
  \operatorname{area}=\pi\Big(\frac1p+\frac1q+\frac1r-1\Big),
\]
so the triangle exists on the sphere exactly when $\frac1p+\frac1q+\frac1r>1$. It exists in the
Euclidean plane when the sum equals $1$, and in the hyperbolic plane when it is less than $1$. The
reflections in its three sides generate a group that tiles the sphere by copies of the triangle. There
are $4\pi/\operatorname{area}$ copies, and the rotations (products of two reflections) form a subgroup
of index $2$:
\[
  \abs{G^+_{p,q,r}}=\frac{2}{\frac1p+\frac1q+\frac1r-1}.
\]
For $(2,3,3)$, $(2,3,4)$, $(2,3,5)$ this gives $12$, $24$, $60$: the rotation groups of the tetrahedron,
the cube (or octahedron), and the icosahedron (or dodecahedron). For $(2,2,n)$ it gives the dihedral
group of order $2n$, the symmetries of a regular $n$-gon viewed as a flat polyhedron (a ``dihedron'').
The integers $p,q,r$ are the orders of the rotation axes: for the icosahedron, $2$ through the edge
midpoints, $3$ through the face centres and $5$ through the vertices. The Euclidean cases $(3,3,3)$,
$(2,4,4)$, $(2,3,6)$ are the reflection groups of the equilateral triangle, the isosceles right
triangle and the $30$--$60$--$90$ triangle, the kaleidoscopes of the hexagonal and square lattices.

Comparing with Proposition~\ref{DYN:prop:detT} we get a curious identity:
\[
  \det C(T_{p,q,r})=\frac{2pqr}{\abs{G^+_{p,q,r}}}\qquad\text{e.g.}\quad
  \det C(E_8)=\frac{2\cdot30}{60}=1,\quad \det C(E_6)=\frac{2\cdot18}{12}=3.
\]
So \emph{the branch diagrams $D_n,E_6,E_7,E_8$ are the dihedra and the Platonic solids}. (The chains
$A_n=T_{1,q,r}$ do not fit this pattern on the nose; Section~\ref{DYN:sec:elsewhere} matches them with
the cyclic groups.) The arms of $E_8$ have lengths $1,2,4$ because the icosahedron
has rotations of orders $2,3,5$. So far this is a coincidence of numbers. Section~\ref{DYN:sec:elsewhere}
explains it: the resolution of the singularity $\C^2/\Gamma$, for $\Gamma$ the binary icosahedral group,
\emph{is} a $T_{2,3,5}$ of curves.

\subsection*{Coxeter's criterion}
There is a third appearance, which puts the first two in one frame. A \emph{Coxeter diagram} has a node for each
generator $s_i$ and an edge labelled $m_{ij}$ when $(s_is_j)^{m_{ij}}=1$ (unlabelled means $3$, no edge
means $2$). Its \emph{Schl\"afli} (or Gram) matrix is
\[
  G_{ij}=-\cos\frac{\pi}{m_{ij}},\qquad G_{ii}=1,
\]
the Gram matrix of unit normals to mirrors meeting at angles $\pi/m_{ij}$. For a crystallographic
diagram, $G$ is the Cartan matrix rescaled ($G=\frac12C$ in the simply laced case).

\begin{theorem}[Coxeter 1934]\label{DYN:thm:coxeter}
Let $\Gamma$ be a connected Coxeter diagram with Schl\"afli matrix $G$. The group
$\langle s_i\mid s_i^2,(s_is_j)^{m_{ij}}\rangle$ is a finite reflection group if and only if $G$ is
positive definite. It acts as a discrete group generated by reflections of a Euclidean space, with a
simplex as fundamental domain, if and only if $G$ is positive semidefinite with a one-dimensional kernel.
The connected positive definite diagrams are $A_n,B_n,D_n,E_6,E_7,E_8,F_4,G_2$, together with the
non-crystallographic $H_3$, $H_4$ and $I_2(m)$.
\end{theorem}

For three mirrors forming a triangle $(p,q,r)$,
$\det G=1-a^2-b^2-c^2-2abc$ with $a=\cos\frac\pi p$, $b=\cos\frac\pi q$, $c=\cos\frac\pi r$. We checked
numerically that this is positive, zero, negative in the spherical, Euclidean, hyperbolic cases
(for example $\det G=0.0955$ for $(2,3,5)$ and $-0.0617$ for $(2,3,7)$). A Platonic solid $\{p,q\}$ has
the linear diagram $\circ\overset{p}{\text{---}}\circ\overset{q}{\text{---}}\circ$, whose determinant is
$1-\cos^2\frac\pi p-\cos^2\frac\pi q$. It is positive if and only if
$\frac1p+\frac1q>\frac12$, the classical condition for $\{p,q\}$ to exist. This is the triangle
$(2,p,q)$ again. So the classification of Platonic solids, of spherical triangle groups and of
branched Dynkin diagrams are three readings of one inequality.

%======================================================================
\section[The icosahedron, seen by Coxeter]{The icosahedron is not made of regular tetrahedra, seen by Coxeter}
\label{DYN:sec:icosa}

Interlude~E (Section~\ref{TET:sec}) proved, with a ring of six balls and the triangle inequality on
the sphere, that $20$ regular tetrahedra glued at a common vertex do not close up into an icosahedron.
Around each of the $12$ edges from the centre there would be $5$ tetrahedra, and $5$ dihedral angles
$\arccos\frac13\approx70.53^\circ$ leave a gap of $7.36^\circ$. Now look at the same configuration
through its Coxeter diagram.

A regular polytope or honeycomb with Schl\"afli symbol $\{p,q,r\}$ has the linear Coxeter diagram
\[
  \circ\overset{p}{\text{---}}\circ\overset{q}{\text{---}}\circ\overset{r}{\text{---}}\circ .
\]
The first three nodes generate the symmetry group of a cell $\{p,q\}$, the last three that of the vertex
figure $\{q,r\}$, and the label $r$ says that $r$ cells surround each edge. ``Twenty tetrahedra around a
vertex, forming an icosahedron'' is the local structure of $\{3,3,5\}$: cell a tetrahedron $\{3,3\}$,
vertex figure an icosahedron $\{3,5\}$, five cells around an edge.

\begin{proposition}\label{DYN:prop:gram335}
The Schl\"afli matrix of $\{p,q,r\}$ has determinant
\[
  \det G_{\{p,q,r\}}=\sin^2\frac\pi p\,\sin^2\frac\pi r-\cos^2\frac\pi q .
\]
For $\{3,3,r\}$ this is $\frac12-\frac34\cos^2\frac\pi r$, and
\[
  \det G_{\{3,3,r\}}\ \gtreqless\ 0\quad\iff\quad r\cdot\arccos\tfrac13\ \lesseqgtr\ 2\pi .
\]
The values are $\frac5{16}$, $\frac18$, $\frac{5-3\tau}{16}\approx0.0091$, $-\frac1{16}$ for
$r=3,4,5,6$, where $\tau=\frac{1+\sqrt5}2$.
\end{proposition}

\begin{proof}
For a tridiagonal matrix with $1$ on the diagonal and off-diagonal entries $-a,-b,-c$, the leading
minors satisfy $D_k=D_{k-1}-(\text{entry})^2D_{k-2}$. So $D_2=1-a^2$, $D_3=1-a^2-b^2$, and
\[
  D_4=D_3-c^2D_2=(1-a^2)(1-c^2)-b^2 .
\]
With $a=b=\frac12$ this is $\frac34(1-c^2)-\frac14=\frac12-\frac34c^2$. It vanishes if and only if
$c^2=\cos^2\frac\pi r=\frac23$, that is, if and only if $\cos\frac{2\pi}r=2c^2-1=\frac13$, that is,
$\frac{2\pi}r=\arccos\frac13$. As $r$ grows, $c^2$ grows, so the sign of $\det G$ is the sign of
$2\pi-r\arccos\frac13$. For $r=5$, $c=\cos\frac\pi5=\frac\tau2$ and $c^2=\frac{\tau+1}4$, which gives
$\frac12-\frac{3(\tau+1)}{16}=\frac{5-3\tau}{16}$.
\end{proof}

So the sign of a $4\times4$ determinant is exactly the sign of the gap left by $r$ regular tetrahedra
around an edge. The critical value is $r=2\pi/\arccos\frac13\approx5.104$, not an integer, so
there is no Euclidean honeycomb of regular tetrahedra (Figure~\ref{DYN:fig:gram}):
\begin{itemize}[nosep]
\item $r=3,4,5$: positive definite, spherical. These are the $5$-cell $\{3,3,3\}$, the $16$-cell
  $\{3,3,4\}$ and the $600$-cell $\{3,3,5\}$, with groups $A_4$, $B_4$, $H_4$.
\item $r=6$: indefinite, hyperbolic. Six tetrahedra around an edge overshoot, and $\{3,3,6\}$ is a
  honeycomb of hyperbolic space with ideal vertices.
\item Euclidean: never.
\end{itemize}
The icosahedral case $r=5$ is spherical, but only just: $(5-3\tau)/16\approx0.0091$, against
$\frac5{16}$ and $\frac18$ for $r=3,4$. This small number is the $7.36^\circ$ of Interlude~E. A flat
icosahedron of regular tetrahedra would need $\det G=0$. On the $3$-sphere the curvature closes the
gap, and $600$ regular spherical tetrahedra form the $600$-cell, the ``ideal'' icosahedral order of
Section~\ref{SM:sec:curved}. Its Schl\"afli determinant is positive, which says that $H_4$ is a finite
reflection group and that the tiling closes up on a sphere. In Schl\"afli's original form the criterion
reads $\sin\frac\pi p\sin\frac\pi r>\cos\frac\pi q$, and for $\{3,3,5\}$ it is
$0.5090>0.5$.

\begin{figure}[ht]
\centering
\begin{tikzpicture}[xscale=1.25,yscale=9]
  \draw[-{Stealth}] (2.6,0)--(8.3,0) node[right,DYNl]{$r$};
  \draw[-{Stealth}] (2.8,-0.15)--(2.8,0.37) node[above,DYNl]{$\det G_{\{3,3,r\}}$};
  \foreach \y in {-0.1,0.1,0.2,0.3} {\draw (2.75,\y)--(2.85,\y); \node[left,DYNl] at (2.75,\y) {$\y$};}
  \draw[thick,addcol,domain=2.9:7.7,samples=120] plot (\x,{0.5-0.75*cos(180/\x)^2});
  \foreach \r/\v in {3/0.3125,4/0.125,5/0.00912,6/-0.0625,7/-0.1088} {\fill (\r,\v) ellipse (0.05 and 0.007); \draw[gray!60] (\r,0)--(\r,-0.006); \node[below,DYNl] at (\r,-0.006) {$\r$};}
  \fill[flagcol] (5.104,0) ellipse (0.06 and 0.008);
  \draw[flagcol,-{Stealth}] (5.9,0.12)--(5.15,0.012);
  \node[flagcol,DYNl,align=left,anchor=west] at (5.9,0.135) {$r=2\pi/\arccos\frac13\approx5.104$\\Euclidean case, not an integer};
  \node[DYNl,anchor=west] at (3.1,0.32) {$5$-cell};
  \node[DYNl,anchor=south west] at (4.05,0.13) {$16$-cell};
  \node[DYNl,anchor=south east] at (4.92,0.012) {$600$-cell};
  \node[DYNl,anchor=south west] at (6.1,-0.055) {hyperbolic};
\end{tikzpicture}
\caption{The Schl\"afli determinant $\frac12-\frac34\cos^2\frac\pi r$ of $\{3,3,r\}$, with $r$ treated as
a real variable. It crosses zero where $r$ regular tetrahedra exactly fill the space around an edge.
The value at $r=5$ is positive but tiny.}
\label{DYN:fig:gram}
\end{figure}

\paragraph{Reading the $600$-cell off its diagram.} The diagram of $H_4$ also counts. By Wythoff's
construction the $k$-faces of $\{p,q,r\}$ correspond to cosets of the subgroup generated by all nodes
except the $(k+1)$-st, so their number is $\abs W$ divided by the order of that subgroup. With
$\abs{W(H_4)}=14400$ (Interlude~C):
\begin{gather*}
  \text{vertices: }\frac{14400}{\abs{W(H_3)}}=\frac{14400}{120}=120,\qquad
  \text{edges: }\frac{14400}{\abs{W(A_1)}\,\abs{W(I_2(5))}}=\frac{14400}{2\cdot10}=720,\\
  \text{triangles: }\frac{14400}{\abs{W(A_2)}\,\abs{W(A_1)}}=\frac{14400}{6\cdot2}=1200,\qquad
  \text{cells: }\frac{14400}{\abs{W(A_3)}}=\frac{14400}{24}=600 .
\end{gather*}
The vertex stabilizer is $H_3$, the symmetry group of the icosahedral vertex figure.

\subsection*{Why a lattice cannot do it}
The label $5$ is also what keeps icosahedral order out of crystals.

\begin{proposition}[Crystallographic restriction]\label{DYN:prop:cryst}
If a rotation of order $m$ maps a lattice $\Lat\subset\R^d$ to itself, then its characteristic
polynomial has integer coefficients. For $d=2,3$ this forces $m\in\{1,2,3,4,6\}$. More generally, a
lattice with a symmetry of order $5$ has dimension at least $4$. If a group isomorphic to the icosahedral
rotation group $A_5$ preserves a lattice in $\R^d$ and acts on some $3$-dimensional subspace as the
rotations of an icosahedron, then $d\ge6$.
\end{proposition}

\begin{proof}
In a basis of $\Lat$ the rotation is an integer matrix. For $d=2$ its trace $2\cos\frac{2\pi}m$ is an
integer, and for $d=3$ the trace $1+2\cos\frac{2\pi}m$ is; either way
$\cos\frac{2\pi}m\in\{0,\pm\frac12,\pm1\}$. An integer matrix of order $5$ has eigenvalues that are fifth
roots of unity, closed under Galois conjugation; a primitive one brings all four primitive ones with it,
so $d\ge4$. For the icosahedral group $A_5$, the two irreducible real representations of dimension $3$ have
characters involving $\sqrt5$ and are swapped by $\sqrt5\mapsto-\sqrt5$. The representation on
$\Lat\otimes\Q$ is rational, so if it contains one of them it contains the other equally often, and
$d\ge3+3$. (The condition on the $3$-dimensional subspace matters: $A_5\subset\Sigma_5=W(A_4)$ preserves the
$4$-dimensional lattice $A_4$, acting through its rational $4$-dimensional representation.)
\end{proof}

For a reflection group the rotation $s_is_j$ has order $m_{ij}$, so the proposition is the
crystallographic table of Lecture~4 again: $4\cos^2\frac\pi m\in\{0,1,2,3\}$ allows
$m\in\{2,3,4,6\}$, and $4\cos^2\frac\pi5=\tau+1$ is irrational. The bounds are attained. The Coxeter
element of $A_4$ has order $h=5$ (Interlude~A, item~3), and Interlude~C (Theorem~\ref{ERG:thm:fold})
realizes $H_3$ inside $D_6$ and $H_4$ inside $E_8$. That is the precise sense in which the $E_8$ lattice
``contains'' the icosahedral order that is impossible in a flat crystal: the $240$ roots project to two
$600$-cells. In three dimensions the best one can do is a quasicrystal.

%======================================================================
\section{The FCC lattice, read off its diagrams}
\label{DYN:sec:fcc}

The FCC lattice is the root lattice $A_3$, the densest packing in $\R^3$ (Hales). We now go through
its features one by one and find each on a diagram.

\subsection*{Two names: $A_3=D_3$}
Lectures~5 and~6 met FCC twice: as $A_3=\{x\in\Z^4:\sum x_i=0\}$ and, in the remark $D_3=A_3$ after the
classification, as $D_3=\{x\in\Z^3:x_1+x_2+x_3\text{ even}\}$.
The diagrams agree: $D_3$ would be a chain of two nodes both joined to a third, which is the chain
$A_3$. That is why Theorem~\ref{L04a:thm:classification} starts $D_n$ at $n=4$. An explicit isometry is
\[
  f(x)=\tfrac12\big(x_1+x_2-x_3-x_4,\ x_1-x_2+x_3-x_4,\ x_1-x_2-x_3+x_4\big).
\]
It sends $e_1,e_2,e_3,e_4$ to the four vertices $\frac12(1,1,1)$, $\frac12(1,-1,-1)$,
$\frac12(-1,1,-1)$, $\frac12(-1,-1,1)$ of a regular tetrahedron inscribed in the cube, and
$e_1-e_2\mapsto(0,1,1)$, $e_1-e_3\mapsto(1,0,1)$, and so on, so the $12$ roots of $A_3$ go to the $12$ roots
$\pm e_a\pm e_b$ of $D_3$. Since the $f(e_i)$ have norm $\frac34$ and inner products $-\frac14$, one checks
$\abs{f(x)}^2=\abs x^2$ whenever $\sum x_i=0$. In Lie theory this is $\mathfrak{so}_6\cong\mathfrak{sl}_4$
(Interlude~A, end of Section~\ref{LIE:sec:classical}). Geometrically, $\SL_4$ acts on
$\Lambda^2\C^4\cong\C^6$ preserving the quadratic form $\omega\mapsto\omega\wedge\omega\in\Lambda^4\C^4\cong\C$.
Its quadric in $\mathbb P^5$ is the Klein quadric, whose points are the lines of $\mathbb P^3$.

Keep an eye on the vectors $f(e_i)=\frac12(\pm1,\pm1,\pm1)$: the coordinate directions of $\R^4$ point
at holes of FCC, as we will see in a moment.

\subsection*{The symmetry group, and the missing inversion}
The Weyl group $W(A_3)=\Sigma_4$, of order $24$, permutes the four coordinates of $\R^4$; in the cube
picture it is the full symmetry group of the inscribed tetrahedron. But the cube has $48$ symmetries,
and they all preserve $D_3$. The extra ones come from the diagram. The automorphism group of an
irreducible root system is $W\rtimes\operatorname{Aut}(\text{diagram})$ (Humphreys, \emph{Introduction
to Lie Algebras}, \S12.2), and for a root lattice it is the same as the automorphism group of the lattice,
since the roots are the minimal vectors and generate.

\begin{proposition}\label{DYN:prop:autD}
$\operatorname{Aut}(D_n)=W(B_n)$, of order $2^nn!$, for every $n\ge2$ except $n=4$, where
$\operatorname{Aut}(D_4)=W(F_4)$ has order $1152=192\cdot6$.
\end{proposition}

\begin{proof}
$W(B_n)$ (all signed permutations) preserves $D_n$, and has order $2^nn!$. The group
$W(D_n)$ (signed permutations with an even number of sign changes) has order $2^{n-1}n!$. The diagram
of $D_n$ has symmetry group $\Z/2$ for $n\ne4$ (for $n=2$, swap the two nodes of $A_1\times A_1$; for
$n=3$, flip the chain $A_3$), and $\Sigma_3$ for $n=4$. So $\abs{\operatorname{Aut}D_n}=2\cdot2^{n-1}n!=2^nn!$
for $n\ne4$, and $\operatorname{Aut}D_n\supseteq W(B_n)$ forces equality. For $n=4$ the order is
$6\cdot192=1152=\abs{W(F_4)}$, and $W(F_4)$ preserves the $D_4$ roots, which are its long roots. (A computer
search over the possible images of the simple roots confirms the orders $48$ for $D_3$ and $1152$ for $D_4$.)
\end{proof}

For FCC this says: point group $O_h=W(B_3)=W(A_3)\times\{\pm1\}$. The inversion $x\mapsto-x$ is not in
$W(A_3)$; it acts on the diagram $A_3$ as the flip. (In general $-1\in W$ exactly when the opposite
involution $-w_0$ is trivial, which fails for $A_n$ with $n\ge2$, for $D_n$ with $n$ odd, and for
$E_6$.) The flip will turn out to swap two kinds of holes.

\subsection*{The holes}
Holes are the vertices of the Voronoi cells (Definition~\ref{L1:def:voronoi}). For the root lattices of
this interlude they are exactly the points locally farthest from the lattice (for a general lattice a
Voronoi vertex need not be a local maximum of the distance).
Around each one, the nearest lattice points span a \emph{Delaunay cell}. In the cube picture FCC has
two kinds of holes (with roots $\pm e_a\pm e_b$, so minimal distance $\sqrt2$):
\begin{itemize}
\item \emph{octahedral holes} at $(1,0,0)+D_3$, at distance $1$ from six lattice points
  $0,\ 2e_1,\ e_1\pm e_2,\ e_1\pm e_3$, which span a regular octahedron of edge $\sqrt2$;
\item \emph{tetrahedral holes} at $\frac12(\pm1,\pm1,\pm1)+D_3$, at distance $\frac{\sqrt3}2$ from four
  lattice points, for example $0,(1,1,0),(1,0,1),(0,1,1)$ around $\frac12(1,1,1)$, which span a regular
  tetrahedron of edge $\sqrt2$.
\end{itemize}
The tetrahedral holes come in two classes, $h+D_3$ and $-h+D_3$ with $h=\frac12(1,1,1)$, whose
tetrahedra point in opposite directions. Together with $D_3$ itself these are the four cosets of $D_3$ in
the dual lattice $D_3^*=\Z^3\cup(\Z^3+h)$:
\[
  D_3^*/D_3=\{[0],[1],[2],[3]\}\cong\Z/4,\qquad [0]=D_3,\ \ [1]=h+D_3,\ \ [2]=e_1+D_3,\ \ [3]=-h+D_3 .
\]
It is cyclic: $h+h=(1,1,1)\in e_1+D_3$, so two steps from a lattice point to a tetrahedral hole
reach an octahedral hole.

\paragraph{Counting cells.} The octahedron has volume $\frac43$ and the tetrahedron $\frac13$, and the
covolume of $D_3$ is $2$. There is one octahedral hole per lattice point (the coset $[2]$ has the same
density as $D_3$) and two tetrahedral holes ($[1]$ and $[3]$), and indeed
\[
  1\cdot\tfrac43+2\cdot\tfrac13=2 .
\]
So FCC space is tiled by octahedra and tetrahedra, and the tetrahedra fill exactly one third of it.
Now put the FCC balls (radius $\frac{\sqrt2}2$) at the lattice points and ask what fraction of each
Delaunay cell they cover. In a tetrahedron this is $4\Omega_T/(4\pi)\cdot\vol\Ball/\vol T$, with
$\Omega_T=3\arccos\frac13-\pi$ the solid angle at a vertex; that is Rogers' $\sigma_3\approx0.7796$
(Section~\ref{L2:sec:rogers}). In the octahedron, with $\Omega_O=4\arcsin\frac13$, it is
$\approx0.7209$. The weighted average
$\frac13(0.7796)+\frac23(0.7209)=0.7405=\pi/\sqrt{18}$ is the FCC density. So the ratio $0.950$ to
Rogers' bound in the table of Section~\ref{L2:sec:rogers} is the price of the octahedra: they take two
thirds of the space, at local density $0.7209$ instead of $0.7796$.

\paragraph{Around an edge.} Each edge of this tiling (a pair of touching balls) lies on two tetrahedra
and two octahedra, alternating. The dihedral angle of the octahedron is $\arccos(-\frac13)\approx109.47^\circ$,
and since
\[
  \arccos\tfrac13+\arccos\big(-\tfrac13\big)=\pi ,
\]
two of each close up exactly (Figure~\ref{DYN:fig:edge}). So the octahedra fill \emph{precisely} the
gap that tetrahedra leave. Five tetrahedra fall $7.36^\circ$ short of $2\pi$, six overshoot by $63.17^\circ$,
and FCC settles for two tetrahedra and two octahedra. At a vertex the same bookkeeping gives $8$ tetrahedra and $6$
octahedra, the $8$ triangles and $6$ squares of the cuboctahedron, which is the vertex figure.

\begin{figure}[ht]
\centering
\begin{tikzpicture}[scale=1.25]
  \begin{scope}
    \foreach \k in {0,...,4} {
      \pgfmathsetmacro{\a}{\k*70.5288}
      \pgfmathsetmacro{\b}{(\k+1)*70.5288}
      \fill[addcol!20] (0,0)--(\a:1.2) arc (\a:\b:1.2)--cycle;
      \draw (0,0)--(\a:1.2);
      \node[DYNl] at ({(\a+\b)/2}:0.8) {T};
    }
    \draw (0,0)--(352.644:1.2);
    \fill[flagcol!50] (0,0)--(352.644:1.2) arc (352.644:360:1.2)--cycle;
    \draw[flagcol,-{Stealth}] (1.65,-0.45)--(1.22,-0.08);
    \node[flagcol,DYNl] at (1.75,-0.6) {gap $7.36^\circ$};
    \node[DYNl] at (0,-1.55) {five tetrahedra: $5\arccos\frac13<2\pi$};
  \end{scope}
  \begin{scope}[xshift=4.3cm]
    \foreach \s/\e/\l/\c in {0/70.5288/T/addcol!20, 70.5288/180/O/gray!20, 180/250.5288/T/addcol!20, 250.5288/360/O/gray!20} {
      \fill[\c] (0,0)--(\s:1.2) arc (\s:\e:1.2)--cycle;
      \draw (0,0)--(\s:1.2);
      \node[DYNl] at ({(\s+\e)/2}:0.8) {\l};
    }
    \node[DYNl] at (0,-1.55) {FCC: $2\arccos\frac13+2\arccos(-\frac13)=2\pi$};
  \end{scope}
  \begin{scope}[xshift=8.6cm]
    \foreach \k in {0,...,5} {
      \pgfmathsetmacro{\a}{\k*60}
      \pgfmathsetmacro{\b}{(\k+1)*60}
      \fill[addcol!20] (0,0)--(\a:1.2) arc (\a:\b:1.2)--cycle;
      \draw (0,0)--(\a:1.2);
      \node[DYNl] at ({(\a+\b)/2}:0.8) {$\triangle$};
    }
    \node[DYNl] at (0,-1.55) {the plane: $6\cdot\frac\pi3=2\pi$};
  \end{scope}
\end{tikzpicture}
\caption{Cross-sections perpendicular to an edge. Left: five regular tetrahedra leave a gap (the
icosahedron problem of Interlude~E). Middle: in the FCC Delaunay tiling, octahedra (O) fill the gaps
between tetrahedra (T). Right: in the plane, six equilateral triangles close up around a vertex, which is why
Rogers' bound is sharp there (Thue).}
\label{DYN:fig:edge}
\end{figure}

\subsection*{The holes on the extended diagram}
Now the diagram. The extended diagram $\tilde A_3$ is a square: the extra node $\alpha_0=-\theta$ is
joined to both ends of the chain. Its four nodes match the four cosets:
\[
  \text{node }0\leftrightarrow[0]\ \text{(lattice points)},\quad
  \text{node }1\leftrightarrow[1],\quad \text{node }2\leftrightarrow[2]\ \text{(octahedral)},\quad
  \text{node }3\leftrightarrow[3].
\]
The reason will be given in Section~\ref{DYN:sec:holes}: the fundamental alcove of the affine Weyl group
has vertices $0,\omega_1,\omega_2,\omega_3$, the $\omega_i$ are representatives of the cosets, and under $f$
they become $h=\frac12(1,1,1)$, $e_1$ and $\frac12(1,1,-1)$. Two symmetries of the square are visible in the lattice.
\begin{itemize}
\item \emph{The flip} fixing nodes $0$ and $2$ and swapping $1$ and $3$ is realized by the inversion $x\mapsto-x$. It
  fixes the lattice and the octahedral holes, and swaps the two classes of tetrahedral holes, ``up'' and
  ``down''.
\item \emph{The rotation} by a quarter turn is translation by $h$: it moves lattice points to
  tetrahedral holes, those to octahedral holes, and so on around the square. It is a symmetry of the whole
  pattern of points and holes, though not of the lattice.
\end{itemize}
The symmetry group of the square, dihedral of order $8$, is $(D_3^*/D_3)\rtimes\operatorname{Aut}(A_3)
=\Z/4\rtimes\Z/2$. This is a general fact (Section~\ref{DYN:sec:holes}).

\paragraph{Chemistry.} Filling some of the four cosets with atoms gives the classical cubic crystal
structures:
\[
  \begin{array}{ll}
  [0]\cup[2]=\Z^3 & \text{rock salt (NaCl): Na on FCC, Cl in the octahedral holes,}\\
  {}[0]\cup[1] & \text{zincblende (ZnS), and diamond when both are carbon,}\\
  {}[0]\cup[1]\cup[3] & \text{fluorite (CaF$_2$): F in all tetrahedral holes, two per Ca,}\\
  {}[0]\cup[1]\cup[2]\cup[3]=D_3^* & \text{body-centred cubic (BCC).}
  \end{array}
\]
The last line says that the dual of FCC is BCC, which is the root-lattice statement $A_3^*\cong$ BCC.
FCC is the densest packing in $\R^3$, and its dual BCC, the lattice of all its points and holes, is the
thinnest \emph{lattice covering} of $\R^3$ (Bambah 1954), with covering density
$\frac{5\sqrt5\pi}{24}\approx1.4635$.

\paragraph{Stacking.} One dimension down, $\tilde A_2$ is a triangle, and the hexagonal lattice has two
classes of (triangular) holes swapped by the flip of $A_2$ (Figure~\ref{DYN:fig:hex}). Lamination
(Section~\ref{L04b:sec:laminated}) places the next layer over one class of holes. The choice between the two
classes at each step is the choice between FCC ($ABCABC$) and the hexagonal close packing ($ABAB$),
and the two choices are exchanged by the symmetry of $\tilde A_2$.

\subsection*{The kissing configuration and the jitterbug}
The $12$ nearest neighbours of a lattice point are the roots of $A_3$, the vertices of the
cuboctahedron. Wythoff's construction reads this off the diagram: the roots form the Weyl group orbit
of the highest root $\theta=\omega_1+\omega_3$, so the cuboctahedron is the $A_3$ diagram with both
end nodes ringed. In the cube picture ($B_3$) it is the middle node ringed, the ``rectified cube''. Its
$8$ triangles face the $8$ tetrahedral holes around the ball and its $6$ squares face the $6$ octahedral
holes. The icosahedron, by contrast, is $\{3,5\}$, the $H_3$ diagram with an end node ringed, and its
diagram needs the label $5$ that no lattice allows. Lecture~11's jitterbug deforms one into the other. It
keeps the $8$ triangles over the tetrahedral holes rigid, rotating them, and folds each square over an
octahedral hole into two triangles: $8+6\cdot2=20$ faces. (In Coxeter's terminology the icosahedron is the
\emph{snub tetrahedron}.) What is lost along the way is the
lattice, since the end point carries a forbidden $5$.

%======================================================================
\section{Seeing the holes: the extended diagram}
\label{DYN:sec:holes}

Everything in the FCC story came from the four nodes of $\tilde A_3$. Here is the general mechanism,
which underlies the systematic treatments of Conway--Sloane and Moody--Patera (gallery).

\subsection*{The alcove and its vertices}
Let $\Delta$ be a simply laced irreducible root system with simple roots $\alpha_1,\dots,\alpha_n$, root
lattice $L$, highest root $\theta=\sum_im_i\alpha_i$, and fundamental weights $\omega_i$
($\ip{\omega_i}{\alpha_j}=\delta_{ij}$). The integers $m_i\ge1$ are the \emph{marks}; put $m_0=1$ for the
extra node $\alpha_0=-\theta$ of the extended diagram. The \emph{fundamental alcove} is the simplex
\[
  \mathcal A=\{x:\ \ip x{\alpha_i}\ge0\ (1\le i\le n),\ \ip x\theta\le1\},
\]
with one wall for each node of the extended diagram. Its vertices are
\[
  v_0=0,\qquad v_i=\frac{\omega_i}{m_i}\quad(1\le i\le n),
\]
since $v_i$ lies on every wall except the wall of $\alpha_i$. For $A_2$ this is the triangle in
Figure~\ref{DYN:fig:hex}. The affine Weyl group, generated by the reflections in the walls of
$\mathcal A$, tiles space by copies of $\mathcal A$ (Interlude~A).

\begin{lemma}\label{DYN:lem:relevant}
For a simply laced root lattice the Voronoi cell of $0$ is
\[
  V(0)=\{x:\ \ip x\alpha\le1\text{ for all roots }\alpha\}=\bigcup_{w\in W}w(\mathcal A).
\]
\end{lemma}

\begin{proof}[Proof for $E_8$]
In general $V(0)=\{x:\ip xv\le\frac12\abs v^2\}$, where $v$ runs over the \emph{Voronoi-relevant} vectors,
and by a theorem of Voronoi, $v$ is relevant if and only if $\pm v$ are the only shortest vectors in its
class $v+2L$ (Conway--Sloane, \emph{Sphere Packings}, Ch.~21). So we sort the short vectors of $E_8$ into
the $2^8=256$ classes of $E_8/2E_8$. Two vectors $v,w$ are in the same class when $\frac12(v-w)\in E_8$,
which needs $\frac14\abs{v-w}^2\in\{0,2,4,\dots\}$, since $E_8$ is even with minimum $2$.
\begin{itemize}[nosep]
\item No vector of norm $2$ or $4$ is in the class of $0$, since $\frac12v$ would have norm $\frac12$ or $1$.
\item Two roots are congruent only if $\frac14(4-2\ip\alpha\beta)\in\{0,2\}$, that is, $\beta=\pm\alpha$. So the
  $240$ roots give $120$ distinct classes.
\item A root $\alpha$ and a vector $v$ of norm $4$ are never congruent: we would need both
  $\frac14(6-2\ip v\alpha)$ and $\frac14(6+2\ip v\alpha)$ (for $v-\alpha$ and $v+\alpha$) in $\{0,2,4,\dots\}$, but
  they add up to $3$.
\item Two vectors of norm $4$ are congruent only if $\frac14(8-2\ip vw)\in\{0,2,4\}$, that is, $w=\pm v$ or
  $w\perp v$. So a class holds at most $8$ orthogonal pairs, $16$ vectors, and the $2160$ vectors of norm
  $4$ (Interlude~B, \eqref{NT:eq:E8theta}) need at least $135$ classes.
\end{itemize}
Since $1+120+135=256$, these are all the classes, each norm-$4$ class has exactly $16$ shortest vectors,
and only the root classes have exactly two. So the relevant vectors are the $240$ roots, with
$\frac12\abs\alpha^2=1$. Finally, on the dominant chamber $\ip x\alpha\le\ip x\theta$ for every root
(since $\theta-\alpha$ is a nonnegative combination of simple roots), so $V(0)$ meets the chamber in
$\mathcal A$, and $W$ permutes the chambers.
\end{proof}

Conway and Sloane (\emph{Sphere Packings}, Ch.~21) prove the same for every root lattice: the relevant
vectors are the roots. Consequently every vertex of a Voronoi cell, that is, every hole, is a $W$-image of
some $v_i$. Which $v_i$ are holes, and what do they look like?

\begin{theorem}[Holes from the extended diagram]\label{DYN:thm:holes}
Let $i\neq0$ be a node of the extended diagram $\tilde\Delta$ with mark $m_i$, and let $W_i$ be the
(finite) reflection group of the diagram $\tilde\Delta\setminus\{i\}$.
\begin{enumerate}[label=(\alph*)]
\item The roots $\alpha_j$, $j\neq i$ (including $\alpha_0$), generate a sublattice of $L$ of index $m_i$.
\item The lattice points nearest to $v_i$ are the $W_i$-orbit of $0$, at distance $\abs{v_i}$.
\item If $\tilde\Delta\setminus\{i\}$ is connected, then $v_i$ is a vertex of $V(0)$ (a hole), and its
  Delaunay cell is the Wythoff polytope of the diagram $\tilde\Delta\setminus\{i\}$ with node $0$ ringed.
  The holes of this type around $0$ number $\abs W/\abs{W(\Delta\setminus\{i\})}$.
\item If $\tilde\Delta\setminus\{i\}$ is disconnected, $v_i$ lies in the relative interior (at the
  centre) of a face of $V(0)$ of positive dimension, so it is not a hole.
\end{enumerate}
\end{theorem}

\begin{proof}[Sketch]
(a) The lattice generated by $\alpha_j$ ($j\ne i$, $j\ge1$) and $\alpha_0=-\sum_jm_j\alpha_j$ is spanned by
the $\alpha_j$ with $j\ne i$ together with $m_i\alpha_i$; its index in $L$ is $m_i$.
(b) The alcoves containing $v_i$ are the images of $\mathcal A$ under $W_i$, the group generated by the
walls through $v_i$. By Lemma~\ref{DYN:lem:relevant}, a lattice point $\lambda$ is nearest to $v_i$ exactly
when $v_i\in V(\lambda)=\lambda+\bigcup_ww(\mathcal A)$, i.e.\ when $v_i$ lies in one of the alcoves
$\lambda+w(\mathcal A)$. The alcoves form a simplicial tiling on which the affine Weyl group acts
preserving the type (node label) of vertices, so $v_i$ can lie in such an alcove only as its vertex of
type $i$, and then $\lambda$ is the alcove's vertex of type $0$. The alcoves with $v_i$ as a vertex are
the $W_i$-images of $\mathcal A$, so these $\lambda$ are the points $W_i\cdot0$.
(c),(d) The group $W_i$ splits as a product over the components of $\tilde\Delta\setminus\{i\}$. The
components not containing node $0$ consist of reflections in linear hyperplanes $\ip x{\alpha_j}=0$, so
they fix $0$ (and $v_i$). The orbit $W_i\cdot0$ is therefore the Wythoff polytope of the component
containing node $0$, with node $0$ ringed, centred at $v_i$; it spans $v_i$ plus the span of that
component's roots, so it is full-dimensional exactly when that component is everything. A point is a
vertex of $V(0)$ exactly when its nearest lattice points (here including $0$) affinely span $\R^n$; this
gives (c). In case (d), $v_i$ lies in the relative interior of the face of $V(0)$ dual to the
lower-dimensional polytope $\conv(W_i\cdot0)$, and the other components act on that face fixing only
$v_i$. The count is the size of the $W$-orbit of $v_i$, whose stabilizer in $W$ is the parabolic
subgroup $W(\Delta\setminus\{i\})$ generated by the simple reflections fixing $\omega_i$.
\end{proof}

Two consequences can be read off immediately. First, the \emph{index} in (a) is the mark: deleting a node
of mark $m$ produces a root sublattice of index $m$ (this is the lattice shadow of the Borel--de Siebenthal
theorem in the gallery). Second, the nodes of mark $1$ are special. For them $v_i=\omega_i$, and the
$\omega_i$ with $m_i=1$, together with $0$, represent the classes of $L^*/L$. So
\[
  \#\{\text{nodes of mark }1\text{ in }\tilde\Delta\}=\abs{L^*/L}=\det C,
\]
and the symmetry group of the extended diagram is $(L^*/L)\rtimes\operatorname{Aut}(\Delta)$: the
translations by the $\omega_i$ ($m_i=1$) rotate the extended diagram, as the quarter turn of $\tilde A_3$ did.
The check: $\tilde A_n$ has $n+1$ nodes of mark $1$; $\tilde D_n$ has $4$; $\tilde E_6$, $\tilde E_7$,
$\tilde E_8$ have $3$, $2$, $1$; and these are the determinants $n+1,4,3,2,1$.

\begin{figure}[ht]
\centering
\begin{tikzpicture}[scale=0.95]
  \begin{scope}
    \clip (-2.3,-1.85) rectangle (2.3,1.85);
    \foreach \a in {-3,...,3} \foreach \b in {-3,...,3} {
      \pgfmathsetmacro{\x}{\a+0.5*\b} \pgfmathsetmacro{\y}{0.8660*\b}
      \fill (\x,\y) circle (0.045);
      \fill[addcol] (\x+0.5,\y+0.2887) circle (0.05);
      \draw[flagcol,thick,fill=white] (\x+1,\y+0.5774) circle (0.05);
    }
  \end{scope}
  \draw[thick] (30:0.5774)--(90:0.5774)--(150:0.5774)--(210:0.5774)--(270:0.5774)--(330:0.5774)--cycle;
  \fill[gray!35,opacity=0.8] (0,0)--(30:0.5774)--(90:0.5774)--cycle;
  \draw[-{Stealth}] (0,0)--(60:1);
  \node[DYNl,anchor=west] at (2.6,1.0) {$\bullet$ lattice $[0]$};
  \node[DYNl,addcol,anchor=west] at (2.6,0.55) {$\bullet$ holes $[1]$};
  \node[DYNl,flagcol,anchor=west] at (2.6,0.1) {$\circ$ holes $[2]$};
  \node[DYNl,align=left,anchor=west] at (2.6,-0.6) {grey: alcove $\mathcal A$\\hexagon: $V(0)=\bigcup_{W}w\mathcal A$};
  \begin{scope}[xshift=7.6cm]
    \node[DYNa] (n0) at (90:0.75) {}; \node[DYNv] (n1) at (210:0.75) {}; \node[DYNv] (n2) at (330:0.75) {};
    \draw (n0)--(n1)--(n2)--(n0);
    \node[DYNl,above=2pt] at (90:0.75) {$0$: lattice};
    \node[DYNl,left=2pt,addcol] at (210:0.75) {$1$: hole $[1]$};
    \node[DYNl,right=2pt,flagcol] at (330:0.75) {$2$: hole $[2]$};
    \node at (0,-1.4) {$\tilde A_2$};
  \end{scope}
\end{tikzpicture}
\caption{The hexagonal lattice $A_2$. The alcove (grey) has vertices $0$, $\omega_1$, $\omega_2$; its six
$W$-images make the Voronoi hexagon, whose vertices are the two classes of holes. The nodes of the
triangle $\tilde A_2$ match the three cosets of $A_2$ in $A_2^*$; the flip of $A_2$ swaps the hole
classes (the $B$/$C$ choice in stacking). The arrow is the root $\theta$.}
\label{DYN:fig:hex}
\end{figure}

\subsection*{Worked examples}

\paragraph{$A_n$.} All marks are $1$ and $\tilde A_n$ is a cycle, so every node is a hole. The hole
$\omega_i$ (class $[i]$) has
\[
  \abs{\omega_i}^2=\frac{i(n+1-i)}{n+1},
\]
and its Delaunay cell is the Wythoff polytope of $A_n$ with the $i$-th node ringed: the \emph{hypersimplex},
the convex hull of the $\binom{n+1}i$ vectors with $i$ ones and $n+1-i$ zeros. For $n=3$:
$i=1,3$ give tetrahedra at distance $\frac{\sqrt3}2$, and $i=2$ gives the octahedron at distance $1$. For
$n=2$ both hole classes are triangles, which is why Rogers' bound is attained by $A_2$. For $n\ge3$ the
hypersimplex with $i=2$ is not a simplex, so Rogers' bound is never attained by $A_n$. The deepest holes
are in the middle, $i=\lfloor\frac{n+1}2\rfloor$.

\paragraph{$D_n$, and a footnote to Lecture~6.} In $\tilde D_n$ ($n\ge4$) the four leaves have mark
$1$ and all other nodes have mark $2$. The inner nodes disconnect the diagram, so the holes are the three
leaves other than $0$:
\begin{itemize}
\item the \emph{vector} hole $\omega_1=e_1$, at distance $1$. Deleting node $1$ leaves $D_n$ with node $0$ at
  the end of the long arm, so the Delaunay cell is the cross-polytope, with $2n$ vertices $0,2e_1,e_1\pm e_j$;
\item the two \emph{spinor} holes $\omega_{n-1},\omega_n=\frac12(1,\dots,1,\pm1)$, at distance $\frac{\sqrt n}2$.
  Now node $0$ sits at a short leaf, and the Delaunay cell is the half-cube, with $2^{n-1}$ vertices.
\end{itemize}
So the deep holes of $D_n$ are the spinor holes $h=\frac12(1,\dots,1)$ for $n\ge5$, as Lecture~6 used
(Section~\ref{L04b:sec:E8}). For $n=3$ the roles reverse: the spinor holes are the tetrahedral holes
of FCC, at distance $\frac{\sqrt3}2<1$, and the deep holes are the octahedral (vector) ones. At $n=4$ all three are
at distance $1$, and all three Delaunay cells are $16$-cells. That tie is triality
(Section~\ref{DYN:sec:triality}).

\paragraph{$E_8$.} Here is the full table for $\tilde E_8=T_{2,3,6}$, in the labelling of Lecture~6. The
extra node $\alpha_0$ is attached to $\alpha_1$, the end of the long arm. For each node we list the mark,
the norm of the alcove vertex $v_i=\omega_i/m_i$, the number of nearest lattice points, and what they span
(all computed numerically).
\begin{center}
\renewcommand{\arraystretch}{1.15}
\small
\begin{tabular}{lcccccccc}
\toprule
node & $\alpha_1$ & $\alpha_2$ & $\alpha_3$ & $\alpha_4$ & $\alpha_5$ & $\alpha_6$ & $\alpha_7$ & $\alpha_8$\\
\midrule
mark $m_i$ & $2$ & $3$ & $4$ & $5$ & $6$ & $3$ & $4$ & $2$\\
$\abs{v_i}^2$ & $\frac12$ & $\frac23$ & $\frac34$ & $\frac45$ & $\frac56$ & $\frac89$ & $\frac78$ & $1$\\
nearest points & $2$ & $3$ & $4$ & $5$ & $6$ & $9$ & $8$ & $16$\\
$\tilde E_8\setminus\{i\}$ & $A_1{\times}E_7$ & $A_2{\times}E_6$ & $A_3{\times}D_5$ & $A_4{\times}A_4$ & $A_5{\times}A_2{\times}A_1$ & $\mathbf{A_8}$ & $A_7{\times}A_1$ & $\mathbf{D_8}$\\
\bottomrule
\end{tabular}
\end{center}
The pattern $\abs{v_i}^2=\frac k{k+1}$ for a cell with $k+1$ points is the squared circumradius of a
regular $k$-simplex of edge $\sqrt2$, which is what the nearest points span in all but the last column.
Only two deletions leave a connected diagram, so $E_8$ has exactly two kinds of holes (Figure~\ref{DYN:fig:ext}):
\begin{itemize}
\item the \emph{shallow} holes, from $\alpha_6$ (mark $3$): Delaunay cell the regular $8$-simplex (the $A_8$
  diagram ringed at an end), circumradius $\sqrt{8/9}$, and $\abs{W(E_8)}/\abs{W(A_7)}=17280$ of them around
  each lattice point;
\item the \emph{deep} holes, from $\alpha_8$ (mark $2$): Delaunay cell the $8$-dimensional cross-polytope
  ($D_8$ ringed at the end of the long arm, $16$ vertices), circumradius $1$, and
  $\abs{W(E_8)}/\abs{W(D_7)}=2160$ of them.
\end{itemize}
So the Voronoi cell of $E_8$ has $17280+2160=19440$ vertices. Both numbers can be found in the theta series
of Interlude~B. The $2160$ deep holes around $0$ are the halves of the $2160$ vectors of norm $4$. The
$17280$ shallow holes are the thirds of the $17520$ vectors of norm $8$, minus the $240$ that are twice a
root ($17520-240=17280$).

\paragraph{Lecture~6 again.} Part (a) of the theorem now explains both constructions of $E_8$. Deleting
$\alpha_8$ (mark $2$) leaves $D_8$, a sublattice of index $2$: this is $E_8=D_8\cup(D_8+h)$. Of the three
nonzero classes of $D_8^*/D_8$, the spinor class of $h$ (distance $\sqrt2$) has been filled, and the vector
class (distance $1$) survives as the deep holes of $E_8$. Deleting $\alpha_6$ (mark $3$) leaves $A_8$, of
index $3$, and gives a construction we did not see in Lecture~6:
\[
  E_8=A_8\cup(A_8+[3])\cup(A_8+[6]),
\]
where $[3]$ is the class of $\omega_3$ in $A_8^*/A_8\cong\Z/9$. It has norm $\frac{3\cdot6}9=2$, so the
glued lattice is even, and it is unimodular since $9/3^2=1$.

\paragraph{How much space is Rogers'?} Per lattice point there are $17280/9=1920$ simplices and
$2160/16=135$ cross-polytopes. With edge $\sqrt2$ their volumes are $\frac1{13440}$ and $\frac{256}{8!}$, so
\[
  1920\cdot\frac1{13440}=\frac17,\qquad 135\cdot\frac{256}{40320}=\frac67 .
\]
One seventh of $\R^8$ is tiled by Rogers' regular simplices, where the local density is
$\sigma_8\approx0.2568$, and six sevenths by cross-polytopes, where it is about $0.2531$. This is why
$E_8$ comes within $1.2\%$ of Rogers' bound (Section~\ref{L2:sec:rogers}): its non-simplex cells,
the cross-polytopes, are almost as efficient as Rogers' simplices. The deep holes
give the covering radius $1$, and the covering density $\vol\Ball^8(1)=\pi^4/24\approx4.06$.

\begin{figure}[ht]
\centering
\begin{tikzpicture}[scale=0.78]
  % A3~ square
  \begin{scope}
    \node[DYNa] (q0) at (0,1) {}; \node[DYNv] (q1) at (1,0) {}; \node[DYNv] (q2) at (0,-1) {}; \node[DYNv] (q3) at (-1,0) {};
    \draw[gray,dashed] (0,0.88)--(0,-0.88);
    \draw (q0)--(q1)--(q2)--(q3)--(q0);
    \node[DYNl,above=2pt] at (0,1) {$[0]$ lattice};
    \node[DYNl,right=2pt] at (1,0) {$[1]$ tetra};
    \node[DYNl,below=2pt] at (0,-1) {$[2]$ octa};
    \node[DYNl,left=2pt] at (-1,0) {tetra $[3]$};
    \draw[-{Stealth},addcol] (70:1.35) arc (70:20:1.35);
    \node[DYNl,addcol] at (45:1.7) {$+h$};
    \node at (0,-2.2) {$\tilde A_3$ (FCC)};
  \end{scope}
  % D4~
  \begin{scope}[xshift=4.3cm]
    \node[DYNv] (c) at (0,0) {}; \node[DYNm,below=2pt] at (0,0) {$2$};
    \node[DYNa] (d0) at (-1,0) {}; \node[DYNl,left=2pt] at (-1,0) {$0$};
    \node[DYNv] (d1) at (60:1) {}; \node[DYNv] (d3) at (0:1) {}; \node[DYNv] (d4) at (-60:1) {};
    \foreach \d in {d0,d1,d3,d4} \draw (c)--(\d);
    \node[DYNl,right=2pt] at (60:1) {$e_1$};
    \node[DYNl,right=2pt] at (0:1) {$\frac12(1,1,1,1)$};
    \node[DYNl,right=2pt] at (-60:1) {$\frac12(1,1,1,-1)$};
    \draw[-{Stealth},addcol] (3.2,0.9) arc (40:-40:1.4);
    \node[DYNl,addcol,anchor=west] at (3.3,0) {triality};
    \node at (0.6,-2.2) {$\tilde D_4$};
  \end{scope}
  % E8~
  \begin{scope}[xshift=10.2cm]
    \foreach \i/\m/\x in {0/1/0,1/2/0.8,2/3/1.6,3/4/2.4,4/5/3.2,5/6/4.0,7/4/4.8,8/2/5.6} {
      \node[DYNv] (e\i) at (\x,0) {}; \node[DYNm,below=2pt] at (\x,0) {$\m$};}
    \node[DYNa] (e0) at (0,0) {};
    \node[DYNv] (e6) at (4.0,0.8) {}; \node[DYNm,right=2pt] at (4.0,0.8) {$3$};
    \draw (e0)--(e1)--(e2)--(e3)--(e4)--(e5)--(e7)--(e8); \draw (e5)--(e6);
    \node[DYNx] at (4.0,0.8) {}; \node[DYNx] at (5.6,0) {};
    \node[DYNl,above=2pt] at (0,0) {$\alpha_0$};
    \node[DYNl,above=2pt] at (0.8,0) {$\alpha_1$};
    \node[DYNl,above=2pt] at (3.2,0) {$\alpha_4$};
    \node[DYNl,left=3pt] at (4.0,0.8) {$\alpha_6$};
    \node[DYNl,above=2pt] at (4.8,0) {$\alpha_7$};
    \node[DYNl,above=4pt] at (5.6,0) {$\alpha_8$};
    \node[DYNl,flagcol,anchor=south] at (4.0,1.05) {delete $\alpha_6$: $A_8$, simplex hole, $r^2=\frac89$};
    \node[DYNl,flagcol,align=center] at (4.9,-1.05) {delete $\alpha_8$: $D_8$,\\ cross-polytope hole, $r^2=1$};
    \node at (2.8,-2.2) {$\tilde E_8=T_{2,3,6}$};
  \end{scope}
\end{tikzpicture}
\caption{Extended diagrams (extra node black, marks in blue). Left: the square $\tilde A_3$ with the four
cosets of FCC in BCC; the dashed axis is the flip (inversion), the arrow the quarter turn (translation by
$h=\frac12(1,1,1)$). Middle: $\tilde D_4$, whose three outer nodes other than $0$ are the three hole
classes of $D_4$, permuted by triality. Right: $\tilde E_8$; the two shaded leaves are the only nodes whose
deletion leaves a connected diagram, and they give the two kinds of holes of $E_8$.}
\label{DYN:fig:ext}
\end{figure}

%======================================================================
\section{Triality}
\label{DYN:sec:triality}

The diagram $D_4$ has a threefold symmetry, and so does $\tilde D_4$ after fixing node $0$
(Figure~\ref{DYN:fig:ext}). On the lattice it acts as follows.

\paragraph{Three classes of holes.} $D_4^*/D_4\cong\Z/2\times\Z/2$ has three nonzero classes, represented
by $e_1$, $\frac12(1,1,1,1)$ and $\frac12(1,1,1,-1)$. Each has exactly $8$ vectors of norm $1$:
\[
  \pm e_i,\qquad \tfrac12(\pm1,\pm1,\pm1,\pm1)\ \text{(even number of $-$)},\qquad
  \tfrac12(\pm1,\pm1,\pm1,\pm1)\ \text{(odd number of $-$)},
\]
and each octet is the vertex set of a $16$-cell (a cross-polytope). By the previous section these are
the deep holes around $0$, all at distance $1$, and each one's Delaunay cell is again a $16$-cell. The
group $W(D_4)$ acts trivially on $D_4^*/D_4$, and $\operatorname{Aut}(D_4)=W(F_4)$ acts on the three classes through
$\Sigma_3$. That is the order count $1152=192\cdot6$ of Proposition~\ref{DYN:prop:autD}.

\paragraph{$F_4$ is $D_4$ together with its holes.} The $24$ holes around $0$ form a $24$-cell, the Voronoi
cell of $D_4$. Compare with Lecture~6's description of $F_4$: the long roots $\pm e_i\pm e_j$ are the
$24$ roots of $D_4$, and the short roots $\pm e_i$, $\frac12(\pm1,\dots,\pm1)$ are exactly the $24$ deep
holes. So
\[
  \Delta(F_4)=\{\text{nearest neighbours of }0\text{ in }D_4\}\cup\{\text{vertices of the Voronoi cell of }0\},
\]
two $24$-cells, polar to each other (the facets of the Voronoi cell are perpendicular to the roots), the
second smaller by a factor $\sqrt2$. The self-duality of the
$24$-cell $\{3,4,3\}$ is the left--right symmetry of its Coxeter diagram
$\circ\text{---}\circ\overset{4}{\text{---}}\circ\text{---}\circ$, which as a Dynkin diagram reverses the
arrow and exchanges long and short roots. And Lecture~6's remark that the $F_4$ root lattice is
$D_4^*$ is the statement that the holes generate it.

\paragraph{Representations.} In Interlude~A's language (Section~\ref{LIE:sec:readoff}, item~8), the three
outer nodes of $D_4$ label the three $8$-dimensional representations of $\mathfrak{so}_8$: the vector
representation and the two half-spin representations. Their weights are exactly the three octets above.
Triality permutes them, and the $\mathfrak{so}_8$-invariant trilinear map
$V\times S_+\times S_-\to\R$ is octonion multiplication in disguise (Baez, \emph{The octonions}). The
subgroup of $\mathrm{Spin}(8)$ fixed by triality is the automorphism group $G_2$ of the octonions.

\paragraph{Folding by triality, on the roots.} Interlude~A listed the folding $D_4\rightsquigarrow G_2$.
Here it is on the $24$ roots. Triality fixes the plane spanned by the central root $\alpha_2$ and
$\alpha_1+\alpha_3+\alpha_4$. Projecting all roots onto this plane, the $6$ roots fixed by triality keep
norm $2$, and the other $18$ fall in $6$ orbits of $3$, each projecting to a vector of norm $\frac23$. The
result is $12$ vectors in two hexagons with squared length ratio $3$: the root system $G_2$ of
Lecture~4.

\paragraph{Lamination.} Laminating $D_4$ means placing a new layer over a deep hole
(Section~\ref{L04b:sec:laminated}). There are three classes to choose from, and triality makes the choices
equivalent. So the laminated lattice $\Lambda_5$ is unique; it is $D_5$.

%======================================================================
\section[Accidental isomorphisms and the $E$-series]{Accidental isomorphisms, and the \texorpdfstring{$E$}{E}-series below \texorpdfstring{$E_6$}{E6}}
\label{DYN:sec:accidents}

The ranges $B_n$ ($n\ge2$), $C_n$ ($n\ge3$), $D_n$ ($n\ge4$) in Theorem~\ref{L04a:thm:classification} hide
coincidences of small diagrams. Interlude~A recorded the Lie algebra isomorphisms. Here is what each says
about lattices and geometry.
\begin{center}
\renewcommand{\arraystretch}{1.25}
\small
\begin{tabular}{p{0.17\textwidth}p{0.25\textwidth}p{0.48\textwidth}}
\toprule
diagrams & groups & lattices and geometry\\
\midrule
$A_1=B_1=C_1$ & $\mathrm{SU}(2)/\pm1=\SO(3)$ & roots $\pm\alpha$; unit quaternions rotate $\R^3$ by conjugation (Interlude~A)\\
$B_2=C_2$ & $\mathrm{Sp}(4)\cong\mathrm{Spin}(5)$ & $B_2$ and $C_2$ differ by a $45^\circ$ rotation and a factor $\sqrt2$ (Lecture~6); root lattices $\Z^2$ and the checkerboard $D_2$\\
$D_2=A_1\times A_1$ & $\SO(4)=\frac{\mathrm{SU}(2)\times\mathrm{SU}(2)}{\pm(1,1)}$ & roots $e_1\pm e_2$ orthogonal; rotations of $\R^4=\mathbb H$ are $x\mapsto qx\bar r$\\
$D_3=A_3$ & $\mathrm{Spin}(6)\cong\mathrm{SU}(4)$ & FCC in two coordinate systems; lines in $\mathbb P^3$ $=$ Klein quadric\\
$E_5=D_5$\newline $E_4=A_4$\newline $E_3=A_2\times A_1$ & $\mathrm{Spin}(10)$\newline $\mathrm{SU}(5)$\newline $\mathrm{SU}(3)\times\mathrm{SU}(2)$ & del Pezzo surfaces of degree $4,5,6$ (below)\\
\bottomrule
\end{tabular}
\end{center}

\subsection*{The $E$-series continued downward}
Delete nodes one at a time from the end of the long arm of $E_8$ ($\alpha_1$, then $\alpha_2$, and so on, in
our labelling). This gives $E_7=T_{2,3,4}$, $E_6=T_{2,3,3}$, and then $E_5=T_{2,3,2}=D_5$, $E_4=T_{2,3,1}=A_4$ and finally
$E_3=A_2\times A_1$, when the branch point is reached and the $T$ falls apart. Lecture~6 found the same
chain inside $E_8$ by collapsing chains (Section~\ref{L04b:sec:E8}). The determinants obey
\eqref{DYN:eq:Edet} all the way down: $\det C(E_n)=9-n$ gives $6,5,4,3,2,1$ for $n=3,\dots,8$, and
indeed $\det(A_2\times A_1)=3\cdot2$, $\det A_4=5$, $\det D_5=4$. Physicists will recognize the chain
$\mathrm{SU}(5)\subset\mathrm{Spin}(10)\subset E_6$ of grand unified theories, with $E_3$ the
$\mathrm{SU}(3)\times\mathrm{SU}(2)$ of the standard model.

The number $9-n$ has a geometric meaning. Blow up the projective plane in $r$ points in general position
($3\le r\le8$). The surface $X_r$ has $H_2(X_r,\Z)=\Z h\oplus\Z e_1\oplus\dots\oplus\Z e_r$ with intersection
form $\operatorname{diag}(1,-1,\dots,-1)$ (a line $h$, and the exceptional curves $e_i$), and canonical class
$K=-3h+e_1+\dots+e_r$, with $K\cdot K=9-r$.

\begin{proposition}\label{DYN:prop:delpezzo}
For $3\le r\le8$ the orthogonal complement $K^\perp\subset H_2(X_r,\Z)$, with the sign of the form reversed,
is the root lattice $E_r$. Its roots are $e_i-e_j$ and $\pm(h-e_i-e_j-e_k)$ (and, for $r\ge6$, more), with
simple roots $e_1-e_2,\dots,e_{r-1}-e_r,\ h-e_1-e_2-e_3$. Its determinant is $K\cdot K=9-r$.
\end{proposition}

\begin{proof}[Sketch]
The listed simple roots are orthogonal to $K$, have square $-2$, and their intersection graph is a chain
$e_1-e_2,\dots,e_{r-1}-e_r$ with $h-e_1-e_2-e_3$ attached to $e_3-e_4$: the diagram $T_{2,3,r-3}=E_r$
(for $r=3$ there is no $e_3-e_4$, and $h-e_1-e_2-e_3$ is orthogonal to the chain: $A_2\times A_1$).
They span $K^\perp$ (exercise). For the determinant, $H_2$ is unimodular and $K$ is primitive, and for a
primitive vector $v$ of a unimodular lattice with $v\cdot v\ne0$ the discriminant of $v^\perp$ is
$\abs{v\cdot v}$. Since $K\cdot K>0$ and the form has signature $(1,r)$, $K^\perp$ is negative definite.
\end{proof}

For $r\le8$ these $X_r$ are the \emph{del Pezzo surfaces} of degree $K\cdot K=9-r$ ($-K$ is ample; for
$r\le6$ it embeds $X_r$ in $\mathbb P^{9-r}$ as a surface of degree $9-r$, the cubic surface for $r=6$). The classes $E$ with $E\cdot E=E\cdot K=-1$ are the lines (exceptional curves), and $W(E_r)$
permutes them transitively. For $r\ge4$ their numbers are \emph{also} read off the diagrams, as the index
$\abs{W(E_r)}/\abs{W(E_{r-1})}$ (a line can be blown down, leaving $X_{r-1}$):
\begin{center}
\renewcommand{\arraystretch}{1.15}
\begin{tabular}{lcccccc}
\toprule
$r$ & $3$ & $4$ & $5$ & $6$ & $7$ & $8$\\
\midrule
degree $9-r$ & $6$ & $5$ & $4$ & $3$ & $2$ & $1$\\
$E_r$ & $A_2\times A_1$ & $A_4$ & $D_5$ & $E_6$ & $E_7$ & $E_8$\\
$\abs{W(E_r)}$ & $12$ & $120$ & $1920$ & $51840$ & $2903040$ & $696729600$\\
lines & $6$ & $10$ & $16$ & $27$ & $56$ & $240$\\
\bottomrule
\end{tabular}
\end{center}
For $r=6$ these are the $27$ lines on a smooth cubic surface (Cayley and Salmon, 1849), whose
configuration has symmetry group $W(E_6)$ of order $51840$. For $r=7$ the $56$ lines come in $28$ pairs that
map to the $28$ bitangents of a plane quartic. For $r=8$ the $240$ exceptional classes $E$ correspond to the
$240$ roots $E+K$ of $E_8$. For $r=4$ the $10$ lines meet in the pattern of the Petersen graph. At $r=9$,
$K\cdot K=0$ and the lattice is the affine $\tilde E_8$, and at $r=10$ we reach $E_{10}$ with $\det=-1$, as in
Section~\ref{DYN:sec:pqr}.

%======================================================================
\section{ADE elsewhere}
\label{DYN:sec:elsewhere}

The simply laced diagrams $A_n,D_n,E_6,E_7,E_8$ classify a surprising number of things. In each case
below we say what the diagram tells us there. Interlude~A (Section~\ref{LIE:sec:readoff}, item~7) already
mentioned McKay, du~Val and Gabriel; here we look more closely.

\subsection*{Finite subgroups of $\mathrm{SU}(2)$ and the icosahedron's $E_8$}
The finite subgroups of $\mathrm{SU}(2)$ are the cyclic groups and the binary dihedral, tetrahedral
($2T$, order $24$), octahedral ($2O$, $48$) and icosahedral ($2I$, $120$) groups, the preimages of the
rotation groups $G^+_{p,q,r}$ of Section~\ref{DYN:sec:pqr}. McKay (1980) attached to $\Gamma$ the graph with
a node for each irreducible representation $\rho_i$ and $a_{ij}$ edges when $\rho_j$ occurs $a_{ij}$ times in
$\C^2\otimes\rho_i$. The result is an extended diagram: $\tilde A_{n-1}$ for $\Z/n$, $\tilde D_{n+2}$ for the
binary dihedral group of order $4n$, and $\tilde E_6$, $\tilde E_7$, $\tilde E_8$ for $2T$, $2O$, $2I$. What
the diagram tells us here:
\begin{itemize}
\item \emph{The marks are the dimensions of the representations.} For $2I$ the nine marks of
  $\tilde E_8$ (Figure~\ref{DYN:fig:ext}) give
  \[
    \textstyle\sum m_i^2=1^2+2^2+3^2+4^2+5^2+6^2+4^2+2^2+3^2=120=\abs{2I}.
  \]
  (The relation
  $2m_i=\sum_{j\sim i}m_j$, which says that the marks are a null vector of the affine Cartan matrix, is the
  dimension count $\dim(\C^2\otimes\rho_i)=2\dim\rho_i$.)
\item \emph{The mark-$1$ nodes are the one-dimensional representations}, so their number $\det C$ is the
  order of the abelianization $\Gamma/[\Gamma,\Gamma]$. We checked by computer that $2T$, $2O$, $2I$ have abelianizations
  of order $3,2,1$, the determinants of $E_6,E_7,E_8$.
\item \emph{The number of nodes is the number of conjugacy classes}: $7,8,9$ for $2T,2O,2I$.
\end{itemize}

\subsection*{Simple singularities}
Klein (1884) showed that the ring of $\Gamma$-invariant polynomials on $\C^2$ is generated by three
invariants with one relation. For $2I$ the relation is
\[
  x^2+y^3+z^5=0,
\]
so $\C^2/2I$ is the surface in $\C^3$ cut out by this equation, singular at the origin. The exponents
$2,3,5$ are again the orders of the rotation axes of the icosahedron. Du Val (1934) resolved these
singularities. For $\C^2/\Gamma$, the minimal resolution replaces the singular point by a configuration
of smooth rational curves $C_1,\dots,C_r$, each with self-intersection $-2$, meeting transversally.
\begin{center}
\emph{The dual graph of the curves is the Dynkin diagram, and the intersection matrix is $-C$.}
\end{center}
For $x^2+y^3+z^5$ we get eight curves in the shape $T_{2,3,5}$, and the lattice they span in
$H_2$ of the resolution is $-E_8$. So the $E_8$ lattice of Lecture~6 \emph{grows out of the icosahedron}.
This also explains the coincidence of Section~\ref{DYN:sec:pqr}. The binary group $\langle p,q,r\rangle$
gives a star-shaped graph with three arms of $p-1$, $q-1$, $r-1$ curves, that is, $T_{p,q,r}$. Algebraic geometry
requires the intersection matrix of exceptional curves to be negative definite (Mumford, Grauert), and
by Proposition~\ref{DYN:prop:detT} this is the condition $\frac1p+\frac1q+\frac1r>1$.

The boundary of a small neighbourhood of the singular point is $\Sph^3/\Gamma$. Its first homology is
$\Gamma/[\Gamma,\Gamma]$, and for the plumbed neighbourhood it is also the cokernel of the intersection
matrix, of order $\abs{\det C}$. (Use the exact sequence $H_2(X)\to H_2(X,\partial X)\to H_1(\partial X)\to0$,
in which the first map is the intersection matrix.) For $E_8$, $\det C=1$, and $\Sph^3/2I$ is a homology
sphere: the \emph{Poincar\'e dodecahedral space}, obtained from a solid dodecahedron by gluing opposite
faces with a twist of $\frac\pi5$. The $E_8$ lattice is its ``filling''.

\paragraph{Functions instead of surfaces.} Arnold (1972) classified the \emph{simple} critical points of
functions: those near which a small deformation produces only finitely many inequivalent singularities
(no continuous moduli). They are
\[
  A_k:\ x^{k+1},\quad D_k:\ x^2y+y^{k-1}\ (k\ge4),\quad E_6:\ x^3+y^4,\quad E_7:\ x^3+xy^3,\quad E_8:\ x^3+y^5,
\]
plus a nondegenerate quadratic form in the remaining variables. The subscript is the Milnor number, the
number of nondegenerate critical points into which the singularity splits under a small perturbation. The
lattice of vanishing cycles (for surfaces) is the negative root lattice, and the monodromy group is the
Weyl group. Lecture~9's Morse lemma (Lemma~\ref{L9:lem:morse}) is the case $A_1$. The birth or death of
a pair of critical points in a one-parameter family of functions is $A_2$: $x^3+tx$. Thom's seven elementary catastrophes are the members of Arnold's list of codimension $1$ to
$4$: $A_2$ (fold), $A_3$ (cusp), $A_4$ (swallowtail), $A_5$ (butterfly), $D_4^\pm$ (hyperbolic and elliptic
umbilic) and $D_5$ (parabolic umbilic).

\subsection*{Representations of quivers and cluster algebras}
A \emph{quiver} is a directed graph, and a representation assigns a vector space to each node and a
linear map to each arrow. Gabriel (1972) proved that a connected quiver has only finitely many
indecomposable representations (up to isomorphism) if and only if its underlying graph is $A_n$, $D_n$,
$E_6$, $E_7$ or $E_8$. The dimension vectors of the indecomposables are exactly the positive roots,
written in the basis of simple roots; Bernstein, Gelfand and Ponomarev (1973) explained this with
reflection functors. So the $E_8$ quiver has $120$ indecomposables. Here the diagram's
\emph{positive roots} become objects.

Fomin and Zelevinsky's cluster algebras (2003) are of finite type exactly for the finite root
systems. The cluster variables correspond to the ``almost positive roots'' (positive roots and negative
simple roots), so type $A_n$ has $\frac{n(n+1)}2+n=\frac{n(n+3)}2$ of them, the diagonals of an
$(n+3)$-gon. Its clusters correspond to the triangulations, counted by a Catalan number. The
\emph{Coxeter number} reappears in Zamolodchikov's periodicity conjecture, proved by Keller (gallery).

\subsection*{Conformal field theory and subfactors}
Cappelli, Itzykson and Zuber (1987) classified the modular invariant partition functions of the
$\mathfrak{sl}_2$ conformal field theories at level $k$. They are labelled by the ADE diagrams with
\emph{Coxeter number} $h=k+2$, and the diagonal terms are given by the \emph{exponents}. So $E_6$, $E_7$,
$E_8$ ($h=12,18,30$) occur at levels $10$, $16$, $28$. In Jones' theory of subfactors (1983), an index less than $4$ must be
of the form $4\cos^2\frac\pi h$. The principal graph of such a subfactor is then an ADE diagram of
Coxeter number $h$, and by the classification of Ocneanu, Izumi, Kawahigashi and others (gallery) the
graphs that occur are exactly $A_n$, $D_{2n}$, $E_6$ and $E_8$, not $D_{\mathrm{odd}}$ or $E_7$. That
restriction is not visible on the diagram alone; it needs the fusion rules. The number $4\cos^2\frac\pi h$ is the largest eigenvalue of the
adjacency matrix of the diagram, squared. It is also Lecture~4's $4\cos^2\theta$, whose possible values
$0,1,2,3$ are the values for $h=2,3,4,6$.

\subsection*{Niemeier lattices: the Coxeter number is a design constant}
Interlude~B (Section~\ref{NT:sec:theta}) used modular forms to count vectors of even unimodular lattices,
and mentioned Niemeier's $24$ lattices in dimension $24$. Their root systems obey a remarkable rule:
\emph{all components have the same Coxeter number $h$, and there are $24h$ roots.} For example
$A_1^{24}$ ($h=2$, $48$ roots), $D_4^6$ ($h=6$, $144$ roots), $E_8^3$ and $D_{16}E_8$ (both $h=30$, $720$ roots).
Here is a proof from two facts already in these notes, with one input from the theory of modular forms.

\begin{proposition}[Venkov]\label{DYN:prop:venkov}
Let $L\subset\R^{24}$ be an even unimodular lattice with root system $\Delta=\Delta_1\cup\dots\cup\Delta_k$.
Then either $\Delta=\emptyset$, or $\Delta$ spans $\R^{24}$ and every $\Delta_j$ has Coxeter number
$h=\abs\Delta/24$.
\end{proposition}

\begin{proof}
\emph{Input from modular forms.} For a harmonic polynomial $P$ of degree $2$, the series
\[
  \Theta_{L,P}=\sum_{x\in L}P(x)\,\mathsf q^{x\cdot x/2}
\]
is a cusp form of weight $12+2=14$ for $\SL_2(\Z)$ (Hecke; see Ebeling, \emph{Lattices and Codes},
Ch.~3), and there are no nonzero cusp forms of weight $14$ (Serre, \emph{A Course in Arithmetic},
Ch.~VII). So the coefficient of $\mathsf q$ vanishes: $\sum_{\alpha\in\Delta}P(\alpha)=0$. Take $P(x)=\ip xy^2-\frac{1}{24}\abs x^2\abs y^2$, which is harmonic. Then
\[
  \sum_{\alpha\in\Delta}\ip\alpha y^2=\frac{\abs\Delta}{12}\abs y^2\qquad\text{for all }y\in\R^{24}.
\]
\emph{Input from Lecture~6.} For an irreducible root system of rank $r$ with roots of norm $2$, eutaxy
(Section~\ref{L04b:sec:rootdense}, item~2) gives $\sum_{\alpha\in\Delta_j}\ip\alpha y^2=\frac{2\abs{\Delta_j}}{r}\abs{y_j}^2$,
where $y_j$ is the projection of $y$ to the span of $\Delta_j$. By Interlude~A (item~3),
$\abs{\Delta_j}=rh_j$, so the constant is $2h_j$. If some $y\ne0$ is orthogonal to all roots, the left
side vanishes and $\abs\Delta=0$. Otherwise take $y$ in the span of $\Delta_j$: $2h_j=\abs\Delta/12$.
\end{proof}

So the Coxeter number, read off from a diagram, is the constant in the $2$-design property of its roots,
and a modular form forces all the constants to agree. Checking the list: $D_{16}E_8$ is allowed because
$h(D_{16})=2\cdot16-2=30=h(E_8)$, and $D_{24}$ has $h=46$ with $1104=24\cdot46$ roots.

\subsection*{Regular polytopes}
Finally the oldest appearance. The symmetry groups of the regular polytopes are Coxeter groups with
linear diagrams, and the labels are the Schl\"afli symbol:
\begin{center}
\renewcommand{\arraystretch}{1.15}
\begin{tabular}{llll}
\toprule
diagram & $\abs W$ & polytopes & what the diagram shows\\
\midrule
$A_3=\{3,3\}$ & $24$ & tetrahedron & symmetric diagram: self-dual\\
$B_3=\{4,3\}$ & $48$ & cube, octahedron & reversing the diagram: duality\\
$H_3=\{5,3\}$ & $120$ & dodecahedron, icosahedron & label $5$: no lattice\\
$A_4=\{3,3,3\}$ & $120$ & $5$-cell & self-dual\\
$B_4=\{4,3,3\}$ & $384$ & tesseract, $16$-cell & $16$-cell $=$ half-tesseract ($D_4$)\\
$F_4=\{3,4,3\}$ & $1152$ & $24$-cell & self-dual; Voronoi cell of $D_4$\\
$H_4=\{5,3,3\}$ & $14400$ & $120$-cell, $600$-cell & $\det G$ barely positive (Section~\ref{DYN:sec:icosa})\\
\bottomrule
\end{tabular}
\end{center}
In dimensions $n\ge5$ only $A_n$ and $B_n$ remain, giving the simplex, the cube and the cross-polytope.
The exceptional diagrams $F_4$ and $H_4$ are why dimension $4$ has six regular polytopes.

%======================================================================
\section{Summary: what the diagram knows}
\label{DYN:sec:summary}

\begin{center}
\renewcommand{\arraystretch}{1.25}
\small
\begin{tabular}{p{0.30\textwidth}p{0.62\textwidth}}
\toprule
feature of the diagram & what it tells us, and where\\
\midrule
$\det C$; $\det C(T_{p,q,r})=pqr(\frac1p+\frac1q+\frac1r-1)$ & density of the root lattice (A, Lecture~6); Platonic solids and triangle groups (\S\ref{DYN:sec:pqr}); $\abs{L^*/L}$, centre of the simply connected group, $\abs{\Gamma^{\mathrm{ab}}}$ (\S\ref{DYN:sec:elsewhere}); degree $9-r$ of a del Pezzo surface (\S\ref{DYN:sec:accidents})\\
sign of the Schl\"afli determinant & spherical, Euclidean or hyperbolic; the icosahedron is not made of regular tetrahedra (\S\ref{DYN:sec:icosa})\\
labels $m_{ij}\in\{2,3,4,6\}$ & crystallographic restriction; a $5$ forces a quasicrystal (\S\ref{DYN:sec:icosa}, C)\\
extended diagram, leaves & holes and their Delaunay cells (\S\ref{DYN:sec:holes})\\
marks $m_i$ & alcove vertices $\omega_i/m_i$; index of maximal-rank root sublattices; dimensions of the representations of $\Gamma\subset\mathrm{SU}(2)$\\
nodes of mark $1$ & classes of $L^*/L$ (FCC: lattice, octahedral and two tetrahedral classes); one-dimensional representations (McKay)\\
symmetries of the diagram & lattice automorphisms beyond $W$ (inversion of FCC, triality of $D_4$); folding (A, C)\\
symmetries of the extended diagram & $(L^*/L)\rtimes\operatorname{Aut}$: translations by holes, as in BCC\\
Coxeter number $h$ & $\abs\Delta=rh$; eutaxy constant $2h$; Niemeier lattices; CIZ level $k=h-2$; subfactor index $4\cos^2\frac\pi h$; periodicity $2(h+h')$\\
exponents & $\abs W=\prod(m_j+1)$ (A); diagonal of the CIZ invariants\\
positive roots & indecomposable quiver representations; cluster variables\\
small coincidences & accidental isomorphisms: $A_3=D_3$ is FCC $=$ checkerboard; $E_5=D_5$, $E_4=A_4$\\
\bottomrule
\end{tabular}
\end{center}

%======================================================================
\section{Gallery of theorems}
\label{DYN:sec:gallery}

\begin{itemize}[leftmargin=1.2em,itemsep=0.45em]
\item \textbf{Finite reflection groups} (H.~S.~M.~Coxeter, \emph{Discrete groups generated by reflections},
  Ann.\ of Math.~35 (1934), 588--621). Theorem~\ref{DYN:thm:coxeter}: the finite groups generated by
  reflections are classified by the positive definite Coxeter diagrams, and the Euclidean ones by the
  extended diagrams. The non-crystallographic $H_3$, $H_4$ are the ones with a $5$.
\item \textbf{Schl\"afli's regular polytopes} (L.~Schl\"afli, \emph{Theorie der vielfachen Kontinuit\"at},
  written 1850--52, published 1901). There are six regular polytopes in $\R^4$ and three in each
  $\R^n$, $n\ge5$. The criterion is the positivity of the determinant in
  Proposition~\ref{DYN:prop:gram335}. Coxeter's \emph{Regular Polytopes} (1948) is the standard account.
\item \textbf{Maximal subgroups of maximal rank} (A.~Borel and J.~de~Siebenthal, Comment.\ Math.\ Helv.~23
  (1949), 200--221). The maximal closed connected subgroups of maximal rank of a compact simple Lie group
  are obtained by deleting a node of the extended diagram whose mark is prime (or a node of mark $1$ from the
  ordinary diagram, giving a centralizer of a circle). Its lattice shadow is Theorem~\ref{DYN:thm:holes}(a):
  $D_8$ and $A_8$ in $E_8$, of index $2$ and $3$.
\item \textbf{Voronoi and Delaunay cells of root lattices} (J.~H.~Conway and N.~J.~A.~Sloane, \emph{Sphere
  Packings, Lattices and Groups}, Ch.~21 (1988), and \emph{The cell structures of certain lattices}, in
  \emph{Miscellanea Mathematica}, Springer 1991; R.~V.~Moody and J.~Patera, J.\ Phys.~A~25 (1992),
  5089--5134). Conway and Sloane determine the Voronoi and Delaunay cells of all root lattices and their
  duals (the relevant vectors of a root lattice are its roots). Moody and Patera classify the faces of
  these cells by decorated Coxeter--Dynkin diagrams, a systematic form of Theorem~\ref{DYN:thm:holes}.
\item \textbf{Thinnest lattice coverings} (R.~P.~Bambah, \emph{On lattice coverings by spheres}, Proc.\
  Nat.\ Inst.\ Sci.\ India~20 (1954), 25--52, for $n=3$; Delone--Ryshkov 1963 for $n=4$;
  Ryshkov--Baranovskii 1975 for $n=5$). The thinnest lattice covering of $\R^n$ for $n\le5$ is $A_n^*$. For
  $n=3$ it is BCC, the union of FCC with all of its holes. For $n\ge6$ the lattice covering problem is
  open.
\item \textbf{McKay correspondence} (J.~McKay, \emph{Graphs, singularities, and finite groups}, Proc.\
  Sympos.\ Pure Math.~37 (1980), 183--186; geometric version: G.~Gonzalez-Sprinberg and J.-L.~Verdier,
  Ann.\ Sci.\ \'Ec.\ Norm.\ Sup.\ (4)~16 (1983), 409--449). The McKay graph of a finite
  $\Gamma\subset\mathrm{SU}(2)$ is an extended ADE diagram. The geometric version matches the
  nontrivial irreducible representations with the exceptional curves of the minimal resolution of
  $\C^2/\Gamma$.
\item \textbf{Simple surface singularities} (F.~Klein, \emph{Vorlesungen \"uber das Ikosaeder}, 1884;
  P.~Du~Val, Proc.\ Cambridge Philos.\ Soc.~30 (1934)). The quotients $\C^2/\Gamma$ are hypersurfaces
  $x^2+y^3+z^5=0$ etc., and the dual graphs of their minimal resolutions are the ADE diagrams, with
  intersection matrix $-C$.
\item \textbf{The $E_8$ manifold} (M.~Freedman, J.\ Differential Geom.~17 (1982), 357--453;
  S.~Donaldson, J.\ Differential Geom.~18 (1983), 279--315). There is a closed simply connected
  topological $4$-manifold with intersection form $E_8$: cap the $E_8$ plumbing, whose boundary is the
  Poincar\'e sphere $\Sph^3/2I$, with a contractible topological manifold. It has no smooth structure
  (Rokhlin: the signature of a smooth spin $4$-manifold is divisible by $16$). Donaldson: the intersection form of a smooth closed
  simply connected $4$-manifold that is definite must be diagonal.
\item \textbf{Simple critical points} (V.~I.~Arnold, Funct.\ Anal.\ Appl.~6 (1972), 254--272). The simple
  germs of functions are exactly $A_k$, $D_k$, $E_6$, $E_7$, $E_8$, with Weyl groups as monodromy groups.
  Morse's nondegenerate critical point of Lecture~9 is $A_1$.
\item \textbf{Gabriel's theorem} (P.~Gabriel, Manuscripta Math.~6 (1972), 71--103; I.~N.~Bernstein,
  I.~M.~Gelfand, V.~A.~Ponomarev, Russian Math.\ Surveys~28 (1973)). A connected quiver has finite
  representation type if and only if its graph is ADE; indecomposables correspond to positive roots.
\item \textbf{Finite type cluster algebras} (S.~Fomin and A.~Zelevinsky, \emph{Cluster algebras II: Finite
  type classification}, Invent.\ Math.~154 (2003), 63--121). Cluster algebras of finite type are classified
  by the Cartan matrices of finite type; the cluster complex is the dual of a generalized associahedron.
\item \textbf{Periodicity of $Y$-systems} (B.~Keller, \emph{The periodicity conjecture for pairs of Dynkin
  diagrams}, Ann.\ of Math.~177 (2013), 111--170; the case $(\Delta,A_1)$: Fomin--Zelevinsky, Ann.\ of
  Math.~158 (2003), 977--1018). For a pair $(\Delta,\Delta')$ of simply laced Dynkin diagrams the
  $Y$-system
  $Y_{i,i',t-1}Y_{i,i',t+1}=\prod_{j\sim i}(1+Y_{j,i',t})\big/\prod_{j'\sim i'}(1+Y_{i,j',t}^{-1})$ is periodic,
  with period dividing $2(h+h')$, twice the sum of the Coxeter numbers. Conjectured by Zamolodchikov
  (for $\Delta'=A_1$) in the thermodynamic Bethe ansatz.
\item \textbf{ADE modular invariants} (A.~Cappelli, C.~Itzykson, J.-B.~Zuber, Comm.\ Math.\ Phys.~113
  (1987), 1--26). The modular invariant partition functions of $\widehat{\mathfrak{sl}}_2$ at level $k$ are
  labelled by ADE diagrams of Coxeter number $k+2$ (with $D$ only for even $k$), with diagonal terms given
  by the exponents.
\item \textbf{Subfactors of small index} (V.~F.~R.~Jones, Invent.\ Math.~72 (1983), 1--25; classification
  announced by A.~Ocneanu (1988); published proofs by M.~Izumi, Publ.\ RIMS~27 (1991), 953--994
  (no $D_{\mathrm{odd}}$, $E_7$) and Y.~Kawahigashi, J.\ Funct.\ Anal.~127 (1995), 63--107 (flatness, and
  the $D_{2n}$ subfactors), with the $E_6$, $E_8$ cases completed by several authors; survey: Jones, Morrison, Snyder, Bull.\ Amer.\
  Math.\ Soc.~51 (2014), 277--327). The index of a subfactor lies in
  $\{4\cos^2\frac\pi n:n\ge3\}\cup[4,\infty]$, and below $4$ the principal graphs are $A_n$, $D_{2n}$,
  $E_6$, $E_8$.
\item \textbf{Niemeier lattices} (H.-V.~Niemeier, J.\ Number Theory~5 (1973), 142--178; B.~B.~Venkov,
  Trudy Mat.\ Inst.\ Steklov~148 (1978), 65--76). Proposition~\ref{DYN:prop:venkov}, and the
  classification: the $24$ even unimodular lattices in $\R^{24}$ are determined by their root systems,
  which are empty (Leech) or of rank $24$ with all components of the same Coxeter number.
\item \textbf{Lines on del Pezzo surfaces} (A.~Cayley and G.~Salmon, 1849; C.~Jordan, \emph{Trait\'e des
  substitutions}, 1870; P.~Du~Val, Proc.\ London Math.\ Soc.~35 (1933); Yu.~I.~Manin, \emph{Cubic Forms},
  1972). Proposition~\ref{DYN:prop:delpezzo}: the lines on a del Pezzo surface of degree $9-r$ are permuted
  by $W(E_r)$, and the $27$ lines on a cubic surface have symmetry group $W(E_6)$ of order $51840$.
\end{itemize}

%======================================================================
\section*{Exercises}
\addcontentsline{toc}{section}{Exercises}

\begin{exercise}
Compute $\det C(T_{1,q,r})$ directly and check that it is $q+r$, the determinant of $A_{q+r-1}$. Show that
$E_{10}=T_{2,3,7}$ has determinant $-1$, and that its Cartan matrix has exactly one negative eigenvalue.
\end{exercise}

\begin{exercise}
Using Proposition~\ref{DYN:prop:gram335}, compute the Schl\"afli determinants of $\{4,3,3\}$, $\{3,4,3\}$,
$\{4,3,4\}$ and $\{5,3,5\}$. Which are spherical, which Euclidean, which hyperbolic? Interpret
$\{4,3,4\}$ geometrically.
\end{exercise}

\begin{exercise}
Count the vertices, edges, $2$-faces and cells of the $24$-cell from its diagram $F_4=\{3,4,3\}$, using
$\abs{W(F_4)}=1152$ and the rule that the stabilizer of a $k$-face is generated by all nodes except the
$(k+1)$-st. Check Euler's relation $V-E+F-C=0$.
\end{exercise}

\begin{exercise}
For the FCC lattice $D_3$, verify that each lattice point lies on $8$ tetrahedra and $6$ octahedra of the
Delaunay tiling, and that each edge lies on $2$ of each. Compute the solid angles at the vertices of the
tetrahedron and the octahedron, and confirm $8\Omega_T+6\Omega_O=4\pi$.
\end{exercise}

\begin{exercise}
Show that the hole $\omega_i$ of $A_n$ has $\abs{\omega_i}^2=i(n+1-i)/(n+1)$ and that its nearest lattice
points form a hypersimplex with $\binom{n+1}i$ vertices. For which $n$ are all Delaunay cells of $A_n$
simplices?
\end{exercise}

\begin{exercise}
In $D_n$, compute the distances from $e_1$ and from $h=\frac12(1,\dots,1)$ to the lattice, and the number of
nearest lattice points. At which $n$ does the deepest hole change type? Relate this to the construction
$E_8=D_8\cup(D_8+h)$.
\end{exercise}

\begin{exercise}
Let $g=\omega_3\in A_8^*$. Show that $A_8\cup(A_8+g)\cup(A_8+2g)$ is an even unimodular lattice with $240$
vectors of norm $2$: $72$ from $A_8$ and $84$ from each nontrivial coset. Conclude that it is $E_8$.
\end{exercise}

\begin{exercise}
Project the $24$ roots of $D_4$ onto the plane fixed by triality, and show that the images form the root
system $G_2$, with $6$ long roots of norm $2$ and $6$ short roots of norm $\frac23$.
\end{exercise}

\begin{exercise}
On the cubic surface $X_6$, list the classes $e_i$, $h-e_i-e_j$, $2h-\sum_{k\ne i}e_k$, and check that each
satisfies $E\cdot E=E\cdot K=-1$. How many are there? Show that the simple roots of
Proposition~\ref{DYN:prop:delpezzo} span $K^\perp$.
\end{exercise}

\begin{exercise}
For the extended diagrams $\tilde E_6$ and $\tilde E_7$, find the marks, check that $\sum m_i^2$ equals
$24$ and $48$, and that the number of nodes of mark $1$ is $\det C$. Which nodes give holes of the
lattices $E_6$ and $E_7$, and what are their Delaunay cells?
\end{exercise}

\section*{Further reading}
\addcontentsline{toc}{section}{Further reading}
\begin{itemize}[nosep]
\item M.~Hazewinkel, W.~Hesselink, D.~Siersma and F.~D.~Veldkamp, \emph{The ubiquity of Coxeter--Dynkin
  diagrams (an introduction to the A--D--E problem)}, Nieuw Arch.\ Wisk.~25 (1977), 257--307. The
  classic survey of where ADE appears.
\item H.~S.~M.~Coxeter, \emph{Regular Polytopes}, 3rd ed., Dover, 1973. Schl\"afli's criterion, Wythoff's
  construction, and the $600$-cell.
\item J.~E.~Humphreys, \emph{Reflection Groups and Coxeter Groups}, Cambridge, 1990. Coxeter's theorem,
  affine Weyl groups and alcoves.
\item J.~H.~Conway and N.~J.~A.~Sloane, \emph{Sphere Packings, Lattices and Groups}, Ch.~4 and 21. Glue,
  holes and Voronoi cells of root lattices.
\item P.~Slodowy, \emph{Simple Singularities and Simple Algebraic Groups}, Lecture Notes in Math.~815,
  Springer, 1980. The bridge between $\C^2/\Gamma$ and Lie theory.
\item J.~C.~Baez, \emph{The octonions}, Bull.\ Amer.\ Math.\ Soc.~39 (2002), 145--205. Triality, $G_2$ and
  the exceptional diagrams.
\item I.~V.~Dolgachev, \emph{Classical Algebraic Geometry: A Modern View}, Cambridge, 2012, Ch.~8--9.
  Del Pezzo surfaces and the root systems $E_r$.
\end{itemize}
% ===== end chapters/DYN.tex =====
\endgroup

\end{document}
